% ############################################### PREAMBLE AND CONTROL PANEL #

% ########################################################### DOCUMENT CLASS #
\documentclass[  %   draft
             % , demo  % black boxes instead of figures
              , 8pt, twoside, openright]{extbook}

% ########################################################## AND CONTROL PANEL #

% true - hmarginration 1:1, urlcolor - darkgreen, version - for searching
\newif\ifscreen 	   \screentrue
% true - no apendix, no author of grammer, no backmatter ISBN, 
\newif\ifPDF               \PDFtrue
% true - rotate pages of morph verb tables
\newif\ifPDFrotate         \PDFrotatetrue
% true - adds white searchable value !
\newif\ifPDFsearchvalue    \PDFsearchvaluetrue
\newif\ifmakecovers        \makecoverstrue
\newif\ifmakefrontmatter   \makefrontmattertrue
\newif\ifinputintroduction \inputintroductiontrue
\newif\ifinputletters      \inputletterstrue
\newif\ifinputphon         \inputphontrue
\newif\ifinputmorpho       \inputmorphotrue
\newif\ifmakebackmatter    \makebackmattertrue
\newif\ifshowcrop           \showcroptrue
% corrections
\newif\ifC                  \Ctrue

\newif\ifPRINTHOUSE 	   \PRINTHOUSEfalse
\newif\ifPDFSEARCH  	   \PDFSEARCHfalse
\newif\ifPDFPRINT 	   \PDFPRINTfalse
\newif\ifPDFtest	   \PDFtesttrue

%versions
%PRINTHOUSE
\ifPRINTHOUSE
\makecoversfalse
\screenfalse
\PDFfalse
\PDFrotatefalse
\PDFsearchvaluefalse
\showcroptrue
\Ctrue
\fi

%PDFSEARCH
\ifPDFSEARCH
\screentrue
\PDFtrue
\PDFrotatetrue
\PDFsearchvaluetrue
\showcropfalse
\Ctrue
\fi

%PDFPRINT
\ifPDFPRINT
\screenfalse
\PDFtrue
\PDFrotatefalse
\PDFsearchvaluefalse
\showcropfalse
\Ctrue
\fi

%PDFtest
\ifPDFtest
\makecoversfalse
\makefrontmatterfalse
\inputintroductionfalse
\inputletterstrue
\inputphonfalse
\inputmorphofalse
\makebackmatterfalse
\screentrue
\PDFtrue
\PDFrotatefalse
\PDFsearchvaluefalse
\showcropfalse
\Cfalse
\fi

% 1. PDF (for printing)  
% a) searchable value, b) no apendix, c) version= PDF d) no ISBN and no ISBN bar
% e) automatic rotation of landscape pages, f) no author of Icelandic Grammer
% 2. PDF (for viewing offline and searching) + screen (horizontal ratio 1:1) 
% 3. PRINTHOUSE
% a) no searchable value, b) apendix, c) version ISBN d) ISBN bar
% e) no rotation of landscape pages, f)urllinkcolor black
% f) A4 paper dimentions, C5 layout, shows cropping marks

% ########################################################## PAGE GEOMETRIES #

\usepackage%[showframe]
           {geometry}
           

\ifshowcrop
% Basic geometry of document
\geometry
  { headsep    =   \baselineskip
  , textwidth  = 42\baselineskip
  , textheight = 60\baselineskip
  , hmarginratio = \ifscreen 1:1\else 2:3\fi
  , vmarginratio = 2:3
  , bindingoffset = 0cm
  , onecolumn
  , a4paper 
  , layout=c5paper,
  , layouthoffset=\dimexpr(\paperwidth-\csname Gm@layoutwidth\endcsname)/2\relax
   , layoutvoffset=\dimexpr(\paperheight-\csname Gm@layoutheight\endcsname)/2\relax
  , showcrop
  }
\else
\geometry
  { headsep    =   \baselineskip
  , textwidth  = 42\baselineskip
  , textheight = 60\baselineskip
  , hmarginratio = \ifscreen 1:1\else 2:3\fi
  , vmarginratio = 2:3
  , bindingoffset = 0cm
  , onecolumn,
  , c5paper
  }
\fi

\ifshowcrop
% Geometry of the dictionary section
\newcommand\dictionarygeometry{\newgeometry
  { textwidth  = 42\baselineskip
  , textheight = 60\baselineskip
  , hmarginratio = \ifscreen 1:1\else 2:3\fi
  , vmarginratio = 3:2
  , bindingoffset = 0cm
  , twocolumn
  , a4paper 
  , layout=c5paper
  , layouthoffset=\dimexpr(\paperwidth-\csname Gm@layoutwidth\endcsname)/2\relax
   , layoutvoffset=\dimexpr(\paperheight-\csname Gm@layoutheight\endcsname)/2\relax
  , showcrop
  }}
   % Geometry of the cover pages
\newcommand\covergeometry{\newgeometry
  { textwidth  = 42\baselineskip
  , textheight = 60\baselineskip
  , hmarginratio = 1:1
  , vmarginratio = 2:3
  , bindingoffset = 0cm
  , onecolumn
  , layout=c5paper
 , layouthoffset=\dimexpr(\paperwidth-\csname Gm@layoutwidth\endcsname)/2\relax
   , layoutvoffset=\dimexpr(\paperheight-\csname Gm@layoutheight\endcsname)/2\relax
  , showcrop
  }} 
\else
% Geometry of the dictionary section
\newcommand\dictionarygeometry{\newgeometry
  { textwidth  = 42\baselineskip
  , textheight = 60\baselineskip
  , hmarginratio = \ifscreen 1:1\else 2:3\fi
  , vmarginratio = 3:2
  , bindingoffset = 0cm
  , twocolumn,
  , c5paper
  }}
   % Geometry of the cover pages
\newcommand\covergeometry{\newgeometry
  { textwidth  = 42\baselineskip
  , textheight = 60\baselineskip
  , hmarginratio = 1:1
  , vmarginratio = 2:3
  , bindingoffset = 0cm
  , onecolumn
  , c5paper
  }} 
\fi

% ################################################################# LANGUAGES #

\usepackage[icelandic, latin, czech]{babel}
\usepackage{csquotes}


% ################################################################## ENCODING #

\usepackage[utf8]{inputenc}

% Smashing uppercase letters with diacritics
% ------------------ THIS IS EVIL AND SHOULD NEVER BE DONE -------------------
% Don't smash anything if we are on the first run to produce the index
\IfFileExists{figures.idx}{
  % Czech uppercase letters with diacritics:
  \DeclareUnicodeCharacter{00C1}{\strut\smash{\' A}} % Á
  \DeclareUnicodeCharacter{010C}{\strut\smash{\v C}} % Č
  \DeclareUnicodeCharacter{010E}{\strut\smash{\v D}} % Ď
  \DeclareUnicodeCharacter{00C9}{\strut\smash{\' E}} % É
  \DeclareUnicodeCharacter{011A}{\strut\smash{\v E}} % Ě
  \DeclareUnicodeCharacter{00CD}{\strut\smash{\' I}} % Í
  \DeclareUnicodeCharacter{0147}{\strut\smash{\v N}} % Ň
  \DeclareUnicodeCharacter{00D3}{\strut\smash{\' O}} % Ó
  \DeclareUnicodeCharacter{0158}{\strut\smash{\v R}} % Ř
  \DeclareUnicodeCharacter{0160}{\strut\smash{\v S}} % Š
  \DeclareUnicodeCharacter{0164}{\strut\smash{\v T}} % Ť
  \DeclareUnicodeCharacter{00DA}{\strut\smash{\' U}} % Ú
  \DeclareUnicodeCharacter{016E}{\strut\smash{\r U}} % Ů
  \DeclareUnicodeCharacter{00DD}{\strut\smash{\' Y}} % Ý
  \DeclareUnicodeCharacter{017D}{\strut\smash{\v Z}} % Ž
  % Icelandic uppercase letters with diacritics: ÁÉÍÓÚÝ (ÐÞÆ) and
  \DeclareUnicodeCharacter{00D6}{\strut\smash{\" O}} % Ö
}{}
% ------------------ THAT WAS EVIL AND SHOULD NEVER BE DONE ------------------

% Non-breakable space
\DeclareUnicodeCharacter{00A0}{~}

% Non-breakable hyphen
\newcommand*{\nobreakhyphen}{\mbox{-}}
\DeclareUnicodeCharacter{2011}{\nobreakhyphen}


% ############################################################ FONTS, SYMBOLS #

\usepackage[T1]{fontenc}

\usepackage{tgpagella}
\usepackage[scaled=0.88]{helvet}      % relative scale of the two fonts

\def\phvfamily{\fontfamily{phv}\selectfont}

%\usepackage{stmaryrd}
% Let's keep this minimal instead of using the package:
\DeclareSymbolFont{stmry}{U}{stmry}{m}{n}
%\SetSymbolFont{stmry}{bold}{U}{stmry}{b}{n}
\DeclareMathSymbol\shortrightarrow\mathrel{stmry}{"01}
\DeclareMathSymbol\shortuparrow\mathrel{stmry}{"02}

% Methinks these are not necessary
%\usepackage{dingbat}
%\usepackage{manfnt}
%\usepackage{latexsym}

\usepackage{pifont}




% #################################################################### LAYOUT #

% \normalsize should be {8pt}{9.6pt}
\def\HUGE {\fontsize{23.887872pt}{3\baselineskip}\selectfont}
\def\Huge {\fontsize{19.906560pt}{3\baselineskip}\selectfont}
\def\huge {\fontsize{16.588800pt}{3\baselineskip}\selectfont}
\def\LARGE{\fontsize{13.824000pt}{2\baselineskip}\selectfont}
\def\Large{\fontsize{11.520000pt}{2\baselineskip}\selectfont}
\def\large{\fontsize{ 9.600000pt}{2\baselineskip}\selectfont}

% Two columns layout ruler
\setlength\columnsep    {2\baselineskip}
\setlength\columnseprule{0.4pt}

% Necessary for baseline alignment
\topskip=\baselineskip
\raggedbottom
\setlength\parskip{0pt} % it's better to avoid glue

% Temporarily suppress warnings
\hbadness=5000
\vbadness=5000

\setlength\emergencystretch{17pt}

% Allow smaller emergencystretch in several cases
\newenvironment{xtolerant}[2]{%
  \par
  \ifx\empty#1\empty\else\tolerance=#1\relax\fi
  \ifx\empty#2\empty\else\emergencystretch=#2\relax\fi
}{%
  \par
}

% ######################################################### ASSORTED PACKAGES #

\usepackage{tikz}
\usetikzlibrary{calc}

\usepackage{tipa}                      % IPA phonetics

\usepackage{fourier-orns}              % ornaments
\usepackage{amsmath}                   % non-breakable dash
\usepackage{hyphenat}                  % no hyphen in abbreviations

\usepackage{hanging}
\usepackage{fix2col}
\usepackage{fixltx2e}                  % subscript
\usepackage{pagecolor}
\usepackage{afterpage}                 % Used in coloring trick for covers
\usepackage{pdfpages}

\usepackage{multirow}
\usepackage{multicol}
\usepackage{rotating}

\usepackage{etoolbox} % http://ctan.org/pkg/etoolbox
%\usepackage{newunicodechar} 		% new definition of middle dot
\usepackage[ISBN=978-80-260-2385-2,SC0]{ean13isbn}

\usepackage{pgffor}

\usepackage{newunicodechar}

\newunicodechar{·}{\chejnikcdot}
\makeatletter
\newcommand{\chejnikcdot}{%
  \ifnum\lastpenalty=10042
    \kern-0.06667em\relax % dots should always be kerned
  \fi
  \textperiodcentered
  \@ifnextchar|{\kern\chejnikkern\relax}{\penalty10042 }%
}

\newcommand{\closeupdotbar}{%
  \renewcommand\chejnikkern{-0.06667em}%
}
\newcommand\chejnikkern{0pt} % default

% Middle dot + | little bit closer, nicer
%\newunicodechar{·}{\advancedcdot}
%\makeatletter
%\newcommand*\advancedcdot{\textperiodcentered
%  \@ifnextchar|{\hspace{-0.03667em}}{\checkifcdotnext}}
%\newcommand*\checkifcdotnext[2]
%  {\def\tmpa{·}\def\tmpb{#1#2}\ifx\tmpa\tmpb\hspace{-0.06667em}\fi#1#2}
%\makeatother

% ############################################################ BASELINE GRID #

\usepackage{eso-pic}
\usepackage{tikzpagenodes}

% Command to draw a baseline grid
\newcommand\drawbaselinegrid{%
  \begin{tikzpicture}[overlay,remember picture]
    \draw [red!30!white, ultra thin, dashed]
      (0,0) grid [ ystep  = \baselineskip, xstep = \textwidth
                 , shift  = (current page text area.north west)
                 , yshift = -\dp\strutbox
                 ] ++(\textwidth,-\textheight);
    \draw [red!30!white, thin]
      (0,0) grid [ step = \baselineskip, xstep = \textwidth
                 , shift=(current page text area.north)
                 ] ++(0.5\textwidth,-\textheight)
      (0,0) grid [ step = \baselineskip, xstep = \textwidth
                 , shift=(current page text area.north)
                 ] ++(-0.5\textwidth,-\textheight);
    \draw[black!10!white]
      (current page text area.north west)
        rectangle (current page text area.south east);
  \end{tikzpicture}}

% Draw a baseline grid on every page
% \AddToShipoutPicture{\drawbaselinegrid}


% ################################################################## METADATA #

\title{%
  \textbf{Islandsko-český studijní slovník}%
  \thanks{Tato kniha byla vytvořena v \LaTeX{}u pod Ubuntu 14.04.%
    Poděkování patří všem autorům, kteří publikují pod svobodnými licencemi.}}
\author{Aleš Chejn, Ján Zaťko, Jón Gíslason}
\date{říjen 2011}


% ############################################################## BIBLIOGRAPHY #

\usepackage[style=authoryear]{biblatex}
\addbibresource{slovnik.bib}


% ##################################################################### INDEX #

\usepackage{imakeidx}

% Index of authors of photographs
\makeindex
  [ name = figures
  , title = {Seznam autorů fotografií a ilustrací}
  , intoc
  , columnseprule
  , columnsep = \columnsep
  , noautomatic                       % We handle this in make.sh
  ]

% icelandic rules start
\usepackage{filecontents}
\def\mygroup#1{{\bfseries#1}}
\begin{filecontents*}{style.xdy}
;; style.xdy
(markup-letter-group :open-head "~n\mygroup{" :close-head "}")
\end{filecontents*}
% icelandic rules end


% ############################################################### PAGE STYLES #

\usepackage{fancyhdr}

% -------------------------------------- BASIC PAGE STYLE
\fancypagestyle{plain}{
  \fancyhf\relax                       % Clear header/footer
  \renewcommand \headrulewidth {0.0pt} % No header rule
  \renewcommand \footrulewidth {0.0pt} % No footer rule
  \fancyfoot[C]{\thepage}              % Page number in footer, centred
}

% -------------------------------------- INDEX PAGE STYLE  
%\fancypagestyle{indexstyle}{
%  \fancyhf\relax                       % Clear header/footer
%  \renewcommand \headrulewidth {0.4pt} % Thin header rule
%  \renewcommand \footrulewidth {0.0pt} % No footer rule
%  \fancyhead[C]{\thepage}              % Page number in footer, centred
%}

% -------------------------------------- DICTIONARY PAGE STYLE
\fancypagestyle{myheadings}{
  \fancyhf\relax                       % Clear header/footer
  \renewcommand \headrulewidth {0.4pt} % Thin header rule
  \fancyhead[CO,CE]{\thepage}          % Page number in header, centred
  % NOTE: the page numbers will be printed when the dictionary is ready
  \fancyhead[LE,LO]{\phvfamily\bfseries\rightmark}
  \fancyhead[RE,RO]{\phvfamily\bfseries\leftmark}
}

\pagestyle{plain}  

% ##################################################################### LISTS #

\usepackage{expdlist}
\usepackage{enumitem}

\setlist{nosep, topsep=\baselineskip}
\setlist[itemize]{leftmargin=*,label=\scalebox{.6}{\textbullet}}


% #################################################################### FLOATS #

\usepackage{caption}
\captionsetup{labelformat=empty}

% Alter some LaTeX defaults for better treatment of figures:
% See p.105 of "TeX Unbound" for suggested values. 
% See pp. 199-200 of Lamport's "LaTeX" book for details.
% ------------------------------------- Parameters for ALL pages:
\renewcommand \topfraction    {0.9}   %   max fraction of floats at top
\renewcommand \bottomfraction {0.8}   %   max fraction of floats at bottom
% ------------------------------------- Parameters for TEXT (non float) pages:
\setcounter{topnumber}    {1}
\setcounter{bottomnumber} {1}
\setcounter{totalnumber}  {1}         %   2 may work better
\setcounter{dbltopnumber} {2}         %   for 2-column pages
\renewcommand \dbltopfraction   {0.9} %   fit big float above 2-col. text
\renewcommand \textfraction     {0.07}%   allow minimal text w. figs
% ------------------------------------- Parameters for FLOAT (non text) pages:
\renewcommand \floatpagefraction{0.7} %   require fuller float pages
% N.B.: floatpagefraction MUST be less than topfraction !!
\renewcommand{\dblfloatpagefraction}{0.7} % require fuller float pages
% remember to use [htp] or [htpb] for placement

\makeatletter
\setlength{\@fptop}{0pt}
\setlength{\@fpbot}{0pt plus 1fil}
\makeatother

% This kills spacing around floats
\setlength \intextsep        {0pt}
\setlength \floatsep         {0pt}
\setlength \textfloatsep     {0pt}
\setlength \dbltextfloatsep  {0pt}
\setlength \dblfloatsep      {0pt}

% This kills spacing around captions
\setlength \abovecaptionskip {0pt}
\setlength \belowcaptionskip {0pt}


% #################################################################### TABLES #

\usepackage{booktabs}
\usepackage{tabularx}
\usepackage{longtable}
\usepackage{ltxtable}
\usepackage{lscape}

\ifPDFrotate
\else
\makeatletter			       
\renewcommand\PLS@Rotate[1]{}          % this command prevents pdflscape 
\makeatother                           % to rotate pages
\fi

%\LetLtxMacro\oldtabularx\tabularx
%\LetLtxMacro\oldendtabularx\endtabularx
%\renewenvironment{tabularx}[1]
%  {\oldtabularx{\columnwidth}{#2}}
%  {\oldendtabularx}

\newcommand\LTXfw[1]{\LTXtable{\columnwidth}{#1}}

%%% \LetLtxMacro\oldtabularx\tabularx
%%% \LetLtxMacro\oldendtabularx\endtabularx
%%% \newenvironment{longtabularx}[2]
%%%   {\oldtabularx{#1}{#2}}
%%%   {\oldendtabularx}

% ######################################################### czech babel vs cmidrule
% ###################################### http://tex.stackexchange.com/a/112001/8380
\begingroup
  \makeatletter
  \catcode`\-=\active
  \AtBeginDocument{
  \def\@@@cmidrule[#1-#2]#3#4{\global\@cmidla#1\relax
    \global\advance\@cmidla\m@ne
    \ifnum\@cmidla>0\global\let\@gtempa\@cmidrulea\else
    \global\let\@gtempa\@cmidruleb\fi
    \global\@cmidlb#2\relax
    \global\advance\@cmidlb-\@cmidla
    \global\@thisrulewidth=#3
    \@setrulekerning{#4}
    \ifnum\@lastruleclass=\z@\vskip \aboverulesep\fi
    \ifnum0=`{\fi}\@gtempa
    \noalign{\ifnum0=`}\fi\futurenonspacelet\@tempa\@xcmidrule}
  }
\endgroup

% macro that supress cmidrule at page brake
% http://tex.stackexchange.com/questions/110548/using-longtable-with-cmidrule-page-breaking-issues
\makeatletter
\def\LT@output{%
  \ifnum\outputpenalty <-\@Mi
    \ifnum\outputpenalty > -\LT@end@pen
      \LT@err{floats and marginpars not allowed in a longtable}\@ehc
    \else
      \setbox\z@\vbox{\unvbox\@cclv}%
      \ifdim \ht\LT@lastfoot>\ht\LT@foot
        \dimen@\pagegoal
        \advance\dimen@-\ht\LT@lastfoot
        \ifdim\dimen@<\ht\z@
          \setbox\@cclv\vbox{\unvbox\z@
\copy\LT@foot\vss}%
          \@makecol
          \@outputpage
          \setbox\z@\vbox{\box\LT@head}%
        \fi
      \fi
      \global\@colroom\@colht
      \global\vsize\@colht
      \vbox
        {\unvbox\z@
\box\ifvoid\LT@lastfoot\LT@foot\else\LT@lastfoot\fi}%
    \fi
  \else
    \setbox\@cclv\vbox{\unvbox\@cclv

\setbox0\lastbox
\ifdim\ht0=0.29999pt \unskip % hope that was a cline we threw away
\else
\nointerlineskip
\box0     % put it back, whatever it was
\fi

\copy\LT@foot\vss}%
    \@makecol
    \@outputpage
      \global\vsize\@colroom
    \copy\LT@head\nobreak
  \fi}

\makeatother

% These values ensure that for top, mid and bottom rules
% space above + rule width + space below = baselineskip
\setlength \abovetopsep    {0.43\baselineskip}
\setlength \heavyrulewidth {0.10\baselineskip}
\setlength \belowbottomsep {0.43\baselineskip}
\setlength \aboverulesep   {0.47\baselineskip}
\setlength \lightrulewidth {0.06\baselineskip}
\setlength \belowrulesep   {0.47\baselineskip}

%\usepackage{letltxmacro}

%%% \LetLtxMacro\oldtabular\tabular
%%% \LetLtxMacro\oldendtabular\endtabular
%%% \renewenvironment{tabular}[2][t]
%%%   {\oldtabular[#1]{#2}}
%%%   {\oldendtabular}

%%% \LetLtxMacro\oldtable\table
%%% \LetLtxMacro\oldendtable\endtable
%%% \renewenvironment{table}[1][h]
%%%   {\leavevmode\oldtable[#1]\vspace{\dp\strutbox}}
%%%   {\vspace{-\dp\strutbox}\oldendtable}
%%%   % This aligns a [t] aligned tab-envo to baseline grid


%%% % declention and conjugation tables
%%% \usepackage{floatrow}
%%% \DeclareFloatFont{footnotesize}{\footnotesize}
%%% % "scriptsize" is defined by floatrow, "tiny" not
%%% \floatsetup[table]{font=footnotesize}

% ################################################################### FIGURES #

\usepackage{adjustbox} % loads also graphicx

\graphicspath{%
  {/home/chejnik/Dokumenty/hvalur.org/images/biolib/full/}%
  {/home/chejnik/Dokumenty/hvalur.org/images/uploaded_files/}}

\newlength\dicfigheight
\newlength\dicfigskip

\newcommand\dicFigure[3]{ % {filename}{caption}{index entry}
  \setlength\fboxsep{0pt}\setlength\fboxrule{0.5pt}
  \def\dicfigbox{\fbox{%
    \adjincludegraphics [ max height = 0.8\columnwidth
                        , max width  = 0.8\columnwidth ] {#1}}}
  \setlength\abovecaptionskip {\dp\strutbox}
  \setlength\belowcaptionskip {\ht\strutbox} 
  \setlength\dicfigheight {\heightof\dicfigbox+\fboxrule}
  \setlength\dicfigskip
    {\baselineskip*\numexpr1+\dicfigheight/\baselineskip\relax-\dicfigheight}
  \begin{figure}[ht]
    \centering\vspace*{\dicfigskip}\dicfigbox\caption{#2}\index[figures]{#3}
  \end{figure}
}


% #################################################################### COLORS #

\usepackage{color}

% http://www.colorschemer.com/online.html
% main color - grammar color
% #66023C
\definecolor {darkgreen}          {rgb} {0.40, 0.01, 0.24}

\definecolor {royalazure}         {rgb} {0.00, 0.22, 0.66}
\definecolor {brown}              {rgb} {0.40, 0.01, 0.24}

% COLORS FOR THUMB INDEXES
\definecolor {babyblueeyes}       {rgb} {0.63, 0.79, 0.95}
\definecolor {unitednationsblue}  {rgb} {0.36, 0.57, 0.90}
\definecolor {blue(ryb)}          {rgb} {0.01, 0.28, 1.00}
\definecolor {darkblue}           {rgb} {0.00, 0.00, 0.55}
\definecolor {screamingreen}      {rgb} {0.46, 1.00, 0.44}
\definecolor {limegreen}          {rgb} {0.20, 0.80, 0.20}
\definecolor {islamicgreen}       {rgb} {0.00, 0.56, 0.00}
\definecolor {upforestgreen}      {rgb} {0.00, 0.27, 0.13}
\definecolor {icterine}           {rgb} {0.99, 0.97, 0.37}
\definecolor {orange(colorwheel)} {rgb} {1.00, 0.50, 0.00}
\definecolor {orange-red}         {rgb} {1.00, 0.27, 0.00}
\definecolor {oucrimsonred}       {rgb} {0.60, 0.00, 0.00}
\definecolor {cottoncandy}        {rgb} {1.00, 0.74, 0.85}
\definecolor {orchid}             {rgb} {0.85, 0.44, 0.84}
\definecolor {vividcerise}        {rgb} {0.85, 0.11, 0.51}
\definecolor {patriarch}          {rgb} {0.50, 0.00, 0.50}

\definecolor {title}              {RGB} {  16,   13,   32}
\definecolor {golden}             {RGB} { 241,  184,   45}
\definecolor {freqstar}           {RGB} {  91,    3,   99}
\definecolor {newdev}             {RGB} {   3,   11,   99}

\definecolor {color1}             {RGB} { 182,   86,    0}
\definecolor {color2}             {RGB} { 143,    9,    6}
\definecolor {color3}             {RGB} {   3,   23,  118}
\definecolor {color4}             {RGB} {   0,   82,  168}
\definecolor {color5}             {RGB} {   0,   85,  142}
\definecolor {color6}             {RGB} {   0,  115,  162}
\definecolor {color7}             {RGB} {  34,  146,  186}
\definecolor {color8}             {RGB} {  40,  159,  153}
\definecolor {color9}             {RGB} {   0,  125,  111}
\definecolor {color10}            {RGB} {   4,  107,   60}
\definecolor {color11}            {RGB} {  71,  134,   81}
\definecolor {color12}            {RGB} { 109,  134,   42}
\definecolor {color13}            {RGB} { 205,  194,   18}
\definecolor {color14}            {RGB} { 204,  162,   24}

%\definecolor {color1}             {RGB} { 239,  222,  205}
%\definecolor {color2}             {RGB} { 251,  206,  177}
%\definecolor {color3}             {RGB} { 251,  185,  142}
%\definecolor {color4}             {RGB} { 248,  166,  112}
%\definecolor {color5}             {RGB} { 216,  231,  246}
%\definecolor {color6}             {RGB} { 186,  215,  244}
%\definecolor {color7}             {RGB} { 152,  197,  243}
%\definecolor {color8}             {RGB} { 118,  181,  245}
%\definecolor {color9}             {RGB} { 181,  249,  185}
%\definecolor {color10}            {RGB} { 160,  248,  167}
%\definecolor {color11}            {RGB} { 128,  247,  136}
%\definecolor {color12}            {RGB} {  87,  248,   98}
%\definecolor {color13}            {RGB} { 250,  241,  138}
%\definecolor {color14}            {RGB} { 250,  240,  111}
%\definecolor {color15}            {RGB} { 244,  239,   85}
%\definecolor {color16}            {RGB} { 245,  239,   57}

\definecolor {lightgrey}          {RGB} { 105,  105,  105}
\definecolor {darkgreen_real}     {rgb} {0.00, 0.50, 0.00}


% ############################################################# THUMB INDEXES #

% Thumb indexes' colors
%\newcommand\BoxColor{%
%  \ifcase\theletternum darkgreen!30\or babyblueeyes\or unitednationsblue\or blue(ryb)\or screamingreen\or limegreen\or islamicgreen\or upforestgreen\or icterine\or orange(colorwheel)\or orange-red%
%  \or oucrimsonred\or cottoncandy\or orchid\or vividcerise\or patriarch\or babyblueeyes\or unitednationsblue\or blue(ryb)\or screamingreen\or limegreen\or islamicgreen\or upforestgreen\or icterine\or orange(colorwheel)\or orange-red%
%  \or oucrimsonred\or cottoncandy\or orchid\or vividcerise\or patriarch\else darkgreen!30\fi}
  
% Thumb indexes' colors
%\newcommand\BoxColor{%
%  \ifcase\theletternum darkgreen!30\or color1\or color2\or color3\or color4\or color5\or color6\or color7\or color8\or color9\or color10%
%  \or color11\or color12\or color13\or color14\or color1\or color2\or color3\or color4\or color5\or color6\or color7\or color8\or color9%
%  \or color10\or color11\or color12\or color13\or color14\or color1\or color2\or color3\or color4\or color5\or color6\else color7\or color8\or color9\or color10\or color11\or color12\fi}

%\def\BoxColor#1{darkgreen!\the\numexpr103-#1-#1-#1\relax!blue}

\newcommand\BoxColor[1]{%
\ifcase#1 darkgreen!30\or color13\or color14\or color1\or color2\or color3\or color4\or color5\or color6\or color7\or color8\or color9\or color10%
\or color11\or color12\or color13\or color14\or color1\or color2\or color3\or color4\or color5\or color6\or color7\or color8\or color9%
\or color10\or color11\or color12\or color13\or color14\or color1\or color2\or color3\or color4\or color5\or color6\else color7\or color8\or color9\or color10\or color11\or color12\fi}

\ifshowcrop
%%%%%%%%%%%%%%%%%%%%%%%%%%%%%%%%%%%%%%%%%%%%%%%%%%%%%%%%%%%%%%%%%%%%%%%%%%%%%%%%

\newcounter{letternum}
\newcounter{lettersum}                 \setcounter{lettersum}{34}
% these margins separate top thumb from top of page and last thumb from
% bottom of page. BEWARE, it's from top of the (cut out) page, not the folio!
\newlength{\thumbtopmargin}            \setlength{\thumbtopmargin}{1cm}
\newlength{\thumbbottommargin}         \setlength{\thumbbottommargin}{1cm}
\newlength{\thumbwidth}                \setlength{\thumbwidth}{2.20cm}
\newlength{\thumbheight}
\pgfmathsetlength{\thumbheight}%
  {(\csname Gm@layoutheight\endcsname-\thumbtopmargin-\thumbbottommargin)/\value{lettersum}}
% inner xsep of letter from box margin
\newlength{\thumbsep}                  \setlength{\thumbsep}{.30cm}
% overlap over cutting line
\newlength{\thumboverlap}              \setlength{\thumboverlap}{3cm}

\tikzset{
  thumb/.style={
    text=white,
    outer sep=0pt,
    inner xsep=\thumbsep,
    minimum height=\thumbheight,
    text width=\thumbwidth-2\thumbsep,
    font=\sffamily\bfseries,},
  thumb east/.style={thumb,
    xshift=\dimexpr(+\thumboverlap-\paperwidth+\csname Gm@layoutwidth\endcsname)/2\relax,
    align=left,anchor=north east,},
  thumb west/.style={thumb,
    xshift=\dimexpr(-\thumboverlap+\paperwidth-\csname Gm@layoutwidth\endcsname)/2\relax,
    align=right,anchor=north west,},
}

\def\thumbnew{}
\def\thumbold{}
\usepackage{everypage}
\AddEverypageHook{\if\relax\thumbnew\relax\xdef\thumbnew{\thumbold}\fi}

\def\ethumbs#1,#2\relax%
  {\if\relax#1\relax\else\eventhumb{#1}\fi%
   \if\relax#2\relax\else\ethumbs#2\relax\fi%
   \gdef\thumbnew{}%
   \gdef\thumbold{#1,}}

\def\othumbs#1,#2\relax%
  {\if\relax#1\relax\else\oddthumb{#1}\fi%
   \if\relax#2\relax\else\othumbs#2\relax\fi%
   \gdef\thumbold{#1,}%
   \gdef\thumbnew{}}

\newcommand{\drawthumb}[2]{%
  \begin{tikzpicture}[remember picture, overlay]
    \node [thumb #2, fill=\BoxColor{#1}]
       at ($(current page.north #2)-%
            (0,\csname Gm@layoutvoffset\endcsname+\thumbtopmargin-\thumbheight+%
               #1*\thumbheight)$) {\csname Let#1\endcsname};
   \end{tikzpicture}}



%%%%%%%%%%%%%%%%%%%%%%%%%%%%%%%%%%%%%%%%%%%%%%%%%%%%%%%%%%%%%%%%%%%%%%%%%%%%%%%%
\else
% THUMB INDEXES
% new counter to hold the current number of the letter to determine the vertical position
\newcounter{letternum}
% newcounter for the sum of all letters to get the right height of a box
\newcounter{lettersum}
\setcounter{lettersum}{34}
% some margin settings
\newlength{\thumbtopmargin}
\setlength{\thumbtopmargin}{\ifshowcrop 3cm\else 1.5cm\fi} %1cm
\newlength{\thumbbottommargin}
\setlength{\thumbbottommargin}{\ifshowcrop 6cm\else 2cm\fi} %2.5cm
% calculate the box height by dividing the page height
\newlength{\thumbheight}
\pgfmathsetlength{\thumbheight}{%
(\paperheight-\thumbtopmargin-\thumbbottommargin)%
/%
\value{lettersum}
}
% box width
\newlength{\thumbwidth}
\setlength{\thumbwidth}{\ifshowcrop 2cm\else 0.5cm\fi} %0.5cm
% style the boxes
\tikzset{
thumb/.style={
   text=white,
   minimum height=\thumbheight,
   text width=\thumbwidth,
   outer sep=0pt,
   font=\sffamily\bfseries,
 }
 }

\def\thumbnew{}
\def\thumbold{}
\usepackage{everypage}
\AddEverypageHook{\if\relax\thumbnew\relax\xdef\thumbnew{\thumbold}\fi}

\def\ethumbs#1,#2\relax{\if\relax#1\relax\else\eventhumb{#1}\fi%
                        \if\relax#2\relax\else\ethumbs#2\relax\fi%
                        \gdef\thumbnew{}%
                        \gdef\thumbold{#1,}%
}

\def\othumbs#1,#2\relax{\if\relax#1\relax\else\oddthumb{#1}\fi%
                        \if\relax#2\relax\else\othumbs#2\relax\fi%
                        \gdef\thumbold{#1,}%
                        \gdef\thumbnew{}%
}

\newcommand{\drawthumb}[2]{%
  % see pgfmanual.pdf for more information about this part
  \begin{tikzpicture}[remember picture, overlay]
    \node [thumb, fill=\BoxColor{#1}, text centered, anchor=north #2]
       at ($(current page.north #2)-%
%            (0,\thumbtopmargin+\value{letternum}*\thumbheight)%
            (0,\thumbtopmargin+#1*\thumbheight)$) {\csname Let#1\endcsname};
   \end{tikzpicture}}
\fi

\newcommand{\oddthumb} [1]{\drawthumb{#1}{west}}
\newcommand{\eventhumb}[1]{\drawthumb{#1}{east}}

\newcommand{\lettergroup}[1]%
  {\refstepcounter{letternum}%
   \expandafter\gdef\csname Let\theletternum\endcsname{#1}%
   \xdef\thumbnew{\theletternum,\thumbnew}%
   \fancyhead[LO]{\phvfamily\bfseries\rightmark%
     \expandafter\ethumbs\thumbnew\relax\relax\relax}%
   \fancyhead[RE]{\phvfamily\bfseries\leftmark%
     \expandafter\othumbs\thumbnew\relax\relax\relax}}
     
% ############################################################## LOCALIZATION #

% renames the index name
%\addto\captionsczech{%
%  \renewcommand\indexname{{Seznam autorů fotografií a ilustrací}}}

% renames the contents name
\addto\captionsczech{%
  \renewcommand\contentsname{Obsah}}

\renewcommand \contentsname {Obsah}
\renewcommand \chaptername  {Kapitola}


% ################################################################ SECTIONING #
\usepackage[explicit]{titlesec}

\setlength{\beforetitleunit}{\baselineskip}
\setlength{\aftertitleunit} {\baselineskip}

\titlespacing*{\chapter}   {0em}{*0}{*3} % Remembere there's a topskip too
\titlespacing*{\section}   {0em}{*1}{*2}
\titlespacing*{\subsection}{0em}{*1}{*2}

\newcommand\ghostdrop[2]{%
  \raisebox{-#1\baselineskip}[0pt][0pt]{#2}}

\titleformat {\chapter} {} {} {0pt} % {#1}
  {\ghostdrop{1.0}{\parbox[t]{\columnwidth}
     {\LARGE\bfseries\centering\MakeUppercase{#1}}}}

\def\longchapterskip{\vspace{2\baselineskip}}
\def\longsectionskip{\vspace{2\baselineskip}}

\titleformat {\section} {\bfseries} {} {0pt}
  {\ghostdrop{1.0}{\parbox[t]{\columnwidth}
     {\Large\bfseries\centering#1}}}

\titleformat {\subsection} {\bfseries} {} {0pt}
  {\ghostdrop{1.0}{\hspace{1em}\parbox[t]{\columnwidth-1em}
     {\large\bfseries#1}}}


\setlength\multicolsep{0pt}
%\BeforeBeginEnvironment{multicols}{\vspace{-\dp\strutbox}}


% ######################################################## TABLE OF CONTENTS #
\usepackage{titletoc}

\titlecontents {chapter} [0pt] {\addvspace\baselineskip\bfseries}
  {} {} {\titlerule*[2\baselineskip]{}\contentspage}

\titlecontents {section} [\baselineskip] {} {} {}
  {\titlerule*[1\baselineskip]{.}\contentspage}


% #################################################### ASSORTED CUSTOM MACROS #

\def\textIS#1{\foreignlanguage{icelandic}{#1}}
\def\textCS#1{\foreignlanguage{czech}{#1}}
\def\textLA#1{\foreignlanguage{latin}{#1}}

\usepackage{accsupp}

\makeatletter                                                                                                          
  \def\ACCSUPP@bdc{\special{pdf:code \ACCSUPP@span BDC}}                                                               
  \def\ACCSUPP@emc{\special{pdf:code EMC}}                                                                             
\makeatother

\usepackage{pgffor}

\newcommand\blspace[1][1]{\vspace{#1\baselineskip}}

\usepackage{xparse}


\newcommand \dicsymFrequent  {\raisebox{-.2ex}{\color{darkgreen}\ding{167}}}
\newcommand \dicsymSee       {$\shortrightarrow$}
\newcommand \dicsymCompare   {$\shortuparrow$}
\newcommand \dicsymIdiom     {$\diamondsuit$}
\newcommand \dicsymExampleIS {$\triangleright$}
\newcommand \dicsymExampleCS {\guilsinglright}
\newcommand \dicsymProverb   {{\color{newdev}\ding{96}}}
\newcommand \dicsymPhoto
  {\includegraphics[keepaspectratio,height=0.65em]{photo_icon}}

%\NewDocumentCommand\dicTerm{mo}
 % {{\phvfamily\bfseries\textIS{#1}}\IfValueT{#2}{, \phvfamily\bfseries\textIS{#2}}}

\NewDocumentCommand\dicEntry{o}
  {\par\hangpara{0.25\baselineskip}{1}%
   \IfValueT{#1}{\markboth{#1}{#1}}%
    \ifPDFsearchvalue
    \settowidth{\eyja}{#1}%
  \makebox[\eyja]{\color{white}#1}\hspace{-\eyja}%
  \fi
   \ignorespaces}
   
   \NewDocumentCommand\dicEntryGuide{o}
  {\par\hangpara{0.29\baselineskip}{1}%
   \IfValueT{#1}{\markboth{#1}{#1}}%
    \ifPDFsearchvalue
    \settowidth{\eyja}{#1}%
  \makebox[\eyja]{\color{white}#1}\hspace{-\eyja}%
  \fi
   \ignorespaces}
   
\newcommand\dicSubEntry
  {\\\hspace*{-\hangindent}}

\NewDocumentCommand\textAlt{mm}
  {\BeginAccSupp{method=pdfstringdef,ActualText={#2}}#1\EndAccSupp{}}

\newcounter{terms}

\NewDocumentCommand\dicStyleIntTerm{m}{\textIS{{\phvfamily\bfseries#1}}}
\NewDocumentCommand\dicStyleExtTerm{m}{\textIS{\textit{#1}}}

   \NewDocumentCommand\dicTerm{mo}
  {{\closeupdotbar \phvfamily\bfseries\textIS{#1}}\IfValueT{#2}{, \phvfamily\bfseries\textIS{#2}}}

  \newcommand\dicTermList[1]
  {\foreach \t [count=\n] in {#1}{\ifnum\n=1\else, \fi\dicTerm{\t}}}
  
\def\dicIPA#1{\textipa{[#1]}}

\NewDocumentCommand\dicPos{smo}
  {{\color{darkgreen}\textbf{\small#2}%
   \IfValueT{#3}{\nobreak\textsubscript{#3}}}}

%\NewDocumentCommand\dicFlx{smo}
 % {{\color{darkgreen}\footnotesize#1\IfValueT{#2}{\textsubscript{#2}}}}

  \NewDocumentCommand\dicFlx{smo}
  {{\color{darkgreen}\footnotesize#2%
   \IfValueT{#3}{\nobreak\textsubscript{#3}}}}
  
\newcommand\dicExampleCS[1]{\dicsymExampleCS~\textit{#1}}
\newcommand\dicExampleIS[1]{\dicsymExampleIS~\textIS{\textit{#1}}}

\NewDocumentCommand\dicLink{sm}
  {\IfBooleanF#1{\dicsymSee\nobreak\hspace{.08667em plus .03333em}}\nobreak\hspace{.03667em}\dicTerm{#2}}
\NewDocumentCommand\dicSynonym{sm}
  {(\IfBooleanF#1{\dicsymSee\nobreak\hspace{.08667em plus .03333em}}\nobreak\hspace{.03667em}\textIS{\textit{#2}})}
\NewDocumentCommand\dicAntonym{sm}
  {(\IfBooleanF#1{\dicsymCompare\nobreak\hspace{.08667em plus .03333em}}\nobreak\hspace{.03667em}\textIS{\textit{#2}})}
\NewDocumentCommand\dicCompare{sm}
  {(\IfBooleanF#1{\dicsymCompare\nobreak\hspace{.08667em plus .03333em}}\nobreak\hspace{.03667em}\textIS{\textit{#2}})}


\newcommand\dicPhraseIS[1]{\textIS{\textbf{#1}}}

\NewDocumentCommand\dicIdiom{mo}
  {\dicSubEntry{\color{newdev}\dicStyleIntTerm{#1} %
     \IfValueTF{#2}{+ \dicStyleIntTerm{#2} \markboth{#1 + #2}{#1 + #2}}{\markboth{#1}{#1}}}\dicsymIdiom}
     
\newcommand\dicProverb
  {\dicSubEntry\dicsymProverb}



\newcommand\dicLangCat[1]{{\footnotesize#1}}
\newcommand\dicFieldCat[1]{{\footnotesize#1}}


\newcommand\dicDirectTranslationCS[1]{#1}
\newcommand\dicIndirectTranslationCS[1]{{\footnotesize#1}}

\usepackage{placeins}

\newcommand*{\dicLetter}[2]{%
  \FloatBarrier\newpage\phantomsection\pdfbookmark[1]{#1}{#2}%
  \noindent\parbox[b][9\baselineskip][c]{\columnwidth}
    {\centering\HUGE\strut\smash{\MakeUppercase{#1}~\MakeLowercase{#1}}}%
      %\par\xdef\tpd{\the\prevdepth}
      %\par\prevdepth\tpd%
  \expandafter\lettergroup{\MakeLowercase{#1}}}

  
%tabulky
\newcommand{\specialcell}[2][l]{%
\begin{tabular}[#1]{@{}l@{}}#2\end{tabular}}


%\newcommand{\devisionguide}[1]{\hspace*{0.3cm}{{{{\foreignlanguage{icelandic}{\color{newdev}{\phvfamily{\textbf{#1}}}}}}}}}
%\newcommand{\devisionguideindent}[1]{\hspace*{-0.1cm}{{{{\foreignlanguage{icelandic}{\color{newdev}{\phvfamily{\textbf{#1}}}}}}}}}

%\newcommand\n{$n$\nobreakdash-\hspace{0pt}} %nonbreakable dash in grammar endings
%\def \nobreakseq {\nobreak \hskip 0pt \hbox} %nolinebreak in case like (-s, -)

% rule line
\newcommand{\HRule}{\rule{\linewidth}{0.1mm}}

% width of pdf found box
\newlength\eyja

%\newcommand\strutsmash[1]{\strut\smash{#1}}


% Dictionary symbols


% Command synonyms
%\let\freqstar\dicsymFrequent
%\let\photoicon\dicsymPhoto

\newcommand{\vstretch}[1]{\vspace*{\stretch{#1}}}

\newcommand*{\addthin}{\hskip0.16667em\relax}
\newcommand*{\addthinS}{\hskip0.06667em\relax}
\newcommand*{\addthinM}{\hskip0.10007em\relax}
\newcommand*{\addthinSS}{\hskip0.00007em\relax}
% #################################################################### COVERS #

\def\dictnameCZ{islandsko-český studijní slovník}
\def\dictnameIS{íslensk-tékknesk stúdentaorðabók}

\newcommand*{\frontcoverimages}{%
  ds_image_blaberjalyng_0_2.jpg,
  ds_image_lundi_0_1.jpg,
  ds_image_hestur_0_1.jpg,
  ds_image_hrafnafifa_0_2.jpg,
  ds_image_baldursbra_0_1.jpg,
  20948.jpg,
  ds_image_hvalur_0_1.jpg,
  ds_image_blodberg_0_1.jpg}

\newcommand*{\backcoverimages}{%
  ds_image_spoi_0_1.jpg,
  ds_image_mosalyng_0_1.jpg,
  ds_image_blaklukkulyng_0_1.jpg,
  20787.jpg,
  ds_image_ufsi_0_2.jpg,
  %ds_image_gullbra_0_1.jpg,
  ds_image_vatnsberi_0_1.jpg,
  ds_image_islandsfifill_0_2.jpg,
  21137.jpg}

\newcommand\makecoverwith[1]{
 \pagecolor{title}\afterpage\nopagecolor
 \begin{center}
 \vspace*{0.3cm}%{\fill}
 {\Huge\scshape\color{white}%
   \dictnameCZ\\\dictnameIS}
 \vspace*{1em}                        % I think this is necessary
 \begin{figure}[ht]\centering%
   \setlength\fboxsep{0pt}\setlength\fboxrule{0.5pt}%
   \foreach \file in #1{
     \fbox{\includegraphics[width=5.5cm]{\file}}}
 \end{figure}
 \vspace*{\fill}
 \end{center}
}


% ###################################################### HYPERREF & PDF INFO #

% unicode neccessary so that national characters in hypersetup appear ok
\usepackage[pdftex, unicode, hyperfootnotes=false]{hyperref}

% hyperlinks in black
\makeatletter
\let\Hy@linktoc\Hy@linktoc@none
\makeatother

\hypersetup
  { pdftitle={Islandsko-český studijní slovník}
  , pdfauthor={Aleš Chejn, Jón Gíslason, Ján Zaťko}
  , pdfsubject={Islandsko-český studijní slovník}
  , pdfkeywords={Islandsko-český studijní slovník%
               , islandština, čeština, slovník}
  , bookmarks=true
  , colorlinks=true
  , citecolor=black
  , urlcolor=\ifscreen darkgreen\else black\fi
  , linkcolor=black}


% ################################################################# MICROTYPE #

\usepackage[final]{microtype}

% ############################################################### HYPHENATION #

% this allows to include brackets inside hyphenation rules f.e (pro)-dis-ku-to-vat
\lccode`\(`\(
\lccode`\)`\)
\lccode`\[`\[
\lccode`\]`\]

\begin{hyphenrules}{icelandic}
\hyphenation{
  frum-semja
  burm-neskur
  ó-breytan-le-gur
  pla-tó-nskt
  skatta-fram-tal
  skrúf-gang-ur
  strand-rauð-viður
  prent-sver-ta
  grun-nur
  uppi-skrop-pa
  ger-vi-mennska
  sö-lu-stað-ur
  sý-na-gó-ga
  mark-mið
  trau-st-ur
  ís-kyggi-leg-ur
  brjó-sta-höld
  kné-krjú-pa
  nel-li-ka
  kaffi-kvörn
  margs-konar
  skrúf-gang-ur
  uppi-skrop-pa
  veð-rað-ur
  vest-ur-átt
  mis-mun-andi
  Bar-ba-do-si
  tung-u-mál
  au-ga-brún
  há-sker-pu-sjón-varp
  rúnn-stykki
  ger-vi-tungl
  yfir-borðs-menn-ska
  gæ-fu-mun-ur
  hlí-fi-skil-di
  á-vís-un
  ó-um-flýj-an-leg-ur
  san-mar-ín-sk-ur
  gæ-ða-stað-all
}
\end{hyphenrules}
\begin{hyphenrules}{czech}
\hyphenation{
  teo-rie
  mi-liardy
  (z)-od-po-věd-nost
  tan-tié-ma
  alian-ce
  pe-dia-tr(ička)
  pa-cienty
  zá-chra-nář(ka)
  dia-man-to-vá
  poe-zie
  ma-te-riál-ní
  egoi-s-ta
  egoist-ka
  nai-vi-ta
  ne-upřím-nost
  pě-ti-úhel-ník
  še-s-ti-úhel-ník
  za-šmo-dr-chat
  ar-cheo-lo-gic-ký
  ne-omlu-vi-tel-ný
  (vy)-pro-du-ko-vat
  dia-gno-s-ti-ko-vat
  teo-lož-ka
  ateist-ka
  oceá-nu
  im-pe-ria-li-s-ta
  (vy)-zdvih-nout
  ne-úměr-ný
  kon-ti-nuál-ní
  neu-t-rum
  ne-opráv-ně-ně
  vig-vam
  pro-zaic-ký
  geo-lož-ka
  seiz-mic-ká
  teo-rém
  olean-dr
  poe-ta
  (vy)-fo-to-gra-fo-vat
  (vy)-ven-ti-lo-vat
  pneu-mo-nie
  Lu-cem-bur-čan-(ka)
  ma-te-riály
  am-bi-cióz-ní
  Myan-ma(r)
  Myan-ma-řan-(ka)
  mol-dav-šti-na
  si-tua-cí
  ne-ukáz-ně-ný
  he-te-ro-trof-ní
  ne-omlu-vi-tel-ná
  ne-upřím-ný
  dvou-bliz-ný
  ne-úmy-sl-ný
  ne-ustá-lý-mi
  ne-ohle-du-pl-ný
  ne-osob-ní
  ne-oso-le-ný
  ne-uvě-ři-tel-ný
  ne-ohro-že-ný
  ne-oby-čej-ný
  ne-útul-ný
  ne-únav-ný
  ne-ože-nit
  po-ukáz-ka
  za-úto-čit
  aero-dy-na-mic-ký
  hru-bián-ství
  (pro)-di-s-ku-to-vat
  (vy)-ča-lou-nit
  (na)-oč-ko-vá-ní
  Bhú-tán-ka
  (za)-šně-ro-vat
  soud-ní
  re-vo-lu-cio-nář(ka)
  á-zer-báj-d-žán
  čtyř-hran-ný
  ze-m-dle-lý
  klou-zač-ka
  Vest-mann-ské
  tří-stran-ný
  (za)-bu-rá-ce-ní
  os-tud-ný
  pro-du-cent-(ka)
  chvá-s-tal-(ka)
  apar-theid
  asis-ten-ce
  asis-tent-(ka)
  (u)-či-nit
  (u)-za-mknout
  past-vi-na
  (za)-zá-lo-ho-vat
  Af-ri-čan-(ka)
  (za)-an-ga-žo-vá-ní
  se-rióz-ní
  opo-nent-(ka)
  pro-tiv-ník
  zvlášt-ní
  (po)-div-ný
  ho-mo-se-xuál-ní
  ko-alič-ní
  si-tua-ce
  si-tua-cí
  egoi-s-mus
  spe-ciál-ní
  spe-cia-li-zo-va-ný
  spe-cia-li-zo-va-né
  sobě-stač-ný
  sa-mo-zvo-le-ný
  sa-mo-čin-ně
 chvá-s-ta-vost
 klient-(ka)
 ku-rióz-ní
 (na)-šrou-bo-vat
 (za)-šrou-bo-vat
 krea-ti-vi-ta
 pneu-ma-ti-ka
 toa-le-ta
 so-cia-list-ka
 (vy)-dez-in-fi-ko-va-ný
 vše-uměl
 vše-uměl-(kyně)
 (pře)-škrt-nout
 (vy)-svlék-nout
 vy-učo-va-cí
 vy-učo-va-cích
 (vy)-chla-s-tat
 (vy)-te-to-vat
 pa-cient-(ka)
 ne-opod-s-tat-ně-ný
 de-se-ti-úhel-ník
 pe-rio-di-kum
 (za)-mum-lá-ní
 ne-smě-lá
 ne-dů-vě-ři-vá
 tee-na-ger-(ka)
 (z)-de-ge-ne-ro-vá-ní
 mo-ci-chti-vý
 bez-úhon-nost
 zne-uctít
 vo-dě-odol-ný
 we-bo-vý
 web-de-sign
 in-že-nýr-(ka)
 zá-pa-do-is-land-ský
 klient-(ka)
 za-uja-tý
 vy-uží-va-né
 vy-uží-vat
 (za)-re-gi-s-t-ro-vat
 (z)-kom-pri-mo-va-ný
 na-cio-na-li-s-mus
 ma-k-ro-eko-no-mic-ká
 bez-bo-lest-ný
 viol-ka
 troj-úhel-ník
 (za)-bu-rá-ce-ní
 ne-ústup-nost
 (vy)-cvi-če-ný
 zprá-va
 ge-ria-tr-(ička)
 osud-ný
 ná-pra-vy
 vstup-ní
 obec-ně
 re-pub-li-ka
 (po)-div-ný
 pa-slun-ce
 vliv-ný
 ba-ham-ský
 pe-dia-tr-(ička)
 ha-sič-(ka)
 vnuč-ka
 bro-s-kvoň
 lo-di-vod-(ka)
 mi-s-tr-(yně)
 in-že-nýr-(ka)
 dr-ži-tel-(ka)
 to-pič-(ka)
 pro-mi-nent-ní
 takt-ní
 frag-men-tár-ní
 po-zo-ru-hod-ný
 úspě-š-ný
 trans-fu-ze
 bur-žoa-zie
 pře-stro-je-ní
 oh-ni-vzdor-ný
 dis-tan-ční
 de-frau-da-ce
 fi-s-kál-ní
 zběh-lý
 kvap-ný
 střed-ní
 sa-mo-hlá-s-ka
 pře-vo-di-tel-ný
 zban-k-ro-to-va-ný
 zá-clo-na
 zná-mý
 zá-klad-ní
 na-uka
 ro-š-ťa-č-ka
 lí-bán-ky
 sou-du
 zkou-š-ka
 mlý-nek
 prá-šiv-ka
 kul-ma
 la-hod-ný
 zá-meč-ni-ce
 ne-vý-znam-ný
 ovzdu-ší
 mno-ho-úhel-ní-ko-vý
 žra-vý
 zcest-ný
 esen-ce
 (při)-nu-ce-ný
 skr-b-lic-ký
 skle-s-lý
 bez-bož-ný
 ne-opráv-ně-ný
 ošu-mě-lý
 pru-der-ní
 chvá-s-ta-vý
 sou-drž-ný
 van-dal-(ka)
 sa-bo-tér-(ka)
 (po)-div-ně
 mo-k-řad
 plyt-ký
 na-spo-du
 izo-la-ce
 le-toun
 za-han-be-ný
 ří-ze-ní
 kom-plex-ní
 pr-kén-ko
 ba-vič-(ka)
 plyt-ký
 ne-mo-tor-ně
 šéf-(ka)
 přes-čas
 zji-s-ti-tel-ný
 ta-ko-vý
 kňou-ral-(ka)
 se-me-s-tr
 po-pra-s-ka-ný
 han-ba
 vzpří-me-ný
 chys-ta-ný
 lik-vi-da-ce
 or-ga-ni-za-ce
 ri-s-kant-ní
 (za)-pís-kat
 (za)-hvíz-dat
 vrou-cí
  }
\end{hyphenrules}


% ############################################################# CZECH TYPOGRAPHY #
% https://math.feld.cvut.cz/tkadlec/ftp/text/textypo.pdf



% ################################################################# DOCUMENT #

\begin{document}
\pagenumbering{gobble}

% -------------------------------------------------------------- Front cover -

\ifmakecovers
\covergeometry
  \makecoverwith\frontcoverimages
\restoregeometry  
\fi


% ============================================================= FRONT MATTER =

\ifmakefrontmatter

% ------------------------------------------------------------------ Patitul -

\cleardoublepage

\begin{center}
  %\vspace*{2cm}{\fontsize{20}{20}\selectfont\scshape\dictnameCZ\\\dictnameIS}
  \vspace*{\stretch 1}
  {\Huge\scshape\dictnameCZ\\\dictnameIS}
  \vspace{\stretch 4}
\end{center}

% -------------------------------------------------------------------- Titul -

\cleardoublepage

\begin{center}
  %\vspace*{2cm}{\fontsize{24}{24}\selectfont\scshape\dictnameCZ\\\dictnameIS}
  \vspace*{\stretch 1}
  {\HUGE\scshape\dictnameCZ\\\dictnameIS}
  \vspace{\stretch 2}
\end{center}

% ----------------------------------------------------------------- Impresum -

\clearpage

\noindent\begin{minipage}[t][\textheight][t]{\textwidth}
Autoři \textit{Islandsko-českého studijního slovníku}:\\
  \textbf{Aleš Chejn, Ján Zaťko, Jón Gíslason}\\[\baselineskip]
\ifPDF
\else
Autor \textit{Islandské gramatiky}:\\
  \textbf{Vojtěch Kupča}\\[\baselineskip]
\fi
Poděkování:\\
  \textbf{Dorotě Nierychlewské-Chejn, Jiřině Hanzlové, Jiřině Chejnové,
          Simoně Petrásové,\\ Petru Mikešovi, Vojtěchu Kupčovi, 
          Renatě Peškové Emilsson, Amiru Mulahumici, \\ Kristýně Antonové, 
          Martině Kašparové, Lucii Brabcové, Samoriele, Adamu Kožouškovi}\\[\baselineskip]
Finanční příspěvek na výtisk exemplářů určených knihovnám\\ v České republice,
Slovenské republice a na Islandu poskytli:\\
  \textbf{Petr Mikeš, Dorota Nierychlewska-Chejn, Jiřina Chejnová, Otta Chejn}
  \ifPDF
  \\[39\baselineskip]
  \else
  \\[36\baselineskip]
  \fi
\textbf{Hvalur.org - 2016}\\
Licence: \textbf{GNU Free Documentation License}\\
Sazba: \textbf{Aleš Chejn} v prostředí  \textbf\LaTeX\\
  Zvláštní poděkování patří \textbf{Paolu al-Imk\=an\=\i\ Brasolin} a uživatelům  \textbf{\TeX\ - \LaTeX\ Stack Exchange}\\[\baselineskip]
\ifPDF
Verze: \textbf{PDF verze} \textit{Islandsko-českého studijního slovníku} vhodná pro \textbf{\ifscreen offline prohlížení a vyhledávání\else vlastní tisk\fi}\\ volně dostupná na \url{www.hvalur.org}\\[\baselineskip]
\else
\textbf\ISBN\\[\baselineskip]
\fi
\ifPDF
\else
Tisk: \textbf{Tiskárna Melmen, Pardubice}\\
\fi
Heslová část vysázena patkovým fontem \TeX\ Gyre Pagella,
hesla bezpatkovým fontem {\phvfamily Helvetica}.
\end{minipage}

% -------------------------------------------------------- Table of contents -

\cleardoublepage 

\pdfbookmark\contentsname{toc}    
\tableofcontents
\thispagestyle{empty}
\cleardoublepage 

\fi

\pagenumbering{arabic}
\setcounter{page}{3}
% ============================================================= INTRODUCTION =

\ifinputintroduction

% ------------------------------------- The study Icelandic-Czech dictionary -

\chapter{O Islandsko-českém studijním slovníku}
Cílem vytvoření\textit{ Islandsko-českého studijního slovníku} je zpřístupnění islandského jazyka všem zájemcům o~tento jazyk. 

\textit{Islandsko-český studijní slovník} je určen především českým a~slovenským studentům islandštiny, začátečníkům i~pokročilým, kteří zde naleznou ucelené fonetické, morfologické a~syntaktické informace, jakož
i~slovní zásobu potřebnou k~aktivnímu i~pasivnímu užití s~rozsahem vhodným ke studiu islandštiny na vysoké škole.
Slovník je rovněž určen všem zájemcům o~islandštinu a~Island, jimž slovník nabízí encyklopedické informace, úvod do fonetiky a~gramatiky a~velké množství fotografií a~ilustrací. 
Třebaže se jedná o~slovník primárně zaměřený na potřeby českých a~slovenských uživatelů, obsahuje také řadu informací, které usnadní práci islandskému uživateli (např. použití více jak 10 000 synonym a~antonym, překlady příkladů ap.).

Slovní zásoba původně vychází z~\textit{Concise Icelandic-English Dictionary} (\cite {ic_en}), značně však byla rozšířena díky přístupu k~webovému projektu \textit{ISLEX} (\cite {int1}). Velké množství hesel bylo
přidáno na podnět uživatelů online verze \textit{Islandsko-českého studijního slovníku}.

Slovník obsahuje 30~575 hesel s~46~725 významy nebo slovními spojeními. Z~celkového počtu 30~575 je 28~097 svébytných hesel obsahujících alespoň jeden význam nebo slovní spojení. 
Zbylá hesla (v~počtu 2~478) jsou varianty, dublety nebo nepravidelné tvary.

Slovník je primárně online elektronický slovník, volně přístupný na internetových stránkách \url{www.hvalur.org}  (\cite {int14}). 
Tištěná verze slovníku je publikována na vlastní náklady v~počtu 30 kusů, přičemž většina exemplářů je určena knihovnám a~institucím v~České republice, Slovenské republice a~na Islandu
a~několik exemplářů je určeno autorům nebo zájemcům z~řad veřejnosti. Dvě PDF verze (vyhledávání\,/\addthin prohlížení offline a~vlastní tisk) jsou dostupné zdarma na internetových stránkách slovníku (viz výše).

\textit{Islandsko-český studijní slovník} je otevřený svobodný projekt, jehož se v~průběhu tvorby zúčastnilo velké množství spoluautorů (viz Historie tvorby IČSS), kterým tímto jménem autorů velmi děkuji.
Děkuji také všem uživatelům online verze slovníku za trpělivost, dobrou vůli a~veškeré připomínky a~návrhy, jmenovitě především Petru Mikešovi, Adamu Kožouškovi a~Samoriele za neutuchající víru v~náš projekt.
Poděkování patří také našim rodinám a~blízkým, díky kterým jsme se mohli soustavně věnovat tvorbě slovníku.
Srdečné poděkování patří mé ženě Dorotce za podporu a~nadšení potřebné k~tvorbě slovníku.

\blspace[5]

{\centering Aleš Chejn, Brwinów, 24. prosince 2015\par}


% ------------ The history of the formation Icelandic-Czech dictionary study -

\chapter{Historie tvorby Islandsko-českého studijního slovníku}
\longchapterskip
Projekt \textit{Islandsko-českého studijního slovníku} (IČSS) vznikl v~roce 2001 na nepřímý podnět doc. PhDr. Heleny Kadečkové, CSc. z~Univerzity Karlovy,
kdy byl vypracován seznam slov. 
Tento seznam obsahoval islandská slova a~jejich překlady, později byly přidány základní morfologické informace o~slovním druhu a~deklinační\,/\addthin konjugační koncovky podle vzoru z~\textit{Concise Icelandic-English Dictionary} (\cite {ic_en}). 
Seznam v~této době obsahoval okolo 1 200 slov. Byl zveřejněn pomocí statických webových stránek a~později byl umístěn na webu Univerzity Karlovy. 

V~roce 2006 byl slovník rozšířen o~nových 1 500 heslových slov a~stávající hesla byla rozšířena o~nové významy, zvláště významy frázových sloves. 
Ke konci roku 2006  byl slovník převeden z programu Excel do databáze MySQL a~byly zahájeny práce na vytvoření aplikace \textit{Dictionary System} (DS) (\cite {int19}), která by podporovala tvorbu IČSS. 
Obsah slovníku byl zveřejněn pod svobodnou licencí na webových stránkách \url{www.hvalur.org} (\cite {int14}), kde se nachází dodnes. 

V~roce 2007 byla vytvořena první verze aplikace DS. V~této době se k~projektu přidala Renata Pešková Emilsson, která se nyní věnuje podpoře výuky češtiny na Islandu. 
Ta slovník obohatila rozsáhlým seznamem slov z~praktického života na Islandu. Dále se slovník rozšiřoval o~heslová slova podle \textit{Concise Icelandic-English Dictionary} (\cite {ic_en}). 
V průběhu následujících třech let bylo doplněno přibližně 15~000 slov. V této době byla hesla rozšířena o~synonyma, příklady, frekvenci, oborovou a~stylovou charakteristiku. K projektu se na Islandu připojila na krátkou dobu Kristýna Antonová, studentka islandštiny na Islandské univerzitě, která 
nás upozornila na nesrovnalosti v~uživatelském rozhraní aplikace DS.

Během léta 2009 byl vytvořen skript na generování skloňování a~časování heslových slov. Výsledky generování byly porovnány s~\textit{Beygingarlýsing íslensks nútímamáls} (\cite {DalvikVM}) a~během následujících let konzultovány
s~Jónem Gíslasonem. Kristín Bjarnadóttir, vedoucí projektu \textit{Beygingarlýsing íslensks nútímamáls}, zodpověděla řadu našich dotazů týkajících se deklinací a~konjugací.

V~roce 2009 začala spolupráce s~\textit{Biolib.cz} (\cite {int5}), jmenovitě s~Ondřejem Zichou. Obdrželi jsme  databázi s~náhledy obrázků, autory a~licencemi a~latinskými názvy botanických a~zoologických druhů.

V~zimním semestru 2009 jsme požádali s~Amirem Mulahumicem, bosenským studentem islandštiny na Islandské univerzitě, Studentský fond Islandské univerzity o~grant na tvorbu zvukové databáze islandských slov. 
Studentský fond návrh přijal. Peníze z~grantu jsme určili pro islandského rodilého mluvčího, Jóna Gíslasona, lektora islandského jazyka na Islandské univerzitě v~Reykjavíku, který nahrál více jak 22~000 islandských slov do elektronické formy. 
S~Dorotou Nierychlewskou-Chejn jsme připravili dva seznamy heslových slov (jeden pro nahrávajícího a~druhý pro zpracování nahrávek). Na podporu nahrávek jsme vytvořili webové stránky nazvané \textit{Databáze s~výslovností islandských slov} (\cite {int13}) a~tyto stránky jsme přeložili do 6~jazyků (Jón Gíslason do islandštiny, Ján Zaťko do slovenštiny, francouzštiny a~angličtiny, Amir Mulamuhic do srbochorvatštiny a~Dorota Nierychlewska-Chejn do polštiny).

V~roce 2010 jsme přidali do databáze pravidla islandské výslovnosti, která jsou uvedena v~knize \textit{Handbók um íslenskan framburð} (\cite {is2}). Z~těchto pravidel jsme automaticky vygenerovali fonetický zápis výslovnosti všech hesel ve slovníku. 
Jón Gíslason ověřil správnost fonetického zápisu u~všech heslových slov.

Na konci roku 2011 jsme spojili webové stránky slovníku se stránkami aplikace a~sjednotili jsme a~vylepšili vzhled webových stránek. Cennými připomínkami se na vytvoření aplikace podílela
RNDr. Ingrid Nagyová, PhD., vysokoškolská pedagožka Ostravské Univerzity. Ján Zaťko přeložil nové webové stránky, Průvodce po slovníku do slovenštiny, francouzštiny a~angličtiny a~celý manuál aplikace do angličtiny. 
Dorota Nierychlewska-Chejn přeložila stránky do polštiny, Jón Gíslason do islandštiny a~Marcos Helena do portugalštiny.

Od prosince 2011 do března 2014 jsme realizovali první čtení, ve kterém jsme doplnili chybějící informace, přidali nová hesla a~opravili pravopisné chyby. 
Obsah slovníku jsme porovnávali především s~\textit{Íslensk orðabók} (\mbox{\cite {is}}) a~\textit{Islandsko-českým slovníkem} (\cite {is7}). S~Vojtěchem Kupčou jsme konzultovali mnohé nesnáze při překladech.

V~letech 2013 až 2015 byl vytvořen morfologický klíč, který rozšiřuje klasický systém zobrazení morfologické informace v~islandských slovnících o~přesné zařazení daného hesla do určité deklinační\,/\addthin konjugační skupiny, jejíž
skloňování\,/\addthin časování je ukázáno na příkladu. V~roce 2014 jsme obdrželi roční grant od Ostravské univerzity v~Ostravě na tvorbu morfologického klíče. S~realizací grantu nám velmi pomohla doc. PhDr. Diana Svobodová, Ph.D. a~Ing. Zuzana Havrlantová.

Od března 2014 do srpna 2015 jsme realizovali druhé finální čtení, ve kterém jsme porovnávali informace s~webovým slovníkem \textit{ISLEX} (\cite {int1}). 
Na požádání jsme obdrželi od Þórdís Úlfarsdóttir (vedoucí projektu \textit{ISLEX}) seznam víceslovných příslovečných frází. Takto poskytnutý seznam nám usnadnil opravy slovníku. 
S~Þórdís Úlfarsdóttir jsme konzultovali různé gramatické jevy a~záležitosti týkající se deklinací\,/\addthin konjugací.
Martina Kašparová, studentka islandštiny na Islandské univerzitě, prošla celý slovník, přičemž se soustředila mimo jiné na stylistické a~pravopisné korektury překladů. 
Lucie Brabcová udělala jazykové a~typografické korektury celého slovníku.

Vojtěch Kupča vydal v~roce 2014 \textit{Islandskou gramatiku} (\cite {is77}) a~souhlasil s~připojením této publikace ke~knižní verzi IČSS.
V~roce 2015 jsme s~přispěním Paola al-Imk\=an\=\i\ Brasolin vysázeli slovník v~typografickém prostředí {\LaTeX}. V~únoru 2016 byl slovník vytištěn v~Tiskárně MELMEN, s.r.o v~Pardubicích a svázán 
v~Knihařství Vojtíšková Hrubá v~Pardubicích. 

\blspace[5]

{\centering Všem spoluautorům děkujeme za spolupráci.\par}


% -------------------------------------------------  Guide to the dictionary -

\chapter{Průvodce po slovníku}
\begin{multicols}{2}
\def\tableA{\noindent\begin{tabularx}{\columnwidth}{>{\bfseries}lll}
  \toprule
  \textit{přípona} & \textit{přípona s~koncovkou} & \textit{příklad} \\
  \midrule
  -ing-   & -ingur   & rugl··ingur \\
          & -ing     & hreyf··ing \\
          & -ingi    & erf··ingi \\
          & -ingar   & lækn··ingar \\
  -ling-  & -lingur  & vett··lingur \\
  -ning-  & -ningur  & flut··ningur \\
          & -ning    & set··ning \\
          & -ningar  & kos··ningar \\
  -að-    & -aður    & mark··aður \\
  -nað-   & -naður   & skil··naður \\
  -ar-    & -ari     & dóm··ari \\
  -skap-  & -skapur  & kveð··skapur \\
  -dóm-   & -dómur   & lær··dómur \\
  -leik-  & -leikur  & sann··leikur \\
          & -leiki   & mögu··leiki \\
  -lát-   & -látur   & laus··látur \\
  -átt-   & -átta    & víð··átta \\
  -sam-   & -samur   & starf··samur \\
  -sem-   & -semi    & hamingju··semi \\
  -nesk-  & -neskur  & tékk··neskur \\
          & -neska   & tékk··neska \\
  -leg-   & -legur   & blíð··legur \\
          & -lega    & blíð··lega \\
  -ung-   & -ungur   & gler··ungur \\
          & -unga    & snopp··unga \\
          & -ungi    & sveit··ungi \\
          & -ung     & nýj··ung \\
  -ótt-   & -óttur   & blett··óttur \\
  -ræn-   & -rænn    & suð··rænn \\
          & -ræna    & nor··ræna \\
  -ern-   & -erni    & fað··erni \\
  -ind-   & -indi    & vís··indi \\
  -and-   & -andi    & gef··andi \\
  -ug-    & -ugur    & göf··ugur \\
          & -ugt     & stöð··ugt \\
  -ust-   & -usta    & þjón··usta \\
  -ald-   & -aldur   & far··aldur \\
  -uð-    & -uður    & söfn··uður \\
  -úð-    & -úð      & dul··úð \\
  -versk- & -verskur & ung··verskur \\
          & -verska  & ung··verska \\
  -ist-   & -isti    & hass··isti \\
  -ism-   & -ismi    & húman··ismi \\
  -ang-   & -angur   & far··angur \\
  \bottomrule
  \end{tabularx}
  \captionof{table}{Seznam přípon označených ve slovníku}
  \label{table1}}

\def\tableB{\noindent\begin{tabularx}{\columnwidth}{>{\bfseries}lX}
  \toprule
  \textbf{symbol} & \textbf{vysvětlivka} \\
  \midrule
  m                   & označuje slovní druh \\
  \textsubscript{1}   & označuje třídu \\
  (-a, -ar)           & označuje deklinační / konjugační koncovky hesla \\
  \textsubscript{8}   & označuje podtřídu\\
  \bottomrule
  \end{tabularx}
  \captionof{table}{Morfologická informace}
  \label{table2}}

\def\tableC{\noindent\begin{tabularx}{\columnwidth}
    {>{\bfseries}l>{\footnotesize\itshape}lllll}
  \toprule
  \multirow{3}{*}{{{\textbf{f{\textsubscript{1}}} \Large{\textbf{7}}}}}
    & nom & kisa & kisan    & kisur & kisurnar \\
    & acc & kisu & kisuna   & kisur & kisurnar \\
    & dat & kisu & kisunni  & kisum & kisunum \\
    & gen & kisu & kisunnar & kisa  & kisanna \\
  \bottomrule
  \end{tabularx}
  \captionof{table}{Skloňování hesla \textit{kisa}}
  \label{table3}}

\def\tableD{\noindent\begin{tabularx}{\columnwidth}{>{\bfseries}lllll}
  \toprule
  \textit{zkratka} & \textit{slovo} & \textit{pád} & \textit{číslo}
    & \specialcell{\textit{předmět}\\\textit{podmět}} \\
  \midrule
  e-að  & eitthvað   & 1. & sg & neživotný \\
  e-ð   & eitthvað   & 4. & sg & neživotný \\
  e-ir  & einhverjir & 1. & pl & životný \\
  e-ja  & einhverja  & 4. & pl & životný \\
  e-jum & einhverjum & 3. & pl & životný \\
  e-m   & einhverjum & 3. & sg & životný \\
  e-n   & einhvern   & 4. & sg & životný \\
  e-r   & einhver    & 1. & sg & životný \\
  e-rra & einhverra  & 2. & pl & životný \\
  e-rs  & einhvers   & 2. & sg & životný \\
  e-s   & einhvers  & 2. & sg & neživotný \\
  e-u   & einhverju  & 3. & sg & neživotný \\
  \bottomrule
  \end{tabularx}
  \captionof{table}{Seznam použitých zkratek neurčitých zájmen}
  \label{table4}}

%#############################################################################

\section{Heslo}

\subsection*{Řazení slov}

Heslová slova jsou zobrazena tučně a~jsou seřazena podle islandské abecedy:

\blspace
{\bfseries\centering
  a~á b c d ð e é f g h i~í j k~l m n\\
  o~ó p q r s~t u~ú v~w x y ý z~þ æ ö\par}
\blspace

K~snadnější orientaci v~heslové části je použit u~každého písmene tzv. palcový index, který se nachází na hřbetě slovníku. Každé písmeno je v~palcovém indexu barevně odlišeno.

\subsection*{Homonyma}

Homonyma (heslová slova stejného reprezentativního tvaru) jsou označena povýšenými arabskými číslicemi uváděnými za heslovým slovem.

\blspace
  \dicEntry \dicTerm{vor\textsuperscript{1}} \dicPos{n} \dicFlx{(-s, -)}
  \dicEntry \dicTerm{vor\textsuperscript{2}} \dicPos{pron} \dicFlx{pers gen pl}
\blspace

Homonyma jsou seřazena podle českého ustáleného řazení slovních druhů. Poté následují předpony a~zkratky.
V~případě dvou podstatných jmen je pořadí mužský rod, ženský rod, střední rod. V~případě dvou podstatných jmen stejného rodu má přednost slovo se slabým skloňováním před slovem se silným skloňováním. V~případě dvou sloves stejného reprezentativního tvaru má přednost sloveso se slabým časováním před slovesem se silným časováním.

\subsection*{Varianty a~dublety}

U~některých heslových slov jsou v~záhlaví hesla uvedeny a) varianty (obměny)

\blspace
  \dicEntry \dicTerm{ærsla··fenginn}, \dicTerm{ærslafullur} \dicPos{adj}
\blspace

nebo b) dublety (dvojtvary).

\blspace
  \dicEntry \dicTerm{{ævin··týri}, \dicTerm{ævintýr}} \dicPos{n}
\blspace

Varianty a~dublety jsou méně frekventované než samotné heslo. Lze je vyhledat ve slovníku jako hesla obsahující fonetické a~morfologické informace s~odkazem na původní heslové slovo.

\blspace
  \dicEntry \dicTerm{ærsla··|fullur} \dicIPA{{a}{i}{\textsubring{r}}{s}{\textsubring{d}}{l}{a}{f}{\textscy}{\textsubring{d}}{l}{\textscy}{\textsubring{r}}} \dicPos{adj} \dicFlx{(comp -fyllri, sup -fyllstur)} \dicLink{ærslafenginn}
  \dicEntry \dicTerm{ævin··týr} \dicIPA{{a}{i}{\textlengthmark}{v}{\textsci}{n}{t\textsuperscript{h}}{i}{\textsubring{r}}} \dicPos{n} \dicFlx{(-s, -)} \dicLink{ævintýri}
\blspace

\subsection*{Vybrané tvary jako hesla}

Ve slovníku jsou uvedeny vybrané tvary heslových slov, u~kterých by mohl mít uživatel potíže s~identifikací. Jedná se o~a) genitiv jednotného čísla ({\textbf{sg}}) a~nominativ množného čísla ({\textbf{pl}}) podstatných jmen, b) ženský a~střední rod přídavných jmen a~2. a~3. stupeň rodu mužského přídavných jmen, c) oznamovací způsob ({\textbf{ind}}) přítomného času ({\textbf{praes}}) a~minulého času ({\textbf{pf}}) 1. osoby ({\textbf{pers}}) sg a~pl a~spojovací způsob ({\textbf{con}}) minulého času 1. osoby sg a~jiné nepravidelné tvary.

\blspace
  \dicEntry \dicTerm{öttum} \dicPos{v} \dicFlx{ind pf pl 1 pers} \dicLink{etja}
\blspace

\subsection*{Dělení heslových slov}

K~dělení heslových slov je použita střední tečka \textbf{'·'}, která odděluje části složeného slova, přípony a~předpony. Platí následující pravidla:

\renewcommand{\labelenumi}{\alph{enumi})}
\begin{enumerate}
\item jediné dělení heslového slova je označeno dvěma středními tečkami \textbf{'··'}. Může se jednat o~dělení částí složeného slova
\blspace
\dicEntry \dicTerm{ráð··hús}
\blspace

nebo o~dělení mezi kořenem a~předponou / příponou.

\blspace
\dicEntry \dicTerm{þjón··usta}
\blspace
\item V~případě, že se v~hesle vyskytuje více dělení, jsou použity dvě střední tečky \textbf{'··'} pro významově poslední dělení a~jedna střední tečka \textbf{'·'} pro zbylé / zbylá dělení.
\blspace
\dicEntry \dicTerm{við·bótar··líf·eyris·sparn·aður}
\blspace
\item V~hesle s~předponou je označena jedna nebo i~více předpon (každá zvlášť).
\blspace
\dicEntry \dicTerm{ó··að·finnan·legur}
\blspace
\item V~hesle s~příponou je označena jedna nebo i~více přípon (každá zvlášť).
\blspace
\dicEntry \dicTerm{létt·úð··ugur}
\blspace
\item V~hesle s~příponou je označena pouze přípona / přípony poslední části složeného slova. Důvodem pro neoznačení přípon jiných částí složených slov je lepší čitelnost hesla. Příklad označení pouze přípony poslední části složeniny:
\blspace
\dicEntry \dicTerm{flutninga··miðl·ari}
\end{enumerate}

Následuje seznam přípon, které jsou ve slovníku označeny:

\tableA

\section{Frekvence slov}

Symbolem \dicsymFrequent\ je označeno 2~825 hesel patřících mezi velmi frekventovaná islandská slova. Tento seznam vychází z~korpusu \textit{Íslenskur Orðasjóður} (\cite{int7}) a~knihy \textit{Íslensk orðtíðnibók} (\cite{is2}).

\blspace
  \dicEntry \dicTerm{koma}\textsuperscript{2} \dicsymFrequent
\blspace

\section{Fonetický zápis}

V~hranatých závorkách v~záhlaví hesla je uveden fonetický zápis výslovnosti hesla v~mezinárodní fonetické abecedě (IPA, angl. International Phonetic Alphabet).

\blspace
  \dicEntry \dicTerm{öldu··rót} \dicIPA{{\oe}{l}{\textsubring{d}}{\textscy}{r}{ou}{\textsubring{d}}}
\blspace

Znak '{\textlengthmark}' označuje dlouhou samohlásku nebo souhlásku.

\blspace
  \dicEntry \dicTerm{ör}\textsuperscript{2} \dicIPA{{\oe}{\textlengthmark}{\textsubring{r}}}
\blspace

\section{Slovní druhy, předpony a~zkratky}

V~záhlaví hesla je uvedena morfologická informace popisující a) slovní druh hesla, b) předponu nebo c) zkratku. Pokud se heslo řadí do více slovních druhů, jsou jednotlivé zkratky pro slovní druhy odděleny lomítkem.

\blspace
  \dicEntry \dicTerm{af} \dicPos{prep / adv}
\blspace

Informace o~slovním druhu, předponě či zkratce je vysázena barevně a~tučně.

\subsection*{Podstatná jména}

Podstatná jména jsou uvedena v~1. pádě (\textbf{nom}) jednotného čísla (\textbf{sg}) nebo v~1. pádě množného čísla (\textbf{pl}), pokud se heslo v~jednotném čísle nevyskytuje. Podstatná jména jsou označena následovně: \textbf{m} (l. masculinum) - podstatné jméno rodu mužského,
\textbf{f} (l. femininum) - podstatné jméno rodu ženského, \textbf{n} (l. neutrum) - podstatné jméno rodu středního, \textbf{subs} (l. substantivum) - podstatné jméno bez jednoznačného mluvnického rodu.

\blspace
  \dicEntry \dicTerm{hestur} \dicPos{m}
  \dicEntry \dicTerm{Hanoí} \dicPos{subs}
\blspace

\subsection*{Přídavná jména}

Přídavná jména jsou uvedena v~1. pádě (\textbf{nom}) rodu mužského (\textbf{m}) jednotného čísla (\textbf{sg}) 1. stupně (\textbf{pos}) (l. positivus) a~jsou označena zkratkou \textbf{adj} (l. adjectivum).

\blspace
  \dicEntry \dicTerm{fallegur} \dicPos{adj}
\blspace

V~případě, že u~přídavného jména dochází v~rodě ženském (\textbf{f}) nebo středním (\textbf{n}) k~přehlásce anebo je tvar druhého stupně (\textbf{comp} (l. comparativus)) či třetího stupně (\textbf{sup} (l. superlativus)) nepravidelný, pak jsou tyto tvary uvedeny v~nominativu jako heslová slova s~odkazem na hlavní heslový tvar

\blspace
  \dicEntry \dicTerm{yngri} \dicPos{adj} \dicFlx{comp m} \dicLink{ungur}
\blspace

nebo jako svébytná hesla s~významy.

\blspace
  \dicEntry \dicTerm{fremri} \dicPos{adj} \dicFlx{comp (sup fremstur)}
\blspace

\subsection*{Zájmena}

Zájmena jsou označena zkratkou \textbf{pron} (l. pronomen) a~jsou tříděna na zájmena ukazovací \textbf{dem} (angl. demonstrative), přivlastňovací \textbf{poss} (angl. possessive), osobní \textbf{pers} (angl. personal), neurčitá \textbf{indef} (angl. indefinite), zvratná \textbf{refl} (angl. reflexive) a~vztažná \textbf{rel} (angl. relative).

\blspace
  \dicEntry \dicTerm{þessi} \dicPos{pron} \dicFlx{dem}
\blspace

U~zájmen přivlastňovacích je v~3. osobě označen rod.

\blspace
  \dicEntry \dicTerm{hans} \dicPos{pron} \dicFlx{poss m}
\blspace

U~zájmen osobních je označen pád a~číslo.

\blspace
  \dicEntry \dicTerm{mig} \dicPos{pron} \dicFlx{pers acc sg}
\blspace

\subsection*{Číslovky}

Číslovky jsou označeny zkratkou \textbf{num} (l. numerale). V~případě vybraných číslovek (1 až 4) jsou ve slovníku uvedeny jako heslová slova také tvary ženského i~středního rodu. V~takovém případě jsou označeny zkratkou \textbf{m} pro mužský rod, zkratkou \textbf{f} pro ženský rod a~zkratkou \textbf{n} pro střední rod. Zkratkou \textbf{ord} (angl. ordinal number) jsou označeny číslovky řadové.

\blspace
  \dicEntry \dicTerm{tut··tugu} \dicPos{num}
  \dicEntry \dicTerm{tveir} \dicPos{num} \dicFlx{pl m}
  \dicEntry \dicTerm{fyrsti} \dicPos{num} \dicFlx{ord}
\blspace

\subsection*{Slovesa}

Slovesa jsou uvedena v~infinitivu a~jsou označena zkratkou \textbf{v} (l. verbum).

\blspace
  \dicEntry \dicTerm{fara} \dicPos{v}
\blspace

Slovesa ve středním rodě jsou označena jako \textbf{refl} (angl. reflexive).

\blspace
  \dicEntry \dicTerm{nálgast} \dicPos{v} \dicFlx{refl}
\blspace

Slovesa, která jsou vždy neosobní, jsou označena \textbf{impers} (angl. impersonal).

\blspace
  \dicEntry \dicTerm{svima} \dicPos{v} \dicFlx{impers}
\blspace

U~sloves je dále uvedeno, s~jakým pádem se pojí. Zkratka \textbf{acc} označuje 4. pád, \textbf{dat} 3. pád, \textbf{gen} 2. pád a~\textbf{nom} 1. pád. V~případě, že se sloveso pojí s~různými pády, je použito lomítko.

\blspace
  \dicEntry \dicTerm{klóra} \dicPos{v} \dicFlx{acc / dat}
\blspace

V~případě, že se sloveso pojí s~více pády, je použito znaménko plus.

\blspace
  \dicEntry \dicTerm{gefa} \dicPos{v} \dicFlx{dat + acc}
\blspace

\subsection*{Příslovce}

Příslovce jsou označena zkratkou \textbf{adv} (l. adverbium). Pokud chybí informace o~stupni, jedná se o~1. stupeň stupňování (\textbf{pos}). 2. stupeň je vždy označen zkratkou \textbf{comp} a~3. stupeň zkratkou \textbf{sup}.

\blspace
  \dicEntry \dicTerm{þar} \dicPos{adv}
  \dicEntry \dicTerm{ofar} \dicPos{adv} \dicFlx{comp (pos uppi, sup efst)}
  \dicEntry \dicTerm{síðast} \dicPos{adv} \dicFlx{sup (comp síðar)}
\blspace

\subsection*{Předložky}

Předložky jsou označeny zkratkou \textbf{prep} (l. praepositio). Pokud se předložka pojí výhradně s~jedním pádem, jsou použity zkratky pro 4. pád (\textbf{acc}), 3. pád (\textbf{dat}), 2. pád (\textbf{gen}).

\blspace
  \dicEntry \dicTerm{af} \dicPos{prep / adv} \dicFlx{dat}
\blspace

V~případě, že se předložka pojí s~různými pády, je použito lomítko.

\blspace
  \dicEntry \dicTerm{fyrir} \dicPos{prep / adv} \dicFlx{acc / dat}
\blspace

\subsection*{Spojky}

Spojky jsou označeny zkratkou \textbf{conj} (l. conjunctio).

\subsection*{Částice}

Částice jsou označeny zkratkou \textbf{part} (l. particula).

\subsection*{Citoslovce}

Citoslovce jsou označeny zkratkou \textbf{inter} (l. interjectio).

\subsection*{Předpony}

Předpony jsou označeny zkratkou \textbf{predp}.

\subsection*{Zkratky}

Zkratky jsou označeny zkratkou \textbf{zkr}.

\section{Skloňování, časování a~stupňování}

Každé heslo obsahuje zápis morfologické informace, která přesně informuje, jak se heslo skloňuje / časuje nebo stupňuje.

V~slovníku jsou použity dva odlišné systémy zápisu morfologické informace, které se vzájemně doplňují:

\begin{enumerate}
\item tradiční systém - do závorky za slovním druhem se zapisují deklinační / konjugační koncovky. Počet koncovek závisí na slovním druhu (viz dále).
\item numerický systém (vychází ze systému Helgi Haraldssona (\cite {is6})) - číslice odkazující na morfologický klíč, kde jsou podrobně vypsány deklinační / konjugační třídy.
Další informace o~numerickém systému lze nalézt na straně \pageref{sec:dec}.
\end{enumerate}

Deklinační / konjugační tvary nebo koncovky podstatných jmen, přídavných jmen, zájmen, číslovek, sloves a~příslovcí jsou uvedeny v~závorkách za slovním druhem.

Hnízdovací šev \textbf{'|'} označuje místo ve slově, kde se deklinační koncovka připojuje ke slovu.

\blspace
  \dicEntry \dicTerm{heils|a{\textsuperscript{1}}} \dicPos{f} \dicFlx{(-u)}
\blspace

Pokud není hnízdovací šev \textbf{'|'} uveden, tak se koncovka připojuje k~samotnému heslu.

\blspace
  \dicEntry \dicTerm{mynd} \dicPos{f} \dicFlx{(-ar, -ir)}
\blspace

Pokud dochází v~ohýbání k~přehlásce nebo je ohýbání nepravidelné, pak jsou uvedeny celé tvary hesla.

\blspace
  \dicEntry \dicTerm{maður} \dicPos{m} \dicFlx{(manns, menn)}
\blspace

V~případě, že heslové slovo má více tvarů, jsou varianty koncovky odděleny lomítkem.

\blspace
  \dicEntry \dicTerm{beð|ur} \dicPos{m} \dicFlx{(-s / -jar, -ir)}
\blspace

\subsection*{Podstatná jména}

U~podstatných jmen je uveden tvar hesla pro 2.~pád jednotného čísla a~1.~pád množného čísla.

\blspace
  \dicEntry \dicTerm{hval|ur} \dicPos{m} \dicFlx{(-s, -ir)}
\blspace

kde '\textit{hvals}' je 2.~pád jednotného čísla a~'\textit{hvalir}' 1.~pád množného čísla.
Pokud je uveden pouze jeden tvar, jedná se o~2.~pád jednotného čísla a~znamená to rovněž, že heslové slovo se neohýbá v~množném čísle.

\blspace
  \dicEntry \dicTerm{heils|a{\textsuperscript{1}}} \dicPos{f} \dicFlx{(-u)}
\blspace

Pokud je použita zkratka pro množné číslo (\textbf{pl}), znamená to, že se heslo ohýbá pouze v~množném čísle.

\blspace
  \dicEntry \dicTerm{afar··kostir} \dicPos{m} \dicFlx{pl}
\blspace

Pokud je použita zkratka \textbf{indecl}, znamená to, že podstatné jméno je nesklonné.

\blspace
  \dicEntry \dicTerm{október} \dicPos{m} \dicFlx{indecl}
\blspace

\subsection*{Přídavná jména}

U~přídavných jmen jsou uvedeny v~závorkách tvary, které poukazují na nepravidelnosti ve skloňování.

Pokud přídavné jméno v~ženském rodě neztrácí koncové '\textit{-r}', je uveden tvar hesla v~ženském rodě.

\blspace
  \dicEntry \dicTerm{hýr} \dicPos{adj} \dicFlx{(f -)}
\blspace

Dále jsou uvedeny tvary ženského rodu, ve~kterých dochází k~přehlásce:

\blspace
  \dicEntry \dicTerm{reglu··|samur} \dicPos{adj} \dicFlx{(f -söm)}
\blspace

Pokud je přídavné jméno nesklonné, je uvedena zkratka \textbf{indecl}.

\blspace
  \dicEntry \dicTerm{tví··mála} \dicPos{adj} \dicFlx{indecl}
\blspace

Pokud má přídavné jméno nepravidelné tvary 2.~a~3.~stupně stupňování, jsou tyto tvary uvedeny v~závorce spolu se zkratkou pro 2.~stupeň (\textbf{comp}) a~3.~stupeň (\textbf{sup}).

\blspace
  \dicEntry \dicTerm{þunnur} \dicPos{adj} \dicFlx{(comp þynnri, sup þynnstur)}
\blspace

V~případě, že je uveden pouze jeden ze dvou zbývajících stupňů, znamená to, že heslo se stupňuje pouze v~uvedených stupních.

\blspace
  \dicEntry \dicTerm{syðri} \dicPos{adj} \dicFlx{comp (sup syðstur)}
\blspace

\subsection*{Slovesa}

Islandská slovesa se dělí na slabá, silná a~nepravidelná. Vzhledem k~tomuto dělení se liší počet slovních tvarů uvedených v~závorce za zkratkou \textbf{v}. Princip uvádění slovních tvarů spočívá v~tom, že uživatel
je schopen (více méně) z~těchto tvarů odvodit časování slovesa. V~tomto slovníku jsme
záměrně umístili některá slabá slovesa mezi silná (místo dvou slovesných tvarů je uvedeno pět tvarů) kvůli tomu, aby byl uživatel přesněji informován o~časování slovesa.

Slabá slovesa, která tvoří 1.~osobu jednotného čísla minulého času oznamovacího způsobu koncovkou '\textit{-aði}', mají v~závorce uvedenu jen jednu koncovku, jmenovitě '\textit{(-aði)}'.

\blspace
  \dicEntry \dicTerm{ætl|a} \dicPos{v} \dicFlx{(-aði)}
\blspace

U~zbývajících skupin slabých sloves se v~závorce nachází koncovka 1.~osoby jednotného čísla minulého času oznamovacího způsobu a~supinum.

\blspace
  \dicEntry \dicTerm{kenn|a} \dicPos{v} \dicFlx{(-di, -t)}
\blspace

Silná nebo nepravidelná slovesa mají v~závorce vždy pět tvarů, jmenovitě 1.~osobu jednotného čísla přítomného času oznamovacího způsobu, 1.~osobu jednotného čísla minulého času oznamovacího způsobu, 1.~osobu množného čísla minulého času oznamovacího způsobu, 1.~osobu jednotného čísla minulého času spojovacího způsobu a~supinum.

\blspace
  \dicEntry \dicTerm{grípa} \dicPos{v} \dicFlx{(gríp, greip, gripum, gripi, gripið)}
\blspace

Silná nebo nepravidelná slovesa vyskytující se pouze v~neosobním užití mají v~závorce 4 tvary, jmenovitě 3.~osobu jednotného čísla přítomného času oznamovacího způsobu, 3.~osobu jednotného čísla minulého času oznamovacího způsobu, 3.~osobu jednotného čísla minulého času spojovacího způsobu a~supinum.

\blspace
  \dicEntry \dicTerm{kross··|bregða} \dicPos{v} \dicFlx{(-bregður, -brá, -brygði, -brugðið)}
\blspace

U~některých významů v~heslovém slově se vyskytují gramatické informace, které popisují chování daného významu.

Zkratka \textbf{impers} označuje, že se daný význam slovesa vyskytuje v~neosobním užití. Pomocí slovního spojení je vyjádřeno užití slovesa.

\blspace
  \dicEntry \dicTerm{batn|a} \dicPos{v} \dicFlx{(-aði)} \dicTerm{e-m batnar} \dicFlx{impers} (kdo) se uzdravuje, (kdo) se zotavuje, (komu) je lépe
\blspace

Zkratka \textbf{refl} označuje mediopasivní tvary slovesa.

\blspace
  \dicEntry \dicTerm{að··lag|a} \dicPos{v} \dicFlx{(-aði)} \dicIdiom{aðlagast} {\dicPhraseIS{aðlagast e-u}} \dicFlx{refl} {přizpůsobit se (čemu)}
\blspace

\subsection*{Příslovce}

U~příslovcí se v~záhlaví nachází tvary stupňování. Tyto tvary jsou uvedeny v~závorce spolu se zkratkou pro 2.~stupeň (\textbf{comp}) a~3.~stupeň (\textbf{sup}).

\blspace
  \dicEntry \dicTerm{vel\textsuperscript{2}} \dicPos{adv} \dicFlx{(comp betur, sup best)}
\blspace

Nepravidelné tvary 2.~a~3.~stupně jsou uvedeny buď jako hesla odkazující na základní tvar v~1.~stupni (\textbf{pos}), anebo jako svébytná hesla. V~obou případech jsou uvedeny v~závorce zbývající stupně.

\blspace
  \dicEntry \dicTerm{síst} \dicPos{adv} \dicFlx{sup (pos varla, comp síður)}
\blspace

Pokud některý ze stupňů uveden není, znamená to, že pro dané slovo neexistuje.

\blspace
  \dicEntry \dicTerm{síðar} \dicPos{adv} \dicFlx{comp (sup síðast)}
\blspace

\section{Morfologická informace} \label{sec:dec}

V~záhlaví hesla se kromě deklinačních / konjugačních koncovek nachází také morfologické numerické informace, které odkazují na morfologický klíč (viz strana \pageref{sec:morpho}), kde jsou podrobně vypsány kompletní deklinační / konjugační tabulky.

\blspace
  \dicEntry \dicTerm{ap|i} \dicPos{m}[1] \dicFlx{(-a, -ar)}[8]
\blspace

Morfologická informace je zapsána podle následujícího schématu: \textbf{m{\textsubscript{1}} (-a, -ar){\textsubscript{8}}}

\blspace
\tableB
\blspace

Deklinační / konjugační třídy a~jejich počet nejsou v~islandštině ustálené a~slouží pouze pro rozlišení odlišných skloňování / časování.

Názvy podtříd slouží pro rozlišení odlišných skloňování / časování.

\subsection*{Práce s~morfologickým klíčem}

Pro vyhledání kompletních deklinačních / konjugačních tabulek slouží následující postup, který demonstrujeme na ukázce:

\blspace
  \dicEntry \dicTerm{kis|a} \dicPos{f}[1] \dicFlx{(-u, -ur)}[7]
\blspace

Nejdříve určíme slovní druh (v~případě podstatných jmen také rod), v~tomto případě podstatné jméno rodu ženského, a~navštívíme Morfologický klíč (str. \pageref{sec:morpho}) pro podstatná jména rodu ženského (str. \pageref{sec:morpho_f}).
Následně dohledáme třídu (v~tomto případě číslo {\textsubscript{\textbf{1}}}) a~následně dohledáme dodatečné numerické informace za závorkou (v~tomto případě číslo {\textsubscript{\textbf{7}}}).

\tableC

\section{Definice}

\subsection*{Základní tvar}

Základním tvarem definice je, že jednomu islandskému slovu odpovídá jeden nebo více českých významů.

\blspace
  \dicEntry \dicTerm{landa··fræð|i} \dicPos{f} \dicFlx{(-i)} \textCS{zeměpis, geografie}
\blspace

V~některých případech je v~české části použita závorka. V~prvním případě je závorka užívaná ke~zkrácení zápisu dvou slov ('\textit{pomoc}', '\textit{výpomoc}').

\blspace
  \dicEntry \dicTerm{að··stoð} \dicPos{f} \dicFlx{(-ar)} \textCS{(vý)pomoc, asistence}
\blspace

V~druhém případě se v~závorce nachází upřesňující informace, která zabraňuje dvojznačnosti.

\blspace
  \dicEntry \dicTerm{fíkj|a} \dicPos{f} \dicFlx{(-u, -ur)} \textCS{fík} {\footnotesize {\textCS{(plod)}}}
\blspace

V~mnoha případech je závorka v~české části použita pro příklad užití českého slova a~tím vymezení jeho významu.

\blspace
  \dicEntry \dicTerm{að·gengi··legur} \dicPos{adj} {přístupný, dostupný (informace ap.)}
\blspace

Lomítko \textbf{'/'} slouží v~české části významu k~oddělení více významů pojících se s~jedním nebo více slovy a~je použito v~české části následujícím způsobem:

\blspace
  \dicEntry \dicTerm{aðal··|gata} \dicPos{f} \dicFlx{(-götu, -götur)} \textCS{hlavní třída / ulice}
\blspace

Význam přečteme jako '\textit{hlavní třída, hlavní ulice}'.

\subsection*{Opis heslového slova}

V~případech, kdy islandské slovo nemá ekvivalent v~českém jazyce, jsou významy islandských slov opsány delším opisem a~zobrazeny menším písmem. Příkladem jsou slova z~lidových pověstí, kulinářské speciality nebo botanické pojmy.

\blspace
  \dicEntry \dicTerm{til··ber|i} \dicPos{m} \dicFlx{(-a, -ar)} {\footnotesize{pov.}} {\footnotesize{stvoření, které dojí krávy a~ovce hospodářů}}
\blspace

\subsection*{Slovní spojení}

V~definici hesla jsou kromě samotných významů uváděna také slovní spojení, která jsou sázena tučně. Po slovním spojení následuje překlad slovního spojení.

\blspace
  \dicEntry \dicTerm{lær|a} \dicPos{v} \dicFlx{(-ði, -t)} \textbf{læra e-ð utanbókar} {učit se (co) zpaměti, memorovat (co)}
\blspace

Závorky jsou použity na tu část slovního spojení, kterou je možné vypustit.

\blspace
  \dicEntry \dicTerm{ár\textsuperscript{1}} \dicPos{f} \dicFlx{(-ar, -ar)} \textbf{taka (of) djúpt í árinni} {\footnotesize{přen.}} {přehnat (co v~tvrzení)}
\blspace

Příklad můžeme přečíst jako '\textit{taka djúpt í árinni}' nebo '\textit{taka of djúpt í árinni}'.

Mezi slovní spojení v~tomto slovníku řadíme a) tvary heslového slova, b) příkladová spojení, c) ustálená slovní spojení a~d) rčení nebo přísloví.

\begin{enumerate}
\item Tvary heslového slova\\
Do této skupiny řadíme tvar množného čísla hesla (zkratka \textbf{pl}).
\blspace
\dicEntry \dicTerm{bót} \dicPos{f} \dicFlx{(-ar, bætur)} \textbf{bætur} \dicFlx{pl} {dávky, přídavky, příspěvky}\\

Dále do této skupiny řadíme tvar slovesa v~středním slovesném rodě (zkratka \textbf{refl}).\\

\dicEntry \dicTerm{finna} \dicPos{v} \dicFlx{(finn, fann, fundum, fyndi, fundið)} \dicIdiom{finnast} {\dicPhraseIS{finnast}} \dicFlx{refl} {být nalezen, najít se, nalézt se, nacházet se}
\blspace
\item Příkladová spojení\\
Příkladová spojení ilustrují komunikační chování heslových slov v~příslušných významech a~v~různých syntagmatech (např. sloveso s~určitou předložkou ap.). Jsou označena tučně modrou barvou a~odsazená pro snadnou orientaci jako svébytné heslo.
\blspace
\dicEntry \dicTerm{koma\textsuperscript{2}} \dicPos{v} \dicFlx{(kem, kom, komum, kæmi, komið)} \dicIdiom{koma}[að] {\dicPhraseIS{koma að e-m}} {přiblížit se ke (komu), přijít ke (komu)}
\blspace
\item Ustálená slovní spojení\\
Tento druh slovních spojení je charakterizován většinou přeneseným významem. Tato spojení jsou uváděna za posledním významem nebo příkladovým spojením a~jsou označená tučně modrou barvou a~odsazená pro snadnou orientaci jako svébytné heslo.
\blspace
\dicEntry \dicTerm{kast} \dicPos{n} \dicFlx{(-s, köst)} \dicIdiom{kast} {\dicPhraseIS{þegar til kastanna kemur}} {\footnotesize{přen.}} {když jde do tuhého}
\blspace
\item Rčení a~přísloví\\
Rčení a~přísloví se nacházejí vždy na konci definice a~jsou odlišená symbolem \dicsymProverb\ a~odsazená pro snadnou orientaci.
\blspace
\dicEntry \dicTerm{æfing} \dicPos{f} \dicFlx{(-ar, -ar)} \dicProverb{} \dicPhraseIS{Æfing skapar meistara.} {\footnotesize{přís.}} \textCS{Cvičení dělá mistra.}
\end{enumerate}

\section{Vícevýznamovost}

\subsection*{Vícevýznamová hesla}

Značné množství hesel má více než jeden význam. Významy jsou odlišeny arabskými číslicemi a~v~některých případech i~islandskými synonymy.

\blspace
  \dicEntry \dicTerm{land} \dicPos{n} \dicFlx{(-s, lönd)} {\textbf{1.}} \dicSynonym{þurrlendi} {souš, pevnina, země} {\textbf{2.}} \dicSynonym{árbakki} {břeh} {\textbf{3.}} \dicSynonym{ríki} {země, stát, vlast} {\textbf{4.}} \dicSynonym{landareign} {země, půda, pozemek}
\blspace

\subsection*{Vícevýznamová slovní spojení}

V~případě, že slovní spojení uvnitř definice má více významů, jsou jednotlivé významy odlišeny písmeny {\textbf{a., b., c.}} atd.

\blspace
  \dicEntry \dicTerm{drag} {\dicPos{n} \dicFlx{(-s, drög)} {\textbf{1.}} {mokřina, podmoklý terén;} {\textbf{2.}} \textbf{drög} {\footnotesize{pl}} {\textbf{a.}} \dicSynonym{uppsprettur} {prameny (řeky ap.)} {\textbf{b.}} \dicSynonym{undirbúningur} {náčrt, návrh.}}
\blspace

\section{Řazení významů a~slovních spojení v~definici}
\longsectionskip

V~definici jsou nejdříve seřazeny významy slova a~odlišeny arabskými číslicemi.

\subsection*{V rámci významu}

Slovní spojení jsou uváděna v~rámci jednotlivého významu, pokud se s~ním významově vážou.

\blspace
  \dicEntry \dicTerm{tím|i} \dicPos{m} \dicFlx{(-a, -ar)} {\textbf{1.}} \dicSynonym{tíð} {čas, doba}; \textbf{í þann tíma} \dicFlx{adv} \dicSynonym{þá\textsuperscript{2}} {pak} {\textbf{2.}} \dicSynonym{rétt stund} {(správný) čas / moment}
\blspace

\subsection*{Za posledním očíslovaným významem}

Pokud se slovní spojení neváže k~žádnému významu, je řazeno za posledním očíslovaným významem. Ustálená slovní spojení jsou uvedena symbolem \dicsymIdiom.

\blspace
  \dicEntry \dicTerm{finna} \dicPos{v} \dicFlx{(finn, fann, fundum, fyndi, fundið)} {\textbf{3.}} \dicSynonym{skynja} {cítit, vnímat} \dicIdiom{finna}[að] {\dicPhraseIS{finna að e-u}} {kritizovat (co), nacházet chyby na (čem)}
\blspace

\subsection*{Způsob řazení slovních spojení}

Slovní spojení, která se nachází za posledním očíslovaným významem, jsou řazena podle předložky / příslovce / zájmena, se kterou se pojí. Pro snadnou orientaci v~rozsáhlých definicích je uvedeno heslové slovo nebo tvar heslového slova plus předložka / příslovce / zájmeno a~barevně odlišeno modrou barvou.

\blspace
  \dicEntry \dicTerm{koma\textsuperscript{2}} \dicPos{v} \dicFlx{(kem, kom, komum, kæmi, komið)} \dicIdiom{koma}[út] {\dicPhraseIS{koma út}} {vyjít (kniha ap.), objevit se (na trhu ap.)}
\blspace

Poté následují slovní spojení v~slovesném středním rodě řazená opět podle předložky / příslovce / zájmena, se kterou se pojí:

\blspace
  \dicEntry \dicTerm{koma\textsuperscript{2}} \dicPos{v} \dicFlx{(kem, kom, komum, kæmi, komið)} \dicIdiom{komast} {\dicPhraseIS{komast}} \dicFlx{refl} {dostat se, dostávat se}
\blspace

Poté následují ustálená slovní spojení, která nebylo možné zařadit do předchozích položek:

\blspace
  \dicEntry \dicTerm{koma\textsuperscript{2}} \dicPos{v} \dicFlx{(kem, kom, komum, kæmi, komið)} \dicIdiom{koma} {\dicPhraseIS{(komdu) sæll (og blessaður)}} {ahoj, nazdar, dobrý den}
\blspace

Nakonec jsou uvedena přísloví a~rčení označená symbolem \dicsymProverb\:

\blspace
  \dicEntry \dicTerm{eik} \dicPos{n} \dicFlx{(-ar / -ur, -ur)} \dicProverb{} \dicPhraseIS{Ekki / eigi fellur eik við fyrsta högg.} {\footnotesize{přís.}} {Jednou ranou dub nespadne.}
\blspace

\section{Synonyma a~antonyma / hesla ke~srovnání}
\longsectionskip

\subsection*{Synonyma}

Synonyma nebo synonymní slovní spojení jsou uvedeny před opisem významu hesla. Jsou umístěna v~závorkách. Islandská synonyma slouží a) islandskému uživateli k~rozlišení českých významů b) českému uživateli jako dodatečná informace o~významu slova.

\blspace
  \dicEntry \dicTerm{staða} \dicPos{f} \dicFlx{(stöðu, stöður)} {\textbf{1.}} \dicSynonym{ástand} {situace, stav, poměry} {\textbf{2.}} \dicSynonym{embætti} {post, (pracovní) místo / pozice (ve firmě ap.)}
\blspace

V~případě, že synonymum lze nalézt jako heslo ve slovníku, je použita šipka '\dicsymSee\ '.

\blspace
  \dicEntry \dicTerm{bygging} {\dicPos{f} \dicFlx{(-ar, -ar)} {\textbf{1.}} \dicSynonym{hús} {budova, stavba, stavení} {\textbf{2.}} \textit{(það að byggja)} {(vý)stavba} }
\blspace

\subsection*{Antonyma / Hesla ke~srovnání}

Antonyma či hesla ke~srovnání jsou uvedena za opisem významu heslového slova a~jsou označena šipkou '$\shortuparrow$'.

\blspace
  \dicEntry \dicTerm{kaldur} \dicPos{adj} \dicFlx{(f köld)} {\textbf{1.}} {studený, chladný} \textit{($\shortuparrow$ heitur)}
\blspace

\section{Oborové a~stylové charakteristiky}

K~upřesnění významu heslového slova slouží oborové a~stylové charakteristiky.

\subsection*{Oborové charakteristiky}

Oborové charakteristiky zařazují islandské slovo do určitého oboru a~naznačují použití daného slova v~jazyce

\blspace
  \dicEntry \dicTerm{berg} \dicPos{n} \dicFlx{(-s, -)} {\textbf{1.}} {\footnotesize{geol.}} {hornina, kámen} {\textbf{2.}} \dicSynonym{klöpp\textsuperscript{1}} {skalní stěna, skalisko}
\blspace

a~vymezují význam českého slova.

\blspace
  \dicEntry \dicTerm{mús} \dicPos{f} \dicFlx{(-ar, mýs)} {\textbf{1.}} {\footnotesize{zool.}} {myš} (\textit{l. Mus}) {\textbf{2.}} {\footnotesize{poč.}} {myš}
\blspace

\subsection*{Stylové charakteristiky}

Stylové charakteristiky jsou uváděny v~případě, kdy význam heslového slova je možné použít pouze při určité příležitosti nebo je význam jazykově zabarven.

\blspace
  \dicEntry \dicTerm{gaur} \dicPos{m} \dicFlx{(-s, -ar)} {\textbf{1.}} {\footnotesize{hovor.}} {chlápek, maník, člověk}
\blspace

\section{Syntax}

Syntaktické informace popisují:

\begin{enumerate}
\item zdali je sloveso tranzitivní nebo intranzitivní a~s~jakým(i) pádem (pády) se pojí, pokud je tranzitivní,
\item zdali je podmět / předmět životný či neživotný,
\item zdali je sloveso osobní či neosobní a~v~jakém pádu je podmět, pokud je sloveso neosobní,
\item zdali je podmět gramatický,
\item zdali je sloveso použito ve středním rodě,
\item nejčastější příklady, které se pojí se slovesem,
\item se kterými příslovci nebo předložkami se sloveso pojí.
\end{enumerate}

Syntaktické informace se ve slovníku vyskytují ve třech formách: a) informace uvedené zkratkou, b) teoretické příklady a~c) praktické příklady.

\subsection*{Syntaktické informace uvedené zkratkou}

Zde máme na mysli informace o~pádě (\textbf{nom, acc, dat, gen}), slovesu v~středním slovesném rodě (\textbf{refl}), neosobním slovesu (\textbf{impers}), čísle (\textbf{sg, pl}), rozkazovacím způsobu (\textbf{imper}) a~spojovacím způsobu (\textbf{con}).

\subsection*{Teoretické příklady}

Teoretickým příkladem nazýváme příklad, kdy podmět nebo předmět je zastoupen zájmenem neurčitým v~pádě, se~kterým se slovo pojí a~sloveso je zpravidla v~infinitivu. V~případě, že podmět není v~1.~pádě (např. u~sloves neosobních), je sloveso uvedeno v~3.~osobě jednotného čísla přítomného času.
Ve slovních spojeních používáme zkratky pro tvary neurčitého zájmena \textbf{einhver} (někdo) a~\textbf{eitthvað} (něco). Tyto zkratky označují a) pád, b) číslo a~c) životnost nebo neživotnost podmětu / předmětu.\\
Následuje seznam zkratek neurčitých zájmen použitých ve slovníku: \\

\tableD
\blspace
Příklad použití zkratek s~neurčitými zájmeny:

\blspace
  \dicEntry \dicTerm{leit|a} \dicPos{v} \dicFlx{(-aði)} \textbf{leita að e-u / e-m} {hledat (co / koho)} \dicExampleIS{leita að lyklunum} \dicExampleCS{hledat klíče}
\blspace

Příklad tedy přečteme jako '\textit{leita að einhverju}' (\textit{hledat co}) a~'\textit{leita að einhverjum}' (\textit{hledat koho}). V~tomto případě se sloveso pojí s~3.~pádem a~předmět může být životný i~neživotný.

\subsection*{Praktické příklady}

Ve slovníku se vyskytuje velké množství praktických příkladů a~jejich překladů. Praktický příklad se liší od příkladu teoretického tím, že místo zájmen neurčitých jsou uvedena frekventovaná slova z~běžného projevu. Praktické příklady slouží k~ilustraci použití
heslového slova.\\
Praktickým příkladem může být a) (složené) slovo, b) slovní spojení nebo c) věta.

\begin{enumerate}
\item Slovo jako praktický příklad:
\blspace
\dicEntry \dicTerm{hett|a} \dicPos{f} \dicFlx{(-u, -ur)} \textbf{1.} {kapuce, kápě} \dicExampleIS{hettuúlpa} \dicExampleCS{bunda s~kapucí}
\blspace
\item Slovní spojení jako praktický příklad:
\blspace
\dicEntry \dicTerm{tím|i} \dicPos{m} \dicFlx{(-a, -ar)} \textbf{3.} \dicSynonym{klukkustund} {hodina} {\footnotesize{(šedesát minut)}} \dicExampleIS{tveggja tíma gangur} \dicExampleCS{dvouhodinový pochod}
\blspace
\item Věta jako praktický příklad:
\blspace
\dicEntry \dicTerm{mæt|a} \dicPos{v} \dicFlx{(-ti, -t)} \dicIdiom{mætast} \dicPhraseIS{mætast} \dicSynonym*{ná saman} \dicFlx{refl} {setkat se, potkat se} \dicExampleIS{Þau mættust á gangstéttinni.} \dicExampleCS{Potkali se na chodníku.}
\end{enumerate}

\section{Fotografie a~ilustrace}

Ve slovníku je publikováno 630 fotografií rostlin, živočichů a~fotografií, která dokumentují místopisně heslová slova (např. kulinářské speciality).
Fotografie jsou publikovány pod veřejnými licencemi. Jména autorů a~názvy licencí lze nalézt v~Seznamu autorů fotografií a~ilustrací na straně \pageref{sec:photo}. U~hesla a~konkrétně u~daného významu, ke kterému se fotografie vztahuje, je použit symbol fotoaparátu \dicsymPhoto. Fotografie je umístěna buď pod heslem, anebo v~horní části stránky, kde se heslo nachází, anebo na stránce (stránkách) následující(ch).

%\par\begin{center}\setlength\fboxsep{0pt}\setlength\fboxrule{0.5pt}
%\fbox{\includegraphics[width=6cm]{ds_image_posthus_0_2.jpg}}\end{center}
%\par\begin{center}\footnotesize {Pósthús}\end{center}

\end{multicols}

% ---------------------------------------- List of abbreviations and symbols -

\chapter{Seznam použitých zkratek a symbolů}
\begin{multicols}{2}
\begin{description}[leftmargin=!,labelwidth=1.5cm]
%[\breaklabel\setleftmargin{40pt}\setlabelstyle{\bfseries}
%\addtolength{\itemsep}{-0.5\baselineskip}]
\item[{abb}] {zkrácený}
\item[{acc}] {4. pád, akuzativ}
\item[{adj}] {přídavné jméno}
\item[{adv}] {příslovce}
\item[{akt}] {rod činný, aktivum}
\item[{anat.}] {anatomie}
\item[{angl.}] {anglicky}
\item[{astro.}] {astronomie, výzkum vesmíru}
\item[{ap.}] {a podobně}
\item[{básn.}] {básnický, poetický výraz}
\item[{biol.}] {biologie}
\item[{bot.}] {botanika}

\item[{comp}] {2. stupeň stupňování, komparativ}
\item[{con}] {spojovací způsob, konjunktiv}
\item[{conj}] {spojka}
\item[{dat}] {3. pád, dativ}
\item[{def}] {určitý tvar}
\item[{dem}] {zájmeno ukazovací}
\item[{dět.}] {dětsky}

\item[{e-að}] {co (sg nom neživotný) (\textit{eitthvað})}
\item[{e-ð}] {co (sg acc neživotný) (\textit{eitthvað})}
\item[{e-ir}] {kdo (pl nom životný) (\textit{einhverjir})}
\item[{e-ja}] {koho (pl acc životný) (\textit{einhverja})}
\item[{e-jum}] {komu (pl dat životný) (\textit{einhverjum})}
\item[{e-m}] {komu (sg dat životný) (\textit{einhverjum})}
\item[{e-n}] {koho (sg acc životný) (\textit{einhvern})}
\item[{e-r}] {kdo (sg nom životný) (\textit{einhver})}
\item[{e-rra}] {koho (pl gen životný) (\textit{einhverra})}
\item[{e-rs}] {koho (sg gen životný) (\textit{einhvers})}
\item[{e-s}] {čeho (sg gen neživotný) (\textit{einhvers})}
\item[{e-u}] {čemu (sg dat neživotný) (\textit{einhverju})}

\item[{ekon.}] {ekonomika, obchod}
\item[{elek.}] {elektřina}
\item[{f}] {rod ženský}
\item[{filos.}] {filosofie, logika}
\item[{form.}] {formálně}
\item[{fyz.}] {fyzika}
\item[{gen}] {2. pád, genitiv}
\item[{geog.}] {zeměpis, geografie}
\item[{geol.}] {geologie}
\item[{han.}] {hanlivě, pejorativně}
\item[{hist.}] {historie}
\item[{hovor.}] {hovorově}
\item[{hrub.}] {hrubě}
\item[{hud.}] {hudba}
\item[{imper}] {rozkazovací způsob, imperativ}
\item[{impers}] {sloveso neosobní}
\item[{ind}] {oznamovací způsob, indikativ}
\item[{indecl}] {nesklonný}
\item[{indef}] {neurčitý tvar, (pron +) zájmeno neurčité}
\item[{inf}] {infinitiv}
\item[{int}] {tázací zájmeno}
\item[{inter}] {citoslovce}
\item[{chem.}] {chemie}
\item[{jaz.}] {jazykověda}
\item[{kulin.}] {kulinářství, vaření}
\item[{l.}] {latinsky}
\item[{let.}] {letectví}
\item[{lit.}] {literatura, vydavatelství}
\item[{m}] {rod mužský}
\item[{mat.}] {matematika}
\item[{med. }] {lékařství}
\item[{meteo.}] {meteorologie}
\item[{myt.}] {mytologie}
\item[{n}] {rod střední}
\item[{náb.}] {náboženství}
\item[{nám.}] {námořnictví, rybolov}
\item[{nom}] {1. pád, nominativ}
\item[{num}] {číslovka}
\item[{ord}] {řadová číslovka}
\item[{p}] {osoba}
\item[{part}] {částice}
\item[{pers}] {(v +) osoba, (pron +) zájmeno osobní}
\item[{pf}] {minulý čas}
\item[{pl}] {množné číslo, plurál}
\item[{poč.}] {informatika}
\item[{pol.}] {politika, politologie}
\item[{poss}] {zájmeno přivlastňovací}
\item[{pos}] {1. stupeň stupňování, pozitiv}
\item[{pov.}] {lidové pověsti, folkloristika}
\item[{pp}] {příčestí minulé}
\item[{praes}] {přítomný čas}
\item[{práv.}] {právnictví, soudnictví}
\item[{predp}] {předpona}
\item[{prep}] {předložka}
\item[{presp}] {příčestí přítomné}
\item[{pron}] {zájmeno}
\item[{prop}] {vlastní jméno, proprium}
\item[{přen.}] {přeneseně}
\item[{přís.}] {přísloví}
\item[{psych.}] {psychologie}
\item[{refl}] {(v +) médium, (pron +) zvratné zájmeno}
\item[{rel}] {vztažné zájmeno}
\item[{sg}] {jednotné číslo, singulár}
\item[{slang.}] {slang}
\item[{sport.}] {sport}
\item[{stav.}] {stavebnictví, architektura}
\item[{subs}] {podstatné jméno (bez rodu)}
\item[{sup}] {3. stupeň stupňování, superlativ}
\item[{supin}] {supinum}
\item[{škol.}] {školství}
\item[{techn.}] {technika, mechanika}
\item[{v}] {sloveso}
\item[{voj.}] {vojenství}
\item[{zast.}] {zastarale}
\item[{zkr}] {zkratka}
\item[{zool.}] {zoologie}

\item[{;}] {středník uprostřed definice odděluje věty; středník na konci definice odděluje významy}

\item[{(...)}] {v kulatých závorkách jsou uvedeny doplňkové informace}
\item[{[...]}] {v hranatých závorkách je uveden fonetický zápis výslovnosti}\footnote{Vysvětlení fonetických symbolů je uvedeno v kapitole Fonetika \ref{sec:phon_phonems} na straně \pageref{sec:phon_phonems}}
\item[{|}] {naznačuje místo, kde se deklinační koncovka připojuje ke slovnímu základu}
\item[{·}] {vedlejší dělení složeného slova}
\item[{··}] {hlavní dělení složeného slova}
\item[{+}] {plus}
\item[{/}] {lomítko označuje více možností}
\item[{\dicsymSee}] {šipka odkazuje na heslo ve slovníku}
\item[{\dicsymCompare}] {šipka odkazuje na heslo ve slovníku ke srovnání}
\item[{\dicsymExampleIS}] {uvozuje islandský příklad}
\item[{\dicsymExampleCS}] {uvozuje český překlad příkladu}
\item[{\dicsymIdiom}] {uvozuje slovní spojení}
\item[{\dicsymFrequent}] {heslo patří mezi frekventované výrazy}

\item[{\dicsymPhoto}] {fotografie, ilustrace}
\item[{\dicsymProverb}] {přísloví, rčení}

\end{description}
\end{multicols}

\chapter{Skrá yfir skammstafanir og tákn}
\begin{multicols}{2}
\begin{description}[leftmargin=!,labelwidth=1.5cm]
%[\breaklabel\setleftmargin{40pt}\setlabelstyle{\bfseries}
%\addtolength{\itemsep}{-0.5\baselineskip}]
\item[{abb}] {stýfður}
\item[{acc}] {þolfall}
\item[{adj}] {lýsingarorð}
\item[{adv}] {atviksorð}
\item[{akt}] {germynd}
\item[{anat.}] {líffærafræði}
\item[{angl.}] {enska}
\item[{astro.}] {geimvísindi, stjörnufræði}
\item[{ap.}] {e.þ.h.}
\item[{básn.}] {ljóðrænn}
\item[{biol.}] {líffræði}
\item[{bot.}] {grasafræði}

\item[{comp}] {miðstig}
\item[{con}] {viðtengingarháttur}
\item[{conj}] {samtenging}
\item[{dat}] {þágufall}
\item[{def}] {með greini}
\item[{dem}] {ábendingarfornafn}
\item[{dět.}] {barnamál}

\item[{e-að}] {(\textit{eitthvað})}
\item[{e-ð}] {(\textit{eitthvað})}
\item[{e-ir}] {(\textit{einhverjir})}
\item[{e-ja}] {(\textit{einhverja})}
\item[{e-jum}] {(\textit{einhverjum})}
\item[{e-m}] {(\textit{einhverjum})}
\item[{e-n}] {(\textit{einhvern})}
\item[{e-r}] { (\textit{einhver})}
\item[{e-rra}] {(\textit{einhverra})}
\item[{e-rs}] {(\textit{einhvers})}
\item[{e-s}] {(\textit{einhvers})}
\item[{e-u}] { (\textit{einhverju})}

\item[{ekon.}] {hagfræði, viðskipti}
\item[{elek.}] {rafmagn}
\item[{f}] {kvenkyn}
\item[{filos.}] {heimspeki, rökfræði}
\item[{form.}] {formlegt mál}
\item[{fyz.}] {eðlisfræði}
\item[{gen}] {eignarfall}
\item[{geog.}] {landafræði}
\item[{geol.}] {jarðfræði}
\item[{han.}] {niðrandi}
\item[{hist.}] {sagnfræði}
\item[{hovor.}] {talmál}
\item[{hrub.}] {gróft}
\item[{hud.}] {tónlist}
\item[{imper}] {boðháttur}
\item[{impers}] {ópersónuleg sögn}
\item[{ind}] {framsöguháttur}
\item[{indecl}] {óbeygjanlegur}
\item[{indef}] {óákveðinn, (pron +) óákveðið fornafn}
\item[{inf}] {nafnháttur}
\item[{int}] {spurnigarfornafn}
\item[{inter}] {upphrópun}
\item[{chem.}] {efnafræði}
\item[{jaz.}] {málfræði, málvísindi}
\item[{kulin.}] {eldhús}
\item[{l.}] {latína}
\item[{let.}] {flug, flugmál}
\item[{lit.}] {skáldskapur}
\item[{m}] {nafnorð, karlkyn}
\item[{mat.}] {stærðfræði}
\item[{med. }] {læknisfræði}
\item[{meteo.}] {veðurfræði}
\item[{myt.}] {goðafræði}
\item[{n}] {nafnorð, hvorugkyn}
\item[{náb.}] {trúfræði}
\item[{nám.}] {sjómennska}
\item[{nom}] {nefnifall}
\item[{num}] {töluorð}
\item[{ord}] {raðtala}
\item[{p}] {persóna}
\item[{part}] {ögn}
\item[{pers}] {(v +) persóna, (pron +) persónufornafn}
\item[{pf}] {þátíð}
\item[{pl}] {fleirtala}
\item[{poč.}] {tölvufræði, tölvur}
\item[{pol.}] {stjórnmál, stjórnmálafræði}
\item[{poss}] {eignarfornafn}
\item[{pos}] {frumstig}
\item[{pov.}] {þjóðtrú}
\item[{pp}] {lýsingarháttur þátíðar}
\item[{praes}] {nútíð}
\item[{práv.}] {lögfræði, dómsmál}
\item[{predp}] {forskeyti}
\item[{prep}] {forsetning}
\item[{presp}] {lýsingarháttur nútíðar}
\item[{pron}] {fornafn}
\item[{prop}] {sérnafn}
\item[{přen.}] {afleidd merking}
\item[{přís.}] {málsháttur}
\item[{psych.}] {sálfræði}
\item[{refl}] {(v +) miðmynd, (pron +) afturbeygt fornafn}
\item[{rel}] {tilvísunarfornafn}
\item[{sg}] {eintala}
\item[{slang.}] {slangur}
\item[{sport.}] {íþróttir}
\item[{stav.}] {húsasmíði, byggingarlist}
\item[{subs}] {nafnorð (án kyns)}
\item[{sup}] {efstastig}
\item[{supin}] {sagnbót}
\item[{škol.}] {skólamál}
\item[{techn.}] {tækni, verkfræði}
\item[{v}] {sagnorð}
\item[{voj.}] {her}
\item[{zast.}] {úrelt mál}
\item[{zkr}] {skammstöfun}
\item[{zool.}] {dýrafræði}

\item[{;}] {semíkomma í miðri skýringu aðgreinir setningu; semíkomma í enda skýringar aðgreinir mismunandi merkingarsvið}

\item[{(...)}] {innan sviga eru viðbótarupplýsingar}
\item[{[...]}] {innan hornklofa er hljóðritaður framburður}
\item[{|}] {merkir stað, þar sem beygingarending bætist við stofn}
\item[{·}] {aukasamskeyti samsetts orðs}
\item[{··}] {aðalsamskeyti samsetts orðs}
\item[{+}] {plús}
\item[{/}] {skástrik táknar fleiri möguleika}
\item[{\dicsymSee}] {ör vísar til uppflettiorðs í orðabókinni}
\item[{\dicsymCompare}] {ör vísar til uppflettiorðs í bókinni til samanburðar}
\item[{\dicsymExampleIS}] {kemur á undan dæmi á íslensku}
\item[{\dicsymExampleCS}] {kemur á undan tékkneskri þýðingu á dæminu}
\item[{\dicsymIdiom}] {orð með háa tíðni}

\item[{\dicsymPhoto}] {ljósmynd}
\item[{\dicsymProverb}] {málsháttur}

\end{description}
\end{multicols}

% ------------------------------------------------- Headword part and alphabet

\chapter{Heslová část}
{\renewcommand{\arraystretch}{2}%
%\topskip0pt
\vstretch{1}
{\setlength{\tabcolsep}{1.5em}
\begin{table}[htpb]
\centering

\begin{tabular}{ccccccc} \hline
A a &  & G g &  & O o &  & V v \\
Á á  &  & H h &  & Ó ó &  & W w \\
B b &  & I i &  & P p &  & X x \\
C c &  & Í í &  & Q q &  & Y y \\
D d &  & J j &  & R r &  & Ý ý \\
Ð ð &  & K k &  & S s &  & Z z \\
E e &  & L l &  & T t &  & Þ þ \\
É é  &  & M m &  & U u &  & Æ æ \\
F f &  & N n &  & Ú ú &  & Ö ö \\
\hline
\end{tabular}

%\caption{\textbf{Islandská abeceda}}
\label{my-islandska-abeceda}
\end{table}
}
\vstretch{1}
{\renewcommand{\arraystretch}{1}%

\clearpage

\fi

% =========================================================== THE DICTIONARY =

\ifinputletters

\cleardoublepage

\dictionarygeometry
\pagestyle{myheadings}

\newcommand*{\alphabetF}{%
  a,aa,b,c,d,e,ee,f,g,h,i,ii,j,k,l,m,n,o,oo,p,r,s1,s2,t,u,uu,v,w,x,y,yy,th,ae,oe}
\ifPDFtest
\dicLetter{v}{letter26}
\dicEntry[vað] \dicTerm{vað} \dicIPA{{v}{a}{\textlengthmark}{\texttheta}} \dicPos{n}[2] \dicFlx{(‑s, vöð)}[8] \dicDirectTranslationCS{brod};  \dicIdiom{vað}{ \dicPhraseIS{hafa vaðið fyrir neðan sig}} \dicLangCat{přen.} \dicDirectTranslationCS{hrát\,/\addthin vsadit na jistotu};  \dicPhraseIS{ríða á vaðið með e‑ð} \dicLangCat{přen.} \dicDirectTranslationCS{být první s~(čím), převzít iniciativu v~(čem)};  \dicPhraseIS{tefla á tæpasta vaðið} \dicLangCat{přen.} \dicDirectTranslationCS{sázet všechno na jednu kartu}
\dicEntry[vaða] \dicTerm{vaða\smash{\textsuperscript{1}}} \dicIPA{{v}{a}{\textlengthmark}{ð}{a}} \dicPos{f}[1] \dicFlx{(vöðu, vöður)}[20] \dicDirectTranslationCS{hejno, houf} \dicIndirectTranslationCS{(ryb, tuleňů, velryb)}
\dicEntry[vaða] \dicTerm{vaða\smash{\textsuperscript{2}}} \dicsymFrequent\  \dicIPA{{v}{a}{\textlengthmark}{ð}{a}} \dicPos{v}[6] \dicFlx{(veð, óð, óðum, væði, vaðið)}[111] \textbf{1.} \dicDirectTranslationCS{brodit (se), brouzdat (se)} \dicExampleIS{vaða yfir ána} \dicExampleCS{brodit se přes řeku}  \textbf{2.} \dicDirectTranslationCS{vrazit, vtrhnout}  \textbf{3.} \dicPhraseIS{láta e‑ð vaða} \dicLangCat{přen.} \dicDirectTranslationCS{říct (co) na rovinu};  \dicPhraseIS{vaða úr einu í annað} \dicDirectTranslationCS{přeskakovat z~jednoho tématu na druhé};  \dicIdiom{vaða}[á]{ \dicPhraseIS{það veður á e‑m}} \dicFlx{impers} \dicLangCat{přen.} \dicDirectTranslationCS{(komu) se pusa nezastaví, (kdo) to sype ze sebe};  \dicIdiom{vaða}[í]{ \dicPhraseIS{vaða í peningum}} \dicLangCat{přen.} \dicDirectTranslationCS{topit se v~penězích};  \dicIdiom{vaða}[ofan]{ \dicPhraseIS{láta ekki vaða ofan í sig}} \dicLangCat{přen.} \dicDirectTranslationCS{nenechat si šlapat po hlavě};  \dicIdiom{vaða}[uppi]{ \dicPhraseIS{vaða uppi}} \dicDirectTranslationCS{dovolovat si, vyskakovat si};  \dicIdiom{vaða}[yfir]{ \dicPhraseIS{vaða yfir e‑n}} \dicDirectTranslationCS{nebrat ohledy na (koho)};  \dicIdiom{vaðast}[út]{ \dicPhraseIS{vaðast út}} \dicFlx{refl} \dicSynonym{blotna} \dicDirectTranslationCS{zmočit se, promočit se, promáčet se};  \dicIdiom{vaða}{ \dicPhraseIS{tunglið veður í skýjum}} \dicDirectTranslationCS{měsíc se prodírá mraky}
\dicEntry[vaðall] \dicTerm{vað|all} \dicIPA{{v}{a}{\textlengthmark}{ð}{a}{\textsubring{d}}{\textsubring{l}}} \dicPos{m}[6] \dicFlx{(‑als, ‑lar)}[28] \textbf{1.} \dicSynonym*{grunnt vatn} \dicDirectTranslationCS{mělčina, mělká voda, brod}  \textbf{2.} \dicSynonym{málæði} \dicDirectTranslationCS{proud řeči, tlachání}
\dicEntry[vaðfugl] \dicTerm{vað··fugl} \dicIPA{{v}{a}{ð}{f}{\textscy}{\r{g}}{\textsubring{l}}} \dicPos{m}[4] \dicFlx{(‑s, ‑ar)}[3] \dicFieldCat{zool.} \dicDirectTranslationCS{bahňák} \textit{(l.~{\textLA{Charadrii}})}  \dicsymPhoto\ 
\dicFigure{ds_image_vadfugl_0_1.jpg}{Vaðfugl}{Vaðfugl - Mdk572, CC BY-SA 3.0}
\dicEntry[vaðið] \dicTerm{vaðið} \dicIPA{{v}{a}{\textlengthmark}{ð}{\textsci}{\texttheta}} \dicPos{v} \dicFlx{supin} \dicLink{vaða\smash{\textsuperscript{2}}}
\dicEntry[vaðmál] \dicTerm{vað··mál} \dicIPA{{v}{a}{ð}{m}{au}{\textsubring{l}}} \dicPos{n}[2] \dicFlx{(‑s)}[2] \dicDirectTranslationCS{ševiot} \dicIndirectTranslationCS{(hrubší, drsnější vlna)}
\dicEntry[vaðstígvél] \dicTerm{vað··stíg·vél} \dicIPA{{v}{a}{ð}{s}{\textsubring{d}}{i}{\textbabygamma}{v}{j}{\textepsilon}{\textsubring{l}}} \dicPos{n}[2] \dicFlx{(‑s, ‑)}[5] \dicDirectTranslationCS{rybářské boty, boty na brodění}
\dicEntry[vaf] \dicTerm{vaf} \dicIPA{{v}{a}{\textlengthmark}{f}} \dicPos{n}[2] \dicFlx{(‑s, vöf)}[8] \dicDirectTranslationCS{zábal, zabalení};  \dicPhraseIS{e‑að er þungt í vöfum} \dicDirectTranslationCS{(co) je těžkopádné\,/\addthin obtížné}
\dicEntry[vafaatriði] \dicTerm{vafa··at·riði} \dicIPA{{v}{a}{\textlengthmark}{v}{a}{a}{\textsubring{d}}{r}{\textsci}{ð}{\textsci}} \dicPos{n}[2] \dicFlx{(‑s, ‑)}[14] \dicDirectTranslationCS{pochybnost, pochyba}
\dicEntry[vafalaus] \dicTerm{vafa··laus} \dicIPA{{v}{a}{\textlengthmark}{v}{a}{l}{\oe i}{s}} \dicPos{adj}[5]\dicFlx{}[-1] \dicSynonym{efalaus} \dicDirectTranslationCS{nepochybný, nesporný} \dicExampleIS{Þetta er vafalaus eign hans.} \dicExampleCS{To je jeho nesporné vlastnictví.}
\dicEntry[vafalaust] \dicTerm{vafa··laust} \dicsymFrequent\  \dicIPA{{v}{a}{\textlengthmark}{v}{a}{l}{\oe i}{s}{\textsubring{d}}} \dicPos{adv} \dicSynonym{efalaust} \dicDirectTranslationCS{nepochybně, nesporně} \dicExampleIS{Hann var vafalaust með hitasótt.} \dicExampleCS{Nepochybně měl horečku.}
\dicEntry[vafamál] \dicTerm{vafa··mál} \dicIPA{{v}{a}{\textlengthmark}{v}{a}{m}{au}{\textsubring{l}}} \dicPos{n}[2] \dicFlx{(‑s, ‑)}[5] \dicSynonym{efamál} \dicDirectTranslationCS{pochybnost, nejistota, otázka} \dicExampleIS{Það er vafamál að hann komi.} \dicExampleCS{To je otázka, jestli přijde.};  \dicPhraseIS{það er ekkert vafamál að (e‑að gerist)} \dicDirectTranslationCS{není nejmenší pochybnosti, že (se (co) stane)}
\dicEntry[vafasamur] \dicTerm{vafa··|samur} \dicIPA{{v}{a}{\textlengthmark}{v}{a}{s}{a}{m}{\textscy}{\textsubring{r}}} \dicPos{adj}[1] \dicFlx{(f ‑söm)}[2] \dicSynonym*{efasamur} \dicDirectTranslationCS{pochybný, nejistý, sporný}
\dicEntry[vafði] \dicTerm{vafði} \dicIPA{{v}{a}{v}{ð}{\textsci}} \dicPos{v} \dicFlx{ind pf sg 1 pers} \dicLink{vefja}
\dicEntry[vaff] \dicTerm{vaff} \dicIPA{{v}{a}{f}{\textlengthmark}} \dicPos{n}[2] \dicFlx{(‑s, vöff)}[8] \dicDirectTranslationCS{název písmene \clqq v\crqq };  \dicPhraseIS{tvöfalt vaff} \dicDirectTranslationCS{dvojité \clqq v\crqq , název písmene \clqq w\crqq }
\dicEntry[vaffla] \dicTerm{vaffla} \dicIPA{{v}{a}{f}{l}{a}} \dicPos{f}[1] \dicFlx{(vöfflu, vöfflur)}[20] \dicFieldCat{kulin.} \dicDirectTranslationCS{vafle, oplatka}
\dicEntry[vafi] \dicTerm{vaf|i} \dicsymFrequent\  \dicIPA{{v}{a}{\textlengthmark}{v}{\textsci}} \dicPos{m}[1] \dicFlx{(‑a)}[3] \dicSynonym{efi} \dicDirectTranslationCS{pochybnost, pochyba, nejistota};  \dicPhraseIS{velkjast í vafa} \dicFlx{refl} \dicDirectTranslationCS{být na pochybách, pochybovat};  \dicPhraseIS{vera í vafa um e‑ð} \dicDirectTranslationCS{pochybovat o~(čem), být na pochybách o~(čem), nebýt si jistý (čím)} \dicExampleIS{Ég er í vafa um hvort þetta er leyfilegt.} \dicExampleCS{Nejsem si jistý, jestli to je povolené.};  \dicPhraseIS{það er enginn vafi að (e‑að gerist)} \dicFlx{impers} \dicDirectTranslationCS{není pochyb o~tom, že (se (co) stane)}
\dicEntry[vafið] \dicTerm{vafið} \dicIPA{{v}{a}{\textlengthmark}{v}{\textsci}{\texttheta}} \dicPos{v} \dicFlx{supin} \dicLink{vefja}
\dicEntry[vafinn] \dicTerm{vafinn} \dicIPA{{v}{a}{\textlengthmark}{v}{\textsci}{\textsubring{n}}} \dicPos{adj}[6]\dicFlx{}[-14] \dicPhraseIS{vera skuldum vafinn} \dicLangCat{přen.} \dicDirectTranslationCS{topit se v~dluzích}
\dicEntry[vafningalaus] \dicTerm{vafninga··laus} \dicIPA{{v}{a}{\textsubring{b}}{n}{i}{\ng}{\r{g}}{a}{l}{\oe i}{s}} \dicPos{adj}[5]\dicFlx{}[-1] \dicSynonym{hiklaus} \dicDirectTranslationCS{okamžitý, pohotový}
\dicEntry[vafningalaust] \dicTerm{vafninga··laust} \dicIPA{{v}{a}{\textsubring{b}}{n}{i}{\ng}{\r{g}}{a}{l}{\oe i}{s}{\textsubring{d}}} \dicPos{adv} \dicDirectTranslationCS{bez váhání, okamžitě, pohotově}
\dicEntry[vafningsjurt] \dicTerm{vafnings··jurt}\dicTerm{, vafningsplanta} \dicIPA{{v}\-{a}\-{\textsubring{b}}\-{n}\-{i}\-{\ng}\-{\r{g}}\-{s}\-{j}\-{\textscy}\-{\textsubring{r}}\-{\textsubring{d}}\-} \dicPos{f}[7] \dicFlx{(‑ar, ‑ir)}[1] \dicDirectTranslationCS{popínavá rostlina, liána}
\dicEntry[vafningsplanta] \dicTerm{vafnings··|planta} \dicIPA{{v}{a}{\textsubring{b}}{n}{i}{\ng}{\r{g}}{s}{p\smash{\textsuperscript{h}}}{l}{a}{\textsubring{n}}{\textsubring{d}}{a}} \dicPos{f}[1] \dicFlx{(‑plöntu, ‑plöntur)}[14] \dicLink{vafningsjurt}
\dicEntry[vafningur] \dicTerm{vafn··ing|ur} \dicIPA{{v}{a}{\textsubring{b}}{n}{i}{\ng}{\r{g}}{\textscy}{\textsubring{r}}} \dicPos{m}[6] \dicFlx{(‑s, ‑ar)}[8] \textbf{1.} \dicSynonym{vaf} \dicDirectTranslationCS{závit, smyčka} \dicIndirectTranslationCS{(něco zavinutého)}  \textbf{2.} \dicSynonym{flétta\smash{\textsuperscript{1}}} \dicDirectTranslationCS{role, rolka, pletenec}  \textbf{3.} \dicPhraseIS{vafningar} \dicFlx{pl} \dicSynonym{undanbrögð} \dicDirectTranslationCS{vytáčky, vykrucování}
\dicEntry[vafra] \dicTerm{vafr|a} \dicIPA{{v}{a}{v}{r}{a}} \dicPos{v}[1] \dicFlx{(‑aði)}[13] \textbf{1.} \dicFieldCat{poč.} \dicDirectTranslationCS{prohlížet\,/\addthin prohledávat (internet), brouzdat (na webu ap.), surfovat (po internetu ap.)}  \textbf{2.} \dicSynonym{reika} \dicDirectTranslationCS{brouzdat se, bloumat, brousit} \dicExampleIS{vafra um} \dicExampleCS{brouzdat se kolem}
\dicEntry[vafrakaka] \dicTerm{vafra··|kaka} \dicIPA{{v}{a}{v}{r}{a}{k\smash{\textsuperscript{h}}}{a}{\r{g}}{a}} \dicPos{f}[1] \dicFlx{(‑köku, ‑kökur)}[20] \dicFieldCat{poč.} \dicDirectTranslationCS{cookie}
\dicEntry[vafrari] \dicTerm{vafr··ar|i}\dicTerm{, vafri} \dicIPA{{v}\-{a}\-{v}\-{r}\-{a}\-{r}\-{\textsci}\-} \dicPos{m}[1] \dicFlx{(‑a, ‑ar)}[11] \dicFieldCat{poč.} \dicSynonym{vefskoðari} \dicDirectTranslationCS{(webový) prohlížeč, browser}
\dicEntry[vafri] \dicTerm{vafr|i} \dicIPA{{v}{a}{v}{r}{\textsci}} \dicPos{m}[1] \dicFlx{(‑a, ‑ar)}[8] \dicLink{vafrari}
\dicEntry[vafstra] \dicTerm{vafstr|a} \dicIPA{{v}{a}{f}{s}{\textsubring{d}}{r}{a}} \dicPos{v}[1] \dicFlx{(‑aði)}[50] \dicSynonym{bauka} \dicDirectTranslationCS{dřít (se), pachtit se, lopotit se}
\dicEntry[vafstur] \dicTerm{vafstur} \dicIPA{{v}{a}{f}{s}{\textsubring{d}}{\textscy}{\textsubring{r}}} \dicPos{n}[2] \dicFlx{(‑s)}[28] \dicDirectTranslationCS{potíž, starost, svízel} \dicExampleIS{Hvaða vafstur er þetta í þér?} \dicExampleCS{Co ti dělá starosti?}
\dicEntry[vaga] \dicTerm{vag|a} \dicIPA{{v}{a}{\textlengthmark}{\textbabygamma}{a}} \dicPos{v}[1] \dicFlx{(‑aði)}[50] \dicDirectTranslationCS{kolébat se (při chůzi), jít kolébavou chůzí (jako kachna ap.)}
\dicEntry[vagga] \dicTerm{vagga\smash{\textsuperscript{1}}} \dicIPA{{v}{a}{\r{g}}{\textlengthmark}{a}} \dicPos{f}[1] \dicFlx{(vöggu, vöggur)}[8] \textbf{1.} \dicSynonym{barnarúm} \dicDirectTranslationCS{kolébka};  \dicPhraseIS{frá vöggu til grafar} \dicFlx{adv} \dicLangCat{přen.} \dicDirectTranslationCS{od kolébky (až) do hrobu}  \textbf{2.} \dicLangCat{přen.} \dicDirectTranslationCS{kolébka, vlast, původ}
\dicEntry[vagga] \dicTerm{vagg|a\smash{\textsuperscript{2}}} \dicIPA{{v}{a}{\r{g}}{\textlengthmark}{a}} \dicPos{v}[1] \dicFlx{(‑aði)}[13] \dicFlx{dat} \textbf{1.} \dicFlx{dat} \dicSynonym{rugga} \dicDirectTranslationCS{kolébat, houpat, pohupovat} \dicExampleIS{vagga barni} \dicExampleCS{kolébat dítě}  \textbf{2.} \dicDirectTranslationCS{kolébat se, houpat se, kývat se}
\dicEntry[vagl] \dicTerm{vagl} \dicIPA{{v}{a}{\r{g}}{\textsubring{l}}} \dicPos{n}[2] \dicFlx{(‑s, vögl)}[8] \dicFieldCat{med. } \dicDirectTranslationCS{šedý zákal, katarakta}
\dicEntry[vagn] \dicTerm{vagn} \dicsymFrequent\  \dicIPA{{v}{a}{\r{g}}{\textsubring{n}}} \dicPos{m}[4] \dicFlx{(‑s, ‑ar)}[4] \textbf{1.} \dicSynonym{ökutæki} \dicDirectTranslationCS{vozidlo, vůz, vagon} \dicExampleIS{í opnum vögnum} \dicExampleCS{v~otevřených vagonech}  \textbf{2.} \dicSynonym{strætisvagn} \dicDirectTranslationCS{(městský) autobus}  \textbf{3.} \dicSynonym{barnavagn} \dicDirectTranslationCS{vozík, kočárek}
\dicEntry[vagnkjálki] \dicTerm{vagn··kjálk|i} \dicIPA{{v}{a}{\r{g}}{\textsubring{n}}{c\smash{\textsuperscript{h}}}{au}{\textsubring{l}}{\r{\textObardotlessj}}{\textsci}} \dicPos{m}[1] \dicFlx{(‑a, ‑ar)}[1] \dicDirectTranslationCS{oj}
\dicEntry[vagnstjóri] \dicTerm{vagn··stjór|i} \dicIPA{{v}{a}{\r{g}}{\textsubring{n}}{s}{\textsubring{d}}{j}{ou}{r}{\textsci}} \dicPos{m}[1] \dicFlx{(‑a, ‑ar)}[1] \dicSynonym{ökumaður} \dicDirectTranslationCS{řidič(ka) (autobusu ap.)}
\dicEntry[vaka] \dicTerm{vaka\smash{\textsuperscript{1}}} \dicIPA{{v}{a}{\textlengthmark}{\r{g}}{a}} \dicPos{f}[1] \dicFlx{(vöku, vökur)}[8] \textbf{1.} \dicSynonym*{það að vaka} \dicDirectTranslationCS{bdění, bdělost};  \dicPhraseIS{halda vöku fyrir e‑m} \dicDirectTranslationCS{udržovat (koho) v~bdělosti}  \textbf{2.} \dicLangCat{zast.} \dicSynonym{kvöldvaka} \dicIndirectTranslationCS{doba od zhasnutí světel v~domácnosti do uložení se ke spánku}  \textbf{3.} \dicSynonym*{varðtími} \dicDirectTranslationCS{hlídka, stráž}
\dicEntry[vaka] \dicTerm{vak|a\smash{\textsuperscript{2}}} \dicsymFrequent\  \dicIPA{{v}{a}{\textlengthmark}{\r{g}}{a}} \dicPos{v}[2] \dicFlx{(‑ti, ‑að)}[87] \dicDirectTranslationCS{bdít, být vzhůru} \dicExampleIS{vaka til morguns} \dicExampleCS{bdít do rána};  \dicIdiom{vaka}[eftir]{ \dicPhraseIS{vaka eftir e‑m}} \dicDirectTranslationCS{čekat\,/\addthin vyčkávat na (koho) (dlouho do noci ap.)};  \dicIdiom{vaka}[fram]{ \dicPhraseIS{vaka fram á nótt}} \dicDirectTranslationCS{zůstat vzhůru dlouho do noci};  \dicIdiom{vaka}[fyrir]{ \dicPhraseIS{e‑að vakir fyrir e‑m}} \dicDirectTranslationCS{(co kdo) zamýšlí, (co) má (kdo) za lubem};  \dicIdiom{vaka}[yfir]{ \dicPhraseIS{vaka yfir e‑m}} \dicDirectTranslationCS{dávat pozor na (koho), bdít nad (kým)} \dicExampleIS{vaka yfir börnum sínum} \dicExampleCS{bdít nad svými dětmi}
\dicEntry[vakandi] \dicTerm{vak··andi} \dicsymFrequent\  \dicIPA{{v}{a}{\textlengthmark}{\r{g}}{a}{n}{\textsubring{d}}{\textsci}} \dicPos{adj}[13] \dicFlx{indecl}[1] \textbf{1.} \dicSynonym*{sem vakir} \dicDirectTranslationCS{bdící, bdělý, (jsoucí) vzhůru}  \textbf{2.} \dicSynonym{opinn} \dicDirectTranslationCS{bdělý, ostražitý};  \dicPhraseIS{hafa vakandi auga með e‑u} \dicDirectTranslationCS{pozorně (co) sledovat};  \dicPhraseIS{vera vakandi fyrir e‑ð} \dicDirectTranslationCS{nespouštět (co) z~očí, bedlivě (co) sledovat}
\dicEntry[vakar] \dicTerm{vakar} \dicIPA{{v}{a}{\textlengthmark}{\r{g}}{a}{\textsubring{r}}} \dicPos{f} \dicFlx{sg gen} \dicLink{vök}
\dicEntry[vaki] \dicTerm{vak|i} \dicIPA{{v}{a}{\textlengthmark}{\r{\textObardotlessj}}{\textsci}} \dicPos{m}[1] \dicFlx{(‑a, ‑ar)}[8] \dicFieldCat{med.} \dicLangCat{zast.} \dicSynonym{hormón} \dicDirectTranslationCS{hormon}
\dicEntry[vakið] \dicTerm{vakið} \dicIPA{{v}{a}{\textlengthmark}{\r{\textObardotlessj}}{\textsci}{\texttheta}} \dicPos{v} \dicFlx{supin} \dicLink{vekja}
\dicEntry[vakinn] \dicTerm{vakinn} \dicIPA{{v}{a}{\textlengthmark}{\r{\textObardotlessj}}{\textsci}{\textsubring{n}}} \dicPos{adj}[6]\dicFlx{}[-12] \textbf{1.} \dicSynonym*{sem vakir} \dicDirectTranslationCS{bdící, bdělý, (jsoucí) vzhůru};  \dicPhraseIS{vera vakinn og sofinn í e‑u} \dicLangCat{přen.} \dicDirectTranslationCS{věnovat se (čemu) od rána do večera, věnovat se (čemu) tělem i~duší}  \textbf{2.} \dicSynonym*{sem blæs} \dicDirectTranslationCS{svěží, čerstvý (vítr ap.)} \dicExampleIS{vakinn vindur} \dicExampleCS{čerstvý vítr}
\dicEntry[vakir] \dicTerm{vakir} \dicIPA{{v}{a}{\textlengthmark}{\r{\textObardotlessj}}{\textsci}{\textsubring{r}}} \dicPos{f} \dicFlx{pl nom} \dicLink{vök}
\dicEntry[vakka] \dicTerm{vakk|a} \dicIPA{{v}{a}{h}{\r{g}}{a}} \dicPos{v}[1] \dicFlx{(‑aði)}[13] \dicDirectTranslationCS{kroužit, courat se, okounět} \dicExampleIS{vakka kringum e‑n} \dicExampleCS{kroužit kolem (koho)}
\dicEntry[vakna] \dicTerm{vakn|a} \dicsymFrequent\  \dicIPA{{v}{a}{h}{\r{g}}{n}{a}} \dicPos{v}[1] \dicFlx{(‑aði)}[43] \textbf{1.} \dicDirectTranslationCS{vzbudit se, (pro)budit se, probouzet se, probrat se} \dicExampleIS{vakna snemma} \dicExampleCS{probudit se časně}  \textbf{2.} \dicSynonym*{lifna við} \dicDirectTranslationCS{(pro)budit se, probouzet se, oživnout, ožívat} \dicExampleIS{Starfsáhuginn vaknar.} \dicExampleCS{Probouzí se chuť do práce.};  \dicPhraseIS{vakna til lífsins} \dicLangCat{přen.} \dicDirectTranslationCS{probudit\,/\addthin probouzet se k~životu, oživnout, ožívat};  \dicIdiom{vakna}[upp]{ \dicPhraseIS{vakna upp}} \dicDirectTranslationCS{probudit se, vzbudit se};  \dicIdiom{vakna}[við]{ \dicPhraseIS{vakna við e‑ð}} \dicDirectTranslationCS{probudit se (čím)} \dicExampleIS{vakna við símhringingu} \dicExampleCS{probudit se zvoněním telefonu}; { \dicPhraseIS{vakna við vondan draum}} \dicLangCat{přen.} \dicDirectTranslationCS{probudit se do reality}
\dicEntry[vaknaður] \dicTerm{vakn··aður} \dicIPA{{v}{a}{h}{\r{g}}{n}{a}{ð}{\textscy}{\textsubring{r}}} \dicPos{adj}[3] \dicFlx{(vöknuð)}[2] \dicDirectTranslationCS{probuzený, vzbuzený, probraný}
\dicEntry[vaknandi] \dicTerm{vakn··andi} \dicIPA{{v}{a}{h}{\r{g}}{n}{a}{n}{\textsubring{d}}{\textsci}} \dicPos{adj}[13] \dicFlx{indecl}[1] \dicDirectTranslationCS{budící se, probouzející se}
\dicEntry[vakning] \dicTerm{vak··ning} \dicIPA{{v}{a}{h}{\r{g}}{n}{i}{\ng}{\r{g}}} \dicPos{f}[4] \dicFlx{(‑ar)}[7] \textbf{1.} \dicSynonym*{það að vekja} \dicDirectTranslationCS{(pro)buzení}  \textbf{2.} \dicSynonym*{trúarvakning} \dicDirectTranslationCS{probuzení, znovuoživení, obrození} \dicExampleIS{félagsleg vakning} \dicExampleCS{společenské obrození}
\dicEntry[vakt] \dicTerm{vakt} \dicIPA{{v}{a}{x}{\textsubring{d}}} \dicPos{f}[7] \dicFlx{(‑ar, ‑ir)}[1] \textbf{1.} \dicSynonym{gæsla} \dicDirectTranslationCS{hlídka, hlídání, střežení}  \textbf{2.} \dicSynonym*{vinnutímabil} \dicDirectTranslationCS{(pracovní) směna};  \dicPhraseIS{vinna í\,/\addthin á vöktum} \dicDirectTranslationCS{pracovat na směny}  \textbf{3.} \dicDirectTranslationCS{(pracovní) služba} \dicExampleIS{á vakt} \dicExampleCS{ve službě}
\dicEntry[vakta] \dicTerm{vakt|a} \dicIPA{{v}{a}{x}{\textsubring{d}}{a}} \dicPos{v}[1] \dicFlx{(‑aði)}[13] \dicFlx{acc} \dicDirectTranslationCS{hlídat, střežit, monitorovat} \dicExampleIS{vakta svæðið} \dicExampleCS{monitorovat oblast}
\dicEntry[vaktavinna] \dicTerm{vakta··vinn|a} \dicIPA{{v}{a}{x}{\textsubring{d}}{a}{v}{\textsci}{n}{a}} \dicPos{f}[1] \dicFlx{(‑u)}[5] \dicDirectTranslationCS{práce na směny}
\dicEntry[vakti] \dicTerm{vakti} \dicIPA{{v}{a}{x}{\textsubring{d}}{\textsci}} \dicPos{v} \dicFlx{ind pf sg 1 pers} \dicLink{vekja}
\dicEntry[vaktmaður] \dicTerm{vakt··|maður} \dicIPA{{v}{a}{x}{\textsubring{d}}{m}{a}{ð}{\textscy}{\textsubring{r}}} \dicPos{m}[13] \dicFlx{(‑manns, ‑menn)}[2] \dicSynonym{varðmaður} \dicDirectTranslationCS{hlídač(ka), dozorce, dozorkyně}
\dicEntry[val] \dicTerm{val} \dicIPA{{v}{a}{\textlengthmark}{\textsubring{l}}} \dicPos{n}[2] \dicFlx{(‑s, völ)}[8] \textbf{1.} \dicSynonym*{það að velja} \dicDirectTranslationCS{volba, výběr} \dicExampleIS{val á e‑u} \dicExampleCS{volba (čeho)};  \dicPhraseIS{að eigin vali} \dicFlx{adv} \dicDirectTranslationCS{podle vlastní volby, podle vlastního výběru};  \dicPhraseIS{eiga ekkert val} \dicDirectTranslationCS{nemít na výběr}  \textbf{2.} \dicFieldCat{škol.} \dicDirectTranslationCS{volitelný předmět}
\dicEntry[vald] \dicTerm{vald} \dicsymFrequent\  \dicIPA{{v}{a}{l}{\textsubring{d}}} \dicPos{n}[2] \dicFlx{(‑s, völd)}[8] \textbf{1.} \dicDirectTranslationCS{moc, síla, kontrola};  \dicPhraseIS{beita valdi} \dicDirectTranslationCS{užít síly};  \dicPhraseIS{gera allt sem í mannlegu valdi stendur} \dicDirectTranslationCS{udělat vše, co je v~lidských silách};  \dicPhraseIS{gefa sig e‑u á vald} \dicDirectTranslationCS{poddat se síle (čeho) (básně ap.)};  \dicPhraseIS{hafa e‑n á valdi sínu} \dicDirectTranslationCS{mít (koho) ve své moci};  \dicPhraseIS{hafa vald á e‑u} \dicDirectTranslationCS{ovládat (co), mít (co) pod kontrolou};  \dicPhraseIS{ná valdi á e‑u} \dicDirectTranslationCS{ovládnout (co), získat kontrolu nad (čím)};  \dicPhraseIS{taka e‑n með valdi} \dicDirectTranslationCS{zmocnit se (koho) násilím}  \textbf{2.} \dicPhraseIS{völd} \dicFlx{pl} \dicDirectTranslationCS{síla, moc};  \dicPhraseIS{af völdum e‑s} \dicDirectTranslationCS{z~důvodu (čeho), příčinou (čeho)};  \dicPhraseIS{e‑að situr (enn) að völdum} \dicDirectTranslationCS{(co) je (stále) u\,/\addthin při moci};  \dicPhraseIS{hrökklast frá völdum} \dicFlx{refl} \dicDirectTranslationCS{přijít o~moc};  \dicPhraseIS{komast til valda} \dicFlx{refl} \dicDirectTranslationCS{dostat se k~moci, přijít k~moci};  \dicPhraseIS{taka völdin} \dicDirectTranslationCS{převzít moc, ujmout se moci};  \dicPhraseIS{vera við völd} \dicDirectTranslationCS{být u\,/\addthin při moci};  \dicProverb\  \dicPhraseIS{Valdið spillir.} \dicLangCat{přís.} \dicDirectTranslationCS{Moc kazí.}
\dicEntry[valda] \dicTerm{vald|a\smash{\textsuperscript{1}}} \dicIPA{{v}{a}{l}{\textsubring{d}}{a}} \dicPos{v}[1] \dicFlx{(‑aði)}[13] \dicFlx{acc} \textbf{1.} \dicDirectTranslationCS{chránit, střežit} \dicIndirectTranslationCS{(v~šachách)} \dicExampleIS{valda peðið með riddaranum} \dicExampleCS{střežit pěšce koněm}  \textbf{2.} \dicFieldCat{mat.} \dicDirectTranslationCS{umocnit} \dicExampleIS{valda tölu} \dicExampleCS{umocnit číslo}
\dicEntry[valda] \dicTerm{valda\smash{\textsuperscript{2}}} \dicsymFrequent\  \dicIPA{{v}{a}{l}{\textsubring{d}}{a}} \dicPos{v}[7] \dicFlx{(veld, olli, ollum, ylli, valdið)}[11] \dicFlx{(dat +) dat} \dicSynonym{orsaka} \dicDirectTranslationCS{způsobit, způsobovat, přivodit, vyvolat, vyvolávat} \dicExampleIS{valda honum áhyggjum} \dicExampleCS{způsobit mu obavy};  \dicProverb\  \dicPhraseIS{Sjaldan veldur einn þá tveir deila.} \dicLangCat{přís.} \dicDirectTranslationCS{K~hádce jsou zapotřebí dva.}
\dicEntry[valda] \dicTerm{valda\smash{\textsuperscript{3}}} \dicsymFrequent\  \dicIPA{{v}{a}{l}{\textsubring{d}}{a}} \dicPos{v}[7] \dicFlx{(veld, olli, ollum, ylli, valdið)}[11] \dicFlx{dat} \dicSynonym*{ráða við} \dicDirectTranslationCS{dokázat, zvládnout} \dicIndirectTranslationCS{(nepoužívá se v~minulém čase, kde je nutná jiná formulace)} \dicExampleIS{valda verkefninu} \dicExampleCS{zvládnout zadání}
\dicEntry[valdabarátta] \dicTerm{valda··bar·átt|a} \dicIPA{{v}{a}{l}{\textsubring{d}}{a}{\textsubring{b}}{a}{r}{au}{h}{\textsubring{d}}{a}} \dicPos{f}[1] \dicFlx{(‑u)}[5] \dicDirectTranslationCS{boj o~moc, mocenský boj}
\dicEntry[valdajafnvægi] \dicTerm{valda··jafn·vægi} \dicIPA{{v}{a}{l}{\textsubring{d}}{a}{j}{a}{\textsubring{b}}{n}{v}{a}{i}{j}{\textsci}} \dicPos{n}[2] \dicFlx{(‑s)}[20] \dicDirectTranslationCS{rovnováha sil\,/\addthin moci}
\dicEntry[valdalaus] \dicTerm{valda··laus} \dicIPA{{v}{a}{l}{\textsubring{d}}{a}{l}{\oe i}{s}} \dicPos{adj}[5]\dicFlx{}[-1] \dicDirectTranslationCS{(jsoucí) bez moci, nemající moc, bezmocný}
\dicEntry[valdamaður] \dicTerm{valda··|maður} \dicIPA{{v}{a}{l}{\textsubring{d}}{a}{m}{a}{ð}{\textscy}{\textsubring{r}}} \dicPos{m}[13] \dicFlx{(‑manns, ‑menn)}[2] \dicDirectTranslationCS{mocnář(ka), vladař(ka)}
\dicEntry[valdarán] \dicTerm{valda··rán} \dicIPA{{v}{a}{l}{\textsubring{d}}{a}{r}{au}{\textsubring{n}}} \dicPos{n}[2] \dicFlx{(‑s, ‑)}[5] \dicDirectTranslationCS{(politický) převrat, puč}
\dicEntry[valdasjúkur] \dicTerm{valda··sjúkur} \dicIPA{{v}{a}{l}{\textsubring{d}}{a}{s}{j}{u}{\r{g}}{\textscy}{\textsubring{r}}} \dicPos{adj}[1]\dicFlx{}[-1] \dicDirectTranslationCS{chtivý moci, mocichtivý}
\dicEntry[valdastreita] \dicTerm{valda··streit|a} \dicIPA{{v}{a}{l}{\textsubring{d}}{a}{s}{\textsubring{d}}{r}{ei}{\textsubring{d}}{a}} \dicPos{f}[1] \dicFlx{(‑u)}[5] \dicDirectTranslationCS{boj o~moc, mocenský boj}
\dicEntry[valdataka] \dicTerm{valda··|taka} \dicIPA{{v}{a}{l}{\textsubring{d}}{a}{t\smash{\textsuperscript{h}}}{a}{\r{g}}{a}} \dicPos{f}[1] \dicFlx{(‑töku, ‑tökur)}[20] \dicDirectTranslationCS{převzetí moci}
\dicEntry[valdbeiting] \dicTerm{vald··beit·ing} \dicIPA{{v}{a}{l}{\textsubring{d}}{\textsubring{b}}{ei}{\textsubring{d}}{i}{\ng}{\r{g}}} \dicPos{f}[4] \dicFlx{(‑ar)}[7] \dicDirectTranslationCS{použití síly\,/\addthin násilí}
\dicEntry[valdi] \dicTerm{valdi} \dicIPA{{v}{a}{l}{\textsubring{d}}{\textsci}} \dicPos{v} \dicFlx{ind pf sg 1 pers} \dicLink{velja}
\dicEntry[valdið] \dicTerm{valdið} \dicIPA{{v}{a}{l}{\textsubring{d}}{\textsci}{\texttheta}} \dicPos{v} \dicFlx{supin} \textbf{1.} \dicLink{valda\smash{\textsuperscript{2}}}  \textbf{2.} \dicLink{valda\smash{\textsuperscript{3}}}
\dicEntry[valdníðsla] \dicTerm{vald··níðsl|a} \dicIPA{{v}{a}{l}{\textsubring{d}}{n}{i}{ð}{s}{\textsubring{d}}{l}{a}} \dicPos{f}[1] \dicFlx{(‑u)}[5] \dicDirectTranslationCS{zneužití moci}
\dicEntry[valdsmaður] \dicTerm{valds··|maður} \dicIPA{{v}{a}{l}{\textsubring{d}}{s}{m}{a}{ð}{\textscy}{\textsubring{r}}} \dicPos{m}[13] \dicFlx{(‑manns, ‑menn)}[2] \dicDirectTranslationCS{mocnář(ka), potentát(ka)}
\dicEntry[valdsmannslegur] \dicTerm{valds·manns··legur} \dicIPA{{v}{a}{l}{\textsubring{d}}{s}{m}{a}{n}{s}{l}{\textepsilon}{\textbabygamma}{\textscy}{\textsubring{r}}} \dicPos{adj}[1]\dicFlx{}[-8] \dicDirectTranslationCS{mocnářský, panovnický, vladařský}
\dicEntry[valdsvið] \dicTerm{vald··svið} \dicIPA{{v}{a}{l}{\textsubring{d}}{s}{v}{\textsci}{\texttheta}} \dicPos{n}[2] \dicFlx{(‑s, ‑)}[5] \dicSynonym{umráðasvæði} \dicDirectTranslationCS{oblast\,/\addthin sféra moci}
\dicEntry[valdur] \dicTerm{valdur} \dicIPA{{v}{a}{l}{\textsubring{d}}{\textscy}{\textsubring{r}}} \dicPos{adj}[2] \dicFlx{(f völd)}[16] \dicPhraseIS{vera valdur að e‑u} \dicDirectTranslationCS{být původcem (čeho), být odpovědný za (co), mít (co) na svědomí} \dicExampleIS{vera valdur að slysi} \dicExampleCS{být odpovědný za nehodu}
\dicEntry[valfrelsi] \dicTerm{val··frelsi} \dicIPA{{v}{a}{l}{f}{r}{\textepsilon}{l}{s}{\textsci}} \dicPos{n}[2] \dicFlx{(‑s)}[20] \dicDirectTranslationCS{svoboda\,/\addthin možnost volby\,/\addthin výběru}
\dicEntry[valfrjáls] \dicTerm{val··frjáls} \dicIPA{{v}{a}{l}{f}{r}{j}{au}{l}{s}} \dicPos{adj}[5]\dicFlx{}[-1] \dicDirectTranslationCS{volitelný, výběrový, fakultativní} \dicExampleIS{Námskeiðið er valfrjálst.} \dicExampleCS{Tento kurz je volitelný.}
\dicEntry[valgrein] \dicTerm{val··grein} \dicIPA{{v}{a}{l}{\r{g}}{r}{ei}{\textsubring{n}}} \dicPos{f}[4] \dicFlx{(‑ar, ‑ar)}[1] \dicFieldCat{škol.} \dicDirectTranslationCS{volitelný předmět}
\dicEntry[valhneta] \dicTerm{val··hnet|a}\dicTerm{, valhnot} \dicIPA{{v}\-{a}\-{l}\-{\textsubring{n}}\-{\textepsilon}\-{\textsubring{d}}\-{a}\-} \dicPos{f}[1] \dicFlx{(‑u, ‑ur)}[7] \dicDirectTranslationCS{vlašský ořech} \dicIndirectTranslationCS{(plod)}
\dicEntry[valhnot] \dicTerm{val··hnot} \dicIPA{{v}{a}{l}{\textsubring{n}}{\textopeno}{\textsubring{d}}} \dicPos{f}[7] \dicFlx{(‑ar, ‑ir)}[1] \dicLink{valhneta}
\dicEntry[valhnota] \dicTerm{val··hnot|a} \dicIPA{{v}{a}{l}{\textsubring{n}}{\textopeno}{\textsubring{d}}{a}} \dicPos{f}[1] \dicFlx{(‑u, ‑ur)}[7] \dicDirectTranslationCS{ořech} \dicIndirectTranslationCS{(dřevo jako stavební materiál)}
\dicEntry[valhnotutré] \dicTerm{val·hnotu··tré} \dicIPA{{v}{a}{l}{\textsubring{n}}{\textopeno}{\textsubring{d}}{\textscy}{t\smash{\textsuperscript{h}}}{r}{j}{\textepsilon}} \dicPos{n}[2] \dicFlx{(‑s, ‑)}[36] \dicFieldCat{bot.} \dicDirectTranslationCS{ořešák, ořešák královský} \textit{(l.~{\textLA{Juglans regia}})}  \dicsymPhoto\ 
\dicFigure{1714.jpg}{Valhnotutré}{Valhnotutré - Zicha Ondřej, Biolib, Copyright/CC-BY-NC}
\dicEntry[valhopp] \dicTerm{val··hopp} \dicIPA{{v}{a}{\textlengthmark}{\textsubring{l}}{h}{\textopeno}{h}{\textsubring{b}}} \dicPos{n}[2] \dicFlx{(‑s, ‑)}[5] \dicDirectTranslationCS{drobný klus}
\dicEntry[Valhöll] \dicTerm{Val··|höll} \dicIPA{{v}{a}{\textlengthmark}{\textsubring{l}}{h}{\oe}{\textsubring{d}}{\textsubring{l}}} \dicPos{f}[7] \dicFlx{(‑hallar)}[21] \dicFieldCat{myt.} \dicDirectTranslationCS{Valhalla} \dicIndirectTranslationCS{(Ódinova síň v~Ásgardu, ve které přebývají padlí bojovníci)}
\dicEntry[valið] \dicTerm{valið} \dicIPA{{v}{a}{\textlengthmark}{l}{\textsci}{\texttheta}} \dicPos{v} \dicFlx{supin} \dicLink{velja}
\dicEntry[valinkunnur] \dicTerm{valin··kunnur} \dicIPA{{v}{a}{\textlengthmark}{l}{\textsci}{n}{k\smash{\textsuperscript{h}}}{\textscy}{n}{\textscy}{\textsubring{r}}} \dicPos{adj}[1]\dicFlx{}[-1] \dicSynonym{heiðarlegur} \dicDirectTranslationCS{vážený, respektovaný} \dicExampleIS{valinkunnur maður} \dicExampleCS{vážený pán}
\dicEntry[valinn] \dicTerm{valinn} \dicIPA{{v}{a}{\textlengthmark}{l}{\textsci}{\textsubring{n}}} \dicPos{adj}[6]\dicFlx{}[-8] \textbf{1.} \dicSynonym*{kosinn} \dicDirectTranslationCS{vybraný, zvolený}  \textbf{2.} \dicSynonym{ágætur} \dicDirectTranslationCS{vybraný, prvotřídní}
\dicEntry[valkostur] \dicTerm{val··kost|ur} \dicIPA{{v}{a}{l}{k\smash{\textsuperscript{h}}}{\textopeno}{s}{\textsubring{d}}{\textscy}{\textsubring{r}}} \dicPos{m}[10] \dicFlx{(‑ar, ‑ir)}[4] \dicSynonym{val} \dicDirectTranslationCS{možnost, alternativa, volba} \dicExampleIS{eiga valkosti} \dicExampleCS{mít možnosti}
\dicEntry[valkyrja] \dicTerm{val··kyrj|a} \dicIPA{{v}{a}{l}{c\smash{\textsuperscript{h}}}{\textsci}{r}{j}{a}} \dicPos{f}[1] \dicFlx{(‑u, ‑ur)}[7] \dicFieldCat{myt.} \dicDirectTranslationCS{valkýra}
\dicEntry[vall] \dicTerm{vall} \dicIPA{{v}{a}{\textsubring{d}}{\textsubring{l}}} \dicPos{v} \dicFlx{ind pf sg 1 pers} \dicLink{vella\smash{\textsuperscript{2}}}
\dicEntry[vallar] \dicTerm{vallar} \dicIPA{{v}{a}{\textsubring{d}}{l}{a}{\textsubring{r}}} \dicPos{m} \dicFlx{sg gen} \dicLink{völlur}
\dicEntry[vallarfoxgras] \dicTerm{vallar··fox·|gras} \dicIPA{{v}{a}{\textsubring{d}}{l}{a}{\textsubring{r}}{f}{\textopeno}{x}{s}{\r{g}}{r}{a}{s}} \dicPos{n}[2] \dicFlx{(‑grass, ‑grös)}[8] \dicFieldCat{bot.} \dicDirectTranslationCS{bojínek luční} \textit{(l.~{\textLA{Phleum pratense}})}
\dicEntry[vallarsveifgras] \dicTerm{vallar··sveif·|gras} \dicIPA{{v}{a}{\textsubring{d}}{l}{a}{\textsubring{r}}{s}{v}{ei}{v}{\r{g}}{r}{a}{s}} \dicPos{n}[2] \dicFlx{(‑grass, ‑grös)}[8] \dicFieldCat{bot.} \dicDirectTranslationCS{lipnice luční} \textit{(l.~{\textLA{Poa pratensis}})}  \dicsymPhoto\ 
\dicFigure{53715.jpg}{Vallarsveifgras}{Vallarsveifgras - Forest \& Kim [[p:2684;Starr]], Biolib, CC-BY}
\dicEntry[vallhumall] \dicTerm{vall··hum|all} \dicIPA{{v}{a}{\textsubring{d}}{\textsubring{l}}{h}{\textscy}{m}{a}{\textsubring{d}}{\textsubring{l}}} \dicPos{m}[6] \dicFlx{(‑als, ‑lar)}[27] \dicFieldCat{bot.} \dicDirectTranslationCS{řebříček, řebříček obecný} \textit{(l.~{\textLA{Achillea millefolium}})}  \dicsymPhoto\ 
\dicFigure{310.jpg}{Vallhumall}{Vallhumall - Zicha Ondřej, Biolib, Copyright/CC-BY-NC}
\dicEntry[vallhæra] \dicTerm{vall··hær|a} \dicIPA{{v}{a}{\textsubring{d}}{\textsubring{l}}{h}{a}{i}{r}{a}} \dicPos{f}[1] \dicFlx{(‑u, ‑ur)}[7] \dicFieldCat{bot.} \dicDirectTranslationCS{bika mnohokvětá} \textit{(l.~{\textLA{Luzula multiflora}})}
\dicEntry[valmenni] \dicTerm{val··menni} \dicIPA{{v}{a}{l}{m}{\textepsilon}{n}{\textsci}} \dicPos{n}[2] \dicFlx{(‑s, ‑)}[14] \dicDirectTranslationCS{ušlechtilý\,/\addthin vznešený člověk}
\dicEntry[valmúgi] \dicTerm{val··múg|i} \dicIPA{{v}{a}{l}{m}{u}{j}{\textsci}} \dicPos{m}[1] \dicFlx{(‑a)}[3] \dicLink{valmúi}
\dicEntry[valmúi] \dicTerm{val··mú|i}\dicTerm{, valmúgi} \dicIPA{{v}\-{a}\-{l}\-{m}\-{u}\-{\textsci}\-} \dicPos{m}[1] \dicFlx{(‑a)}[3] \dicFieldCat{bot.} \dicSynonym{draumsóley} \dicDirectTranslationCS{mák} \textit{(l.~{\textLA{Papaver}})}  \dicsymPhoto\ 
\dicFigure{60684.jpg}{Valmúi}{Valmúi - Zicha Ondřej, Biolib, Copyright/CC-BY-NC}
\dicEntry[valmynd] \dicTerm{val··mynd} \dicIPA{{v}{a}{l}{m}{\textsci}{n}{\textsubring{d}}} \dicPos{f}[7] \dicFlx{(‑ar, ‑ir)}[1] \dicFieldCat{poč.} \dicDirectTranslationCS{nabídka, menu}
\dicEntry[valmöguleiki] \dicTerm{val··mögu·leik|i} \dicIPA{{v}{a}{l}{m}{\oe}{\textbabygamma}{\textscy}{l}{ei}{\r{\textObardotlessj}}{\textsci}} \dicPos{m}[1] \dicFlx{(‑a, ‑ar)}[1] \dicDirectTranslationCS{možnost výběru\,/\addthin volby}
\dicEntry[valrönd] \dicTerm{val··|rönd} \dicIPA{{v}{a}{l}{r}{\oe}{n}{\textsubring{d}}} \dicPos{f}[8] \dicFlx{(‑randar, ‑rendur\,/\addthin ‑randir)}[9] \dicFieldCat{poč.} \dicDirectTranslationCS{menu, panel nabídek}
\dicEntry[vals] \dicTerm{vals} \dicIPA{{v}{a}{l}{s}} \dicPos{m}[4] \dicFlx{(‑, ‑ar)}[29] \textbf{1.} \dicFieldCat{hud.} \dicDirectTranslationCS{valčík}  \textbf{2.} \dicFieldCat{techn.} \dicDirectTranslationCS{váleček} \dicExampleIS{vals í ritvél} \dicExampleCS{váleček v~psacím stroji}
\dicEntry[valslöngva] \dicTerm{val··slöngv|a} \dicIPA{{v}{a}{l}{s}{\textsubring{d}}{l}{\oe i}{\ng}{\r{g}}{v}{a}} \dicPos{f}[1] \dicFlx{(‑u, ‑ur)}[7] \dicDirectTranslationCS{katapult}
\dicEntry[valt] \dicTerm{valt} \dicIPA{{v}{a}{\textsubring{l}}{\textsubring{d}}} \dicPos{v} \dicFlx{ind pf sg 1 pers} \dicLink{velta\smash{\textsuperscript{3}}}
\dicEntry[valtari] \dicTerm{valt··ar|i} \dicIPA{{v}{a}{\textsubring{l}}{\textsubring{d}}{a}{r}{\textsci}} \dicPos{m}[1] \dicFlx{(‑a, ‑ar)}[10] \dicDirectTranslationCS{válec} \dicIndirectTranslationCS{(stavební stroj)}
\dicEntry[valtur] \dicTerm{valtur} \dicIPA{{v}{a}{\textsubring{l}}{\textsubring{d}}{\textscy}{\textsubring{r}}} \dicPos{adj}[1] \dicFlx{(f völt)}[11] \dicSynonym{óstöðugur} \dicDirectTranslationCS{vratký, nepevný, nestabilní};  \dicPhraseIS{e‑að er valt í sessi} \dicLangCat{přen.} \dicDirectTranslationCS{(co) nestojí na pevných základech, (co) má vratkou půdu pod nohama};  \dicPhraseIS{e‑að stendur völtum fótum} \dicLangCat{přen.} \dicDirectTranslationCS{(co) stojí na vratkých základech};  \dicPhraseIS{vera valtur á fótunum} \dicDirectTranslationCS{mít vratkou chůzi}
\dicEntry[valur] \dicTerm{val|ur\smash{\textsuperscript{1}}} \dicIPA{{v}{a}{\textlengthmark}{l}{\textscy}{\textsubring{r}}} \dicPos{m}[9] \dicFlx{(‑s, ‑ir)}[9] \dicFieldCat{zool.} \dicSynonym{fálki} \dicDirectTranslationCS{raroh lovecký} \textit{(l.~{\textLA{Falco rusticolus}})}
\dicEntry[valur] \dicTerm{val|ur\smash{\textsuperscript{2}}} \dicIPA{{v}{a}{\textlengthmark}{l}{\textscy}{\textsubring{r}}} \dicPos{m}[9] \dicFlx{(‑s, ‑ir)}[9] \dicDirectTranslationCS{padlý} \dicIndirectTranslationCS{(v~bitvě)}
\dicEntry[vambar] \dicTerm{vambar} \dicIPA{{v}{a}{m}{\textsubring{b}}{a}{\textsubring{r}}} \dicPos{f} \dicFlx{sg gen} \dicLink{vömb}
\dicEntry[vamm] \dicTerm{vamm} \dicIPA{{v}{a}{m}{\textlengthmark}} \dicPos{n}[2] \dicFlx{(‑s)}[2] \dicPhraseIS{mega ekki vamm sitt vita} \dicLangCat{přen.} \dicDirectTranslationCS{mít pověst bez poskvrny}
\dicEntry[vammlaus] \dicTerm{vamm··laus} \dicIPA{{v}{a}{m}{l}{\oe i}{s}} \dicPos{adj}[5]\dicFlx{}[-1] \dicDirectTranslationCS{neposkvrněný, bezúhonný}
\dicEntry[vammleysi] \dicTerm{vamm··leysi} \dicIPA{{v}{a}{m}{l}{ei}{s}{\textsci}} \dicPos{n}[2] \dicFlx{(‑s)}[20] \dicDirectTranslationCS{neposkvrněnost, bezúhonnost}
\dicEntry[vampíra] \dicTerm{vampír|a}\dicTerm{, vampýra} \dicIPA{{v}\-{a}\-{\textsubring{m}}\-{\textsubring{b}}\-{i}\-{r}\-{a}\-} \dicPos{f}[1] \dicFlx{(‑u, ‑ur)}[7] \textbf{1.} \dicFieldCat{pov.} \dicDirectTranslationCS{vampýr, upír}  \textbf{2.} \dicFieldCat{zool.} \dicDirectTranslationCS{vampýr}
\dicEntry[vampýra] \dicTerm{vampýr|a} \dicIPA{{v}{a}{\textsubring{m}}{\textsubring{b}}{i}{r}{a}} \dicPos{f}[1] \dicFlx{(‑u, ‑ur)}[7] \dicLink{vampíra}
\dicEntry[van] \dicTerm{van-} \dicIPA{{v}{a}{\textsubring{n}}} \dicPos{predp} \dicDirectTranslationCS{ne-} \dicExampleIS{vankunnátta} \dicExampleCS{neznalost}
\dicEntry[vana] \dicTerm{van|a} \dicIPA{{v}{a}{\textlengthmark}{n}{a}} \dicPos{v}[1] \dicFlx{(‑aði)}[13] \dicFlx{acc} \dicSynonym{gelda} \dicDirectTranslationCS{(vy)kastrovat}
\dicEntry[vanabindandi] \dicTerm{vana··bind·andi} \dicIPA{{v}{a}{\textlengthmark}{n}{a}{\textsubring{b}}{\textsci}{n}{\textsubring{d}}{a}{n}{\textsubring{d}}{\textsci}} \dicPos{adj}[13] \dicFlx{indecl}[1] \dicDirectTranslationCS{návykový} \dicExampleIS{vanabindandi efni} \dicExampleCS{návyková látka}
\dicEntry[vanabundinn] \dicTerm{vana··bundinn} \dicIPA{{v}{a}{\textlengthmark}{n}{a}{\textsubring{b}}{\textscy}{n}{\textsubring{d}}{\textsci}{\textsubring{n}}} \dicPos{adj}[6]\dicFlx{}[-2] \dicSynonym{hefðbundinn} \dicDirectTranslationCS{navyklý, obvyklý, tradiční}
\dicEntry[vanadín] \dicTerm{vanadín}\dicTerm{, vanadíum} \dicIPA{{v}\-{a}\-{\textlengthmark}\-{n}\-{a}\-{\textsubring{d}}\-{i}\-{\textsubring{n}}\-} \dicPos{n}[2] \dicFlx{(‑s)}[2] \dicFieldCat{chem.} \dicDirectTranslationCS{vanad} \textit{(l.~{\textLA{V, Vanadium}})}
\dicEntry[vanadíum] \dicTerm{vanadíum} \dicIPA{{v}{a}{\textlengthmark}{n}{a}{\textsubring{d}}{i}{j}{\textscy}{\textsubring{m}}} \dicPos{n}[2] \dicFlx{(‑s)}[2] \dicLink{vanadín}
\dicEntry[vanaður] \dicTerm{van··aður} \dicIPA{{v}{a}{\textlengthmark}{n}{a}{ð}{\textscy}{\textsubring{r}}} \dicPos{adj}[3] \dicFlx{(f vönuð)}[2] \dicDirectTranslationCS{(vy)kastrovaný}
\dicEntry[vanafastur] \dicTerm{vana··|fastur} \dicIPA{{v}{a}{\textlengthmark}{n}{a}{f}{a}{s}{\textsubring{d}}{\textscy}{\textsubring{r}}} \dicPos{adj}[1] \dicFlx{(f ‑föst)}[11] \dicDirectTranslationCS{mající pevné zvyky}
\dicEntry[vanagangur] \dicTerm{vana··gang|ur} \dicIPA{{v}{a}{\textlengthmark}{n}{a}{\r{g}}{au}{\ng}{\r{g}}{\textscy}{\textsubring{r}}} \dicPos{m}[6] \dicFlx{(‑s)}[7] \dicPhraseIS{e‑að gengur sinn vanagang} \dicLangCat{přen.} \dicDirectTranslationCS{(co) jde po zaběhlých\,/\addthin vyšlapaných kolejích}
\dicEntry[Vanaheimur] \dicTerm{Vana··heim|ur} \dicIPA{{v}{a}{\textlengthmark}{n}{a}{h}{ei}{m}{\textscy}{\textsubring{r}}} \dicPos{m}[6] \dicFlx{(‑s)}[3] \dicFieldCat{myt.} \dicDirectTranslationCS{Vanaheim} \dicIndirectTranslationCS{(sídlo Vanů)}
\dicEntry[vanahugsun] \dicTerm{vana··hugs|un} \dicIPA{{v}{a}{\textlengthmark}{n}{a}{h}{\textscy}{x}{s}{\textscy}{\textsubring{n}}} \dicPos{f}[7] \dicFlx{(‑unar)}[9] \dicDirectTranslationCS{tradiční\,/\addthin konvenční\,/\addthin konzervativní myšlení}
\dicEntry[vanalega] \dicTerm{vana··lega} \dicIPA{{v}{a}{\textlengthmark}{n}{a}{l}{\textepsilon}{\textbabygamma}{a}} \dicPos{adv} \dicSynonym{venjulega} \dicDirectTranslationCS{obvykle, zpravidla, tradičně}
\dicEntry[vanalegur] \dicTerm{vana··legur} \dicIPA{{v}{a}{\textlengthmark}{n}{a}{l}{\textepsilon}{\textbabygamma}{\textscy}{\textsubring{r}}} \dicPos{adj}[1]\dicFlx{}[-8] \dicSynonym{venjulegur} \dicDirectTranslationCS{obvyklý, normální, tradiční}
\dicEntry[vananir] \dicTerm{vananir} \dicIPA{{v}{a}{\textlengthmark}{n}{a}{n}{\textsci}{\textsubring{r}}} \dicPos{f} \dicFlx{pl nom} \dicLink{vönun}
\dicEntry[vanbúnaður] \dicTerm{van··bú·nað|ur} \dicIPA{{v}{a}{n}{\textsubring{b}}{u}{n}{a}{ð}{\textscy}{\textsubring{r}}} \dicPos{m}[10] \dicFlx{(‑ar)}[9] \dicDirectTranslationCS{nedostatečná příprava, nepřipravenost};  \dicPhraseIS{e‑m er ekkert að vanbúnaði} \dicFlx{impers} \dicDirectTranslationCS{(komu) nic nebrání, (komu) nic nechybí} \dicExampleIS{Nú er okkur ekkert að vanbúnaði að leggja af stað í ferðina.} \dicExampleCS{A~nyní nám nic nebrání vyrazit na cestu.}
\dicEntry[vand] \dicTerm{vand-} \dicIPA{{v}{a}{n}{\textsubring{d}}} \dicPos{predp} \dicDirectTranslationCS{obtížně, komplikovaně}
\dicEntry[vanda] \dicTerm{vand|a} \dicIPA{{v}{a}{n}{\textsubring{d}}{a}} \dicPos{v}[1] \dicFlx{(‑aði)}[13] \dicFlx{acc} \textbf{1.} \dicDirectTranslationCS{dát si záležet, věnovat velkou péči\,/\addthin úsilí} \dicExampleIS{vanda verk sitt} \dicExampleCS{dát si záležet na své práci};  \dicPhraseIS{vanda sig} \dicDirectTranslationCS{snažit se, vynasnažit se, dát si záležet}  \textbf{2.} \dicPhraseIS{vanda um við e‑n} \dicDirectTranslationCS{(po)kárat (koho), napomenout\,/\addthin napomínat (koho)};  \dicIdiom{vandast}{ \dicPhraseIS{(nú) vandast málið}} \dicFlx{refl} \dicDirectTranslationCS{(nyní) se to komplikuje, (nyní) nastávají potíže}
\dicEntry[vandaður] \dicTerm{vand··aður} \dicIPA{{v}{a}{n}{\textsubring{d}}{a}{ð}{\textscy}{\textsubring{r}}} \dicPos{adj}[3] \dicFlx{(f vönduð)}[1] \textbf{1.} \dicSynonym*{vel gerður} \dicDirectTranslationCS{kvalitní, pečlivě udělaný} \dicExampleIS{vandaður gripur} \dicExampleCS{kvalitní věc}  \textbf{2.} \dicSynonym{valinkunnur} \dicDirectTranslationCS{počestný, poctivý} \dicExampleIS{vandaður til orðs og æðis} \dicExampleCS{počestný ve slovech a~činech}
\dicEntry[vandalaus] \dicTerm{vanda··laus} \dicIPA{{v}{a}{n}{\textsubring{d}}{a}{l}{\oe i}{s}} \dicPos{adj}[5]\dicFlx{}[-4] \textbf{1.} \dicSynonym{auðveldur} \dicDirectTranslationCS{jednoduchý, bezproblémový, snadný} \dicExampleIS{vandalaust verk} \dicExampleCS{jednoduchá práce};  \dicPhraseIS{það er vandalaust að (gera e‑ð)} \dicDirectTranslationCS{není problém (udělat (co))}  \textbf{2.} \dicSynonym{óskyldur} \dicDirectTranslationCS{nepříbuzný, cizí} \dicExampleIS{hjá vandalausum} \dicExampleCS{u~cizích}
\dicEntry[vandalítill] \dicTerm{vanda··|lítill} \dicIPA{{v}{a}{n}{\textsubring{d}}{a}{l}{i}{\textsubring{d}}{\textsci}{\textsubring{d}}{\textsubring{l}}} \dicPos{adj}[11] \dicFlx{(comp ‑minni, sup ‑minn\-st\-ur)}[10] \dicSynonym{vandalaus} \dicDirectTranslationCS{snadný, jednoduchý, bezproblémový}
\begin{xtolerant}{}{1pt}
\dicEntry[vandamaður] \dicTerm{vanda··|maður}\addthinS\dicIPA{{v}{a}{n}{\textsubring{d}}{a}{m}{a}{ð}{\textscy}{\textsubring{r}}}\addthinS\dicPos{m}[13]\addthinS\dicFlx{(‑manns, ‑menn)}[2] \dicSynonym{skyldmenni} \dicDirectTranslationCS{příbuzný, příbuzná} \dicExampleIS{Hún bauð aðeins nánum vandamönnum til veislunnar.} \dicExampleCS{Na oslavu pozvala jen nejbližší příbuzné.}
\end{xtolerant}
\dicEntry[vandamál] \dicTerm{vanda··mál} \dicsymFrequent\  \dicIPA{{v}{a}{n}{\textsubring{d}}{a}{m}{au}{\textsubring{l}}} \dicPos{n}[2] \dicFlx{(‑s, ‑)}[5] \dicDirectTranslationCS{problém, potíž, nesnáz};  \dicPhraseIS{það er ekki\,/\addthin ekkert vandamál að (gera e‑ð)} \dicDirectTranslationCS{není problém (udělat (co))} \dicExampleIS{Það er ekki vandamál að læra íslensku.} \dicExampleCS{Není problém se naučit islandštinu.}
\dicEntry[vandar] \dicTerm{vandar} \dicIPA{{v}{a}{n}{\textsubring{d}}{a}{\textsubring{r}}} \dicPos{m} \dicFlx{sg gen} \dicLink{vöndur}
\dicEntry[vandasamur] \dicTerm{vanda··|samur} \dicIPA{{v}{a}{n}{\textsubring{d}}{a}{s}{a}{m}{\textscy}{\textsubring{r}}} \dicPos{adj}[1] \dicFlx{(f ‑söm)}[2] \dicDirectTranslationCS{obtížný, problematický, těžký} \dicExampleIS{vandasamt verkefni} \dicExampleCS{obtížný úkol};  \dicPhraseIS{það er vandasamt að (gera e‑ð)} \dicDirectTranslationCS{je obtížné (udělat (co))}
\dicEntry[vandaverk] \dicTerm{vanda··verk} \dicIPA{{v}{a}{n}{\textsubring{d}}{a}{v}{\textepsilon}{\textsubring{r}}{\r{g}}} \dicPos{n}[2] \dicFlx{(‑s, ‑)}[5] \dicDirectTranslationCS{pracný úkol, namáhavá práce}
\dicEntry[vandfarinn] \dicTerm{vand··farinn} \dicIPA{{v}{a}{n}{\textsubring{d}}{f}{a}{r}{\textsci}{\textsubring{n}}} \dicPos{adj}[6]\dicFlx{}[-3] \dicDirectTranslationCS{problematický, ošidný, ošemetný} \dicExampleIS{vandfarinn vegur} \dicExampleCS{ošidná cesta}
\dicEntry[vandfenginn] \dicTerm{vand··fenginn} \dicIPA{{v}{a}{n}{\textsubring{d}}{f}{ei}{\textltailn}{\r{\textObardotlessj}}{\textsci}{\textsubring{n}}} \dicPos{adj}[6]\dicFlx{}[-2] \dicDirectTranslationCS{obtížně dostupný} \dicExampleIS{vandfengin vara} \dicExampleCS{obtížně dostupné zboží}
\dicEntry[vandfundinn] \dicTerm{vand··fundinn} \dicIPA{{v}{a}{n}{\textsubring{d}}{f}{\textscy}{n}{\textsubring{d}}{\textsci}{\textsubring{n}}} \dicPos{adj}[6]\dicFlx{}[-2] \dicDirectTranslationCS{obtížný\,/\addthin těžký k~nalezení} \dicExampleIS{Betri samstarfsmaður er vandfundinn.} \dicExampleCS{Lepšího spolupracovníka se dá těžko najít.}
\dicEntry[vandfýsi] \dicTerm{vand··fýs|i} \dicIPA{{v}{a}{n}{\textsubring{d}}{f}{i}{s}{\textsci}} \dicPos{f}[3] \dicFlx{(‑i)}[3] \dicDirectTranslationCS{vybíravost, náročnost}
\dicEntry[vandfýsinn] \dicTerm{vand··fýsinn} \dicIPA{{v}{a}{n}{\textsubring{d}}{f}{i}{s}{\textsci}{\textsubring{n}}} \dicPos{adj}[6]\dicFlx{}[-2] \dicDirectTranslationCS{vybíravý, náročný (divák ap.)} \dicExampleIS{vandfýsnir áhorfendur} \dicExampleCS{nároční diváci}
\dicEntry[vandi] \dicTerm{vand|i\smash{\textsuperscript{1}}} \dicsymFrequent\  \dicIPA{{v}{a}{n}{\textsubring{d}}{\textsci}} \dicPos{m}[1] \dicFlx{(‑a)}[3] \textbf{1.} \dicSynonym*{erfitt verkefni} \dicDirectTranslationCS{obtíž, obtížný úkol, náročná práce}  \textbf{2.} \dicSynonym{vandamál} \dicDirectTranslationCS{potíž, problém};  \dicPhraseIS{e‑m er vandi að höndum} \dicFlx{impers} \dicDirectTranslationCS{(kdo) stojí před problémem};  \dicPhraseIS{ráða fram úr vandanum} \dicDirectTranslationCS{vyřešit problémy};  \dicPhraseIS{setja e‑n í vanda} \dicDirectTranslationCS{dostat (koho) do potíží, způsobit (komu) potíže};  \dicPhraseIS{vanda ber að höndum} \dicFlx{impers} \dicDirectTranslationCS{nastávají potíže};  \dicPhraseIS{vera í vanda staddur} \dicDirectTranslationCS{mít trable};  \dicPhraseIS{það er hægur vandi að (gera e‑ð)} \dicDirectTranslationCS{je to snadná záležitost (udělat (co))}
\dicEntry[vandi] \dicTerm{vand|i\smash{\textsuperscript{2}}} \dicIPA{{v}{a}{n}{\textsubring{d}}{\textsci}} \dicPos{m}[1] \dicFlx{(‑a)}[3] \dicSynonym{vani} \dicDirectTranslationCS{zvyk, obyčej, tendence};  \dicPhraseIS{að vanda} \dicFlx{adv} \dicDirectTranslationCS{ze zvyku};  \dicPhraseIS{eiga vanda til að (gera e‑ð)} \dicDirectTranslationCS{mít ve zvyku (dělat (co))}
\dicEntry[vandi] \dicTerm{vandi\smash{\textsuperscript{3}}} \dicIPA{{v}{a}{n}{\textsubring{d}}{\textsci}} \dicPos{v} \dicFlx{ind pf sg 1 pers} \dicLink{venja}
\dicEntry[vandkvæði] \dicTerm{vand··kvæði} \dicIPA{{v}{a}{n}{\textsubring{d}}{k\smash{\textsuperscript{h}}}{v}{a}{i}{ð}{\textsci}} \dicPos{n}[2] \dicFlx{pl}[19] \dicSynonym*{erfiðleikar} \dicDirectTranslationCS{obtíže, nesnáze, komplikace};  \dicPhraseIS{það er vandkvæðum bundið að (gera e‑ð)} \dicDirectTranslationCS{(dělání (čeho)) je spojeno s~obtížemi};  \dicPhraseIS{það eru vandkvæði á að (gera e‑ð)} \dicDirectTranslationCS{je těžké\,/\addthin obtížné (udělat (co))}
\dicEntry[vandlátur] \dicTerm{vand··látur} \dicIPA{{v}{a}{n}{\textsubring{d}}{l}{au}{\textsubring{d}}{\textscy}{\textsubring{r}}} \dicPos{adj}[1]\dicFlx{}[-1] \dicSynonym{vandfýsinn} \dicDirectTranslationCS{náročný, vybíravý} \dicExampleIS{vera vandlátur á föt} \dicExampleCS{být náročný při výběru oblečení}
\dicEntry[vandlega] \dicTerm{vand··lega} \dicsymFrequent\  \dicIPA{{v}{a}{n}{\textsubring{d}}{l}{\textepsilon}{\textbabygamma}{a}} \dicPos{adv} \dicDirectTranslationCS{důkladně, pozorně, pečlivě} \dicExampleIS{velta þessu vandlega fyrir sér} \dicExampleCS{důkladně si to promyslet}
\dicEntry[vandlegur] \dicTerm{vand··legur} \dicIPA{{v}{a}{n}{\textsubring{d}}{l}{\textepsilon}{\textbabygamma}{\textscy}{\textsubring{r}}} \dicPos{adj}[1]\dicFlx{}[-8] \dicSynonym{rækilegur} \dicDirectTranslationCS{důkladný, pozorný, pečlivý} \dicExampleIS{vandleg athugun} \dicExampleCS{důkladné šetření}
\dicEntry[vandlæting] \dicTerm{vand··læt·ing} \dicIPA{{v}{a}{n}{\textsubring{d}}{l}{a}{i}{\textsubring{d}}{i}{\ng}{\r{g}}} \dicPos{f}[4] \dicFlx{(‑ar, ‑ar)}[5] \dicDirectTranslationCS{důtka, (po)kárání, moralizování}
\dicEntry[vandmeðfarinn] \dicTerm{vand··með·farinn} \dicIPA{{v}{a}{n}{\textsubring{d}}{m}{\textepsilon}{ð}{f}{a}{r}{\textsci}{\textsubring{n}}} \dicPos{adj}[6]\dicFlx{}[-3] \dicDirectTranslationCS{ožehavý, choulostivý, delikátní} \dicExampleIS{vandmeðfarið efni} \dicExampleCS{ožehavé téma}
\dicEntry[vandrataður] \dicTerm{vand··|rat·aður} \dicIPA{{v}{a}{n}{\textsubring{d}}{r}{a}{\textsubring{d}}{a}{ð}{\textscy}{\textsubring{r}}} \dicPos{adj}[3] \dicFlx{(f ‑rötuð)}[1] \dicDirectTranslationCS{těžký\,/\addthin obtížný na orientaci} \dicExampleIS{vandrataður vegur} \dicExampleCS{cesta obtížná na orientaci}
\dicEntry[vandræðabarn] \dicTerm{vand·ræða··|barn} \dicIPA{{v}{a}{n}{\textsubring{d}}{r}{a}{i}{ð}{a}{\textsubring{b}}{a}{r}{\textsubring{d}}{\textsubring{n}}} \dicPos{n}[2] \dicFlx{(‑barns, ‑börn)}[8] \dicDirectTranslationCS{problémové dítě}
\dicEntry[vandræðagemsi] \dicTerm{vand·ræða··gems|i} \dicIPA{{v}{a}{n}{\textsubring{d}}{r}{a}{i}{ð}{a}{\r{\textObardotlessj}}{\textepsilon}{m}{s}{\textsci}} \dicPos{m}[1] \dicFlx{(‑a, ‑ar)}[1] \dicDirectTranslationCS{potížista, potížistka, rebelant(ka)}
\dicEntry[vandræðalegur] \dicTerm{vand·ræða··legur} \dicsymFrequent\  \dicIPA{{v}{a}{n}{\textsubring{d}}{r}{a}{i}{ð}{a}{l}{\textepsilon}{\textbabygamma}{\textscy}{\textsubring{r}}} \dicPos{adj}[1]\dicFlx{}[-8] \textbf{1.} \dicDirectTranslationCS{rozpačitý, (jsoucí) v~rozpacích, (jsoucí) nesvůj}  \textbf{2.} \dicDirectTranslationCS{rozpačitý, trapný, prekérní (situace ap.)} \dicExampleIS{vandræðaleg þögn} \dicExampleCS{trapné ticho}

\dicEntry[vandræðamaður] \textls[25]{\dicTerm{vand·ræða··|maður} \dicIPA{{v}{a}{n}{\textsubring{d}}{r}{a}{i}{ð}{a}{m}{a}{ð}{\textscy}{\textsubring{r}}} \dicPos{m}[13] \dicFlx{(‑manns, ‑menn)}[2]} \textls[-5]{\dicDirectTranslationCS{problémový člověk, potížista, potížistka, bu\-řič\-(ka)}}

\dicEntry[vandræðamál] \dicTerm{vand·ræða··mál} \dicIPA{{v}{a}{n}{\textsubring{d}}{r}{a}{i}{ð}{a}{m}{au}{\textsubring{l}}} \dicPos{n}[2] \dicFlx{(‑s, ‑)}[5] \dicDirectTranslationCS{ožehavá\,/\addthin problematická  záležitost}
\dicEntry[vandræði] \dicTerm{vand··ræði} \dicsymFrequent\  \dicIPA{{v}{a}{n}{\textsubring{d}}{r}{a}{i}{ð}{\textsci}} \dicPos{n}[2] \dicFlx{pl}[19] \dicDirectTranslationCS{potíže, nesnáze, problémy, trable} \dicExampleIS{Var kötturinn til vandræða?} \dicExampleCS{Byly s~tou kočičkou potíže?};  \dicPhraseIS{lenda í vandræðum} \dicDirectTranslationCS{dostat se do potíží};  \dicPhraseIS{vera í vandræðum með e‑ð} \dicDirectTranslationCS{mít s~(čím) potíže}
\dicEntry[vandséður] \dicTerm{vand··séður} \dicIPA{{v}{a}{n}{\textsubring{d}}{s}{j}{\textepsilon}{ð}{\textscy}{\textsubring{r}}} \dicPos{adj}[2]\dicFlx{}[-12] \dicPhraseIS{e‑að er vandséð} \dicDirectTranslationCS{(co) není zřejmé, (co) je obtížné předpovědět}
\dicEntry[vandskilinn] \dicTerm{vand··skilinn} \dicIPA{{v}{a}{n}{\textsubring{d}}{s}{\r{\textObardotlessj}}{\textsci}{l}{\textsci}{\textsubring{n}}} \dicPos{adj}[6]\dicFlx{}[-7] \dicSynonym{torskilinn} \dicDirectTranslationCS{obtížně srozumitelný}
\dicEntry[vandur] \dicTerm{vandur} \dicIPA{{v}{a}{n}{\textsubring{d}}{\textscy}{\textsubring{r}}} \dicPos{adj}[2] \dicFlx{(f vönd)}[15] \dicPhraseIS{eiga úr vöndu að ráða} \dicDirectTranslationCS{být bezradný, nebýt schopen se rozhodnout};  \dicPhraseIS{vera vandur að virðingu sinni} \dicDirectTranslationCS{dbát o~svou čest, být svědomitý}
\dicEntry[vandvirkni] \dicTerm{vand··virkn|i} \dicIPA{{v}{a}{n}{\textsubring{d}}{v}{\textsci}{\textsubring{r}}{\r{g}}{n}{\textsci}} \dicPos{f}[3] \dicFlx{(‑i)}[3] \dicDirectTranslationCS{důkladnost, pečlivost, preciznost}
\dicEntry[vandvirkur] \dicTerm{vand··virkur} \dicIPA{{v}{a}{n}{\textsubring{d}}{v}{\textsci}{\textsubring{r}}{\r{g}}{\textscy}{\textsubring{r}}} \dicPos{adj}[1]\dicFlx{}[-1] \dicDirectTranslationCS{důkladný, pečlivý, precizní}
\dicEntry[vanefni] \dicTerm{van··efni} \dicIPA{{v}{a}{\textlengthmark}{n}{\textepsilon}{\textsubring{b}}{n}{\textsci}} \dicPos{n}[2] \dicFlx{(‑s, ‑)}[14] \dicDirectTranslationCS{nedostatek, nedostatečné prostředky (finanční ap.)}
\dicEntry[vanfær] \dicTerm{van··fær} \dicIPA{{v}{a}{n}{f}{a}{i}{\textsubring{r}}} \dicPos{adj}[5] \dicFlx{(f ‑)}[10] \textbf{1.} \dicFlx{f} \dicSynonym{barnshafandi} \dicDirectTranslationCS{těhotná, gravidní}  \textbf{2.} \dicSynonym{ófær} \dicDirectTranslationCS{nevhodný, nezpůsobilý (k~těžké práci ap.)}
\dicEntry[vanga] \dicTerm{vang|a} \dicIPA{{v}{au}{\ng}{\r{g}}{a}} \dicPos{v}[1] \dicFlx{(‑aði)}[13] \dicFlx{acc} \dicDirectTranslationCS{tancovat\,/\addthin tančit ploužák} \dicExampleIS{vanga stúlku} \dicExampleCS{tancovat ploužák s~dívkou}
\dicEntry[vangadans] \dicTerm{vanga··dans} \dicIPA{{v}{au}{\ng}{\r{g}}{a}{\textsubring{d}}{a}{n}{s}} \dicPos{m}[4] \dicFlx{(‑, ‑ar)}[29] \dicDirectTranslationCS{ploužák}
\dicEntry[vangamynd] \dicTerm{vanga··mynd} \dicIPA{{v}{au}{\ng}{\r{g}}{a}{m}{\textsci}{n}{\textsubring{d}}} \dicPos{f}[7] \dicFlx{(‑ar, ‑ir)}[1] \dicDirectTranslationCS{profilový portrét\,/\addthin obrázek}
\dicEntry[vangaskegg] \dicTerm{vanga··skegg} \dicIPA{{v}{au}{\ng}{\r{g}}{a}{s}{\r{\textObardotlessj}}{\textepsilon}{\r{g}}} \dicPos{n}[2] \dicFlx{(‑s, ‑)}[13] \dicDirectTranslationCS{kotlety, licousy}
\dicEntry[vangaveltur] \dicTerm{vanga··veltur} \dicsymFrequent\  \dicIPA{{v}{au}{\ng}{\r{g}}{a}{v}{\textepsilon}{\textsubring{l}}{\textsubring{d}}{\textscy}{\textsubring{r}}} \dicPos{f}[1] \dicFlx{pl}[24] \dicDirectTranslationCS{přemítání, hloubání, spekulování};  \dicPhraseIS{vera með vangaveltur út af e‑u} \dicDirectTranslationCS{přemítat o~(čem)}
\dicEntry[vangá] \dicTerm{van··gá} \dicIPA{{v}{a}{n}{\r{g}}{au}} \dicPos{f}[4] \dicFlx{(‑r)}[19] \dicSynonym{yfirsjón} \dicDirectTranslationCS{přehlédnutí, nepozornost, nedbalost};  \dicPhraseIS{af vangá} \dicFlx{adv} \dicDirectTranslationCS{přehlédnutím, z~nepozornosti\,/\addthin nedbalosti}
\dicEntry[vangefinn] \dicTerm{van··gefinn} \dicIPA{{v}{a}{n}{\r{\textObardotlessj}}{\textepsilon}{v}{\textsci}{\textsubring{n}}} \dicPos{adj}[6]\dicFlx{}[-6] \dicSynonym{þroskaheftur} \dicDirectTranslationCS{(duševně) zaostalý\,/\addthin postižený, retardovaný}
\dicEntry[vangeta] \dicTerm{van··get|a} \dicIPA{{v}{a}{n}{\r{\textObardotlessj}}{\textepsilon}{\textsubring{d}}{a}} \dicPos{f}[1] \dicFlx{(‑u)}[5] \dicSynonym{getuleysi} \dicDirectTranslationCS{neschopnost, nemohoucnost}
\dicEntry[vangi] \dicTerm{vang|i} \dicsymFrequent\  \dicIPA{{v}{au}{\textltailn}{\r{\textObardotlessj}}{\textsci}} \dicPos{m}[1] \dicFlx{(‑a, ‑ar)}[8] \dicSynonym{kinn} \dicDirectTranslationCS{tvář, líce} \dicExampleIS{kyssa hana á vangann} \dicExampleCS{políbit ji na tvář};  \dicPhraseIS{vera rjóður í vöngum} \dicDirectTranslationCS{mít červené tváře}
\dicEntry[vangoldinn] \dicTerm{van··goldinn}\dicTerm{, vangreiddur} \dicIPA{{v}\-{a}\-{n}\-{\r{g}}\-{\textopeno}\-{l}\-{\textsubring{d}}\-{\textsci}\-{\textsubring{n}}\-} \dicPos{adj}[6]\dicFlx{}[-2] \dicDirectTranslationCS{nezaplacený, nesplacený}
\dicEntry[vangreiddur] \dicTerm{van··greiddur} \dicIPA{{v}{a}{n}{\r{g}}{r}{ei}{\textsubring{d}}{\textscy}{\textsubring{r}}} \dicPos{adj}[2]\dicFlx{}[-21] \dicLink{vangoldinn}
\dicEntry[vangæfur] \dicTerm{van··gæfur} \dicIPA{{v}{a}{n}{\r{\textObardotlessj}}{a}{i}{v}{\textscy}{\textsubring{r}}} \dicPos{adj}[1]\dicFlx{}[-1] \dicSynonym{veill} \dicDirectTranslationCS{neduživý, chatrný};  \dicPhraseIS{vera vangæfur til heilsunnar} \dicDirectTranslationCS{mít chatrné zdraví}
\dicEntry[vanhaga] \dicTerm{van··hag|a} \dicIPA{{v}{a}{\textlengthmark}{\textsubring{n}}{h}{a}{\textbabygamma}{a}} \dicPos{v}[1] \dicFlx{(‑aði)}[79] \dicFlx{impers} \dicPhraseIS{e‑n vanhagar um e‑ð} \dicDirectTranslationCS{(komu co) schází, (komu) se (čeho) nedostává} \dicExampleIS{Hvað vanhagar þig helst um?} \dicExampleCS{Co ti nejvíce schází?}
\dicEntry[vanheill] \dicTerm{van··heill} \dicIPA{{v}{a}{\textlengthmark}{\textsubring{n}}{h}{ei}{\textsubring{d}}{\textsubring{l}}} \dicPos{adj}[8]\dicFlx{}[-1] \dicSynonym{sjúkur} \dicDirectTranslationCS{churavý, stonavý};  \dicPhraseIS{vera andlega vanheill} \dicDirectTranslationCS{být duševně nemocný}
\dicEntry[vanheilsa] \dicTerm{van··heils|a} \dicIPA{{v}{a}{\textlengthmark}{\textsubring{n}}{h}{ei}{l}{s}{a}} \dicPos{f}[1] \dicFlx{(‑u)}[5] \dicDirectTranslationCS{chatrné zdraví, churavost} \dicExampleIS{eiga við vanheilsu að stríða} \dicExampleCS{potýkat se s~chatrným zdravím}
\dicEntry[vanhelga] \dicTerm{van··helg|a} \dicIPA{{v}{a}{\textlengthmark}{\textsubring{n}}{h}{\textepsilon}{l}{\r{g}}{a}} \dicPos{v}[1] \dicFlx{(‑aði)}[1] \dicFlx{acc} \dicSynonym{saurga} \dicDirectTranslationCS{zneuctít, znesvětit (památku ap.)} \dicExampleIS{vanhelga kirkju} \dicExampleCS{znesvětit kostel}
\dicEntry[vanhelgun] \dicTerm{van··helg|un} \dicIPA{{v}{a}{\textlengthmark}{\textsubring{n}}{h}{\textepsilon}{l}{\r{g}}{\textscy}{\textsubring{n}}} \dicPos{f}[7] \dicFlx{(‑unar)}[9] \dicDirectTranslationCS{zneuctění, znesvěcení}
\dicEntry[vanhirða] \dicTerm{van··hirð|a} \dicIPA{{v}{a}{\textlengthmark}{\textsubring{n}}{h}{\textsci}{r}{ð}{a}} \dicPos{f}[1] \dicFlx{(‑u)}[5] \dicSynonym{hirðuleysi} \dicDirectTranslationCS{zanedbání, zanedbávání, opominutí, opomíjení}
\dicEntry[vanhugsaður] \dicTerm{van··hugs·|aður} \dicIPA{{v}{a}{\textlengthmark}{\textsubring{n}}{h}{\textscy}{x}{s}{a}{ð}{\textscy}{\textsubring{r}}} \dicPos{adj}[3] \dicFlx{(f ‑uð)}[3] \dicDirectTranslationCS{neuvážený, nepromyšlený, nepředložený} \dicExampleIS{vanhugsuð áform} \dicExampleCS{nepromyšlené plány}
\dicEntry[vanhyggja] \dicTerm{van··hyggj|a} \dicIPA{{v}{a}{\textlengthmark}{\textsubring{n}}{h}{\textsci}{\r{\textObardotlessj}}{a}} \dicPos{f}[1] \dicFlx{(‑u)}[5] \dicDirectTranslationCS{neopatrnost, nepozornost}
\dicEntry[vanhæfi] \dicTerm{van··hæfi} \dicIPA{{v}{a}{\textlengthmark}{\textsubring{n}}{h}{a}{i}{v}{\textsci}} \dicPos{n}[2] \dicFlx{(‑s)}[20] \dicFieldCat{práv.} \dicDirectTranslationCS{nezpůsobilost} \dicAntonym{gerhæfi}
\dicEntry[vanhæfni] \dicTerm{van··hæfn|i} \dicIPA{{v}{a}{\textlengthmark}{\textsubring{n}}{h}{a}{i}{\textsubring{b}}{n}{\textsci}} \dicPos{f}[3] \dicFlx{(‑i)}[3] \textbf{1.} \dicDirectTranslationCS{nekvalifikovanost, nekompetentnost}  \textbf{2.} \dicSynonym{vanhæfi} \dicDirectTranslationCS{nezpůsobilost}
\dicEntry[vanhæfur] \dicTerm{van··hæfur} \dicIPA{{v}{a}{\textlengthmark}{\textsubring{n}}{h}{a}{i}{v}{\textscy}{\textsubring{r}}} \dicPos{adj}[1]\dicFlx{}[-1] \textbf{1.} \dicSynonym{óhæfur} \dicDirectTranslationCS{nezpůsobilý, nekvalifikovaný, nekompetentní}  \textbf{2.} \dicFieldCat{práv.} \dicDirectTranslationCS{nezpůsobilý}
\dicEntry[vanhöld] \dicTerm{van··höld} \dicIPA{{v}{a}{\textlengthmark}{\textsubring{n}}{h}{\oe}{l}{\textsubring{d}}} \dicPos{n}[2] \dicFlx{pl}[9] \dicSynonym{tjón} \dicDirectTranslationCS{ztráta, škoda, úbytek (na dobytku ap.)};  \dicPhraseIS{það eru vanhöld á e‑u} \dicFlx{impers} \dicDirectTranslationCS{jsou ztráty na (čem)}
\dicEntry[vani] \dicTerm{van|i} \dicIPA{{v}{a}{\textlengthmark}{n}{\textsci}} \dicPos{m}[1] \dicFlx{(‑a, ‑ar)}[8] \dicSynonym{venja\smash{\textsuperscript{1}}} \dicDirectTranslationCS{zvyk, obyčej};  \dicPhraseIS{e‑að kemst upp í vana} \dicFlx{refl} \dicDirectTranslationCS{zvykne se na (co)};  \dicPhraseIS{bregða vana sínum} \dicDirectTranslationCS{změnit své zvyky};  \dicPhraseIS{leggja e‑ð í vana sinn} \dicDirectTranslationCS{mít (co) ve zvyku}
\dicEntry[vanið] \dicTerm{vanið} \dicIPA{{v}{a}{\textlengthmark}{n}{\textsci}{\texttheta}} \dicPos{v} \dicFlx{supin} \dicLink{venja}
\dicEntry[vanilja] \dicTerm{vanilj|a} \dicIPA{{v}{a}{\textlengthmark}{n}{\textsci}{l}{j}{a}} \dicPos{f}[1] \dicFlx{(‑u)}[5] \dicLink{vanilla}
\dicEntry[vanilla] \dicTerm{vanill|a}\dicTerm{, vanilja} \dicIPA{{v}\-{a}\-{\textlengthmark}\-{n}\-{\textsci}\-{l}\-{a}\-} \dicPos{f}[1] \dicFlx{(‑u)}[5] \dicDirectTranslationCS{vanilka}
\dicEntry[vankaður] \dicTerm{vank··aður} \dicIPA{{v}{au}{\r{\ng}}{\r{g}}{a}{ð}{\textscy}{\textsubring{r}}} \dicPos{adj}[3] \dicFlx{(f vönkuð)}[1] \dicDirectTranslationCS{mající závrať, trpící závratí}
\dicEntry[vankantur] \dicTerm{van··kant|ur} \dicIPA{{v}{a}{\ng}{k\smash{\textsuperscript{h}}}{a}{\textsubring{n}}{\textsubring{d}}{\textscy}{\textsubring{r}}} \dicPos{m}[6] \dicFlx{(‑s, ‑ar)}[5] \dicSynonym{misfella} \dicDirectTranslationCS{nedostatek, nedokonalost} \dicExampleIS{Á frumvarpinu eru alvarlegir vankantar.} \dicExampleCS{V~návrhu se vyskytují závažné nedostatky.}
\dicEntry[vankunnandi] \dicTerm{van··kunn·andi} \dicIPA{{v}{a}{\ng}{k\smash{\textsuperscript{h}}}{\textscy}{n}{a}{n}{\textsubring{d}}{\textsci}} \dicPos{adj}[13] \dicFlx{indecl}[1] \dicDirectTranslationCS{neznalý, neinformovaný, (jsoucí) ne v~obraze}
\dicEntry[vankunnátta] \dicTerm{van··kunn·átt|a} \dicIPA{{v}{a}{\ng}{k\smash{\textsuperscript{h}}}{\textscy}{n}{au}{h}{\textsubring{d}}{a}} \dicPos{f}[1] \dicFlx{(‑u)}[5] \dicDirectTranslationCS{neznalost, neinformovanost, nevědomost} \dicExampleIS{vankunnátta í e‑u} \dicExampleCS{neznalost (čeho)}
\dicEntry[vanlíðan] \dicTerm{van··líðan}\dicTerm{, vanlíðun} \dicIPA{{v}\-{a}\-{n}\-{l}\-{i}\-{ð}\-{a}\-{\textsubring{n}}\-} \dicPos{f}[4] \dicFlx{(‑ar)}[3] \dicSynonym*{slæm líðan} \dicDirectTranslationCS{nelibost, malátnost, nepříjemné pocity (zdravotní obtíže ap.)} \dicExampleIS{þjást af vanlíðan} \dicExampleCS{cítit se nepříjemně}
\dicEntry[vanlíðun] \dicTerm{van··líð|un} \dicIPA{{v}{a}{n}{l}{i}{ð}{\textscy}{\textsubring{n}}} \dicPos{f}[7] \dicFlx{(‑unar)}[9] \dicLink{vanlíðan}
\dicEntry[vanmáttarkennd] \dicTerm{van·máttar··kennd} \dicIPA{{v}{a}{n}{m}{au}{h}{\textsubring{d}}{a}{\textsubring{r}}{c\smash{\textsuperscript{h}}}{\textepsilon}{n}{\textsubring{d}}} \dicPos{f}[7] \dicFlx{(‑ar)}[3] \dicFieldCat{psych.} \dicDirectTranslationCS{komplex méněcennosti}
\dicEntry[vanmáttur] \dicTerm{van··mátt|ur} \dicIPA{{v}{a}{n}{m}{au}{h}{\textsubring{d}}{\textscy}{\textsubring{r}}} \dicPos{m}[12] \dicFlx{(‑ar)}[8] \dicDirectTranslationCS{nemohoucnost, bezmocnost} \dicExampleIS{finna til vanmáttar síns andspænis náttúruöflunum} \dicExampleCS{cítit svou bezmocnost vůči silám přírody}
\dicEntry[vanmeta] \dicTerm{van··|meta} \dicIPA{{v}{a}{n}{m}{\textepsilon}{\textsubring{d}}{a}} \dicPos{v}[6] \dicFlx{(‑met, ‑mat, ‑mátum, ‑mæti, ‑metið)}[16] \dicFlx{acc} \dicDirectTranslationCS{podcenit, podceňovat}
\dicEntry[vanmetinn] \dicTerm{van··metinn} \dicIPA{{v}{a}{n}{m}{\textepsilon}{\textsubring{d}}{\textsci}{\textsubring{n}}} \dicPos{adj}[6]\dicFlx{}[-2] \dicDirectTranslationCS{podceněný, nedoceněný, zneuznaný}
\dicEntry[vann] \dicTerm{vann} \dicIPA{{v}{a}{n}{\textlengthmark}} \dicPos{v} \dicFlx{ind pf sg 1 pers} \dicLink{vinna\smash{\textsuperscript{2}}}
\dicEntry[vannærður] \dicTerm{van··nærður} \dicIPA{{v}{a}{n}{\textlengthmark}{a}{i}{r}{ð}{\textscy}{\textsubring{r}}} \dicPos{adj}[2]\dicFlx{}[-1] \dicDirectTranslationCS{podvyživený} \dicExampleIS{vannærð börn} \dicExampleCS{podvyživené děti}
\dicEntry[vannæring] \dicTerm{van··nær·ing} \dicIPA{{v}{a}{n}{\textlengthmark}{a}{i}{r}{i}{\ng}{\r{g}}} \dicPos{f}[4] \dicFlx{(‑ar)}[7] \dicSynonym{næringarskortur} \dicDirectTranslationCS{podvýživa}
\dicEntry[vanrækja] \dicTerm{van··ræk|ja} \dicIPA{{v}{a}{n}{r}{a}{i}{\r{\textObardotlessj}}{a}} \dicPos{v}[2] \dicFlx{(‑ti, ‑t)}[25] \dicFlx{acc} \dicDirectTranslationCS{opomenout, opomíjet, zanedbat, zanedbávat} \dicExampleIS{vanrækja fjölskylduna} \dicExampleCS{zanedbávat rodinu}
\dicEntry[vanræksla] \dicTerm{van··ræksl|a} \dicIPA{{v}{a}{n}{r}{a}{i}{x}{s}{\textsubring{d}}{l}{a}} \dicPos{f}[1] \dicFlx{(‑u)}[5] \dicDirectTranslationCS{zanedbání, zanedbávání, opominutí, opomíjení} \dicExampleIS{ófyrirgefanleg vanræksla} \dicExampleCS{neomluvitelné zanedbání}
\dicEntry[vansalaus] \dicTerm{vansa··laus} \dicIPA{{v}{a}{n}{s}{a}{l}{\oe i}{s}} \dicPos{adj}[5]\dicFlx{}[-1] \dicDirectTranslationCS{(jsoucí) bez ostudy} \dicIndirectTranslationCS{(nejčastěji v~záporu)} \dicExampleIS{Það er ekki vansalaust.} \dicExampleCS{Je to ostudné.}
\dicEntry[vansi] \dicTerm{vans|i} \dicIPA{{v}{a}{n}{s}{\textsci}} \dicPos{m}[1] \dicFlx{(‑a)}[3] \dicSynonym{skömm\smash{\textsuperscript{1}}} \dicDirectTranslationCS{ostuda, hanba};  \dicPhraseIS{e‑að er e‑m til vansa} \dicDirectTranslationCS{(co) je (komu) na ostudu} \dicExampleIS{Svona framkoma er honum til vansa.} \dicExampleCS{Takové chování mu je na ostudu.}
\dicEntry[vanskapaður] \dicTerm{van··|skap·aður} \dicIPA{{v}{a}{n}{s}{\r{g}}{a}{\textsubring{b}}{a}{ð}{\textscy}{\textsubring{r}}} \dicPos{adj}[3] \dicFlx{(f ‑sköpuð)}[1] \dicDirectTranslationCS{znetvořený, (z)deformovaný}
\dicEntry[vanskil] \dicTerm{van··skil} \dicIPA{{v}{a}{n}{s}{\r{\textObardotlessj}}{\textsci}{\textsubring{l}}} \dicPos{n}[2] \dicFlx{pl}[1] \textbf{1.} \dicDirectTranslationCS{neplacení, nesplácení};  \dicPhraseIS{lenda í vanskilum með e‑ð} \dicDirectTranslationCS{zpozdit se se splácením}  \textbf{2.} \dicDirectTranslationCS{platební neschopnost\,/\addthin potíže}
\dicEntry[vanstilling] \dicTerm{van··still·ing} \dicIPA{{v}{a}{n}{s}{\textsubring{d}}{\textsci}{\textsubring{d}}{l}{i}{\ng}{\r{g}}} \dicPos{f}[4] \dicFlx{(‑ar)}[7] \textbf{1.} \dicSynonym*{lítil sjálfstjórn} \dicDirectTranslationCS{nedostatek sebeovládání, neukázněnost}  \textbf{2.} \dicSynonym*{röng stilling} \dicDirectTranslationCS{špatné nastavení}
\dicEntry[vanstilltur] \dicTerm{van··stilltur} \dicIPA{{v}{a}{n}{s}{\textsubring{d}}{\textsci}{\textsubring{l}}{\textsubring{d}}{\textscy}{\textsubring{r}}} \dicPos{adj}[1]\dicFlx{}[-10] \dicSynonym{óstilltur} \dicDirectTranslationCS{nevyrovnaný, neukázněný, neovládající se}
\dicEntry[vansvefta] \dicTerm{van··svefta} \dicIPA{{v}{a}{n}{s}{v}{\textepsilon}{f}{\textsubring{d}}{a}} \dicPos{adj}[13] \dicFlx{indecl}[1] \dicDirectTranslationCS{nevyspalý, nedospalý}
\dicEntry[vansæll] \dicTerm{van··sæll} \dicIPA{{v}{a}{n}{s}{a}{i}{\textsubring{d}}{\textsubring{l}}} \dicPos{adj}[8]\dicFlx{}[-1] \dicDirectTranslationCS{nešťastný, nespokojený}
\dicEntry[vansæmd] \dicTerm{van··sæmd} \dicIPA{{v}{a}{n}{s}{a}{i}{m}{\textsubring{d}}} \dicPos{f}[7] \dicFlx{(‑ar)}[3] \dicSynonym{skömm\smash{\textsuperscript{1}}} \dicDirectTranslationCS{ostuda, han\-ba, potupa} \dicExampleIS{Þetta er honum til vansæmdar.} \dicExampleCS{Je mu to na ostudu.}
\dicEntry[vant] \dicTerm{vant} \dicIPA{{v}{a}{\textsubring{n}}{\textsubring{d}}} \dicPos{adj}[13]\dicFlx{}[-4] \dicFlx{n} \dicPhraseIS{e‑m verður orðs vant} \dicFlx{impers} \dicDirectTranslationCS{(kdo) nemá slov, (kdo) nenachází vhodná slova};  \dicPhraseIS{vera vant við látinn} \dicDirectTranslationCS{mít napilno, být zaneprázdněný}
\dicEntry[vanta] \dicTerm{vant|a} \dicsymFrequent\  \dicIPA{{v}{a}{\textsubring{n}}{\textsubring{d}}{a}} \dicPos{v}[1] \dicFlx{(‑aði)}[67] \dicFlx{impers} \dicPhraseIS{e‑n vantar e‑ð} \dicDirectTranslationCS{(komu co) chybí, (kdo co) potřebuje, (komu co) schází, (kdo co) postrádá} \dicExampleIS{Mig vantar peninga.} \dicExampleCS{Schází mi peníze.};  \dicPhraseIS{klukkuna vantar kortér í fimm} \dicDirectTranslationCS{je tři čtvrtě na pět};  \dicPhraseIS{það vantar eitthvað í e‑n} \dicDirectTranslationCS{(kdo) není úplně normální};  \dicPhraseIS{það vantar herslumuninn} \dicDirectTranslationCS{chybí jen málo, schází jen kousek}
\dicEntry[vantalaður] \dicTerm{van··|tal·aður} \dicIPA{{v}{a}{n}{t\smash{\textsuperscript{h}}}{a}{l}{a}{ð}{\textscy}{\textsubring{r}}} \dicPos{adj}[3] \dicFlx{(f ‑töluð)}[1] \dicDirectTranslationCS{nevyříkaný, neprojednaný};  \dicPhraseIS{eiga e‑ð vantalað við e‑n} \dicDirectTranslationCS{mít s~(kým co) nevyříkané, mít s~(kým co) na projednání}
\dicEntry[vantraust] \dicTerm{van··traust} \dicIPA{{v}{a}{n}{t\smash{\textsuperscript{h}}}{r}{\oe i}{s}{\textsubring{d}}} \dicPos{n}[2] \dicFlx{(‑s)}[2] \textbf{1.} \dicSynonym{ótrú} \dicDirectTranslationCS{nedůvěra}  \textbf{2.} \dicSynonym*{vantraustsyfirlýsing} \dicDirectTranslationCS{vyjádření\,/\addthin projev nedůvěry};  \dicPhraseIS{bera fram vantraust á e‑ð} \dicDirectTranslationCS{vyjadřovat nedůvěru (čemu) (vládě ap.)};  \dicPhraseIS{lýsa vantrausti á e‑n} \dicDirectTranslationCS{projevovat nedůvěru vůči (komu)}
\dicEntry[vantreysta] \dicTerm{van··treyst|a} \dicIPA{{v}{a}{n}{t\smash{\textsuperscript{h}}}{r}{ei}{s}{\textsubring{d}}{a}} \dicPos{v}[2] \dicFlx{(‑i, ‑)}[17] \dicFlx{dat} \dicDirectTranslationCS{mít nedůvěru, nedůvěřovat, nevěřit} \dicExampleIS{vantreysta orðum hans} \dicExampleCS{nedůvěřovat jeho slovům}
\dicEntry[vantrú] \dicTerm{van··trú} \dicIPA{{v}{a}{n}{t\smash{\textsuperscript{h}}}{r}{u}} \dicPos{f}[4] \dicFlx{(‑ar)}[27] \dicSynonym{vantraust} \dicDirectTranslationCS{nevíra, skepse};  \dicPhraseIS{hafa vantrú á e‑m} \dicDirectTranslationCS{nevěřit (komu)}
\dicEntry[vantrúaður] \dicTerm{van··trú·|aður} \dicIPA{{v}{a}{n}{t\smash{\textsuperscript{h}}}{r}{u}{a}{ð}{\textscy}{\textsubring{r}}} \dicPos{adj}[3] \dicFlx{(f ‑uð)}[3] \textbf{1.} \dicSynonym{tortrygginn} \dicDirectTranslationCS{nevěřící, skeptický};  \dicPhraseIS{vantrúaður á e‑ð} \dicDirectTranslationCS{nevěřící (čemu), skeptický k~(čemu)} \dicExampleIS{Ég er vantrúaður á að hann standi við orð sín.} \dicExampleCS{Nevěřím tomu, že dodrží své slovo.}  \textbf{2.} \dicFieldCat{náb.} \dicSynonym{trúlaus} \dicDirectTranslationCS{ateistický, nevěřící (v~Boha)}
\dicEntry[vanur] \dicTerm{van|ur\smash{\textsuperscript{1}}} \dicIPA{{v}{a}{\textlengthmark}{n}{\textscy}{\textsubring{r}}} \dicPos{m}[9] \dicFlx{(‑s, ‑ir)}[5] \dicFieldCat{myt.} \dicDirectTranslationCS{Van} \dicIndirectTranslationCS{(příslušník starší skupiny bohů v~severské mytologii)} \dicAntonym{ás\smash{\textsuperscript{2}}}
\dicEntry[vanur] \dicTerm{vanur\smash{\textsuperscript{2}}} \dicsymFrequent\  \dicIPA{{v}{a}{\textlengthmark}{n}{\textscy}{\textsubring{r}}} \dicPos{adj}[1] \dicFlx{(f vön)}[2] \dicDirectTranslationCS{zvyklý, navyklý} \dicExampleIS{vanur að ganga} \dicExampleCS{zvyklý chodit};  \dicPhraseIS{vera vanur að (gera e‑ð)} \dicDirectTranslationCS{být zvyklý (dělat (co))};  \dicPhraseIS{vera vanur e‑u} \dicDirectTranslationCS{být zvyklý na (co)};  \dicIdiom{vanur}{ \dicPhraseIS{eins og vant er}} \dicFlx{adv} \dicDirectTranslationCS{jako obvykle\,/\addthin obyčejně}
\dicEntry[Vanúatú] \dicTerm{Vanúatú} \dicIPA{{v}{a}{\textlengthmark}{n}{u}{a}{\textsubring{d}}{u}} \dicPos{n}[4] \dicFlx{indecl}[2] \dicFieldCat{geog.} \dicDirectTranslationCS{Vanuatu}
\dicEntry[Vanúatúi] \dicTerm{Vanúatú|i} \dicIPA{{v}{a}{\textlengthmark}{n}{u}{a}{\textsubring{d}}{u}{\textsci}} \dicPos{m}[1] \dicFlx{(‑a, ‑ar)}[1] \dicDirectTranslationCS{Vanuaťan(ka)}
\dicEntry[vanúatúskur] \dicTerm{vanúatúskur} \dicIPA{{v}{a}{\textlengthmark}{n}{u}{a}{\textsubring{d}}{u}{s}{\r{g}}{\textscy}{\textsubring{r}}} \dicPos{adj}[1]\dicFlx{}[-6] \dicDirectTranslationCS{vanuatský}
\dicEntry[vanvirða] \dicTerm{van··virð|a\smash{\textsuperscript{1}}} \dicIPA{{v}{a}{n}{v}{\textsci}{r}{ð}{a}} \dicPos{f}[1] \dicFlx{(‑u)}[5] \dicDirectTranslationCS{neúcta, pohrdání, despekt}
\dicEntry[vanvirða] \dicTerm{van··vir|ða\smash{\textsuperscript{2}}} \dicIPA{{v}{a}{n}{v}{\textsci}{r}{ð}{a}} \dicPos{v}[2] \dicFlx{(‑ti, ‑t)}[47] \dicFlx{acc} \dicSynonym{óvirða} \dicDirectTranslationCS{zneuctít, (z)hanobit, pohanět} \dicExampleIS{vanvirða minningu hennar} \dicExampleCS{zneuctít její památku}
\dicEntry[vanvirðing] \dicTerm{van··virð·ing} \dicIPA{{v}{a}{n}{v}{\textsci}{r}{ð}{i}{\ng}{\r{g}}} \dicPos{f}[4] \dicFlx{(‑ar)}[7] \dicDirectTranslationCS{zneuctění, (z)hanobení, pohanění}
\dicEntry[vanþakka] \dicTerm{van··þakk|a} \dicIPA{{v}{a}{n}{\texttheta}{a}{h}{\r{g}}{a}} \dicPos{v}[1] \dicFlx{(‑aði)}[13] \dicFlx{acc} \dicDirectTranslationCS{být nevděčný, nepoděkovat} \dicExampleIS{vanþakka hjálpina} \dicExampleCS{být nevděčný za pomoc}
\dicEntry[vanþakklátur] \dicTerm{van··þakk·látur} \dicIPA{{v}{a}{n}{\texttheta}{a}{h}{\r{g}}{l}{au}{\textsubring{d}}{\textscy}{\textsubring{r}}} \dicPos{adj}[1]\dicFlx{}[-1] \textbf{1.} \dicDirectTranslationCS{nevděčný} \dicIndirectTranslationCS{(k\nobreak terý neprojevuje vděčnost)} \dicExampleIS{Mér þykir hann vanþakklátur.} \dicExampleCS{Připadá mi, že je nevděčný.}  \textbf{2.} \dicDirectTranslationCS{nevděčný, nepříjemný (práce ap.)} \dicExampleIS{vanþakklátt starf} \dicExampleCS{nevděčná práce}
\dicEntry[vanþakklæti] \dicTerm{van··þakk·læti} \dicIPA{{v}{a}{n}{\texttheta}{a}{h}{\r{g}}{l}{a}{i}{\textsubring{d}}{\textsci}} \dicPos{n}[2] \dicFlx{(‑s)}[20] \dicDirectTranslationCS{nevděčnost, nevděk}
\dicEntry[vanþekking] \dicTerm{van··þekk·ing} \dicIPA{{v}{a}{n}{\texttheta}{\textepsilon}{h}{\r{\textObardotlessj}}{i}{\ng}{\r{g}}} \dicPos{f}[4] \dicFlx{(‑ar)}[7] \dicSynonym{þekkingarleysi} \dicDirectTranslationCS{neznalost, nevědomost}
\dicEntry[vanþóknun] \dicTerm{van··þókn|un} \dicIPA{{v}{a}{n}{\texttheta}{ou}{h}{\r{g}}{n}{\textscy}{\textsubring{n}}} \dicPos{f}[7] \dicFlx{(‑unar)}[9] \dicDirectTranslationCS{nelibost, nechuť, nevole} \dicExampleIS{vanþóknun á framkomu hennar} \dicExampleCS{nelibost z~jejího chování}
\begin{xtolerant}{}{1pt}
\dicEntry[vanþroska] \dicTerm{van··þroska}\addthinS\dicIPA{{v}{a}{n}{\texttheta}{r}{\textopeno}{s}{\r{g}}{a}}\addthinS\dicPos{adj}[1\hspace{-0.06667em}3]\addthinS\dicFlx{indecl}[1]\addthinS\dicSynonym{óþroskaður}\addthinS\dicDirectTranslationCS{málo\hspace{-0.10667em}\,/\addthin nedostatečně vyvinutý (dítě ap.)}
\end{xtolerant}
\dicEntry[vanþroski] \dicTerm{van··þrosk|i} \dicIPA{{v}{a}{n}{\texttheta}{r}{\textopeno}{s}{\r{\textObardotlessj}}{\textsci}} \dicPos{m}[1] \dicFlx{(‑a)}[3] \dicDirectTranslationCS{nedostatečná vyvinutost, nevyvinutost, nevyspělost}
\dicEntry[vanþróaður] \dicTerm{van··þró·|aður} \dicIPA{{v}{a}{n}{\texttheta}{r}{ou}{a}{ð}{\textscy}{\textsubring{r}}} \dicPos{adj}[3] \dicFlx{(f ‑uð)}[3] \dicDirectTranslationCS{nevyvinutý, nerozvinutý, nevyspělý}
\dicEntry[vanþörf] \dicTerm{van··|þörf} \dicIPA{{v}{a}{n}{\texttheta}{\oe}{r}{f}} \dicPos{f}[7] \dicFlx{(‑þarfar)}[19] \dicSynonym*{þarfleysa} \dicDirectTranslationCS{zbytečnost, nadbytečnost};  \dicPhraseIS{það er ekki vanþörf á e‑u} \dicDirectTranslationCS{je silná potřeba (čeho) (nových bytů ap.)}
\dicEntry[vappa] \dicTerm{vapp|a} \dicIPA{{v}{a}{h}{\textsubring{b}}{a}} \dicPos{v}[1] \dicFlx{(‑aði)}[13] \dicSynonym{trítla} \dicDirectTranslationCS{cupitat, batolit se}
\dicEntry[var] \dicTerm{var\smash{\textsuperscript{1}}} \dicsymFrequent\  \dicIPA{{v}{a}{\textlengthmark}{\textsubring{r}}} \dicPos{adj}[5] \dicFlx{(f vör)}[9] \textbf{1.} \dicSynonym{athugull} \dicDirectTranslationCS{vědomý, uvědomující si};  \dicPhraseIS{verða e‑s var, verða var við e‑ð} \dicDirectTranslationCS{být si (čeho) vědom, uvědomovat si (co)};  \dicPhraseIS{gera vart við sig} {\textbf{a.}} \dicDirectTranslationCS{upozornit na sebe, dát o~sobě vědět (zavoláním ap.)};  {\textbf{b.}} \dicSynonym*{birtast} \dicDirectTranslationCS{objevit se, ukázat se}  \textbf{2.} \dicSynonym{gætinn} \dicDirectTranslationCS{pozorný, ostražitý, bedlivý};  \dicPhraseIS{vera var um sig} \dicDirectTranslationCS{mít se na pozoru, být opatrný}
\dicEntry[var] \dicTerm{var\smash{\textsuperscript{2}}} \dicIPA{{v}{a}{\textlengthmark}{\textsubring{r}}} \dicPos{v} \dicFlx{ind pf sg 1 pers} \dicLink{vera\smash{\textsuperscript{2}}}
\dicEntry[vara] \dicTerm{vara\smash{\textsuperscript{1}}} \dicsymFrequent\  \dicIPA{{v}{a}{\textlengthmark}{r}{a}} \dicPos{f}[1] \dicFlx{(vöru, vörur)}[8] \dicDirectTranslationCS{zboží, výrobek, tovar} \dicExampleIS{aðflutt vara} \dicExampleCS{importované zboží}
\dicEntry[vara] \dicTerm{var|a\smash{\textsuperscript{2}}} \dicsymFrequent\  \dicIPA{{v}{a}{\textlengthmark}{r}{a}} \dicPos{v}[1] \dicFlx{(‑aði)}[13] \dicFlx{acc} \dicDirectTranslationCS{varovat, upozornit, upozorňovat};  \dicPhraseIS{vara e‑n við e‑u} \dicDirectTranslationCS{varovat (koho) před (čím)} \dicExampleIS{vara hana við hættunni} \dicExampleCS{varovat ji před nebezpečím};  \dicIdiom{vara}[sig]{ \dicPhraseIS{vara sig}} \dicDirectTranslationCS{dávat (si) pozor} \dicExampleIS{Varaðu þig!} \dicExampleCS{Dávej pozor!};  \dicIdiom{varast}{ \dicPhraseIS{varast}} \dicFlx{refl acc} \dicDirectTranslationCS{vyvarovat se, vystříhat se} \dicExampleIS{varast þessar aðstæður} \dicExampleCS{vyvarovat se těchto podmínek}
\dicEntry[vara] \dicTerm{var|a\smash{\textsuperscript{3}}} \dicsymFrequent\  \dicIPA{{v}{a}{\textlengthmark}{r}{a}} \dicPos{v}[1] \dicFlx{(‑aði)}[16] \dicSynonym{endast} \dicDirectTranslationCS{(pře)trvat, vydržet} \dicExampleIS{Hvað varir þetta lengi?} \dicExampleCS{Jak dlouho to vydrží?}
\dicEntry[vara] \dicTerm{var|a\smash{\textsuperscript{4}}} \dicsymFrequent\  \dicIPA{{v}{a}{\textlengthmark}{r}{a}} \dicPos{v}[2] \dicFlx{(‑ði, ‑að)}[127] \dicFlx{impers} \dicPhraseIS{fyrr en e‑n varir} \dicDirectTranslationCS{dříve, než by to (kdo) očekával} \dicExampleIS{Þetta getur breyst fyrr en mig varir.} \dicExampleCS{Může se to změnit dříve, než bych očekával.};  \dicPhraseIS{þegar minnst varir} \dicFlx{adv} \dicDirectTranslationCS{náhle, znenadání}
\dicEntry[vara] \dicTerm{vara-\smash{\textsuperscript{5}}} \dicIPA{{v}{a}{\textlengthmark}{r}{a}} \dicPos{predp} \textbf{1.} \dicDirectTranslationCS{náhradní, zásobní, záložní, rezervní} \dicExampleIS{varabúnaður} \dicExampleCS{náhradní vybavení}  \textbf{2.} \dicSynonym{staðgengill} \dicDirectTranslationCS{vice-, místo-} \dicExampleIS{vararíkisstjóri} \dicExampleCS{viceguvernér}
\begin{xtolerant}{}{1pt}
\dicEntry[varaáætlun] \dicTerm{vara··á·ætl|un}\addthinS\dicIPA{{v}{a}{\textlengthmark}{r}{a}{au}{a}{i}{h}{\textsubring{d}}{l}{\textscy}{\textsubring{n}}}\addthinS\dicPos{f}[7]\addthinS\dicFlx{(‑unar,\addthinSS‑anir)}[8] \dicDirectTranslationCS{náhradní\,/\addthin záložní plán}
\end{xtolerant}
\dicEntry[varadekk] \dicTerm{vara··dekk} \dicIPA{{v}{a}{\textlengthmark}{r}{a}{\textsubring{d}}{\textepsilon}{h}{\r{g}}} \dicPos{n}[2] \dicFlx{(‑s, ‑)}[13] \dicDirectTranslationCS{rezerva, rezervní kolo s~pneumatikou}
\dicEntry[varaforði] \dicTerm{vara··forð|i} \dicIPA{{v}{a}{\textlengthmark}{r}{a}{f}{\textopeno}{r}{ð}{\textsci}} \dicPos{m}[1] \dicFlx{(‑a)}[3] \dicDirectTranslationCS{zásoba, rezerva}
\dicEntry[varaformaður] \dicTerm{vara··for·|maður} \dicIPA{{v}{a}{\textlengthmark}{r}{a}{f}{\textopeno}{r}{m}{a}{ð}{\textscy}{\textsubring{r}}} \dicPos{m}[13] \dicFlx{(‑manns, ‑menn)}[2] \dicDirectTranslationCS{místopředseda, místopředsedkyně}
\dicEntry[varaforseti] \dicTerm{vara··for·set|i} \dicIPA{{v}{a}{\textlengthmark}{r}{a}{f}{\textopeno}{\textsubring{r}}{s}{\textepsilon}{\textsubring{d}}{\textsci}} \dicPos{m}[1] \dicFlx{(‑a, ‑ar)}[1] \dicDirectTranslationCS{viceprezident(ka)}
\dicEntry[varahljóð] \dicTerm{vara··hljóð} \dicIPA{{v}{a}{\textlengthmark}{r}{a}{\textsubring{l}}{j}{ou}{\texttheta}} \dicPos{n}[2] \dicFlx{(‑s, ‑)}[5] \dicFieldCat{jaz.} \dicDirectTranslationCS{retná souhláska, retnice, labiála}
\dicEntry[varahlutur] \dicTerm{vara··hlut|ur} \dicIPA{{v}{a}{\textlengthmark}{r}{a}{\textsubring{l}}{\textscy}{\textsubring{d}}{\textscy}{\textsubring{r}}} \dicPos{m}[10] \dicFlx{(‑ar, ‑ir)}[13] \dicDirectTranslationCS{náhradní díl}
\dicEntry[varalið] \dicTerm{vara··lið} \dicIPA{{v}{a}{\textlengthmark}{r}{a}{l}{\textsci}{\texttheta}} \dicPos{n}[2] \dicFlx{(‑s, ‑)}[5] \textbf{1.} \dicFieldCat{voj.} \dicSynonym{herlið} \dicDirectTranslationCS{záložní oddíl}  \textbf{2.} \dicFieldCat{sport.} \dicSynonym*{keppnishópur} \dicDirectTranslationCS{záložní mužstvo, rezerva}
\dicEntry[varalitur] \dicTerm{vara··lit|ur} \dicIPA{{v}{a}{\textlengthmark}{r}{a}{l}{\textsci}{\textsubring{d}}{\textscy}{\textsubring{r}}} \dicPos{m}[10] \dicFlx{(‑ar, ‑ir)}[13] \dicDirectTranslationCS{rtěnka}
\dicEntry[varamaður] \dicTerm{vara··|maður} \dicIPA{{v}{a}{\textlengthmark}{r}{a}{m}{a}{ð}{\textscy}{\textsubring{r}}} \dicPos{m}[13] \dicFlx{(‑manns, ‑menn)}[2] \dicDirectTranslationCS{náhradník, náhradnice} \dicIndirectTranslationCS{(v~představenstvu, sportu, divadle ap.)}
\dicEntry[varamæltur] \dicTerm{vara··mæltur} \dicIPA{{v}{a}{\textlengthmark}{r}{a}{m}{a}{i}{\textsubring{l}}{\textsubring{d}}{\textscy}{\textsubring{r}}} \dicPos{adj}[1]\dicFlx{}[-13] \dicFieldCat{jaz.} \dicDirectTranslationCS{retný, labiální}
\dicEntry[varanlega] \dicTerm{varan··lega} \dicIPA{{v}{a}{\textlengthmark}{r}{a}{n}{l}{\textepsilon}{\textbabygamma}{a}} \dicPos{adv} \dicDirectTranslationCS{trvale, trvanlivě, permanentně}
\dicEntry[varanlegur] \dicTerm{varan··legur} \dicIPA{{v}{a}{\textlengthmark}{r}{a}{n}{l}{\textepsilon}{\textbabygamma}{\textscy}{\textsubring{r}}} \dicPos{adj}[1]\dicFlx{}[-8] \dicSynonym*{sem endist} \dicDirectTranslationCS{trvalý, trvanlivý, permanentní} \dicExampleIS{varanlegur friður} \dicExampleCS{trvalý mír}
\dicEntry[varanleiki] \dicTerm{varan··leik|i} \dicIPA{{v}{a}{\textlengthmark}{r}{a}{n}{l}{ei}{\r{\textObardotlessj}}{\textsci}} \dicPos{m}[1] \dicFlx{(‑a)}[3] \dicSynonym{ending} \dicDirectTranslationCS{trvalost, trvanlivost}
\dicEntry[varar] \dicTerm{varar} \dicIPA{{v}{a}{\textlengthmark}{r}{a}{\textsubring{r}}} \dicPos{f} \dicFlx{sg gen} \dicLink{vör\smash{\textsuperscript{1}}}
\dicEntry[varasamur] \dicTerm{vara··|samur} \dicIPA{{v}{a}{\textlengthmark}{r}{a}{s}{a}{m}{\textscy}{\textsubring{r}}} \dicPos{adj}[1] \dicFlx{(f ‑söm)}[2] \dicSynonym{hættulegur} \dicDirectTranslationCS{ošidný, ošemetný, zrádný} \dicExampleIS{varasamt tiltæki} \dicExampleCS{ošemetná výstřednost}
\dicEntry[varasjóður] \dicTerm{vara··sjóð|ur} \dicIPA{{v}{a}{\textlengthmark}{r}{a}{s}{j}{ou}{ð}{\textscy}{\textsubring{r}}} \dicPos{m}[9] \dicFlx{(‑s, ‑ir)}[6] \dicFieldCat{ekon.} \dicDirectTranslationCS{rezervní fond}
\dicEntry[varaskeifa] \dicTerm{vara··skeif|a} \dicIPA{{v}{a}{\textlengthmark}{r}{a}{s}{\r{\textObardotlessj}}{ei}{v}{a}} \dicPos{f}[1] \dicFlx{(‑u, ‑ur)}[19] \textbf{1.} \dicSynonym*{aukaskeifa} \dicDirectTranslationCS{náhradní podkova}  \textbf{2.} \dicDirectTranslationCS{náhrada, náhradník, náhradnice} \dicIndirectTranslationCS{(v~hanlivém významu)};  \dicPhraseIS{hafa e‑n fyrir varaskeifu} \dicDirectTranslationCS{používat (koho) jako nouzové řešení}
\dicEntry[varaþurrkur] \dicTerm{vara··þurrk|ur} \dicIPA{{v}{a}{\textlengthmark}{r}{a}{\texttheta}{\textscy}{\textsubring{r}}{\r{g}}{\textscy}{\textsubring{r}}} \dicPos{m}[6] \dicFlx{(‑s, ‑ar)}[4] \dicDirectTranslationCS{suché\,/\addthin po\-pra\-s\-ka\-né rty}
\dicEntry[varð] \dicTerm{varð} \dicIPA{{v}{a}{r}{\texttheta}} \dicPos{v} \dicFlx{ind pf sg 1 pers} \dicLink{verða}
\dicEntry[varða] \dicTerm{varð|a\smash{\textsuperscript{1}}} \dicIPA{{v}{a}{r}{ð}{a}} \dicPos{v}[1] \dicFlx{(‑aði)}[56] \dicFlx{acc} \dicSynonym*{hlaða vörður} \dicDirectTranslationCS{zasadit\,/\addthin sázet patníky} \dicExampleIS{varða veg} \dicExampleCS{sázet patníky u~cesty}
\dicEntry[varða] \dicTerm{varð|a\smash{\textsuperscript{2}}} \dicsymFrequent\  \dicIPA{{v}{a}{r}{ð}{a}} \dicPos{v}[1] \dicFlx{(‑aði)}[55] \dicFlx{dat} \textbf{1.} \dicSynonym{snerta} \dicDirectTranslationCS{týkat se, souviset (s)};  \dicPhraseIS{hvað e‑ð varðar} \dicFlx{impers} \dicDirectTranslationCS{co se týče (čeho), pokud jde o~(co)};  \dicPhraseIS{hvað varðar e‑n um það?} \dicFlx{impers} \dicDirectTranslationCS{co je (komu) po tom?};  \dicPhraseIS{e‑n varðar ekkert um e‑ð} \dicFlx{impers} \dicDirectTranslationCS{(co) se (koho) netýká}  \textbf{2.} \dicFlx{dat} \dicSynonym{gæta} \dicDirectTranslationCS{hájit, bránit, střežit};  \dicPhraseIS{varða e‑u} \dicDirectTranslationCS{chránit (co), hájit (co) (zájmy ap.)}  \textbf{3.} \dicFlx{dat} \dicDirectTranslationCS{mít za následek, (vy)ústit};  \dicPhraseIS{e‑að varðar e‑u} \dicDirectTranslationCS{(co) má za následek (co), (co) ústí v~(co) (trest ap.)};  \dicPhraseIS{e‑að varðar við lög} \dicDirectTranslationCS{(co) je v~rozporu se zákonem}
\dicEntry[varðandi] \dicTerm{varð··andi} \dicIPA{{v}{a}{r}{ð}{a}{n}{\textsubring{d}}{\textsci}} \dicPos{adj}[13] \dicFlx{indecl}[1] \dicPhraseIS{varðandi e‑ð} \dicDirectTranslationCS{ohledně (čeho), týkající se (čeho), související s~(čím)} \dicExampleIS{varðandi þetta atriði} \dicExampleCS{ohledně tohoto bodu}
\dicEntry[varðar] \dicTerm{varðar} \dicIPA{{v}{a}{r}{ð}{a}{\textsubring{r}}} \dicPos{m} \dicFlx{sg gen} \dicLink{vörður}
\dicEntry[varðeldur] \dicTerm{varð··eld|ur} \dicIPA{{v}{a}{r}{ð}{\textepsilon}{l}{\textsubring{d}}{\textscy}{\textsubring{r}}} \dicPos{m}[6] \dicFlx{(‑s, ‑ar)}[4] \dicDirectTranslationCS{(táborový) oheň, táborák, vatra}
\dicEntry[varðengill] \dicTerm{varð··eng|ill} \dicIPA{{v}{a}{r}{ð}{ei}{\textltailn}{\r{\textObardotlessj}}{\textsci}{\textsubring{d}}{\textsubring{l}}} \dicPos{m}[6] \dicFlx{(‑ils, ‑lar)}[35] \dicFieldCat{náb.} \dicDirectTranslationCS{anděl strážný}
\dicEntry[varðhald] \dicTerm{varð··|hald} \dicIPA{{v}{a}{r}{\texttheta}{h}{a}{l}{\textsubring{d}}} \dicPos{n}[2] \dicFlx{(‑halds, ‑höld)}[8] \dicFieldCat{práv.} \dicDirectTranslationCS{vazba} \dicExampleIS{sitja í varðhaldi} \dicExampleCS{sedět ve vazbě};  \dicPhraseIS{setja e‑n í varðhald} \dicDirectTranslationCS{umístit (koho) do vazby}
\dicEntry[varðhundur] \dicTerm{varð··hund|ur} \dicIPA{{v}{a}{r}{\texttheta}{h}{\textscy}{n}{\textsubring{d}}{\textscy}{\textsubring{r}}} \dicPos{m}[6] \dicFlx{(‑s, ‑ar)}[4] \dicDirectTranslationCS{hlídací pes, hlídač}
\dicEntry[varði] \dicTerm{varð|i\smash{\textsuperscript{1}}} \dicIPA{{v}{a}{r}{ð}{\textsci}} \dicPos{m}[1] \dicFlx{(‑a, ‑ar)}[8] \dicSynonym{minnismerki} \dicDirectTranslationCS{pomník, památník}
\dicEntry[varði] \dicTerm{varði\smash{\textsuperscript{2}}} \dicIPA{{v}{a}{r}{ð}{\textsci}} \dicPos{v} \dicFlx{ind pf sg 1 pers} \textbf{1.} \dicLink{verja\smash{\textsuperscript{2}}}  \textbf{2.} \dicLink{verja\smash{\textsuperscript{3}}}
\dicEntry[varðmaður] \dicTerm{varð··|maður} \dicIPA{{v}{a}{r}{ð}{m}{a}{ð}{\textscy}{\textsubring{r}}} \dicPos{m}[13] \dicFlx{(‑manns, ‑menn)}[2] \dicSynonym{gæslumaður} \dicDirectTranslationCS{strážce, strážkyně, hlídač(ka), pochůzkář(ka)}
\dicEntry[varðskip] \dicTerm{varð··skip} \dicIPA{{v}{a}{r}{ð}{s}{\r{\textObardotlessj}}{\textsci}{\textsubring{b}}} \dicPos{n}[2] \dicFlx{(‑s, ‑)}[5] \dicFieldCat{nám.} \dicDirectTranslationCS{hlídková loď}
\dicEntry[varðstjóri] \dicTerm{varð··stjór|i} \dicIPA{{v}{a}{r}{ð}{s}{\textsubring{d}}{j}{ou}{r}{\textsci}} \dicPos{m}[1] \dicFlx{(‑a, ‑ar)}[1] \dicDirectTranslationCS{důstojník\,/\addthin důstojnice ve službě}
\dicEntry[varðstofa] \dicTerm{varð··stof|a} \dicIPA{{v}{a}{r}{ð}{s}{\textsubring{d}}{\textopeno}{v}{a}} \dicPos{f}[1] \dicFlx{(‑u, ‑ur)}[7] \dicDirectTranslationCS{strážnice, policejní stanoviště}
\dicEntry[varðtunna] \dicTerm{varð··tunn|a} \dicIPA{{v}{a}{r}{ð}{t\smash{\textsuperscript{h}}}{\textscy}{n}{a}} \dicPos{f}[1] \dicFlx{(‑u, ‑ur)}[7] \dicFieldCat{nám.} \dicDirectTranslationCS{strážní koš}
\dicEntry[varðturn] \dicTerm{varð··turn} \dicIPA{{v}{a}{r}{ð}{t\smash{\textsuperscript{h}}}{\textscy}{r}{\textsubring{d}}{\textsubring{n}}} \dicPos{m}[4] \dicFlx{(‑s, ‑ar)}[3] \dicDirectTranslationCS{strážní věž, hláska}
\dicEntry[varðveisla] \dicTerm{varð··veisl|a} \dicIPA{{v}{a}{r}{ð}{v}{ei}{s}{\textsubring{d}}{l}{a}} \dicPos{f}[1] \dicFlx{(‑u)}[5] \textbf{1.} \dicSynonym{geymsla} \dicDirectTranslationCS{úschova, uchování, opatrování} \dicExampleIS{fá honum sjóðinn til varðveislu} \dicExampleCS{svěřit mu hotovost do úschovy}  \textbf{2.} \dicDirectTranslationCS{konzervace} \dicIndirectTranslationCS{(způsob ochrany)}
\dicEntry[varðveita] \dicTerm{varð··veit|a} \dicsymFrequent\  \dicIPA{{v}{a}{r}{ð}{v}{ei}{\textsubring{d}}{a}} \dicPos{v}[2] \dicFlx{(‑ti, ‑t)}[54] \dicFlx{acc} \dicDirectTranslationCS{zachovat, zachovávat, uchovat, uchovávat, udržet, udržovat, (za)chránit (přírodu ap.)} \dicExampleIS{varðveita minningarnar} \dicExampleCS{uchovat vzpomínky};  \dicPhraseIS{varðveita e‑ð eins og sjáaldur auga síns} \dicLangCat{přen.} \dicDirectTranslationCS{střežit (co) jako oko v~hlavě};  \dicIdiom{varðveitast}{ \dicPhraseIS{varðveitast}} \dicFlx{refl} \dicDirectTranslationCS{zachovat se, být zachován (v~dobrém stavu ap.)}
\dicEntry[varfærinn] \dicTerm{var··færinn} \dicIPA{{v}{a}{\textsubring{r}}{f}{a}{i}{r}{\textsci}{\textsubring{n}}} \dicPos{adj}[6]\dicFlx{}[-2] \dicSynonym*{varúðarfullur} \dicDirectTranslationCS{rozvážný, obezřetný} \dicExampleIS{vera varfærinn í fjármálum} \dicExampleCS{být obezřetný ve finančních záležitostech}
\dicEntry[varfærni] \dicTerm{var··færn|i} \dicIPA{{v}{a}{\textsubring{r}}{f}{a}{i}{r}{\textsubring{d}}{n}{\textsci}} \dicPos{f}[3] \dicFlx{(‑i)}[3] \dicDirectTranslationCS{rozvážnost, obezřetnost}
\dicEntry[vargaskapur] \dicTerm{varga··skap|ur} \dicIPA{{v}{a}{r}{\r{g}}{a}{s}{\r{g}}{a}{\textsubring{b}}{\textscy}{\textsubring{r}}} \dicPos{m}[10] \dicFlx{(‑ar)}[15] \dicDirectTranslationCS{vzteklost, zvlčilost}
\dicEntry[vargur] \dicTerm{varg|ur} \dicIPA{{v}{a}{r}{\r{g}}{\textscy}{\textsubring{r}}} \dicPos{m}[6] \dicFlx{(‑s, ‑ar)}[11] \textbf{1.} \dicSynonym*{skaðlegt dýr} \dicDirectTranslationCS{škodná (zvěř)}  \textbf{2.} \dicSynonym*{ofsafenginn maður} \dicDirectTranslationCS{výbušný\,/\addthin prudký člověk}
\dicEntry[vargúlfur] \dicTerm{varg··úlf|ur} \dicIPA{{v}{a}{r}{\r{g}}{u}{l}{v}{\textscy}{\textsubring{r}}} \dicPos{m}[6] \dicFlx{(‑s, ‑ar)}[4] \dicLink{varúlfur}
\dicEntry[varhugaverður] \dicTerm{var·huga··verður} \dicIPA{{v}{a}{\textlengthmark}{\textsubring{r}}{h}{\textscy}{\textbabygamma}{a}{v}{\textepsilon}{r}{ð}{\textscy}{\textsubring{r}}} \dicPos{adj}[2]\dicFlx{}[-1] \dicDirectTranslationCS{ošemetný, riskantní, vyžadující opatrnost} \dicExampleIS{Innsiglingin getur verið varhugaverð í vondu veðri.} \dicExampleCS{Vplutí do přístavu může být při špatném počasí riskantní.}
\dicEntry[varhugi] \dicTerm{var··hug|i} \dicIPA{{v}{a}{\textlengthmark}{\textsubring{r}}{h}{\textscy}{i}{j}{\textsci}} \dicPos{m}[1] \dicFlx{(‑a)}[3] \dicSynonym{varúð} \dicDirectTranslationCS{ostražitost, opatrnost};  \dicPhraseIS{gjalda varhuga við e‑u} \dicDirectTranslationCS{mít se na pozoru před (čím)}
\dicEntry[vari] \dicTerm{var|i} \dicIPA{{v}{a}{\textlengthmark}{r}{\textsci}} \dicPos{m}[1] \dicFlx{(‑a, ‑ar)}[8] \textbf{1.} \dicSynonym{varkárni} \dicDirectTranslationCS{opatrnost, obezřetnost};  \dicPhraseIS{hafa vara(nn) á (sér)} \dicDirectTranslationCS{být opatrný, dát\,/\addthin dávat (si) pozor}  \textbf{2.} \dicLangCat{zast.} \dicSynonym{viðvörun} \dicDirectTranslationCS{varování, upozornění};  \dicPhraseIS{taka e‑m vara fyrir e‑u} \dicDirectTranslationCS{varovat (koho) před (čím)};  \dicIdiom{vari}{ \dicPhraseIS{til vara}} \dicFlx{adv} \dicDirectTranslationCS{v~rezervě\,/\addthin záloze}; { \dicPhraseIS{til vonar og vara}} \dicFlx{adv} \dicSynonym*{til öryggis} \dicDirectTranslationCS{pro jistotu};  \dicProverb\  \dicPhraseIS{Allur er varinn góður.} \dicLangCat{přís.} \dicDirectTranslationCS{Opatrnost matka rozumu.}
\dicEntry[varið] \dicTerm{varið} \dicIPA{{v}{a}{\textlengthmark}{r}{\textsci}{\texttheta}} \dicPos{v} \dicFlx{supin} \textbf{1.} \dicLink{verja\smash{\textsuperscript{2}}}  \textbf{2.} \dicLink{verja\smash{\textsuperscript{3}}}
\dicEntry[varinn] \dicTerm{varinn} \dicIPA{{v}{a}{\textlengthmark}{r}{\textsci}{\textsubring{n}}} \dicPos{adj}[6]\dicFlx{}[-14] \dicDirectTranslationCS{chráněný, bráněný}
\dicEntry[varir] \dicTerm{varir} \dicIPA{{v}{a}{\textlengthmark}{r}{\textsci}{\textsubring{r}}} \dicPos{f} \dicFlx{pl nom} \dicLink{vör\smash{\textsuperscript{1}}}
\dicEntry[varkár] \dicTerm{var··kár} \dicIPA{{v}{a}{\textsubring{r}}{k\smash{\textsuperscript{h}}}{au}{\textsubring{r}}} \dicPos{adj}[5] \dicFlx{(f ‑)}[8] \dicSynonym{gætinn} \dicDirectTranslationCS{opatrný, ostražitý, obezřetný} \dicExampleIS{vera varkár í orðum} \dicExampleCS{být obezřetný ve volbě slov}
\dicEntry[varkárni] \dicTerm{var··kárn|i} \dicIPA{{v}{a}{\textsubring{r}}{k\smash{\textsuperscript{h}}}{au}{r}{\textsubring{d}}{n}{\textsci}} \dicPos{f}[3] \dicFlx{(‑i)}[3] \dicDirectTranslationCS{opatrnost, ostražitost, obezřetnost} \dicExampleIS{fjalla um málið af varkárni} \dicExampleCS{mluvit o~té věci obezřetně}
\dicEntry[varla] \dicTerm{varla} \dicsymFrequent\  \dicIPA{{v}{a}{r}{\textsubring{d}}{l}{a}} \dicPos{adv} \dicFlx{(comp síður, sup síst)} \dicSynonym{tæplega} \dicDirectTranslationCS{sotva, stěží} \dicExampleIS{Hann tók varla eftir honum.} \dicExampleCS{Sotva si ho všiml.}
\dicEntry[varlega] \dicTerm{var··leg|a} \dicsymFrequent\  \dicIPA{{v}{a}{r}{l}{\textepsilon}{\textbabygamma}{a}} \dicPos{adv} \dicFlx{(comp ‑ar, sup ‑ast)} \dicSynonym{gætilega} \dicDirectTranslationCS{opatrně, obezřetně} \dicExampleIS{Hann blés varlega ryk af myndinni.} \dicExampleCS{Opatrně sfoukl prach z~obrazu.}
\dicEntry[varmaaflfræði] \dicTerm{varma··afl·fræð|i} \dicIPA{{v}{a}{r}{m}{a}{a}{\textsubring{b}}{\textsubring{l}}{f}{r}{a}{i}{ð}{\textsci}} \dicPos{f}[3] \dicFlx{(‑i)}[3] \dicLink{varmafræði}
\dicEntry[varmaeining] \dicTerm{varma··ein·ing} \dicIPA{{v}{a}{r}{m}{a}{ei}{n}{i}{\ng}{\r{g}}} \dicPos{f}[4] \dicFlx{(‑ar, ‑ar)}[5] \dicFieldCat{fyz.} \dicSynonym{kaloría} \dicDirectTranslationCS{kalorie}
\dicEntry[varmafræði] \dicTerm{varma··fræð|i}\dicTerm{, varmaaflfræði} \dicIPA{{v}\-{a}\-{r}\-{m}\-{a}\-{f}\-{r}\-{a}\-{i}\-{ð}\-{\textsci}\-} \dicPos{f}[3] \dicFlx{(‑i)}[3] \dicFieldCat{fyz.} \dicDirectTranslationCS{termodynamika}
\dicEntry[varmaorka] \dicTerm{varma··ork|a} \dicIPA{{v}{a}{r}{m}{a}{\textopeno}{\textsubring{r}}{\r{g}}{a}} \dicPos{f}[1] \dicFlx{(‑u)}[5] \dicFieldCat{fyz.} \dicDirectTranslationCS{teplo, tepelná energie}
\dicEntry[varmaskiptir] \dicTerm{varma··skipt|ir} \dicIPA{{v}{a}{r}{m}{a}{s}{\r{\textObardotlessj}}{\textsci}{f}{\textsubring{d}}{\textsci}{\textsubring{r}}} \dicPos{m}[7] \dicFlx{(‑is)}[2] \dicDirectTranslationCS{výměník tepla}
\dicEntry[varmenni] \dicTerm{var··menni} \dicIPA{{v}{a}{r}{m}{\textepsilon}{n}{\textsci}} \dicPos{n}[2] \dicFlx{(‑s, ‑)}[14] \dicSynonym{illmenni} \dicDirectTranslationCS{ničema, zlosyn(ka)}
\dicEntry[varmi] \dicTerm{varm|i} \dicIPA{{v}{a}{r}{m}{\textsci}} \dicPos{m}[1] \dicFlx{(‑a)}[3] \textbf{1.} \dicSynonym{hiti} \dicDirectTranslationCS{teplo}  \textbf{2.} \dicFieldCat{fyz.} \dicSynonym{varmaorka} \dicDirectTranslationCS{teplo}
\dicEntry[varmur] \dicTerm{varmur} \dicIPA{{v}{a}{r}{m}{\textscy}{\textsubring{r}}} \dicPos{adj}[1] \dicFlx{(f vörm)}[2] \dicPhraseIS{að vörmu spori} \dicFlx{adv} \dicSynonym{strax} \dicDirectTranslationCS{(i)hned, okamžitě}
\dicEntry[varna] \dicTerm{varn|a} \dicIPA{{v}{a}{r}{\textsubring{d}}{n}{a}} \dicPos{v}[1] \dicFlx{(‑aði)}[13] \dicFlx{dat (+ gen)} \dicSynonym{afstýra} \dicDirectTranslationCS{(za)bránit, zamezit, zamezovat};  \dicPhraseIS{varna e‑m e‑s} \dicDirectTranslationCS{zabránit (komu) v~(čem), zamezit (komu co)} \dicExampleIS{varna honum inngöngu} \dicExampleCS{zabránit mu ve vstupu};  \dicPhraseIS{varna e‑m máls} \dicDirectTranslationCS{bránit (komu) v~řeči}
\dicEntry[varnaðarorð] \dicTerm{varnaðar··orð} \dicIPA{{v}{a}{r}{\textsubring{d}}{n}{a}{ð}{a}{r}{\textopeno}{r}{\texttheta}} \dicPos{n}[2] \dicFlx{(‑s, ‑)}[5] \dicDirectTranslationCS{varování, výstraha}
\dicEntry[varnaður] \dicTerm{varn··að|ur} \dicIPA{{v}{a}{r}{\textsubring{d}}{n}{a}{ð}{\textscy}{\textsubring{r}}} \dicPos{m}[10] \dicFlx{(‑ar)}[9] \textbf{1.} \dicSynonym{viðvörun} \dicDirectTranslationCS{výstraha, varování};  \dicPhraseIS{e‑að er víti til varnaðar} \dicLangCat{přen.} \dicDirectTranslationCS{(co) je odstrašující příklad}  \textbf{2.} \dicSynonym{vernd} \dicDirectTranslationCS{ostraha, ochrana}
\dicEntry[varnagli] \dicTerm{var··nagl|i} \dicIPA{{v}{a}{r}{n}{a}{\r{g}}{l}{\textsci}} \dicPos{m}[1] \dicFlx{(‑a, ‑ar)}[8] \dicSynonym{fyrirvari} \dicDirectTranslationCS{pojistka, výhrada};  \dicPhraseIS{slá varnagla(nn) við e‑u} \dicLangCat{přen.} \dicDirectTranslationCS{pojistit se před (čím), zabezpečit se proti (čemu)}
\dicEntry[varnar] \dicTerm{varnar} \dicIPA{{v}{a}{r}{\textsubring{d}}{n}{a}{\textsubring{r}}} \dicPos{f} \dicFlx{sg gen} \dicLink{vörn}
\dicEntry[varnarbandalag] \dicTerm{varnar··banda·|lag} \dicIPA{{v}{a}{r}{\textsubring{d}}{n}{a}{r}{\textsubring{b}}{a}{n}{\textsubring{d}}{a}{l}{a}{x}} \dicPos{n}[2] \dicFlx{(‑lags, ‑lög)}[8] \dicFieldCat{voj.} \dicDirectTranslationCS{obranná aliance}
\dicEntry[varnargarður] \dicTerm{varnar··garð|ur} \dicIPA{{v}{a}{r}{\textsubring{d}}{n}{a}{r}{\r{g}}{a}{r}{ð}{\textscy}{\textsubring{r}}} \dicPos{m}[6] \dicFlx{(‑s, ‑ar)}[5] \dicDirectTranslationCS{ochranná hráz, ochranný násep}
\dicEntry[varnarlaus] \dicTerm{varnar··laus} \dicIPA{{v}{a}{r}{\textsubring{d}}{n}{a}{r}{l}{\oe i}{s}} \dicPos{adj}[5]\dicFlx{}[-1] \dicDirectTranslationCS{bezbranný, nechráněný} \dicExampleIS{vera alveg varnarlaus gagnvart þessum hótunum} \dicExampleCS{být úplně bezbranný vůči těmto výhrůžkám}
\begin{xtolerant}{}{9pt}
\dicEntry[varnarleikmaður] \textls[15]{\dicTerm{varnar·leik··|maður} \dicIPA{{v}{a}{r}{\textsubring{d}}{n}{a}{r}{l}{ei}{\r{g}}{m}{a}{ð}{\textscy}{\textsubring{r}}}} \dicPos{m}[13] \dicFlx{(‑manns, ‑menn)}[2] \dicFieldCat{sport.} \dicDirectTranslationCS{obránce, obránkyně}
\end{xtolerant}
\dicEntry[varnarleikur] \dicTerm{varnar··leik|ur} \dicIPA{{v}{a}{r}{\textsubring{d}}{n}{a}{r}{l}{ei}{\r{g}}{\textscy}{\textsubring{r}}} \dicPos{m}[6] \dicFlx{(‑s)}[17] \dicFieldCat{sport.} \dicDirectTranslationCS{obrana, obranná hra}
\dicEntry[varnarleysi] \dicTerm{varnar··leysi} \dicIPA{{v}{a}{r}{\textsubring{d}}{n}{a}{r}{l}{ei}{s}{\textsci}} \dicPos{n}[2] \dicFlx{(‑s)}[20] \dicDirectTranslationCS{bezbrannost}
\dicEntry[varnarlið] \dicTerm{varnar··lið} \dicIPA{{v}{a}{r}{\textsubring{d}}{n}{a}{r}{l}{\textsci}{\texttheta}} \dicPos{n}[2] \dicFlx{(‑s, ‑)}[5] \dicFieldCat{voj.} \dicDirectTranslationCS{obrana, obranné síly}
\dicEntry[varnarmaður] \dicTerm{varnar··|maður} \dicIPA{{v}{a}{r}{\textsubring{d}}{n}{a}{r}{m}{a}{ð}{\textscy}{\textsubring{r}}} \dicPos{m}[13] \dicFlx{(‑manns, ‑menn)}[2] \dicFieldCat{sport.} \dicDirectTranslationCS{obránce, obránkyně}
\dicEntry[varnarmálaráðherra] \dicTerm{varnar·mála··ráð·herr|a} \dicIPA{{v}\-{a}\-{r}\-{\textsubring{d}}\-{n}\-{a}\-{r}\-{m}\-{au}\-{l}\-{a}\-{r}\-{au}\-{\texttheta}\-{h}\-{\textepsilon}\-{r}\-{a}\-} \dicPos{m}[1] \dicFlx{(‑a, ‑ar)}[17] \dicDirectTranslationCS{ministr(yně) obrany}
\dicEntry[varnarmálaráðuneyti] \dicTerm{varnar·mála··ráðu·neyti} \dicIPA{{v}\-{a}\-{r}\-{\textsubring{d}}\-{n}\-{a}\-{r}\-{m}\-{au}\-{l}\-{a}\-{r}\-{au}\-{ð}\-{\textscy}\-{n}\-{ei}\-{\textsubring{d}}\-{\textsci}\-} \dicPos{n}[2] \dicFlx{(‑s, ‑)}[14] \dicDirectTranslationCS{ministerstvo obrany}
\begin{xtolerant}{}{1pt}
\dicEntry[varnarstaða] \dicTerm{varnar··|staða}\addthinS\dicIPA{{v}{a}{r}{\textsubring{d}}{n}{a}{\textsubring{r}}{s}{\textsubring{d}}{a}{ð}{a}}\addthinS\dicPos{f}[1]\addthinS\dicFlx{(‑stöðu)}[2]\addthinS\dicDirectTranslationCS{obrana, obranná pozice};  \dicPhraseIS{fara í varnarstöðu} \dicLangCat{přen.} \dicDirectTranslationCS{přejít do obrany}
\end{xtolerant}
\dicEntry[varnarveggur] \dicTerm{varnar··vegg|ur} \dicIPA{{v}{a}{r}{\textsubring{d}}{n}{a}{r}{v}{\textepsilon}{\r{g}}{\textscy}{\textsubring{r}}} \dicPos{m}[10] \dicFlx{(‑jar, ‑ir)}[23] \dicDirectTranslationCS{ochranná zeď\,/\addthin bariéra}
\dicEntry[varningur] \dicTerm{varn··ing|ur} \dicIPA{{v}{a}{r}{\textsubring{d}}{n}{i}{\ng}{\r{g}}{\textscy}{\textsubring{r}}} \dicPos{m}[6] \dicFlx{(‑s)}[9] \dicSynonym{vara\smash{\textsuperscript{1}}} \dicDirectTranslationCS{zboží, výrobek, produkt} \dicExampleIS{flytja varning til landsins} \dicExampleCS{dovážet zboží do země}
\dicEntry[varnir] \dicTerm{varnir} \dicIPA{{v}{a}{r}{\textsubring{d}}{n}{\textsci}{\textsubring{r}}} \dicPos{f} \dicFlx{pl nom} \dicLink{vörn}
\dicEntry[varp] \dicTerm{varp\smash{\textsuperscript{1}}} \dicIPA{{v}{a}{\textsubring{r}}{\textsubring{b}}} \dicPos{n}[2] \dicFlx{(‑s, vörp)}[8] \textbf{1.} \dicSynonym{kast} \dicDirectTranslationCS{vrh, vrhnutí}  \textbf{2.} \dicSynonym*{fuglavarp} \dicDirectTranslationCS{kladení vajec}  \textbf{3.} \dicSynonym*{varpsvæði} \dicDirectTranslationCS{hnízdiště}
\dicEntry[varp] \dicTerm{varp\smash{\textsuperscript{2}}} \dicIPA{{v}{a}{\textsubring{r}}{\textsubring{b}}} \dicPos{v} \dicFlx{ind pf sg 1 pers} \dicLink{verpa\smash{\textsuperscript{2}}}
\dicEntry[varpa] \dicTerm{varpa\smash{\textsuperscript{1}}} \dicIPA{{v}{a}{\textsubring{r}}{\textsubring{b}}{a}} \dicPos{f}[1] \dicFlx{(vörpu, vörpur)}[20] \dicFieldCat{nám.} \dicDirectTranslationCS{vlečná (rybářská) síť, nevod}
\dicEntry[varpa] \dicTerm{varp|a\smash{\textsuperscript{2}}} \dicsymFrequent\  \dicIPA{{v}{a}{\textsubring{r}}{\textsubring{b}}{a}} \dicPos{v}[1] \dicFlx{(‑aði)}[13] \dicFlx{dat} \dicDirectTranslationCS{vrhnout, vrhat, hodit, házet} \dicExampleIS{varpa steininum} \dicExampleCS{hodit kamenem};  \dicIdiom{varpa}[á]{ \dicPhraseIS{varpa ljósi á e‑ð}} \dicLangCat{přen.} \dicSynonym{útskýra} \dicDirectTranslationCS{objasnit (co), vysvětlit (co)};  \dicIdiom{varpa}[fram]{ \dicPhraseIS{varpa fram e‑u}} \dicDirectTranslationCS{nadhodit (co) (otázku ap.)} \dicExampleIS{varpa fram spurningu} \dicExampleCS{nadhodit otázku};  \dicIdiom{varpa}{ \dicPhraseIS{varpa akkerum}} \dicFieldCat{nám.} \dicDirectTranslationCS{vyhodit kotvu}; { \dicPhraseIS{varpa öndinni léttar}} \dicDirectTranslationCS{vydechnout si, oddechnout si}
\dicEntry[varpanir] \dicTerm{varpanir} \dicIPA{{v}{a}{\textsubring{r}}{\textsubring{b}}{a}{n}{\textsci}{\textsubring{r}}} \dicPos{f} \dicFlx{pl nom} \dicLink{vörpun}
\dicEntry[varpfugl] \dicTerm{varp··fugl} \dicIPA{{v}{a}{\textsubring{r}}{\textsubring{b}}{f}{\textscy}{\r{g}}{\textsubring{l}}} \dicPos{m}[4] \dicFlx{(‑s, ‑ar)}[3] \dicDirectTranslationCS{hnízdící pták, pták kladoucí vejce}
\dicEntry[varptími] \dicTerm{varp··tím|i} \dicIPA{{v}{a}{\textsubring{r}}{\textsubring{b}}{t\smash{\textsuperscript{h}}}{i}{m}{\textsci}} \dicPos{m}[1] \dicFlx{(‑a)}[3] \dicDirectTranslationCS{(doba) hnízdění, doba kladení vajec}
\dicEntry[Varsjá] \dicTerm{Var··sjá} \dicIPA{{v}{a}{\textsubring{r}}{s}{j}{au}} \dicPos{f}[4] \dicFlx{(‑r)}[20] \dicFieldCat{geog.} \dicDirectTranslationCS{Varšava} \dicIndirectTranslationCS{(hlavní město Polska)}
\dicEntry[varsla] \dicTerm{varsla} \dicIPA{{v}{a}{\textsubring{r}}{s}{\textsubring{d}}{l}{a}} \dicPos{f}[1] \dicFlx{(vörslu, vörslur)}[8] \dicSynonym{gæsla} \dicDirectTranslationCS{opatrování, střežení, dohled} \dicExampleIS{Þýfið er í vörslu lögreglunnar.} \dicExampleCS{Ukradené zboží je pod dohledem policie.}
\dicEntry[vart] \dicTerm{vart} \dicsymFrequent\  \dicIPA{{v}{a}{\textsubring{r}}{\textsubring{d}}} \dicPos{adv} \dicSynonym{varla} \dicDirectTranslationCS{stěží, sotva, (jen) taktak} \dicExampleIS{Við getum vart gengið lengra.} \dicExampleCS{Stěží můžeme jít dál.}
\dicEntry[varta] \dicTerm{varta} \dicIPA{{v}{a}{\textsubring{r}}{\textsubring{d}}{a}} \dicPos{f}[1] \dicFlx{(vörtu, vörtur)}[20] \dicDirectTranslationCS{bradavice}
\dicEntry[vartappi] \dicTerm{var··tapp|i} \dicIPA{{v}{a}{\textsubring{r}}{t\smash{\textsuperscript{h}}}{a}{h}{\textsubring{b}}{\textsci}} \dicPos{m}[1] \dicFlx{(‑a, ‑ar)}[8] \dicFieldCat{elek.} \dicSynonym{öryggi} \dicDirectTranslationCS{pojistka}
\dicEntry[varúð] \dicTerm{var··úð} \dicIPA{{v}{a}{\textlengthmark}{r}{u}{\texttheta}} \dicPos{f}[4] \dicFlx{(‑ar)}[3] \textbf{1.} \dicSynonym{aðvörun} \dicDirectTranslationCS{výstraha, varování}  \textbf{2.} \dicSynonym{aðgætni} \dicDirectTranslationCS{opatrnost, ostražitost}
\dicEntry[varúlfur] \dicTerm{var··úlf|ur}\dicTerm{, verúlfur}\dicTerm{,  vargúlfur} \dicIPA{{v}\-{a}\-{\textlengthmark}\-{r}\-{u}\-{l}\-{v}\-{\textscy}\-{\textsubring{r}}\-} \dicPos{m}[6] \dicFlx{(‑s, ‑ar)}[4] \dicFieldCat{pov.} \dicDirectTranslationCS{vlkodlak}
\dicEntry[vasabók] \dicTerm{vasa··|bók} \dicIPA{{v}{a}{\textlengthmark}{s}{a}{\textsubring{b}}{ou}{\r{g}}} \dicPos{f}[8] \dicFlx{(‑bókar, ‑bækur)}[5] \textbf{1.} \dicSynonym*{lítil skrifbók} \dicDirectTranslationCS{notes, notýsek, blok}  \textbf{2.} \dicSynonym*{lítil bók} \dicDirectTranslationCS{knížka do kapsy}
\dicEntry[vasahnífur] \dicTerm{vasa··hníf|ur} \dicIPA{{v}{a}{\textlengthmark}{s}{a}{\textsubring{n}}{i}{v}{\textscy}{\textsubring{r}}} \dicPos{m}[6] \dicFlx{(‑s, ‑ar)}[22] \dicDirectTranslationCS{kapesní nůž, kudla}
\dicEntry[vasaklútur] \dicTerm{vasa··klút|ur} \dicsymFrequent\  \dicIPA{{v}{a}{\textlengthmark}{s}{a}{k\smash{\textsuperscript{h}}}{l}{u}{\textsubring{d}}{\textscy}{\textsubring{r}}} \dicPos{m}[6] \dicFlx{(‑s, ‑ar)}[22] \dicDirectTranslationCS{kapesník} \dicExampleIS{velktur vasaklútur} \dicExampleCS{zmuchlaný kapesník}
\dicEntry[vasaljós] \dicTerm{vasa··ljós} \dicIPA{{v}{a}{\textlengthmark}{s}{a}{l}{j}{ou}{s}} \dicPos{n}[2] \dicFlx{(‑s, ‑)}[5] \dicDirectTranslationCS{baterka, kapesní svítilna}
\dicEntry[vasaminni] \dicTerm{vasa··minni} \dicIPA{{v}{a}{\textlengthmark}{s}{a}{m}{\textsci}{n}{\textsci}} \dicPos{n}[2] \dicFlx{(‑s, ‑)}[14] \dicFieldCat{poč.} \dicDirectTranslationCS{flash (paměť)}
\dicEntry[vasaorðabók] \dicTerm{vasa··orða·|bók} \dicIPA{{v}{a}{\textlengthmark}{s}{a}{\textopeno}{r}{ð}{a}{\textsubring{b}}{ou}{\r{g}}} \dicPos{f}[8] \dicFlx{(‑bókar, ‑bækur)}[5] \dicDirectTranslationCS{kapesní slovník}
\dicEntry[vasapeli] \dicTerm{vasa··pel|i} \dicIPA{{v}{a}{\textlengthmark}{s}{a}{p\smash{\textsuperscript{h}}}{\textepsilon}{l}{\textsci}} \dicPos{m}[1] \dicFlx{(‑a, ‑ar)}[1] \dicDirectTranslationCS{kapesní láhev, placatka}
\dicEntry[vasapeningar] \dicTerm{vasa··pen·ingar} \dicIPA{{v}{a}{\textlengthmark}{s}{a}{p\smash{\textsuperscript{h}}}{\textepsilon}{n}{i}{\ng}{\r{g}}{a}{\textsubring{r}}} \dicPos{m}[6] \dicFlx{pl}[2] \dicDirectTranslationCS{kapesné}
\dicEntry[vasareiknir] \dicTerm{vasa··reikn|ir} \dicIPA{{v}{a}{\textlengthmark}{s}{a}{r}{ei}{h}{\r{g}}{n}{\textsci}{\textsubring{r}}} \dicPos{m}[7] \dicFlx{(‑is)}[2] \dicDirectTranslationCS{kalkulačka}
\dicEntry[vasatölva] \dicTerm{vasa··tölv|a} \dicIPA{{v}{a}{\textlengthmark}{s}{a}{t\smash{\textsuperscript{h}}}{\oe}{l}{v}{a}} \dicPos{f}[1] \dicFlx{(‑u, ‑ur)}[7] \dicFieldCat{poč.} \dicDirectTranslationCS{kapesní počítač}
\dicEntry[vasaúr] \dicTerm{vasa··úr} \dicIPA{{v}{a}{\textlengthmark}{s}{au}{\textsubring{r}}} \dicPos{n}[2] \dicFlx{(‑s, ‑)}[5] \dicDirectTranslationCS{kapesní hodinky}
\dicEntry[vasaútgáfa] \dicTerm{vasa··út·gáf|a} \dicIPA{{v}{a}{\textlengthmark}{s}{au}{\textsubring{d}}{\r{g}}{au}{v}{a}} \dicPos{f}[1] \dicFlx{(‑u, ‑ur)}[13] \textbf{1.} \dicDirectTranslationCS{kapesní vydání}  \textbf{2.} \dicLangCat{hovor.} \dicSynonym*{smærri eftirmynd} \dicDirectTranslationCS{menší podobizna} \dicExampleIS{Hann er vasaútgáfa af frænda sínum.} \dicExampleCS{Je takovou menší podobiznou svého strýčka.}
\dicEntry[vasaþjófur] \dicTerm{vasa··þjóf|ur} \dicIPA{{v}{a}{\textlengthmark}{s}{a}{\texttheta}{j}{ou}{v}{\textscy}{\textsubring{r}}} \dicPos{m}[6] \dicFlx{(‑s, ‑ar)}[22] \dicDirectTranslationCS{kapsář(ka), kapesní zloděj(ka)}
\dicEntry[vasi] \dicTerm{vas|i} \dicsymFrequent\  \dicIPA{{v}{a}{\textlengthmark}{s}{\textsci}} \dicPos{m}[1] \dicFlx{(‑a, ‑ar)}[8] \textbf{1.} \dicDirectTranslationCS{kapsa} \dicExampleIS{taka lyklana upp úr vasanum} \dicExampleCS{vytáhnout klíče z~kapsy};  \dicPhraseIS{borga úr eigin vasa} \dicLangCat{přen.} \dicDirectTranslationCS{zaplatit z~vlastní kapsy};  \dicPhraseIS{hafa e‑n í vasanum} \dicLangCat{přen.} \dicDirectTranslationCS{mít (koho) v~kapse}  \textbf{2.} \dicSynonym{krukka\smash{\textsuperscript{1}}} \dicDirectTranslationCS{váza} \dicExampleIS{blómavasi} \dicExampleCS{květinová váza}
\dicEntry[vasilín] \dicTerm{vasilín}\dicTerm{, vassilín} \dicIPA{{v}\-{a}\-{\textlengthmark}\-{s}\-{\textsci}\-{l}\-{i}\-{\textsubring{n}}\-} \dicPos{n}[2] \dicFlx{(‑s)}[2] \dicDirectTranslationCS{vazelína}
\dicEntry[vaska] \dicTerm{vask|a} \dicIPA{{v}{a}{s}{\r{g}}{a}} \dicPos{v}[1] \dicFlx{(‑aði)}[13] \dicPhraseIS{vaska upp} \dicDirectTranslationCS{(u)mýt nádobí};  \dicPhraseIS{vaska upp e‑ð} \dicDirectTranslationCS{(u)mýt (co) (talíře ap.)}
\dicEntry[vaskafat] \dicTerm{vaska··|fat} \dicIPA{{v}{a}{s}{\r{g}}{a}{f}{a}{\textsubring{d}}} \dicPos{n}[2] \dicFlx{(‑fats, ‑föt)}[8] \dicDirectTranslationCS{lavor, (plastové) umyvadlo, umývací nádoba}
\dicEntry[vaski] \dicTerm{vask|i} \dicIPA{{v}{a}{s}{\r{\textObardotlessj}}{\textsci}} \dicPos{m}[1] \dicFlx{(‑a, ‑ar)}[8] \dicLink{vaskur\smash{\textsuperscript{1}}}
\dicEntry[vasklega] \dicTerm{vask··lega} \dicIPA{{v}{a}{s}{\r{g}}{l}{\textepsilon}{\textbabygamma}{a}} \dicPos{adv} \dicDirectTranslationCS{energicky, rázně}
\dicEntry[vasklegur] \dicTerm{vask··legur} \dicIPA{{v}{a}{s}{\r{g}}{l}{\textepsilon}{\textbabygamma}{\textscy}{\textsubring{r}}} \dicPos{adj}[1]\dicFlx{}[-8] \dicSynonym{hvatlegur} \dicDirectTranslationCS{energický, rázný} \dicExampleIS{vaskleg framganga} \dicExampleCS{energický průběh}
\dicEntry[vaskleiki] \dicTerm{vask··leik|i} \dicIPA{{v}{a}{s}{\r{g}}{l}{ei}{\r{\textObardotlessj}}{\textsci}} \dicPos{m}[1] \dicFlx{(‑a)}[3] \dicSynonym{röskleiki} \dicDirectTranslationCS{ráznost, elán}
\dicEntry[vaskur] \dicTerm{vask|ur\smash{\textsuperscript{1}}}\dicTerm{, vaski} \dicIPA{{v}\-{a}\-{s}\-{\r{g}}\-{\textscy}\-{\textsubring{r}}\-} \dicPos{m}[6] \dicFlx{(‑s, ‑ar)}[5] \textbf{1.} \dicSynonym*{baðvaskur} \dicDirectTranslationCS{umyvadlo}  \textbf{2.} \dicSynonym*{uppþvottavaskur} \dicDirectTranslationCS{dřez};  \dicIdiom{vaskur}{ \dicPhraseIS{e‑að fer í vaskinn}} \dicLangCat{přen.} \dicDirectTranslationCS{(co) jde od desíti k~pěti}
\dicEntry[vaskur] \dicTerm{vaskur\smash{\textsuperscript{2}}} \dicIPA{{v}{a}{s}{\r{g}}{\textscy}{\textsubring{r}}} \dicPos{adj}[1] \dicFlx{(f vösk)}[2] \dicDirectTranslationCS{udatný, srdnatý, zmužilý}
\dicEntry[vassilín] \dicTerm{vassilín} \dicIPA{{v}{a}{s}{\textlengthmark}{\textsci}{l}{i}{\textsubring{n}}} \dicPos{n}[2] \dicFlx{(‑s)}[2] \dicLink{vasilín}
\dicEntry[Vatíkanið] \dicTerm{Vatíkan|ið} \dicIPA{{v}{a}{\textlengthmark}{\textsubring{d}}{i}{\r{g}}{a}{n}{\textsci}{\texttheta}} \dicPos{n}[2] \dicFlx{(‑sins)}[3] \dicFlx{def} \dicFieldCat{geog.} \dicSynonym{Páfagarður} \dicDirectTranslationCS{Vatikán}
\dicEntry[vatíkanskur] \dicTerm{vatíkanskur} \dicIPA{{v}{a}{\textlengthmark}{\textsubring{d}}{i}{\r{g}}{a}{n}{s}{\r{g}}{\textscy}{\textsubring{r}}} \dicPos{adj}[1] \dicFlx{(f vatíkönsk)}[3] \dicDirectTranslationCS{vatikánský}
\dicEntry[vatn] \dicTerm{vatn} \dicsymFrequent\  \dicIPA{{v}{a}{h}{\textsubring{d}}{\textsubring{n}}} \dicPos{n}[2] \dicFlx{(‑s, vötn)}[8] \textbf{1.} \dicSynonym{lögur} \dicDirectTranslationCS{voda} \dicExampleIS{glas af vatni} \dicExampleCS{sklenice vody};  \dicPhraseIS{eimað vatn} \dicDirectTranslationCS{destilovaná voda}  \textbf{2.} \dicSynonym{stöðuvatn} \dicDirectTranslationCS{jezero}  \textbf{3.} \dicSynonym{fljót} \dicDirectTranslationCS{řeka, vodní tok};  \dicIdiom{vatn}{ \dicPhraseIS{dýfa ekki hendi í kalt vatn}} \dicLangCat{přen.} \dicDirectTranslationCS{nezašpinit si ruce, nesmočit ruku v~studené vodě} \dicIndirectTranslationCS{(vyhýbat se obtížnému a~nepříjemnému)};  \dicPhraseIS{e‑að er vatn á myllu e‑rs} \dicLangCat{přen.} \dicDirectTranslationCS{(co) je voda na (čí) mlýn};  \dicPhraseIS{e‑að heldur ekki vatni} \dicLangCat{přen.} \dicDirectTranslationCS{(co) má (logické) mezery, (co) se nedrží logiky}; { \dicPhraseIS{e‑að svífur yfir vötnunum}} \dicLangCat{přen.} \dicDirectTranslationCS{(co) visí ve vzduchu};  \dicPhraseIS{eins og að drekka vatn} \dicLangCat{přen.} \dicDirectTranslationCS{jako facka (lehké ap.)};  \dicPhraseIS{eins og að skvetta vatni á gæs} \dicLangCat{přen.} \dicDirectTranslationCS{házet hrách na zeď} \dicIndirectTranslationCS{(o~něčem bezúčelném)}; { \dicPhraseIS{eins og fiskur í vatni}} \dicLangCat{přen.} \dicDirectTranslationCS{jako ryba ve vodě};  \dicPhraseIS{e‑r fær vatn í munninn} \dicLangCat{přen.} \dicDirectTranslationCS{(komu) se sbíhají sliny};  \dicPhraseIS{e‑r heldur hvorki vatni né vindi} \dicLangCat{přen.} \dicDirectTranslationCS{(kdo) je ukecaný, (kdo) všechno vyžvaní, (kdo) neudrží žádné tajemství};  \dicPhraseIS{fiska í gruggugu vatni} \dicLangCat{přen.} \dicDirectTranslationCS{pustit se na tenký led} \dicIndirectTranslationCS{(pokoušet se o~úspěch riskantní cestou)}; { \dicPhraseIS{kasta af sér vatni}} \dicSynonym{pissa} \dicDirectTranslationCS{čůrat, vyčůrat se, močit, vymočit se};  \dicPhraseIS{mega ekki vatni halda af e‑u} \dicLangCat{přen.} \dicDirectTranslationCS{neubránit se slzám (čeho) (radosti ap.)};  \dicPhraseIS{mikið vatn hefur runnið til sjávar (síðan þetta gerðist)} \dicLangCat{přen.} \dicDirectTranslationCS{uběhlo mnoho vody (od té doby, co se to stalo)};  \dicPhraseIS{sækja vatn yfir lækinn} \dicLangCat{přen.} \dicDirectTranslationCS{chodit pro vodu přes potok} \dicIndirectTranslationCS{(uplatňovat značně obtížnější řešení)};  \dicPhraseIS{þetta kemur allt með kalda vatninu} \dicLangCat{přen.} \dicDirectTranslationCS{všechno má svůj čas}
\dicEntry[vatna] \dicTerm{vatn|a} \dicIPA{{v}{a}{h}{\textsubring{d}}{n}{a}} \dicPos{v}[1] \dicFlx{(‑aði)}[49] \dicFlx{dat} \textbf{1.} \dicDirectTranslationCS{napojit, napájet, dát\,/\addthin dávat napít} \dicExampleIS{vatna hestunum} \dicExampleCS{napojit koně}  \textbf{2.} \dicDirectTranslationCS{zalít, zalévat, zavlažit, zavlažovat} \dicExampleIS{vatna blómunum} \dicExampleCS{zalévat květiny};  \dicIdiom{vatna}{ \dicPhraseIS{vatna músum}} \dicSynonym{gráta} \dicDirectTranslationCS{plakat, ronit slzy}
\dicEntry[vatnafræði] \dicTerm{vatna··fræð|i} \dicIPA{{v}{a}{h}{\textsubring{d}}{n}{a}{f}{r}{a}{i}{ð}{\textsci}} \dicPos{f}[3] \dicFlx{(‑i)}[3] \dicDirectTranslationCS{hydrologie}
\dicEntry[vatnafræðingur] \dicTerm{vatna·fræð··ing|ur} \dicIPA{{v}{a}{h}{\textsubring{d}}{n}{a}{f}{r}{a}{i}{ð}{i}{\ng}{\r{g}}{\textscy}{\textsubring{r}}} \dicPos{m}[6] \dicFlx{(‑s, ‑ar)}[8] \dicDirectTranslationCS{hydrolog, hydroložka}
\dicEntry[Vatnajökull] \dicTerm{Vatna··jök|ull} \dicIPA{{v}{a}{h}{\textsubring{d}}{n}{a}{j}{\oe}{\r{g}}{\textscy}{\textsubring{d}}{\textsubring{l}}} \dicPos{m}[6] \dicFlx{(‑uls)}[30] \dicFieldCat{geog.} \dicDirectTranslationCS{Vatna\-jökull} \dicIndirectTranslationCS{(největší ledovec Islandu)}
\dicEntry[vatnakarpi] \dicTerm{vatna··karp|i} \dicIPA{{v}{a}{h}{\textsubring{d}}{n}{a}{k\smash{\textsuperscript{h}}}{a}{\textsubring{r}}{\textsubring{b}}{\textsci}} \dicPos{m}[1] \dicFlx{(‑a, ‑ar)}[8] \dicFieldCat{zool.} \dicDirectTranslationCS{kapr obecný} \textit{(l.~{\textLA{Cyprinus carpio}})}  \dicsymPhoto\ 
\dicFigure{63055.jpg}{Vatnakarpi}{Vatnakarpi - Spurný Vladimír, Biolib, PD}
\dicEntry[vatnakrabbi] \dicTerm{vatna··krabb|i} \dicIPA{{v}{a}{h}{\textsubring{d}}{n}{a}{k\smash{\textsuperscript{h}}}{r}{a}{\textsubring{b}}{\textsci}} \dicPos{m}[1] \dicFlx{(‑a, ‑ar)}[8] \dicFieldCat{zool.} \dicDirectTranslationCS{rak, rak říční} \textit{(l.~{\textLA{Astacus astacus}})}  \dicsymPhoto\ 
\dicFigure{1373.jpg}{Vatnakrabbi}{Vatnakrabbi - Zicha Ondřej, Biolib, Copyright/CC-BY-NC}
\dicEntry[vatnalilja] \dicTerm{vatna··lilj|a} \dicIPA{{v}{a}{h}{\textsubring{d}}{n}{a}{l}{\textsci}{l}{j}{a}} \dicPos{f}[1] \dicFlx{(‑u, ‑ur)}[7] \dicFieldCat{bot.} \dicDirectTranslationCS{leknín} \textit{(l.~{\textLA{Nymphaea}})}  \dicsymPhoto\ 
\dicFigure{9205.jpg}{Vatnalilja}{Vatnalilja - Deml Miroslav, Biolib, Copyright/CC-BY}
\dicEntry[vatnalíffræði] \dicTerm{vatna··líf·fræð|i} \dicIPA{{v}{a}{h}{\textsubring{d}}{n}{a}{l}{i}{f}{r}{a}{i}{ð}{\textsci}} \dicPos{f}[3] \dicFlx{(‑i)}[3] \dicDirectTranslationCS{hydrobiologie}
\dicEntry[vatnamælingar] \dicTerm{vatna··mæl·ingar} \dicIPA{{v}{a}{h}{\textsubring{d}}{n}{a}{m}{a}{i}{l}{i}{\ng}{\r{g}}{a}{\textsubring{r}}} \dicPos{f}[4] \dicFlx{pl}[6] \dicDirectTranslationCS{hydrografie}
\dicEntry[vatnaskil] \dicTerm{vatna··skil} \dicIPA{{v}{a}{h}{\textsubring{d}}{n}{a}{s}{\r{\textObardotlessj}}{\textsci}{\textsubring{l}}} \dicPos{n}[2] \dicFlx{pl}[1] \textbf{1.} \dicFieldCat{geol.} \dicDirectTranslationCS{vodní předěl, rozvodí}  \textbf{2.} \dicSynonym{tímamót} \dicDirectTranslationCS{(časový) předěl, mezník}
\dicEntry[vatnaskíði] \dicTerm{vatna··skíði} \dicIPA{{v}{a}{h}{\textsubring{d}}{n}{a}{s}{\r{\textObardotlessj}}{i}{ð}{\textsci}} \dicPos{n}[2] \dicFlx{(‑s, ‑)}[14] \dicFieldCat{sport.} \dicDirectTranslationCS{vodní lyže}
\dicEntry[vatnsafl] \dicTerm{vatns··afl} \dicIPA{{v}{a}{s}{\textlengthmark}{a}{\textsubring{b}}{\textsubring{l}}} \dicPos{n}[2] \dicFlx{(‑s)}[2] \dicSynonym{vatnsorka} \dicDirectTranslationCS{vodní síla\,/\addthin energie}
\dicEntry[vatnsaflsvirkjun] \dicTerm{vatns·afls··virkj|un} \dicIPA{{v}{a}{s}{\textlengthmark}{a}{\textsubring{b}}{\textsubring{l}}{s}{v}{\textsci}{\textsubring{r}}{\r{\textObardotlessj}}{\textscy}{\textsubring{n}}} \dicPos{f}[7] \dicFlx{(‑unar, ‑anir)}[8] \dicDirectTranslationCS{využívání vodní energie na výrobu elektřiny}
\dicEntry[vatnsberi] \dicTerm{vatns··ber|i} \dicIPA{{v}{a}{s}{\textsubring{b}}{\textepsilon}{r}{\textsci}} \dicPos{m}[1] \dicFlx{(‑a, ‑ar)}[1] \textbf{1.} \dicDirectTranslationCS{nosič(ka) vody}  \textbf{2.} \dicDirectTranslationCS{Vodnář} \dicIndirectTranslationCS{(znamení zvěrokruhu)} \textit{(l.~{\textLA{Aquarius}})}  \dicsymPhoto\   \textbf{3.} \dicFieldCat{bot.} \dicDirectTranslationCS{orlíček} \textit{(l.~{\textLA{Aquilegia}})}
\dicFigure{ds_image_vatnsberi_0_1.jpg}{Vatnsberi}{Vatnsberi - Danny Steaven, CC BY 1.0}
\dicEntry[vatnsborð] \dicTerm{vatns··borð} \dicIPA{{v}{a}{s}{\textsubring{b}}{\textopeno}{r}{\texttheta}} \dicPos{n}[2] \dicFlx{(‑s)}[2] \dicDirectTranslationCS{(vodní) hladina}
\dicEntry[vatnsból] \dicTerm{vatns··ból} \dicIPA{{v}{a}{s}{\textsubring{b}}{ou}{\textsubring{l}}} \dicPos{n}[2] \dicFlx{(‑s, ‑)}[5] \dicDirectTranslationCS{vodní zdroj, napajedlo}
\dicEntry[vatnsbyssa] \dicTerm{vatns··byss|a} \dicIPA{{v}{a}{s}{\textsubring{b}}{\textsci}{s}{a}} \dicPos{f}[1] \dicFlx{(‑u, ‑ur)}[19] \dicDirectTranslationCS{vodní pistole} \dicIndirectTranslationCS{(hračka)}
\dicEntry[vatnsfall] \dicTerm{vatns··|fall} \dicIPA{{v}{a}{s}{f}{a}{\textsubring{d}}{\textsubring{l}}} \dicPos{n}[2] \dicFlx{(‑falls, ‑föll)}[8] \dicSynonym{elfur} \dicDirectTranslationCS{řeka, (vodní) tok}
\dicEntry[vatnsfælni] \dicTerm{vatns··fæln|i} \dicIPA{{v}{a}{s}{f}{a}{i}{l}{n}{\textsci}} \dicPos{f}[3] \dicFlx{(‑i)}[3] \dicFieldCat{psych.} \dicDirectTranslationCS{hydrofobie}
\dicEntry[vatnsgeymir] \dicTerm{vatns··geym|ir} \dicIPA{{v}{a}{s}{\r{\textObardotlessj}}{ei}{m}{\textsci}{\textsubring{r}}} \dicPos{m}[7] \dicFlx{(‑is, ‑ar)}[1] \dicDirectTranslationCS{nádrž na vodu, vodní nádrž, cisterna s~vodou}
\dicEntry[vatnsglas] \dicTerm{vatns··|glas} \dicIPA{{v}{a}{s}{\r{g}}{l}{a}{s}} \dicPos{n}[2] \dicFlx{(‑glass, ‑glös)}[8] \textbf{1.} \dicSynonym*{glas fyrir vatn} \dicDirectTranslationCS{sklenice na vodu}  \textbf{2.} \dicSynonym*{glas af vatni} \dicDirectTranslationCS{sklenice vody}
\dicEntry[vatnsgufa] \dicTerm{vatns··guf|a} \dicIPA{{v}{a}{s}{\r{g}}{\textscy}{v}{a}} \dicPos{f}[1] \dicFlx{(‑u, ‑ur)}[7] \dicDirectTranslationCS{vodní pára}
\dicEntry[vatnsheldur] \dicTerm{vatns··heldur} \dicIPA{{v}{a}{s}{h}{\textepsilon}{l}{\textsubring{d}}{\textscy}{\textsubring{r}}} \dicPos{adj}[2]\dicFlx{}[-14] \dicDirectTranslationCS{nepromokavý, voděodolný, vodotěsný} \dicExampleIS{Úrið er vatnshelt.} \dicExampleCS{Hodinky jsou vodotěsné.}
\dicEntry[vatnshræddur] \dicTerm{vatns··hræddur} \dicIPA{{v}{a}{s}{\textsubring{r}}{a}{i}{\textsubring{d}}{\textscy}{\textsubring{r}}} \dicPos{adj}[2]\dicFlx{}[-18] \dicDirectTranslationCS{mající strach z~vody, bojící se vody}
\dicEntry[vatnshræðsla] \dicTerm{vatns··hræðsl|a} \dicIPA{{v}{a}{s}{\textsubring{r}}{a}{i}{ð}{s}{\textsubring{d}}{l}{a}} \dicPos{f}[1] \dicFlx{(‑u)}[5] \dicDirectTranslationCS{strach z~vody}
\dicEntry[vatnskassi] \dicTerm{vatns··kass|i} \dicIPA{{v}{a}{s}{k\smash{\textsuperscript{h}}}{a}{s}{\textsci}} \dicPos{m}[1] \dicFlx{(‑a, ‑ar)}[8] \textbf{1.} \dicDirectTranslationCS{(splachovací) nádržka}  \textbf{2.} \dicDirectTranslationCS{chladič} \dicIndirectTranslationCS{(v~autě)}
\dicEntry[vatnskrani] \dicTerm{vatns··kran|i} \dicIPA{{v}{a}{s}{k\smash{\textsuperscript{h}}}{r}{a}{n}{\textsci}} \dicPos{m}[1] \dicFlx{(‑a, ‑ar)}[8] \dicDirectTranslationCS{vodovodní kohoutek}
\dicEntry[vatnskæling] \dicTerm{vatns··kæl·ing} \dicIPA{{v}{a}{s}{c\smash{\textsuperscript{h}}}{a}{i}{l}{i}{\ng}{\r{g}}} \dicPos{f}[4] \dicFlx{(‑ar)}[7] \dicDirectTranslationCS{vodní chlazení}
\dicEntry[vatnslaus] \dicTerm{vatns··laus} \dicIPA{{v}{a}{s}{l}{\oe i}{s}} \dicPos{adj}[5]\dicFlx{}[-1] \dicDirectTranslationCS{(jsoucí) bez vody}
\dicEntry[vatnslás] \dicTerm{vatns··lás} \dicIPA{{v}{a}{s}{l}{au}{s}} \dicPos{m}[4] \dicFlx{(‑s, ‑ar)}[25] \textbf{1.} \dicFieldCat{techn.} \dicDirectTranslationCS{(instalační) sifon}  \textbf{2.} \dicFieldCat{techn.} \dicDirectTranslationCS{ventil} \dicIndirectTranslationCS{(zařízení řídící nebo zastavující průtok v~potrubí)}
\dicEntry[vatnsleiðsla] \dicTerm{vatns··leiðsl|a} \dicIPA{{v}{a}{s}{l}{ei}{ð}{s}{\textsubring{d}}{l}{a}} \dicPos{f}[1] \dicFlx{(‑u, ‑ur)}[19] \dicDirectTranslationCS{vodovod} \dicIndirectTranslationCS{(zařízení na rozvod vody)}
\dicEntry[vatnsleysi] \dicTerm{vatns··leysi} \dicIPA{{v}{a}{s}{l}{ei}{s}{\textsci}} \dicPos{n}[2] \dicFlx{(‑s)}[20] \dicDirectTranslationCS{nedostatek vody}
\dicEntry[vatnslitamynd] \dicTerm{vatns·lita··mynd} \dicIPA{{v}{a}{s}{l}{\textsci}{\textsubring{d}}{a}{m}{\textsci}{n}{\textsubring{d}}} \dicPos{f}[7] \dicFlx{(‑ar, ‑ir)}[1] \dicDirectTranslationCS{akvarel} \dicIndirectTranslationCS{(obraz malovaný vodovými barvami)}
\dicEntry[vatnslitur] \dicTerm{vatns··lit|ur} \dicIPA{{v}{a}{s}{l}{\textsci}{\textsubring{d}}{\textscy}{\textsubring{r}}} \dicPos{m}[10] \dicFlx{(‑ar, ‑ir)}[13] \dicDirectTranslationCS{akvarelová\,/\addthin vodová barva, vodovka}
\dicEntry[vatnsmelóna] \dicTerm{vatns··melón|a} \dicIPA{{v}{a}{s}{m}{\textepsilon}{l}{ou}{n}{a}} \dicPos{f}[1] \dicFlx{(‑u, ‑ur)}[7] \dicDirectTranslationCS{vodní meloun}
\dicEntry[vatnsmerki] \dicTerm{vatns··merki} \dicIPA{{v}{a}{s}{m}{\textepsilon}{\textsubring{r}}{\r{\textObardotlessj}}{\textsci}} \dicPos{n}[2] \dicFlx{(‑s, ‑)}[16] \dicDirectTranslationCS{vodoznak, vodotisk}
\dicEntry[vatnsmylla] \dicTerm{vatns··myll|a} \dicIPA{{v}{a}{s}{m}{\textsci}{l}{\textlengthmark}{a}} \dicPos{f}[1] \dicFlx{(‑u, ‑ur)}[19] \dicDirectTranslationCS{vodní mlýn}
\dicEntry[vatnsmælir] \dicTerm{vatns··mæl|ir} \dicIPA{{v}{a}{s}{m}{a}{i}{l}{\textsci}{\textsubring{r}}} \dicPos{m}[7] \dicFlx{(‑is, ‑ar)}[1] \dicDirectTranslationCS{vodoměr}
\dicEntry[vatnsnotkun] \dicTerm{vatns··notk|un} \dicIPA{{v}{a}{s}{n}{\textopeno}{\textsubring{d}}{\r{g}}{\textscy}{\textsubring{n}}} \dicPos{f}[7] \dicFlx{(‑unar)}[9] \dicDirectTranslationCS{spotřeba vody}
\dicEntry[vatnsorka] \dicTerm{vatns··ork|a} \dicIPA{{v}{a}{s}{\textlengthmark}{\textopeno}{\textsubring{r}}{\r{g}}{a}} \dicPos{f}[1] \dicFlx{(‑u)}[5] \dicSynonym{vatnsafl} \dicDirectTranslationCS{vodní energie}
\dicEntry[vatnsósa] \dicTerm{vatns··ósa} \dicIPA{{v}{a}{s}{\textlengthmark}{ou}{s}{a}} \dicPos{adj}[13] \dicFlx{indecl}[1] \dicDirectTranslationCS{promočený, promáčený, podmáčený} \dicExampleIS{Jarðvegurinn er vatnsósa.} \dicExampleCS{Země je podmáčená.}
\dicEntry[vatnsrennsli] \dicTerm{vatns··rennsli} \dicIPA{{v}{a}{s}{r}{\textepsilon}{n}{s}{\textsubring{d}}{l}{\textsci}} \dicPos{n}[2] \dicFlx{(‑s)}[20] \dicDirectTranslationCS{(prů)tok vody}
\dicEntry[vatnssalerni] \textls[21]{\dicTerm{vatns··salerni} \dicIPA{{v}{a}{s}{a}{l}{\textepsilon}{r}{\textsubring{d}}{n}{\textsci}} \dicPos{n}[2] \dicFlx{(‑s, ‑)}[14] \dicDirectTranslationCS{(splachovací) \mbox{záchod}}}
\dicEntry[vatnsskorpa] \dicTerm{vatns··skorp|a} \dicIPA{{v}{a}{s}{\r{g}}{\textopeno}{\textsubring{r}}{\textsubring{b}}{a}} \dicPos{f}[1] \dicFlx{(‑u)}[5] \dicDirectTranslationCS{vodní hladina}
\dicEntry[vatnssopi] \dicTerm{vatns··sop|i} \dicIPA{{v}{a}{s}{\textopeno}{\textsubring{b}}{\textsci}} \dicPos{m}[1] \dicFlx{(‑a, ‑ar)}[1] \dicDirectTranslationCS{doušek vody}
\dicEntry[vatnsveita] \dicTerm{vatns··veit|a} \dicIPA{{v}{a}{s}{v}{ei}{\textsubring{d}}{a}} \dicPos{f}[1] \dicFlx{(‑u, ‑ur)}[13] \textbf{1.} \dicSynonym{vatnsleiðsla} \dicDirectTranslationCS{vodovodní systém}  \textbf{2.} \dicSynonym*{vatnsveitufyrirtæki} \dicDirectTranslationCS{vodárna} \dicIndirectTranslationCS{(závod dodávající vodu)}
\dicEntry[vatnsvirkjun] \dicTerm{vatns··virkj|un} \dicIPA{{v}{a}{s}{v}{\textsci}{\textsubring{r}}{\r{\textObardotlessj}}{\textscy}{\textsubring{n}}} \dicPos{f}[7] \dicFlx{(‑unar, ‑anir)}[8] \dicDirectTranslationCS{vodní elektrárna, hydroelektrárna}
\dicEntry[vatnsþéttur] \dicTerm{vatns··þéttur} \dicIPA{{v}{a}{s}{\texttheta}{j}{\textepsilon}{h}{\textsubring{d}}{\textscy}{\textsubring{r}}} \dicPos{adj}[1]\dicFlx{}[-10] \dicSynonym{vatnsheldur} \dicDirectTranslationCS{vodotěsný, nepromokavý} \dicExampleIS{Þakið er vatnsþétt.} \dicExampleCS{Střecha je nepromokavá.}
\dicEntry[vatnsþró] \dicTerm{vatns··|þró} \dicIPA{{v}{a}{s}{\texttheta}{r}{ou}} \dicPos{f}[4] \dicFlx{(‑þróar\,/\addthin ‑þrór, ‑þrór\,/\addthin ‑þrær)}[29] \dicDirectTranslationCS{(vodní) cisterna\,/\addthin nádrž}
\dicEntry[vatnsþrýstingur] \dicTerm{vatns··þrýst·ing|ur} \dicIPA{{v}{a}{s}{\texttheta}{r}{i}{s}{\textsubring{d}}{i}{\ng}{\r{g}}{\textscy}{\textsubring{r}}} \dicPos{m}[6] \dicFlx{(‑s)}[9] \dicDirectTranslationCS{tlak vody, hydrostatický tlak}
\dicEntry[vatt] \dicTerm{vatt\smash{\textsuperscript{1}}} \dicIPA{{v}{a}{h}{\textsubring{d}}} \dicPos{n}[2] \dicFlx{(‑s, vött)}[8] \textbf{1.} \dicFieldCat{fyz.} \dicDirectTranslationCS{watt}  \textbf{2.} \dicDirectTranslationCS{vatelín} \dicIndirectTranslationCS{(řídká osnovní pletená tkanina)}
\dicEntry[vatt] \dicTerm{vatt\smash{\textsuperscript{2}}} \dicIPA{{v}{a}{h}{\textsubring{d}}} \dicPos{v} \dicFlx{ind pf sg 1 pers} \dicLink{vinda\smash{\textsuperscript{2}}}
\dicEntry[vax] \dicTerm{vax} \dicIPA{{v}{a}{x}{s}} \dicPos{n}[2] \dicFlx{(‑)}[24] \dicDirectTranslationCS{vosk}
\dicEntry[vaxa] \dicTerm{vaxa} \dicsymFrequent\  \dicIPA{{v}{a}{x}{s}{a}} \dicPos{v}[6] \dicFlx{(vex, óx, uxum, yxi, vaxið)}[117] \textbf{1.} \dicSynonym{stækka} \dicDirectTranslationCS{(vy)růst, vyrůstat, narůst, narůstat} \dicExampleIS{Strákurinn er enn að vaxa.} \dicExampleCS{Chlapec stále roste.}  \textbf{2.} \dicSynonym*{aukast} \dicDirectTranslationCS{růst, vzrůstat, zvedat se} \dicExampleIS{Áin vex.} \dicExampleCS{Voda v~řece roste.};  \dicIdiom{vaxa}[í]{ \dicPhraseIS{e‑m vex e‑að í augum}} \dicFlx{impers} \dicLangCat{přen.} \dicDirectTranslationCS{(co) se (komu) jeví větší, než ve skutečnosti je (problém ap.)};  \dicIdiom{vaxa}[upp úr]{ \dicPhraseIS{vaxa upp úr e‑u}} \dicDirectTranslationCS{vyrůst z~(čeho)}
\dicEntry[vaxandi] \dicTerm{vax··andi} \dicsymFrequent\  \dicIPA{{v}{a}{x}{s}{a}{n}{\textsubring{d}}{\textsci}} \dicPos{adj}[13] \dicFlx{indecl}[1] \dicDirectTranslationCS{rostoucí, narůstající, vzrůstající};  \dicPhraseIS{e‑að fer vaxandi} \dicDirectTranslationCS{(co) narůstá, (co) roste (zájem ap.)};  \dicPhraseIS{vaxandi tungl} \dicDirectTranslationCS{dorůstající\,/\addthin přibývající měsíc}
\dicEntry[vaxdúkur] \dicTerm{vax··dúk|ur} \dicIPA{{v}{a}{x}{s}{\textsubring{d}}{u}{\r{g}}{\textscy}{\textsubring{r}}} \dicPos{m}[6] \dicFlx{(‑s, ‑ar)}[22] \dicDirectTranslationCS{voskové\,/\addthin voskované plátno}
\dicEntry[vaxið] \dicTerm{vaxið} \dicIPA{{v}{a}{x}{s}{\textsci}{\texttheta}} \dicPos{v} \dicFlx{supin} \dicLink{vaxa}
\dicEntry[vaxinn] \dicTerm{vaxinn} \dicIPA{{v}{a}{x}{s}{\textsci}{\textsubring{n}}} \dicPos{adj}[6]\dicFlx{}[-4] \textbf{1.} \dicSynonym{gróinn} \dicDirectTranslationCS{porostlý, zarostlý}  \textbf{2.} \dicDirectTranslationCS{vyrostlý, urostlý, vzrostlý};  \dicPhraseIS{vera vel vaxinn} \dicDirectTranslationCS{být urostlý, být pěkně rostlý};  \dicPhraseIS{vera illa vaxinn} \dicDirectTranslationCS{být špatně rostlý, mít ošklivou postavu};  \dicPhraseIS{vera vaxinn upp úr e‑u} {\textbf{a.}} \dicDirectTranslationCS{vyrůst z~(čeho) (z~oblečení ap.)};  {\textbf{b.}} \dicLangCat{přen.} \dicDirectTranslationCS{vyrůst z~(čeho) (klukovin ap.)}
\dicEntry[vaxkaka] \dicTerm{vax··|kaka} \dicIPA{{v}{a}{x}{s}{k\smash{\textsuperscript{h}}}{a}{\r{g}}{a}} \dicPos{f}[1] \dicFlx{(‑köku, ‑kökur)}[20] \dicDirectTranslationCS{plástev, voština}
\dicEntry[vaxkerti] \dicTerm{vax··kerti} \dicIPA{{v}{a}{x}{s}{c\smash{\textsuperscript{h}}}{\textepsilon}{\textsubring{r}}{\textsubring{d}}{\textsci}} \dicPos{n}[2] \dicFlx{(‑s, ‑)}[14] \dicDirectTranslationCS{vosková svíce}
\dicEntry[vaxlitur] \dicTerm{vax··lit|ur} \dicIPA{{v}{a}{x}{s}{l}{\textsci}{\textsubring{d}}{\textscy}{\textsubring{r}}} \dicPos{m}[10] \dicFlx{(‑ar, ‑ir)}[13] \dicDirectTranslationCS{vosková barva, voskovka}
\dicEntry[vaxmynd] \dicTerm{vax··mynd} \dicIPA{{v}{a}{x}{s}{m}{\textsci}{n}{\textsubring{d}}} \dicPos{f}[7] \dicFlx{(‑ar, ‑ir)}[1] \dicDirectTranslationCS{vosková figurína}
\dicEntry[vaxmyndasafn] \dicTerm{vax·mynda··|safn} \dicIPA{{v}{a}{x}{s}{m}{\textsci}{n}{\textsubring{d}}{a}{s}{a}{\textsubring{b}}{\textsubring{n}}} \dicPos{n}[2] \dicFlx{(‑safns, ‑söfn)}[8] \dicDirectTranslationCS{muzeum voskových figurín}
\dicEntry[vaxtabyrði] \dicTerm{vaxta··byrð|i} \dicIPA{{v}{a}{x}{s}{\textsubring{d}}{a}{\textsubring{b}}{\textsci}{r}{ð}{\textsci}} \dicPos{f}[2] \dicFlx{(‑i\,/\addthin ‑ar)}[4] \dicFieldCat{ekon.} \dicDirectTranslationCS{úroková zátěž}
\dicEntry[vaxtabætur] \dicTerm{vaxta··bætur} \dicIPA{{v}{a}{x}{s}{\textsubring{d}}{a}{\textsubring{b}}{a}{i}{\textsubring{d}}{\textscy}{\textsubring{r}}} \dicPos{f}[8] \dicFlx{pl}[6] \dicFieldCat{ekon.} \dicDirectTranslationCS{úleva z~úroků}
\dicEntry[vaxtagreiðsla] \dicTerm{vaxta··greiðsl|a} \dicIPA{{v}{a}{x}{s}{\textsubring{d}}{a}{\r{g}}{r}{ei}{ð}{s}{\textsubring{d}}{l}{a}} \dicPos{f}[1] \dicFlx{(‑u, ‑ur)}[13] \dicFieldCat{ekon.} \dicDirectTranslationCS{úroková platba}
\dicEntry[vaxtakjör] \dicTerm{vaxta··kjör} \dicIPA{{v}{a}{x}{s}{\textsubring{d}}{a}{c\smash{\textsuperscript{h}}}{\oe}{\textsubring{r}}} \dicPos{n}[2] \dicFlx{pl}[9] \dicFieldCat{ekon.} \dicDirectTranslationCS{úrokové podmínky}
\dicEntry[vaxtakostnaður] \dicTerm{vaxta··kost·nað|ur} \dicIPA{{v}{a}{x}{s}{\textsubring{d}}{a}{k\smash{\textsuperscript{h}}}{\textopeno}{s}{\textsubring{d}}{n}{a}{ð}{\textscy}{\textsubring{r}}} \dicPos{m}[10] \dicFlx{(‑ar)}[9] \dicFieldCat{ekon.} \dicDirectTranslationCS{úrokové náklady}
\dicEntry[vaxtalaus] \dicTerm{vaxta··laus} \dicIPA{{v}{a}{x}{s}{\textsubring{d}}{a}{l}{\oe i}{s}} \dicPos{adj}[5]\dicFlx{}[-4] \dicFieldCat{ekon.} \dicDirectTranslationCS{bezúročný} \dicExampleIS{vaxtalaust lán} \dicExampleCS{bezúročná půjčka}
\dicEntry[vaxtamunur] \dicTerm{vaxta··mun|ur} \dicIPA{{v}{a}{x}{s}{\textsubring{d}}{a}{m}{\textscy}{n}{\textscy}{\textsubring{r}}} \dicPos{m}[10] \dicFlx{(‑ar)}[12] \dicFieldCat{ekon.} \dicDirectTranslationCS{úroková marže}
\dicEntry[vaxtar] \dicTerm{vaxtar} \dicIPA{{v}{a}{x}{s}{\textsubring{d}}{a}{\textsubring{r}}} \dicPos{m} \dicFlx{sg gen} \dicLink{vöxtur}
\dicEntry[vaxtarbroddur] \textls[10]{\dicTerm{vaxtar··brodd|ur} \dicIPA{{v}{a}{x}{s}{\textsubring{d}}{a}{r}{\textsubring{b}}{r}{\textopeno}{\textsubring{d}}{\textscy}{\textsubring{r}}} \dicPos{m}[6] \dicFlx{(‑s, ‑ar)}[4]} \dicDirectTranslationCS{vrchol růstu\,/\addthin rozvoje, plný rozkvět} \dicExampleIS{vaxtarbroddur vestrænnar menningar} \dicExampleCS{vrchol růstu západní kultury}
\dicEntry[vaxtargeta] \dicTerm{vaxtar··get|a} \dicIPA{{v}{a}{x}{s}{\textsubring{d}}{a}{r}{\r{\textObardotlessj}}{\textepsilon}{\textsubring{d}}{a}} \dicPos{f}[1] \dicFlx{(‑u)}[5] \dicDirectTranslationCS{růstový potenciál}
\dicEntry[vaxtarhormón] \dicTerm{vaxtar··hormón} \dicIPA{{v}{a}{x}{s}{\textsubring{d}}{a}{\textsubring{r}}{h}{\textopeno}{r}{m}{ou}{\textsubring{n}}} \dicPos{n}[2] \dicFlx{(‑s, ‑)}[5] \dicFieldCat{biol.} \dicDirectTranslationCS{růstový hormon}
\dicEntry[vaxtarhraði] \textls[10]{\dicTerm{vaxtar··hrað|i} \dicIPA{{v}{a}{x}{s}{\textsubring{d}}{a}{\textsubring{r}}{a}{ð}{\textsci}} \dicPos{m}[1] \dicFlx{(‑a)}[3] \dicDirectTranslationCS{rychlost\,/\addthin tempo růstu}}
\dicEntry[vaxtarskeið] \textls[15]{\dicTerm{vaxtar··skeið} \dicIPA{{v}{a}{x}{s}{\textsubring{d}}{a}{\textsubring{r}}{s}{\r{\textObardotlessj}}{ei}{\texttheta}} \dicPos{n}[2] \dicFlx{(‑s, ‑)}[5] \dicDirectTranslationCS{období růstu}} \dicExampleIS{langt vaxtarskeið eftir kreppuna} \dicExampleCS{dlouhé období růstu po krizi}
\dicEntry[vaxtarskilyrði] \dicTerm{vaxtar··skil·yrði} \dicIPA{{v}{a}{x}{s}{\textsubring{d}}{a}{\textsubring{r}}{s}{\r{\textObardotlessj}}{\textsci}{l}{\textsci}{r}{ð}{\textsci}} \dicPos{n}[2] \dicFlx{(‑s, ‑)}[14] \dicDirectTranslationCS{růstové podmínky}
\dicEntry[vaxtavextir] \dicTerm{vaxta··vextir} \dicIPA{{v}{a}{x}{s}{\textsubring{d}}{a}{v}{\textepsilon}{x}{s}{\textsubring{d}}{\textsci}{\textsubring{r}}} \dicPos{m}[11] \dicFlx{pl}[1] \dicFieldCat{ekon.} \dicDirectTranslationCS{složený úrok, úrok z~úroku}
\dicEntry[vá] \dicTerm{vá\smash{\textsuperscript{1}}} \dicIPA{{v}{au}{\textlengthmark}} \dicPos{f}[4] \dicFlx{(‑r\,/\addthin ‑ar)}[22] \dicSynonym{hætta\smash{\textsuperscript{1}}} \dicDirectTranslationCS{ohrožení, nebezpečí};  \dicPhraseIS{það er vá fyrir dyrum} \dicLangCat{přen.} \dicDirectTranslationCS{nebezpečí číhá za dveřmi}
\dicEntry[vá] \dicTerm{vá\smash{\textsuperscript{2}}} \dicIPA{{v}{au}{\textlengthmark}} \dicPos{inter} \dicDirectTranslationCS{ty jo!, pane jo!, wow!}
\dicEntry[vátrygging] \dicTerm{vá··trygg·ing} \dicIPA{{v}{au}{\textlengthmark}{t\smash{\textsuperscript{h}}}{r}{\textsci}{\r{\textObardotlessj}}{i}{\ng}{\r{g}}} \dicPos{f}[4] \dicFlx{(‑ar, ‑ar)}[5] \dicFieldCat{ekon.} \dicSynonym{trygging} \dicDirectTranslationCS{pojištění}
\dicEntry[vátryggingariðgjald] \dicTerm{vá·tryggingar··ið·|gjald} \dicIPA{{v}\-{au}\-{\textlengthmark}\-{t\smash{\textsuperscript{h}}}\-{r}\-{\textsci}\-{\r{\textObardotlessj}}\-{i}\-{\ng}\-{\r{g}}\-{a}\-{r}\-{\textsci}\-{ð}\-{\r{\textObardotlessj}}\-{a}\-{l}\-{\textsubring{d}}\-} \dicPos{n}[2] \dicFlx{(‑gjalds, ‑gjöld)}[8] \dicFieldCat{ekon.} \dicDirectTranslationCS{pojistná prémie}
\dicEntry[vátryggingarskírteini] \dicTerm{vá·tryggingar··skír·teini} \dicIPA{{v}\-{au}\-{\textlengthmark}\-{t\smash{\textsuperscript{h}}}\-{r}\-{\textsci}\-{\r{\textObardotlessj}}\-{i}\-{\ng}\-{\r{g}}\-{a}\-{\textsubring{r}}\-{s}\-{\r{\textObardotlessj}}\-{i}\-{\textsubring{r}}\-{t\smash{\textsuperscript{h}}}\-{ei}\-{n}\-{\textsci}\-} \dicPos{n}[2] \dicFlx{(‑s, ‑)}[14] \dicFieldCat{ekon.} \dicDirectTranslationCS{pojistka, pojistná smlouva}
\dicEntry[vátryggingartaki] \dicTerm{vá·tryggingar··tak|i} \dicIPA{{v}{au}{\textlengthmark}{t\smash{\textsuperscript{h}}}{r}{\textsci}{\r{\textObardotlessj}}{i}{\ng}{\r{g}}{a}{\textsubring{r}}{t\smash{\textsuperscript{h}}}{a}{\r{\textObardotlessj}}{\textsci}} \dicPos{m}[1] \dicFlx{(‑a, ‑ar)}[8] \dicDirectTranslationCS{pojištěnec, pojištěnka}
\dicEntry[vátryggingarupphæð] \dicTerm{vá·tryggingar··upp·hæð} \dicIPA{{v}\-{au}\-{\textlengthmark}\-{t\smash{\textsuperscript{h}}}\-{r}\-{\textsci}\-{\r{\textObardotlessj}}\-{i}\-{\ng}\-{\r{g}}\-{a}\-{r}\-{\textscy}\-{h}\-{\textsubring{b}}\-{h}\-{a}\-{i}\-{\texttheta}\-} \dicPos{f}[7] \dicFlx{(‑ar, ‑ir)}[1] \dicFieldCat{ekon.} \dicDirectTranslationCS{pojistná částka}
\dicEntry[vátryggja] \dicTerm{vá··trygg|ja} \dicIPA{{v}{au}{\textlengthmark}{t\smash{\textsuperscript{h}}}{r}{\textsci}{\r{\textObardotlessj}}{a}} \dicPos{v}[2] \dicFlx{(‑ði, ‑t)}[89] \dicFlx{acc} \dicDirectTranslationCS{pojistit, pojišťovat} \dicExampleIS{vátryggja húsið} \dicExampleCS{pojistit dům}
\dicEntry[vátryggjandi] \dicTerm{vá··tryggj·|andi} \dicIPA{{v}{au}{\textlengthmark}{t\smash{\textsuperscript{h}}}{r}{\textsci}{\r{\textObardotlessj}}{a}{n}{\textsubring{d}}{\textsci}} \dicPos{m}[2] \dicFlx{(‑anda, ‑endur)}[1] \dicFieldCat{ekon.} \dicDirectTranslationCS{pojistitel(ka), pojišťovatel(ka)}
\dicEntry[vb.] \dicTerm{vb.} \dicPos{zkr} \dicPhraseIS{veik beyging} \dicFieldCat{jaz.} \dicDirectTranslationCS{slabé skloňování\,/\addthin časování}
\dicEntry[veð] \dicTerm{veð\smash{\textsuperscript{1}}} \dicIPA{{v}{\textepsilon}{\textlengthmark}{\texttheta}} \dicPos{n}[2] \dicFlx{(‑s, ‑)}[5] \textbf{1.} \dicDirectTranslationCS{zástava, hypotéka} \dicExampleIS{lán gegn veði í fasteign} \dicExampleCS{půjčka oproti zástavě nemovitosti};  \dicPhraseIS{setja e‑ð að veði} \dicDirectTranslationCS{zastavit (co)}  \textbf{2.} \dicSynonym{hætta\smash{\textsuperscript{1}}} \dicDirectTranslationCS{sázka, risk, riziko};  \dicPhraseIS{e‑að er í veði} \dicDirectTranslationCS{(co) je v~sázce};  \dicPhraseIS{leggja e‑ð að veði} \dicDirectTranslationCS{dát\,/\addthin dávat (co) v~sázku, riskovat (co)}
\dicEntry[veð] \dicTerm{veð\smash{\textsuperscript{2}}} \dicIPA{{v}{\textepsilon}{\textlengthmark}{\texttheta}} \dicPos{v} \dicFlx{ind praes sg 1 pers} \dicLink{vaða\smash{\textsuperscript{2}}}
\dicEntry[veðbanki] \dicTerm{veð··bank|i} \dicIPA{{v}{\textepsilon}{ð}{\textsubring{b}}{au}{\r{\textltailn}}{\r{\textObardotlessj}}{\textsci}} \dicPos{m}[1] \dicFlx{(‑a, ‑ar)}[8] \dicFieldCat{ekon.} \dicDirectTranslationCS{hypoteční banka}
\dicEntry[veðbók] \dicTerm{veð··|bók} \dicIPA{{v}{\textepsilon}{ð}{\textsubring{b}}{ou}{\r{g}}} \dicPos{f}[8] \dicFlx{(‑bókar, ‑bækur)}[5] \dicFieldCat{ekon.} \dicDirectTranslationCS{hypoteční registr}
\dicEntry[veðbókarvottorð] \dicTerm{veð·bókar··vott·orð} \dicIPA{{v}{\textepsilon}{ð}{\textsubring{b}}{ou}{\r{g}}{a}{r}{v}{\textopeno}{h}{\textsubring{d}}{\textopeno}{r}{\texttheta}} \dicPos{n}[2] \dicFlx{(‑s, ‑)}[5] \dicFieldCat{ekon.} \dicDirectTranslationCS{výpis z~hypotečního registru}
\dicEntry[veðbréf] \dicTerm{veð··bréf} \dicIPA{{v}{\textepsilon}{ð}{\textsubring{b}}{r}{j}{\textepsilon}{f}} \dicPos{n}[2] \dicFlx{(‑s, ‑)}[5] \dicFieldCat{ekon.} \dicDirectTranslationCS{hypoteční list}
\dicEntry[veðdeild] \dicTerm{veð··deild} \dicIPA{{v}{\textepsilon}{ð}{\textsubring{d}}{ei}{l}{\textsubring{d}}} \dicPos{f}[7] \dicFlx{(‑ar, ‑ir)}[1] \dicDirectTranslationCS{hypoteční oddělení}
\dicEntry[veðhafi] \dicTerm{veð··haf|i} \dicIPA{{v}{\textepsilon}{\textlengthmark}{\texttheta}{h}{a}{v}{\textsci}} \dicPos{m}[1] \dicFlx{(‑a, ‑ar)}[8] \dicFieldCat{práv.} \dicDirectTranslationCS{hypoteční věřitel(ka)}
\dicEntry[veðhlaup] \dicTerm{veð··hlaup} \dicIPA{{v}{\textepsilon}{ð}{\textsubring{l}}{\oe i}{\textsubring{b}}} \dicPos{n}[2] \dicFlx{(‑s, ‑)}[5] \dicDirectTranslationCS{(sázkový) dostih}
\dicEntry[veðhlaupahestur] \dicTerm{veð·hlaupa··hest|ur} \dicIPA{{v}{\textepsilon}{ð}{\textsubring{l}}{\oe i}{\textsubring{b}}{a}{h}{\textepsilon}{s}{\textsubring{d}}{\textscy}{\textsubring{r}}} \dicPos{m}[6] \dicFlx{(‑s, ‑ar)}[4] \dicDirectTranslationCS{dostihový kůň}
\dicEntry[veðhæfur] \dicTerm{veð··hæfur} \dicIPA{{v}{\textepsilon}{\textlengthmark}{\texttheta}{h}{a}{i}{v}{\textscy}{\textsubring{r}}} \dicPos{adj}[1]\dicFlx{}[-1] \dicFieldCat{práv.} \dicDirectTranslationCS{zastavitelný, použitelný jako zástava}
\dicEntry[veðja] \dicTerm{veðj|a} \dicIPA{{v}{\textepsilon}{ð}{j}{a}} \dicPos{v}[1] \dicFlx{(‑aði)}[1] \dicDirectTranslationCS{vsadit (se), sázet (se)};  \dicIdiom{veðja}[á]{ \dicPhraseIS{veðja e‑ð á e‑ð}} \dicDirectTranslationCS{vsadit (co) na (co)} \dicExampleIS{veðja stórfé á hross} \dicExampleCS{vsadit balík na koně};  \dicIdiom{veðja}[um]{ \dicPhraseIS{veðja um e‑ð}} \dicDirectTranslationCS{vsadit se o~(co)};  \dicIdiom{veðja}[við]{ \dicPhraseIS{veðja við e‑n um e‑ð}} \dicDirectTranslationCS{vsadit se s~(kým) o~(co)} \dicExampleIS{veðja við hana að ég verð fljótari en hún} \dicExampleCS{vsadit se s~ní, že budu rychlejší než ona};  \dicIdiom{veðjast}[við]{ \dicPhraseIS{veðjast við e‑n um e‑ð}} \dicFlx{refl} \dicDirectTranslationCS{vsadit se s~(kým) o~(co)}
\dicEntry[veðlán] \dicTerm{veð··lán} \dicIPA{{v}{\textepsilon}{ð}{l}{au}{\textsubring{n}}} \dicPos{n}[2] \dicFlx{(‑s, ‑)}[5] \dicFieldCat{ekon.} \dicDirectTranslationCS{zajištěná půjčka}
\dicEntry[veðlánari] \dicTerm{veð··lán·ar|i} \dicIPA{{v}{\textepsilon}{ð}{l}{au}{n}{a}{r}{\textsci}} \dicPos{m}[1] \dicFlx{(‑a, ‑ar)}[13] \dicDirectTranslationCS{zastavárník, zastavárnice}
\dicEntry[veðmangari] \dicTerm{veð··mang·ar|i} \dicIPA{{v}{\textepsilon}{ð}{m}{au}{\ng}{\r{g}}{a}{r}{\textsci}} \dicPos{m}[1] \dicFlx{(‑a, ‑ar)}[10] \dicDirectTranslationCS{bookmaker(ka)}
\dicEntry[veðmál] \dicTerm{veð··mál} \dicIPA{{v}{\textepsilon}{ð}{m}{au}{\textsubring{l}}} \dicPos{n}[2] \dicFlx{(‑s, ‑)}[5] \dicSynonym{veð\smash{\textsuperscript{1}}} \dicDirectTranslationCS{sázka} \dicExampleIS{Hún vann veðmálið.} \dicExampleCS{Vyhrála sázku.}
\dicEntry[veðra] \dicTerm{veðr|a} \dicIPA{{v}{\textepsilon}{ð}{r}{a}} \dicPos{v}[1] \dicFlx{(‑aði)}[1] \dicLangCat{zast.} \dicSynonym{hnusa} \dicDirectTranslationCS{čichat (vzduch ap.)};  \dicIdiom{veðrast}{ \dicPhraseIS{veðrast}} \dicFlx{refl} \dicFieldCat{geol.} \dicDirectTranslationCS{zvětrat, zvětrávat, (z)erodovat}
\dicEntry[veðraður] \dicTerm{veðr··|aður} \dicIPA{{v}{\textepsilon}{ð}{r}{a}{ð}{\textscy}{\textsubring{r}}} \dicPos{adj}[3] \dicFlx{(f ‑uð)}[3] \textbf{1.} \dicDirectTranslationCS{zvětralý, erodovaný}  \textbf{2.} \dicDirectTranslationCS{poničený\,/\addthin zničený špatným počasím a~větrem} \dicExampleIS{veðruð kirkja} \dicExampleCS{kostel poničený špatným počasím}
\dicEntry[veðrasamur] \dicTerm{veðra··|samur} \dicIPA{{v}{\textepsilon}{ð}{r}{a}{s}{a}{m}{\textscy}{\textsubring{r}}} \dicPos{adj}[1] \dicFlx{(f ‑söm)}[2] \dicDirectTranslationCS{vystavený špatnému počasí, bouřlivý} \dicIndirectTranslationCS{(o~místě)}
\dicEntry[veðrast] \dicTerm{veðr|ast} \dicIPA{{v}{\textepsilon}{ð}{r}{a}{s}{\textsubring{d}}} \dicPos{v}[1] \dicFlx{(‑aðist)}[95] \dicFlx{refl} \dicPhraseIS{veðrast upp við e‑ð} \dicDirectTranslationCS{zpychnout po (čem)} \dicExampleIS{veðrast upp við hrós} \dicExampleCS{zpychnout po pochvale}
\dicEntry[veðrátta] \dicTerm{veðr··átt|a} \dicIPA{{v}{\textepsilon}{ð}{r}{au}{h}{\textsubring{d}}{a}} \dicPos{f}[1] \dicFlx{(‑u)}[5] \dicSynonym{veðurfar} \dicDirectTranslationCS{podnebí, klima, počasí}
\dicEntry[veðreiðahestur] \dicTerm{veð·reiða··hest|ur} \dicIPA{{v}{\textepsilon}{ð}{r}{ei}{ð}{a}{h}{\textepsilon}{s}{\textsubring{d}}{\textscy}{\textsubring{r}}} \dicPos{m}[6] \dicFlx{(‑s, ‑ar)}[4] \dicDirectTranslationCS{dostihový kůň}
\dicEntry[veðreiðar] \dicTerm{veð··reiðar} \dicIPA{{v}{\textepsilon}{ð}{r}{ei}{ð}{a}{\textsubring{r}}} \dicPos{f}[7] \dicFlx{pl}[2] \dicDirectTranslationCS{dostihy}
\dicEntry[veðréttur] \dicTerm{veð··rétt|ur} \dicIPA{{v}{\textepsilon}{ð}{r}{j}{\textepsilon}{h}{\textsubring{d}}{\textscy}{\textsubring{r}}} \dicPos{m}[10] \dicFlx{(‑ar)}[7] \textbf{1.} \dicFieldCat{práv.} \dicSynonym{forgangsréttur} \dicDirectTranslationCS{zástavní právo} \dicIndirectTranslationCS{(věcné právo k~cizí věci)}  \textbf{2.} \dicFieldCat{práv.} \dicSynonym*{safn réttarreglna} \dicDirectTranslationCS{zástavní právo} \dicIndirectTranslationCS{(sbírka právních předpisů)}
\dicEntry[veðrun] \dicTerm{veðr|un} \dicIPA{{v}{\textepsilon}{ð}{r}{\textscy}{\textsubring{n}}} \dicPos{f}[7] \dicFlx{(‑unar)}[9] \dicFieldCat{geol.} \dicDirectTranslationCS{zvětrávání, eroze}
\dicEntry[veðset] \dicTerm{veð··set} \dicIPA{{v}{\textepsilon}{ð}{s}{\textepsilon}{\textsubring{d}}} \dicPos{v} \dicFlx{ind praes sg 1 pers} \dicLink{veðsetja}
\dicEntry[veðsetja] \dicTerm{veð··|setja} \dicIPA{{v}{\textepsilon}{ð}{s}{\textepsilon}{\textsubring{d}}{j}{a}} \dicPos{v}[4] \dicFlx{(‑set, ‑setti, ‑settum, ‑setti, ‑sett)}[26] \dicFlx{acc} \dicDirectTranslationCS{zastavit, zatížit hypotékou} \dicIndirectTranslationCS{(dát nějakou věc jako záruku za půjčku ap.)} \dicExampleIS{veðsetja hús fyrir láni} \dicExampleCS{zatížit dům hypotékou kvůli půjčce}
\dicEntry[veðsetning] \dicTerm{veð··set·ning} \dicIPA{{v}{\textepsilon}{ð}{s}{\textepsilon}{\textsubring{d}}{n}{i}{\ng}{\r{g}}} \dicPos{f}[4] \dicFlx{(‑ar, ‑ar)}[5] \dicDirectTranslationCS{zastavení, zatížení hypotékou}
\dicEntry[veðsett] \dicTerm{veð··sett} \dicIPA{{v}{\textepsilon}{ð}{s}{\textepsilon}{h}{\textsubring{d}}} \dicPos{v} \dicFlx{supin} \dicLink{veðsetja}
\dicEntry[veðsetti] \dicTerm{veð··setti} \dicIPA{{v}{\textepsilon}{ð}{s}{\textepsilon}{h}{\textsubring{d}}{\textsci}} \dicPos{v} \dicFlx{ind\,/\addthin con pf sg 1 pers} \dicLink{veðsetja}
\dicEntry[veðsettum] \dicTerm{veð··settum} \dicIPA{{v}{\textepsilon}{ð}{s}{\textepsilon}{h}{\textsubring{d}}{\textscy}{\textsubring{m}}} \dicPos{v} \dicFlx{ind pf pl 1 pers} \dicLink{veðsetja}
\dicEntry[veðskuld] \dicTerm{veð··skuld} \dicIPA{{v}{\textepsilon}{ð}{s}{\r{g}}{\textscy}{l}{\textsubring{d}}} \dicPos{f}[7] \dicFlx{(‑ar, ‑ir)}[1] \dicFieldCat{ekon.} \dicDirectTranslationCS{hypoteční dluh}
\dicEntry[veðskuldabréf] \dicTerm{veð··skulda··bréf} \dicIPA{{v}{\textepsilon}{ð}{s}{\r{g}}{\textscy}{l}{\textsubring{d}}{a}{\textsubring{b}}{r}{j}{\textepsilon}{f}} \dicPos{n}[2] \dicFlx{(‑s, ‑)}[5] \dicFieldCat{ekon.} \dicDirectTranslationCS{hypoteční dluhopis, hypoteční zástavní list}
\dicEntry[veðtryggja] \dicTerm{veð··trygg|ja} \dicIPA{{v}{\textepsilon}{ð}{t\smash{\textsuperscript{h}}}{r}{\textsci}{\r{\textObardotlessj}}{a}} \dicPos{v}[2] \dicFlx{(‑ði, ‑t)}[92] \dicFlx{acc} \dicFieldCat{ekon.} \dicDirectTranslationCS{pojistit zástavou}
\dicEntry[veður] \dicTerm{veður} \dicsymFrequent\  \dicIPA{{v}{\textepsilon}{\textlengthmark}{ð}{\textscy}{\textsubring{r}}} \dicPos{n}[2] \dicFlx{(‑s, ‑)}[25] \dicDirectTranslationCS{počasí} \dicExampleIS{gott veður} \dicExampleCS{dobré počasí};  \dicIdiom{veður}{ \dicPhraseIS{fá\,/\addthin hafa veður af e‑u}} \dicDirectTranslationCS{doslechnout se (co)};  \dicPhraseIS{gera veður út af e‑u} \dicDirectTranslationCS{dělat kvůli (čemu) mnoho povyku};  \dicPhraseIS{láta e‑ð í veðri vaka} \dicDirectTranslationCS{naznačit (co), naznačovat (co)};  \dicPhraseIS{sækja í sig veðrið} \dicDirectTranslationCS{sbírat síly};  \dicPhraseIS{vita hvaðan á sig stendur veðrið} \dicLangCat{přen.} \dicDirectTranslationCS{vědět, na čem člověk je};  \dicPhraseIS{vita ekki hvaðan á sig stendur veðrið} \dicLangCat{přen.} \dicDirectTranslationCS{nevědět, která bije};  \dicPhraseIS{það er allra veðra von} \dicLangCat{přen.} \dicDirectTranslationCS{člověk se může dočkat všeho možného}
\dicEntry[veðurathugun] \dicTerm{veður··at·hug|un} \dicIPA{{v}{\textepsilon}{\textlengthmark}{ð}{\textscy}{r}{a}{\textsubring{d}}{h}{\textscy}{\textbabygamma}{\textscy}{\textsubring{n}}} \dicPos{f}[7] \dicFlx{(‑unar, ‑anir)}[8] \dicDirectTranslationCS{sledování počasí, meteorologické pozorování}
\dicEntry[veðurathugunarstöð] \dicTerm{veður·at·hugunar··stöð} \dicIPA{{v}\-{\textepsilon}\-{\textlengthmark}\-{ð}\-{\textscy}\-{r}\-{a}\-{\textsubring{d}}\-{h}\-{\textscy}\-{\textbabygamma}\-{\textscy}\-{n}\-{a}\-{\textsubring{r}}\-{s}\-{\textsubring{d}}\-{\oe}\-{\texttheta}\-} \dicPos{f}[6] \dicFlx{(‑var, ‑var)}[1] \dicFieldCat{meteo.} \dicDirectTranslationCS{meteorologická stanice}
\dicEntry[veðurbarinn] \dicTerm{veður··barinn}\dicTerm{, veðurbitinn} \dicIPA{{v}\-{\textepsilon}\-{\textlengthmark}\-{ð}\-{\textscy}\-{r}\-{\textsubring{b}}\-{a}\-{r}\-{\textsci}\-{\textsubring{n}}\-} \dicPos{adj}[6]\dicFlx{}[-14] \textbf{1.} \dicSynonym{veðraður} \dicDirectTranslationCS{poničený\,/\addthin zničený vlivem počasí} \dicExampleIS{Jörð er veðurbarin.} \dicExampleCS{Země je zničená vlivem počasí.}  \textbf{2.} \dicDirectTranslationCS{ošlehaný větrem} \dicIndirectTranslationCS{(o~člověku)}
\dicEntry[veðurbitinn] \dicTerm{veður··bitinn} \dicIPA{{v}{\textepsilon}{\textlengthmark}{ð}{\textscy}{r}{\textsubring{b}}{\textsci}{\textsubring{d}}{\textsci}{\textsubring{n}}} \dicPos{adj}[6]\dicFlx{}[-2] \dicLink{veðurbarinn}
\dicEntry[veðurblíða] \dicTerm{veður··blíð|a} \dicIPA{{v}{\textepsilon}{\textlengthmark}{ð}{\textscy}{r}{\textsubring{b}}{l}{i}{ð}{a}} \dicPos{f}[1] \dicFlx{(‑u)}[5] \dicDirectTranslationCS{vlídné\,/\addthin přívětivé počasí}
\dicEntry[veðurbreyting] \dicTerm{veður··breyt·ing} \dicIPA{{v}{\textepsilon}{ð}{\textscy}{r}{\textsubring{b}}{r}{ei}{\textsubring{d}}{i}{\ng}{\r{g}}} \dicPos{f}[4] \dicFlx{(‑ar, ‑ar)}[5] \textbf{1.} \dicDirectTranslationCS{změna počasí}  \textbf{2.} \dicDirectTranslationCS{klimatická změna, změna klimatu}
\dicEntry[veðurdagur] \dicTerm{veður··dag|ur} \dicIPA{{v}{\textepsilon}{\textlengthmark}{ð}{\textscy}{r}{\textsubring{d}}{a}{\textbabygamma}{\textscy}{\textsubring{r}}} \dicPos{m}[6] \dicFlx{(‑s, ‑ar)}[62] \dicPhraseIS{einn góðan veðurdag} \dicFlx{adv} \dicDirectTranslationCS{jednou, jednoho dne} \dicExampleIS{Einn góðan veðurdag ætla ég að hlaupa svo hratt að \dots{}} \dicExampleCS{Jednou bych chtěla běžet tak rychle, až\dots{}}
\dicEntry[veðurfar] \dicTerm{veður··far} \dicIPA{{v}{\textepsilon}{\textlengthmark}{ð}{\textscy}{\textsubring{r}}{f}{a}{\textsubring{r}}} \dicPos{n}[2] \dicFlx{(‑s)}[2] \dicSynonym{loftslag} \dicDirectTranslationCS{podnebí, klima}
\dicEntry[veðurfarsbreyting] \dicTerm{veður·fars··breyt·ing} \dicIPA{{v}\-{\textepsilon}\-{\textlengthmark}\-{ð}\-{\textscy}\-{\textsubring{r}}\-{f}\-{a}\-{\textsubring{r}}\-{s}\-{\textsubring{b}}\-{r}\-{ei}\-{\textsubring{d}}\-{i}\-{\ng}\-{\r{g}}\-} \dicPos{f}[4] \dicFlx{(‑ar, ‑ar)}[5] \dicDirectTranslationCS{klimatická změna} \dicIndirectTranslationCS{(zvláště v~pl)}
\dicEntry[veðurfarsfræði] \dicTerm{veður·fars··fræð|i} \dicIPA{{v}{\textepsilon}{\textlengthmark}{ð}{\textscy}{\textsubring{r}}{f}{a}{\textsubring{r}}{s}{f}{r}{a}{i}{ð}{\textsci}} \dicPos{f}[3] \dicFlx{(‑i)}[3] \dicSynonym{loftslagsfræði} \dicDirectTranslationCS{klimatologie}
\dicEntry[veðurfrétt] \dicTerm{veður··frétt} \dicIPA{{v}{\textepsilon}{\textlengthmark}{ð}{\textscy}{\textsubring{r}}{f}{r}{j}{\textepsilon}{h}{\textsubring{d}}} \dicPos{f}[7] \dicFlx{(‑ar, ‑ir)}[1] \dicDirectTranslationCS{zpráva o~počasí} \dicIndirectTranslationCS{(zvláště v~pl)}
\dicEntry[veðurfræði] \dicTerm{veður··fræð|i} \dicIPA{{v}{\textepsilon}{\textlengthmark}{ð}{\textscy}{\textsubring{r}}{f}{r}{a}{i}{ð}{\textsci}} \dicPos{f}[3] \dicFlx{(‑i)}[3] \dicDirectTranslationCS{meteorologie}
\dicEntry[veðurfræðingur] \dicTerm{veður·fræð··ing|ur} \dicIPA{{v}{\textepsilon}{\textlengthmark}{ð}{\textscy}{\textsubring{r}}{f}{r}{a}{i}{ð}{i}{\ng}{\r{g}}{\textscy}{\textsubring{r}}} \dicPos{m}[6] \dicFlx{(‑s, ‑ar)}[8] \dicDirectTranslationCS{meteorolog, meteoroložka}
\dicEntry[veðurglöggur] \dicTerm{veður··|glöggur} \dicIPA{{v}{\textepsilon}{\textlengthmark}{ð}{\textscy}{r}{\r{g}}{l}{\oe}{\r{g}}{\textscy}{\textsubring{r}}} \dicPos{adj}[10] \dicFlx{(comp ‑gleggri, sup ‑gleggstur)}[1] \dicDirectTranslationCS{znalý počasí}
\dicEntry[veðurhorfur] \dicTerm{veður··horfur} \dicIPA{{v}{\textepsilon}{\textlengthmark}{ð}{\textscy}{\textsubring{r}}{h}{\textopeno}{r}{v}{\textscy}{\textsubring{r}}} \dicPos{f}[1] \dicFlx{pl}[12] \dicDirectTranslationCS{prognóza počasí}
\dicEntry[veðurhæð] \dicTerm{veður··hæð} \dicIPA{{v}{\textepsilon}{\textlengthmark}{ð}{\textscy}{\textsubring{r}}{h}{a}{i}{\texttheta}} \dicPos{f}[7] \dicFlx{(‑ar)}[3] \dicFieldCat{meteo.} \dicSynonym{vindhraði} \dicDirectTranslationCS{rychlost\,/\addthin síla větru}
\dicEntry[veðurkort] \dicTerm{veður··kort} \dicIPA{{v}{\textepsilon}{\textlengthmark}{ð}{\textscy}{\textsubring{r}}{k\smash{\textsuperscript{h}}}{\textopeno}{\textsubring{r}}{\textsubring{d}}} \dicPos{n}[2] \dicFlx{(‑s, ‑)}[5] \dicDirectTranslationCS{povětrnostní mapa}
\dicEntry[veðurlag] \textls[10]{\dicTerm{veður··lag} \dicIPA{{v}{\textepsilon}{\textlengthmark}{ð}{\textscy}{r}{l}{a}{x}} \dicPos{n}[2] \dicFlx{(‑s)}[2] \dicSynonym{veðrátta} \dicDirectTranslationCS{podnebí, klima, počasí}}
\dicEntry[veðurlýsing] \textls[10]{\dicTerm{veður··lýs·ing} \dicIPA{{v}{\textepsilon}{\textlengthmark}{ð}{\textscy}{r}{l}{i}{s}{i}{\ng}{\r{g}}} \dicPos{f}[4] \dicFlx{(‑ar, ‑ar)}[5] \dicDirectTranslationCS{informace\,/\addthin zpráva o~počasí}}
\dicEntry[veðurspá] \dicTerm{veður··spá} \dicIPA{{v}{\textepsilon}{\textlengthmark}{ð}{\textscy}{\textsubring{r}}{s}{\textsubring{b}}{au}} \dicPos{f}[4] \dicFlx{(‑r, ‑r)}[18] \dicDirectTranslationCS{předpověď počasí}
\dicEntry[veðurstofa] \dicTerm{veður··stof|a} \dicIPA{{v}{\textepsilon}{\textlengthmark}{ð}{\textscy}{\textsubring{r}}{s}{\textsubring{d}}{\textopeno}{v}{a}} \dicPos{f}[1] \dicFlx{(‑u, ‑ur)}[7] \dicDirectTranslationCS{meteorologický institut}
\dicEntry[veðurstöð] \dicTerm{veður··stöð} \dicIPA{{v}{\textepsilon}{\textlengthmark}{ð}{\textscy}{\textsubring{r}}{s}{\textsubring{d}}{\oe}{\texttheta}} \dicPos{f}[6] \dicFlx{(‑var, ‑var)}[1] \dicDirectTranslationCS{povětrnostní\,/\addthin meteorologická stanice}
\dicEntry[veðurtepptur] \dicTerm{veður··tepptur} \dicIPA{{v}{\textepsilon}{\textlengthmark}{ð}{\textscy}{\textsubring{r}}{t\smash{\textsuperscript{h}}}{\textepsilon}{f}{\textsubring{d}}{\textscy}{\textsubring{r}}} \dicPos{adj}[1]\dicFlx{}[-13] \dicDirectTranslationCS{zdržený\,/\addthin uvězněný kvůli špatnému počasí} \dicExampleIS{sitja þar veðurtepptur í tvo daga} \dicExampleCS{být tam uvězněný kvůli špatnému počasí dva dny}
\dicEntry[veðurtungl] \dicTerm{veður··tungl} \dicIPA{{v}{\textepsilon}{\textlengthmark}{ð}{\textscy}{\textsubring{r}}{t\smash{\textsuperscript{h}}}{u}{\ng}{\textsubring{l}}} \dicPos{n}[2] \dicFlx{(‑s, ‑)}[5] \dicDirectTranslationCS{meteorologická družice}
\dicEntry[veðurtunglamynd] \dicTerm{veður·tungla··mynd} \dicIPA{{v}{\textepsilon}{\textlengthmark}{ð}{\textscy}{\textsubring{r}}{t\smash{\textsuperscript{h}}}{u}{\ng}{l}{a}{m}{\textsci}{n}{\textsubring{d}}} \dicPos{f}[7] \dicFlx{(‑ar, ‑ir)}[1] \dicDirectTranslationCS{satelitní\,/\addthin družicový snímek}
\dicEntry[vef] \dicTerm{vef} \dicIPA{{v}{\textepsilon}{\textlengthmark}{f}} \dicPos{v} \dicFlx{ind praes sg 1 pers} \textbf{1.} \dicLink{vefa}  \textbf{2.} \dicLink{vefja}
\dicEntry[vefa] \dicTerm{vefa} \dicIPA{{v}{\textepsilon}{\textlengthmark}{v}{a}} \dicPos{v}[6] \dicFlx{(vef, óf, ófum, væfi, ofið)}[123] \dicFlx{acc} \textbf{1.} \dicDirectTranslationCS{utkat, tkát} \dicExampleIS{vefa teppi} \dicExampleCS{utkat koberec}  \textbf{2.} \dicFieldCat{poč.} \dicDirectTranslationCS{tvořit\,/\addthin vytvářet webové stránky}
\dicEntry[vefari] \dicTerm{vef··ar|i} \dicIPA{{v}{\textepsilon}{\textlengthmark}{v}{a}{r}{\textsci}} \dicPos{m}[1] \dicFlx{(‑a, ‑ar)}[13] \textbf{1.} \dicDirectTranslationCS{tkadlec, tkadlena}  \textbf{2.} \dicFieldCat{poč.} \dicDirectTranslationCS{správce\,/\addthin správkyně webových stránek, webmaster(ka)}
\dicEntry[vefði] \dicTerm{vefði} \dicIPA{{v}{\textepsilon}{v}{ð}{\textsci}} \dicPos{v} \dicFlx{con pf sg 1 pers} \dicLink{vefja}
\dicEntry[veffang] \dicTerm{vef··|fang} \dicIPA{{v}{\textepsilon}{f}{\textlengthmark}{au}{\ng}{\r{g}}} \dicPos{n}[2] \dicFlx{(‑fangs, ‑föng)}[8] \dicFieldCat{poč.} \dicSynonym{slóð} \dicDirectTranslationCS{internetová\,/\addthin webová adresa}
\dicEntry[vefgátt] \dicTerm{vef··|gátt} \dicIPA{{v}{\textepsilon}{v}{\r{g}}{au}{h}{\textsubring{d}}} \dicPos{f}[7] \dicFlx{(‑gáttar, ‑gáttir\,/\addthin ‑gættir)}[25] \dicFieldCat{poč.} \dicDirectTranslationCS{webový portál}
\dicEntry[vefhönnuður] \dicTerm{vef··hönn·uð|ur} \dicIPA{{v}{\textepsilon}{\textlengthmark}{f}{h}{\oe}{n}{\textscy}{ð}{\textscy}{\textsubring{r}}} \dicPos{m}[10] \dicFlx{(‑ar, ‑ir)}[4] \dicDirectTranslationCS{návrhář(ka) webu, webový designér, webová designérka}
\dicEntry[vefhönnun] \dicTerm{vef··hönn|un} \dicIPA{{v}{\textepsilon}{\textlengthmark}{f}{h}{\oe}{n}{\textscy}{\textsubring{n}}} \dicPos{f}[7] \dicFlx{(‑unar)}[12] \dicDirectTranslationCS{webdesign, webový design}
\dicEntry[vefja] \dicTerm{vefja} \dicsymFrequent\  \dicIPA{{v}{\textepsilon}{v}{j}{a}} \dicPos{v}[4] \dicFlx{(vef, vafði, vöfðum, vefði, vafið)}[10] \dicFlx{acc\,/\addthin dat} \dicDirectTranslationCS{ovinout, ovíjet, omotat, omotávat, obtočit} \dicExampleIS{vefja teppi utan um hana} \dicExampleCS{ovinout ji přikrývkou};  \dicPhraseIS{vefja e‑n örmum} \dicDirectTranslationCS{obejmout (koho)};  \dicIdiom{vefja}[inn]{ \dicPhraseIS{vefja e‑ð inn í e‑ð}} \dicDirectTranslationCS{zabalit (co) do (čeho), ovinout (co čím)} \dicExampleIS{vefja gjöfina inn í pappír} \dicExampleCS{zabalit dárek do papíru};  \dicIdiom{vefja}[um]{ \dicPhraseIS{vefja e‑m um fingur sér}} \dicLangCat{přen.} \dicDirectTranslationCS{omotat si (koho) kolem prstu};  \dicIdiom{vefja}[upp]{ \dicPhraseIS{vefja e‑ð upp}} \dicDirectTranslationCS{ovinout (co), omotat (co)};  \dicIdiom{vefjast}{ \dicPhraseIS{e‑að vefst fyrir e‑m}} \dicFlx{refl} \dicDirectTranslationCS{(co) se (komu) plete}; { \dicPhraseIS{það vefst fyrir e‑m að (gera e‑ð)}} \dicFlx{refl impers} \dicDirectTranslationCS{(kdo) má potíže (udělat (co))}
\dicEntry[vefjafræði] \dicTerm{vefja··fræð|i} \dicIPA{{v}{\textepsilon}{v}{j}{a}{f}{r}{a}{i}{ð}{\textsci}} \dicPos{f}[3] \dicFlx{(‑i)}[3] \dicFieldCat{biol.} \dicDirectTranslationCS{histologie}
\dicEntry[vefjarhöttur] \dicTerm{vefjar··|höttur} \dicIPA{{v}{\textepsilon}{v}{j}{a}{\textsubring{r}}{h}{\oe}{h}{\textsubring{d}}{\textscy}{\textsubring{r}}} \dicPos{m}[11] \dicFlx{(‑hattar, ‑hettir)}[5] \dicDirectTranslationCS{turban}
\dicEntry[vefjarsýni] \dicTerm{vefjar··sýni} \dicIPA{{v}{\textepsilon}{v}{j}{a}{\textsubring{r}}{s}{i}{n}{\textsci}} \dicPos{n}[2] \dicFlx{(‑s, ‑)}[14] \dicDirectTranslationCS{vzorek tkáně, tkáňový vzorek}
\dicEntry[vefjaskemmd] \dicTerm{vefja··skemmd} \dicIPA{{v}{\textepsilon}{v}{j}{a}{s}{\r{\textObardotlessj}}{\textepsilon}{m}{\textsubring{d}}} \dicPos{f}[7] \dicFlx{(‑ar, ‑ir)}[1] \dicDirectTranslationCS{poškození tkáně}
\dicEntry[vefmyndavél] \dicTerm{vef··mynda·vél} \dicIPA{{v}{\textepsilon}{v}{m}{\textsci}{n}{\textsubring{d}}{a}{v}{j}{\textepsilon}{\textsubring{l}}} \dicPos{f}[4] \dicFlx{(‑ar, ‑ar)}[1] \dicFieldCat{poč.} \dicDirectTranslationCS{webkamera, webová kamera}
\dicEntry[vefnaðariðnaður] \dicTerm{vefnaðar··iðn·að|ur} \dicIPA{{v}{\textepsilon}{\textsubring{b}}{n}{a}{ð}{a}{r}{\textsci}{ð}{n}{a}{ð}{\textscy}{\textsubring{r}}} \dicPos{m}[10] \dicFlx{(‑ar)}[9] \dicDirectTranslationCS{textilní průmysl}
\dicEntry[vefnaðarvara] \dicTerm{vefnaðar··|vara} \dicIPA{{v}{\textepsilon}{\textsubring{b}}{n}{a}{ð}{a}{r}{v}{a}{r}{a}} \dicPos{f}[1] \dicFlx{(‑vöru, ‑vörur)}[8] \dicDirectTranslationCS{textilní výrobek, tkanina}
\dicEntry[vefnaðarvörubúð] \dicTerm{vefnaðar·vöru··búð} \dicIPA{{v}{\textepsilon}{\textsubring{b}}{n}{a}{ð}{a}{r}{v}{\oe}{r}{\textscy}{\textsubring{b}}{u}{\texttheta}} \dicPos{f}[7] \dicFlx{(‑ar, ‑ir)}[1] \dicDirectTranslationCS{obchod s~textilním zbožím, galanterie}
\dicEntry[vefnaður] \dicTerm{vef··nað|ur} \dicIPA{{v}{\textepsilon}{\textsubring{b}}{n}{a}{ð}{\textscy}{\textsubring{r}}} \dicPos{m}[10] \dicFlx{(‑ar)}[7] \textbf{1.} \dicSynonym*{það að vefa} \dicDirectTranslationCS{tkaní}  \textbf{2.} \dicSynonym{voð} \dicDirectTranslationCS{textil, textilie, tkaný výrobek}
\dicEntry[vefpóstur] \dicTerm{vef··póst|ur} \dicIPA{{v}{\textepsilon}{v}{p\smash{\textsuperscript{h}}}{ou}{s}{\textsubring{d}}{\textscy}{\textsubring{r}}} \dicPos{m}[6] \dicFlx{(‑s)}[7] \dicFieldCat{poč.} \dicDirectTranslationCS{webová pošta, webmail}
\dicEntry[vefrænn] \dicTerm{vef··rænn} \dicIPA{{v}{\textepsilon}{v}{r}{a}{i}{\textsubring{d}}{\textsubring{n}}} \dicPos{adj}[7]\dicFlx{}[-2] \dicFieldCat{med. } \dicSynonym{líkamlegur} \dicDirectTranslationCS{somatický, tělesný} \dicExampleIS{vefrænn sjúkdómur} \dicExampleCS{somatické onemocnění}
\dicEntry[vefsetur] \dicTerm{vef··setur} \dicIPA{{v}{\textepsilon}{f}{s}{\textepsilon}{\textsubring{d}}{\textscy}{\textsubring{r}}} \dicPos{n}[2] \dicFlx{(‑s, ‑)}[25] \dicFieldCat{poč.} \dicDirectTranslationCS{webová\,/\addthin internetová stránka}
\dicEntry[vefsíða] \dicTerm{vef··síð|a} \dicIPA{{v}{\textepsilon}{f}{s}{i}{ð}{a}} \dicPos{f}[1] \dicFlx{(‑u, ‑ur)}[13] \dicFieldCat{poč.} \dicDirectTranslationCS{webová stránka}
\dicEntry[vefskoðari] \dicTerm{vef··skoð·ar|i} \dicIPA{{v}{\textepsilon}{f}{s}{\r{g}}{\textopeno}{ð}{a}{r}{\textsci}} \dicPos{m}[1] \dicFlx{(‑a, ‑ar)}[13] \dicFieldCat{poč.} \dicSynonym{vafri} \dicDirectTranslationCS{(internetový) prohlížeč}
\dicEntry[vefslóð] \dicTerm{vef··slóð} \dicIPA{{v}{\textepsilon}{f}{s}{\textsubring{d}}{l}{ou}{\texttheta}} \dicPos{f}[7] \dicFlx{(‑ar, ‑ir)}[1] \dicFieldCat{poč.} \dicDirectTranslationCS{webová\,/\addthin internetová adresa}
\dicEntry[vefstjóri] \dicTerm{vef··stjór|i} \dicIPA{{v}{\textepsilon}{f}{s}{\textsubring{d}}{j}{ou}{r}{\textsci}} \dicPos{m}[1] \dicFlx{(‑a, ‑ar)}[1] \dicFieldCat{poč.} \dicDirectTranslationCS{webmaster(ka), webový správce, webová správkyně}
\dicEntry[vefstóll] \dicTerm{vef··stól|l} \dicIPA{{v}{\textepsilon}{f}{s}{\textsubring{d}}{ou}{\textsubring{d}}{\textsubring{l}}} \dicPos{m}[6] \dicFlx{(‑s, ‑ar)}[48] \dicDirectTranslationCS{(tkalcovský) stav}
\dicEntry[vefur] \dicTerm{vef|ur} \dicIPA{{v}{\textepsilon}{\textlengthmark}{v}{\textscy}{\textsubring{r}}} \dicPos{m}[9] \dicFlx{(‑jar\,/\addthin ‑s, ‑ir)}[26] \textbf{1.} \dicSynonym*{ofið efni} \dicDirectTranslationCS{tkanivo, pletivo}  \textbf{2.} \dicSynonym{kóngulóarvefur} \dicDirectTranslationCS{pavučina, (pavoučí) síť}  \textbf{3.} \dicFieldCat{biol.} \dicDirectTranslationCS{tkáň}  \textbf{4.} \dicSynonym{vefsetur} \dicDirectTranslationCS{web, webová stránka}
\dicEntry[vefverslun] \dicTerm{vef··versl|un} \dicIPA{{v}{\textepsilon}{v}{\textlengthmark}{\textepsilon}{\textsubring{r}}{s}{\textsubring{d}}{l}{\textscy}{\textsubring{n}}} \dicPos{f}[7] \dicFlx{(‑unar, ‑anir)}[8] \dicDirectTranslationCS{internetový\,/\addthin elektronický obchod, e‑shop}
\dicEntry[vefþjónn] \dicTerm{vef··þjón|n} \dicIPA{{v}{\textepsilon}{f}{\texttheta}{j}{ou}{\textsubring{d}}{\textsubring{n}}} \dicPos{m}[6] \dicFlx{(‑s, ‑ar)}[42] \dicFieldCat{poč.} \dicDirectTranslationCS{webový server}
\dicEntry[veg] \dicTerm{veg} \dicIPA{{v}{\textepsilon}{\textlengthmark}{x}} \dicPos{v} \dicFlx{ind praes sg 1 pers} \textbf{1.} \dicLink{vega\smash{\textsuperscript{1}}}  \textbf{2.} \dicLink{vega\smash{\textsuperscript{2}}}
\dicEntry[vega] \dicTerm{vega\smash{\textsuperscript{1}}} \dicsymFrequent\  \dicIPA{{v}{\textepsilon}{\textlengthmark}{\textbabygamma}{a}} \dicPos{v}[6] \dicFlx{(veg, vó, vógum, vægi, vegið)}[108] \dicFlx{acc} \textbf{1.} \dicSynonym{vigta} \dicDirectTranslationCS{(z)vážit} \dicIndirectTranslationCS{(zjistit hmotnost)} \dicExampleIS{vega hlutinn} \dicExampleCS{vážit věc}  \textbf{2.} \dicSynonym*{vera að þyngd} \dicDirectTranslationCS{vážit} \dicIndirectTranslationCS{(mít nějakou hmotnost)} \dicExampleIS{Pokinn vegur 20 kíló.} \dicExampleCS{Pytel váží dvacet kilo.}  \textbf{3.} \dicPhraseIS{e‑að vegur þungt} \dicLangCat{přen.} \dicDirectTranslationCS{(co) má váhu (slova papeže ap.)};  \dicIdiom{vega}[sig upp]{ \dicPhraseIS{vega sig upp á e‑ð}} \dicDirectTranslationCS{vytáhnout se na (co) (zeď ap.)};  \dicIdiom{vega}[upp]{ \dicPhraseIS{vega upp á móti e‑u}} \dicDirectTranslationCS{vyvážit (co),  vykompenzovat (co)} \dicExampleIS{Afraksturinn vegur upp á móti tilkostnaðinum.} \dicExampleCS{Zisk vyváží výdaje.}; { \dicPhraseIS{vega upp e‑ð}} \dicDirectTranslationCS{nahradit (co), kompenzovat (co)}
\dicEntry[vega] \dicTerm{vega\smash{\textsuperscript{2}}} \dicIPA{{v}{\textepsilon}{\textlengthmark}{\textbabygamma}{a}} \dicPos{v}[6] \dicFlx{(veg, vó, vógum, vægi, vegið)}[108] \dicFlx{acc} \dicSynonym{drepa} \dicDirectTranslationCS{zabít, zabíjet} \dicExampleIS{vega mann} \dicExampleCS{zabít člověka};  \dicIdiom{vega}[að]{ \dicPhraseIS{vega að e‑m}} \dicDirectTranslationCS{(za)útočit na (koho) (slovně ap.)};  \dicIdiom{vegast}{ \dicPhraseIS{vegast}} \dicFlx{refl} \dicSynonym*{berjast} \dicDirectTranslationCS{bojovat, válčit}
\dicEntry[vegabréf] \dicTerm{vega··bréf} \dicIPA{{v}{\textepsilon}{\textlengthmark}{\textbabygamma}{a}{\textsubring{b}}{r}{j}{\textepsilon}{f}} \dicPos{n}[2] \dicFlx{(‑s, ‑)}[5] \dicSynonym{passi} \dicDirectTranslationCS{(cestovní) pas} \dicExampleIS{gilt vegabréf} \dicExampleCS{platný cestovní pas}
\dicEntry[vegabréfaeftirlit] \dicTerm{vega·bréfa··eftir·lit} \dicIPA{{v}\-{\textepsilon}\-{\textlengthmark}\-{\textbabygamma}\-{a}\-{\textsubring{b}}\-{r}\-{j}\-{\textepsilon}\-{v}\-{a}\-{\textepsilon}\-{f}\-{\textsubring{d}}\-{\textsci}\-{r}\-{l}\-{\textsci}\-{\textsubring{d}}\-} \dicPos{n}[2] \dicFlx{(‑s)}[2] \dicDirectTranslationCS{pasová kontrola}
\dicEntry[vegabréfaskoðun] \dicTerm{vega·bréfa··skoð|un} \dicIPA{{v}{\textepsilon}{\textlengthmark}{\textbabygamma}{a}{\textsubring{b}}{r}{j}{\textepsilon}{v}{a}{s}{\r{g}}{\textopeno}{ð}{\textscy}{\textsubring{n}}} \dicPos{f}[7] \dicFlx{(‑unar)}[9] \dicDirectTranslationCS{pasová kontrola}
\dicEntry[vegabréfsáritun] \dicTerm{vega·bréfs··á·rit|un} \dicIPA{{v}{\textepsilon}{\textlengthmark}{\textbabygamma}{a}{\textsubring{b}}{r}{j}{\textepsilon}{f}{s}{au}{r}{\textsci}{\textsubring{d}}{\textscy}{\textsubring{n}}} \dicPos{f}[7] \dicFlx{(‑unar, ‑anir)}[8] \dicDirectTranslationCS{vízum}
\dicEntry[vegagerð] \dicTerm{vega··gerð} \dicIPA{{v}{\textepsilon}{\textlengthmark}{\textbabygamma}{a}{\r{\textObardotlessj}}{\textepsilon}{r}{\texttheta}} \dicPos{f}[7] \dicFlx{(‑ar)}[3] \dicDirectTranslationCS{stavba silnice}
\dicEntry[vegahandbók] \dicTerm{vega··hand·|bók} \dicIPA{{v}{\textepsilon}{\textlengthmark}{\textbabygamma}{a}{h}{a}{n}{\textsubring{d}}{\textsubring{b}}{ou}{\r{g}}} \dicPos{f}[8] \dicFlx{(‑bókar, ‑bækur)}[5] \dicDirectTranslationCS{motoristický\,/\addthin silniční průvodce} \dicIndirectTranslationCS{(kniha)}
\dicEntry[vegakerfi] \dicTerm{vega··kerfi} \dicIPA{{v}{\textepsilon}{\textlengthmark}{\textbabygamma}{a}{c\smash{\textsuperscript{h}}}{\textepsilon}{r}{v}{\textsci}} \dicPos{n}[2] \dicFlx{(‑s, ‑)}[14] \dicDirectTranslationCS{silniční systém}
\dicEntry[vegakort] \dicTerm{vega··kort} \dicIPA{{v}{\textepsilon}{\textlengthmark}{\textbabygamma}{a}{k\smash{\textsuperscript{h}}}{\textopeno}{\textsubring{r}}{\textsubring{d}}} \dicPos{n}[2] \dicFlx{(‑s, ‑)}[5] \dicDirectTranslationCS{automapa, silniční mapa}
\dicEntry[vegalaus] \dicTerm{vega··laus} \dicIPA{{v}{\textepsilon}{\textlengthmark}{\textbabygamma}{a}{l}{\oe i}{s}} \dicPos{adj}[5]\dicFlx{}[-1] \textbf{1.} \dicSynonym*{án vegar} \dicDirectTranslationCS{(jsoucí) bez silnic\,/\addthin komunikace} \dicExampleIS{vegalaus sveit} \dicExampleCS{venkov bez silnic}  \textbf{2.} \dicSynonym*{athvarfslaus} \dicDirectTranslationCS{bezprizorní (dítě ap.)}
\dicEntry[vegalengd] \dicTerm{vega··lengd} \dicIPA{{v}{\textepsilon}{\textlengthmark}{\textbabygamma}{a}{l}{\textepsilon}{\ng}{\textsubring{d}}} \dicPos{f}[7] \dicFlx{(‑ar, ‑ir)}[1] \dicDirectTranslationCS{vzdálenost} \dicIndirectTranslationCS{(mezi dvěma místy)}
\dicEntry[vegalögregla] \dicTerm{vega··lög·regl|a} \dicIPA{{v}{\textepsilon}{\textlengthmark}{\textbabygamma}{a}{l}{\oe}{\textbabygamma}{r}{\textepsilon}{\r{g}}{l}{a}} \dicPos{f}[1] \dicFlx{(‑u)}[5] \dicDirectTranslationCS{dopravní policie}
\dicEntry[vegamál] \dicTerm{vega··mál} \dicIPA{{v}{\textepsilon}{\textlengthmark}{\textbabygamma}{a}{m}{au}{\textsubring{l}}} \dicPos{n}[2] \dicFlx{pl}[1] \dicDirectTranslationCS{silniční záležitosti, otázky silnic}
\dicEntry[vegamerki] \dicTerm{vega··merki} \dicIPA{{v}{\textepsilon}{\textlengthmark}{\textbabygamma}{a}{m}{\textepsilon}{\textsubring{r}}{\r{\textObardotlessj}}{\textsci}} \dicPos{n}[2] \dicFlx{(‑s, ‑)}[16] \dicSynonym{umferðarmerki} \dicDirectTranslationCS{dopravní značka}
\dicEntry[vegamót] \textls[10]{\dicTerm{vega··mót} \dicIPA{{v}{\textepsilon}{\textlengthmark}{\textbabygamma}{a}{m}{ou}{\textsubring{d}}} \dicPos{n}[2] \dicFlx{pl}[1] \dicDirectTranslationCS{křižovatka, rozcestí}} \dicExampleIS{koma að vegamótum} \dicExampleCS{přijet na křižovatku}
\dicEntry[veganesti] \dicTerm{vega··nesti} \dicIPA{{v}{\textepsilon}{\textlengthmark}{\textbabygamma}{a}{n}{\textepsilon}{s}{\textsubring{d}}{\textsci}} \dicPos{n}[2] \dicFlx{(‑s)}[20] \dicDirectTranslationCS{vínek, výbava na cestu (životem ap.)} \dicExampleIS{hafa góðar gáfur í veganesti} \dicExampleCS{dostat dobré vlohy do vínku}
\dicEntry[vegarfi] \dicTerm{veg··arf|i} \dicIPA{{v}{\textepsilon}{\textlengthmark}{\textbabygamma}{a}{r}{v}{\textsci}} \dicPos{m}[1] \dicFlx{(‑a)}[3] \dicFieldCat{bot.} \dicDirectTranslationCS{rožec prameništní} \textit{(l.~{\textLA{Cerastium fontanum}})}  \dicsymPhoto\ 
\dicFigure{ds_image_vegarfi_0_1.jpg}{Vegarfi}{Vegarfi - Me, PD}
\dicEntry[vegarhelmingur] \dicTerm{vegar··helm·ing|ur} \dicIPA{{v}{\textepsilon}{\textlengthmark}{\textbabygamma}{a}{\textsubring{r}}{h}{\textepsilon}{l}{m}{i}{\ng}{\r{g}}{\textscy}{\textsubring{r}}} \dicPos{m}[6] \dicFlx{(‑s, ‑ar)}[8] \dicDirectTranslationCS{(jízdní) pruh}
\dicEntry[vegarkantur] \dicTerm{vegar··kant|ur}\dicTerm{, vegkantur} \dicIPA{{v}\-{\textepsilon}\-{\textbabygamma}\-{a}\-{\textsubring{r}}\-{k\smash{\textsuperscript{h}}}\-{a}\-{\textsubring{n}}\-{\textsubring{d}}\-{\textscy}\-{\textsubring{r}}\-} \dicPos{m}[6] \dicFlx{(‑s, ‑ar)}[5] \dicSynonym{vegbrún} \dicDirectTranslationCS{krajnice, okraj silnice\,/\addthin cesty}
\dicEntry[vegasalt] \dicTerm{vega··|salt} \dicIPA{{v}{\textepsilon}{\textlengthmark}{\textbabygamma}{a}{s}{a}{\textsubring{l}}{\textsubring{d}}} \dicPos{n}[2] \dicFlx{(‑salts, ‑sölt)}[8] \dicDirectTranslationCS{(podepřená) houpačka}
\dicEntry[vegaskilti] \dicTerm{vega··skilti} \dicIPA{{v}{\textepsilon}{\textlengthmark}{\textbabygamma}{a}{s}{\r{\textObardotlessj}}{\textsci}{\textsubring{l}}{\textsubring{d}}{\textsci}} \dicPos{n}[2] \dicFlx{(‑s, ‑)}[14] \dicDirectTranslationCS{(dopravní) ukazatel, směrovka}
\dicEntry[vegavilltur] \dicTerm{vega··villtur}\dicTerm{, vegvilltur} \dicIPA{{v}\-{\textepsilon}\-{\textlengthmark}\-{\textbabygamma}\-{a}\-{v}\-{\textsci}\-{\textsubring{l}}\-{\textsubring{d}}\-{\textscy}\-{\textsubring{r}}\-} \dicPos{adj}[1]\dicFlx{}[-10] \dicSynonym{villtur} \dicDirectTranslationCS{ztracený} \dicIndirectTranslationCS{(ten, který nezná další směr cesty)}
\dicEntry[vegavinna] \dicTerm{vega··vinn|a} \dicIPA{{v}{\textepsilon}{\textlengthmark}{\textbabygamma}{a}{v}{\textsci}{n}{a}} \dicPos{f}[1] \dicFlx{(‑u)}[5] \dicDirectTranslationCS{práce na silnici, silniční práce}
\dicEntry[vegbrún] \dicTerm{veg··brún} \dicIPA{{v}{\textepsilon}{\textbabygamma}{\textsubring{b}}{r}{u}{\textsubring{n}}} \dicPos{f}[7] \dicFlx{(‑ar, ‑ir)}[1] \dicSynonym{vegarkantur} \dicDirectTranslationCS{krajnice, okraj silnice\,/\addthin cesty}
\dicEntry[vegfarandi] \dicTerm{veg··far·|andi} \dicIPA{{v}{\textepsilon}{x}{f}{a}{r}{a}{n}{\textsubring{d}}{\textsci}} \dicPos{m}[2] \dicFlx{(‑anda, ‑endur)}[1] \dicSynonym{ferðamaður} \dicDirectTranslationCS{cestující, pocestný}
\dicEntry[vegferð] \dicTerm{veg··ferð} \dicIPA{{v}{\textepsilon}{x}{f}{\textepsilon}{r}{\texttheta}} \dicPos{f}[7] \dicFlx{(‑ar)}[3] \dicSynonym{för\smash{\textsuperscript{1}}} \dicDirectTranslationCS{cesta, putování, pouť (životem ap.)}
\dicEntry[veggfóðra] \dicTerm{vegg··fóðr|a} \dicIPA{{v}{\textepsilon}{\r{g}}{f}{ou}{ð}{r}{a}} \dicPos{v}[1] \dicFlx{(‑aði)}[1] \dicFlx{acc} \dicDirectTranslationCS{(vy)tapetovat}
\dicEntry[veggfóðrari] \dicTerm{vegg··fóðr·ar|i} \dicIPA{{v}{\textepsilon}{\r{g}}{f}{ou}{ð}{r}{a}{r}{\textsci}} \dicPos{m}[1] \dicFlx{(‑a, ‑ar)}[13] \dicDirectTranslationCS{tapetář(ka)}
\dicEntry[veggfóður] \dicTerm{vegg··fóður} \dicIPA{{v}{\textepsilon}{\r{g}}{f}{ou}{ð}{\textscy}{\textsubring{r}}} \dicPos{n}[2] \dicFlx{(‑s)}[28] \dicDirectTranslationCS{tapeta} \dicIndirectTranslationCS{(papírový ap. potah stěn)}
\dicEntry[veggjakrot] \dicTerm{veggja··krot} \dicIPA{{v}{\textepsilon}{\r{\textObardotlessj}}{\textlengthmark}{a}{k\smash{\textsuperscript{h}}}{r}{\textopeno}{\textsubring{d}}} \dicPos{n}[2] \dicFlx{(‑s)}[2] \dicDirectTranslationCS{graffiti}
\dicEntry[veggjalús] \dicTerm{veggja··|lús} \dicIPA{{v}{\textepsilon}{\r{\textObardotlessj}}{\textlengthmark}{a}{l}{u}{s}} \dicPos{f}[9] \dicFlx{(‑lúsar, ‑lýs)}[6] \dicFieldCat{zool.} \dicDirectTranslationCS{štěnice domácí} \textit{(l.~{\textLA{Cimex lectularius}})}
\dicEntry[veggklukka] \dicTerm{vegg··klukk|a} \dicIPA{{v}{\textepsilon}{\r{g}}{k\smash{\textsuperscript{h}}}{l}{\textscy}{h}{\r{g}}{a}} \dicPos{f}[1] \dicFlx{(‑u, ‑ur)}[13] \dicDirectTranslationCS{nástěnné hodiny}
\dicEntry[veggklæðning] \dicTerm{vegg··klæð·ning} \dicIPA{{v}{\textepsilon}{\r{g}}{k\smash{\textsuperscript{h}}}{l}{a}{i}{ð}{n}{i}{\ng}{\r{g}}} \dicPos{f}[4] \dicFlx{(‑ar)}[7] \dicDirectTranslationCS{obklad\,/\addthin obložení stěn}
\dicEntry[vegglampi] \dicTerm{vegg··lamp|i} \dicIPA{{v}{\textepsilon}{\r{g}}{l}{a}{\textsubring{m}}{\textsubring{b}}{\textsci}} \dicPos{m}[1] \dicFlx{(‑a, ‑ar)}[8] \dicDirectTranslationCS{nástěnná lampa}
\dicEntry[veggmynd] \dicTerm{vegg··mynd} \dicIPA{{v}{\textepsilon}{\r{g}}{m}{\textsci}{n}{\textsubring{d}}} \dicPos{f}[7] \dicFlx{(‑ar, ‑ir)}[1] \dicDirectTranslationCS{nástěnný obraz}
\dicEntry[veggskreyting] \dicTerm{vegg··skreyt·ing} \dicIPA{{v}{\textepsilon}{\r{g}}{s}{\r{g}}{r}{ei}{\textsubring{d}}{i}{\ng}{\r{g}}} \dicPos{f}[4] \dicFlx{(‑ar, ‑ar)}[5] \dicDirectTranslationCS{nástěnná výzdoba\,/\addthin dekorace}
\dicEntry[veggspjald] \dicTerm{vegg··|spjald} \dicIPA{{v}{\textepsilon}{\r{g}}{s}{\textsubring{b}}{j}{a}{l}{\textsubring{d}}} \dicPos{n}[2] \dicFlx{(‑spjalds, ‑spjöld)}[8] \dicSynonym{plakat} \dicDirectTranslationCS{plakát}
\dicEntry[veggteppi] \dicTerm{vegg··teppi} \dicIPA{{v}{\textepsilon}{\r{g}}{t\smash{\textsuperscript{h}}}{\textepsilon}{h}{\textsubring{b}}{\textsci}} \dicPos{n}[2] \dicFlx{(‑s, ‑)}[14] \dicDirectTranslationCS{gobelín, goblén, tapiserie, nástěnný koberec}
\dicEntry[veggur] \dicTerm{vegg|ur} \dicsymFrequent\  \dicIPA{{v}{\textepsilon}{\r{g}}{\textlengthmark}{\textscy}{\textsubring{r}}} \dicPos{m}[9] \dicFlx{(‑jar\,/\addthin ‑s, ‑ir)}[26] \textbf{1.} \dicSynonym*{hluti byggingar} \dicDirectTranslationCS{stěna, zeď (pokoje ap.)} \dicExampleIS{Myndin hangir á veggnum.} \dicExampleCS{Obraz visí na stěně.}  \textbf{2.} \dicSynonym*{utan um húsagarð} \dicDirectTranslationCS{zeď};  \dicIdiom{veggur}{ \dicPhraseIS{fara með veggjum}} \dicLangCat{přen.} \dicDirectTranslationCS{držet se při zdi} \dicIndirectTranslationCS{(postupovat opatrně, neriskovat)}; { \dicPhraseIS{ganga á vegg, reka sig á vegg}} \dicLangCat{přen.} \dicDirectTranslationCS{narazit na zeď} \dicIndirectTranslationCS{(setkat se s~odporem)}; { \dicPhraseIS{stilla e‑m upp við vegg}} \dicLangCat{přen.} \dicDirectTranslationCS{přitlačit (koho) ke zdi}
\dicEntry[veggöng] \dicTerm{veg··göng} \dicIPA{{v}{\textepsilon}{\r{g}}{\textlengthmark}{\oe i}{\ng}{\r{g}}} \dicPos{n}[2] \dicFlx{pl}[9] \dicDirectTranslationCS{(silniční) tunel}
\dicEntry[veghefill] \dicTerm{veg··hef|ill} \dicIPA{{v}{\textepsilon}{\textlengthmark}{x}{h}{\textepsilon}{v}{\textsci}{\textsubring{d}}{\textsubring{l}}} \dicPos{m}[6] \dicFlx{(‑ils, ‑lar)}[35] \dicDirectTranslationCS{srovnávač} \dicIndirectTranslationCS{(stavební stroj)}
\dicEntry[vegið] \dicTerm{vegið} \dicIPA{{v}{ei}{\textsci}{\texttheta}} \dicPos{v} \dicFlx{supin} \textbf{1.} \dicLink{vega\smash{\textsuperscript{1}}}  \textbf{2.} \dicLink{vega\smash{\textsuperscript{2}}}
\dicEntry[veginn] \dicTerm{veginn} \dicIPA{{v}{ei}{\textsci}{\textsubring{n}}} \dicPos{adj}[6]\dicFlx{}[-6] \textbf{1.} \dicDirectTranslationCS{zabitý}  \textbf{2.} \dicDirectTranslationCS{zvážený (na váze ap.)}
\dicEntry[vegkantur] \dicTerm{veg··kant|ur} \dicIPA{{v}{\textepsilon}{\textbabygamma}{k\smash{\textsuperscript{h}}}{a}{\textsubring{n}}{\textsubring{d}}{\textscy}{\textsubring{r}}} \dicPos{m}[6] \dicFlx{(‑s, ‑ar)}[5] \dicLink{vegarkantur}
\dicEntry[vegklæðning] \dicTerm{veg··klæð·ning} \dicIPA{{v}{\textepsilon}{\textbabygamma}{k\smash{\textsuperscript{h}}}{l}{a}{i}{ð}{n}{i}{\ng}{\r{g}}} \dicPos{f}[4] \dicFlx{(‑ar)}[7] \dicDirectTranslationCS{kryt vozovky} \dicIndirectTranslationCS{(asfaltový, cementobetonový ap.)}
\dicEntry[veglaus] \dicTerm{veg··laus} \dicIPA{{v}{\textepsilon}{\textbabygamma}{l}{\oe i}{s}} \dicPos{adj}[5]\dicFlx{}[-1] \dicSynonym{vegalaus} \dicDirectTranslationCS{(jsoucí) bez silnice\,/\addthin cesty}
\dicEntry[veglegur] \dicTerm{veg··legur} \dicIPA{{v}{\textepsilon}{\textbabygamma}{l}{\textepsilon}{\textbabygamma}{\textscy}{\textsubring{r}}} \dicPos{adj}[1]\dicFlx{}[-8] \textbf{1.} \dicSynonym{glæsilegur} \dicDirectTranslationCS{honosný, velkolepý}  \textbf{2.} \dicSynonym*{sem heiður er að} \dicDirectTranslationCS{důstojný} \dicExampleIS{veglegt sæti} \dicExampleCS{důstojné místo}
\dicEntry[vegleysa] \dicTerm{veg··leys|a} \dicIPA{{v}{\textepsilon}{\textbabygamma}{l}{ei}{s}{a}} \dicPos{f}[1] \dicFlx{(‑u, ‑ur)}[13] \dicDirectTranslationCS{terén bez silnic\,/\addthin cest}
\dicEntry[veglyndi] \dicTerm{veg··lyndi} \dicIPA{{v}{\textepsilon}{\textbabygamma}{l}{\textsci}{n}{\textsubring{d}}{\textsci}} \dicPos{n}[2] \dicFlx{(‑s)}[20] \dicSynonym{drengskapur} \dicDirectTranslationCS{šlechetnost, ušlechtilost}
\dicEntry[veglyndur] \dicTerm{veg··lyndur} \dicIPA{{v}{\textepsilon}{\textbabygamma}{l}{\textsci}{n}{\textsubring{d}}{\textscy}{\textsubring{r}}} \dicPos{adj}[2]\dicFlx{}[-14] \dicSynonym{göfugur} \dicDirectTranslationCS{šlechetný, ušlechtilý} \dicExampleIS{vera veglyndur og góðgjarn} \dicExampleCS{být ušlechtilý a~vlídný}
\dicEntry[vegmóður] \dicTerm{veg··móður} \dicIPA{{v}{\textepsilon}{\textbabygamma}{m}{ou}{ð}{\textscy}{\textsubring{r}}} \dicPos{adj}[2]\dicFlx{}[-6] \dicDirectTranslationCS{unavený z~pochodu\,/\addthin cesty}
\dicEntry[vegna] \dicTerm{vegn|a\smash{\textsuperscript{1}}} \dicIPA{{v}{\textepsilon}{\r{g}}{n}{a}} \dicPos{v}[1] \dicFlx{(‑aði)}[69] \dicFlx{impers} \dicPhraseIS{e‑m vegnar (vel)} \dicDirectTranslationCS{(komu) se daří\,/\addthin vede (dobře), (kdo) prospívá}
\dicEntry[vegna] \dicTerm{vegna\smash{\textsuperscript{2}}} \dicsymFrequent\  \dicIPA{{v}{\textepsilon}{\r{g}}{n}{a}} \dicPos{prep\,/\addthin adv} \dicFlx{gen} \dicSynonym{sökum} \dicDirectTranslationCS{kvůli, z~důvodu} \dicExampleIS{Lokað vegna veðurs.} \dicExampleCS{Zavřeno kvůli počasí.};  \dicPhraseIS{allra hluta vegna} \dicFlx{adv} \dicDirectTranslationCS{ve všech ohledech, v~každém směru};  \dicPhraseIS{beggja vegna} \dicFlx{adv} \dicDirectTranslationCS{z~obou stran, po obou stranách};  \dicPhraseIS{beggja vegna e‑s} \dicFlx{prep} \dicDirectTranslationCS{po obou stranách (čeho)};  \dicPhraseIS{hvers vegna} \dicFlx{adv} {\textbf{a.}} \dicDirectTranslationCS{proč, kvůli čemu} \dicIndirectTranslationCS{(v~přímé otázce)};  {\textbf{b.}} \dicDirectTranslationCS{proč, kvůli čemu} \dicIndirectTranslationCS{(ve vedlejší větě)};  \dicPhraseIS{mín vegna} \dicDirectTranslationCS{co se mě týká, pro mě za mě};  \dicPhraseIS{vegna þess, þess vegna} \dicFlx{adv} \dicDirectTranslationCS{proto, kvůli tomu};  \dicPhraseIS{vegna þess að} \dicFlx{conj} \dicSynonym*{af því að} \dicDirectTranslationCS{protože; kvůli tomu, že}
\dicEntry[vegrið] \dicTerm{veg··rið} \dicIPA{{v}{\textepsilon}{\textbabygamma}{r}{\textsci}{\texttheta}} \dicPos{n}[2] \dicFlx{(‑s, ‑)}[5] \dicDirectTranslationCS{svodidlo}
\dicEntry[vegsama] \dicTerm{veg··sam|a} \dicIPA{{v}{\textepsilon}{x}{s}{a}{m}{a}} \dicPos{v}[1] \dicFlx{(‑aði)}[13] \dicFlx{acc} \dicSynonym{lofa} \dicDirectTranslationCS{velebit, opěvovat, oslavovat} \dicExampleIS{vegsama guð} \dicExampleCS{velebit Boha}
\dicEntry[vegsauki] \dicTerm{vegs··auk|i} \dicIPA{{v}{\textepsilon}{x}{s}{\oe i}{\r{\textObardotlessj}}{\textsci}} \dicPos{m}[1] \dicFlx{(‑a)}[3] \dicSynonym{heiður\smash{\textsuperscript{1}}} \dicDirectTranslationCS{pocta, vyznamenání, hold, čest} \dicExampleIS{Þessi viðurkenning er vegsauki fyrir hana.} \dicExampleCS{Toto ocenění je pro ni čest.}
\dicEntry[vegsemd] \dicTerm{veg··semd} \dicIPA{{v}{\textepsilon}{x}{s}{\textepsilon}{m}{\textsubring{d}}} \dicPos{f}[7] \dicFlx{(‑ar, ‑ir)}[1] \dicSynonym{heiður\smash{\textsuperscript{1}}} \dicDirectTranslationCS{čest, pocta, vyznamenání} \dicExampleIS{Þetta eru viðurkennd verðlaun og því mikil vegsemd að hljóta þau.} \dicExampleCS{Toto je uznávaná cena, a~je proto velkou ctí ji obdržet.}
\dicEntry[vegsummerki] \dicTerm{vegs··um·merki} \dicIPA{{v}{\textepsilon}{x}{s}{\textscy}{m}{\textlengthmark}{\textepsilon}{\textsubring{r}}{\r{\textObardotlessj}}{\textsci}} \dicPos{n}[2] \dicFlx{pl}[17] \dicSynonym{verksummerki} \dicDirectTranslationCS{(viditelné) stopy\,/\addthin známky (po zemětřesení ap.)} \dicExampleIS{Það sjást enn vegsummerki.} \dicExampleCS{Stopy jsou ještě vidět.}
\dicEntry[vegu] \dicTerm{vegu} \dicIPA{{v}{\textepsilon}{\textlengthmark}{\textbabygamma}{\textscy}} \dicPos{adv} \textbf{1.} \dicPhraseIS{á tvo vegu} \dicFlx{adv} \dicDirectTranslationCS{dvěma způsoby};  \dicPhraseIS{á marga vegu} \dicFlx{adv} \dicDirectTranslationCS{mnoha způsoby};  \dicPhraseIS{á ýmsa vegu} \dicFlx{adv} \dicDirectTranslationCS{různými způsoby}  \textbf{2.} \dicPhraseIS{á þrjá vegu} \dicFlx{adv} \dicDirectTranslationCS{ze třech stran}
\dicEntry[vegur] \dicTerm{veg|ur} \dicsymFrequent\  \dicIPA{{v}{\textepsilon}{\textlengthmark}{\textbabygamma}{\textscy}{\textsubring{r}}} \dicPos{m}[10] \dicFlx{(‑ar\,/\addthin ‑s, ‑ir)}[29] \textbf{1.} \dicSynonym*{leið} \dicDirectTranslationCS{cesta, silnice} \dicExampleIS{opinber vegur} \dicExampleCS{veřejná cesta}  \textbf{2.} \dicSynonym{háttur\smash{\textsuperscript{1}}} \dicDirectTranslationCS{cesta, způsob} \dicExampleIS{á ýmsan veg} \dicExampleCS{různým způsobem};  \dicIdiom{vegur}{ \dicPhraseIS{á eigin vegum}} \dicFlx{adv} \dicDirectTranslationCS{samostatně, individuálně};  \dicPhraseIS{e‑m eru allir vegir færir} \dicFlx{impers} \dicLangCat{přen.} \dicDirectTranslationCS{(kdo) má všechny dveře otevřené}; { \dicPhraseIS{engan veginn}} \dicFlx{adv} \dicDirectTranslationCS{rozhodně\,/\addthin vůbec ne}; { \dicPhraseIS{fara villur vegar}} \dicDirectTranslationCS{(za)bloudit};  \dicPhraseIS{færa e‑ð til betri vegar} \dicLangCat{přen.} \dicDirectTranslationCS{obrátit (co) do pozitivna, vidět (co) v~lepším světle};  \dicPhraseIS{færa e‑ð til verri vegar} \dicLangCat{přen.} \dicDirectTranslationCS{obrátit (co) do negativna, vidět (co) v~horším světle};  \dicPhraseIS{koma í veg fyrir e‑ð} \dicDirectTranslationCS{předejít (čemu), zabránit (čemu)};  \dicPhraseIS{nokkurn veginn} \dicFlx{adv} \dicDirectTranslationCS{prakticky, doslova};  \dicPhraseIS{ryðja e‑m úr vegi} \dicLangCat{přen.} \dicDirectTranslationCS{odstranit (koho) z~cesty};  \dicPhraseIS{segja\,/\addthin vísa e‑m til vegar} \dicDirectTranslationCS{ukázat (komu) cestu\,/\addthin směr};  \dicPhraseIS{standa í vegi fyrir e‑u} \dicLangCat{přen.} \dicDirectTranslationCS{stát (čemu) v~cestě, bránit (čemu)};  \dicPhraseIS{vera (vel) á veg kominn með e‑ð} \dicLangCat{přen.} \dicDirectTranslationCS{být na dobré cestě s~(čím), pokročit s~(čím) (úkolem ap.)}; { \dicPhraseIS{vera í þann veginn að (gera e‑ð)}} \dicDirectTranslationCS{chystat se (udělat (co)), být na dobré cestě (udělat (co))};  \dicPhraseIS{vera vel\,/\addthin illa á vegi staddur í e‑u} \dicLangCat{přen.} \dicDirectTranslationCS{stát si dobře\,/\addthin špatně v~(čem)};  \dicPhraseIS{verða á vegi e‑rs} \dicLangCat{přen.} \dicDirectTranslationCS{zkřížit (komu) cestu, stát (komu) v~cestě};  \dicPhraseIS{virða e‑ð til betri vegar} \dicDirectTranslationCS{vidět (co) pozitivně, vidět (co) v~lepším světle};  \dicPhraseIS{það er ekki úr vegi að (gera e‑ð)} \dicDirectTranslationCS{není od věci (udělat (co))};  \dicPhraseIS{það er ekki vinnandi vegur} \dicDirectTranslationCS{to je zhola nemožné};  \dicPhraseIS{það má til sanns vegar færa} \dicDirectTranslationCS{je to fakt}
\dicEntry[vegvilltur] \dicTerm{veg··villtur} \dicIPA{{v}{\textepsilon}{\textbabygamma}{v}{\textsci}{\textsubring{l}}{\textsubring{d}}{\textscy}{\textsubring{r}}} \dicPos{adj}[1]\dicFlx{}[-10] \dicLink{vegavilltur}
\dicEntry[vegvísir] \dicTerm{veg··vís|ir} \dicIPA{{v}{\textepsilon}{\textbabygamma}{v}{i}{s}{\textsci}{\textsubring{r}}} \dicPos{m}[7] \dicFlx{(‑is, ‑ar)}[1] \dicSynonym{leiðarvísir} \dicDirectTranslationCS{směrovka, ukazatel, šipka} \dicExampleIS{Það er enga vegvísa að finna.} \dicExampleCS{Nikde nejsou žádné ukazatele.}
\dicEntry[veiða] \dicTerm{vei|ða} \dicsymFrequent\  \dicIPA{{v}{ei}{\textlengthmark}{ð}{a}} \dicPos{v}[2] \dicFlx{(‑ddi, ‑tt)}[167] \dicFlx{acc} \textbf{1.} \dicDirectTranslationCS{(na)chytat, (na)lovit, ulovit, rybařit} \dicIndirectTranslationCS{(o~rybolovu)} \dicExampleIS{veiða fisk} \dicExampleCS{rybařit}  \textbf{2.} \dicDirectTranslationCS{(u)lovit, nalovit} \dicExampleIS{veiða gæsir} \dicExampleCS{lovit husy};  \dicIdiom{veiða}[upp]{ \dicPhraseIS{veiða e‑ð upp úr e‑m}} \dicLangCat{přen.} \dicDirectTranslationCS{vyloudit (co) z~(koho), vymámit (co) z~(koho) (tajemství ap.)};  \dicIdiom{veiðast}{ \dicPhraseIS{veiðast}} \dicFlx{refl} \dicDirectTranslationCS{lovit se, chytat se (ryby ap.)}; { \dicPhraseIS{það veiðist e‑að}} \dicFlx{refl impers} \dicDirectTranslationCS{(co) se chytá, (co) se loví (makrela ap.)}; { \dicPhraseIS{það veiðist (vel)}} \dicFlx{refl impers} \dicDirectTranslationCS{je (dobrý) úlovek, (dobře) se chytá\,/\addthin loví (ryby ap.)}
\dicEntry[veiðarfæri] \dicTerm{veiðar··færi} \dicIPA{{v}{ei}{ð}{a}{\textsubring{r}}{f}{a}{i}{r}{\textsci}} \dicPos{n}[2] \dicFlx{(‑s, ‑)}[14] \dicDirectTranslationCS{rybářská výstroj, rybářské náčiní}
\dicEntry[veiði] \dicTerm{veið|i} \dicsymFrequent\  \dicIPA{{v}{ei}{\textlengthmark}{ð}{\textsci}} \dicPos{f}[2] \dicFlx{(‑i\,/\addthin ‑ar, ‑ar)}[3] \textbf{1.} \dicSynonym{afli} \dicDirectTranslationCS{úlovek} \dicIndirectTranslationCS{(zvláště ryb)} \dicExampleIS{fín veiði} \dicExampleCS{pěkný úlovek}  \textbf{2.} \dicPhraseIS{veiðar} \dicFlx{pl} {\textbf{a.}} \dicDirectTranslationCS{lov, honba};  {\textbf{b.}} \dicDirectTranslationCS{rybaření, rybolov};  \dicIdiom{veiði}{ \dicPhraseIS{nú ber vel í veiði}} \dicLangCat{přen.} \dicDirectTranslationCS{nyní je ideální možnost\,/\addthin příležitost}; { \dicPhraseIS{þetta er sýnd veiði en ekki gefinn}} \dicLangCat{přen.} \dicDirectTranslationCS{není to tak snadné, jak to vypadá}
\dicEntry[veiðiá] \dicTerm{veiði··á} \dicIPA{{v}{ei}{\textlengthmark}{ð}{\textsci}{au}} \dicPos{f}[4] \dicFlx{(‑r, ‑r)}[18] \dicDirectTranslationCS{řeka vhodná k~rybaření}
\dicEntry[veiðibjalla] \dicTerm{veiði··|bjalla} \dicIPA{{v}{ei}{\textlengthmark}{ð}{\textsci}{\textsubring{b}}{j}{a}{\textsubring{d}}{l}{a}} \dicPos{f}[1] \dicFlx{(‑bjöllu, ‑bjöllur)}[20] \dicFieldCat{zool.} \dicDirectTranslationCS{racek mořský} \textit{(l.~{\textLA{Larus marinus}})}  \dicsymPhoto\ 
\dicFigure{ds_image_veidibjalla_0_1.jpg}{Veiðibjalla}{Veiðibjalla - Andreas Trepte, CC BY-SA 2.5}
\dicEntry[veiðidýr] \dicTerm{veiði··dýr} \dicIPA{{v}{ei}{\textlengthmark}{ð}{\textsci}{\textsubring{d}}{i}{\textsubring{r}}} \dicPos{n}[2] \dicFlx{(‑s, ‑)}[5] \dicDirectTranslationCS{lovné zvíře, lovná zvěř}
\dicEntry[veiðiferð] \dicTerm{veiði··ferð} \dicIPA{{v}{ei}{\textlengthmark}{ð}{\textsci}{f}{\textepsilon}{r}{\texttheta}} \dicPos{f}[7] \dicFlx{(‑ar, ‑ir)}[1] \dicDirectTranslationCS{vyjížďka na lov\,/\addthin rybolov}
\dicEntry[veiðihár] \dicTerm{veiði··hár} \dicIPA{{v}{ei}{\textlengthmark}{ð}{\textsci}{h}{au}{\textsubring{r}}} \dicPos{n}[2] \dicFlx{(‑s, ‑)}[5] \dicDirectTranslationCS{vous (kočičí ap.)} \dicIndirectTranslationCS{(na čumáku některých zvířat)}
\dicEntry[veiðihundur] \dicTerm{veiði··hund|ur} \dicIPA{{v}{ei}{\textlengthmark}{ð}{\textsci}{h}{\textscy}{n}{\textsubring{d}}{\textscy}{\textsubring{r}}} \dicPos{m}[6] \dicFlx{(‑s, ‑ar)}[4] \dicDirectTranslationCS{lovecký pes}
\dicEntry[veiðihús] \dicTerm{veiði··hús} \dicIPA{{v}{ei}{\textlengthmark}{ð}{\textsci}{h}{u}{s}} \dicPos{n}[2] \dicFlx{(‑s, ‑)}[5] \dicDirectTranslationCS{rybářská chata, rybářský srub}
\dicEntry[veiðikvóti] \dicTerm{veiði··kvót|i} \dicIPA{{v}{ei}{\textlengthmark}{ð}{\textsci}{k\smash{\textsuperscript{h}}}{v}{ou}{\textsubring{d}}{\textsci}} \dicPos{m}[1] \dicFlx{(‑a, ‑ar)}[1] \dicDirectTranslationCS{rybolovná\,/\addthin lov\-ná kvóta}
\dicEntry[veiðileyfi] \dicTerm{veiði··leyfi} \dicIPA{{v}{ei}{\textlengthmark}{ð}{\textsci}{l}{ei}{v}{\textsci}} \dicPos{n}[2] \dicFlx{(‑s, ‑)}[14] \dicDirectTranslationCS{rybářské\,/\addthin lovecké povolení}
\dicEntry[veiðimaður] \dicTerm{veiði··|maður} \dicIPA{{v}{ei}{\textlengthmark}{ð}{\textsci}{m}{a}{ð}{\textscy}{\textsubring{r}}} \dicPos{m}[13] \dicFlx{(‑manns, ‑menn)}[2] \textbf{1.} \dicDirectTranslationCS{rybář(ka)}  \textbf{2.} \dicDirectTranslationCS{lovec, lovkyně}  \textbf{3.} \dicSynonym*{sérstakt spil} \dicIndirectTranslationCS{druh karetní hry}  \textbf{4.} \dicPhraseIS{Veiðimaðurinn} \dicFlx{def} \dicFieldCat{astro.} \dicSynonym*{Óríon} \dicDirectTranslationCS{Orion}
\dicEntry[veiðimálastjóri] \dicTerm{veiði·mála··stjór|i} \dicIPA{{v}{ei}{\textlengthmark}{ð}{\textsci}{m}{au}{l}{a}{s}{\textsubring{d}}{j}{ou}{r}{\textsci}} \dicPos{m}[1] \dicFlx{(‑a, ‑ar)}[1] \dicDirectTranslationCS{ředitel(ka) sladkovodního rybolovu}
\dicEntry[veiðiréttur] \dicTerm{veiði··rétt|ur} \dicIPA{{v}{ei}{\textlengthmark}{ð}{\textsci}{r}{j}{\textepsilon}{h}{\textsubring{d}}{\textscy}{\textsubring{r}}} \dicPos{m}[10] \dicFlx{(‑ar)}[7] \dicDirectTranslationCS{rybářské právo}
\dicEntry[veiðiskip] \dicTerm{veiði··skip} \dicIPA{{v}{ei}{\textlengthmark}{ð}{\textsci}{s}{\r{\textObardotlessj}}{\textsci}{\textsubring{b}}} \dicPos{n}[2] \dicFlx{(‑s, ‑)}[5] \dicDirectTranslationCS{rybářský člun, rybářská loď}
\dicEntry[veiðistöng] \dicTerm{veiði··|stöng} \dicIPA{{v}{ei}{\textlengthmark}{ð}{\textsci}{s}{\textsubring{d}}{\oe i}{\ng}{\r{g}}} \dicPos{f}[8] \dicFlx{(‑stangar, ‑stengur\,/\addthin ‑stangir)}[9] \dicDirectTranslationCS{rybářský prut}
\dicEntry[veiðisvæði] \dicTerm{veiði··svæði} \dicIPA{{v}{ei}{\textlengthmark}{ð}{\textsci}{s}{v}{a}{i}{ð}{\textsci}} \dicPos{n}[2] \dicFlx{(‑s, ‑)}[14] \dicDirectTranslationCS{loviště, revír, honitba}
\dicEntry[veiðitími] \dicTerm{veiði··tím|i} \dicIPA{{v}{ei}{\textlengthmark}{ð}{\textsci}{t\smash{\textsuperscript{h}}}{i}{m}{\textsci}} \dicPos{m}[1] \dicFlx{(‑a)}[3] \dicDirectTranslationCS{období (rybo)lovu}
\dicEntry[veiðiþjófnaður] \dicTerm{veiði··þjóf·nað|ur} \dicIPA{{v}{ei}{\textlengthmark}{ð}{\textsci}{\texttheta}{j}{ou}{\textsubring{b}}{n}{a}{ð}{\textscy}{\textsubring{r}}} \dicPos{m}[10] \dicFlx{(‑ar, ‑ir)}[8] \dicDirectTranslationCS{pytlačení, pytláctví}
\dicEntry[veiðiþjófur] \dicTerm{veiði··þjóf|ur} \dicIPA{{v}{ei}{\textlengthmark}{ð}{\textsci}{\texttheta}{j}{ou}{v}{\textscy}{\textsubring{r}}} \dicPos{m}[6] \dicFlx{(‑s, ‑ar)}[22] \dicDirectTranslationCS{pytlák, pytlačka}
\dicEntry[veif] \dicTerm{veif\smash{\textsuperscript{1}}} \dicIPA{{v}{ei}{\textlengthmark}{f}} \dicPos{f}[4] \dicFlx{(‑ar, ‑ar)}[1] \dicSynonym{veifa\smash{\textsuperscript{1}}} \dicDirectTranslationCS{praporek, vlaječka}
\dicEntry[veif] \dicTerm{veif\smash{\textsuperscript{2}}} \dicIPA{{v}{ei}{\textlengthmark}{f}} \dicPos{n}[2] \dicFlx{(‑s)}[2] \dicPhraseIS{annað veifið} \dicFlx{adv} \dicDirectTranslationCS{sem tam, občas, čas od času}
\dicEntry[veifa] \dicTerm{veif|a\smash{\textsuperscript{1}}} \dicIPA{{v}{ei}{\textlengthmark}{v}{a}} \dicPos{f}[1] \dicFlx{(‑u, ‑ur)}[19] \dicDirectTranslationCS{praporek, vlaječka}
\dicEntry[veifa] \dicTerm{veif|a\smash{\textsuperscript{2}}} \dicsymFrequent\  \dicIPA{{v}{ei}{\textlengthmark}{v}{a}} \dicPos{v}[1] \dicFlx{(‑aði)}[1] \dicFlx{dat} \textbf{1.} \dicFlx{dat} \dicDirectTranslationCS{(za)mávat, máchat (šátkem ap.)} \dicExampleIS{veifa fána} \dicExampleCS{mávat vlajkou}  \textbf{2.} \dicDirectTranslationCS{(za)mávat (na pozdrav ap.)};  \dicPhraseIS{veifa til e‑rs} \dicDirectTranslationCS{mávat (komu)} \dicExampleIS{veifa til fólksins} \dicExampleCS{mávat lidem}
\dicEntry[veigalítill] \dicTerm{veiga··|lítill} \dicIPA{{v}{ei}{\textlengthmark}{\textbabygamma}{a}{l}{i}{\textsubring{d}}{\textsci}{\textsubring{d}}{\textsubring{l}}} \dicPos{adj}[11] \dicFlx{(comp ‑minni, sup ‑minn\-st\-ur)}[10] \dicDirectTranslationCS{druhořadý, vedlejší, podružný} \dicExampleIS{veigalítið atriði} \dicExampleCS{druhořadý detail}
\dicEntry[veigamikill] \dicTerm{veiga··|mikill} \dicIPA{{v}{ei}{\textlengthmark}{\textbabygamma}{a}{m}{\textsci}{\r{\textObardotlessj}}{\textsci}{\textsubring{d}}{\textsubring{l}}} \dicPos{adj}[11] \dicFlx{(comp ‑meiri, sup ‑mestur)}[12] \dicSynonym{mikilvægur} \dicDirectTranslationCS{prvořadý, zásadní, hlavní} \dicExampleIS{veigamikil breyting} \dicExampleCS{zásadní změna}
\dicEntry[veigra] \dicTerm{veigr|a} \dicIPA{{v}{ei}{\textbabygamma}{r}{a}} \dicPos{v}[1] \dicFlx{(‑aði)}[1] \dicFlx{dat} \dicPhraseIS{veigra sér við e‑u} \dicDirectTranslationCS{zdráhat se (co), rozpakovat se (co)} \dicExampleIS{veigra sér við að fara þangað} \dicExampleCS{zdráhat se tam jít}
\dicEntry[veigur] \dicTerm{veig|ur} \dicIPA{{v}{ei}{\textlengthmark}{\textbabygamma}{\textscy}{\textsubring{r}}} \dicPos{m}[10] \dicFlx{(‑s\,/\addthin ‑ar)}[35] \dicSynonym{afl} \dicDirectTranslationCS{(vnitřní) síla, energie} \dicExampleIS{Það er veigur í sögunni.} \dicExampleCS{Ten příběh má v~sobě vnitřní sílu.}
\dicEntry[veikburða] \textls[15]{\dicTerm{veik··burða} \dicIPA{{v}{ei}{\textlengthmark}{\r{g}}{\textsubring{b}}{\textscy}{r}{ð}{a}} \dicPos{adj}[13] \dicFlx{indecl}[1] \textbf{1.} \dicSynonym{lasburða}} \dicDirectTranslationCS{slabý, zesláblý, neduživý}  \textbf{2.} \dicSynonym{kraftlítill} \dicDirectTranslationCS{slabý, chabý, chatrný (argument ap.)} \dicExampleIS{veikburða tilraun} \dicExampleCS{chabý pokus}
\dicEntry[veikbyggður] \dicTerm{veik··byggður} \dicIPA{{v}{ei}{\textlengthmark}{\r{g}}{\textsubring{b}}{\textsci}{\textbabygamma}{ð}{\textscy}{\textsubring{r}}} \dicPos{adj}[2]\dicFlx{}[-1] \dicSynonym{kraftlítill} \dicDirectTranslationCS{slabý, křehký, nepevný} \dicExampleIS{vera lítill og veikbyggður} \dicExampleCS{být malý a~křehký}
\dicEntry[veikgeðja] \dicTerm{veik··geðja} \dicIPA{{v}{ei}{\textlengthmark}{\r{g}}{\r{\textObardotlessj}}{\textepsilon}{ð}{j}{a}} \dicPos{adj}[13] \dicFlx{indecl}[1] \dicSynonym{ístöðulaus} \dicDirectTranslationCS{slabý (povahou), nepevný, (jsoucí) se slabou vůlí} \dicExampleIS{veikgeðja unglingur} \dicExampleCS{mladý člověk se slabou vůlí}
\dicEntry[veiki] \dicTerm{veik|i} \dicIPA{{v}{ei}{\textlengthmark}{\r{\textObardotlessj}}{\textsci}} \dicPos{f}[3] \dicFlx{(‑i, ‑ir)}[2] \dicSynonym{sjúkdómur} \dicDirectTranslationCS{nemoc, onemocnění} \dicExampleIS{Það kom upp alvarleg veiki í skólanum.} \dicExampleCS{Ve škole se objevilo závažné onemocnění.}
\dicEntry[veikindadagur] \dicTerm{veikinda··dag|ur} \dicIPA{{v}{ei}{\textlengthmark}{\r{\textObardotlessj}}{\textsci}{n}{\textsubring{d}}{a}{\textsubring{d}}{a}{\textbabygamma}{\textscy}{\textsubring{r}}} \dicPos{m}[6] \dicFlx{(‑s, ‑ar)}[62] \dicDirectTranslationCS{den nemocenské}
\dicEntry[veikindafrí] \dicTerm{veikinda··frí} \dicIPA{{v}{ei}{\textlengthmark}{\r{\textObardotlessj}}{\textsci}{n}{\textsubring{d}}{a}{f}{r}{i}} \dicPos{n}[2] \dicFlx{(‑s, ‑)}[5] \dicDirectTranslationCS{nemocenská (dovolená), zdravotní dovolená}
\dicEntry[veikindalegur] \dicTerm{veikinda··legur} \dicIPA{{v}{ei}{\textlengthmark}{\r{\textObardotlessj}}{\textsci}{n}{\textsubring{d}}{a}{l}{\textepsilon}{\textbabygamma}{\textscy}{\textsubring{r}}} \dicPos{adj}[1]\dicFlx{}[-8] \dicSynonym*{laslegur} \dicDirectTranslationCS{nemoc\-ný, churavý}
\dicEntry[veikindavottorð] \dicTerm{veikinda··vott·orð} \dicIPA{{v}{ei}{\textlengthmark}{\r{\textObardotlessj}}{\textsci}{n}{\textsubring{d}}{a}{v}{\textopeno}{h}{\textsubring{d}}{\textopeno}{r}{\texttheta}} \dicPos{n}[2] \dicFlx{(‑s, ‑)}[5] \dicDirectTranslationCS{potvrzení o~nemoci}
\dicEntry[veikindi] \dicTerm{veik··indi} \dicsymFrequent\  \dicIPA{{v}{ei}{\textlengthmark}{\r{\textObardotlessj}}{\textsci}{n}{\textsubring{d}}{\textsci}} \dicPos{n}[2] \dicFlx{pl}[19] \dicSynonym{sjúkdómur} \dicDirectTranslationCS{onemocnění, nemoc} \dicExampleIS{Hann tekur mikið af meðulum og lyfjum vegna veikinda í höfði.} \dicExampleCS{Bere hodně léků a~prášků kvůli onemocnění hlavy.}
\dicEntry[veikja] \dicTerm{veik|ja} \dicIPA{{v}{ei}{\textlengthmark}{\r{\textObardotlessj}}{a}} \dicPos{v}[2] \dicFlx{(‑ti, ‑t)}[25] \dicFlx{acc} \dicSynonym*{gera veikari} \dicDirectTranslationCS{oslabit, oslabovat, ochabnout, ochabovat} \dicExampleIS{Áfengið veikir vöðvakraftinn.} \dicExampleCS{Alkohol oslabuje sílu svalů.};  \dicIdiom{veikjast}{ \dicPhraseIS{veikjast}} \dicFlx{refl} {\textbf{a.}} \dicSynonym*{fá sjúkdóm} \dicDirectTranslationCS{onemocnět, ochořet} \dicExampleIS{veikjast af lungnabólgu} \dicExampleCS{onemocnět zápalem plic};  {\textbf{b.}} \dicSynonym*{verða haldminni} \dicDirectTranslationCS{ochabnout, ochabovat, ztratit\,/\addthin ztrácet sílu}
\dicEntry[veikla] \dicTerm{veikl|a} \dicIPA{{v}{ei}{h}{\r{g}}{l}{a}} \dicPos{v}[1] \dicFlx{(‑aði)}[1] \dicSynonym{veikja} \dicDirectTranslationCS{oslabit, ochabnout};  \dicIdiom{veiklast}{ \dicPhraseIS{veiklast}} \dicFlx{refl} \dicDirectTranslationCS{oslabit se, být oslaben, zeslábnout} \dicExampleIS{Heilsan veiklast.} \dicExampleCS{Zdraví je oslabeno.}
\dicEntry[veiklaður] \dicTerm{veikl··|aður} \dicIPA{{v}{ei}{h}{\r{g}}{l}{a}{ð}{\textscy}{\textsubring{r}}} \dicPos{adj}[3] \dicFlx{(f ‑uð)}[3] \dicSynonym{máttvana} \dicDirectTranslationCS{oslabený, zesláblý, vysílený} \dicExampleIS{veiklaður líkami} \dicExampleCS{oslabené tělo}
\dicEntry[veikleikamerki] \dicTerm{veik·leika··merki} \dicIPA{{v}{ei}{\textlengthmark}{\r{g}}{l}{ei}{\r{g}}{a}{m}{\textepsilon}{\textsubring{r}}{\r{\textObardotlessj}}{\textsci}} \dicPos{n}[2] \dicFlx{(‑s, ‑)}[16] \dicDirectTranslationCS{známka\,/\addthin náznak slabosti}
\dicEntry[veikleiki] \dicTerm{veik··leik|i} \dicIPA{{v}{ei}{\textlengthmark}{\r{g}}{l}{ei}{\r{\textObardotlessj}}{\textsci}} \dicPos{m}[1] \dicFlx{(‑a, ‑ar)}[1] \textbf{1.} \dicSynonym{veila} \dicDirectTranslationCS{slabost, nedokonalost, nedostatečnost}  \textbf{2.} \dicSynonym{breyskleiki} \dicDirectTranslationCS{slabost, slabina, slabé místo (charakteru ap.)}
\dicEntry[veiklulegur] \dicTerm{veiklu··legur} \dicIPA{{v}{ei}{h}{\r{g}}{l}{\textscy}{l}{\textepsilon}{\textbabygamma}{\textscy}{\textsubring{r}}} \dicPos{adj}[1]\dicFlx{}[-8] \dicSynonym{veikindalegur} \dicDirectTranslationCS{malátný, vypadající nemocně, nemocný (od pohledu)}
\dicEntry[veiklun] \dicTerm{veikl|un} \dicIPA{{v}{ei}{h}{\r{g}}{l}{\textscy}{\textsubring{n}}} \dicPos{f}[7] \dicFlx{(‑unar)}[9] \textbf{1.} \dicSynonym{hnignun} \dicDirectTranslationCS{oslabení, zeslabení (demokracie ap.)}  \textbf{2.} \dicSynonym{veikleiki} \dicDirectTranslationCS{oslabení, slabost, podlomení (zdraví)}
\dicEntry[veiklundaður] \dicTerm{veik··lund·|aður}\dicTerm{, veiklyndur} \dicIPA{{v}\-{ei}\-{\textlengthmark}\-{\r{g}}\-{l}\-{\textscy}\-{n}\-{\textsubring{d}}\-{a}\-{ð}\-{\textscy}\-{\textsubring{r}}\-} \dicPos{adj}[3] \dicFlx{(f ‑uð)}[3] \dicDirectTranslationCS{slabý (povahou), nepevný, mající slabou vůli}
\dicEntry[veiklyndi] \dicTerm{veik··lyndi} \dicIPA{{v}{ei}{\textlengthmark}{\r{g}}{l}{\textsci}{n}{\textsubring{d}}{\textsci}} \dicPos{n}[2] \dicFlx{(‑s)}[20] \dicSynonym{ístöðuleysi} \dicDirectTranslationCS{(povahová) slabost, slabá vůle}
\dicEntry[veiklyndur] \dicTerm{veik··lyndur} \dicIPA{{v}{ei}{\textlengthmark}{\r{g}}{l}{\textsci}{n}{\textsubring{d}}{\textscy}{\textsubring{r}}} \dicPos{adj}[2]\dicFlx{}[-14] \dicLink{veiklundaður}
\dicEntry[veikur] \dicTerm{veikur} \dicsymFrequent\  \dicIPA{{v}{ei}{\textlengthmark}{\r{g}}{\textscy}{\textsubring{r}}} \dicPos{adj}[1]\dicFlx{}[-1] \textbf{1.} \dicSynonym{lasinn} \dicDirectTranslationCS{nemocný} \dicExampleIS{verða veikur} \dicExampleCS{onemocnět}  \textbf{2.} \dicSynonym{kraftlítill} \dicDirectTranslationCS{slabý} \dicIndirectTranslationCS{(málo intenzivní)} \dicExampleIS{svara e‑m veikum rómi} \dicExampleCS{odpovědět (komu) slabým hlasem}  \textbf{3.} \dicDirectTranslationCS{slabý, tenký} \dicIndirectTranslationCS{(který má malou tloušťku)} \dicExampleIS{veikur ís} \dicExampleCS{slabý led}  \textbf{4.} \dicDirectTranslationCS{slabý} \dicIndirectTranslationCS{(málo odolný)} \dicExampleIS{veik von} \dicExampleCS{slabá naděje}  \textbf{5.} \dicFieldCat{jaz.} \dicDirectTranslationCS{slabý} \dicIndirectTranslationCS{(náležející do určité skloňovací\,/\addthin časovací třídy)} \dicAntonym{sterkur};  \dicPhraseIS{veik sögn} \dicFieldCat{jaz.} \dicDirectTranslationCS{slabé sloveso};  \dicIdiom{veikur}{ \dicPhraseIS{vera veikur fyrir e‑u}} \dicLangCat{přen.} \dicDirectTranslationCS{mít slabost pro (co)};  \dicPhraseIS{vera veikur í e‑ð} \dicLangCat{přen.} \dicDirectTranslationCS{být blázen do (čeho)}
\dicEntry[veila] \dicTerm{veil|a} \dicIPA{{v}{ei}{\textlengthmark}{l}{a}} \dicPos{f}[1] \dicFlx{(‑u, ‑ur)}[19] \dicSynonym{galli\smash{\textsuperscript{2}}} \dicDirectTranslationCS{vada, závada, defekt} \dicExampleIS{veila í persónuleika} \dicExampleCS{charakterová vada}
\dicEntry[veill] \dicTerm{veill} \dicIPA{{v}{ei}{\textsubring{d}}{\textsubring{l}}} \dicPos{adj}[8]\dicFlx{}[-1] \dicDirectTranslationCS{slabý, chatrný} \dicExampleIS{vera veill fyrir hjarta} \dicExampleCS{být slabý na srdce};  \dicPhraseIS{vera veill á geðsmunum} \dicDirectTranslationCS{být psychicky labilní}
\dicEntry[veimiltíta] \dicTerm{veimil··tít|a} \dicIPA{{v}{ei}{\textlengthmark}{m}{\textsci}{l}{t\smash{\textsuperscript{h}}}{i}{\textsubring{d}}{a}} \dicPos{f}[1] \dicFlx{(‑u, ‑ur)}[7] \textbf{1.} \dicSynonym*{kjarkleysingi} \dicDirectTranslationCS{slaboch, slaboška, měkkota}  \textbf{2.} \dicFieldCat{zool.} \dicDirectTranslationCS{jespák malý} \textit{(l.~{\textLA{Calidris minuta}})}  \dicsymPhoto\ 
\dicFigure{ds_image_veimiltita_0_2.jpg}{Veimiltíta}{Veimiltíta - Ken Billington, CC BY-SA 3.0}
\dicEntry[vein] \dicTerm{vein} \dicIPA{{v}{ei}{\textlengthmark}{\textsubring{n}}} \dicPos{n}[2] \dicFlx{(‑s, ‑)}[5] \dicSynonym{óp} \dicDirectTranslationCS{jekot, vřískot, skřek} \dicExampleIS{hræðilegt vein} \dicExampleCS{strašný vřískot}
\dicEntry[veina] \dicTerm{vein|a} \dicIPA{{v}{ei}{\textlengthmark}{n}{a}} \dicPos{v}[1] \dicFlx{(‑aði)}[1] \dicDirectTranslationCS{(za)ječet, (za)vřeštět} \dicExampleIS{veina af sársauka} \dicExampleCS{zaječet bolestí}
\dicEntry[veinandi] \dicTerm{vein··andi} \dicIPA{{v}{ei}{\textlengthmark}{n}{a}{n}{\textsubring{d}}{\textsci}} \dicPos{adj}[13] \dicFlx{indecl}[1] \dicDirectTranslationCS{ječící, vřeštící}
\dicEntry[veira] \dicTerm{veir|a} \dicIPA{{v}{ei}{\textlengthmark}{r}{a}} \dicPos{f}[1] \dicFlx{(‑u, ‑ur)}[7] \textbf{1.} \dicFieldCat{med.} \dicDirectTranslationCS{vir, virus}  \textbf{2.} \dicFieldCat{poč.} \dicSynonym{tölvuveira} \dicDirectTranslationCS{(počítačový) virus\,/\addthin vir}
\dicEntry[veirufræði] \dicTerm{veiru··fræð|i} \dicIPA{{v}{ei}{\textlengthmark}{r}{\textscy}{f}{r}{a}{i}{ð}{\textsci}} \dicPos{f}[3] \dicFlx{(‑i)}[3] \dicFieldCat{med. } \dicDirectTranslationCS{virologie}
\dicEntry[veirufræðingur] \dicTerm{veiru·fræð··ing|ur} \dicIPA{{v}{ei}{\textlengthmark}{r}{\textscy}{f}{r}{a}{i}{ð}{i}{\ng}{\r{g}}{\textscy}{\textsubring{r}}} \dicPos{m}[6] \dicFlx{(‑s, ‑ar)}[8] \dicFieldCat{med. } \dicDirectTranslationCS{virolog, viroložka}
\dicEntry[veirusjúkdómur] \dicTerm{veiru··sjúk·dóm|ur} \dicIPA{{v}{ei}{\textlengthmark}{r}{\textscy}{s}{j}{u}{\r{g}}{\textsubring{d}}{ou}{m}{\textscy}{\textsubring{r}}} \dicPos{m}[6] \dicFlx{(‑s, ‑ar)}[10] \dicFieldCat{med. } \dicDirectTranslationCS{virové onemocnění}
\dicEntry[veisla] \dicTerm{veisl|a} \dicsymFrequent\  \dicIPA{{v}{ei}{s}{\textsubring{d}}{l}{a}} \dicPos{f}[1] \dicFlx{(‑u, ‑ur)}[13] \dicSynonym{gildi} \dicDirectTranslationCS{oslava, večírek, posvícení, hostina} \dicExampleIS{halda veislu} \dicExampleCS{pořádat večírek}
\dicEntry[veisluborð] \dicTerm{veislu··borð} \dicIPA{{v}{ei}{s}{\textsubring{d}}{l}{\textscy}{\textsubring{b}}{\textopeno}{r}{\texttheta}} \dicPos{n}[2] \dicFlx{(‑s, ‑)}[5] \dicDirectTranslationCS{slavnostní stůl}
\dicEntry[veislusalur] \dicTerm{veislu··sal|ur} \dicIPA{{v}{ei}{s}{\textsubring{d}}{l}{\textscy}{s}{a}{l}{\textscy}{\textsubring{r}}} \dicPos{m}[10] \dicFlx{(‑ar, ‑ir)}[14] \dicDirectTranslationCS{hodovní síň\,/\addthin sál}
\dicEntry[veit] \dicTerm{veit} \dicIPA{{v}{ei}{\textlengthmark}{\textsubring{d}}} \dicPos{v} \dicFlx{ind praes sg 1 pers} \dicLink{vita\smash{\textsuperscript{1}}}
\dicEntry[veita] \dicTerm{veit|a\smash{\textsuperscript{1}}} \dicIPA{{v}{ei}{\textlengthmark}{\textsubring{d}}{a}} \dicPos{f}[1] \dicFlx{(‑u, ‑ur)}[13] \dicDirectTranslationCS{rozvodný systém} \dicExampleIS{rafveita} \dicExampleCS{elektrický rozvodný systém}
\dicEntry[veita] \dicTerm{veit|a\smash{\textsuperscript{2}}} \dicsymFrequent\  \dicIPA{{v}{ei}{\textlengthmark}{\textsubring{d}}{a}} \dicPos{v}[2] \dicFlx{(‑ti, ‑t)}[55] \dicFlx{dat (+ acc)} \textbf{1.} \dicFlx{dat + acc} \dicDirectTranslationCS{udělit, udělovat, dát, dávat, přiznat, přiznávat, poskytnout, poskytovat} \dicExampleIS{veita e‑m e‑ð} \dicExampleCS{udělit (komu co)};  \dicPhraseIS{veita e‑u athygli\,/\addthin eftirtekt} \dicDirectTranslationCS{věnovat (čemu) pozornost};  \dicPhraseIS{veita sér e‑ð} \dicDirectTranslationCS{dovolit si (co), dopřát si (co)}  \textbf{2.} \dicFlx{dat + acc} \dicDirectTranslationCS{podat, podávat, předložit, předkládat, servírovat} \dicExampleIS{veita gestunum pinnamat} \dicExampleCS{podávat hostům jednohubky}  \textbf{3.} \dicFlx{dat} \dicDirectTranslationCS{zavlažit, zavlažovat, zavodnit, zavodňovat}  \textbf{4.} \dicPhraseIS{veita viðnám\,/\addthin mótspyrnu} \dicDirectTranslationCS{klást odpor}  \textbf{5.} \dicPhraseIS{e‑m veitir vel\,/\addthin illa} \dicFlx{impers} \dicDirectTranslationCS{(komu) to jde dobře\,/\addthin špatně};  \dicIdiom{veita}[af]{ \dicPhraseIS{það veitir ekki af e‑u}} \dicFlx{impers} \dicDirectTranslationCS{(co) je zapotřebí};  \dicIdiom{veitast}{ \dicPhraseIS{e‑m veitist e‑að}} \dicFlx{refl impers} \dicDirectTranslationCS{(kdo) má (co), (komu) připadne (co) (čest ap.)}; { \dicPhraseIS{e‑m veitist e‑að auðvelt\,/\addthin erfitt}} \dicFlx{refl impers} \dicDirectTranslationCS{(komu co) připadá snadné\,/\addthin těžké};  \dicIdiom{veitast}[að]{ \dicPhraseIS{veitast að e‑m}} \dicFlx{refl} \dicDirectTranslationCS{zaútočit na (koho), napadnout (koho)}
\dicEntry[veitandi] \dicTerm{veit··|andi} \dicIPA{{v}{ei}{\textlengthmark}{\textsubring{d}}{a}{n}{\textsubring{d}}{\textsci}} \dicPos{m}[2] \dicFlx{(‑anda, ‑endur)}[1] \textbf{1.} \dicSynonym*{sá sem veitir} \dicDirectTranslationCS{poskytovatel(ka) (služeb ap.)}  \textbf{2.} \dicSynonym{veitingamaður} \dicDirectTranslationCS{hostinský, hostinská}
\dicEntry[veiting] \dicTerm{veit··ing} \dicIPA{{v}{ei}{\textlengthmark}{\textsubring{d}}{i}{\ng}{\r{g}}} \dicPos{f}[4] \dicFlx{(‑ar, ‑ar)}[5] \textbf{1.} \dicDirectTranslationCS{poskytnutí, udělení (pracovního místa ap.)} \dicExampleIS{veiting embætta} \dicExampleCS{udělení pozice}  \textbf{2.} \dicPhraseIS{veitingar} \dicFlx{pl} \dicDirectTranslationCS{občerstvení, pohoštění} \dicExampleIS{panta veitingar} \dicExampleCS{objednat občerstvení};  \dicPhraseIS{léttar veitingar} \dicDirectTranslationCS{lehké občerstvení}
\dicEntry[veitingahús] \dicTerm{veitinga··hús} \dicsymFrequent\  \dicIPA{{v}{ei}{\textlengthmark}{\textsubring{d}}{i}{\ng}{\r{g}}{a}{h}{u}{s}} \dicPos{n}[2] \dicFlx{(‑s, ‑)}[5] \dicDirectTranslationCS{restaurace, hostinec} \dicExampleIS{eigandi lítils veitingahúss} \dicExampleCS{majitel malé restaurace}
\dicEntry[veitingamaður] \dicTerm{veitinga··|maður} \dicIPA{{v}{ei}{\textlengthmark}{\textsubring{d}}{i}{\ng}{\r{g}}{a}{m}{a}{ð}{\textscy}{\textsubring{r}}} \dicPos{m}[13] \dicFlx{(‑manns, ‑menn)}[2] \dicDirectTranslationCS{hostinský, hostinská, restauratér(ka)}
\dicEntry[veitingasalur] \dicTerm{veitinga··sal|ur} \dicIPA{{v}{ei}{\textlengthmark}{\textsubring{d}}{i}{\ng}{\r{g}}{a}{s}{a}{l}{\textscy}{\textsubring{r}}} \dicPos{m}[10] \dicFlx{(‑ar, ‑ir)}[14] \dicDirectTranslationCS{jídelna, restaurace}
\dicEntry[veitingastaður] \dicTerm{veitinga··stað|ur} \dicIPA{{v}{ei}{\textlengthmark}{\textsubring{d}}{i}{\ng}{\r{g}}{a}{s}{\textsubring{d}}{a}{ð}{\textscy}{\textsubring{r}}} \dicPos{m}[10] \dicFlx{(‑ar, ‑ir)}[14] \dicSynonym{veitingahús} \dicDirectTranslationCS{restaurace, hostinec}
\dicEntry[vek] \dicTerm{vek} \dicIPA{{v}{\textepsilon}{\textlengthmark}{\r{g}}} \dicPos{v} \dicFlx{ind praes sg 1 pers} \dicLink{vekja}
\dicEntry[vekja] \dicTerm{vekja} \dicsymFrequent\  \dicIPA{{v}{\textepsilon}{\textlengthmark}{\r{\textObardotlessj}}{a}} \dicPos{v}[4] \dicFlx{(vek, vakti, vöktum, vekti, vakið)}[46] \dicFlx{acc} \textbf{1.} \dicDirectTranslationCS{vzbudit, (pro)budit, probouzet, probrat, probírat} \dicIndirectTranslationCS{(přivést ze spánku ap.)} \dicExampleIS{vekja hann kl. 7} \dicExampleCS{probudit ho v~sedm hodin};  \dicPhraseIS{vekja e‑n upp} \dicDirectTranslationCS{probudit (koho)};  \dicPhraseIS{vekja upp draug} \dicDirectTranslationCS{vyvolat ducha}  \textbf{2.} \dicDirectTranslationCS{vzbudit, (pro)budit, probouzet, vyvolat, vyvolávat} \dicIndirectTranslationCS{(způsobit vznik něčeho)} \dicExampleIS{vekja athygli} \dicExampleCS{vzbudit pozornost};  \dicPhraseIS{vekja upp e‑ð} \dicDirectTranslationCS{probudit (co), vyvolat (co) (vzpomínky ap.)} \dicExampleIS{Myndirnar vekja upp minningar.} \dicExampleCS{Fotografie probouzejí vzpomínky.}
\dicEntry[vekjandi] \dicTerm{vekj··andi} \dicIPA{{v}{\textepsilon}{\textlengthmark}{\r{\textObardotlessj}}{a}{n}{\textsubring{d}}{\textsci}} \dicPos{adj}[13] \dicFlx{indecl}[1] \dicDirectTranslationCS{podnětný, inspirující, inspirativní}
\dicEntry[vekjaraklukka] \dicTerm{vekjara··klukk|a} \dicIPA{{v}{\textepsilon}{\textlengthmark}{\r{\textObardotlessj}}{a}{r}{a}{k\smash{\textsuperscript{h}}}{l}{\textscy}{h}{\r{g}}{a}} \dicPos{f}[1] \dicFlx{(‑u, ‑ur)}[13] \dicDirectTranslationCS{budík}
\dicEntry[vekjari] \dicTerm{vekj··ar|i} \dicIPA{{v}{\textepsilon}{\textlengthmark}{\r{\textObardotlessj}}{a}{r}{\textsci}} \dicPos{m}[1] \dicFlx{(‑a, ‑ar)}[13] \dicSynonym{vekjaraklukka} \dicDirectTranslationCS{budík}
\dicEntry[vekti] \dicTerm{vekti} \dicIPA{{v}{\textepsilon}{x}{\textsubring{d}}{\textsci}} \dicPos{v} \dicFlx{con pf sg 1 pers} \dicLink{vekja}
\dicEntry[vektor] \dicTerm{vektor} \dicIPA{{v}{\textepsilon}{x}{\textsubring{d}}{\textopeno}{\textsubring{r}}} \dicPos{m}[4] \dicFlx{(‑s, ‑ar)}[14] \dicFieldCat{mat.} \dicDirectTranslationCS{vektor}
\dicEntry[vel] \dicTerm{vel\smash{\textsuperscript{1}}} \dicIPA{{v}{\textepsilon}{\textlengthmark}{\textsubring{l}}} \dicPos{v} \dicFlx{ind praes sg 1 pers} \dicLink{velja}
\dicEntry[vel] \dicTerm{vel\smash{\textsuperscript{2}}} \dicsymFrequent\  \dicIPA{{v}{\textepsilon}{\textlengthmark}{\textsubring{l}}} \dicPos{adv} \dicFlx{(comp betur, sup best)} \dicDirectTranslationCS{dobře} \dicExampleIS{Allt gengur vel.} \dicExampleCS{Všechno jde dobře.};  \dicPhraseIS{e‑m er vel við e‑ð} \dicFlx{impers} \dicDirectTranslationCS{(kdo) má (co) rád, (komu) je dobře s~(čím)};  \dicPhraseIS{gerðu svo vel!} \dicDirectTranslationCS{prosím!, račte!};  \dicPhraseIS{vel á minnst} \dicFlx{adv} \dicDirectTranslationCS{mimochodem, abych nezapomněl, à propos};  \dicPhraseIS{vera vel að sér í e‑u} \dicDirectTranslationCS{být dobrý\,/\addthin zběhlý v~(čem)}
\dicEntry[veld] \dicTerm{veld} \dicIPA{{v}{\textepsilon}{l}{\textsubring{d}}} \dicPos{v} \dicFlx{ind praes sg 1 pers} \textbf{1.} \dicLink{valda\smash{\textsuperscript{2}}}  \textbf{2.} \dicLink{valda\smash{\textsuperscript{3}}}
\dicEntry[veldi] \dicTerm{veldi\smash{\textsuperscript{1}}} \dicIPA{{v}{\textepsilon}{l}{\textsubring{d}}{\textsci}} \dicPos{n}[2] \dicFlx{(‑s, ‑)}[14] \textbf{1.} \dicSynonym{ríki} \dicDirectTranslationCS{mocnost, velmoc, říše, impérium}  \textbf{2.} \dicSynonym{vald} \dicDirectTranslationCS{moc, nadvláda}  \textbf{3.} \dicFieldCat{mat.} \dicDirectTranslationCS{mocnina} \dicExampleIS{fimm í þriðja veldi} \dicExampleCS{pět na třetí};  \dicIdiom{veldi}{ \dicPhraseIS{það er veldi á e‑m}} \dicFlx{impers} \dicDirectTranslationCS{(kdo) má moc a~peníze}; { \dicPhraseIS{í öllu sínu veldi}} \dicFlx{adv} \dicDirectTranslationCS{celou svou silou}
\dicEntry[veldi] \dicTerm{veldi\smash{\textsuperscript{2}}} \dicIPA{{v}{\textepsilon}{l}{\textsubring{d}}{\textsci}} \dicPos{v} \dicFlx{con pf sg 1 pers} \dicLink{velja}
\dicEntry[veldissproti] \dicTerm{veldis··sprot|i} \dicIPA{{v}{\textepsilon}{l}{\textsubring{d}}{\textsci}{s}{\textsubring{b}}{r}{\textopeno}{\textsubring{d}}{\textsci}} \dicPos{m}[1] \dicFlx{(‑a, ‑ar)}[1] \dicDirectTranslationCS{žezlo} \dicExampleIS{veldissproti konungs} \dicExampleCS{královské žezlo}
\dicEntry[veldisstofn] \dicTerm{veldis··stofn}\dicTerm{, veldistala} \dicIPA{{v}\-{\textepsilon}\-{l}\-{\textsubring{d}}\-{\textsci}\-{s}\-{\textsubring{d}}\-{\textopeno}\-{\textsubring{b}}\-{\textsubring{n}}\-} \dicPos{m}[4] \dicFlx{(‑s, ‑ar)}[3] \dicFieldCat{mat.} \dicDirectTranslationCS{základ mocniny, mocněnec}
\dicEntry[veldistala] \dicTerm{veldis··|tala} \dicIPA{{v}{\textepsilon}{l}{\textsubring{d}}{\textsci}{s}{t\smash{\textsuperscript{h}}}{a}{l}{a}} \dicPos{f}[1] \dicFlx{(‑tölu, ‑tölur)}[14] \dicLink{veldisstofn}
\dicEntry[veldisvísir] \dicTerm{veldis··vís|ir} \dicIPA{{v}{\textepsilon}{l}{\textsubring{d}}{\textsci}{s}{v}{i}{s}{\textsci}{\textsubring{r}}} \dicPos{m}[7] \dicFlx{(‑is, ‑ar)}[1] \dicFieldCat{mat.} \dicDirectTranslationCS{exponent, mocnitel}
\dicEntry[velfarnaður] \dicTerm{vel··far·nað|ur} \dicIPA{{v}{\textepsilon}{l}{f}{a}{r}{\textsubring{d}}{n}{a}{ð}{\textscy}{\textsubring{r}}} \dicPos{m}[10] \dicFlx{(‑ar)}[9] \dicSynonym{velgengni} \dicDirectTranslationCS{zdar, úspěch, sukces} \dicExampleIS{óska honum velfarnaðar} \dicExampleCS{popřát mu úspěch}
\dicEntry[velferð] \dicTerm{vel··ferð} \dicIPA{{v}{\textepsilon}{l}{f}{\textepsilon}{r}{\texttheta}} \dicPos{f}[7] \dicFlx{(‑ar)}[3] \dicDirectTranslationCS{blaho, blahobyt, prospěch, prosperita} \dicExampleIS{tryggja velferð þjóðarinnar} \dicExampleCS{zaručit prospěch národa}
\dicEntry[velferðarkerfi] \dicTerm{vel·ferðar··kerfi} \dicIPA{{v}{\textepsilon}{l}{f}{\textepsilon}{r}{ð}{a}{\textsubring{r}}{c\smash{\textsuperscript{h}}}{\textepsilon}{r}{v}{\textsci}} \dicPos{n}[2] \dicFlx{(‑s, ‑)}[14] \dicDirectTranslationCS{sociální systém, systém sociálního zabezpečení}
\dicEntry[velferðarráðherra] \dicTerm{vel·ferðar··ráð·herr|a} \dicIPA{{v}\-{\textepsilon}\-{l}\-{f}\-{\textepsilon}\-{r}\-{ð}\-{a}\-{r}\-{au}\-{\texttheta}\-{h}\-{\textepsilon}\-{r}\-{a}\-} \dicPos{m}[1] \dicFlx{(‑a, ‑ar)}[17] \dicDirectTranslationCS{ministr(yně) sociálních věcí}
\dicEntry[velferðarráðuneyti] \dicTerm{vel·ferðar··ráðu·neyti} \dicIPA{{v}\-{\textepsilon}\-{l}\-{f}\-{\textepsilon}\-{r}\-{ð}\-{a}\-{r}\-{au}\-{ð}\-{\textscy}\-{n}\-{ei}\-{\textsubring{d}}\-{\textsci}\-} \dicPos{n}[2] \dicFlx{(‑s, ‑)}[14] \dicDirectTranslationCS{ministerstvo sociálních věcí}
\dicEntry[velferðarríki] \dicTerm{vel·ferðar··ríki} \dicIPA{{v}{\textepsilon}{l}{f}{\textepsilon}{r}{ð}{a}{r}{i}{\r{\textObardotlessj}}{\textsci}} \dicPos{n}[2] \dicFlx{(‑s, ‑)}[16] \dicDirectTranslationCS{sociální stát}
\dicEntry[velgengni] \dicTerm{vel··gengn|i} \dicIPA{{v}{\textepsilon}{l}{\r{\textObardotlessj}}{ei}{\ng}{n}{\textsci}} \dicPos{f}[3] \dicFlx{(‑i)}[3] \dicDirectTranslationCS{zdar, úspěch, sukces} \dicExampleIS{eiga velgengni að fagna í e‑u} \dicExampleCS{slavit úspěch s~(čím)};  \dicPhraseIS{njóta velgengni} \dicDirectTranslationCS{těšit se z~úspěchu}
\dicEntry[velgerðamaður] \dicTerm{vel·gerða··|maður} \dicIPA{{v}{\textepsilon}{l}{\r{\textObardotlessj}}{\textepsilon}{r}{ð}{a}{m}{a}{ð}{\textscy}{\textsubring{r}}} \dicPos{m}[13] \dicFlx{(‑manns, ‑menn)}[2] \dicLink{velgerðarmaður}
\begin{xtolerant}{}{5pt}
\dicEntry[velgerðarmaður] \dicTerm{vel·gerðar··|maður}\dicTerm{, velgerðamaður}\addthinSS\textls[-13]{\dicIPA{{v}\-{\textepsilon}\-{l}\-{\r{\textObardotlessj}}\-{\textepsilon}\-{r}\-{ð}\-{a}\-{r}\-{m}\-{a}\-{ð}\-{\textscy}\-{\textsubring{r}}}}\addthinS\dicPos{m}[1\hspace{-0.06667em}3] \dicFlx{(‑manns, ‑menn)}[2] \dicDirectTranslationCS{dobrodinec, dobrodinka, mecenáš(ka), patron(ka)}
\end{xtolerant}
\dicEntry[velgerningur] \dicTerm{vel··ger·ning|ur}\dicTerm{, velgjörningur} \dicIPA{{v}\-{\textepsilon}\-{l}\-{\r{\textObardotlessj}}\-{\textepsilon}\-{r}\-{\textsubring{d}}\-{n}\-{i}\-{\ng}\-{\r{g}}\-{\textscy}\-{\textsubring{r}}\-} \dicPos{m}[6] \dicFlx{(‑s, ‑ar)}[8] \dicDirectTranslationCS{dobrodiní}
\dicEntry[velgja] \dicTerm{velgj|a\smash{\textsuperscript{1}}} \dicIPA{{v}{\textepsilon}{l}{\r{\textObardotlessj}}{a}} \dicPos{f}[1] \dicFlx{(‑u, ‑ur)}[25] \textbf{1.} \dicSynonym{ógleði} \dicDirectTranslationCS{nevolnost} \dicExampleIS{finna fyrir velgju} \dicExampleCS{cítit nevolnost}  \textbf{2.} \dicSynonym{hlýja\smash{\textsuperscript{1}}} \dicDirectTranslationCS{teplo}
\dicEntry[velgja] \dicTerm{velg|ja\smash{\textsuperscript{2}}} \dicIPA{{v}{\textepsilon}{l}{\r{\textObardotlessj}}{a}} \dicPos{v}[2] \dicFlx{(‑di, ‑t)}[138] \dicFlx{acc} \dicDirectTranslationCS{ohřát, ohřívat, zahřát, zahřívat} \dicExampleIS{velgja vatnið} \dicExampleCS{ohřát vodu};  \dicIdiom{velgja}[við]{ \dicPhraseIS{e‑n velgir við e‑u}} \dicFlx{impers} \dicDirectTranslationCS{(komu) se zvedá žaludek z~(čeho)} \dicExampleIS{Mig velgir við matnum.} \dicExampleCS{Zvedá se mi z~toho jídla žaludek.}
\dicEntry[velgjulegur] \dicTerm{velgju··legur} \dicIPA{{v}{\textepsilon}{l}{\r{\textObardotlessj}}{\textscy}{l}{\textepsilon}{\textbabygamma}{\textscy}{\textsubring{r}}} \dicPos{adj}[1]\dicFlx{}[-8] \dicDirectTranslationCS{nechutný, působící nevolnost} \dicExampleIS{velgjuleg rómantík} \dicExampleCS{nechutná romantika}
\dicEntry[velgjörðamaður] \dicTerm{vel·gjörða··|maður} \dicIPA{{v}{\textepsilon}{l}{\r{\textObardotlessj}}{\oe}{r}{ð}{a}{m}{a}{ð}{\textscy}{\textsubring{r}}} \dicPos{m}[13] \dicFlx{(‑manns, ‑menn)}[2] \dicLink{velgerðarmaður}
\begin{xtolerant}{}{1pt}
\dicEntry[velgjörðarmaður] \dicTerm{vel·gjörðar··|maður} \textls[-13]{\dicIPA{{v}{\textepsilon}{l}{\r{\textObardotlessj}}{\oe}{r}{ð}{a}{r}{m}{a}{ð}{\textscy}{\textsubring{r}}}} \dicPos{m}[13] \dicFlx{(‑manns, ‑menn)}[2] \dicLink{velgerðarmaður}
\end{xtolerant}
\dicEntry[velgjörðarsendiherra] \textls[15]{\dicTerm{vel·gjörðar··sendi·herr|a} \dicIPA{{v}\-{\textepsilon}\-{l}\-{\r{\textObardotlessj}}\-{\oe}\-{r}\-{ð}\-{a}\-{\textsubring{r}}\-{s}\-{\textepsilon}\-{n}\-{\textsubring{d}}\-{\textsci}\-{h}\-{\textepsilon}\-{r}\-{a}\-} \dicPos{m}[1] \dicFlx{(‑a, ‑ar)}[17]} \dicDirectTranslationCS{velvyslanec\,/\addthin velvyslankyně dobré vůle}
\dicEntry[velgjörningur] \dicTerm{vel··gjör·ning|ur} \dicIPA{{v}{\textepsilon}{l}{\r{\textObardotlessj}}{\oe}{r}{\textsubring{d}}{n}{i}{\ng}{\r{g}}{\textscy}{\textsubring{r}}} \dicPos{m}[6] \dicFlx{(‑s)}[9] \dicLink{velgerningur}
\dicEntry[velheppnaður] \dicTerm{vel··heppn·|aður} \dicIPA{{v}{\textepsilon}{\textlengthmark}{\textsubring{l}}{h}{\textepsilon}{h}{\textsubring{b}}{n}{a}{ð}{\textscy}{\textsubring{r}}} \dicPos{adj}[3] \dicFlx{(f ‑uð)}[4] \dicDirectTranslationCS{úspěšný, zdárný, zdařilý}
\dicEntry[velja] \dicTerm{velja} \dicsymFrequent\  \dicIPA{{v}{\textepsilon}{l}{j}{a}} \dicPos{v}[4] \dicFlx{(vel, valdi, völdum, veldi, valið)}[3] \dicFlx{(dat +) acc} \textbf{1.} \dicSynonym{kjósa} \dicDirectTranslationCS{vybrat, vybírat, (z)volit} \dicExampleIS{velja lit} \dicExampleCS{vybrat barvu};  \dicPhraseIS{velja á milli e‑s} \dicDirectTranslationCS{vybrat si mezi (čím)};  \dicPhraseIS{velja sér e‑ð} \dicDirectTranslationCS{vybrat si (co)} \dicExampleIS{velja sér föt} \dicExampleCS{vybrat si oblečení};  \dicPhraseIS{velja úr e‑u} \dicDirectTranslationCS{vybrat (si) z~(čeho)};  \dicPhraseIS{velja um e‑ð} \dicDirectTranslationCS{vybrat (si) z~(čeho)}  \textbf{2.} \dicDirectTranslationCS{vybrat si, zvolit si, rozhodnout se};  \dicIdiom{veljast}{ \dicPhraseIS{veljast}} \dicFlx{refl} \dicDirectTranslationCS{být vybrán\,/\addthin zvolen}
\dicEntry[velkja] \dicTerm{velk|ja} \dicIPA{{v}{\textepsilon}{\textsubring{l}}{\r{\textObardotlessj}}{a}} \dicPos{v}[2] \dicFlx{(‑ti, ‑t)}[26] \dicFlx{acc\,/\addthin dat} \textbf{1.} \dicSynonym{krumpa} \dicDirectTranslationCS{(po)mačkat, (z)muchlat} \dicExampleIS{velkja bókina} \dicExampleCS{pomačkat knížku}  \textbf{2.} \dicPhraseIS{e‑n velkir} \dicFlx{impers} \dicDirectTranslationCS{(kdo) je zmítán, (kdo) se zmítá} \dicExampleIS{Skipverja velkti lengi í hafi.} \dicExampleCS{Posádka se dlouho zmítala na moři.};  \dicIdiom{velkja}[fyrir]{ \dicPhraseIS{velkja e‑u fyrir sér}} \dicDirectTranslationCS{promyslet si (co)};  \dicIdiom{velkjast}{ \dicPhraseIS{velkjast}} \dicFlx{refl} \dicDirectTranslationCS{zmítat se, převalovat se}
\dicEntry[velkominn] \dicTerm{vel··kominn} \dicsymFrequent\  \dicIPA{{v}{\textepsilon}{l}{k\smash{\textsuperscript{h}}}{\textopeno}{m}{\textsci}{\textsubring{n}}} \dicPos{adj}[6]\dicFlx{}[-2] \dicDirectTranslationCS{vítaný};  \dicPhraseIS{e‑að er e‑m velkomið} \dicDirectTranslationCS{(kdo) může (co) (bez rozpaků)};  \dicPhraseIS{vertu velkominn!} \dicDirectTranslationCS{buď vítán!, vítej!}
\dicEntry[velktur] \dicTerm{velktur} \dicIPA{{v}{\textepsilon}{\textsubring{l}}{\textsubring{d}}{\textscy}{\textsubring{r}}} \dicPos{adj}[1]\dicFlx{}[-10] \dicDirectTranslationCS{zmačkaný, zmuchlaný} \dicExampleIS{velktir peningaseðlar} \dicExampleCS{zmačkané bankovky}
\dicEntry[vell] \dicTerm{vell} \dicIPA{{v}{\textepsilon}{\textsubring{d}}{\textsubring{l}}} \dicPos{v} \dicFlx{ind praes sg 1 pers} \dicLink{vella\smash{\textsuperscript{2}}}
\dicEntry[vella] \dicTerm{vell|a\smash{\textsuperscript{1}}} \dicIPA{{v}{\textepsilon}{\textsubring{d}}{l}{a}} \dicPos{f}[1] \dicFlx{(‑u, ‑ur)}[19] \textbf{1.} \dicSynonym{lasleiki} \dicDirectTranslationCS{malátnost, nevolnost}  \textbf{2.} \dicSynonym{væmni} \dicDirectTranslationCS{sentimentalita, sentimentální řeč}  \textbf{3.} \dicSynonym{suða\smash{\textsuperscript{1}}} \dicDirectTranslationCS{kypění, bublání}
\dicEntry[vella] \dicTerm{vella\smash{\textsuperscript{2}}} \dicIPA{{v}{\textepsilon}{\textsubring{d}}{l}{a}} \dicPos{v}[6] \dicFlx{(vell, vall, ullum, ylli, ollið)}[2] \textbf{1.} \dicDirectTranslationCS{přehnat se (bouře ap.)}  \textbf{2.} \dicSynonym{sjóða} \dicDirectTranslationCS{bublat, kypět}  \textbf{3.} \dicDirectTranslationCS{zpívat} \dicIndirectTranslationCS{(o~zvuku, který vydává koliha)}
\dicEntry[vellauðugur] \dicTerm{vell··auð·ugur} \dicIPA{{v}{\textepsilon}{\textsubring{d}}{l}{\oe i}{ð}{\textscy}{\textbabygamma}{\textscy}{\textsubring{r}}} \dicPos{adj}[1]\dicFlx{}[-6] \dicDirectTranslationCS{bohatý, zámožný}
\dicEntry[vellingur] \dicTerm{vell··ing|ur} \dicIPA{{v}{\textepsilon}{\textsubring{d}}{l}{i}{\ng}{\r{g}}{\textscy}{\textsubring{r}}} \dicPos{m}[6] \dicFlx{(‑s)}[9] \dicDirectTranslationCS{kaše, nákyp} \dicIndirectTranslationCS{(ovesná, rýžová ap. zvláště s~mlékem)}
\dicEntry[vellir] \dicTerm{vellir} \dicIPA{{v}{\textepsilon}{\textsubring{d}}{l}{\textsci}{\textsubring{r}}} \dicPos{m} \dicFlx{pl nom} \dicLink{völlur}
\dicEntry[vellíðan] \dicTerm{vel··líðan}\dicTerm{, vellíðun} \dicIPA{{v}\-{\textepsilon}\-{l}\-{\textlengthmark}\-{i}\-{ð}\-{a}\-{\textsubring{n}}\-} \dicPos{f}[4] \dicFlx{(‑ar)}[3] \dicSynonym*{góð líðan} \dicDirectTranslationCS{pocit zdraví, pohoda}
\dicEntry[vellíðun] \dicTerm{vel··líð|un} \dicIPA{{v}{\textepsilon}{l}{\textlengthmark}{i}{ð}{\textscy}{\textsubring{n}}} \dicPos{f}[7] \dicFlx{(‑unar)}[9] \dicLink{vellíðan}
\dicEntry[vellulegur] \dicTerm{vellu··legur} \dicIPA{{v}{\textepsilon}{\textsubring{d}}{l}{\textscy}{l}{\textepsilon}{\textbabygamma}{\textscy}{\textsubring{r}}} \dicPos{adj}[1]\dicFlx{}[-8] \textbf{1.} \dicSynonym{veiklulegur} \dicDirectTranslationCS{malátný, slabý}  \textbf{2.} \dicSynonym{væminn} \dicDirectTranslationCS{dojemný, sentimentální}
\dicEntry[vellyktandi] \dicTerm{vel··lykt·andi} \dicIPA{{v}{\textepsilon}{l}{\textlengthmark}{\textsci}{x}{\textsubring{d}}{a}{n}{\textsubring{d}}{\textsci}} \dicPos{n}[2] \dicFlx{(‑s, ‑)}[14] \dicLangCat{zast.} \dicSynonym{ilmvatn} \dicDirectTranslationCS{parfém, voňavka}
\dicEntry[velmegandi] \dicTerm{vel··meg·andi} \dicIPA{{v}{\textepsilon}{l}{m}{\textepsilon}{\textbabygamma}{a}{n}{\textsubring{d}}{\textsci}} \dicPos{adj}[13] \dicFlx{indecl}[1] \dicSynonym{efnaður} \dicDirectTranslationCS{bohatý, zámožný, prosperující}
\dicEntry[velmegun] \dicTerm{vel··meg|un} \dicIPA{{v}{\textepsilon}{l}{m}{\textepsilon}{\textbabygamma}{\textscy}{\textsubring{n}}} \dicPos{f}[7] \dicFlx{(‑unar)}[9] \dicSynonym{hagsæld} \dicDirectTranslationCS{prosperita, blahobyt} \dicExampleIS{lifa í velmegun} \dicExampleCS{žít v~blahobytu}
\dicEntry[velmetinn] \dicTerm{vel··metinn} \dicIPA{{v}{\textepsilon}{l}{m}{\textepsilon}{\textsubring{d}}{\textsci}{\textsubring{n}}} \dicPos{adj}[6]\dicFlx{}[-6] \dicDirectTranslationCS{oceňovaný, ceněný, uznávaný}
\dicEntry[velska] \dicTerm{velsk|a} \dicIPA{{v}{\textepsilon}{l}{s}{\r{g}}{a}} \dicPos{f}[1] \dicFlx{(‑u)}[5] \dicDirectTranslationCS{waleština, welština, velština}
\dicEntry[velskur] \dicTerm{velskur} \dicIPA{{v}{\textepsilon}{l}{s}{\r{g}}{\textscy}{\textsubring{r}}} \dicPos{adj}[1]\dicFlx{}[-1] \dicDirectTranslationCS{waleský, velšský}
\dicEntry[velsæla] \dicTerm{vel··sæl|a} \dicIPA{{v}{\textepsilon}{l}{s}{a}{i}{l}{a}} \dicPos{f}[1] \dicFlx{(‑u)}[5] \dicLink{velsæld}
\dicEntry[velsæld] \dicTerm{vel··sæld}\dicTerm{, velsæla} \dicIPA{{v}\-{\textepsilon}\-{l}\-{s}\-{a}\-{i}\-{l}\-{\textsubring{d}}\-} \dicPos{f}[7] \dicFlx{(‑ar)}[3] \dicSynonym{velmegun} \dicDirectTranslationCS{prosperita, blahobyt} \dicExampleIS{njóta velsældar} \dicExampleCS{těšit se z~blahobytu}
\dicEntry[velsæmi] \dicTerm{vel··sæmi} \dicIPA{{v}{\textepsilon}{l}{s}{a}{i}{m}{\textsci}} \dicPos{n}[2] \dicFlx{(‑s)}[20] \dicDirectTranslationCS{slušnost, způsobnost, dob\-ré způsoby}
\dicEntry[velsæmisbrot] \dicTerm{vel·sæmis··brot} \dicIPA{{v}{\textepsilon}{l}{s}{a}{i}{m}{\textsci}{s}{\textsubring{b}}{r}{\textopeno}{\textsubring{d}}} \dicPos{n}[2] \dicFlx{(‑s, ‑)}[5] \dicDirectTranslationCS{neslušnost, nepřístojnost, porušení slušného chování}
\dicEntry[velt] \dicTerm{velt} \dicIPA{{v}{\textepsilon}{\textsubring{l}}{\textsubring{d}}} \dicPos{v} \dicFlx{ind praes sg 1 pers} \dicLink{velta\smash{\textsuperscript{3}}}
\dicEntry[velta] \dicTerm{velt|a\smash{\textsuperscript{1}}} \dicIPA{{v}{\textepsilon}{\textsubring{l}}{\textsubring{d}}{a}} \dicPos{f}[1] \dicFlx{(‑u, ‑ur)}[19] \textbf{1.} \dicSynonym*{það að velta} \dicDirectTranslationCS{přetočení, převalení, obrácení}  \textbf{2.} \dicFieldCat{ekon.} \dicSynonym*{viðskiptavelta} \dicDirectTranslationCS{obrat}  \textbf{3.} \dicFieldCat{sport.} \dicDirectTranslationCS{salto, přemet}
\dicEntry[velta] \dicTerm{velt|a\smash{\textsuperscript{2}}} \dicsymFrequent\  \dicIPA{{v}{\textepsilon}{\textsubring{l}}{\textsubring{d}}{a}} \dicPos{v}[2] \dicFlx{(‑i, ‑)}[14] \dicFlx{dat} \textbf{1.} \dicSynonym{sveifla\smash{\textsuperscript{2}}} \dicDirectTranslationCS{(od)valit, válet} \dicExampleIS{velta steininum frá dyrunum} \dicExampleCS{odvalit kámen ze dveří};  \dicPhraseIS{velta vöngum} {\textbf{a.}} \dicDirectTranslationCS{kroutit hlavou};  {\textbf{b.}} \dicLangCat{přen.} \dicDirectTranslationCS{kroutit hlavou (nad situací ap.)} \dicExampleIS{velta vöngum yfir e‑u} \dicExampleCS{kroutit nad (čím) hlavou};  \dicPhraseIS{velta sér} \dicDirectTranslationCS{obrátit se, otočit se} \dicIndirectTranslationCS{ (v~leže)} \dicExampleIS{velta sér á magann} \dicExampleCS{obrátit se na břicho}  \textbf{2.} \dicFieldCat{ekon.} \dicDirectTranslationCS{mít obrat} \dicExampleIS{Fyrirtækið veltir stórfé.} \dicExampleCS{Firma má velký obrat.};  \dicIdiom{velta}[fyrir]{ \dicPhraseIS{velta e‑u fyrir sér}} \dicDirectTranslationCS{rozvážit si (co)};  \dicIdiom{veltast}{ \dicPhraseIS{veltast}} \dicFlx{refl} \dicDirectTranslationCS{válet se, převalovat se (v~trávě ap.)}
\dicEntry[velta] \dicTerm{velta\smash{\textsuperscript{3}}} \dicsymFrequent\  \dicIPA{{v}{\textepsilon}{\textsubring{l}}{\textsubring{d}}{a}} \dicPos{v}[6] \dicFlx{(velt, valt, ultum, ylti, oltið)}[1] \textbf{1.} \dicSynonym{rugga} \dicDirectTranslationCS{houpat se, kolébat se} \dicExampleIS{velta í sjóganginum} \dicExampleCS{houpat se ve vlnách}  \textbf{2.} \dicSynonym*{leggjast á hliðina} \dicDirectTranslationCS{převrhnout se, převrátit se} \dicExampleIS{Bíllinn valt.} \dicExampleCS{Auto se převrátilo.}  \textbf{3.} \dicSynonym{rúlla\smash{\textsuperscript{2}}} \dicDirectTranslationCS{valit se, točit se, kutálet se} \dicExampleIS{Hjólið veltur áfram.} \dicExampleCS{Kolo se točí dál.};  \dicIdiom{velta}[á]{ \dicPhraseIS{e‑að veltur á e‑u}} \dicDirectTranslationCS{(co) záleží na (čem)}
\dicEntry[veltihurð] \dicTerm{velti··hurð} \dicIPA{{v}{\textepsilon}{\textsubring{l}}{\textsubring{d}}{\textsci}{h}{\textscy}{r}{\texttheta}} \dicPos{f}[7] \dicFlx{(‑ar, ‑ir)}[1] \dicDirectTranslationCS{otáčecí dveře}
\dicEntry[veltufé] \dicTerm{veltu··|fé} \dicIPA{{v}{\textepsilon}{\textsubring{l}}{\textsubring{d}}{\textscy}{f}{j}{\textepsilon}} \dicPos{n}[3] \dicFlx{(‑fjár)}[5] \dicFieldCat{ekon.} \dicDirectTranslationCS{oběžný kapitál}
\dicEntry[veltuskattur] \dicTerm{veltu··skatt|ur} \dicIPA{{v}{\textepsilon}{\textsubring{l}}{\textsubring{d}}{\textscy}{s}{\r{g}}{a}{h}{\textsubring{d}}{\textscy}{\textsubring{r}}} \dicPos{m}[6] \dicFlx{(‑s, ‑ar)}[5] \dicFieldCat{ekon.} \dicDirectTranslationCS{daň z~obratu}
\dicEntry[velunnari] \dicTerm{vel··unn·ar|i} \dicIPA{{v}{\textepsilon}{\textlengthmark}{l}{\textscy}{n}{a}{r}{\textsci}} \dicPos{m}[1] \dicFlx{(‑a, ‑ar)}[13] \dicDirectTranslationCS{příznivec, příznivkyně, sympatizant(ka)}
\dicEntry[velvild] \dicTerm{vel··vild} \dicIPA{{v}{\textepsilon}{l}{v}{\textsci}{l}{\textsubring{d}}} \dicPos{f}[4] \dicFlx{(‑ar)}[3] \dicSynonym{góðvild} \dicDirectTranslationCS{dobrá vůle, dobrotivost} \dicExampleIS{mæta alls staðar velvild} \dicExampleCS{shledávat se všude s~dobrou vůlí}
\dicEntry[velviljaður] \dicTerm{vel··vilj·|aður} \dicIPA{{v}{\textepsilon}{l}{v}{\textsci}{l}{j}{a}{ð}{\textscy}{\textsubring{r}}} \dicPos{adj}[3] \dicFlx{(f ‑uð)}[3] \dicDirectTranslationCS{přející, mající dobrou vůli} \dicExampleIS{vera henni velviljaður} \dicExampleCS{přát jí to nejlepší}
\dicEntry[velvirðing] \textls[15]{\dicTerm{vel··virð·ing} \dicIPA{{v}{\textepsilon}{l}{v}{\textsci}{r}{ð}{i}{\ng}{\r{g}}} \dicPos{f}[4] \dicFlx{(‑ar)}[7] \dicDirectTranslationCS{shovívavost};}  \dicPhraseIS{biðja e‑n velvirðingar á e‑u} \dicDirectTranslationCS{požádat (koho) o~shovívavost v~(čem)}
\dicEntry[velvirkur] \dicTerm{vel··virkur} \dicIPA{{v}{\textepsilon}{l}{v}{\textsci}{\textsubring{r}}{\r{g}}{\textscy}{\textsubring{r}}} \dicPos{adj}[1]\dicFlx{}[-6] \dicSynonym{vandvirkur} \dicDirectTranslationCS{pečlivý, důkladný} \dicExampleIS{velvirkur smiður} \dicExampleCS{pečlivý řemeslník}
\dicEntry[velþóknun] \dicTerm{vel··þókn|un} \dicIPA{{v}{\textepsilon}{l}{\texttheta}{ou}{h}{\r{g}}{n}{\textscy}{\textsubring{n}}} \dicPos{f}[7] \dicFlx{(‑unar)}[9] \dicSynonym*{hugþokki} \dicDirectTranslationCS{zalíbení, libost, gusto} \dicExampleIS{horfa á myndirnar með velþóknun} \dicExampleCS{prohlížet si fotky se zalíbením};  \dicPhraseIS{hafa velþóknun á e‑u} \dicDirectTranslationCS{mít v~(čem) zalíbení}
\dicEntry[ven] \dicTerm{ven} \dicIPA{{v}{\textepsilon}{\textlengthmark}{\textsubring{n}}} \dicPos{v} \dicFlx{ind praes sg 1 pers} \dicLink{venja}
\dicEntry[venda] \dicTerm{ven|da} \dicIPA{{v}{\textepsilon}{n}{\textsubring{d}}{a}} \dicPos{v}[2] \dicFlx{(‑ti, ‑t)}[37] \dicFlx{dat} \dicFieldCat{nám.} \dicDirectTranslationCS{křižovat} \dicIndirectTranslationCS{(plachtit klikatou drahou proti směru větru)} \dicExampleIS{venda skipi} \dicExampleCS{křižovat s~lodí};  \dicIdiom{venda}{ \dicPhraseIS{eiga í ekkert hús að venda}} \dicLangCat{přen.} \dicDirectTranslationCS{nemít se už na koho obrátit, nemít útočiště} \dicIndirectTranslationCS{(nemít jinou možnost)}
\dicEntry[vendi] \dicTerm{vendi} \dicIPA{{v}{\textepsilon}{n}{\textsubring{d}}{\textsci}} \dicPos{v} \dicFlx{con pf sg 1 pers} \dicLink{venja}
\dicEntry[vendihnappur] \dicTerm{vendi··hnapp|ur} \dicIPA{{v}{\textepsilon}{n}{\textsubring{d}}{\textsci}{\textsubring{n}}{a}{h}{\textsubring{b}}{\textscy}{\textsubring{r}}} \dicPos{m}[6] \dicFlx{(‑s, ‑ar)}[25] \dicFieldCat{poč.} \dicDirectTranslationCS{klávesa Enter}
\dicEntry[vending] \dicTerm{vend··ing} \dicIPA{{v}{\textepsilon}{n}{\textsubring{d}}{i}{\ng}{\r{g}}} \dicPos{f}[4] \dicFlx{(‑ar, ‑ar)}[5] \dicFieldCat{nám.} \dicDirectTranslationCS{křižování} \dicIndirectTranslationCS{(plavba klikatou drahou proti směru větru)}
\dicEntry[vendir] \dicTerm{vendir} \dicIPA{{v}{\textepsilon}{n}{\textsubring{d}}{\textsci}{\textsubring{r}}} \dicPos{m} \dicFlx{pl nom} \dicLink{vöndur}
\dicEntry[Venesúela] \dicTerm{Venesúela} \dicIPA{{v}{\textepsilon}{\textlengthmark}{n}{\textepsilon}{s}{u}{\textepsilon}{l}{a}} \dicPos{n}[4] \dicFlx{indecl}[2] \dicFieldCat{geog.} \dicDirectTranslationCS{Venezuela}

\dicEntry[Venesúelamaður] \textls[24]{\dicTerm{Venesúela··|maður} \dicIPA{{v}{\textepsilon}{\textlengthmark}{n}{\textepsilon}{s}{u}{\textepsilon}{l}{a}{m}{a}{ð}{\textscy}{\textsubring{r}}} \dicPos{m}[13] \dicFlx{(‑manns, ‑menn)}[2]} \dicDirectTranslationCS{Venezuelec, Venezuelka}

\dicEntry[venesúelskur] \dicTerm{venesúelskur} \dicIPA{{v}{\textepsilon}{\textlengthmark}{n}{\textepsilon}{s}{u}{\textepsilon}{l}{s}{\r{g}}{\textscy}{\textsubring{r}}} \dicPos{adj}[1]\dicFlx{}[-6] \dicDirectTranslationCS{venezuelský}
\dicEntry[venja] \dicTerm{venj|a\smash{\textsuperscript{1}}} \dicsymFrequent\  \dicIPA{{v}{\textepsilon}{n}{j}{a}} \dicPos{f}[1] \dicFlx{(‑u, ‑ur)}[7] \dicSynonym{vani} \dicDirectTranslationCS{zvyk, obyčej, tradice};  \dicPhraseIS{að venju} \dicFlx{adv} \dicDirectTranslationCS{jako obvykle\,/\addthin vždy};  \dicPhraseIS{bregða út af venjunni} \dicDirectTranslationCS{porušit tradici, zlomit zvyk};  \dicPhraseIS{venju fremur} \dicFlx{adv} \dicDirectTranslationCS{více než obvykle, neobvykle}
\dicEntry[venja] \dicTerm{venja\smash{\textsuperscript{2}}} \dicsymFrequent\  \dicIPA{{v}{\textepsilon}{n}{j}{a}} \dicPos{v}[4] \dicFlx{(ven, vandi, vöndum, vendi, vanið)}[3] \dicFlx{acc} \textbf{1.} \dicDirectTranslationCS{zvyknout, zvykat,  přivyknout, přivykat};  \dicPhraseIS{venja e‑n við e‑ð} \dicDirectTranslationCS{přivyknout (koho) na (co)};  \dicPhraseIS{venja sig á e‑ð} \dicDirectTranslationCS{zvyknout si na (co)} \dicExampleIS{venja sig á að lesa upphátt} \dicExampleCS{zvyknout si číst nahlas};  \dicPhraseIS{venja sig af e‑u} \dicDirectTranslationCS{odvyknout si od (čeho)}  \textbf{2.} \dicDirectTranslationCS{(vy)cvičit (psa ap.)} \dicExampleIS{venja hund} \dicExampleCS{cvičit psa};  \dicIdiom{venjast}{ \dicPhraseIS{venjast e‑u}} \dicFlx{refl} \dicDirectTranslationCS{zvyknout si na (co), přivyknout si na (co)};  \dicPhraseIS{venjast á að (gera e‑ð)} \dicFlx{refl} \dicDirectTranslationCS{zvyknout si (dělat (co))};  \dicPhraseIS{e‑að venst} \dicFlx{refl} \dicDirectTranslationCS{na (co) se zvykne}
\dicEntry[venjubundinn] \dicTerm{venju··bundinn} \dicIPA{{v}{\textepsilon}{n}{j}{\textscy}{\textsubring{b}}{\textscy}{n}{\textsubring{d}}{\textsci}{\textsubring{n}}} \dicPos{adj}[6]\dicFlx{}[-2] \dicDirectTranslationCS{obvyklý, tradiční, běžný} \dicExampleIS{venjubundnar athafnir} \dicExampleCS{tradiční obřady}
\dicEntry[venjulega] \dicTerm{venju··leg|a} \dicsymFrequent\  \dicIPA{{v}{\textepsilon}{n}{j}{\textscy}{l}{\textepsilon}{\textbabygamma}{a}} \dicPos{adv} \dicFlx{(sup ‑ast)} \dicSynonym{vanalega} \dicDirectTranslationCS{obyčejně, obvykle, běžně} \dicExampleIS{Ég beið eins og venjulega þolinmóð úti í bíl.} \dicExampleCS{Jako obvykle jsem trpělivě čekala venku v~autě.}
\dicEntry[venjulegur] \dicTerm{venju··legur} \dicsymFrequent\  \dicIPA{{v}{\textepsilon}{n}{j}{\textscy}{l}{\textepsilon}{\textbabygamma}{\textscy}{\textsubring{r}}} \dicPos{adj}[1]\dicFlx{}[-8] \dicSynonym{vanalegur} \dicDirectTranslationCS{obyčejný, obvyklý, běžný} \dicExampleIS{Þetta er enginn venjulegur bíll.} \dicExampleCS{To není žádné obyčejné auto.}
\dicEntry[vensl] \dicTerm{vensl} \dicIPA{{v}{\textepsilon}{n}{s}{\textsubring{d}}{\textsubring{l}}} \dicPos{n}[2] \dicFlx{pl}[1] \textbf{1.} \dicSynonym{tengdir} \dicDirectTranslationCS{příbuznost}  \textbf{2.} \dicFieldCat{mat.} \dicSynonym*{mengi raðtvennda} \dicDirectTranslationCS{relace}
\dicEntry[venslaður] \dicTerm{vensl··|aður} \dicIPA{{v}{\textepsilon}{n}{s}{\textsubring{d}}{l}{a}{ð}{\textscy}{\textsubring{r}}} \dicPos{adj}[3] \dicFlx{(f ‑uð)}[3] \dicSynonym{tengdur} \dicDirectTranslationCS{příbuzný, spřízněný (svatbou)} \dicExampleIS{venslaður e‑m} \dicExampleCS{spřízněný s~(kým)}
\dicEntry[venslamaður] \dicTerm{vensla··|maður} \dicIPA{{v}{\textepsilon}{n}{s}{\textsubring{d}}{l}{a}{m}{a}{ð}{\textscy}{\textsubring{r}}} \dicPos{m}[13] \dicFlx{(‑manns, ‑menn)}[2] \dicDirectTranslationCS{příbuzný, příbuzná} \dicIndirectTranslationCS{(člověk spřízněný sňatkem)}
\dicEntry[ventill] \dicTerm{vent|ill} \dicIPA{{v}{\textepsilon}{\textsubring{n}}{\textsubring{d}}{\textsci}{\textsubring{d}}{\textsubring{l}}} \dicPos{m}[6] \dicFlx{(‑ils, ‑lar)}[35] \dicDirectTranslationCS{ventil}
\dicEntry[Venus] \dicTerm{Venus} \dicIPA{{v}{\textepsilon}{\textlengthmark}{n}{\textscy}{s}} \dicPos{f}[5] \dicFlx{(‑ar)}[8] \dicFieldCat{astro.} \dicDirectTranslationCS{Venuše}
\dicEntry[vepja] \dicTerm{vepj|a} \dicIPA{{v}{\textepsilon}{\textlengthmark}{\textsubring{b}}{j}{a}} \dicPos{f}[1] \dicFlx{(‑u, ‑ur)}[7] \dicFieldCat{zool.} \dicDirectTranslationCS{čejka chocholatá} \textit{(l.~{\textLA{Vanellus vanellus}})}  \dicsymPhoto\ 
\dicFigure{122096.jpg}{Vepja}{Vepja - Došlý Martin, Biolib, Copyright/PD}
\dicEntry[vepjulilja] \dicTerm{vepju··lilj|a} \dicIPA{{v}{\textepsilon}{\textlengthmark}{\textsubring{b}}{j}{\textscy}{l}{\textsci}{l}{j}{a}} \dicPos{f}[1] \dicFlx{(‑u, ‑ur)}[7] \dicFieldCat{bot.} \dicDirectTranslationCS{řebčík chocholatý} \textit{(l.~{\textLA{Fritillaria meleagris}})}  \dicsymPhoto\ 
\dicFigure{ds_image_vepjulilja_0_1.jpg}{Vepjulilja}{Vepjulilja - Michael Apel, CC BY 3.0}
\dicEntry[ver] \dicTerm{ver\smash{\textsuperscript{1}}} \dicIPA{{v}{\textepsilon}{\textlengthmark}{\textsubring{r}}} \dicPos{n}[2] \dicFlx{(‑s, ‑)}[5] \textbf{1.} \dicSynonym*{veiðistöð} \dicDirectTranslationCS{rybářská stanice}  \textbf{2.} \dicDirectTranslationCS{výrobna, továrna, fabrika} \dicIndirectTranslationCS{(jako poslední část složeniny)} \dicExampleIS{álver} \dicExampleCS{hliníkárna}  \textbf{3.} \dicDirectTranslationCS{povlak, povlečení, cícha (na polštář ap.)}
\dicEntry[ver] \dicTerm{ver\smash{\textsuperscript{2}}} \dicIPA{{v}{\textepsilon}{\textlengthmark}{\textsubring{r}}} \dicPos{v} \dicFlx{ind praes sg 1 pers} \textbf{1.} \dicLink{verja\smash{\textsuperscript{2}}}  \textbf{2.} \dicLink{verja\smash{\textsuperscript{3}}}
\dicEntry[vera] \dicTerm{ver|a\smash{\textsuperscript{1}}} \dicsymFrequent\  \dicIPA{{v}{\textepsilon}{\textlengthmark}{r}{a}} \dicPos{f}[1] \dicFlx{(‑u, ‑ur)}[7] \textbf{1.} \dicSynonym*{lifandi vera} \dicDirectTranslationCS{bytost, tvor, stvoření} \dicExampleIS{lifandi vera} \dicExampleCS{živý tvor}  \textbf{2.} \dicFieldCat{filos.} \dicSynonym{tilvist} \dicDirectTranslationCS{bytí}  \textbf{3.} \dicSynonym{dvöl} \dicDirectTranslationCS{pobyt, pobývání};  \dicPhraseIS{í raun og veru} \dicFlx{adv} \dicDirectTranslationCS{ve skutečnosti};  \dicPhraseIS{í verunni} \dicFlx{adv} \dicDirectTranslationCS{ve skutečnosti}
\dicEntry[vera] \dicTerm{vera\smash{\textsuperscript{2}}} \dicIPA{{v}{\textepsilon}{\textlengthmark}{r}{a}} \dicPos{v}[8] \dicFlx{(er, var, vorum, væri, verið)}[1] \textbf{1.} \dicDirectTranslationCS{být}  \textbf{2.} \dicDirectTranslationCS{být, ležet, nacházet se} \dicIndirectTranslationCS{(o~pozici)}  \textbf{3.} \dicSynonym*{dveljast} \dicDirectTranslationCS{být, pobývat, přebývat, zůstat, zůstávat, setrvat, setrvávat} \dicIndirectTranslationCS{(o~pobytu)}  \textbf{4.} \dicIndirectTranslationCS{pomocné sloveso -- vyjadřuje průběhový čas} \dicExampleIS{Ég er að lesa bókina.} \dicExampleCS{Právě si čtu knížku.}  \textbf{5.} \dicIndirectTranslationCS{pomocné sloveso -- vyjadřuje trpný rod} \dicExampleIS{Hann er barinn.} \dicExampleCS{Je bit.};  \dicIdiom{vera}[að]{ \dicPhraseIS{vera (vel) að sér um e‑ð}} \dicDirectTranslationCS{vyznat se v~(čem), mít (dobré) znalosti o~(čem)}; { \dicPhraseIS{vera alltaf að}} \dicDirectTranslationCS{být činorodý, stále něco dělat};  \dicIdiom{vera}[á]{ \dicPhraseIS{vera á því}} \dicDirectTranslationCS{myslet si, domnívat se, být toho názoru} \dicExampleIS{Ég var á því að það væri ekki hægt.} \dicExampleCS{Myslel jsem si, že to není možné.};  \dicIdiom{vera}[á eftir]{ \dicPhraseIS{vera á eftir}} \dicDirectTranslationCS{být\,/\addthin zůstat pozadu (v~školní docházce ap.)}; { \dicPhraseIS{vera á eftir e‑m}} \dicDirectTranslationCS{sledovat (koho), jít za (kým)};  \dicIdiom{vera}[á móti]{ \dicPhraseIS{vera á móti e‑u}} \dicDirectTranslationCS{být proti (čemu)};  \dicIdiom{vera}[eftir]{ \dicPhraseIS{e‑að er eftir}} \dicDirectTranslationCS{(co) zbývá, (co) zůstává (jídlo ap.)};  \dicIdiom{vera}[frá]{ \dicPhraseIS{vera frá í e‑u}} \dicDirectTranslationCS{být na tom zle s~(čím) (se zády ap.)}; { \dicPhraseIS{vera frá sér}} \dicDirectTranslationCS{být mimo, bláznit, zbláznit se} \dicExampleIS{Ertu alveg frá þér?} \dicExampleCS{Zbláznil ses?};  \dicIdiom{vera}[fyrir]{ \dicPhraseIS{e‑r er fyrir e‑m}} \dicDirectTranslationCS{(kdo komu) překáží, (kdo komu) vadí} \dicExampleIS{Hún leit á hann eins og hann væri fyrir henni.} \dicExampleCS{Podívala se na něho, jako by jí překážel.};  \dicIdiom{vera}[hjá]{ \dicPhraseIS{vera hjá e‑m}} \dicDirectTranslationCS{být v~péči (koho), chodit\,/\addthin docházet ke (komu) (k~lékaři ap.)};  \dicIdiom{vera}[í]{ \dicPhraseIS{e‑r er í e‑u}} {\textbf{a.}} \dicDirectTranslationCS{(kdo) má (co) na sobě, (kdo) je oblečen v~(čem)} \dicExampleIS{vera í sokkum} \dicExampleCS{mít na sobě ponožky};  {\textbf{b.}} \dicDirectTranslationCS{(kdo) je v~(čem) (práci ap.)} \dicExampleIS{vera í vinnu} \dicExampleCS{být v~práci};  {\textbf{c.}} \dicDirectTranslationCS{(kdo) v~(čem) lítá (v~drogách ap.)} \dicExampleIS{vera í dópi} \dicExampleCS{lítat v~drogách};  {\textbf{d.}} \dicDirectTranslationCS{(kdo) hraje (co) (fotbal ap.)};  {\textbf{e.}} \dicDirectTranslationCS{(kdo) chodí na (co) (na kurzy karate ap.)}; { \dicPhraseIS{vera í því}} \dicDirectTranslationCS{být v~tom} \dicIndirectTranslationCS{(být opilý)};  \dicIdiom{vera}[með]{ \dicPhraseIS{vera með}} \dicSynonym*{taka þátt í} \dicDirectTranslationCS{(z)účastnit se, zapojit se}; { \dicPhraseIS{vera með e‑ð}} \dicSynonym{eiga\smash{\textsuperscript{2}}} \dicDirectTranslationCS{mít (co)} \dicExampleIS{vera með falleg augu} \dicExampleCS{mít hezké oči}; { \dicPhraseIS{vera með e‑n}} \dicDirectTranslationCS{mít (koho) s~sebou}; { \dicPhraseIS{vera með e‑m}} \dicSynonym{fylgja\smash{\textsuperscript{2}}} \dicDirectTranslationCS{chodit s~(kým) (s~dívkou ap.)} \dicExampleIS{Er hann með henni?} \dicExampleCS{Chodí s~ní?};  \dicIdiom{vera}[saman]{ \dicPhraseIS{vera saman}} {\textbf{a.}} \dicDirectTranslationCS{být spolu};  {\textbf{b.}} \dicDirectTranslationCS{chodit spolu, být spolu} \dicIndirectTranslationCS{(jako pár)};  {\textbf{c.}} \dicDirectTranslationCS{být spolu} \dicIndirectTranslationCS{(jako spoluhráči)};  \dicIdiom{vera}[til]{ \dicPhraseIS{vera til}} {\textbf{a.}} \dicSynonym*{vera á lífi} \dicDirectTranslationCS{existovat, žít};  {\textbf{b.}} \dicSynonym*{vera tilbúinn} \dicDirectTranslationCS{být připravený\,/\addthin hotový} \dicExampleIS{Ertu til?} \dicExampleCS{Jsi hotový?}; { \dicPhraseIS{vera til í að (gera e‑ð)}} \dicDirectTranslationCS{být připraven k~(dělání (čeho))};  \dicIdiom{vera}[um]{ \dicPhraseIS{e‑m er ekki\,/\addthin ekkert um e‑ð}} \dicFlx{impers} \dicDirectTranslationCS{(kdo co) nemusí, (kdo) nemá (co) rád};  \dicIdiom{vera}[uppi]{ \dicPhraseIS{vera uppi}} {\textbf{a.}} \dicDirectTranslationCS{být nahoře (v~prvním patře ap.)};  {\textbf{b.}} \dicDirectTranslationCS{žít} \dicExampleIS{meðan við erum uppi} \dicExampleCS{zatímco žijeme};  \dicIdiom{vera}[úr]{ \dicPhraseIS{e‑að er úr e‑u}} \dicDirectTranslationCS{(co) je z~(čeho) (stůl ze dřeva ap.)}; { \dicPhraseIS{vera úr leik}} \dicLangCat{přen.} \dicDirectTranslationCS{být mimo hru};  \dicIdiom{vera}[út úr]{ \dicPhraseIS{e‑að er út úr myndinni}} \dicDirectTranslationCS{(co) už není aktuální}; { \dicPhraseIS{vera út úr}} \dicDirectTranslationCS{ležet\,/\addthin stát stranou};  \dicIdiom{vera}[við]{ \dicPhraseIS{vera við}} \dicDirectTranslationCS{být přítomný, být na místě} \dicExampleIS{vera við milli 4 og 5} \dicExampleCS{být na místě mezi čtvrtou a~pátou};  \dicIdiom{vera}[yfir]{ \dicPhraseIS{vera yfir e‑u}} \dicDirectTranslationCS{řídit (co), vést (co) (pobočku ap.)};  \dicIdiom{vera}{ \dicPhraseIS{er það?}} \dicDirectTranslationCS{opravdu?, skutečně?}; { \dicPhraseIS{láta e‑ð\,/\addthin e‑n vera}} \dicDirectTranslationCS{nechat (co\,/\addthin koho) být, nechat (co\,/\addthin koho) na pokoji};  \dicPhraseIS{vera búinn að vera} {\textbf{a.}} \dicDirectTranslationCS{být skončený\,/\addthin dokončený (závod ap.)};  {\textbf{b.}} \dicDirectTranslationCS{být vyřízený, být na pokraji sil}
\dicEntry[veraldarhyggja] \dicTerm{ver·aldar··hyggj|a} \dicIPA{{v}{\textepsilon}{\textlengthmark}{r}{a}{l}{\textsubring{d}}{a}{\textsubring{r}}{h}{\textsci}{\r{\textObardotlessj}}{a}} \dicPos{f}[1] \dicFlx{(‑u)}[5] \dicDirectTranslationCS{sekularismus}
\dicEntry[veraldarsaga] \dicTerm{ver·aldar··|saga} \dicIPA{{v}{\textepsilon}{\textlengthmark}{r}{a}{l}{\textsubring{d}}{a}{\textsubring{r}}{s}{a}{\textbabygamma}{a}} \dicPos{f}[1] \dicFlx{(‑sögu)}[2] \dicSynonym{mannkynssaga} \dicDirectTranslationCS{světové dějiny, historie světa}
\dicEntry[veraldarvanur] \dicTerm{ver·aldar··|vanur} \dicIPA{{v}{\textepsilon}{\textlengthmark}{r}{a}{l}{\textsubring{d}}{a}{r}{v}{a}{n}{\textscy}{\textsubring{r}}} \dicPos{adj}[1] \dicFlx{(f ‑vön)}[2] \dicDirectTranslationCS{světaznalý, kosmopolitní}
\dicEntry[veraldlegur] \dicTerm{ver·ald··legur} \dicIPA{{v}{\textepsilon}{\textlengthmark}{r}{a}{l}{\textsubring{d}}{l}{\textepsilon}{\textbabygamma}{\textscy}{\textsubring{r}}} \dicPos{adj}[1]\dicFlx{}[-8] \textbf{1.} \dicSynonym{jarðneskur} \dicDirectTranslationCS{pozemský}  \textbf{2.} \dicSynonym*{ekki kirkjulegur} \dicDirectTranslationCS{světský} \dicExampleIS{veraldlegur leiðtogi} \dicExampleCS{světský vůdce}
\dicEntry[verbúð] \dicTerm{ver··búð} \dicIPA{{v}{\textepsilon}{r}{\textsubring{b}}{u}{\texttheta}} \dicPos{f}[7] \dicFlx{(‑ar, ‑ir)}[1] \dicDirectTranslationCS{rybářská chata\,/\addthin bouda}
\dicEntry[verð] \dicTerm{verð\smash{\textsuperscript{1}}} \dicsymFrequent\  \dicIPA{{v}{\textepsilon}{r}{\texttheta}} \dicPos{n}[2] \dicFlx{(‑s, ‑)}[5] \dicDirectTranslationCS{cena} \dicIndirectTranslationCS{(v~penězích)} \dicExampleIS{verð á e‑u} \dicExampleCS{cena (čeho)};  \dicPhraseIS{koma e‑u í verð} \dicDirectTranslationCS{prodat (co)};  \dicPhraseIS{e‑að var dýru verði keypt} \dicLangCat{přen.} \dicDirectTranslationCS{(co) přišlo draho}
\dicEntry[verð] \dicTerm{verð\smash{\textsuperscript{2}}} \dicIPA{{v}{\textepsilon}{r}{\texttheta}} \dicPos{v} \dicFlx{ind praes sg 1 pers} \dicLink{verða}
\dicEntry[verða] \dicTerm{verða} \dicsymFrequent\  \dicIPA{{v}{\textepsilon}{r}{ð}{a}} \dicPos{v}[6] \dicFlx{(verð, varð, urðum, yrði, orðið)}[37] \textbf{1.} \dicIndirectTranslationCS{pomocné sloveso -- vyjadřuje budoucí čas} \dicExampleIS{Það verður gaman að tala við hana.} \dicExampleCS{Bude to legrace s~ní mluvit.}  \textbf{2.} \dicDirectTranslationCS{muset, musit} \dicExampleIS{verða nauðsynlega að fara} \dicExampleCS{muset nutně odejít}  \textbf{3.} \dicSynonym*{gerast} \dicDirectTranslationCS{stát se (lékařem ap.)} \dicExampleIS{verða kennari} \dicExampleCS{stát se učitelem}  \textbf{4.} \dicSynonym*{dveljast} \dicDirectTranslationCS{zůstat, zůstávat, setrvat, setrvávat} \dicExampleIS{verða heima} \dicExampleCS{zůstat doma}  \textbf{5.} \dicPhraseIS{e‑m verður heitt\,/\addthin kalt} \dicFlx{impers} \dicDirectTranslationCS{(komu) začíná být horko\,/\addthin zima};  \dicPhraseIS{e‑m verður illt} \dicFlx{impers} \dicDirectTranslationCS{(komu) se dělá špatně, (komu) začíná být špatně};  \dicIdiom{verða}[að]{ \dicPhraseIS{e‑að verður að engu}} \dicDirectTranslationCS{nic z~(čeho) nebude, (co) se nestane skutečností}; { \dicPhraseIS{verða að e‑u}} \dicDirectTranslationCS{proměnit se v~(co), stát se (čím)} \dicExampleIS{Ísinn verður að vatni.} \dicExampleCS{Led se promění ve vodu.}; { \dicPhraseIS{verði þér\,/\addthin ykkur að góðu!}} \dicFlx{imper} \dicDirectTranslationCS{dobrou chuť!};  \dicIdiom{verða}[af]{ \dicPhraseIS{e‑m verður meint af e‑u}} \dicFlx{impers} \dicDirectTranslationCS{(komu) se udělá zle z~(čeho) (z~jídla ap.)}; { \dicPhraseIS{það verður ekkert af e‑u}} \dicFlx{impers} \dicDirectTranslationCS{nic z~(čeho) nebude, (co) se neuskuteční};  \dicIdiom{verða}[á]{ \dicPhraseIS{e‑m verður á e‑ð}} \dicFlx{impers} \dicDirectTranslationCS{(komu) se (co) stává, (komu) se přihodí (co) (chyba ap.)};  \dicIdiom{verða}[eftir]{ \dicPhraseIS{verða eftir}} {\textbf{a.}} \dicDirectTranslationCS{zůstat (doma ap.)} \dicIndirectTranslationCS{(nikam nejít)};  {\textbf{b.}} \dicDirectTranslationCS{zůstat, být ponechán\,/\addthin zapomenut (kufr ve vlaku ap.)};  \dicIdiom{verða}[fyrir]{ \dicPhraseIS{verða fyrir e‑u}} \dicDirectTranslationCS{zažít (co), zakusit (co), být vystaven (čemu) (šoku ap.)} \dicExampleIS{Fólk varð fyrir áfalli.} \dicExampleCS{Lidé zažili šok.}; { \dicPhraseIS{verða fyrir bíl}} \dicDirectTranslationCS{být sražen autem, být přejet (autem)};  \dicIdiom{verða}[ofan á]{ \dicPhraseIS{e‑að verður ofan á}} \dicDirectTranslationCS{(co) je přijato};  \dicIdiom{verða}[til]{ \dicPhraseIS{e‑að verður til e‑s}} \dicDirectTranslationCS{(co) způsobuje (co), (co) je příčinou (čeho)}; { \dicPhraseIS{verða til}} \dicDirectTranslationCS{vzniknout, vznikat, být vytvořen, povstat, povstávat} \dicExampleIS{Ljóðið varð til.} \dicExampleCS{Vznikla báseň.};  \dicIdiom{verða}[um]{ \dicPhraseIS{e‑m verður mikið um}} \dicFlx{impers} \dicDirectTranslationCS{(kdo) je šokován, (kdo) je v~šoku}; { \dicPhraseIS{hvað verður um e‑n?}} \dicDirectTranslationCS{co bude s~(kým)?, co se s~(kým) stane?};  \dicIdiom{verða}[undir]{ \dicPhraseIS{verða undir}} \dicDirectTranslationCS{prohrát, podlehnout, být poražen\,/\addthin pokořen};  \dicIdiom{verða}[úr]{ \dicPhraseIS{e‑að verður úr}} \dicDirectTranslationCS{výsledkem je (co)}; { \dicPhraseIS{það verður eitthvað úr e‑m}} \dicFlx{impers} \dicLangCat{přen.} \dicDirectTranslationCS{něco z~(koho) bude}; { \dicPhraseIS{það verður ekkert úr e‑u}} \dicFlx{impers} \dicDirectTranslationCS{nic z~(čeho) nebude (z~deště ap.)} \dicExampleIS{Það varð ekkert úr rigningunni.} \dicExampleCS{Z~deště nic nebylo.};  \dicIdiom{verða}[úti]{ \dicPhraseIS{e‑að verður illa úti}} \dicDirectTranslationCS{(co) je poškozeno\,/\addthin zničeno}; { \dicPhraseIS{verða úti}} \dicDirectTranslationCS{zahynout, umrznout};  \dicIdiom{verða}[úti um]{ \dicPhraseIS{verða sér úti um e‑ð}} \dicDirectTranslationCS{sehnat (co), získat (co)};  \dicIdiom{verða}[við]{ \dicPhraseIS{verða við e‑u}} \dicDirectTranslationCS{vyhovět (čemu), vyjít vstříc (čemu)}
\dicEntry[verðandi] \dicTerm{verð··and|i\smash{\textsuperscript{1}}} \dicIPA{{v}{\textepsilon}{r}{ð}{a}{n}{\textsubring{d}}{\textsci}} \dicPos{f}[3] \dicFlx{(‑i)}[3] \dicDirectTranslationCS{nastávání, uskutečňování}
\dicEntry[Verðandi] \dicTerm{Verð··and|i\smash{\textsuperscript{2}}} \dicIPA{{v}{\textepsilon}{r}{ð}{a}{n}{\textsubring{d}}{\textsci}} \dicPos{f}[3] \dicFlx{(‑i)}[4] \dicFlx{prop} \dicFieldCat{myt.} \dicDirectTranslationCS{Verdandi} \dicIndirectTranslationCS{(sudička)}
\dicEntry[verðandi] \dicTerm{verð··andi\smash{\textsuperscript{3}}} \dicIPA{{v}{\textepsilon}{r}{\texttheta}{a}{n}{\textsubring{d}}{\textsci}} \dicPos{adj}[13] \dicFlx{indecl}[1] \dicDirectTranslationCS{nastávající, nadcházející} \dicExampleIS{verðandi móðir} \dicExampleCS{nastávající matka}
\dicEntry[verðbólga] \dicTerm{verð··bólg|a} \dicIPA{{v}{\textepsilon}{r}{ð}{\textsubring{b}}{ou}{l}{\r{g}}{a}} \dicPos{f}[1] \dicFlx{(‑u)}[5] \dicFieldCat{ekon.} \dicDirectTranslationCS{inflace}
\dicEntry[verðbólgustig] \dicTerm{verð·bólgu··stig} \dicIPA{{v}{\textepsilon}{r}{ð}{\textsubring{b}}{ou}{l}{\r{g}}{\textscy}{s}{\textsubring{d}}{\textsci}{x}} \dicPos{n}[2] \dicFlx{(‑s, ‑)}[5] \dicFieldCat{ekon.} \dicDirectTranslationCS{míra inflace}
\dicEntry[verðbréf] \dicTerm{verð··bréf} \dicIPA{{v}{\textepsilon}{r}{ð}{\textsubring{b}}{r}{j}{\textepsilon}{f}} \dicPos{n}[2] \dicFlx{(‑s, ‑)}[5] \dicFieldCat{ekon.} \dicDirectTranslationCS{cenný papír, cenina}
\dicEntry[verðbréfamarkaður] \dicTerm{verð·bréfa··mark·að|ur} \dicIPA{{v}\-{\textepsilon}\-{r}\-{ð}\-{\textsubring{b}}\-{r}\-{j}\-{\textepsilon}\-{v}\-{a}\-{m}\-{a}\-{\textsubring{r}}\-{\r{g}}\-{a}\-{ð}\-{\textscy}\-{\textsubring{r}}\-} \dicPos{m}[10] \dicFlx{(‑ar\,/\addthin ‑s, ‑ir)}[31] \dicFieldCat{ekon.} \dicDirectTranslationCS{trh\,/\addthin burza cenných papírů}
\dicEntry[verðbréfamiðlari] \textls[15]{\dicTerm{verð·bréfa··miðl·ar|i} \dicIPA{{v}\-{\textepsilon}\-{r}\-{ð}\-{\textsubring{b}}\-{r}\-{j}\-{\textepsilon}\-{v}\-{a}\-{m}\-{\textsci}\-{ð}\-{l}\-{a}\-{r}\-{\textsci}\-} \dicPos{m}[1] \dicFlx{(‑a, ‑ar)}[13]} \dicFieldCat{ekon.} \dicDirectTranslationCS{(burzovní) makléř(ka)}
\dicEntry[verðbréfasjóður] \dicTerm{verð·bréfa··sjóð|ur} \dicIPA{{v}{\textepsilon}{r}{ð}{\textsubring{b}}{r}{j}{\textepsilon}{v}{a}{s}{j}{ou}{ð}{\textscy}{\textsubring{r}}} \dicPos{m}[9] \dicFlx{(‑s, ‑ir)}[6] \dicFieldCat{ekon.} \dicDirectTranslationCS{investiční\,/\addthin podílový fond}
\dicEntry[verðbætur] \dicTerm{verð··bætur} \dicIPA{{v}{\textepsilon}{r}{ð}{\textsubring{b}}{a}{i}{\textsubring{d}}{\textscy}{\textsubring{r}}} \dicPos{f}[8] \dicFlx{pl}[6] \dicFieldCat{ekon.} \dicDirectTranslationCS{indexace}
\dicEntry[verðfall] \dicTerm{verð··fall} \dicIPA{{v}{\textepsilon}{r}{ð}{f}{a}{\textsubring{d}}{\textsubring{l}}} \dicPos{n}[2] \dicFlx{(‑s)}[2] \dicDirectTranslationCS{cenový pád, pád ceny}
\dicEntry[verðflokkur] \dicTerm{verð··flokk|ur} \dicIPA{{v}{\textepsilon}{r}{ð}{f}{l}{\textopeno}{h}{\r{g}}{\textscy}{\textsubring{r}}} \dicPos{m}[6] \dicFlx{(‑s, ‑ar)}[8] \dicDirectTranslationCS{cenová kategorie\,/\addthin skupina}
\dicEntry[verðgildi] \dicTerm{verð··gildi} \dicIPA{{v}{\textepsilon}{r}{ð}{\r{\textObardotlessj}}{\textsci}{l}{\textsubring{d}}{\textsci}} \dicPos{n}[2] \dicFlx{(‑s)}[20] \dicDirectTranslationCS{(peněžní) hodnota}
\dicEntry[verðhjöðnun] \dicTerm{verð··hjöðn|un} \dicIPA{{v}{\textepsilon}{r}{\texttheta}{\c{c}}{\oe}{ð}{n}{\textscy}{\textsubring{n}}} \dicPos{f}[7] \dicFlx{(‑unar)}[12] \dicFieldCat{ekon.} \dicDirectTranslationCS{deflace}
\dicEntry[verðhrun] \dicTerm{verð··hrun} \dicIPA{{v}{\textepsilon}{r}{\texttheta}{\textsubring{r}}{\textscy}{\textsubring{n}}} \dicPos{n}[2] \dicFlx{(‑s, ‑)}[5] \dicDirectTranslationCS{cenový propad\,/\addthin kolaps, propad\,/\addthin kolaps ceny}
\dicEntry[verðhækkun] \dicTerm{verð··hækk|un} \dicIPA{{v}{\textepsilon}{r}{\texttheta}{h}{a}{i}{h}{\r{g}}{\textscy}{\textsubring{n}}} \dicPos{f}[7] \dicFlx{(‑unar, ‑anir)}[8] \dicDirectTranslationCS{zdražení, nárůst\,/\addthin zvýšení ceny}
\dicEntry[verði] \dicTerm{verði} \dicIPA{{v}{\textepsilon}{r}{ð}{\textsci}} \dicPos{v} \dicFlx{con pf sg 1 pers} \textbf{1.} \dicLink{verja\smash{\textsuperscript{2}}}  \textbf{2.} \dicLink{verja\smash{\textsuperscript{3}}}
\dicEntry[verðir] \dicTerm{verðir} \dicIPA{{v}{\textepsilon}{r}{ð}{\textsci}{\textsubring{r}}} \dicPos{m} \dicFlx{pl nom} \dicLink{vörður}
\dicEntry[verðjöfnun] \dicTerm{verð··jöfn|un} \dicIPA{{v}{\textepsilon}{r}{ð}{j}{\oe}{\textsubring{b}}{n}{\textscy}{\textsubring{n}}} \dicPos{f}[7] \dicFlx{(‑unar)}[12] \dicFieldCat{ekon.} \dicDirectTranslationCS{cenové vyrovnání}
\dicEntry[verðlag] \dicTerm{verð··lag} \dicIPA{{v}{\textepsilon}{r}{ð}{l}{a}{x}} \dicPos{n}[2] \dicFlx{(‑s)}[2] \dicFieldCat{ekon.} \dicDirectTranslationCS{cena, cenová úroveň\,/\addthin hladina} \dicExampleIS{lágt verðlag} \dicExampleCS{nízká cenová hladina}
\dicEntry[verðlagning] \dicTerm{verð··lag·ning} \dicIPA{{v}{\textepsilon}{r}{ð}{l}{a}{\r{g}}{n}{i}{\ng}{\r{g}}} \dicPos{f}[4] \dicFlx{(‑ar)}[7] \dicDirectTranslationCS{stanovení\,/\addthin určení ceny}
\dicEntry[verðlagseftirlit] \dicTerm{verð·lags··eftir·lit} \dicIPA{{v}{\textepsilon}{r}{ð}{l}{a}{x}{s}{\textepsilon}{f}{\textsubring{d}}{\textsci}{r}{l}{\textsci}{\textsubring{d}}} \dicPos{n}[2] \dicFlx{(‑s)}[2] \dicFieldCat{ekon.} \dicDirectTranslationCS{kontrola cenové hladiny}

\dicEntry[verðlagsgrundvöllur] \textls[15]{\dicTerm{verð·lags··grund·|völlur} \dicIPA{{v}\-{\textepsilon}\-{r}\-{ð}\-{l}\-{a}\-{x}\-{s}\-{\r{g}}\-{r}\-{\textscy}\-{n}\-{\textsubring{d}}\-{v}\-{\oe}\-{\textsubring{d}}\-{l}\-{\textscy}\-{\textsubring{r}}} \dicPos{m}[11] \dicFlx{(‑vall\-ar)}[6]} \dicFieldCat{ekon.} \dicDirectTranslationCS{cenový základ}

\dicEntry[verðlagsráð] \dicTerm{verð·lags··ráð} \dicIPA{{v}{\textepsilon}{r}{ð}{l}{a}{x}{s}{r}{au}{\texttheta}} \dicPos{n}[2] \dicFlx{(‑s, ‑)}[5] \dicFieldCat{ekon.} \dicDirectTranslationCS{cenová komise}
\dicEntry[verðlagsvísitala] \dicTerm{verð·lags··vísi·|tala} \dicIPA{{v}\-{\textepsilon}\-{r}\-{ð}\-{l}\-{a}\-{x}\-{s}\-{v}\-{i}\-{s}\-{\textsci}\-{t\smash{\textsuperscript{h}}}\-{a}\-{l}\-{a}\-} \dicPos{f}[1] \dicFlx{(‑tölu, ‑tölur)}[14] \dicFieldCat{ekon.} \dicDirectTranslationCS{index cenové úrovně}
\dicEntry[verðlaun] \dicTerm{verð··laun} \dicIPA{{v}{\textepsilon}{r}{ð}{l}{\oe i}{\textsubring{n}}} \dicPos{n}[2] \dicFlx{pl}[1] \textbf{1.} \dicDirectTranslationCS{cena, ocenění (v~soutěži ap.)}  \textbf{2.} \dicFieldCat{sport.} \dicDirectTranslationCS{medaile};  \dicPhraseIS{veita e‑m verðlaun} \dicDirectTranslationCS{udělit (komu) cenu\,/\addthin medaili}
\dicEntry[verðlauna] \dicTerm{verð··laun|a} \dicIPA{{v}{\textepsilon}{r}{ð}{l}{\oe i}{n}{a}} \dicPos{v}[1] \dicFlx{(‑aði)}[1] \dicFlx{acc} \dicDirectTranslationCS{udělit cenu\,/\addthin medaili} \dicExampleIS{verðlauna hana fyrir árangurinn} \dicExampleCS{udělit jí cenu za výsledek}
\dicEntry[verðlaunagripur] \dicTerm{verð·launa··grip|ur} \dicIPA{{v}{\textepsilon}{r}{ð}{l}{\oe i}{n}{a}{\r{g}}{r}{\textsci}{\textsubring{b}}{\textscy}{\textsubring{r}}} \dicPos{m}[9] \dicFlx{(‑s, ‑ir)}[8] \dicDirectTranslationCS{medaile, trofej, cena (pro vítěze)}
\dicEntry[verðlaunahafi] \dicTerm{verð·launa··haf|i} \dicIPA{{v}{\textepsilon}{r}{ð}{l}{\oe i}{n}{a}{h}{a}{v}{\textsci}} \dicPos{m}[1] \dicFlx{(‑a, ‑ar)}[8] \dicDirectTranslationCS{medailista, medailistka, držitel(ka) medaile\,/\addthin trofeje}
\dicEntry[verðlaus] \dicTerm{verð··laus} \dicIPA{{v}{\textepsilon}{r}{ð}{l}{\oe i}{s}} \dicPos{adj}[5]\dicFlx{}[-1] \dicDirectTranslationCS{bezcenný, nehodnotný} \dicIndirectTranslationCS{(který nemá cenu)} \dicExampleIS{Krónan er orðin verðlaus.} \dicExampleCS{Koruna je již bezcenná.}
\dicEntry[verðleggja] \dicTerm{verð··|leggja} \dicIPA{{v}{\textepsilon}{r}{ð}{l}{\textepsilon}{\r{\textObardotlessj}}{a}} \dicPos{v}[4] \dicFlx{(‑legg, ‑lagði, ‑lögðum, ‑legði, ‑lagt)}[22] \dicFlx{acc} \dicDirectTranslationCS{zhodnotit, ocenit, stanovit\,/\addthin určit cenu} \dicExampleIS{verðleggja húsið} \dicExampleCS{stanovit cenu domu}
\dicEntry[verðleiki] \dicTerm{verð··leik|i} \dicIPA{{v}{\textepsilon}{r}{ð}{l}{ei}{\r{\textObardotlessj}}{\textsci}} \dicPos{m}[1] \dicFlx{(‑a, ‑ar)}[1] \textbf{1.} \dicSynonym{kostur} \dicDirectTranslationCS{klad, přednost}  \textbf{2.} \dicDirectTranslationCS{zásluha};  \dicPhraseIS{að verðleikum} \dicFlx{adv} \dicDirectTranslationCS{po zásluze, zaslouženě}
\dicEntry[verðlisti] \dicTerm{verð··list|i} \dicIPA{{v}{\textepsilon}{r}{ð}{l}{\textsci}{s}{\textsubring{d}}{\textsci}} \dicPos{m}[1] \dicFlx{(‑a, ‑ar)}[1] \dicDirectTranslationCS{ceník}
\dicEntry[verðlækkun] \dicTerm{verð··lækk|un} \dicIPA{{v}{\textepsilon}{r}{ð}{l}{a}{i}{h}{\r{g}}{\textscy}{\textsubring{n}}} \dicPos{f}[7] \dicFlx{(‑unar, ‑anir)}[8] \dicDirectTranslationCS{snížení\,/\addthin pokles ceny}
\dicEntry[verðmerkja] \dicTerm{verð··merk|ja} \dicIPA{{v}{\textepsilon}{r}{ð}{m}{\textepsilon}{\textsubring{r}}{\r{\textObardotlessj}}{a}} \dicPos{v}[2] \dicFlx{(‑ti, ‑t)}[28] \dicFlx{acc} \dicDirectTranslationCS{označit cenou}
\dicEntry[verðmiði] \dicTerm{verð··mið|i} \dicIPA{{v}{\textepsilon}{r}{ð}{m}{\textsci}{ð}{\textsci}} \dicPos{m}[1] \dicFlx{(‑a, ‑ar)}[1] \dicDirectTranslationCS{cenovka, cenový štítek}
\dicEntry[verðmunur] \dicTerm{verð··mun|ur} \dicIPA{{v}{\textepsilon}{r}{ð}{m}{\textscy}{n}{\textscy}{\textsubring{r}}} \dicPos{m}[10] \dicFlx{(‑ar)}[12] \dicDirectTranslationCS{cenový rozdíl}
\dicEntry[verðmyndun] \dicTerm{verð··mynd|un} \dicIPA{{v}{\textepsilon}{r}{ð}{m}{\textsci}{n}{\textsubring{d}}{\textscy}{\textsubring{n}}} \dicPos{f}[7] \dicFlx{(‑unar)}[9] \dicFieldCat{ekon.} \dicDirectTranslationCS{stanovení ceny}
\dicEntry[verðmætamat] \dicTerm{verð·mæta··mat} \dicIPA{{v}{\textepsilon}{r}{ð}{m}{a}{i}{\textsubring{d}}{a}{m}{a}{\textsubring{d}}} \dicPos{n}[2] \dicFlx{(‑s)}[2] \dicDirectTranslationCS{hodnoty, hodnotový žebříček}
\dicEntry[verðmæti] \dicTerm{verð··mæti} \dicsymFrequent\  \dicIPA{{v}{\textepsilon}{r}{ð}{m}{a}{i}{\textsubring{d}}{\textsci}} \dicPos{n}[2] \dicFlx{(‑s, ‑)}[14] \textbf{1.} \dicDirectTranslationCS{hodnota (vnitřní ap.)} \dicExampleIS{andlegt verðmæti} \dicExampleCS{duševní hodnota}  \textbf{2.} \dicFieldCat{ekon.} \dicSynonym{gildi} \dicDirectTranslationCS{hodnota} \dicExampleIS{auka verðmæti vörunnar} \dicExampleCS{zvýšit hodnotu výrobku}
\begin{xtolerant}{}{1pt}
\dicEntry[verðmætur] \dicTerm{verð··mætur}\addthinS\dicIPA{{v}{\textepsilon}{r}{ð}{m}{a}{i}{\textsubring{d}}{\textscy}{\textsubring{r}}}\addthinS\dicPos{adj}[1]\dicFlx{}[-1]\addthinS\dicDirectTranslationCS{cenný, hodnotný} \dicExampleIS{verðmætir gripir} \dicExampleCS{cenné věci}
\end{xtolerant}
\dicEntry[verðskrá] \dicTerm{verð··skrá} \dicIPA{{v}{\textepsilon}{r}{ð}{s}{\r{g}}{r}{au}} \dicPos{f}[4] \dicFlx{(‑r\,/\addthin ‑ar, ‑r)}[21] \dicDirectTranslationCS{tarif, sazba}
\dicEntry[verðskulda] \dicTerm{verð··skuld|a} \dicIPA{{v}{\textepsilon}{r}{ð}{s}{\r{g}}{\textscy}{l}{\textsubring{d}}{a}} \dicPos{v}[1] \dicFlx{(‑aði)}[1] \dicFlx{acc} \dicDirectTranslationCS{zasloužit si, zasluhovat si} \dicExampleIS{verðskulda refsingu} \dicExampleCS{zasloužit si potrestání}
\dicEntry[verðskuldaður] \dicTerm{verð··skuld·|aður} \dicIPA{{v}{\textepsilon}{r}{ð}{s}{\r{g}}{\textscy}{l}{\textsubring{d}}{a}{ð}{\textscy}{\textsubring{r}}} \dicPos{adj}[3] \dicFlx{(f ‑uð)}[3] \dicDirectTranslationCS{zasloužený} \dicExampleIS{verðskuldaður sigur} \dicExampleCS{zasloužené vítězství}
\dicEntry[verðstöðvun] \dicTerm{verð··stöðv|un} \dicIPA{{v}{\textepsilon}{r}{ð}{s}{\textsubring{d}}{\oe}{ð}{v}{\textscy}{\textsubring{n}}} \dicPos{f}[7] \dicFlx{(‑unar, ‑anir)}[8] \dicFieldCat{ekon.} \dicDirectTranslationCS{zmrazení cen}
\dicEntry[verðsveifla] \dicTerm{verð··sveifl|a} \dicIPA{{v}{\textepsilon}{r}{ð}{s}{v}{ei}{\textsubring{b}}{l}{a}} \dicPos{f}[1] \dicFlx{(‑u, ‑ur)}[19] \dicDirectTranslationCS{fluktuace\,/\addthin kolísání ceny}
\dicEntry[verðtilboð] \dicTerm{verð··til·boð} \dicIPA{{v}{\textepsilon}{r}{ð}{t\smash{\textsuperscript{h}}}{\textsci}{l}{\textsubring{b}}{\textopeno}{\texttheta}} \dicPos{n}[2] \dicFlx{(‑s, ‑)}[5] \dicDirectTranslationCS{cenová nabídka}
\dicEntry[verðtollur] \dicTerm{verð··toll|ur} \dicIPA{{v}{\textepsilon}{r}{ð}{t\smash{\textsuperscript{h}}}{\textopeno}{\textsubring{d}}{l}{\textscy}{\textsubring{r}}} \dicPos{m}[6] \dicFlx{(‑s, ‑ar)}[4] \dicFieldCat{ekon.} \dicDirectTranslationCS{hodnotové clo}
\dicEntry[verðtryggður] \dicTerm{verð··tryggður} \dicIPA{{v}{\textepsilon}{r}{ð}{t\smash{\textsuperscript{h}}}{r}{\textsci}{\textbabygamma}{ð}{\textscy}{\textsubring{r}}} \dicPos{adj}[2]\dicFlx{}[-1] \dicFieldCat{ekon.} \dicDirectTranslationCS{indexovaný}
\dicEntry[verðtrygging] \dicTerm{verð··trygg·ing} \dicIPA{{v}{\textepsilon}{r}{ð}{t\smash{\textsuperscript{h}}}{r}{\textsci}{\r{\textObardotlessj}}{i}{\ng}{\r{g}}} \dicPos{f}[4] \dicFlx{(‑ar, ‑ar)}[5] \dicFieldCat{ekon.} \dicDirectTranslationCS{indexování}
\dicEntry[verðtryggja] \dicTerm{verð··trygg|ja} \dicIPA{{v}{\textepsilon}{r}{ð}{t\smash{\textsuperscript{h}}}{r}{\textsci}{\r{\textObardotlessj}}{a}} \dicPos{v}[2] \dicFlx{(‑ði, ‑t)}[89] \dicFlx{acc} \dicFieldCat{ekon.} \dicDirectTranslationCS{indexovat}
\dicEntry[verðugur] \dicTerm{verð··ugur} \dicIPA{{v}{\textepsilon}{r}{ð}{\textscy}{\textbabygamma}{\textscy}{\textsubring{r}}} \dicPos{adj}[1]\dicFlx{}[-8] \textbf{1.} \dicSynonym*{sem verðskuldar e‑ð} \dicDirectTranslationCS{zasloužilý, hodný (čeho)};  \dicPhraseIS{vera verðugur e‑s} \dicDirectTranslationCS{zasloužit si (co), být hoden (čeho)}  \textbf{2.} \dicSynonym*{mikils virði} \dicDirectTranslationCS{záslužný, chvályhodný}
\dicEntry[verður] \dicTerm{verður\smash{\textsuperscript{1}}} \dicIPA{{v}{\textepsilon}{r}{ð}{\textscy}{\textsubring{r}}} \dicPos{adj}[2]\dicFlx{}[-1] \textbf{1.} \dicSynonym*{sem er e‑s virði} \dicDirectTranslationCS{hodnotný, cenný} \dicExampleIS{lítils verður} \dicExampleCS{nemající velkou hodnotu}  \textbf{2.} \dicSynonym{maklegur} \dicDirectTranslationCS{zasluhující si, hodný (čeho)} \dicExampleIS{Verður er verkamaðurinn launanna.} \dicExampleCS{Dělník si zasluhuje plat.};  \dicIdiom{verður}{ \dicPhraseIS{e‑að er mikils um vert}} \dicDirectTranslationCS{(co) je důležité, (co) má značný význam}; { \dicPhraseIS{það er ekki vert að (gera e‑ð)}} \dicDirectTranslationCS{nevyplatí se (udělat (co)), nestojí to za to (udělat (co))}
\dicEntry[verður] \dicTerm{verður\smash{\textsuperscript{2}}} \dicIPA{{v}{\textepsilon}{r}{ð}{\textscy}{\textsubring{r}}} \dicPos{v} \dicFlx{ind praes sg 3 pers} \dicLink{verða}
\dicEntry[verðvísitala] \dicTerm{verð··vísi·|tala} \dicIPA{{v}{\textepsilon}{r}{ð}{v}{i}{s}{\textsci}{t\smash{\textsuperscript{h}}}{a}{l}{a}} \dicPos{f}[1] \dicFlx{(‑tölu, ‑tölur)}[14] \dicFieldCat{ekon.} \dicDirectTranslationCS{cenový ukazatel}
\dicEntry[vergangur] \dicTerm{ver··gang|ur} \dicIPA{{v}{\textepsilon}{r}{\r{g}}{au}{\ng}{\r{g}}{\textscy}{\textsubring{r}}} \dicPos{m}[6] \dicFlx{(‑s)}[7] \dicSynonym{betl} \dicDirectTranslationCS{žebrání, žebrota} \dicExampleIS{lenda á vergangi} \dicExampleCS{stát se žebrákem}
\dicEntry[vergirni] \dicTerm{ver··girn|i} \dicIPA{{v}{\textepsilon}{r}{\r{\textObardotlessj}}{\textsci}{r}{\textsubring{d}}{n}{\textsci}} \dicPos{f}[3] \dicFlx{(‑i)}[3] \dicFieldCat{med.} \dicDirectTranslationCS{nymfomanie}
\dicEntry[vergjarn] \dicTerm{ver··|gjarn} \dicIPA{{v}{\textepsilon}{r}{\r{\textObardotlessj}}{a}{r}{\textsubring{d}}{\textsubring{n}}} \dicPos{adj}[5] \dicFlx{(f ‑gjörn)}[6] \dicFieldCat{med. } \dicDirectTranslationCS{nymfomanský} \dicExampleIS{vergjörn kona} \dicExampleCS{nymfomanka}
\dicEntry[vergur] \dicTerm{vergur} \dicIPA{{v}{\textepsilon}{r}{\r{g}}{\textscy}{\textsubring{r}}} \dicPos{adj}[1]\dicFlx{}[-1] \dicFieldCat{ekon.} \dicSynonym{brúttó} \dicDirectTranslationCS{hrubý, brutto} \dicExampleIS{vergar tekjur} \dicExampleCS{hrubé příjmy}
\dicEntry[verið] \dicTerm{verið} \dicIPA{{v}{\textepsilon}{\textlengthmark}{r}{\textsci}{\texttheta}} \dicPos{v} \dicFlx{supin} \dicLink{vera\smash{\textsuperscript{2}}}
\dicEntry[verja] \dicTerm{verj|a\smash{\textsuperscript{1}}} \dicIPA{{v}{\textepsilon}{r}{j}{a}} \dicPos{f}[1] \dicFlx{(‑u, ‑ur)}[7] \textbf{1.} \dicSynonym{hlíf} \dicDirectTranslationCS{ochrana, chránění}  \textbf{2.} \dicSynonym{getnaðarvörn} \dicDirectTranslationCS{ochrana, kondom, prezervativ}
\dicEntry[verja] \dicTerm{verja\smash{\textsuperscript{2}}} \dicsymFrequent\  \dicIPA{{v}{\textepsilon}{r}{j}{a}} \dicPos{v}[4] \dicFlx{(ver, varði, vörðum, verði, varið)}[13] \dicFlx{acc} \textbf{1.} \dicSynonym{vernda} \dicDirectTranslationCS{(o)chránit, ochraňovat, (u)bránit};  \dicPhraseIS{verja e‑ð fyrir e‑u} \dicDirectTranslationCS{chránit (co) před (čím)} \dicExampleIS{verja landið fyrir innrásarhernum} \dicExampleCS{bránit zemi před agresory};  \dicPhraseIS{verja e‑n falli} \dicDirectTranslationCS{zachránit (koho) před pádem};  \dicPhraseIS{verja sig} \dicDirectTranslationCS{chránit\,/\addthin bránit se}  \textbf{2.} \dicFieldCat{práv.} \dicDirectTranslationCS{hájit, obhajovat};  \dicPhraseIS{verja mál} \dicDirectTranslationCS{hájit případ (u~soudu ap.)} \dicExampleIS{verja málið fyrir rétti} \dicExampleCS{hájit případ před soudem};  \dicIdiom{verjast}{ \dicPhraseIS{verjast e‑u\,/\addthin e‑m}} \dicFlx{refl} \dicDirectTranslationCS{bránit se (čemu\,/\addthin komu)} \dicExampleIS{verjast árás} \dicExampleCS{bránit se útoku}; { \dicPhraseIS{geta ekki varist e‑u}} \dicFlx{refl} \dicDirectTranslationCS{nemoct se ubránit (čemu) (smíchu ap.)}
\dicEntry[verja] \dicTerm{verja\smash{\textsuperscript{3}}} \dicIPA{{v}{\textepsilon}{r}{j}{a}} \dicPos{v}[4] \dicFlx{(ver, varði, vörðum, verði, varið)}[13] \dicFlx{dat} \dicSynonym{nýta} \dicDirectTranslationCS{užít, užívat, spotřebovat, spotřebovávat, utratit, utrácet} \dicExampleIS{verja peningum sínum vel} \dicExampleCS{užívat své peníze dobře};  \dicIdiom{verja}[í]{ \dicPhraseIS{verja tíma í e‑ð}} \dicDirectTranslationCS{strávit čas (čím)} \dicExampleIS{Ég hef varið miklum tíma í að lesa bókina.} \dicExampleCS{Strávil jsem hodně času čtením knížky.};  \dicIdiom{verja}[til]{ \dicPhraseIS{verja peningum sínum til e‑s}} \dicDirectTranslationCS{použít své peníze na (co), utratit své peníze za (co)}
\dicEntry[verjandi] \dicTerm{verj··|andi\smash{\textsuperscript{1}}} \dicIPA{{v}{\textepsilon}{r}{j}{a}{n}{\textsubring{d}}{\textsci}} \dicPos{m}[2] \dicFlx{(‑anda, ‑endur)}[1] \textbf{1.} \dicSynonym*{sá sem ver} \dicDirectTranslationCS{ochránce, ochránkyně, obránce, obránkyně}  \textbf{2.} \dicFieldCat{práv.} \dicDirectTranslationCS{obhájce, obhájkyně}
\dicEntry[verjandi] \dicTerm{verj··andi\smash{\textsuperscript{2}}} \dicIPA{{v}{\textepsilon}{r}{j}{a}{n}{\textsubring{d}}{\textsci}} \dicPos{adj}[13] \dicFlx{indecl}[1] \dicDirectTranslationCS{omluvitelný, ospravedlnitelný} \dicIndirectTranslationCS{(často se záporem)} \dicExampleIS{e‑að er ekki verjandi} \dicExampleCS{(co) je neomluvitelné}
\dicEntry[verk] \dicTerm{verk} \dicsymFrequent\  \dicIPA{{v}{\textepsilon}{\textsubring{r}}{\r{g}}} \dicPos{n}[2] \dicFlx{(‑s, ‑)}[5] \textbf{1.} \dicSynonym{vinna\smash{\textsuperscript{1}}} \dicDirectTranslationCS{práce, činnost};  \dicPhraseIS{koma sér að verki} \dicDirectTranslationCS{dát\,/\addthin pustit se do práce};  \dicPhraseIS{kunna vel til verka} \dicDirectTranslationCS{být zručný ve své práci};  \dicPhraseIS{segja e‑m fyrir verkum} \dicDirectTranslationCS{dávat (komu) instrukce};  \dicPhraseIS{standa e‑n að verki} \dicDirectTranslationCS{přistihnout (koho) při činu, nachytat (koho)};  \dicPhraseIS{vera í verki með e‑m} \dicDirectTranslationCS{spolupracovat s~(kým)}  \textbf{2.} \dicSynonym{listaverk} \dicDirectTranslationCS{dílo, výtvor, práce}  \textbf{3.} \dicSynonym{gangverk} \dicDirectTranslationCS{strojek, mechanismus};  \dicIdiom{verk}{ \dicPhraseIS{e‑m fellur aldrei verk úr hendi}} \dicFlx{impers} \dicDirectTranslationCS{(kdo) nikdy nezahálí, (kdo) nikdy nesedí s~rukama v~klíně}; { \dicPhraseIS{e‑að gerir það að verkum að (e‑að gerist)}} \dicDirectTranslationCS{(co) způsobuje, že se ((co) děje)}; { \dicPhraseIS{e‑m verður ekkert úr verki}} \dicFlx{impers} \dicDirectTranslationCS{(kdo) nic neudělá, (komu) se nepodaří nic udělat};  \dicPhraseIS{ganga hreint til verks} \dicLangCat{přen.} \dicDirectTranslationCS{jít rovnou na věc};  \dicPhraseIS{taka viljann fyrir verkið} \dicLangCat{přen.} \dicDirectTranslationCS{dobrá vůle se počítá} \dicIndirectTranslationCS{(ne výsledek)}
\dicEntry[verka] \dicTerm{verk|a} \dicsymFrequent\  \dicIPA{{v}{\textepsilon}{\textsubring{r}}{\r{g}}{a}} \dicPos{v}[1] \dicFlx{(‑aði)}[1] \dicFlx{acc} \textbf{1.} \dicDirectTranslationCS{zpracovat, opracovat (úlovek ap.)}  \textbf{2.} \dicSynonym{orka\smash{\textsuperscript{2}}} \dicDirectTranslationCS{(za)působit, účinkovat};  \dicPhraseIS{verka á e‑n} \dicDirectTranslationCS{působit na (koho), účinkovat na (koho), mít efekt\,/\addthin účinek na (koho)} \dicExampleIS{Þetta verkar undarlega á mig.} \dicExampleCS{To na mě divně působí.}
\dicEntry[verkafólk] \dicTerm{verka··fólk} \dicIPA{{v}{\textepsilon}{\textsubring{r}}{\r{g}}{a}{f}{ou}{\textsubring{l}}{\r{g}}} \dicPos{n}[2] \dicFlx{(‑s)}[2] \dicSynonym{starfsfólk} \dicDirectTranslationCS{pracující, pracovníci, dělníci}
\dicEntry[verkahringur] \dicTerm{verka··hring|ur} \dicIPA{{v}{\textepsilon}{\textsubring{r}}{\r{g}}{a}{\textsubring{r}}{i}{\ng}{\r{g}}{\textscy}{\textsubring{r}}} \dicPos{m}[6] \dicFlx{(‑s)}[17] \dicSynonym*{starfssvæði} \dicDirectTranslationCS{působnost, (pracovní) náplň};  \dicPhraseIS{vera í verkahring e‑rs} \dicDirectTranslationCS{být v~(čí) působnosti} \dicExampleIS{Það er ekki í mínum verkahring að kalla hópinn saman.} \dicExampleCS{Není v~mojí pracovní náplni svolávat skupinu.}
\dicEntry[verkakona] \dicTerm{verka··kon|a} \dicIPA{{v}{\textepsilon}{\textsubring{r}}{\r{g}}{a}{k\smash{\textsuperscript{h}}}{\textopeno}{n}{a}} \dicPos{f}[1] \dicFlx{(‑u, ‑ur)}[27] \dicDirectTranslationCS{dělnice, pracovnice}
\dicEntry[verkakvennafélag] \dicTerm{verka·kvenna··fé·|lag} \dicIPA{{v}\-{\textepsilon}\-{\textsubring{r}}\-{\r{g}}\-{a}\-{k\smash{\textsuperscript{h}}}\-{v}\-{\textepsilon}\-{n}\-{a}\-{f}\-{j}\-{\textepsilon}\-{l}\-{a}\-{x}\-} \dicPos{n}[2] \dicFlx{(‑lags, ‑lög)}[8] \dicFieldCat{pol.} \dicDirectTranslationCS{ženská odborová organizace}
\dicEntry[verkalaun] \dicTerm{verka··laun} \dicIPA{{v}{\textepsilon}{\textsubring{r}}{\r{g}}{a}{l}{\oe i}{\textsubring{n}}} \dicPos{n}[2] \dicFlx{pl}[1] \dicSynonym{kaup} \dicDirectTranslationCS{výdělek, mzda}
\dicEntry[verkalýðsbarátta] \dicTerm{verka·lýðs··bar·átt|a} \dicIPA{{v}\-{\textepsilon}\-{\textsubring{r}}\-{\r{g}}\-{a}\-{l}\-{i}\-{ð}\-{s}\-{\textsubring{b}}\-{a}\-{r}\-{au}\-{h}\-{\textsubring{d}}\-{a}\-} \dicPos{f}[1] \dicFlx{(‑u)}[5] \dicDirectTranslationCS{dělnický boj, boj odborových svazů}
\dicEntry[verkalýðsfélag] \dicTerm{verka·lýðs··fé·|lag} \dicIPA{{v}{\textepsilon}{\textsubring{r}}{\r{g}}{a}{l}{i}{ð}{s}{f}{j}{\textepsilon}{l}{a}{x}} \dicPos{n}[2] \dicFlx{(‑lags, ‑lög)}[8] \dicDirectTranslationCS{odborová organizace}
\dicEntry[verkalýðsflokkur] \dicTerm{verka·lýðs··flokk|ur} \dicIPA{{v}{\textepsilon}{\textsubring{r}}{\r{g}}{a}{l}{i}{ð}{s}{f}{l}{\textopeno}{h}{\r{g}}{\textscy}{\textsubring{r}}} \dicPos{m}[6] \dicFlx{(‑s, ‑ar)}[8] \dicFieldCat{pol.} \dicDirectTranslationCS{dělnická strana}
\dicEntry[verkalýðshreyfing] \dicTerm{verka·lýðs··hreyf·ing} \dicIPA{{v}\-{\textepsilon}\-{\textsubring{r}}\-{\r{g}}\-{a}\-{l}\-{i}\-{ð}\-{s}\-{\textsubring{r}}\-{ei}\-{v}\-{i}\-{\ng}\-{\r{g}}\-} \dicPos{f}[4] \dicFlx{(‑ar, ‑ar)}[5] \dicDirectTranslationCS{dělnické\,/\addthin odborové hnutí}
\dicEntry[verkalýðssamtök] \dicTerm{verka·lýðs··sam·tök} \dicIPA{{v}{\textepsilon}{\textsubring{r}}{\r{g}}{a}{l}{i}{ð}{s}{a}{m}{t\smash{\textsuperscript{h}}}{\oe}{\r{g}}} \dicPos{n}[2] \dicFlx{pl}[9] \dicDirectTranslationCS{odbory, odborová organizace, odborové sdružení}
\dicEntry[verkalýður] \dicTerm{verka··lýð|ur} \dicIPA{{v}{\textepsilon}{\textsubring{r}}{\r{g}}{a}{l}{i}{ð}{\textscy}{\textsubring{r}}} \dicPos{m}[6] \dicFlx{(‑s)}[17] \dicDirectTranslationCS{dělnická\,/\addthin pracující třída, dělnictvo, proletariát}
\dicEntry[verkamaður] \dicTerm{verka··|maður} \dicsymFrequent\  \dicIPA{{v}{\textepsilon}{\textsubring{r}}{\r{g}}{a}{m}{a}{ð}{\textscy}{\textsubring{r}}} \dicPos{m}[13] \dicFlx{(‑manns, ‑menn)}[2] \dicDirectTranslationCS{dělník, dělnice, pracovník, pracovnice} \dicExampleIS{útlendir verkamenn} \dicExampleCS{zahraniční dělníci}
\dicEntry[verkamannaflokkur] \dicTerm{verka·manna··flokk|ur} \dicIPA{{v}\-{\textepsilon}\-{\textsubring{r}}\-{\r{g}}\-{a}\-{m}\-{a}\-{n}\-{a}\-{f}\-{l}\-{\textopeno}\-{h}\-{\r{g}}\-{\textscy}\-{\textsubring{r}}\-} \dicPos{m}[6] \dicFlx{(‑s, ‑ar)}[8] \dicFieldCat{pol.} \dicDirectTranslationCS{dělnická strana}
\dicEntry[verkamannasamband] \textls[27]{\dicTerm{verka·manna··sam·|band} \dicIPA{{v}\-{\textepsilon}\-{\textsubring{r}}\-{\r{g}}\-{a}\-{m}\-{a}\-{n}\-{a}\-{s}\-{a}\-{m}\-{\textsubring{b}}\-{a}\-{n}\-{\textsubring{d}}}} \dicPos{n}[2] \dicFlx{(‑bands, ‑bönd)}[8] \dicDirectTranslationCS{odborový svaz}
\dicEntry[verkamannavinna] \dicTerm{verka·manna··vinn|a} \dicIPA{{v}{\textepsilon}{\textsubring{r}}{\r{g}}{a}{m}{a}{n}{a}{v}{\textsci}{n}{a}} \dicPos{f}[1] \dicFlx{(‑u)}[5] \dicDirectTranslationCS{nekvalifikovaná práce}
\dicEntry[verkaskipting] \dicTerm{verka··skipt·ing} \dicIPA{{v}{\textepsilon}{\textsubring{r}}{\r{g}}{a}{s}{\r{\textObardotlessj}}{\textsci}{f}{\textsubring{d}}{i}{\ng}{\r{g}}} \dicPos{f}[4] \dicFlx{(‑ar)}[7] \dicDirectTranslationCS{dělba práce}
\dicEntry[verkefnabók] \dicTerm{verk·efna··|bók} \dicIPA{{v}{\textepsilon}{\textsubring{r}}{\r{g}}{\textepsilon}{\textsubring{b}}{n}{a}{\textsubring{b}}{ou}{\r{g}}} \dicPos{f}[8] \dicFlx{(‑bókar, ‑bækur)}[5] \dicDirectTranslationCS{cvičebnice}
\dicEntry[verkefni] \dicTerm{verk··efni} \dicsymFrequent\  \dicIPA{{v}{\textepsilon}{\textsubring{r}}{\r{g}}{\textepsilon}{\textsubring{b}}{n}{\textsci}} \dicPos{n}[2] \dicFlx{(‑s, ‑)}[14] \dicDirectTranslationCS{úkol, úloha, zadání, projekt} \dicExampleIS{gefa honum verkefni} \dicExampleCS{dát mu úkol}
\dicEntry[verkefnisstjóri] \dicTerm{verk·efnis··stjór|i} \dicIPA{{v}{\textepsilon}{\textsubring{r}}{\r{g}}{\textepsilon}{\textsubring{b}}{n}{\textsci}{s}{\textsubring{d}}{j}{ou}{r}{\textsci}} \dicPos{m}[1] \dicFlx{(‑a, ‑ar)}[1] \dicDirectTranslationCS{vedoucí projektu, projektový manažer, projektová manažerka}
\dicEntry[verkeyðandi] \dicTerm{verk··eyð·andi} \dicIPA{{v}{\textepsilon}{\textsubring{r}}{\r{g}}{ei}{ð}{a}{n}{\textsubring{d}}{\textsci}} \dicPos{adj}[13] \dicFlx{indecl}[1] \dicDirectTranslationCS{analgetický, tlumící bolest} \dicExampleIS{verkeyðandi dropar} \dicExampleCS{kapičky tlumící bolest}
\dicEntry[verkfall] \dicTerm{verk··|fall} \dicIPA{{v}{\textepsilon}{\textsubring{r}}{\r{g}}{f}{a}{\textsubring{d}}{\textsubring{l}}} \dicPos{n}[2] \dicFlx{(‑falls, ‑föll)}[8] \dicDirectTranslationCS{stávka} \dicExampleIS{hóta verkfalli} \dicExampleCS{hrozit stávkou}
\dicEntry[verkfallsbrjótur] \dicTerm{verk·falls··brjót|ur} \dicIPA{{v}{\textepsilon}{\textsubring{r}}{\r{g}}{f}{a}{l}{s}{\textsubring{b}}{r}{j}{ou}{\textsubring{d}}{\textscy}{\textsubring{r}}} \dicPos{m}[6] \dicFlx{(‑s, ‑ar)}[10] \dicDirectTranslationCS{stávkokaz(ka)}
\dicEntry[verkfallsréttur] \dicTerm{verk·falls··rétt|ur} \dicIPA{{v}{\textepsilon}{\textsubring{r}}{\r{g}}{f}{a}{l}{s}{r}{j}{\textepsilon}{h}{\textsubring{d}}{\textscy}{\textsubring{r}}} \dicPos{m}[10] \dicFlx{(‑ar)}[7] \dicFieldCat{práv.} \dicDirectTranslationCS{právo na stávku}
\dicEntry[verkfræði] \dicTerm{verk··fræð|i} \dicIPA{{v}{\textepsilon}{\textsubring{r}}{\r{g}}{f}{r}{a}{i}{ð}{\textsci}} \dicPos{f}[3] \dicFlx{(‑i)}[3] \dicDirectTranslationCS{inženýrství}
\dicEntry[verkfræðideild] \dicTerm{verk·fræði··deild} \dicIPA{{v}{\textepsilon}{\textsubring{r}}{\r{g}}{f}{r}{a}{i}{ð}{\textsci}{\textsubring{d}}{ei}{l}{\textsubring{d}}} \dicPos{f}[7] \dicFlx{(‑ar, ‑ir)}[1] \dicFieldCat{škol.} \dicDirectTranslationCS{fakulta strojního inženýrství}
\dicEntry[verkfræðingur] \dicTerm{verk·fræð··ing|ur} \dicIPA{{v}{\textepsilon}{\textsubring{r}}{\r{g}}{f}{r}{a}{i}{ð}{i}{\ng}{\r{g}}{\textscy}{\textsubring{r}}} \dicPos{m}[6] \dicFlx{(‑s, ‑ar)}[8] \dicDirectTranslationCS{inženýr(ka)}
\dicEntry[verkfærakassi] \dicTerm{verk·færa··kass|i} \dicIPA{{v}{\textepsilon}{\textsubring{r}}{\r{g}}{f}{a}{i}{r}{a}{k\smash{\textsuperscript{h}}}{a}{s}{\textsci}} \dicPos{m}[1] \dicFlx{(‑a, ‑ar)}[8] \dicDirectTranslationCS{(přenosná) skříňka na nářadí,  bedna s~nářadím}
\dicEntry[verkfæri] \dicTerm{verk··færi} \dicIPA{{v}{\textepsilon}{\textsubring{r}}{\r{g}}{f}{a}{i}{r}{\textsci}} \dicPos{n}[2] \dicFlx{(‑s, ‑)}[14] \dicSynonym{áhald} \dicDirectTranslationCS{nářadí, nástroj} \dicExampleIS{kunna að halda rétt á verkfæri} \dicExampleCS{umět správně zacházet s~nástrojem}
\dicEntry[verkfærisfall] \dicTerm{verk·færis··|fall} \dicIPA{{v}{\textepsilon}{\textsubring{r}}{\r{g}}{f}{a}{i}{r}{\textsci}{s}{f}{a}{\textsubring{d}}{\textsubring{l}}} \dicPos{n}[2] \dicFlx{(‑falls, ‑föll)}[8] \dicFieldCat{jaz.} \dicSynonym{tækisfall} \dicDirectTranslationCS{7. pád, instrumentál}
\dicEntry[verkja] \dicTerm{verkj|a} \dicIPA{{v}{\textepsilon}{\textsubring{r}}{\r{\textObardotlessj}}{a}} \dicPos{v}[1] \dicFlx{(‑aði)}[66] \dicFlx{impers} \dicPhraseIS{e‑n verkjar í e‑ð} \dicDirectTranslationCS{(koho co) bolí}
\dicEntry[verkjakast] \dicTerm{verkja··|kast} \dicIPA{{v}{\textepsilon}{\textsubring{r}}{\r{\textObardotlessj}}{a}{k\smash{\textsuperscript{h}}}{a}{s}{\textsubring{d}}} \dicPos{n}[2] \dicFlx{(‑kasts, ‑köst)}[8] \dicDirectTranslationCS{záchvat bolesti}
\dicEntry[verkjalaus] \dicTerm{verkja··laus} \dicIPA{{v}{\textepsilon}{\textsubring{r}}{\r{\textObardotlessj}}{a}{l}{\oe i}{s}} \dicPos{adj}[5]\dicFlx{}[-1] \dicDirectTranslationCS{bezbolestný, (jsoucí) bez bolesti}
\dicEntry[verkjalyf] \dicTerm{verkja··lyf} \dicIPA{{v}{\textepsilon}{\textsubring{r}}{\r{\textObardotlessj}}{a}{l}{\textsci}{f}} \dicPos{n}[2] \dicFlx{(‑s, ‑)}[13] \dicDirectTranslationCS{lék proti bolesti, analgetikum}
\dicEntry[verkjastillandi] \dicTerm{verkja··still·andi} \dicIPA{{v}{\textepsilon}{\textsubring{r}}{\r{\textObardotlessj}}{a}{s}{\textsubring{d}}{\textsci}{\textsubring{d}}{l}{a}{n}{\textsubring{d}}{\textsci}} \dicPos{adj}[13] \dicFlx{indecl}[1] \dicDirectTranslationCS{tišící\,/\addthin mírnící bolest, analgetický} \dicExampleIS{verkjastillandi lyf} \dicExampleCS{lék tišící bolest}
\dicEntry[verkjatafla] \dicTerm{verkja··|tafla} \dicIPA{{v}{\textepsilon}{\textsubring{r}}{\r{\textObardotlessj}}{a}{t\smash{\textsuperscript{h}}}{a}{\textsubring{b}}{l}{a}} \dicPos{f}[1] \dicFlx{(‑töflu, ‑töflur)}[20] \dicDirectTranslationCS{prášek\,/\addthin pilulka proti bolesti}
\dicEntry[verkkunnátta] \dicTerm{verk··kunn·átt|a} \dicIPA{{v}{\textepsilon}{\textsubring{r}}{\r{g}}{k\smash{\textsuperscript{h}}}{\textscy}{n}{au}{h}{\textsubring{d}}{a}} \dicPos{f}[1] \dicFlx{(‑u)}[5] \dicDirectTranslationCS{(pracovní) dovednost\,/\addthin znalost, know-how}
\dicEntry[verklaginn] \dicTerm{verk··laginn} \dicIPA{{v}{\textepsilon}{\textsubring{r}}{\r{g}}{l}{a}{i}{j}{\textsci}{\textsubring{n}}} \dicPos{adj}[6]\dicFlx{}[-3] \dicDirectTranslationCS{zručný, dovedný}
\dicEntry[verklagni] \dicTerm{verk··lagn|i} \dicIPA{{v}{\textepsilon}{\textsubring{r}}{\r{g}}{l}{a}{\r{g}}{n}{\textsci}} \dicPos{f}[3] \dicFlx{(‑i)}[3] \dicDirectTranslationCS{zručnost, dovednost}
\dicEntry[verklaus] \dicTerm{verk··laus} \dicIPA{{v}{\textepsilon}{\textsubring{r}}{\r{g}}{l}{\oe i}{s}} \dicPos{adj}[5]\dicFlx{}[-1] \dicDirectTranslationCS{netečný, pasivní}
\dicEntry[verklegur] \dicTerm{verk··legur} \dicsymFrequent\  \dicIPA{{v}{\textepsilon}{\textsubring{r}}{\r{g}}{l}{\textepsilon}{\textbabygamma}{\textscy}{\textsubring{r}}} \dicPos{adj}[1]\dicFlx{}[-8] \textbf{1.} \dicSynonym*{sem felst í verki} \dicDirectTranslationCS{praktický, pracovní} \dicExampleIS{bóklegt og verklegt námskeið} \dicExampleCS{teoretický a~praktický kurz}  \textbf{2.} \dicSynonym{sterklegur} \dicDirectTranslationCS{bytelný, solidní}
\dicEntry[verklýsing] \dicTerm{verk··lýs·ing} \dicIPA{{v}{\textepsilon}{\textsubring{r}}{\r{g}}{l}{i}{s}{i}{\ng}{\r{g}}} \dicPos{f}[4] \dicFlx{(‑ar, ‑ar)}[5] \dicDirectTranslationCS{popis práce}
\dicEntry[verkmenning] \dicTerm{verk··menn·ing} \dicIPA{{v}{\textepsilon}{\textsubring{r}}{\r{g}}{m}{\textepsilon}{n}{i}{\ng}{\r{g}}} \dicPos{f}[4] \dicFlx{(‑ar)}[7] \dicDirectTranslationCS{pracovní znalosti\,/\addthin dovednosti}
\dicEntry[verkmenntaskóli] \dicTerm{verk··mennta·skól|i} \dicIPA{{v}{\textepsilon}{\textsubring{r}}{\r{g}}{m}{\textepsilon}{\textsubring{n}}{\textsubring{d}}{a}{s}{\r{g}}{ou}{l}{\textsci}} \dicPos{m}[1] \dicFlx{(‑a, ‑ar)}[1] \dicFieldCat{škol.} \dicDirectTranslationCS{(střední odborné) učiliště}
\dicEntry[verknaður] \dicTerm{verk··nað|ur} \dicIPA{{v}{\textepsilon}{\textsubring{r}}{\r{g}}{n}{a}{ð}{\textscy}{\textsubring{r}}} \dicPos{m}[10] \dicFlx{(‑ar, ‑ir)}[8] \dicSynonym{athöfn} \dicDirectTranslationCS{skutek, čin, akt} \dicExampleIS{illur verknaður} \dicExampleCS{špatný čin}
\dicEntry[verknám] \dicTerm{verk··nám} \dicIPA{{v}{\textepsilon}{\textsubring{r}}{\r{g}}{n}{au}{\textsubring{m}}} \dicPos{n}[2] \dicFlx{(‑s)}[2] \dicFieldCat{škol.} \dicDirectTranslationCS{praktické studium}
\dicEntry[verksamningur] \dicTerm{verk··samn·ing|ur} \dicIPA{{v}{\textepsilon}{\textsubring{r}}{\r{g}}{s}{a}{m}{n}{i}{\ng}{\r{g}}{\textscy}{\textsubring{r}}} \dicPos{m}[6] \dicFlx{(‑s, ‑ar)}[8] \dicDirectTranslationCS{pracovní smlouva, kontrakt}
\dicEntry[verksmiðja] \dicTerm{verk··smiðj|a} \dicsymFrequent\  \dicIPA{{v}{\textepsilon}{\textsubring{r}}{\r{g}}{s}{m}{\textsci}{ð}{j}{a}} \dicPos{f}[1] \dicFlx{(‑u, ‑ur)}[7] \dicDirectTranslationCS{továrna, výrobna, závod} \dicExampleIS{Verksmiðjan framleiðir bíla.} \dicExampleCS{Továrna vyrábí auta.}
\dicEntry[verksmiðjumengun] \dicTerm{verk·smiðju··meng|un} \dicIPA{{v}{\textepsilon}{\textsubring{r}}{\r{g}}{s}{m}{\textsci}{ð}{j}{\textscy}{m}{ei}{\ng}{\r{g}}{\textscy}{\textsubring{n}}} \dicPos{f}[7] \dicFlx{(‑unar)}[9] \dicDirectTranslationCS{tovární znečištění\,/\addthin znečišťování}
\dicEntry[verkstjóri] \dicTerm{verk··stjór|i} \dicsymFrequent\  \dicIPA{{v}{\textepsilon}{\textsubring{r}}{\r{g}}{s}{\textsubring{d}}{j}{ou}{r}{\textsci}} \dicPos{m}[1] \dicFlx{(‑a, ‑ar)}[1] \dicDirectTranslationCS{vedoucí práce, mistr(ová) (vedoucí dílny ap.)} \dicExampleIS{starfa sem verkstjóri} \dicExampleCS{pracovat jako vedoucí práce}
\dicEntry[verkstjórn] \dicTerm{verk··stjórn} \dicIPA{{v}{\textepsilon}{\textsubring{r}}{\r{g}}{s}{\textsubring{d}}{j}{ou}{r}{\textsubring{d}}{\textsubring{n}}} \dicPos{f}[7] \dicFlx{(‑ar)}[3] \dicDirectTranslationCS{řízení práce}
\dicEntry[verkstæði] \dicTerm{verk··stæði} \dicIPA{{v}{\textepsilon}{\textsubring{r}}{\r{g}}{s}{\textsubring{d}}{a}{i}{ð}{\textsci}} \dicPos{n}[2] \dicFlx{(‑s, ‑)}[14] \dicDirectTranslationCS{dílna, opravna, servisní závod (automobilů ap.)}
\begin{xtolerant}{}{1pt}
\dicEntry[verksummerki] \dicTerm{verks··um·merki}\addthinS\dicIPA{{v}{\textepsilon}{\textsubring{r}}{x}{s}{\textscy}{m}{\textlengthmark}{\textepsilon}{\textsubring{r}}{\r{\textObardotlessj}}{\textsci}}\addthinS\dicPos{n}[2]\addthinS\dicFlx{pl}[17] \dicDirectTranslationCS{(viditelné) stopy\,/\addthin známky} \dicExampleIS{verksummerki eftir jarðskjálftann} \dicExampleCS{stopy po zemětřesení}
\end{xtolerant}
\dicEntry[verksvið] \dicTerm{verk··svið} \dicIPA{{v}{\textepsilon}{\textsubring{r}}{\r{g}}{s}{v}{\textsci}{\texttheta}} \dicPos{n}[2] \dicFlx{(‑s, ‑)}[5] \dicSynonym{verkahringur} \dicDirectTranslationCS{pole působnosti, kompetence}
\dicEntry[verktaki] \dicTerm{verk··tak|i} \dicIPA{{v}{\textepsilon}{\textsubring{r}}{\r{g}}{t\smash{\textsuperscript{h}}}{a}{\r{\textObardotlessj}}{\textsci}} \dicPos{m}[1] \dicFlx{(‑a, ‑ar)}[8] \dicDirectTranslationCS{podnikatel(ka), živnostník, živnostnice}
\dicEntry[verktækni] \dicTerm{verk··tækn|i} \dicIPA{{v}{\textepsilon}{\textsubring{r}}{\r{g}}{t\smash{\textsuperscript{h}}}{a}{i}{h}{\r{g}}{n}{\textsci}} \dicPos{f}[3] \dicFlx{(‑i)}[3] \dicDirectTranslationCS{pracovní postup, technologie}
\dicEntry[verkun] \dicTerm{verk|un} \dicIPA{{v}{\textepsilon}{\textsubring{r}}{\r{g}}{\textscy}{\textsubring{n}}} \dicPos{f}[7] \dicFlx{(‑unar, ‑anir)}[8] \textbf{1.} \dicSynonym{áhrif} \dicDirectTranslationCS{účinek, účinkování}  \textbf{2.} \dicSynonym{vinnsla} \dicDirectTranslationCS{zpracování, úprava (ryb ap.)} \dicExampleIS{verkun fisks} \dicExampleCS{úprava ryby}
\dicEntry[verkur] \dicTerm{verk|ur} \dicsymFrequent\  \dicIPA{{v}{\textepsilon}{\textsubring{r}}{\r{g}}{\textscy}{\textsubring{r}}} \dicPos{m}[10] \dicFlx{(‑jar, ‑ir)}[23] \dicSynonym{sársauki} \dicDirectTranslationCS{bolest} \dicExampleIS{Það dregur úr verknum.} \dicExampleCS{Bolest ustupuje.}
\dicEntry[verkþáttur] \dicTerm{verk··|þáttur} \dicIPA{{v}{\textepsilon}{\textsubring{r}}{\r{g}}{\texttheta}{au}{h}{\textsubring{d}}{\textscy}{\textsubring{r}}} \dicPos{m}[12] \dicFlx{(‑þáttar, ‑þættir)}[7] \dicDirectTranslationCS{pracovní etapa\,/\addthin fáze}
\dicEntry[verma] \dicTerm{verm|a} \dicIPA{{v}{\textepsilon}{r}{m}{a}} \dicPos{v}[2] \dicFlx{(‑di, ‑t)}[141] \dicFlx{acc} \dicSynonym{hlýja\smash{\textsuperscript{1}}} \dicDirectTranslationCS{zahřát, zahřívat, vyhřát, vyhřívat} \dicExampleIS{verma sig} \dicExampleCS{zahřát se}
\dicEntry[vermir] \dicTerm{verm|ir} \dicIPA{{v}{\textepsilon}{r}{m}{\textsci}{\textsubring{r}}} \dicPos{m}[7] \dicFlx{(‑is)}[2] \dicSynonym{varmi} \dicDirectTranslationCS{zahřátí, zahřívání};  \dicPhraseIS{e‑að er skammgóður vermir} \dicLangCat{přen.} \dicDirectTranslationCS{(co) je jen chvilková úleva}
\dicEntry[vermireitur] \dicTerm{vermi··reit|ur} \dicIPA{{v}{\textepsilon}{r}{m}{\textsci}{r}{ei}{\textsubring{d}}{\textscy}{\textsubring{r}}} \dicPos{m}[9] \dicFlx{(‑s, ‑ir)}[8] \dicDirectTranslationCS{pařeniště}
\dicEntry[vernd] \dicTerm{vernd} \dicsymFrequent\  \dicIPA{{v}{\textepsilon}{r}{n}{\textsubring{d}}} \dicPos{f}[4] \dicFlx{(‑ar)}[3] \dicSynonym{vörn} \dicDirectTranslationCS{ochrana, chránění, obrana} \dicExampleIS{vernd gegn hættu} \dicExampleCS{ochrana před nebezpečím}
\dicEntry[vernda] \dicTerm{vernd|a} \dicsymFrequent\  \dicIPA{{v}{\textepsilon}{r}{n}{\textsubring{d}}{a}} \dicPos{v}[1] \dicFlx{(‑aði)}[1] \dicFlx{acc} \dicSynonym{skýla\smash{\textsuperscript{2}}} \dicDirectTranslationCS{(o)chránit, ochraňovat, (u)bránit} \dicExampleIS{vernda náttúruna gegn mengun} \dicExampleCS{chránit přírodu před znečištěním}
\dicEntry[verndaður] \dicTerm{vernd··|aður} \dicIPA{{v}{\textepsilon}{r}{n}{\textsubring{d}}{a}{ð}{\textscy}{\textsubring{r}}} \dicPos{adj}[3] \dicFlx{(f ‑uð)}[3] \dicPhraseIS{verndaður vinnustaður} \dicDirectTranslationCS{chráněná dílna}
\dicEntry[verndardýrlingur] \dicTerm{verndar··dýr·ling|ur} \dicIPA{{v}{\textepsilon}{r}{n}{\textsubring{d}}{a}{r}{\textsubring{d}}{i}{r}{l}{i}{\ng}{\r{g}}{\textscy}{\textsubring{r}}} \dicPos{m}[6] \dicFlx{(‑s, ‑ar)}[8] \dicDirectTranslationCS{patron(ka), svatý ochránce, svatá ochránkyně}
\dicEntry[verndarengill] \dicTerm{verndar··eng|ill} \dicIPA{{v}{\textepsilon}{r}{n}{\textsubring{d}}{a}{r}{ei}{\textltailn}{\r{\textObardotlessj}}{\textsci}{\textsubring{d}}{\textsubring{l}}} \dicPos{m}[6] \dicFlx{(‑ils, ‑lar)}[35] \textbf{1.} \dicSynonym{engill} \dicDirectTranslationCS{anděl strážný}  \textbf{2.} \dicLangCat{přen.} \dicDirectTranslationCS{anděl strážný}
\dicEntry[verndargripur] \dicTerm{verndar··grip|ur} \dicIPA{{v}{\textepsilon}{r}{n}{\textsubring{d}}{a}{r}{\r{g}}{r}{\textsci}{\textsubring{b}}{\textscy}{\textsubring{r}}} \dicPos{m}[9] \dicFlx{(‑s, ‑ir)}[8] \dicDirectTranslationCS{talisman, amulet}
\begin{xtolerant}{}{1pt}
\dicEntry[verndarhönd] \dicTerm{verndar··|hönd}\addthinS\dicIPA{{v}{\textepsilon}{r}{n}{\textsubring{d}}{a}{\textsubring{r}}{h}{\oe}{n}{\textsubring{d}}}\addthinS\dicPos{f}[8]\addthinS\dicFlx{(‑handar)}[13] \dicPhraseIS{halda verndarhendi yfir e‑m} \dicLangCat{přen.} \dicDirectTranslationCS{držet nad (kým) ochrannou ruku}
\end{xtolerant}
\dicEntry[verndari] \dicTerm{vernd··ar|i} \dicIPA{{v}{\textepsilon}{r}{n}{\textsubring{d}}{a}{r}{\textsci}} \dicPos{m}[1] \dicFlx{(‑a, ‑ar)}[13] \textbf{1.} \dicDirectTranslationCS{patron(ka), ochránce, ochránkyně} \dicIndirectTranslationCS{(světec jako ochránce lidí ap.)}  \textbf{2.} \dicDirectTranslationCS{patron(ka), mecenáš(ka)} \dicIndirectTranslationCS{(člověk poskytující hmotnou, mravní pomoc)}
\dicEntry[verndarlaus] \dicTerm{verndar··laus} \dicIPA{{v}{\textepsilon}{r}{n}{\textsubring{d}}{a}{r}{l}{\oe i}{s}} \dicPos{adj}[5]\dicFlx{}[-1] \dicDirectTranslationCS{nechráněný, neochraňovaný}
\dicEntry[verndarstefna] \dicTerm{verndar··stefn|a} \dicIPA{{v}{\textepsilon}{r}{n}{\textsubring{d}}{a}{\textsubring{r}}{s}{\textsubring{d}}{\textepsilon}{\textsubring{b}}{n}{a}} \dicPos{f}[1] \dicFlx{(‑u)}[5] \dicFieldCat{ekon.} \dicDirectTranslationCS{protekcionismus, ochranářství}
\dicEntry[verndarsvæði] \dicTerm{verndar··svæði} \dicIPA{{v}{\textepsilon}{r}{n}{\textsubring{d}}{a}{\textsubring{r}}{s}{v}{a}{i}{ð}{\textsci}} \dicPos{n}[2] \dicFlx{(‑s, ‑)}[14] \textbf{1.} \dicSynonym*{náttúruverndarsvæði} \dicDirectTranslationCS{(přírodní) rezervace}  \textbf{2.} \dicDirectTranslationCS{rezervace} \dicIndirectTranslationCS{(území vyhrazené domorodcům)}  \textbf{3.} \dicDirectTranslationCS{protektorát}
\begin{xtolerant}{}{1pt}
\dicEntry[verndartollur] \dicTerm{verndar··toll|ur}\addthinS\dicIPA{{v}{\textepsilon}{r}{n}{\textsubring{d}}{a}{\textsubring{r}}{t\smash{\textsuperscript{h}}}{\textopeno}{\textsubring{d}}{l}{\textscy}{\textsubring{r}}}\addthinS\dicPos{m}[6]\addthinS\dicFlx{(‑s, ‑ar)}[4]\addthinS\dicFieldCat{ekon.} \dicDirectTranslationCS{ochran\-né clo}
\end{xtolerant}
\dicEntry[verndarvængur] \dicTerm{verndar··væng|ur} \dicIPA{{v}{\textepsilon}{r}{n}{\textsubring{d}}{a}{r}{v}{a}{i}{\ng}{\r{g}}{\textscy}{\textsubring{r}}} \dicPos{m}[9] \dicFlx{(‑s\,/\addthin ‑jar, ‑ir)}[26] \dicPhraseIS{taka e‑n undir sinn verndarvæng} \dicLangCat{přen.} \dicDirectTranslationCS{vzít (koho) pod svá ochranná křídla}
\dicEntry[verndun] \dicTerm{vernd|un} \dicIPA{{v}{\textepsilon}{r}{n}{\textsubring{d}}{\textscy}{\textsubring{n}}} \dicPos{f}[7] \dicFlx{(‑unar)}[9] \dicSynonym{friðun} \dicDirectTranslationCS{ochrana, ochraňování, chránění, prezervace} \dicExampleIS{verndun fugla} \dicExampleCS{ochrana ptáků}
\dicEntry[verp] \dicTerm{verp} \dicIPA{{v}{\textepsilon}{\textsubring{r}}{\textsubring{b}}} \dicPos{v} \dicFlx{ind praes sg 1 pers} \dicLink{verpa\smash{\textsuperscript{2}}}
\dicEntry[verpa] \dicTerm{verp|a\smash{\textsuperscript{1}}}\dicTerm{, verpa\smash{\textsuperscript{2}}} \dicIPA{{v}\-{\textepsilon}\-{\textsubring{r}}\-{\textsubring{b}}\-{a}\-} \dicPos{v}[2] \dicFlx{(‑ti, ‑t)}[76] \dicFlx{dat} \dicDirectTranslationCS{(na)klást, snést, snášet} \dicIndirectTranslationCS{(vajíčka)} \dicExampleIS{verpa eggjum} \dicExampleCS{snášet vajíčka}
\dicEntry[verpa] \dicTerm{verpa\smash{\textsuperscript{2}}} \dicIPA{{v}{\textepsilon}{\textsubring{r}}{\textsubring{b}}{a}} \dicPos{v}[6] \dicFlx{(verp, varp, urpum, yrpi, orpið)}[34] \dicFlx{dat} \dicLink{verpa\smash{\textsuperscript{1}}}
\dicEntry[verr] \dicTerm{verr} \dicIPA{{v}{\textepsilon}{r}{\textlengthmark}} \dicPos{adv} \dicFlx{comp (pos illa, sup verst)} \dicDirectTranslationCS{hůř(e)};  \dicPhraseIS{því er nú verr (og miður)} \dicFlx{adv} \dicDirectTranslationCS{bohužel}
\dicEntry[verri] \dicTerm{verri} \dicsymFrequent\  \dicIPA{{v}{\textepsilon}{r}{\textlengthmark}{\textsci}} \dicPos{adj} \dicFlx{comp} \textbf{1.} \dicLink{vondur}  \textbf{2.} \dicLink{illur}  \textbf{3.} \dicLink{slæmur}
\dicEntry[vers] \dicTerm{vers} \dicIPA{{v}{\textepsilon}{\textsubring{r}}{s}} \dicPos{n}[2] \dicFlx{(‑, ‑)}[22] \textbf{1.} \dicFieldCat{náb.} \dicDirectTranslationCS{verš (v~bibli ap.)};  \dicPhraseIS{syngja sitt síðasta vers} \dicLangCat{přen.} \dicDirectTranslationCS{zemřít, skonat}  \textbf{2.} \dicFieldCat{lit.} \dicSynonym{vísuorð} \dicDirectTranslationCS{verš}
\dicEntry[versla] \dicTerm{versl|a} \dicIPA{{v}{\textepsilon}{\textsubring{r}}{s}{\textsubring{d}}{l}{a}} \dicPos{v}[1] \dicFlx{(‑aði)}[34] \textbf{1.} \dicSynonym*{skipta við} \dicDirectTranslationCS{obchodovat};  \dicPhraseIS{versla með e‑ð} \dicDirectTranslationCS{obchodovat s~(čím)} \dicExampleIS{Kaupmaðurinn verslar með fatnað.} \dicExampleCS{Obchodník obchoduje s~oděvy.}  \textbf{2.} \dicSynonym*{gera innkaup} \dicDirectTranslationCS{nakupovat};  \dicPhraseIS{versla við e‑n} \dicDirectTranslationCS{nakupovat u~(koho)}
\dicEntry[verslanamiðstöð] \textls[15]{\dicTerm{verslana··mið·stöð} \dicIPA{{v}{\textepsilon}{\textsubring{r}}{s}{\textsubring{d}}{l}{a}{n}{a}{m}{\textsci}{ð}{s}{\textsubring{d}}{\oe}{\texttheta}} \dicPos{f}[6] \dicFlx{(‑var, ‑var)}[1]} \dicLink{verslunarmiðstöð}
\dicEntry[Versló] \dicTerm{Versló} \dicIPA{{v}{\textepsilon}{\textsubring{r}}{s}{\textsubring{d}}{l}{ou}} \dicPos{m}[15] \dicFlx{indecl}[2] \dicDirectTranslationCS{Versló} \dicIndirectTranslationCS{(zkratka názvu Islandské obchodní akademie (celý název -- Verslunarskóli Íslands))}
\dicEntry[verslun] \dicTerm{versl|un} \dicsymFrequent\  \dicIPA{{v}{\textepsilon}{\textsubring{r}}{s}{\textsubring{d}}{l}{\textscy}{\textsubring{n}}} \dicPos{f}[7] \dicFlx{(‑unar, ‑anir)}[8] \textbf{1.} \dicSynonym{búð} \dicDirectTranslationCS{obchod, prodejna} \dicExampleIS{Í næstu götu eru verslanir.} \dicExampleCS{V~další ulici jsou obchody.}  \textbf{2.} \dicSynonym{viðskipti} \dicDirectTranslationCS{obchod, obchodování}
\dicEntry[verslunarborg] \dicTerm{verslunar··borg} \dicIPA{{v}{\textepsilon}{\textsubring{r}}{s}{\textsubring{d}}{l}{\textscy}{n}{a}{r}{\textsubring{b}}{\textopeno}{r}{\r{g}}} \dicPos{f}[7] \dicFlx{(‑ar, ‑ir)}[1] \dicDirectTranslationCS{obchodní centrum, centrum obchodu} \dicIndirectTranslationCS{(o~městě)}
\dicEntry[verslunarbréf] \dicTerm{verslunar··bréf} \dicIPA{{v}{\textepsilon}{\textsubring{r}}{s}{\textsubring{d}}{l}{\textscy}{n}{a}{r}{\textsubring{b}}{r}{j}{\textepsilon}{f}} \dicPos{n}[2] \dicFlx{(‑s, ‑)}[5] \dicDirectTranslationCS{obchodní dopis}
\dicEntry[verslunareinokun] \dicTerm{verslunar··ein·ok|un} \dicIPA{{v}{\textepsilon}{\textsubring{r}}{s}{\textsubring{d}}{l}{\textscy}{n}{a}{r}{ei}{n}{\textopeno}{\r{g}}{\textscy}{\textsubring{n}}} \dicPos{f}[7] \dicFlx{(‑unar)}[9] \dicDirectTranslationCS{obchodní monopol}
\dicEntry[verslunarfloti] \dicTerm{verslunar··flot|i} \dicIPA{{v}{\textepsilon}{\textsubring{r}}{s}{\textsubring{d}}{l}{\textscy}{n}{a}{\textsubring{r}}{f}{l}{\textopeno}{\textsubring{d}}{\textsci}} \dicPos{m}[1] \dicFlx{(‑a)}[3] \dicFieldCat{nám.} \dicDirectTranslationCS{obchodní loďstvo}
\dicEntry[verslunarfrelsi] \dicTerm{verslunar··frelsi} \dicIPA{{v}{\textepsilon}{\textsubring{r}}{s}{\textsubring{d}}{l}{\textscy}{n}{a}{\textsubring{r}}{f}{r}{\textepsilon}{l}{s}{\textsci}} \dicPos{n}[2] \dicFlx{(‑s)}[20] \dicDirectTranslationCS{svoboda obchodování}
\dicEntry[verslunarfulltrúi] \dicTerm{verslunar··full·trú|i} \dicIPA{{v}\-{\textepsilon}\-{\textsubring{r}}\-{s}\-{\textsubring{d}}\-{l}\-{\textscy}\-{n}\-{a}\-{\textsubring{r}}\-{f}\-{\textscy}\-{\textsubring{d}}\-{\textsubring{l}}\-{t\smash{\textsuperscript{h}}}\-{r}\-{u}\-{\textsci}\-} \dicPos{m}[1] \dicFlx{(‑a, ‑ar)}[1] \dicDirectTranslationCS{obchodní atašé}
\dicEntry[verslunarfyrirtæki] \dicTerm{verslunar··fyrir·tæki} \dicIPA{{v}\-{\textepsilon}\-{\textsubring{r}}\-{s}\-{\textsubring{d}}\-{l}\-{\textscy}\-{n}\-{a}\-{\textsubring{r}}\-{f}\-{\textsci}\-{r}\-{\textsci}\-{\textsubring{r}}\-{t\smash{\textsuperscript{h}}}\-{a}\-{i}\-{\r{\textObardotlessj}}\-{\textsci}\-} \dicPos{n}[2] \dicFlx{(‑s, ‑)}[16] \dicDirectTranslationCS{obchodní společnost}
\dicEntry[verslunargata] \dicTerm{verslunar··|gata} \dicIPA{{v}{\textepsilon}{\textsubring{r}}{s}{\textsubring{d}}{l}{\textscy}{n}{a}{r}{\r{g}}{a}{\textsubring{d}}{a}} \dicPos{f}[1] \dicFlx{(‑götu, ‑götur)}[14] \dicDirectTranslationCS{obchodní\,/\addthin nákupní třída}
\dicEntry[verslunarhalli] \dicTerm{verslunar··hall|i} \dicIPA{{v}{\textepsilon}{\textsubring{r}}{s}{\textsubring{d}}{l}{\textscy}{n}{a}{\textsubring{r}}{h}{a}{\textsubring{d}}{l}{\textsci}} \dicPos{m}[1] \dicFlx{(‑a)}[3] \dicFieldCat{ekon.} \dicSynonym{viðskiptahalli} \dicDirectTranslationCS{deficit obchodní bilance}
\dicEntry[verslunarháskóli] \dicTerm{verslunar··há·skól|i} \dicIPA{{v}{\textepsilon}{\textsubring{r}}{s}{\textsubring{d}}{l}{\textscy}{n}{a}{\textsubring{r}}{h}{au}{s}{\r{g}}{ou}{l}{\textsci}} \dicPos{m}[1] \dicFlx{(‑a, ‑ar)}[1] \dicFieldCat{škol.} \dicDirectTranslationCS{vysoká škola obchodní}
\dicEntry[verslunarhús] \dicTerm{verslunar··hús} \dicIPA{{v}{\textepsilon}{\textsubring{r}}{s}{\textsubring{d}}{l}{\textscy}{n}{a}{\textsubring{r}}{h}{u}{s}} \dicPos{n}[2] \dicFlx{(‑s, ‑)}[5] \dicDirectTranslationCS{obchodní\,/\addthin nákupní dům}
\begin{xtolerant}{}{1pt}
\dicEntry[verslunarjöfnuður] \dicTerm{verslunar··|jöfn·uður} \dicIPA{{v}\-{\textepsilon}\-{\textsubring{r}}\-{s}\-{\textsubring{d}}\-{l}\-{\textscy}\-{n}\-{a}\-{r}\-{j}\-{\oe}\-{\textsubring{b}}\-{n}\-{\textscy}\-{ð}\-{\textscy}\-{\textsubring{r}}}\addthinS\dicPos{m}[10] \dicFlx{(‑jafn\-að\-ar\,/\addthin ‑jöfnuðar)}[43] \dicFieldCat{ekon.} \dicDirectTranslationCS{obchodní bilance}
\end{xtolerant}
\begin{xtolerant}{}{8pt}
\dicEntry[verslunarkeðja] \dicTerm{verslunar··keðj|a} \dicIPA{{v}{\textepsilon}{\textsubring{r}}{s}{\textsubring{d}}{l}{\textscy}{n}{a}{\textsubring{r}}{c\smash{\textsuperscript{h}}}{\textepsilon}{ð}{j}{a}} \dicPos{f}[1] \dicFlx{(‑u, ‑ur)}[7] \dicDirectTranslationCS{obchod\-ní řetězec}
\end{xtolerant}
\dicEntry[verslunarleyfi] \dicTerm{verslunar··leyfi} \dicIPA{{v}{\textepsilon}{\textsubring{r}}{s}{\textsubring{d}}{l}{\textscy}{n}{a}{r}{l}{ei}{v}{\textsci}} \dicPos{n}[2] \dicFlx{(‑s, ‑)}[14] \dicFieldCat{ekon.} \dicDirectTranslationCS{obchodní licence}
\dicEntry[verslunarmaður] \textls[29]{\dicTerm{verslunar··|maður} \dicIPA{{v}{\textepsilon}{\textsubring{r}}{s}{\textsubring{d}}{l}{\textscy}{n}{a}{r}{m}{a}{ð}{\textscy}{\textsubring{r}}} \dicPos{m}[13] \dicFlx{(‑manns, ‑menn)}[2]} \textbf{1.} \dicSynonym{afgreiðslumaður} \dicDirectTranslationCS{prodejce, prodejkyně}  \textbf{2.} \dicDirectTranslationCS{obchodník, obchodnice}

\dicEntry[verslunarmannafélag] \textls[15]{\dicTerm{verslunar·manna··fé·|lag} \dicIPA{{v}\-{\textepsilon}\-{\textsubring{r}}\-{s}\-{\textsubring{d}}\-{l}\-{\textscy}\-{n}\-{a}\-{r}\-{m}\-{a}\-{n}\-{a}\-{f}\-{j}\-{\textepsilon}\-{l}\-{a}\-{x}} \dicPos{n}[2] \dicFlx{(‑lags, ‑lög)}[8]} \dicDirectTranslationCS{sdružení obchodníků}

\dicEntry[verslunarmiðstöð] \dicTerm{verslunar··mið·stöð}\dicTerm{, verslanamiðstöð} \dicIPA{{v}\-{\textepsilon}\-{\textsubring{r}}\-{s}\-{\textsubring{d}}\-{l}\-{\textscy}\-{n}\-{a}\-{r}\-{m}\-{\textsci}\-{ð}\-{s}\-{\textsubring{d}}\-{\oe}\-{\texttheta}\-} \dicPos{f}[6] \dicFlx{(‑var, ‑var)}[1] \dicDirectTranslationCS{nákupní centrum\,/\addthin středisko}
\dicEntry[verslunarráð] \dicTerm{verslunar··ráð} \dicIPA{{v}{\textepsilon}{\textsubring{r}}{s}{\textsubring{d}}{l}{\textscy}{n}{a}{r}{au}{\texttheta}} \dicPos{n}[2] \dicFlx{(‑s, ‑)}[5] \dicFieldCat{ekon.} \dicDirectTranslationCS{obchodní komora}
\dicEntry[verslunarrekstur] \dicTerm{verslunar··rekst|ur} \dicIPA{{v}{\textepsilon}{\textsubring{r}}{s}{\textsubring{d}}{l}{\textscy}{n}{a}{r}{\textepsilon}{x}{s}{\textsubring{d}}{\textscy}{\textsubring{r}}} \dicPos{m}[5] \dicFlx{(‑rar\,/\addthin ‑urs)}[10] \dicDirectTranslationCS{obchodování, řízení\,/\addthin vedení obchodu}
\dicEntry[verslunarsamband] \dicTerm{verslunar··sam·|band} \dicIPA{{v}{\textepsilon}{\textsubring{r}}{s}{\textsubring{d}}{l}{\textscy}{n}{a}{\textsubring{r}}{s}{a}{m}{\textsubring{b}}{a}{n}{\textsubring{d}}} \dicPos{n}[2] \dicFlx{(‑bands, ‑bönd)}[8] \dicDirectTranslationCS{obchodní spojení\,/\addthin  vztah}
\dicEntry[verslunarsamningur] \dicTerm{verslunar··samn·ing|ur} \dicIPA{{v}\-{\textepsilon}\-{\textsubring{r}}\-{s}\-{\textsubring{d}}\-{l}\-{\textscy}\-{n}\-{a}\-{\textsubring{r}}\-{s}\-{a}\-{m}\-{n}\-{i}\-{\ng}\-{\r{g}}\-{\textscy}\-{\textsubring{r}}\-} \dicPos{m}[6] \dicFlx{(‑s, ‑ar)}[8] \dicDirectTranslationCS{obchodní dohoda\,/\addthin smlouva}
\dicEntry[verslunarskóli] \dicTerm{verslunar··skól|i} \dicIPA{{v}{\textepsilon}{\textsubring{r}}{s}{\textsubring{d}}{l}{\textscy}{n}{a}{\textsubring{r}}{s}{\r{g}}{ou}{l}{\textsci}} \dicPos{m}[1] \dicFlx{(‑a, ‑ar)}[1] \dicFieldCat{škol.} \dicDirectTranslationCS{obchodní akademie\,/\addthin škola}
\dicEntry[verslunarstaður] \dicTerm{verslunar··stað|ur} \dicIPA{{v}{\textepsilon}{\textsubring{r}}{s}{\textsubring{d}}{l}{\textscy}{n}{a}{\textsubring{r}}{s}{\textsubring{d}}{a}{ð}{\textscy}{\textsubring{r}}} \dicPos{m}[10] \dicFlx{(‑ar, ‑ir)}[14] \textbf{1.} \dicSynonym{bær} \dicDirectTranslationCS{obchodní centrum\,/\addthin středisko} \dicIndirectTranslationCS{(o~městě)}  \textbf{2.} \dicSynonym{búð} \dicDirectTranslationCS{obchodní centrum\,/\addthin středisko} \dicIndirectTranslationCS{(skupina obchodů)}
\dicEntry[verslunarstjóri] \dicTerm{verslunar··stjór|i} \dicIPA{{v}{\textepsilon}{\textsubring{r}}{s}{\textsubring{d}}{l}{\textscy}{n}{a}{\textsubring{r}}{s}{\textsubring{d}}{j}{ou}{r}{\textsci}} \dicPos{m}[1] \dicFlx{(‑a, ‑ar)}[1] \dicDirectTranslationCS{vedoucí prodejny\,/\addthin obchodu}
\dicEntry[verslunarvara] \dicTerm{verslunar··|vara} \dicIPA{{v}{\textepsilon}{\textsubring{r}}{s}{\textsubring{d}}{l}{\textscy}{n}{a}{r}{v}{a}{r}{a}} \dicPos{f}[1] \dicFlx{(‑vöru, ‑vörur)}[8] \dicDirectTranslationCS{obchodní zboží}
\dicEntry[versna] \dicTerm{versn|a} \dicIPA{{v}{\textepsilon}{\textsubring{r}}{s}{\textsubring{d}}{n}{a}} \dicPos{v}[1] \dicFlx{(‑aði)}[74] \textbf{1.} \dicPhraseIS{e‑ð versnar} \dicDirectTranslationCS{(co) se kazí, (co) se zhoršuje} \dicExampleIS{Veðrið hefur versnað.} \dicExampleCS{Počasí se zhoršilo.}  \textbf{2.} \dicPhraseIS{e‑m versnar} \dicFlx{impers} \dicDirectTranslationCS{(komu) je hůř, (komu) se dělá hůř}
\dicEntry[versnandi] \dicTerm{versn··andi} \dicIPA{{v}{\textepsilon}{\textsubring{r}}{s}{\textsubring{d}}{n}{a}{n}{\textsubring{d}}{\textsci}} \dicPos{adj}[13] \dicFlx{indecl}[1] \dicDirectTranslationCS{horšící se, zhoršující se}
\dicEntry[verst] \dicTerm{verst} \dicIPA{{v}{\textepsilon}{\textsubring{r}}{s}{\textsubring{d}}} \dicPos{adv} \dicFlx{sup} \dicLink{illa}
\dicEntry[verstur] \dicTerm{verstur} \dicIPA{{v}{\textepsilon}{\textsubring{r}}{s}{\textsubring{d}}{\textscy}{\textsubring{r}}} \dicPos{adj} \dicFlx{m sg nom sup} \textbf{1.} \dicLink{illur}  \textbf{2.} \dicLink{slæmur}  \textbf{3.} \dicLink{vondur}
\dicEntry[verstöð] \dicTerm{ver··stöð} \dicIPA{{v}{\textepsilon}{\textsubring{r}}{s}{\textsubring{d}}{\oe}{\texttheta}} \dicPos{f}[6] \dicFlx{(‑var, ‑var)}[1] \dicSynonym{ver\smash{\textsuperscript{1}}} \dicDirectTranslationCS{rybářská\,/\addthin rybolovná stanice}
\dicEntry[vert] \dicTerm{vert} \dicIPA{{v}{\textepsilon}{\textsubring{r}}{\textsubring{d}}} \dicPos{m}[4] \dicFlx{(‑s, ‑ar)}[22] \dicSynonym{veitingamaður} \dicDirectTranslationCS{hostinský, hostinská, restauratér(ka)}
\dicEntry[vertíð] \dicTerm{ver··tíð} \dicsymFrequent\  \dicIPA{{v}{\textepsilon}{\textsubring{r}}{t\smash{\textsuperscript{h}}}{i}{\texttheta}} \dicPos{f}[7] \dicFlx{(‑ar, ‑ir)}[1] \dicDirectTranslationCS{rybářská sezóna} \dicExampleIS{yfirstandandi vertíð} \dicExampleCS{stávající rybářská sezóna}
\dicEntry[verufall] \dicTerm{veru··|fall} \dicIPA{{v}{\textepsilon}{\textlengthmark}{r}{\textscy}{f}{a}{\textsubring{d}}{\textsubring{l}}} \dicPos{n}[2] \dicFlx{(‑falls, ‑föll)}[8] \dicFieldCat{jaz.} \dicDirectTranslationCS{essiv}
\dicEntry[verufræði] \dicTerm{veru··fræð|i} \dicIPA{{v}{\textepsilon}{\textlengthmark}{r}{\textscy}{f}{r}{a}{i}{ð}{\textsci}} \dicPos{f}[3] \dicFlx{(‑i)}[3] \dicFieldCat{filos.} \dicDirectTranslationCS{ontologie}
\dicEntry[verulega] \dicTerm{veru··lega} \dicsymFrequent\  \dicIPA{{v}{\textepsilon}{\textlengthmark}{r}{\textscy}{l}{\textepsilon}{\textbabygamma}{a}} \dicPos{adv} \textbf{1.} \dicSynonym{raunverulega} \dicDirectTranslationCS{skutečně, reálně, fakticky, opravdu}  \textbf{2.} \dicDirectTranslationCS{skutečně, notně, značně} \dicExampleIS{Hann dró verulega úr hraðanum.} \dicExampleCS{Značně snížil rychlost.}
\dicEntry[verulegur] \dicTerm{veru··legur} \dicsymFrequent\  \dicIPA{{v}{\textepsilon}{\textlengthmark}{r}{\textscy}{l}{\textepsilon}{\textbabygamma}{\textscy}{\textsubring{r}}} \dicPos{adj}[1]\dicFlx{}[-8] \textbf{1.} \dicSynonym{raunverulegur} \dicDirectTranslationCS{skutečný, reálný, faktický, opravdový}  \textbf{2.} \dicSynonym{talsverður} \dicDirectTranslationCS{skutečný, notný, značný} \dicExampleIS{veruleg lækkun} \dicExampleCS{značný pokles}
\dicEntry[veruleiki] \dicTerm{veru··leik|i} \dicIPA{{v}{\textepsilon}{\textlengthmark}{r}{\textscy}{l}{ei}{\r{\textObardotlessj}}{\textsci}} \dicPos{m}[1] \dicFlx{(‑a)}[3] \dicDirectTranslationCS{skutečnost, realita}
\dicEntry[verund] \dicTerm{ver··und} \dicIPA{{v}{\textepsilon}{\textlengthmark}{r}{\textscy}{n}{\textsubring{d}}} \dicPos{f}[7] \dicFlx{(‑ar, ‑ir)}[1] \dicFieldCat{filos.} \dicDirectTranslationCS{bytí}
\dicEntry[verustaður] \dicTerm{veru··stað|ur} \dicIPA{{v}{\textepsilon}{\textlengthmark}{r}{\textscy}{s}{\textsubring{d}}{a}{ð}{\textscy}{\textsubring{r}}} \dicPos{m}[10] \dicFlx{(‑ar, ‑ir)}[14] \dicDirectTranslationCS{příbytek, obydlí} \dicIndirectTranslationCS{(místo, kde někdo přebývá)}
\dicEntry[verúlfur] \dicTerm{ver··úlf|ur} \dicIPA{{v}{\textepsilon}{\textlengthmark}{r}{u}{l}{v}{\textscy}{\textsubring{r}}} \dicPos{m}[6] \dicFlx{(‑s, ‑ar)}[4] \dicLink{varúlfur}
\dicEntry[veröld] \dicTerm{ver··|öld} \dicsymFrequent\  \dicIPA{{v}{\textepsilon}{\textlengthmark}{r}{\oe}{l}{\textsubring{d}}} \dicPos{f}[7] \dicFlx{(‑aldar, ‑aldir)}[16] \dicSynonym{heimur} \dicDirectTranslationCS{svět};  \dicPhraseIS{á hjara veraldar} \dicFlx{adv} \dicDirectTranslationCS{na okraji světa};  \dicPhraseIS{hin sjö undur veraldar} \dicDirectTranslationCS{sedm divů světa};  \dicPhraseIS{um víða veröld} \dicFlx{adv} \dicDirectTranslationCS{po celém světě} \dicExampleIS{Hann er frægur um víða veröld.} \dicExampleCS{Je slavný po celém světě.}
\dicEntry[verönd] \dicTerm{ver|önd} \dicIPA{{v}{\textepsilon}{\textlengthmark}{r}{\oe}{n}{\textsubring{d}}} \dicPos{f}[7] \dicFlx{(‑andar, ‑andir)}[16] \dicDirectTranslationCS{veranda}
\dicEntry[vesaldarlegur] \dicTerm{vesaldar··legur}\dicTerm{, vesældarlegur} \dicIPA{{v}\-{\textepsilon}\-{\textlengthmark}\-{s}\-{a}\-{l}\-{\textsubring{d}}\-{a}\-{r}\-{l}\-{\textepsilon}\-{\textbabygamma}\-{\textscy}\-{\textsubring{r}}\-} \dicPos{adj}[1]\dicFlx{}[-8] \dicSynonym{eymdarlegur} \dicDirectTranslationCS{ubohý, žalostný} \dicExampleIS{vera vesaldarlegur á svipinn} \dicExampleCS{mít žalostný výraz}
\dicEntry[vesaldómur] \dicTerm{vesal··dóm|ur} \dicIPA{{v}{\textepsilon}{\textlengthmark}{s}{a}{l}{\textsubring{d}}{ou}{m}{\textscy}{\textsubring{r}}} \dicPos{m}[6] \dicFlx{(‑s)}[12] \textbf{1.} \dicSynonym*{bágt ástand} \dicDirectTranslationCS{nouze, tíseň} \dicExampleIS{lifa í vesaldómi} \dicExampleCS{žít v~nouzi}  \textbf{2.} \dicSynonym{ódugnaður} \dicDirectTranslationCS{slabost, nedostatek vůle}
\dicEntry[vesalingur] \dicTerm{vesal··ing|ur}\dicTerm{, veslingur} \dicIPA{{v}\-{\textepsilon}\-{\textlengthmark}\-{s}\-{a}\-{l}\-{i}\-{\ng}\-{\r{g}}\-{\textscy}\-{\textsubring{r}}\-} \dicPos{m}[6] \dicFlx{(‑s, ‑ar)}[8] \dicSynonym{aumingi} \dicDirectTranslationCS{nešťastník, nešťastnice, ubožák, ubožačka, chudák, chudačka};  \dicPhraseIS{vesalings e‑r} \dicDirectTranslationCS{chudák (kdo), nešťastný (kdo), ubohý (kdo)}
\dicEntry[vesall] \dicTerm{ves|all} \dicIPA{{v}{\textepsilon}{\textlengthmark}{s}{a}{\textsubring{d}}{\textsubring{l}}} \dicPos{adj}[8] \dicFlx{(f ‑öl)}[6] \dicLink{vesæll}
\dicEntry[vesen] \dicTerm{vesen} \dicIPA{{v}{\textepsilon}{\textlengthmark}{s}{\textepsilon}{\textsubring{n}}} \dicPos{n}[2] \dicFlx{(‑s)}[2] \dicLangCat{hovor.} \dicSynonym{vafstur} \dicDirectTranslationCS{trable, potíže, problémy};  \dicPhraseIS{vera með vesen} \dicDirectTranslationCS{dělat potíže, způsobit problémy};  \dicPhraseIS{það er vesen á e‑m} \dicDirectTranslationCS{s~(kým) jsou problémy\,/\addthin trable}
\dicEntry[veski] \dicTerm{veski} \dicIPA{{v}{\textepsilon}{s}{\r{\textObardotlessj}}{\textsci}} \dicPos{n}[2] \dicFlx{(‑s, ‑)}[16] \textbf{1.} \dicSynonym{handtaska} \dicDirectTranslationCS{kabelka, (náprsní) taška}  \textbf{2.} \dicSynonym{peningaveski} \dicDirectTranslationCS{peněženka}
\dicEntry[veslingur] \dicTerm{vesl··ing|ur} \dicIPA{{v}{\textepsilon}{s}{\textsubring{d}}{l}{i}{\ng}{\r{g}}{\textscy}{\textsubring{r}}} \dicPos{m}[6] \dicFlx{(‑s, ‑ar)}[8] \dicLink{vesalingur}
\dicEntry[vespa] \dicTerm{vesp|a} \dicIPA{{v}{\textepsilon}{s}{\textsubring{b}}{a}} \dicPos{f}[1] \dicFlx{(‑u, ‑ur)}[19] \textbf{1.} \dicFieldCat{zool.} \dicSynonym{húsageitungur} \dicDirectTranslationCS{vosa} \textit{(l.~{\textLA{Vespa}})}  \dicsymPhoto\   \textbf{2.} \dicDirectTranslationCS{skútr}
\dicFigure{90405.jpg}{Vespa}{Vespa - Krejčík Stanislav, Biolib, Copyright/CC-BY}
\dicEntry[vessaæð] \dicTerm{vessa··æð} \dicIPA{{v}{\textepsilon}{s}{\textlengthmark}{a}{a}{i}{\texttheta}} \dicPos{f}[4] \dicFlx{(‑ar, ‑ar)}[1] \dicFieldCat{anat.} \dicDirectTranslationCS{lymfatická céva}
\dicEntry[vessi] \dicTerm{vess|i} \dicIPA{{v}{\textepsilon}{s}{\textlengthmark}{\textsci}} \dicPos{m}[1] \dicFlx{(‑a, ‑ar)}[1] \textbf{1.} \dicFieldCat{biol.} \dicDirectTranslationCS{výměšek, sekret}  \textbf{2.} \dicFieldCat{anat.} \dicDirectTranslationCS{míza, lymfa}
\dicEntry[vestan] \dicTerm{vestan} \dicsymFrequent\  \dicIPA{{v}{\textepsilon}{s}{\textsubring{d}}{a}{\textsubring{n}}} \dicPos{prep\,/\addthin adv} \dicFlx{adv} \dicDirectTranslationCS{ze západu};  \dicPhraseIS{vestan e‑s} \dicFlx{prep gen} \dicDirectTranslationCS{na západ od (čeho), západně od (čeho)};  \dicIdiom{að}[vestan]{ \dicPhraseIS{að vestan}} \dicFlx{adv} {\textbf{a.}} \dicDirectTranslationCS{ze západu};  {\textbf{b.}} \dicDirectTranslationCS{po\,/\addthin na západní straně};  \dicIdiom{vestan}[að]{ \dicPhraseIS{vestan að}} \dicFlx{adv} \dicDirectTranslationCS{od\,/\addthin ze západu};  \dicIdiom{vestan}[á]{ \dicPhraseIS{vestan á e‑u}} \dicFlx{prep} \dicDirectTranslationCS{ze západní strany (čeho), na západní straně (čeho)};  \dicIdiom{vestan}[frá]{ \dicPhraseIS{vestan frá}} \dicFlx{adv} \dicDirectTranslationCS{ze\,/\addthin od západu};  \dicIdiom{vestan}[í móti]{ \dicPhraseIS{vestan í móti}} \dicFlx{adv} \dicDirectTranslationCS{směrem na západ};  \dicIdiom{vestan}[megin]{ \dicPhraseIS{vestanmegin, vestan megin}} \dicFlx{adv} \dicDirectTranslationCS{na západní straně, v~západní části};  \dicIdiom{vestan}[til]{ \dicPhraseIS{vestan til}} \dicFlx{adv} \dicDirectTranslationCS{na západě, v~západní části};  \dicIdiom{vestan}[undir]{ \dicPhraseIS{vestan undir e‑u}} \dicFlx{prep} \dicDirectTranslationCS{západně pod (čím), ze západní strany (čeho)};  \dicIdiom{vestan}[við]{ \dicPhraseIS{vestan við e‑ð}} \dicFlx{prep} \dicDirectTranslationCS{na západní straně (čeho), západně od (čeho)};  \dicIdiom{vestan}[yfir]{ \dicPhraseIS{vestan yfir e‑ð}} \dicFlx{prep} \dicDirectTranslationCS{ze západu přes (co)}
\dicEntry[vestanátt] \dicTerm{vestan··átt} \dicIPA{{v}{\textepsilon}{s}{\textsubring{d}}{a}{n}{au}{h}{\textsubring{d}}} \dicPos{f}[7] \dicFlx{(‑ar, ‑ir)}[1] \dicDirectTranslationCS{západní vítr}
\dicEntry[vestanverður] \dicTerm{vestan··verður} \dicIPA{{v}{\textepsilon}{s}{\textsubring{d}}{a}{n}{v}{\textepsilon}{r}{ð}{\textscy}{\textsubring{r}}} \dicPos{adj}[2]\dicFlx{}[-4] \dicDirectTranslationCS{západně polože\-ný};  \dicPhraseIS{að vestanverðu} \dicFlx{adv} \dicDirectTranslationCS{po\,/\addthin na západní straně}
\dicEntry[vestar] \dicTerm{vestar} \dicIPA{{v}{\textepsilon}{s}{\textsubring{d}}{a}{\textsubring{r}}} \dicPos{adv} \dicFlx{comp} \dicLink{vestur\smash{\textsuperscript{2}}}
\dicEntry[vestari] \dicTerm{vestari}\dicTerm{, vestri\smash{\textsuperscript{2}}} \dicIPA{{v}\-{\textepsilon}\-{s}\-{\textsubring{d}}\-{a}\-{r}\-{\textsci}\-} \dicPos{adj}[12] \dicFlx{comp (sup vestastur)}[2] \dicDirectTranslationCS{západnější}
\dicEntry[vestast] \dicTerm{vestast} \dicIPA{{v}{\textepsilon}{s}{\textsubring{d}}{a}{s}{\textsubring{d}}} \dicPos{adv} \dicFlx{sup} \dicLink{vestur\smash{\textsuperscript{2}}}
\dicEntry[vestastur] \dicTerm{vestastur} \dicIPA{{v}{\textepsilon}{s}{\textsubring{d}}{a}{s}{\textsubring{d}}{\textscy}{\textsubring{r}}} \dicPos{adj} \dicFlx{m sg nom sup} \dicLink{vestari}
\dicEntry[Vestfirðingur] \dicTerm{Vest··firð·ing|ur} \dicIPA{{v}{\textepsilon}{s}{\textsubring{d}}{f}{\textsci}{r}{ð}{i}{\ng}{\r{g}}{\textscy}{\textsubring{r}}} \dicPos{m}[6] \dicFlx{(‑s, ‑ar)}[8] \dicDirectTranslationCS{obyvatel(ka) Západních fjordů}
\dicEntry[Vestfirðir] \dicTerm{Vest··firðir} \dicIPA{{v}{\textepsilon}{s}{\textsubring{d}}{f}{\textsci}{r}{ð}{\textsci}{\textsubring{r}}} \dicPos{m}[12] \dicFlx{pl}[1] \dicFieldCat{geog.} \dicDirectTranslationCS{Západní fjordy} \dicIndirectTranslationCS{(severozápadní část Islandu)}
\dicEntry[vestfirska] \dicTerm{vest··firsk|a} \dicIPA{{v}{\textepsilon}{s}{\textsubring{d}}{f}{\textsci}{\textsubring{r}}{s}{\r{g}}{a}} \dicPos{f}[1] \dicFlx{(‑u)}[5] \dicFieldCat{jaz.} \dicIndirectTranslationCS{nářečí islandštiny používané v~Západních fjordech}
\dicEntry[vestfirskur] \dicTerm{vest··firskur} \dicIPA{{v}{\textepsilon}{s}{\textsubring{d}}{f}{\textsci}{\textsubring{r}}{s}{\r{g}}{\textscy}{\textsubring{r}}} \dicPos{adj}[1]\dicFlx{}[-1] \dicDirectTranslationCS{(jsoucí) ze Západních fjordů}
\dicEntry[vesti] \dicTerm{vesti} \dicIPA{{v}{\textepsilon}{s}{\textsubring{d}}{\textsci}} \dicPos{n}[2] \dicFlx{(‑s, ‑)}[14] \dicDirectTranslationCS{vesta}
\dicEntry[vestlægur] \dicTerm{vest··lægur} \dicIPA{{v}{\textepsilon}{s}{\textsubring{d}}{l}{a}{i}{\textbabygamma}{\textscy}{\textsubring{r}}} \dicPos{adj}[1]\dicFlx{}[-1] \dicDirectTranslationCS{západní, (jsoucí) od západu} \dicExampleIS{vestlægur vindur} \dicExampleCS{vítr od západu} \dicAntonym{austlægur};  \dicPhraseIS{vestlæg lengd} \dicFieldCat{geog.} \dicDirectTranslationCS{západní šířka}
\dicEntry[Vestmannaeyjar] \dicTerm{Vest·manna··eyjar} \dicIPA{{v}{\textepsilon}{s}{\textsubring{d}}{m}{a}{n}{a}{ei}{j}{a}{\textsubring{r}}} \dicPos{f}[4] \dicFlx{pl}[15] \dicFieldCat{geog.} \dicDirectTranslationCS{Vestmannské ostrovy} \dicIndirectTranslationCS{(ostrovy jižně od Islandu)}
\dicEntry[vestri] \dicTerm{vestr|i\smash{\textsuperscript{1}}} \dicIPA{{v}{\textepsilon}{s}{\textsubring{d}}{r}{\textsci}} \dicPos{m}[1] \dicFlx{(‑a, ‑ar)}[1] \dicDirectTranslationCS{western, kovbojka}
\dicEntry[vestri] \dicTerm{vestri\smash{\textsuperscript{2}}} \dicIPA{{v}{\textepsilon}{s}{\textsubring{d}}{r}{\textsci}} \dicPos{adj}[12] \dicFlx{comp}[5] \dicLink{vestari}
\dicEntry[vestrænn] \dicTerm{vest··rænn} \dicIPA{{v}{\textepsilon}{s}{\textsubring{d}}{r}{a}{i}{\textsubring{d}}{\textsubring{n}}} \dicPos{adj}[7]\dicFlx{}[-1] \textbf{1.} \dicSynonym{vestlægur} \dicDirectTranslationCS{západní} \dicExampleIS{vestrænn vindur} \dicExampleCS{západní vítr}  \textbf{2.} \dicDirectTranslationCS{západní, (jsoucí) ze západu} \dicIndirectTranslationCS{(původem)} \dicExampleIS{vestræn áhrif} \dicExampleCS{západní vliv}
\dicEntry[vestur] \dicTerm{vestur\smash{\textsuperscript{1}}} \dicsymFrequent\  \dicIPA{{v}{\textepsilon}{s}{\textsubring{d}}{\textscy}{\textsubring{r}}} \dicPos{n}[2] \dicFlx{(‑s)}[28] \dicDirectTranslationCS{západ};  \dicPhraseIS{af vestri} \dicFlx{adv} \dicDirectTranslationCS{ze západu};  \dicPhraseIS{í vestri} \dicFlx{adv} \dicDirectTranslationCS{na západě};  \dicPhraseIS{í vestur} \dicFlx{adv} \dicDirectTranslationCS{na západ};  \dicPhraseIS{úr vestri} \dicFlx{adv} \dicDirectTranslationCS{ze západu}
\dicEntry[vestur] \dicTerm{vest|ur\smash{\textsuperscript{2}}} \dicsymFrequent\  \dicIPA{{v}{\textepsilon}{s}{\textsubring{d}}{\textscy}{\textsubring{r}}} \dicPos{adv} \dicFlx{(comp ‑ar, sup ‑ast)} \textbf{1.} \dicSynonym*{í vestur} \dicDirectTranslationCS{na západ, západně}  \textbf{2.} \dicSynonym*{um Vestfirði} \dicDirectTranslationCS{do Západních fjordů, na Západní fjordy} \dicExampleIS{búa vestur á fjörðum} \dicExampleCS{bydlet na Západních fjordech};  \dicIdiom{vestur}[af]{ \dicPhraseIS{vestur af e‑u}} \dicFlx{prep} \dicDirectTranslationCS{na západní straně (čeho), západně od (čeho)}; { \dicPhraseIS{vestur af}} \dicFlx{adv} \dicDirectTranslationCS{západně, (směrem) na západ};  \dicIdiom{vestur}[eftir]{ \dicPhraseIS{vestur eftir}} \dicFlx{adv} \dicDirectTranslationCS{(směrem) k~západu, (směrem) na západ};  \dicIdiom{vestur}[frá]{ \dicPhraseIS{vestur frá}} \dicFlx{adv} \dicDirectTranslationCS{na západě, na západ};  \dicIdiom{vestur}[fyrir]{ \dicPhraseIS{vestur fyrir e‑ð}} \dicFlx{prep} \dicDirectTranslationCS{na západ od (čeho)};  \dicIdiom{vestur}[með]{ \dicPhraseIS{vestur með e‑u}} \dicFlx{prep} \dicDirectTranslationCS{na západ podél (čeho)};  \dicIdiom{vestur}[um]{ \dicPhraseIS{vestur um e‑ð}} \dicFlx{prep} \dicDirectTranslationCS{na západ (čeho)}; { \dicPhraseIS{vestur um}} \dicFlx{adv} \dicDirectTranslationCS{na západ, západně};  \dicIdiom{vestur}[undir]{ \dicPhraseIS{vestur undir e‑ð}} \dicFlx{prep} \dicDirectTranslationCS{západně pod (co)} \dicIndirectTranslationCS{(o~pohybu)}; { \dicPhraseIS{vestur undir e‑u}} \dicFlx{prep} \dicDirectTranslationCS{západně pod (čím)} \dicIndirectTranslationCS{(o~umístění)};  \dicIdiom{vestur}[úr]{ \dicPhraseIS{vestur úr e‑u}} \dicFlx{prep} \dicDirectTranslationCS{západně z~(čeho)}
\dicEntry[Vestur-Íslendingur] \dicTerm{Vestur-Ís·lend··ing|ur} \dicIPA{{v}\-{\textepsilon}\-{s}\-{\textsubring{d}}\-{\textscy}\-{r}\-{i}\-{s}\-{l}\-{\textepsilon}\-{n}\-{\textsubring{d}}\-{i}\-{\ng}\-{\r{g}}\-{\textscy}\-{\textsubring{r}}\-} \dicPos{m}[6] \dicFlx{(‑s, ‑ar)}[8] \dicIndirectTranslationCS{Islanďan, který emigroval do Severní Ameriky v~období 1870--1914}
\dicEntry[Vestur-Kongó] \dicTerm{Vestur-Kongó} \dicIPA{{v}{\textepsilon}{s}{\textsubring{d}}{\textscy}{\textsubring{r}}{k\smash{\textsuperscript{h}}}{\textopeno}{\ng}{\r{g}}{ou}} \dicPos{n}[4] \dicFlx{indecl}[2] \dicFieldCat{geog.} \dicDirectTranslationCS{Republika Kongo, Kongo-Brazzaville}

\dicEntry[Vestur-Kongómaður] \textls[25]{\dicTerm{Vestur-Kongó··|maður} \dicIPA{{v}{\textepsilon}{s}{\textsubring{d}}{\textscy}{\textsubring{r}}{k\smash{\textsuperscript{h}}}{\textopeno}{\ng}{\r{g}}{ou}{m}{a}{ð}{\textscy}{\textsubring{r}}} \dicPos{m}[13]} \dicFlx{(‑manns, ‑menn)}[2] \dicDirectTranslationCS{obyvatel(ka) Republiky Kongo\,/\addthin Konga-Brazza\-ville}

\dicEntry[Vesturálfa] \dicTerm{Vestur··álf|a} \dicIPA{{v}{\textepsilon}{s}{\textsubring{d}}{\textscy}{r}{au}{l}{v}{a}} \dicPos{f}[1] \dicFlx{(‑u)}[5] \dicFieldCat{geog.} \dicSynonym{Ameríka} \dicDirectTranslationCS{(Severní) Amerika}
\dicEntry[vesturátt] \dicTerm{vestur··átt} \dicIPA{{v}{\textepsilon}{s}{\textsubring{d}}{\textscy}{r}{au}{h}{\textsubring{d}}} \dicPos{f}[7] \dicFlx{(‑ar)}[3] \dicDirectTranslationCS{západní směr} \dicExampleIS{í vesturátt} \dicExampleCS{západním směrem}
\dicEntry[vesturfari] \dicTerm{vestur··far|i} \dicIPA{{v}{\textepsilon}{s}{\textsubring{d}}{\textscy}{\textsubring{r}}{f}{a}{r}{\textsci}} \dicPos{m}[1] \dicFlx{(‑a, ‑ar)}[8] \dicFieldCat{hist.} \dicIndirectTranslationCS{emigrant(ka) směřující z~Islandu na západ do Severní Ameriky}
\dicEntry[vesturheimskur] \dicTerm{vestur··heimskur} \dicIPA{{v}{\textepsilon}{s}{\textsubring{d}}{\textscy}{\textsubring{r}}{h}{ei}{m}{s}{\r{g}}{\textscy}{\textsubring{r}}} \dicPos{adj}[1]\dicFlx{}[-6] \dicDirectTranslationCS{(jsoucí) ze Severní Ameriky (islandský emigrant ap.)}
\dicEntry[Vesturheimur] \dicTerm{Vestur··heim|ur} \dicIPA{{v}{\textepsilon}{s}{\textsubring{d}}{\textscy}{\textsubring{r}}{h}{ei}{m}{\textscy}{\textsubring{r}}} \dicPos{m}[6] \dicFlx{(‑s)}[3] \dicFieldCat{hist.} \dicDirectTranslationCS{Severní Amerika} \dicIndirectTranslationCS{(cíl islandských emigrantů v~období 1870--1914)}
\dicEntry[vesturhvel] \dicTerm{vestur··hvel} \dicIPA{{v}{\textepsilon}{s}{\textsubring{d}}{\textscy}{\textsubring{r}}{k\smash{\textsuperscript{h}}}{v}{\textepsilon}{\textsubring{l}}} \dicPos{n}[2] \dicFlx{(‑s, ‑)}[5] \dicFieldCat{geog.} \dicDirectTranslationCS{západní polokoule}
\dicEntry[vesturíslenskur] \textls[10]{\dicTerm{vestur··ís·lenskur} \dicIPA{{v}{\textepsilon}{s}{\textsubring{d}}{\textscy}{r}{i}{s}{l}{\textepsilon}{n}{s}{\r{g}}{\textscy}{\textsubring{r}}}} \dicPos{adj}[1]\dicFlx{}[-6] \dicDirectTranslationCS{západoislandský} \dicIndirectTranslationCS{(o~člověku, který emigroval z~Islandu do Severní Ameriky v~období 1870--1914)}
\dicEntry[vesturkongóskur] \dicTerm{vestur··kongóskur} \dicIPA{{v}{\textepsilon}{s}{\textsubring{d}}{\textscy}{\textsubring{r}}{k\smash{\textsuperscript{h}}}{\textopeno}{\ng}{\r{g}}{ou}{s}{\r{g}}{\textscy}{\textsubring{r}}} \dicPos{adj}[1]\dicFlx{}[-6] \dicDirectTranslationCS{týkající se Republiky Kongo\,/\addthin Konga-Brazzaville}
\dicEntry[Vesturland] \dicTerm{Vestur··land} \dicIPA{{v}{\textepsilon}{s}{\textsubring{d}}{\textscy}{r}{l}{a}{n}{\textsubring{d}}} \dicPos{n}[2] \dicFlx{(‑s)}[4] \dicFieldCat{geog.} \dicDirectTranslationCS{západní Island, západní část Islandu}
\dicEntry[Vesturlönd] \dicTerm{Vestur··lönd} \dicIPA{{v}{\textepsilon}{s}{\textsubring{d}}{\textscy}{r}{l}{\oe}{n}{\textsubring{d}}} \dicPos{n}[2] \dicFlx{pl}[9] \dicFieldCat{geog.} \dicDirectTranslationCS{Západ} \dicIndirectTranslationCS{(západní Evropa a~státy Severní Ameriky)}
\dicEntry[vesturströnd] \textls[10]{\dicTerm{vestur··|strönd} \dicIPA{{v}{\textepsilon}{s}{\textsubring{d}}{\textscy}{\textsubring{r}}{s}{\textsubring{d}}{r}{\oe}{n}{\textsubring{d}}} \dicPos{f}[8] \dicFlx{(‑strandar, ‑strendur\,/\addthin ‑strand\-ir)}[9]} \dicDirectTranslationCS{západní pobřeží}
\dicEntry[vesturveldi] \dicTerm{vestur··veldi} \dicIPA{{v}{\textepsilon}{s}{\textsubring{d}}{\textscy}{r}{v}{\textepsilon}{l}{\textsubring{d}}{\textsci}} \dicPos{n}[2] \dicFlx{(‑s, ‑)}[14] \dicDirectTranslationCS{Západ, západní mocnosti}
\dicEntry[vesæld] \dicTerm{vesæld} \dicIPA{{v}{\textepsilon}{\textlengthmark}{s}{a}{i}{l}{\textsubring{d}}} \dicPos{f}[4] \dicFlx{(‑ar)}[3] \dicDirectTranslationCS{bída, nouze}
\dicEntry[vesældarlegur] \dicTerm{vesældar··legur} \dicIPA{{v}{\textepsilon}{\textlengthmark}{s}{a}{i}{l}{\textsubring{d}}{a}{r}{l}{\textepsilon}{\textbabygamma}{\textscy}{\textsubring{r}}} \dicPos{adj}[1]\dicFlx{}[-8] \dicLink{vesaldarlegur}
\dicEntry[vesæll] \dicTerm{vesæll}\dicTerm{, vesall} \dicIPA{{v}\-{\textepsilon}\-{\textlengthmark}\-{s}\-{a}\-{i}\-{\textsubring{d}}\-{\textsubring{l}}\-} \dicPos{adj}[8]\dicFlx{}[-1] \dicDirectTranslationCS{ubohý, bídný, zbědovaný}
\dicEntry[vesöld] \dicTerm{ves|öld} \dicIPA{{v}{\textepsilon}{\textlengthmark}{s}{\oe}{l}{\textsubring{d}}} \dicPos{f}[7] \dicFlx{(‑aldar)}[19] \textbf{1.} \dicSynonym{eymd} \dicDirectTranslationCS{bída, nouze}  \textbf{2.} \dicSynonym{ódugnaður} \dicDirectTranslationCS{splín, apatie}
\dicEntry[vetfang] \dicTerm{vet··fang}\dicTerm{, vetfangur} \dicIPA{{v}\-{\textepsilon}\-{\textlengthmark}\-{\textsubring{d}}\-{f}\-{au}\-{\ng}\-{\r{g}}\-} \dicPos{n}[2] \dicFlx{(‑s)}[2] \dicPhraseIS{í einu vetfangi} \dicFlx{adv} \dicDirectTranslationCS{v~mžiku, v~okamžiku} \dicExampleIS{Lögreglan kom að í einu vetfangi.} \dicExampleCS{Policie v~okamžiku přijela.}
\dicEntry[vetfangur] \dicTerm{vet··fang|ur} \dicIPA{{v}{\textepsilon}{\textlengthmark}{\textsubring{d}}{f}{au}{\ng}{\r{g}}{\textscy}{\textsubring{r}}} \dicPos{m}[6] \dicFlx{(‑s)}[7] \dicLink{vetfang}
\dicEntry[vetni] \dicTerm{vetni} \dicIPA{{v}{\textepsilon}{h}{\textsubring{d}}{n}{\textsci}} \dicPos{n}[2] \dicFlx{(‑s)}[20] \dicFieldCat{chem.} \dicDirectTranslationCS{vodík} \textit{(l.~{\textLA{H, Hydrogenium}})}
\dicEntry[vetnissprengja] \dicTerm{vetnis··sprengj|a} \dicIPA{{v}{\textepsilon}{h}{\textsubring{d}}{n}{\textsci}{s}{\textsubring{b}}{r}{ei}{\textltailn}{\r{\textObardotlessj}}{a}} \dicPos{f}[1] \dicFlx{(‑u, ‑ur)}[25] \dicFieldCat{voj.} \dicDirectTranslationCS{vodíková bomba}
\dicEntry[vetra] \dicTerm{vetr|a} \dicIPA{{v}{\textepsilon}{\textlengthmark}{\textsubring{d}}{r}{a}} \dicPos{v}[1] \dicFlx{(‑aði)}[83] \dicFlx{impers} \dicPhraseIS{það vetrar} \dicDirectTranslationCS{začíná zima, zima je na cestě};  \dicPhraseIS{það er farið að vetra} \dicDirectTranslationCS{blíží se zima, nastává\,/\addthin přichází zima}
\dicEntry[vetrar] \dicTerm{vetrar} \dicIPA{{v}{\textepsilon}{\textlengthmark}{\textsubring{d}}{r}{a}{\textsubring{r}}} \dicPos{m} \dicFlx{sg gen} \dicLink{vetur}
\dicEntry[vetraráætlun] \dicTerm{vetrar··á·ætl|un} \dicIPA{{v}{\textepsilon}{\textlengthmark}{\textsubring{d}}{r}{a}{r}{au}{a}{i}{h}{\textsubring{d}}{l}{\textscy}{\textsubring{n}}} \dicPos{f}[7] \dicFlx{(‑unar, ‑anir)}[8] \dicDirectTranslationCS{jízdní řád pro zimní období}
\dicEntry[vetrarblóm] \dicTerm{vetrar··blóm} \dicIPA{{v}{\textepsilon}{\textlengthmark}{\textsubring{d}}{r}{a}{r}{\textsubring{b}}{l}{ou}{\textsubring{m}}} \dicPos{n}[2] \dicFlx{(‑s, ‑)}[5] \dicFieldCat{bot.} \dicDirectTranslationCS{lomikámen vstřícnolistý} \textit{(l.~{\textLA{Saxifraga oppositifolia}})}  \dicsymPhoto\ 
\dicFigure{ds_image_vetrarblom_0_1.jpg}{Vetrarblóm}{Vetrarblóm - Michael Haferkamp, GFDL}
\dicEntry[vetrarbraut] \dicTerm{vetrar··braut} \dicIPA{{v}{\textepsilon}{\textlengthmark}{\textsubring{d}}{r}{a}{r}{\textsubring{b}}{r}{\oe i}{\textsubring{d}}} \dicPos{f}[7] \dicFlx{(‑ar, ‑ir)}[1] \dicFieldCat{astro.} \dicDirectTranslationCS{galaxie};  \dicPhraseIS{Vetrarbrautin} \dicFlx{def} \dicFieldCat{astro.} \dicDirectTranslationCS{Mléčná dráha}
\dicEntry[vetrarbúningur] \dicTerm{vetrar··bú·ning|ur} \dicIPA{{v}{\textepsilon}{\textlengthmark}{\textsubring{d}}{r}{a}{r}{\textsubring{b}}{u}{n}{i}{\ng}{\r{g}}{\textscy}{\textsubring{r}}} \dicPos{m}[6] \dicFlx{(‑s, ‑ar)}[8] \textbf{1.} \dicLangCat{přen.} \dicDirectTranslationCS{zimní kabát\,/\addthin plášť}  \textbf{2.} \dicSynonym*{vetrargervi} \dicDirectTranslationCS{zimní opeření (ptáků ap.)}
\dicEntry[vetrardagur] \dicTerm{vetrar··dag|ur} \dicIPA{{v}{\textepsilon}{\textlengthmark}{\textsubring{d}}{r}{a}{r}{\textsubring{d}}{a}{\textbabygamma}{\textscy}{\textsubring{r}}} \dicPos{m}[6] \dicFlx{(‑s, ‑ar)}[62] \dicDirectTranslationCS{zimní den};  \dicPhraseIS{fyrsti vetrardagur} \dicDirectTranslationCS{první zimní den} \dicIndirectTranslationCS{(sobota v~období 21.--27. října)}
\dicEntry[vetrardekk] \dicTerm{vetrar··dekk} \dicIPA{{v}{\textepsilon}{\textlengthmark}{\textsubring{d}}{r}{a}{r}{\textsubring{d}}{\textepsilon}{h}{\r{g}}} \dicPos{n}[2] \dicFlx{(‑s, ‑)}[13] \dicDirectTranslationCS{zimní pneumatika} \dicAntonym{sumardekk}
\dicEntry[vetrarfrakki] \dicTerm{vetrar··frakk|i} \dicIPA{{v}{\textepsilon}{\textlengthmark}{\textsubring{d}}{r}{a}{\textsubring{r}}{f}{r}{a}{h}{\r{\textObardotlessj}}{\textsci}} \dicPos{m}[1] \dicFlx{(‑a, ‑ar)}[8] \dicDirectTranslationCS{zimník, zimní kabát}
\dicEntry[vetraríþrótt] \dicTerm{vetrar··í·þrótt} \dicIPA{{v}{\textepsilon}{\textlengthmark}{\textsubring{d}}{r}{a}{r}{i}{\texttheta}{r}{ou}{h}{\textsubring{d}}} \dicPos{f}[7] \dicFlx{(‑ar, ‑ir)}[1] \dicDirectTranslationCS{zimní sport} \dicIndirectTranslationCS{(zvláště v~pl)}
\dicEntry[vetrarlag] \dicTerm{vetrar··lag} \dicIPA{{v}{\textepsilon}{\textlengthmark}{\textsubring{d}}{r}{a}{r}{l}{a}{x}} \dicPos{n}[2] \dicFlx{(‑s)}[2] \dicDirectTranslationCS{zimní období, zima};  \dicPhraseIS{að vetrarlagi} \dicFlx{adv} \dicDirectTranslationCS{v~zimě, během zimy}
\dicEntry[vetrarlangt] \dicTerm{vetrar··langt} \dicIPA{{v}{\textepsilon}{\textlengthmark}{\textsubring{d}}{r}{a}{r}{l}{au}{\r{\ng}}{\textsubring{d}}} \dicPos{adv} \dicDirectTranslationCS{celou zimu, přes zimu}
\dicEntry[vetrarlegur] \dicTerm{vetrar··legur} \dicIPA{{v}{\textepsilon}{\textlengthmark}{\textsubring{d}}{r}{a}{r}{l}{\textepsilon}{\textbabygamma}{\textscy}{\textsubring{r}}} \dicPos{adj}[1]\dicFlx{}[-8] \dicDirectTranslationCS{zimní, připomínající zimu}
\dicEntry[vetrarsólhvörf] \dicTerm{vetrar··sól·hvörf}\dicTerm{, vetrarsólstöður} \dicIPA{{v}\-{\textepsilon}\-{\textlengthmark}\-{\textsubring{d}}\-{r}\-{a}\-{\textsubring{r}}\-{s}\-{ou}\-{\textsubring{l}}\-{k\smash{\textsuperscript{h}}}\-{v}\-{\oe}\-{r}\-{f}\-} \dicPos{n}[2] \dicFlx{pl}[9] \dicDirectTranslationCS{zimní slunovrat}
\dicEntry[vetrarsólstöður] \dicTerm{vetrar··sól·stöður} \dicIPA{{v}{\textepsilon}{\textlengthmark}{\textsubring{d}}{r}{a}{\textsubring{r}}{s}{ou}{\textsubring{l}}{s}{\textsubring{d}}{\oe}{ð}{\textscy}{\textsubring{r}}} \dicPos{f}[1] \dicFlx{pl}[9] \dicLink{vetrarsólhvörf}
\dicEntry[vetrarveður] \dicTerm{vetrar··veður} \dicIPA{{v}{\textepsilon}{\textlengthmark}{\textsubring{d}}{r}{a}{r}{v}{\textepsilon}{ð}{\textscy}{\textsubring{r}}} \dicPos{n}[2] \dicFlx{(‑s, ‑)}[25] \dicDirectTranslationCS{zimní počasí}
\dicEntry[vetrarvertíð] \dicTerm{vetrar··ver·tíð} \dicIPA{{v}{\textepsilon}{\textlengthmark}{\textsubring{d}}{r}{a}{r}{v}{\textepsilon}{\textsubring{r}}{t\smash{\textsuperscript{h}}}{i}{\texttheta}} \dicPos{f}[7] \dicFlx{(‑ar, ‑ir)}[1] \dicDirectTranslationCS{zimní rybářská sezóna}
\dicEntry[vetrungur] \dicTerm{vetr··ung|ur} \dicIPA{{v}{\textepsilon}{\textlengthmark}{\textsubring{d}}{r}{u}{\ng}{\r{g}}{\textscy}{\textsubring{r}}} \dicPos{m}[6] \dicFlx{(‑s, ‑ar)}[8] \dicDirectTranslationCS{jednoroček, jednoroční zvíře}
\dicEntry[vettlingur] \dicTerm{vett··ling|ur} \dicIPA{{v}{\textepsilon}{h}{\textsubring{d}}{l}{i}{\ng}{\r{g}}{\textscy}{\textsubring{r}}} \dicPos{m}[6] \dicFlx{(‑s, ‑ar)}[8] \dicSynonym{hanski} \dicDirectTranslationCS{rukavice, palčáky}
\dicEntry[vettvangskönnun] \dicTerm{vett·vangs··|könnun} \dicIPA{{v}{\textepsilon}{h}{\textsubring{d}}{v}{au}{\ng}{\r{g}}{s}{k\smash{\textsuperscript{h}}}{\oe}{n}{\textscy}{\textsubring{n}}} \dicPos{f}[7] \dicFlx{(‑könn\-un\-ar, ‑kann\-an\-ir)}[11] \dicDirectTranslationCS{průzkum\,/\addthin šetření v~terénu}
\dicEntry[vettvangur] \dicTerm{vett··vang|ur} \dicsymFrequent\  \dicIPA{{v}{\textepsilon}{h}{\textsubring{d}}{v}{au}{\ng}{\r{g}}{\textscy}{\textsubring{r}}} \dicPos{m}[6] \dicFlx{(‑s)}[26] \textbf{1.} \dicDirectTranslationCS{místo, dějiště (zločinu ap.)}  \textbf{2.} \dicDirectTranslationCS{pole, scéna};  \dicPhraseIS{á innlendum\,/\addthin erlendum vettvangi} \dicFlx{adv} \dicDirectTranslationCS{na domácí\,/\addthin zahraniční scéně}
\dicEntry[vetur] \dicTerm{vet|ur} \dicsymFrequent\  \dicIPA{{v}{\textepsilon}{\textlengthmark}{\textsubring{d}}{\textscy}{\textsubring{r}}} \dicPos{m}[13] \dicFlx{(‑rar, ‑ur)}[6] \dicDirectTranslationCS{zima} \dicIndirectTranslationCS{(roční období)};  \dicPhraseIS{á veturna} \dicFlx{adv} \dicDirectTranslationCS{v~zimě, během zimy};  \dicPhraseIS{í fyrra vetur} \dicFlx{adv} \dicDirectTranslationCS{minulou zimu};  \dicPhraseIS{í vetur} \dicFlx{adv} \dicDirectTranslationCS{v~zimě};  \dicPhraseIS{um veturinn} \dicFlx{adv} \dicDirectTranslationCS{v~zimě, přes zimu}
\dicEntry[veturgamall] \dicTerm{vetur··|gamall} \dicIPA{{v}{\textepsilon}{\textsubring{d}}{\textscy}{r}{\r{g}}{a}{m}{a}{\textsubring{d}}{\textsubring{l}}} \dicPos{adj}[11] \dicFlx{(f ‑gömul)}[15] \dicDirectTranslationCS{jednoroční, rok starý, starý jednu zimu}
\dicEntry[veturseta] \dicTerm{vetur··set|a} \dicIPA{{v}{\textepsilon}{\textsubring{d}}{\textscy}{\textsubring{r}}{s}{\textepsilon}{\textsubring{d}}{a}} \dicPos{f}[1] \dicFlx{(‑u, ‑ur)}[7] \dicDirectTranslationCS{zimoviště};  \dicPhraseIS{hafa vetursetu e‑s staðar} \dicDirectTranslationCS{přezimovat (kde)}
\dicEntry[vex] \dicTerm{vex} \dicIPA{{v}{\textepsilon}{x}{s}} \dicPos{v} \dicFlx{ind praes sg 1 pers} \dicLink{vaxa}
\dicEntry[vextir] \dicTerm{vextir\smash{\textsuperscript{1}}} \dicIPA{{v}{\textepsilon}{x}{s}{\textsubring{d}}{\textsci}{\textsubring{r}}} \dicPos{m}[11] \dicFlx{pl}[1] \dicFieldCat{ekon.} \dicDirectTranslationCS{úroky} \dicExampleIS{vextir af láni} \dicExampleCS{úroky z~půjčky}
\dicEntry[vextir] \dicTerm{vextir\smash{\textsuperscript{2}}} \dicIPA{{v}{\textepsilon}{x}{s}{\textsubring{d}}{\textsci}{\textsubring{r}}} \dicPos{m} \dicFlx{pl nom} \dicLink{vöxtur}
\dicEntry[vé] \dicTerm{vé} \dicIPA{{v}{j}{\textepsilon}{\textlengthmark}} \dicPos{n}[2] \dicFlx{(‑s, ‑)}[5] \dicDirectTranslationCS{(pohanská) svatyně}
\dicEntry[véfrétt] \dicTerm{vé··frétt} \dicIPA{{v}{j}{\textepsilon}{\textlengthmark}{f}{r}{j}{\textepsilon}{h}{\textsubring{d}}} \dicPos{f}[7] \dicFlx{(‑ar, ‑ir)}[1] \textbf{1.} \dicDirectTranslationCS{věštec, věštkyně}  \textbf{2.} \dicDirectTranslationCS{věštírna, orákulum} \dicExampleIS{véfréttin í Delfí} \dicExampleCS{Delfská věštírna}  \textbf{3.} \dicDirectTranslationCS{věštba, orákulum}
\dicEntry[vék] \dicTerm{vék} \dicIPA{{v}{j}{\textepsilon}{\textlengthmark}{\r{g}}} \dicPos{v} \dicFlx{ind pf sg 1 pers} \dicLink{víkja}
\dicEntry[vél] \dicTerm{vél\smash{\textsuperscript{1}}} \dicsymFrequent\  \dicIPA{{v}{j}{\textepsilon}{\textlengthmark}{\textsubring{l}}} \dicPos{f}[4] \dicFlx{(‑ar, ‑ar)}[1] \textbf{1.} \dicSynonym{tæki\smash{\textsuperscript{1}}} \dicDirectTranslationCS{stroj, přístroj} \dicExampleIS{Vélin er orðin ónýt.} \dicExampleCS{Stroj už je nepoužitelný.}  \textbf{2.} \dicSynonym{hreyfill} \dicDirectTranslationCS{motor}  \textbf{3.} \dicSynonym{flugvél} \dicDirectTranslationCS{letadlo}
\dicEntry[vél] \dicTerm{vél\smash{\textsuperscript{2}}} \dicIPA{{v}{j}{\textepsilon}{\textlengthmark}{\textsubring{l}}} \dicPos{f}[4] \dicFlx{(‑ar, ‑ar)}[1] \dicSynonym{svik} \dicDirectTranslationCS{pikle, pleticha} \dicIndirectTranslationCS{(zvláště v~pl)};  \dicPhraseIS{draga vélar að e‑m} \dicSynonym{svíkja} \dicDirectTranslationCS{podvést (koho)}
\dicEntry[véla] \dicTerm{vél|a} \dicIPA{{v}{j}{\textepsilon}{\textlengthmark}{l}{a}} \dicPos{v}[1] \dicFlx{(‑aði)}[1] \dicFlx{acc} \dicSynonym{ginna} \dicDirectTranslationCS{(z)manipulovat, (z)lákat};  \dicPhraseIS{láta véla sig til e‑s} \dicDirectTranslationCS{nechat se k~(čemu) zmanipulovat}
\dicEntry[vélamaður] \dicTerm{véla··|maður} \dicIPA{{v}{j}{\textepsilon}{\textlengthmark}{l}{a}{m}{a}{ð}{\textscy}{\textsubring{r}}} \dicPos{m}[13] \dicFlx{(‑manns, ‑menn)}[2] \dicDirectTranslationCS{mechanik, mechanička, strojník, strojnice}
\dicEntry[vélamál] \dicTerm{véla··mál}\dicTerm{, vélarmál} \dicIPA{{v}\-{j}\-{\textepsilon}\-{\textlengthmark}\-{l}\-{a}\-{m}\-{au}\-{\textsubring{l}}\-} \dicPos{n}[2] \dicFlx{(‑s, ‑)}[5] \dicFieldCat{poč.} \dicDirectTranslationCS{strojový kód}
\dicEntry[vélamenning] \dicTerm{véla··menn·ing} \dicIPA{{v}{j}{\textepsilon}{\textlengthmark}{l}{a}{m}{\textepsilon}{n}{i}{\ng}{\r{g}}} \dicPos{f}[4] \dicFlx{(‑ar)}[7] \dicDirectTranslationCS{technická\,/\addthin technologická kultura}
\dicEntry[vélarbilun] \dicTerm{vélar··bil|un} \dicIPA{{v}{j}{\textepsilon}{\textlengthmark}{l}{a}{r}{\textsubring{b}}{\textsci}{l}{\textscy}{\textsubring{n}}} \dicPos{f}[7] \dicFlx{(‑unar, ‑anir)}[8] \dicDirectTranslationCS{porucha stroje\,/\addthin motoru}
\dicEntry[vélarhlíf] \dicTerm{vélar··hlíf} \dicIPA{{v}{j}{\textepsilon}{\textlengthmark}{l}{a}{\textsubring{r}}{\textsubring{l}}{i}{f}} \dicPos{f}[4] \dicFlx{(‑ar, ‑ar)}[1] \dicDirectTranslationCS{kapota, kryt (motoru ap.)}
\dicEntry[vélarmál] \dicTerm{vélar··mál} \dicIPA{{v}{j}{\textepsilon}{\textlengthmark}{l}{a}{r}{m}{au}{\textsubring{l}}} \dicPos{n}[2] \dicFlx{(‑s, ‑)}[5] \dicLink{vélamál}
\dicEntry[vélarmálskóði] \dicTerm{vélar·máls··kóð|i} \dicIPA{{v}{j}{\textepsilon}{\textlengthmark}{l}{a}{r}{m}{au}{l}{s}{k\smash{\textsuperscript{h}}}{ou}{ð}{\textsci}} \dicPos{m}[1] \dicFlx{(‑a, ‑ar)}[1] \dicFieldCat{poč.} \dicDirectTranslationCS{strojový kód}
\dicEntry[vélarrúm] \dicTerm{vélar··rúm} \dicIPA{{v}{j}{\textepsilon}{\textlengthmark}{l}{a}{r}{u}{\textsubring{m}}} \dicPos{n}[2] \dicFlx{(‑s)}[2] \dicDirectTranslationCS{strojovna}
\dicEntry[vélasalur] \dicTerm{véla··sal|ur} \dicIPA{{v}{j}{\textepsilon}{\textlengthmark}{l}{a}{s}{a}{l}{\textscy}{\textsubring{r}}} \dicPos{m}[10] \dicFlx{(‑ar, ‑ir)}[14] \dicDirectTranslationCS{strojovna (v~elektrárně ap.)}
\dicEntry[vélasamstæða] \dicTerm{véla··sam·stæð|a} \dicIPA{{v}{j}{\textepsilon}{\textlengthmark}{l}{a}{s}{a}{m}{s}{\textsubring{d}}{a}{i}{ð}{a}} \dicPos{f}[1] \dicFlx{(‑u, ‑ur)}[13] \dicDirectTranslationCS{(strojní) jednotka, agregát (strojů)}
\dicEntry[vélaverkfræði] \dicTerm{véla··verk·fræð|i} \dicIPA{{v}{j}{\textepsilon}{\textlengthmark}{l}{a}{v}{\textepsilon}{\textsubring{r}}{\r{g}}{f}{r}{a}{i}{ð}{\textsci}} \dicPos{f}[3] \dicFlx{(‑i)}[3] \dicDirectTranslationCS{strojní inženýrství}
\dicEntry[vélaverkfræðingur] \dicTerm{véla··verk·fræð·ing|ur} \dicIPA{{v}\-{j}\-{\textepsilon}\-{\textlengthmark}\-{l}\-{a}\-{v}\-{\textepsilon}\-{\textsubring{r}}\-{\r{g}}\-{f}\-{r}\-{a}\-{i}\-{ð}\-{i}\-{\ng}\-{\r{g}}\-{\textscy}\-{\textsubring{r}}\-} \dicPos{m}[6] \dicFlx{(‑s, ‑ar)}[8] \dicSynonym{vélfræðingur} \dicDirectTranslationCS{strojní inženýr(ka)}
\begin{xtolerant}{}{1pt}
\dicEntry[vélaverkstæði] \dicTerm{véla··verk·stæði}\addthinS\dicIPA{{v}{j}{\textepsilon}{\textlengthmark}{l}{a}{v}{\textepsilon}{\textsubring{r}}{\r{g}}{s}{\textsubring{d}}{a}{i}{ð}{\textsci}}\addthinS\dicPos{n}[2]\addthinS\dicFlx{(‑s, ‑)}[14] \dicDirectTranslationCS{dílna, oprav\-na (motorů ap.)}
\end{xtolerant}
\dicEntry[vélaöld] \dicTerm{véla··|öld} \dicIPA{{v}{j}{\textepsilon}{\textlengthmark}{l}{a}{\oe}{l}{\textsubring{d}}} \dicPos{f}[7] \dicFlx{(‑aldar)}[19] \dicFieldCat{hist.} \dicDirectTranslationCS{století strojů}
\dicEntry[vélbátur] \dicTerm{vél··bát|ur} \dicIPA{{v}{j}{\textepsilon}{l}{\textsubring{b}}{au}{\textsubring{d}}{\textscy}{\textsubring{r}}} \dicPos{m}[6] \dicFlx{(‑s, ‑ar)}[22] \dicFieldCat{nám.} \dicDirectTranslationCS{motorový člun, motorová loď}
\dicEntry[vélbúnaður] \dicTerm{vél··bú·nað|ur} \dicIPA{{v}{j}{\textepsilon}{l}{\textsubring{b}}{u}{n}{a}{ð}{\textscy}{\textsubring{r}}} \dicPos{m}[10] \dicFlx{(‑ar)}[9] \textbf{1.} \dicDirectTranslationCS{strojní vybavení\,/\addthin zařízení}  \textbf{2.} \dicFieldCat{poč.} \dicDirectTranslationCS{hardware} \dicIndirectTranslationCS{(technické vybavení počítače)} \dicAntonym{hugbúnaður}
\dicEntry[vélbyssa] \dicTerm{vél··byss|a} \dicIPA{{v}{j}{\textepsilon}{l}{\textsubring{b}}{\textsci}{s}{a}} \dicPos{f}[1] \dicFlx{(‑u, ‑ur)}[19] \dicFieldCat{voj.} \dicDirectTranslationCS{kulomet, strojní puška}
\dicEntry[vélfræði] \dicTerm{vél··fræð|i} \dicIPA{{v}{j}{\textepsilon}{l}{f}{r}{a}{i}{ð}{\textsci}} \dicPos{f}[3] \dicFlx{(‑i)}[3] \dicDirectTranslationCS{strojnictví}
\dicEntry[vélfræðingur] \dicTerm{vél·fræð··ing|ur} \dicIPA{{v}{j}{\textepsilon}{l}{f}{r}{a}{i}{ð}{i}{\ng}{\r{g}}{\textscy}{\textsubring{r}}} \dicPos{m}[6] \dicFlx{(‑s, ‑ar)}[8] \dicSynonym{vélaverkfræðingur} \dicDirectTranslationCS{strojní inženýr(ka)}
\dicEntry[vélgengur] \dicTerm{vél··gengur} \dicIPA{{v}{j}{\textepsilon}{l}{\r{\textObardotlessj}}{ei}{\ng}{\r{g}}{\textscy}{\textsubring{r}}} \dicPos{adj}[1]\dicFlx{}[-6] \dicSynonym{vélrænn} \dicDirectTranslationCS{mechanický} \dicExampleIS{vélgeng leikföng} \dicExampleCS{mechanické hračky}
\dicEntry[vélhjól] \dicTerm{vél··hjól} \dicIPA{{v}{j}{\textepsilon}{l}{\c{c}}{ou}{\textsubring{l}}} \dicPos{n}[2] \dicFlx{(‑s, ‑)}[5] \dicDirectTranslationCS{motocykl, motorka}
\dicEntry[vélinda] \dicTerm{vélind|a} \dicIPA{{v}{j}{\textepsilon}{\textlengthmark}{l}{\textsci}{n}{\textsubring{d}}{a}} \dicPos{n}[1] \dicFlx{(‑a, ‑u)}[1] \dicFieldCat{anat.} \dicDirectTranslationCS{jícen}
\dicEntry[vélknúinn] \dicTerm{vél··knúinn} \dicIPA{{v}{j}{\textepsilon}{l}{k\smash{\textsuperscript{h}}}{n}{u}{\textsci}{\textsubring{n}}} \dicPos{adj}[6]\dicFlx{}[-6] \dicDirectTranslationCS{poháněný motorem} \dicExampleIS{vélknúinn bátur} \dicExampleCS{člun poháněný motorem}
\dicEntry[vélmenni] \dicTerm{vél··menni} \dicIPA{{v}{j}{\textepsilon}{l}{m}{\textepsilon}{n}{\textsci}} \dicPos{n}[2] \dicFlx{(‑s, ‑)}[14] \dicSynonym{róbot} \dicDirectTranslationCS{robot}
\dicEntry[vélráð] \dicTerm{vél··ráð} \dicIPA{{v}{j}{\textepsilon}{l}{r}{au}{\texttheta}} \dicPos{n}[2] \dicFlx{(‑s, ‑)}[5] \dicLink{vélræði}
\dicEntry[vélrita] \dicTerm{vél··rit|a} \dicIPA{{v}{j}{\textepsilon}{l}{r}{\textsci}{\textsubring{d}}{a}} \dicPos{v}[1] \dicFlx{(‑aði)}[1] \dicFlx{acc} \dicDirectTranslationCS{napsat, psát (na (psacím) stroji\,/\addthin počítači)}
\dicEntry[vélritari] \dicTerm{vél··rit·ar|i} \dicIPA{{v}{j}{\textepsilon}{l}{r}{\textsci}{\textsubring{d}}{a}{r}{\textsci}} \dicPos{m}[1] \dicFlx{(‑a, ‑ar)}[13] \dicDirectTranslationCS{písař(ka) (na (psacím) stroji\,/\addthin počítači)}
\dicEntry[vélritun] \dicTerm{vél··rit|un} \dicIPA{{v}{j}{\textepsilon}{l}{r}{\textsci}{\textsubring{d}}{\textscy}{\textsubring{n}}} \dicPos{f}[7] \dicFlx{(‑unar)}[9] \dicDirectTranslationCS{psaní na (psacím) stroji\,/\addthin počítači}
\dicEntry[vélræði] \dicTerm{vél··ræði}\dicTerm{, vélráð} \dicIPA{{v}\-{j}\-{\textepsilon}\-{l}\-{r}\-{a}\-{i}\-{ð}\-{\textsci}\-} \dicPos{n}[2] \dicFlx{(‑s, ‑)}[14] \dicDirectTranslationCS{podvod, úskok}
\dicEntry[vélrænn] \dicTerm{vél··rænn} \dicIPA{{v}{j}{\textepsilon}{l}{r}{a}{i}{\textsubring{d}}{\textsubring{n}}} \dicPos{adj}[7]\dicFlx{}[-1] \textbf{1.} \dicSynonym*{mekanískur} \dicDirectTranslationCS{strojový, mechanický} \dicExampleIS{vélrænar hreyfingar} \dicExampleCS{mechanické pohyby}  \textbf{2.} \dicSynonym*{gerður með tölvu} \dicDirectTranslationCS{strojový, počítačový} \dicExampleIS{vélrænar þýðingar} \dicExampleCS{strojové překlady}
\dicEntry[vélrænt] \dicTerm{vél··rænt} \dicIPA{{v}{j}{\textepsilon}{l}{r}{a}{i}{\textsubring{n}}{\textsubring{d}}} \dicPos{adv} \dicDirectTranslationCS{strojově, mechanicky}
\dicEntry[vélskófla] \dicTerm{vél··skófl|a} \dicIPA{{v}{j}{\textepsilon}{l}{s}{\r{g}}{ou}{\textsubring{b}}{l}{a}} \dicPos{f}[1] \dicFlx{(‑u, ‑ur)}[19] \dicSynonym{skurðgrafa} \dicDirectTranslationCS{rypadlo, bagr}
\dicEntry[vélskóli] \dicTerm{vél··skól|i} \dicIPA{{v}{j}{\textepsilon}{l}{s}{\r{g}}{ou}{l}{\textsci}} \dicPos{m}[1] \dicFlx{(‑a, ‑ar)}[1] \dicFieldCat{škol.} \dicDirectTranslationCS{(střední) průmyslová škola, strojní průmyslovka}
\dicEntry[vélsleði] \dicTerm{vél··sleð|i} \dicIPA{{v}{j}{\textepsilon}{l}{s}{\textsubring{d}}{l}{\textepsilon}{ð}{\textsci}} \dicPos{m}[1] \dicFlx{(‑a, ‑ar)}[1] \dicDirectTranslationCS{sněžný skútr}
\dicEntry[vélsmiður] \dicTerm{vél··smið|ur} \dicIPA{{v}{j}{\textepsilon}{l}{s}{m}{\textsci}{ð}{\textscy}{\textsubring{r}}} \dicPos{m}[9] \dicFlx{(‑s, ‑ir)}[8] \dicDirectTranslationCS{mechanik, mechanička, montér(ka)}
\dicEntry[vélstjóri] \dicTerm{vél··stjór|i} \dicsymFrequent\  \dicIPA{{v}{j}{\textepsilon}{l}{s}{\textsubring{d}}{j}{ou}{r}{\textsci}} \dicPos{m}[1] \dicFlx{(‑a, ‑ar)}[1] \dicDirectTranslationCS{strojník, strojnice, strojvůdce, strojvůdkyně, strojvedoucí} \dicExampleIS{verkfall vélstjóra} \dicExampleCS{stávka strojníků}
\dicEntry[vélsög] \dicTerm{vél··|sög} \dicIPA{{v}{j}{\textepsilon}{l}{s}{\oe}{x}} \dicPos{f}[7] \dicFlx{(‑sagar, ‑sagir)}[16] \dicDirectTranslationCS{motorová\,/\addthin strojní pila}
\dicEntry[véltækni] \dicTerm{vél··tækn|i} \dicIPA{{v}{j}{\textepsilon}{l}{t\smash{\textsuperscript{h}}}{a}{i}{h}{\r{g}}{n}{\textsci}} \dicPos{f}[3] \dicFlx{(‑i)}[3] \dicDirectTranslationCS{strojní technologie}
\dicEntry[vélvirki] \dicTerm{vél··virk|i} \dicIPA{{v}{j}{\textepsilon}{l}{v}{\textsci}{\textsubring{r}}{\r{\textObardotlessj}}{\textsci}} \dicPos{m}[1] \dicFlx{(‑ja, ‑jar)}[14] \dicSynonym{vélsmiður} \dicDirectTranslationCS{montér(ka), mechanik, mechanička}
\dicEntry[vélvæða] \dicTerm{vél··væ|ða} \dicIPA{{v}{j}{\textepsilon}{l}{v}{a}{i}{ð}{a}} \dicPos{v}[2] \dicFlx{(‑ddi, ‑tt)}[171] \dicFlx{acc} \dicDirectTranslationCS{(z)mechanizovat} \dicExampleIS{vélvæða starfsgrein} \dicExampleCS{zmechanizovat pracovní odvětví};  \dicIdiom{vélvæðast}{ \dicPhraseIS{vélvæðast}} \dicFlx{refl} \dicDirectTranslationCS{být (z)mechanizován}
\dicEntry[vélvæðing] \dicTerm{vél··væð·ing} \dicIPA{{v}{j}{\textepsilon}{l}{v}{a}{i}{ð}{i}{\ng}{\r{g}}} \dicPos{f}[4] \dicFlx{(‑ar)}[7] \dicDirectTranslationCS{mechanizace, mechanizování}
\dicEntry[vér] \dicTerm{vér} \dicIPA{{v}{j}{\textepsilon}{\textlengthmark}{\textsubring{r}}} \dicPos{pron}[17] \dicFlx{pers} \dicFlx{nom pl} \dicLangCat{kniž.} \dicDirectTranslationCS{my}
\dicEntry[vh.] \dicTerm{vh.} \dicPos{zkr} \dicPhraseIS{viðtengingarháttur} \dicFieldCat{jaz.} \dicDirectTranslationCS{spojovací způsob}
\dicEntry[við] \dicTerm{við\smash{\textsuperscript{1}}} \dicIPA{{v}{\textsci}{\textlengthmark}{\texttheta}} \dicPos{pron}[17] \dicFlx{pers} \dicFlx{nom pl} \dicDirectTranslationCS{my} \dicExampleIS{Við ætlum að fara heim.} \dicExampleCS{Chceme jít domů.}
\dicEntry[við] \dicTerm{við\smash{\textsuperscript{2}}} \dicsymFrequent\  \dicIPA{{v}{\textsci}{\textlengthmark}{\texttheta}} \dicPos{prep} \dicFlx{dat\,/\addthin acc} \textbf{1.} \dicFlx{acc} \dicDirectTranslationCS{u} \dicIndirectTranslationCS{(o~věcech nebo místě)} \dicExampleIS{Hann sat við borðið.} \dicExampleCS{Seděl u~stolu.}  \textbf{2.} \dicFlx{dat} \dicDirectTranslationCS{na, k(e), (na)proti} \dicIndirectTranslationCS{(o~směru)} \dicExampleIS{Hún brosti við mér.} \dicExampleCS{Usmála se na mě.}  \textbf{3.} \dicFlx{acc} \dicDirectTranslationCS{s(e)} \dicIndirectTranslationCS{(o~společenství nebo zapojení)} \dicExampleIS{samvinna við lögregluna} \dicExampleCS{spolupráce s~policií}  \textbf{4.} \dicFlx{dat} \dicSynonym{með} \dicDirectTranslationCS{s(e)} \dicExampleIS{éta smjör við brauði} \dicExampleCS{jíst chléb s~máslem}  \textbf{5.} \dicFlx{acc} \dicDirectTranslationCS{s(e), při} \dicIndirectTranslationCS{(o~čase nebo něčem v~časové posloupnosti)} \dicExampleIS{við sólsetur} \dicExampleCS{se západem slunce}  \textbf{6.} \dicFlx{dat} \dicDirectTranslationCS{na} \dicExampleIS{viðbrögðin við gagnrýninni} \dicExampleCS{reakce na kritiku};  \dicIdiom{við}{ \dicPhraseIS{við og við}} \dicFlx{adv} \dicDirectTranslationCS{sem tam, občas}
\dicEntry[viða] \dicTerm{við|a} \dicIPA{{v}{\textsci}{\textlengthmark}{ð}{a}} \dicPos{v}[1] \dicFlx{(‑aði)}[1] \dicPhraseIS{viða að sér e‑u} \dicSynonym{safna} \dicDirectTranslationCS{sbírat (co), shromažďovat (co) (prameny ap.)} \dicExampleIS{viða að sér heimildum í ritgerðina} \dicExampleCS{shromažďovat písemné prameny na esej}
\dicEntry[viðamikill] \dicTerm{viða··|mikill} \dicIPA{{v}{\textsci}{\textlengthmark}{ð}{a}{m}{\textsci}{\r{\textObardotlessj}}{\textsci}{\textsubring{d}}{\textsubring{l}}} \dicPos{adj}[11] \dicFlx{(comp ‑meiri, sup ‑mestur)}[12] \dicSynonym{umfangsmikill} \dicDirectTranslationCS{rozsáhlý, (jsoucí) ve velkém měřítku} \dicExampleIS{viðamikið efni} \dicExampleCS{rozsáhlé téma}
\dicEntry[viðarbolur] \dicTerm{viðar··bol|ur} \dicIPA{{v}{\textsci}{\textlengthmark}{ð}{a}{r}{\textsubring{b}}{\textopeno}{l}{\textscy}{\textsubring{r}}} \dicPos{m}[9] \dicFlx{(‑s, ‑ir)}[8] \dicSynonym{trjábolur} \dicDirectTranslationCS{kmen stromu}
\dicEntry[viðarbútur] \dicTerm{viðar··bút|ur} \dicIPA{{v}{\textsci}{\textlengthmark}{ð}{a}{r}{\textsubring{b}}{u}{\textsubring{d}}{\textscy}{\textsubring{r}}} \dicPos{m}[6] \dicFlx{(‑s, ‑ar)}[15] \dicDirectTranslationCS{špalek, poleno}
\dicEntry[viðarkol] \dicTerm{viðar··kol} \dicIPA{{v}{\textsci}{\textlengthmark}{ð}{a}{\textsubring{r}}{k\smash{\textsuperscript{h}}}{\textopeno}{\textsubring{l}}} \dicPos{n}[2] \dicFlx{pl}[1] \dicDirectTranslationCS{dřevěné uhlí}
\dicEntry[viðarvörn] \dicTerm{viðar··|vörn} \dicIPA{{v}{\textsci}{\textlengthmark}{ð}{a}{r}{v}{\oe}{r}{\textsubring{d}}{\textsubring{n}}} \dicPos{f}[7] \dicFlx{(‑varnar, ‑varnir)}[16] \dicDirectTranslationCS{impregnace dřeva}
\dicEntry[viðauki] \dicTerm{við··auk|i} \dicIPA{{v}{\textsci}{\textlengthmark}{ð}{\oe i}{\r{\textObardotlessj}}{\textsci}} \dicPos{m}[1] \dicFlx{(‑a, ‑ar)}[1] \dicDirectTranslationCS{dodatek, apendix} \dicExampleIS{viðauki bókar} \dicExampleCS{knižní dodatek}
\dicEntry[viðbára] \dicTerm{við··bár|a} \dicIPA{{v}{\textsci}{ð}{\textsubring{b}}{au}{r}{a}} \dicPos{f}[1] \dicFlx{(‑u, ‑ur)}[7] \dicDirectTranslationCS{výmluva, vytáčka}
\dicEntry[viðbein] \dicTerm{við··bein} \dicIPA{{v}{\textsci}{ð}{\textsubring{b}}{ei}{\textsubring{n}}} \dicPos{n}[2] \dicFlx{(‑s, ‑)}[5] \dicFieldCat{anat.} \dicDirectTranslationCS{klíční kost}
\dicEntry[viðbeinsbrot] \dicTerm{við·beins··brot} \dicIPA{{v}{\textsci}{ð}{\textsubring{b}}{ei}{n}{s}{\textsubring{b}}{r}{\textopeno}{\textsubring{d}}} \dicPos{n}[2] \dicFlx{(‑s, ‑)}[5] \dicDirectTranslationCS{zlomenina klíční kosti}
\dicEntry[viðbjóðslega] \dicTerm{við·bjóðs··lega} \dicIPA{{v}{\textsci}{ð}{\textsubring{b}}{j}{ou}{ð}{s}{l}{\textepsilon}{\textbabygamma}{a}} \dicPos{adv} \dicDirectTranslationCS{odporně, nechutně, hnusně}
\dicEntry[viðbjóðslegur] \dicTerm{við·bjóðs··legur} \dicIPA{{v}{\textsci}{ð}{\textsubring{b}}{j}{ou}{ð}{s}{l}{\textepsilon}{\textbabygamma}{\textscy}{\textsubring{r}}} \dicPos{adj}[1]\dicFlx{}[-8] \dicDirectTranslationCS{odporný, nechutný, hnusný} \dicExampleIS{viðbjóðslegur óþefur} \dicExampleCS{odporný zápach}
\dicEntry[viðbjóður] \textls[15]{\dicTerm{við··bjóð|ur} \dicIPA{{v}{\textsci}{ð}{\textsubring{b}}{j}{ou}{ð}{\textscy}{\textsubring{r}}} \dicPos{m}[9] \dicFlx{(‑s, ‑ir)}[8] \dicDirectTranslationCS{odpor, hnus, nechutnost};  \dicPhraseIS{hafa viðbjóð á e‑u} \dicDirectTranslationCS{mít odpor k~(čemu), \mbox{hnusit si (co)}}}
\dicEntry[viðbót] \dicTerm{við··|bót} \dicsymFrequent\  \dicIPA{{v}{\textsci}{ð}{\textsubring{b}}{ou}{\textsubring{d}}} \dicPos{f}[8] \dicFlx{(‑bótar, ‑bætur)}[5] \dicSynonym{auki} \dicDirectTranslationCS{přídavek, nádavek};  \dicPhraseIS{í viðbót} \dicFlx{adv} \dicDirectTranslationCS{navíc, navrch, nadto};  \dicPhraseIS{til viðbótar} \dicFlx{adv} \dicDirectTranslationCS{navíc, navrch, nádavkem}
\dicEntry[viðbótarkostnaður] \dicTerm{við·bótar··kost·nað|ur} \dicIPA{{v}\-{\textsci}\-{ð}\-{\textsubring{b}}\-{ou}\-{\textsubring{d}}\-{a}\-{\textsubring{r}}\-{k\smash{\textsuperscript{h}}}\-{\textopeno}\-{s}\-{\textsubring{d}}\-{n}\-{a}\-{ð}\-{\textscy}\-{\textsubring{r}}\-} \dicPos{m}[10] \dicFlx{(‑ar)}[9] \dicDirectTranslationCS{dodatečné náklady}
\dicEntry[viðbótarlífeyrissparnaður] \dicTerm{við·bótar··líf·eyris·spar·nað|ur} \dicIPA{{v}\-{\textsci}\-{ð}\-{\textsubring{b}}\-{ou}\-{\textsubring{d}}\-{a}\-{r}\-{l}\-{i}\-{v}\-{ei}\-{r}\-{\textsci}\-{s}\-{\textsubring{b}}\-{a}\-{r}\-{\textsubring{d}}\-{n}\-{a}\-{ð}\-{\textscy}\-{\textsubring{r}}\-} \dicPos{m}[10] \dicFlx{(‑ar)}[9] \dicFieldCat{ekon.} \dicDirectTranslationCS{penzijní připojištění}
\dicEntry[viðbragð] \dicTerm{við··|bragð} \dicsymFrequent\  \dicIPA{{v}{\textsci}{ð}{\textsubring{b}}{r}{a}{\textbabygamma}{\texttheta}} \dicPos{n}[2] \dicFlx{(‑bragðs, ‑brögð)}[8] \textbf{1.} \dicSynonym*{snöggur kippur} \dicDirectTranslationCS{reakce, reagování} \dicExampleIS{snöggt viðbragð} \dicExampleCS{pohotová reakce}  \textbf{2.} \dicSynonym{andsvar} \dicDirectTranslationCS{reakce, odezva} \dicIndirectTranslationCS{(zvláště v~pl)}
\dicEntry[viðbragðstími] \dicTerm{við·bragðs··tím|i} \dicIPA{{v}{\textsci}{ð}{\textsubring{b}}{r}{a}{\textbabygamma}{ð}{s}{t\smash{\textsuperscript{h}}}{i}{m}{\textsci}} \dicPos{m}[1] \dicFlx{(‑a)}[3] \dicDirectTranslationCS{reakční čas\,/\addthin doba}
\dicEntry[viðbrenndur] \dicTerm{við··brenndur} \dicIPA{{v}{\textsci}{ð}{\textsubring{b}}{r}{\textepsilon}{n}{\textsubring{d}}{\textscy}{\textsubring{r}}} \dicPos{adj}[2]\dicFlx{}[-14] \dicSynonym{sangur} \dicDirectTranslationCS{připálený, spálený (omáčka ap.)}
\dicEntry[viðbrigði] \dicTerm{við··brigði} \dicIPA{{v}{\textsci}{ð}{\textsubring{b}}{r}{\textsci}{\textbabygamma}{ð}{\textsci}} \dicPos{n}[2] \dicFlx{pl}[19] \dicDirectTranslationCS{změna, obrat}
\dicEntry[viðbrunninn] \dicTerm{við··brunninn} \dicIPA{{v}{\textsci}{ð}{\textsubring{b}}{r}{\textscy}{n}{\textsci}{\textsubring{n}}} \dicPos{adj}[6]\dicFlx{}[-15] \dicSynonym{sangur} \dicDirectTranslationCS{připálený, spálený}
\dicEntry[viðbrögð] \dicTerm{við··brögð} \dicIPA{{v}{\textsci}{ð}{\textsubring{b}}{r}{\oe}{\textbabygamma}{\texttheta}} \dicPos{n}[2] \dicFlx{pl}[9] \dicLink{viðbragð}
\dicEntry[viðbundinn] \textls[15]{\dicTerm{við··bundinn} \dicIPA{{v}{\textsci}{ð}{\textsubring{b}}{\textscy}{n}{\textsubring{d}}{\textsci}{\textsubring{n}}} \dicPos{adj}[6]\dicFlx{}[-6] \dicSynonym{upptekinn}} \dicDirectTranslationCS{zaneprázdněný}
\dicEntry[viðburðaríkur] \dicTerm{við·burða··ríkur} \dicIPA{{v}{\textsci}{ð}{\textsubring{b}}{\textscy}{r}{ð}{a}{r}{i}{\r{g}}{\textscy}{\textsubring{r}}} \dicPos{adj}[1]\dicFlx{}[-1] \dicDirectTranslationCS{bohatý na události, rušný (den ap.)}
\dicEntry[viðburður] \dicTerm{við··burð|ur} \dicIPA{{v}{\textsci}{ð}{\textsubring{b}}{\textscy}{r}{ð}{\textscy}{\textsubring{r}}} \dicPos{m}[10] \dicFlx{(‑ar, ‑ir)}[4] \textbf{1.} \dicSynonym{atburður} \dicDirectTranslationCS{událost, příhoda} \dicExampleIS{Þetta er heimssögulegur viðburður.} \dicExampleCS{Toto je historická událost.}  \textbf{2.} \dicSynonym*{e‑að óvanalegt} \dicDirectTranslationCS{vzácnost, rarita}
\dicEntry[viðbúinn] \dicTerm{við··búinn} \dicIPA{{v}{\textsci}{ð}{\textsubring{b}}{u}{\textsci}{\textsubring{n}}} \dicPos{adj}[6]\dicFlx{}[-2] \dicSynonym{tilbúinn} \dicDirectTranslationCS{připravený, přichystaný, nachystaný};  \dicPhraseIS{e‑að er viðbúið} \dicDirectTranslationCS{(co) se dá očekávat, (co) je pravděpodobné};  \dicPhraseIS{vera viðbúinn e‑u} \dicDirectTranslationCS{být připravený na (co), počítat s~(čím)}
\dicEntry[viðbúnaður] \dicTerm{við··bú·nað|ur} \dicIPA{{v}{\textsci}{ð}{\textsubring{b}}{u}{n}{a}{ð}{\textscy}{\textsubring{r}}} \dicPos{m}[10] \dicFlx{(‑ar)}[9] \dicSynonym{undirbúningur} \dicDirectTranslationCS{příprava, přichystání}
\dicEntry[viðbygging] \dicTerm{við··bygg·ing} \dicIPA{{v}{\textsci}{ð}{\textsubring{b}}{\textsci}{\r{\textObardotlessj}}{i}{\ng}{\r{g}}} \dicPos{f}[4] \dicFlx{(‑ar, ‑ar)}[5] \dicDirectTranslationCS{nástavba, přístavba}
\dicEntry[viðbætir] \dicTerm{við··bæt|ir} \dicIPA{{v}{\textsci}{ð}{\textsubring{b}}{a}{i}{\textsubring{d}}{\textsci}{\textsubring{r}}} \dicPos{m}[7] \dicFlx{(‑is, ‑ar)}[1] \dicSynonym{viðbót} \dicDirectTranslationCS{dodatek, apendix}
\dicEntry[viðdvöl] \dicTerm{við··|dvöl} \dicIPA{{v}{\textsci}{ð}{\textsubring{d}}{v}{\oe}{\textsubring{l}}} \dicPos{f}[7] \dicFlx{(‑dvalar, ‑dvalir)}[16] \dicSynonym{dvöl} \dicDirectTranslationCS{zastávka, zastavení} \dicExampleIS{stutt viðdvöl í Svíþjóð} \dicExampleCS{krátká zastávka ve Švédsku};  \dicPhraseIS{hafa viðdvöl e‑s staðar} \dicDirectTranslationCS{zastavit se (kde), udělat si (kde) zastávku}
\dicEntry[viðeigandi] \dicTerm{við··eig·andi} \dicsymFrequent\  \dicIPA{{v}{\textsci}{\textlengthmark}{ð}{ei}{\textbabygamma}{a}{n}{\textsubring{d}}{\textsci}} \dicPos{adj}[13] \dicFlx{indecl}[1] \dicDirectTranslationCS{odpovídající, příslušný, patřičný} \dicExampleIS{Setjið orðin á viðeigandi stað á myndina.} \dicExampleCS{Umístěte slova na patřičné místo na obrázku.}
\dicEntry[viðfangsefni] \dicTerm{við·fangs··efni} \dicsymFrequent\  \dicIPA{{v}{\textsci}{ð}{f}{au}{\ng}{s}{\textepsilon}{\textsubring{b}}{n}{\textsci}} \dicPos{n}[2] \dicFlx{(‑s, ‑)}[14] \textbf{1.} \dicSynonym{verkefni} \dicDirectTranslationCS{zadání, úloha, úkol} \dicExampleIS{mjög flókið viðfangsefni} \dicExampleCS{velmi složité zadání}  \textbf{2.} \dicSynonym{þema} \dicDirectTranslationCS{námět, téma, tematika}
\dicEntry[viðfelldinn] \dicTerm{við··felldinn} \dicIPA{{v}{\textsci}{ð}{f}{\textepsilon}{l}{\textsubring{d}}{\textsci}{\textsubring{n}}} \dicPos{adj}[6]\dicFlx{}[-2] \dicSynonym{geðfelldur} \dicDirectTranslationCS{přívětivý, příjemný, sympatický}
\dicEntry[viðgangur] \dicTerm{við··gang|ur} \dicIPA{{v}{\textsci}{ð}{\r{g}}{au}{\ng}{\r{g}}{\textscy}{\textsubring{r}}} \dicPos{m}[6] \dicFlx{(‑s)}[7] \dicSynonym{framgangur} \dicDirectTranslationCS{rozvoj, rozvíjení se} \dicExampleIS{tryggja viðgang verslunar} \dicExampleCS{zajistit rozvoj obchodu}
\dicEntry[viðgerð] \dicTerm{við··gerð} \dicIPA{{v}{\textsci}{ð}{\r{\textObardotlessj}}{\textepsilon}{r}{\texttheta}} \dicPos{f}[7] \dicFlx{(‑ar, ‑ir)}[1] \dicDirectTranslationCS{oprava, opravení, náprava (poškozené věci ap.)} \dicExampleIS{viðgerð á e‑u} \dicExampleCS{oprava (čeho)};  \dicPhraseIS{e‑að fer í viðgerð} \dicDirectTranslationCS{(co) jde na opravu}
\dicEntry[viðgerðarkostnaður] \dicTerm{við·gerðar··kost·nað|ur} \dicIPA{{v}\-{\textsci}\-{ð}\-{\r{\textObardotlessj}}\-{\textepsilon}\-{r}\-{ð}\-{a}\-{\textsubring{r}}\-{k\smash{\textsuperscript{h}}}\-{\textopeno}\-{s}\-{\textsubring{d}}\-{n}\-{a}\-{ð}\-{\textscy}\-{\textsubring{r}}\-} \dicPos{m}[10] \dicFlx{(‑ar)}[9] \dicDirectTranslationCS{náklady na opravu}
\dicEntry[viðgerðarmaður] \dicTerm{við·gerðar··|maður} \dicIPA{{v}{\textsci}{ð}{\r{\textObardotlessj}}{\textepsilon}{r}{ð}{a}{r}{m}{a}{ð}{\textscy}{\textsubring{r}}} \dicPos{m}[13] \dicFlx{(‑manns, ‑menn)}[2] \dicDirectTranslationCS{opravář(ka), mechanik, mechanička}
\dicEntry[viðgerðarþjónusta] \dicTerm{við·gerðar··þjón·ust|a} \dicIPA{{v}\-{\textsci}\-{ð}\-{\r{\textObardotlessj}}\-{\textepsilon}\-{r}\-{ð}\-{a}\-{\textsubring{r}}\-{\texttheta}\-{j}\-{ou}\-{n}\-{\textscy}\-{s}\-{\textsubring{d}}\-{a}\-} \dicPos{f}[1] \dicFlx{(‑u)}[5] \dicDirectTranslationCS{servis, servisní stanice, opravna}
\dicEntry[viðhafa] \dicTerm{við··|hafa} \dicIPA{{v}{\textsci}{\textlengthmark}{\texttheta}{h}{a}{v}{a}} \dicPos{v}[4] \dicFlx{(‑hef, ‑hafði, ‑höfðum, ‑hefði, ‑haft)}[18] \dicFlx{acc} \dicDirectTranslationCS{(po)užít, (po)užívat} \dicExampleIS{viðhafa ljót orð} \dicExampleCS{používat sprostá slova}
\dicEntry[viðhafði] \dicTerm{við··hafði} \dicIPA{{v}{\textsci}{\textlengthmark}{\texttheta}{h}{a}{v}{ð}{\textsci}} \dicPos{v} \dicFlx{ind pf sg 1 pers} \dicLink{viðhafa}
\dicEntry[viðhafnarútgáfa] \dicTerm{við·hafnar··út·gáf|a} \dicIPA{{v}{\textsci}{\textlengthmark}{\texttheta}{h}{a}{\textsubring{b}}{n}{a}{r}{u}{\textsubring{d}}{\r{g}}{au}{v}{a}} \dicPos{f}[1] \dicFlx{(‑u, ‑ur)}[13] \dicDirectTranslationCS{luxusní edice}
\dicEntry[viðhaft] \dicTerm{við··haft} \dicIPA{{v}{\textsci}{\textlengthmark}{\texttheta}{h}{a}{f}{\textsubring{d}}} \dicPos{v} \dicFlx{supin} \dicLink{viðhafa}
\dicEntry[viðhald] \dicTerm{við··|hald} \dicIPA{{v}{\textsci}{\textlengthmark}{\texttheta}{h}{a}{l}{\textsubring{d}}} \dicPos{n}[2] \dicFlx{(‑halds, ‑höld)}[8] \textbf{1.} \dicSynonym{umhirðing} \dicDirectTranslationCS{údržba, udržování, servis} \dicExampleIS{viðhald húsa} \dicExampleCS{údržba domů}  \textbf{2.} \dicSynonym*{elskhugi\,/\addthin ástkona} \dicDirectTranslationCS{milenec, milenka} \dicExampleIS{Hún átti viðhald úti í bæ.} \dicExampleCS{Měla ve městě milence.}
\dicEntry[viðhalda] \dicTerm{við··|halda} \dicIPA{{v}{\textsci}{\textlengthmark}{\texttheta}{h}{a}{l}{\textsubring{d}}{a}} \dicPos{v}[6] \dicFlx{(‑held, ‑hélt, ‑héldum, ‑héldi, ‑haldið)}[59] \dicFlx{dat} \dicSynonym*{halda við} \dicDirectTranslationCS{udržet, udržovat, zachovat, zachovávat} \dicExampleIS{viðhalda sambandi við hann} \dicExampleCS{udržovat s~ním kontakt}
\dicEntry[viðhaldið] \dicTerm{við··haldið} \dicIPA{{v}{\textsci}{\textlengthmark}{\texttheta}{h}{a}{l}{\textsubring{d}}{\textsci}{\texttheta}} \dicPos{v} \dicFlx{supin} \dicLink{viðhalda}
\dicEntry[viðhaldskostnaður] \dicTerm{við·halds··kost·nað|ur} \dicIPA{{v}\-{\textsci}\-{\textlengthmark}\-{\texttheta}\-{h}\-{a}\-{l}\-{\textsubring{d}}\-{s}\-{k\smash{\textsuperscript{h}}}\-{\textopeno}\-{s}\-{\textsubring{d}}\-{n}\-{a}\-{ð}\-{\textscy}\-{\textsubring{r}}\-} \dicPos{m}[10] \dicFlx{(‑ar)}[9] \dicDirectTranslationCS{náklady na údržbu}
\dicEntry[viðhef] \dicTerm{við··hef} \dicIPA{{v}{\textsci}{\textlengthmark}{\texttheta}{h}{\textepsilon}{f}} \dicPos{v} \dicFlx{ind praes sg 1 pers} \dicLink{viðhafa}
\dicEntry[viðhefði] \dicTerm{við··hefði} \dicIPA{{v}{\textsci}{\textlengthmark}{\texttheta}{h}{\textepsilon}{v}{ð}{\textsci}} \dicPos{v} \dicFlx{con pf sg 1 pers} \dicLink{viðhafa}
\dicEntry[viðheld] \dicTerm{við··held} \dicIPA{{v}{\textsci}{\textlengthmark}{\texttheta}{h}{\textepsilon}{l}{\textsubring{d}}} \dicPos{v} \dicFlx{ind praes sg 1 pers} \dicLink{viðhalda}
\dicEntry[viðhengi] \dicTerm{við··hengi} \dicIPA{{v}{\textsci}{\textlengthmark}{\texttheta}{h}{ei}{\textltailn}{\r{\textObardotlessj}}{\textsci}} \dicPos{n}[2] \dicFlx{(‑s, ‑)}[16] \textbf{1.} \dicFieldCat{poč.} \dicSynonym{fylgiskjal} \dicDirectTranslationCS{příloha (k~e‑mailu ap.)}  \textbf{2.} \dicDirectTranslationCS{přívěsek, přívěšek}
\dicEntry[viðhéldi] \dicTerm{við··héldi} \dicIPA{{v}{\textsci}{\textlengthmark}{\texttheta}{h}{j}{\textepsilon}{l}{\textsubring{d}}{\textsci}} \dicPos{v} \dicFlx{con pf sg 1 pers} \dicLink{viðhalda}
\dicEntry[viðhéldum] \dicTerm{við··héldum} \dicIPA{{v}{\textsci}{\textlengthmark}{\texttheta}{h}{j}{\textepsilon}{l}{\textsubring{d}}{\textscy}{\textsubring{m}}} \dicPos{v} \dicFlx{ind pf pl 1 pers} \dicLink{viðhalda}
\dicEntry[viðhélt] \dicTerm{við··hélt} \dicIPA{{v}{\textsci}{\textlengthmark}{\texttheta}{h}{j}{\textepsilon}{\textsubring{l}}{\textsubring{d}}} \dicPos{v} \dicFlx{ind pf sg 1 pers} \dicLink{viðhalda}
\dicEntry[viðhlítandi] \dicTerm{við··hlít·andi} \dicIPA{{v}{\textsci}{ð}{\textsubring{l}}{i}{\textsubring{d}}{a}{n}{\textsubring{d}}{\textsci}} \dicPos{adj}[13] \dicFlx{indecl}[1] \textbf{1.} \dicSynonym{fullnægjandi} \dicDirectTranslationCS{dostatečný, dostačující, postačující} \dicExampleIS{viðhlítandi árangur} \dicExampleCS{dostatečný výsledek}  \textbf{2.} \dicSynonym{viðeigandi} \dicDirectTranslationCS{příslušný, patřičný}
\dicEntry[viðhorf] \dicTerm{við··horf} \dicsymFrequent\  \dicIPA{{v}{\textsci}{\textlengthmark}{\texttheta}{h}{\textopeno}{r}{f}} \dicPos{n}[2] \dicFlx{(‑s, ‑)}[5] \dicSynonym{afstaða} \dicDirectTranslationCS{hledisko, stanovisko, postoj};  \dicPhraseIS{viðhorf e‑rs til e‑s} \dicDirectTranslationCS{(čí) postoj k~(čemu)}
\dicEntry[viðhöfðum] \dicTerm{við··höfðum} \dicIPA{{v}{\textsci}{\textlengthmark}{\texttheta}{h}{\oe}{v}{ð}{\textscy}{\textsubring{m}}} \dicPos{v} \dicFlx{ind pf pl 1 pers} \dicLink{viðhafa}
\dicEntry[viðhöfn] \dicTerm{við··|höfn} \dicIPA{{v}{\textsci}{\textlengthmark}{\texttheta}{h}{\oe}{\textsubring{b}}{\textsubring{n}}} \dicPos{f}[7] \dicFlx{(‑hafnar)}[19] \dicDirectTranslationCS{okázalost, honosnost, pompa} \dicExampleIS{Brúðkaupið fór fram með mikilli viðhöfn.} \dicExampleCS{Svatba proběhla s~velkou pompou.}
\dicEntry[viðkem] \dicTerm{við··kem} \dicIPA{{v}{\textsci}{ð}{k\smash{\textsuperscript{h}}}{\textepsilon}{\textsubring{m}}} \dicPos{v} \dicFlx{ind praes sg 1 pers} \dicLink{viðkoma\smash{\textsuperscript{2}}}
\dicEntry[viðkom] \dicTerm{við··kom} \dicIPA{{v}{\textsci}{ð}{k\smash{\textsuperscript{h}}}{\textopeno}{\textsubring{m}}} \dicPos{v} \dicFlx{ind pf sg 1 pers} \dicLink{viðkoma\smash{\textsuperscript{2}}}
\dicEntry[viðkoma] \dicTerm{við··kom|a\smash{\textsuperscript{1}}} \dicIPA{{v}{\textsci}{ð}{k\smash{\textsuperscript{h}}}{\textopeno}{m}{a}} \dicPos{f}[1] \dicFlx{(‑u)}[5] \textbf{1.} \dicSynonym{stans} \dicDirectTranslationCS{zastávka, zastavení (na cestě ap.)} \dicExampleIS{hafa viðkomu í borginni} \dicExampleCS{zastavit se ve městě}  \textbf{2.} \dicSynonym{snerting} \dicDirectTranslationCS{dotek, dotyk, dotknutí};  \dicPhraseIS{e‑að er mjúkt\,/\addthin hart viðkomu} \dicDirectTranslationCS{(co) je měkké\,/\addthin tvrdé na dotek} \dicExampleIS{Teppið er mjúkt viðkomu.} \dicExampleCS{Koberec je měkký na dotek.}  \textbf{3.} \dicSynonym{fjölgun} \dicDirectTranslationCS{plodnost, fertilita}
\dicEntry[viðkoma] \dicTerm{við··|koma\smash{\textsuperscript{2}}} \dicIPA{{v}{\textsci}{ð}{k\smash{\textsuperscript{h}}}{\textopeno}{m}{a}} \dicPos{v}[7] \dicFlx{(‑kem, ‑kom, ‑komum, ‑kæmi, ‑kom\-ið)}[8] \dicFlx{dat} \dicSynonym{varða\smash{\textsuperscript{2}}} \dicDirectTranslationCS{týkat se, souviset (s)};  \dicPhraseIS{hvað viðkemur e‑u} \dicDirectTranslationCS{co se (čeho) týká, co souvisí s~(čím)}
\dicEntry[viðkomandi] \dicTerm{við··kom·|andi\smash{\textsuperscript{1}}} \dicIPA{{v}{\textsci}{ð}{k\smash{\textsuperscript{h}}}{\textopeno}{m}{a}{n}{\textsubring{d}}{\textsci}} \dicPos{m}[2] \dicFlx{(‑anda, ‑endur)}[1] \dicDirectTranslationCS{dotyčný, dotyčná} \dicIndirectTranslationCS{(dříve uvedená osoba)} \dicExampleIS{Viðkomandi var ekki heima.} \dicExampleCS{Dotyčný nebyl doma.}
\dicEntry[viðkomandi] \dicTerm{við··kom·andi\smash{\textsuperscript{2}}} \dicIPA{{v}{\textsci}{ð}{k\smash{\textsuperscript{h}}}{\textopeno}{m}{a}{n}{\textsubring{d}}{\textsci}} \dicPos{adj}[13] \dicFlx{indecl}[1] \textbf{1.} \dicDirectTranslationCS{týkající se, relevantní}  \textbf{2.} \dicDirectTranslationCS{dotyčný, příslušný}
\dicEntry[viðkomið] \dicTerm{við··komið} \dicIPA{{v}{\textsci}{ð}{k\smash{\textsuperscript{h}}}{\textopeno}{m}{\textsci}{\texttheta}} \dicPos{v} \dicFlx{supin} \dicLink{viðkoma\smash{\textsuperscript{2}}}
\dicEntry[viðkomum] \dicTerm{við··komum} \dicIPA{{v}{\textsci}{ð}{k\smash{\textsuperscript{h}}}{\textopeno}{m}{\textscy}{\textsubring{m}}} \dicPos{v} \dicFlx{ind pf pl 1 pers} \dicLink{viðkoma\smash{\textsuperscript{2}}}
\dicEntry[viðkomustaður] \dicTerm{við·komu··stað|ur} \dicIPA{{v}{\textsci}{ð}{k\smash{\textsuperscript{h}}}{\textopeno}{m}{\textscy}{s}{\textsubring{d}}{a}{ð}{\textscy}{\textsubring{r}}} \dicPos{m}[10] \dicFlx{(‑ar, ‑ir)}[14] \dicDirectTranslationCS{zastávka, zastavení (na cestě ap.)}
\dicEntry[viðkunnanlegur] \dicTerm{við·kunnan··legur} \dicIPA{{v}{\textsci}{ð}{k\smash{\textsuperscript{h}}}{\textscy}{n}{a}{n}{l}{\textepsilon}{\textbabygamma}{\textscy}{\textsubring{r}}} \dicPos{adj}[1]\dicFlx{}[-8] \dicSynonym{geðþekkur} \dicDirectTranslationCS{sympatický, milý, líbivý} \dicExampleIS{viðkunnanlegt bros} \dicExampleCS{milý úsměv}
\dicEntry[viðkvæði] \dicTerm{við··kvæði} \dicIPA{{v}{\textsci}{ð}{k\smash{\textsuperscript{h}}}{v}{a}{i}{ð}{\textsci}} \dicPos{n}[2] \dicFlx{(‑s, ‑)}[14] \textbf{1.} \dicFieldCat{lit.} \dicSynonym*{stef í kvæði} \dicDirectTranslationCS{refrén, opakovaný verš}  \textbf{2.} \dicLangCat{přen.} \dicDirectTranslationCS{ohraná písnička, ta samá písnička} \dicExampleIS{Það var alltaf sama viðkvæðið.} \dicExampleCS{Byla to pořád ta samá písnička.}
\dicEntry[viðkvæmni] \dicTerm{við··kvæmn|i} \dicIPA{{v}{\textsci}{ð}{k\smash{\textsuperscript{h}}}{v}{a}{i}{m}{n}{\textsci}} \dicPos{f}[3] \dicFlx{(‑i)}[3] \textbf{1.} \dicSynonym{tilfinninganæmi} \dicDirectTranslationCS{citlivost, senzibilita}  \textbf{2.} \dicSynonym*{líkamleg viðkvæmni} \dicDirectTranslationCS{citlivost, choulostivost, přecitlivělost}
\dicEntry[viðkvæmur] \dicTerm{við··kvæmur} \dicsymFrequent\  \dicIPA{{v}{\textsci}{ð}{k\smash{\textsuperscript{h}}}{v}{a}{i}{m}{\textscy}{\textsubring{r}}} \dicPos{adj}[1]\dicFlx{}[-1] \textbf{1.} \dicSynonym{tilfinninganæmur} \dicDirectTranslationCS{citlivý, senzibilní}  \textbf{2.} \dicSynonym{næmur} \dicDirectTranslationCS{citlivý, choulostivý, přecitlivělý};  \dicPhraseIS{vera viðkvæmur fyrir e‑u} \dicDirectTranslationCS{být přecitlivělý na (co) (na infekci ap.)}  \textbf{3.} \dicSynonym{vandmeðfarinn} \dicDirectTranslationCS{citlivý, choulostivý, ožehavý (téma ap.)} \dicExampleIS{viðkvæmt mál} \dicExampleCS{ožehavé téma}
\dicEntry[viðkynning] \textls[10]{\dicTerm{við··kynn·ing} \dicIPA{{v}{\textsci}{ð}{c\smash{\textsuperscript{h}}}{\textsci}{n}{i}{\ng}{\r{g}}} \dicPos{f}[4] \dicFlx{(‑ar)}[7] \dicSynonym{kunningsskapur} \dicDirectTranslationCS{zná\-most, kontakt, vztah} \dicExampleIS{viðkynning hennar við nágrannana} \dicExampleCS{její známosti se sousedy}}
\dicEntry[viðkæmi] \dicTerm{við··kæmi} \dicIPA{{v}{\textsci}{ð}{c\smash{\textsuperscript{h}}}{a}{i}{m}{\textsci}} \dicPos{v} \dicFlx{con pf sg 1 pers} \dicLink{viðkoma\smash{\textsuperscript{2}}}
\dicEntry[viðlag] \dicTerm{við··|lag} \dicIPA{{v}{\textsci}{ð}{l}{a}{x}} \dicPos{n}[2] \dicFlx{(‑lags, ‑lög)}[8] \dicFieldCat{lit.} \dicSynonym{stef} \dicDirectTranslationCS{refrén}
\dicEntry[viðlagasjóður] \dicTerm{við·laga··sjóð|ur} \dicIPA{{v}{\textsci}{ð}{l}{a}{\textbabygamma}{a}{s}{j}{ou}{ð}{\textscy}{\textsubring{r}}} \dicPos{m}[9] \dicFlx{(‑s, ‑ir)}[6] \dicFieldCat{ekon.} \dicDirectTranslationCS{podpůrný fond}
\dicEntry[viðlátinn] \dicTerm{við··látinn} \dicIPA{{v}{\textsci}{ð}{l}{au}{\textsubring{d}}{\textsci}{\textsubring{n}}} \dicPos{adj}[6]\dicFlx{}[-6] \dicSynonym{viðstaddur} \dicDirectTranslationCS{přítomný, (jsoucí) k~zastižení} \dicExampleIS{vera viðlátinn á skrifstofunni} \dicExampleCS{být k~zastižení v~kanceláři}
\dicEntry[viðleitni] \dicTerm{við··leitn|i} \dicIPA{{v}{\textsci}{ð}{l}{ei}{h}{\textsubring{d}}{n}{\textsci}} \dicPos{f}[3] \dicFlx{(‑i)}[3] \dicSynonym{tilraun} \dicDirectTranslationCS{snaha, úsilí};  \dicPhraseIS{sýna viðleitni til e‑s} \dicDirectTranslationCS{snažit se o~(co), vykazovat snahu v~(čem)}
\dicEntry[viðlit] \dicTerm{við··lit} \dicIPA{{v}{\textsci}{ð}{l}{\textsci}{\textsubring{d}}} \dicPos{n}[2] \dicFlx{(‑s)}[2] \dicPhraseIS{virða e‑n ekki viðlits} \dicDirectTranslationCS{nevěnovat (komu) pozornost, ignorovat (koho)};  \dicPhraseIS{það er ekki viðlit að (gera e‑ð)} \dicDirectTranslationCS{nepřichází v~úvahu ((co) udělat)}
\dicEntry[viðlíka] \dicTerm{við··líka\smash{\textsuperscript{1}}} \dicIPA{{v}{\textsci}{ð}{l}{i}{\r{g}}{a}} \dicPos{adj}[13] \dicFlx{indecl}[1] \dicSynonym{ámóta\smash{\textsuperscript{1}}} \dicDirectTranslationCS{podobný, obdobný}
\dicEntry[viðlíka] \dicTerm{við··líka\smash{\textsuperscript{2}}} \dicIPA{{v}{\textsci}{ð}{l}{i}{\r{g}}{a}} \dicPos{adv} \dicDirectTranslationCS{podobně, obdobně} \dicExampleIS{viðlíka fljótur} \dicExampleCS{podobně rychlý}
\dicEntry[viðlíking] \dicTerm{við··lík·ing} \dicIPA{{v}{\textsci}{ð}{l}{i}{\r{\textObardotlessj}}{i}{\ng}{\r{g}}} \dicPos{f}[4] \dicFlx{(‑ar, ‑ar)}[5] \dicFieldCat{lit.} \dicDirectTranslationCS{přirovnání}
\dicEntry[viðloðandi] \dicTerm{við··loð·andi} \dicIPA{{v}{\textsci}{ð}{l}{\textopeno}{ð}{a}{n}{\textsubring{d}}{\textsci}} \dicPos{adj}[13] \dicFlx{indecl}[1] \textbf{1.} \dicDirectTranslationCS{přilnavý, adhezivní}  \textbf{2.} \dicSynonym{viðriðinn} \dicDirectTranslationCS{zapojený, zainteresovaný, zaangažovaný}  \textbf{3.} \dicDirectTranslationCS{přináležející, držící se (myši domu ap.)}
\dicEntry[viðloðun] \dicTerm{við··loð|un} \dicIPA{{v}{\textsci}{ð}{l}{\textopeno}{ð}{\textscy}{\textsubring{n}}} \dicPos{f}[7] \dicFlx{(‑unar)}[9] \dicDirectTranslationCS{přilnavost, adheze}
\dicEntry[viðmið] \dicTerm{við··mið} \dicIPA{{v}{\textsci}{ð}{m}{\textsci}{\texttheta}} \dicPos{n}[2] \dicFlx{(‑s, ‑)}[5] \dicDirectTranslationCS{(srovnávací) měřítko, kritérium, norma}
\dicEntry[viðmiðun] \dicTerm{við··mið|un} \dicIPA{{v}{\textsci}{ð}{m}{\textsci}{ð}{\textscy}{\textsubring{n}}} \dicPos{f}[7] \dicFlx{(‑unar, ‑anir)}[8] \dicDirectTranslationCS{měřítko, kritérium} \dicExampleIS{til viðmiðunar} \dicExampleCS{jako kritérium}
\dicEntry[viðmót] \dicTerm{við··mót} \dicIPA{{v}{\textsci}{ð}{m}{ou}{\textsubring{d}}} \dicPos{n}[2] \dicFlx{(‑s)}[2] \textbf{1.} \dicSynonym{framkoma} \dicDirectTranslationCS{chování, vystupování} \dicExampleIS{þægilegur í viðmóti} \dicExampleCS{s~příjemným chováním}  \textbf{2.} \dicFieldCat{poč.} \dicDirectTranslationCS{uživatelské rozhraní}
\dicEntry[viðmótsþýður] \dicTerm{við·móts··þýður} \dicIPA{{v}{\textsci}{ð}{m}{ou}{\textsubring{d}}{s}{\texttheta}{i}{ð}{\textscy}{\textsubring{r}}} \dicPos{adj}[2]\dicFlx{}[-6] \dicSynonym{notalegur} \dicDirectTranslationCS{příjemný, milý, zdvořilý}
\dicEntry[viðmælandi] \dicTerm{við··mæl·|andi\smash{\textsuperscript{1}}} \dicIPA{{v}{\textsci}{ð}{m}{a}{i}{l}{a}{n}{\textsubring{d}}{\textsci}} \dicPos{m}[2] \dicFlx{(‑anda, ‑endur)}[1] \dicDirectTranslationCS{dotazovaný, dotazovaná} \dicIndirectTranslationCS{(člověk, kterému byla položena otázka)} \dicAntonym{mælandi}
\dicEntry[viðmælandi] \textls[10]{\dicTerm{við··mæl·andi\smash{\textsuperscript{2}}} \dicIPA{{v}{\textsci}{ð}{m}{a}{i}{l}{a}{n}{\textsubring{d}}{\textsci}}} \dicPos{adj}[13] \dicFlx{indecl}[1] \dicDirectTranslationCS{přítomný (k~rozmluvě)} \dicExampleIS{Hann er ekki viðmælandi.} \dicExampleCS{Nelze s~ním mluvit.}
\dicEntry[viðnám] \dicTerm{við··nám} \dicIPA{{v}{\textsci}{ð}{n}{au}{\textsubring{m}}} \dicPos{n}[2] \dicFlx{(‑s)}[2] \textbf{1.} \dicSynonym{mótspyrna} \dicDirectTranslationCS{odpor, vzdor};  \dicPhraseIS{veita e‑u viðnám} \dicDirectTranslationCS{klást (čemu) odpor};  \dicPhraseIS{viðnám gegn e‑u} \dicDirectTranslationCS{odpor vůči (čemu)}  \textbf{2.} \dicFieldCat{elek.} \dicSynonym*{rafmagnsviðnám} \dicDirectTranslationCS{odpor, rezistor}
\dicEntry[viðnámsþrek] \dicTerm{við·náms··þrek} \dicIPA{{v}{\textsci}{ð}{n}{au}{m}{s}{\texttheta}{r}{\textepsilon}{\r{g}}} \dicPos{n}[2] \dicFlx{(‑s)}[2] \dicDirectTranslationCS{odolnost}
\dicEntry[viðra] \dicTerm{viðr|a} \dicIPA{{v}{\textsci}{ð}{r}{a}} \dicPos{v}[1] \dicFlx{(‑aði)}[7] \dicFlx{acc} \textbf{1.} \dicDirectTranslationCS{(vy)větrat, provětrat} \dicExampleIS{viðra sængurfötin} \dicExampleCS{vyvětrat peřiny};  \dicPhraseIS{viðra hundinn} \dicDirectTranslationCS{(vy)venčit psa};  \dicPhraseIS{viðra sig} \dicDirectTranslationCS{provětrat se, vyvětrat se} \dicIndirectTranslationCS{(vyjít ven na čerstvý vzduch)}  \textbf{2.} \dicSynonym{kynna} \dicDirectTranslationCS{ventilovat, prezentovat} \dicExampleIS{viðra skoðanir sínar} \dicExampleCS{ventilovat své názory}  \textbf{3.} \dicPhraseIS{það viðrar vel\,/\addthin illa} \dicFlx{impers} \dicDirectTranslationCS{je hezky\,/\addthin ošklivo} \dicIndirectTranslationCS{(o~počasí)};  \dicIdiom{viðra}[sig upp]{ \dicPhraseIS{viðra sig upp við e‑n}} \dicDirectTranslationCS{vetřít se ke (komu)}
\dicEntry[viðráðanlegur] \dicTerm{við·ráðan··legur} \dicIPA{{v}{\textsci}{ð}{r}{au}{ð}{a}{n}{l}{\textepsilon}{\textbabygamma}{\textscy}{\textsubring{r}}} \dicPos{adj}[1]\dicFlx{}[-8] \dicDirectTranslationCS{zvládnutelný, kontrolovatelný (smích ap.)}
\dicEntry[viðreisn] \dicTerm{við··reisn} \dicIPA{{v}{\textsci}{ð}{r}{ei}{s}{\textsubring{d}}{\textsubring{n}}} \dicPos{f}[7] \dicFlx{(‑ar)}[3] \dicSynonym{endurreisn} \dicDirectTranslationCS{zotavení, opětovný vzestup};  \dicPhraseIS{eiga sér ekki viðreisnar von} \dicLangCat{přen.} \dicDirectTranslationCS{být v~bezvýchodné situaci}
\dicEntry[viðriðinn] \dicTerm{við··riðinn} \dicIPA{{v}{\textsci}{ð}{r}{\textsci}{ð}{\textsci}{\textsubring{n}}} \dicPos{adj}[6]\dicFlx{}[-6] \dicPhraseIS{vera viðriðinn e‑ð} \dicDirectTranslationCS{být do (čeho) zapojený, účastnit se (čeho)}
\dicEntry[viðrini] \dicTerm{viðrini} \dicIPA{{v}{\textsci}{ð}{r}{\textsci}{n}{\textsci}} \dicPos{n}[2] \dicFlx{(‑s, ‑)}[14] \textbf{1.} \dicFieldCat{biol.} \dicSynonym*{tvíkynja skepna} \dicDirectTranslationCS{hermafrodit}  \textbf{2.} \dicDirectTranslationCS{podivín(ka), bizarní osoba}
\dicEntry[viðrinislegur] \dicTerm{viðrinis··legur} \dicIPA{{v}{\textsci}{ð}{r}{\textsci}{n}{\textsci}{s}{l}{\textepsilon}{\textbabygamma}{\textscy}{\textsubring{r}}} \dicPos{adj}[1]\dicFlx{}[-8] \dicSynonym{afkáralegur} \dicDirectTranslationCS{podivínský, bizarní}
\dicEntry[viðræða] \dicTerm{við··ræð|a} \dicIPA{{v}{\textsci}{ð}{r}{a}{i}{ð}{a}} \dicPos{f}[1] \dicFlx{(‑u, ‑ur)}[19] \textbf{1.} \dicSynonym*{samræða} \dicDirectTranslationCS{rozmluva, konverzace, dialog}  \textbf{2.} \dicPhraseIS{viðræður} \dicFlx{pl} \dicDirectTranslationCS{vyjednávání, domlouvání}
\dicEntry[viðsemjandi] \dicTerm{við··semj·|andi} \dicIPA{{v}{\textsci}{ð}{s}{\textepsilon}{m}{j}{a}{n}{\textsubring{d}}{\textsci}} \dicPos{m}[2] \dicFlx{(‑anda, ‑endur)}[1] \dicDirectTranslationCS{smluvní strana}
\dicEntry[viðsjá] \dicTerm{við··sjá} \dicIPA{{v}{\textsci}{ð}{s}{j}{au}} \dicPos{f}[4] \dicFlx{(‑r, ‑r)}[18] \textbf{1.} \dicSynonym{varúð} \dicDirectTranslationCS{obezřetnost, opatrnost, ostražitost}  \textbf{2.} \dicPhraseIS{viðsjár} \dicFlx{pl} \dicSynonym*{deilur} \dicDirectTranslationCS{rozepře, spory, konflikty} \dicExampleIS{Voru miklar viðsjár með mönnum.} \dicExampleCS{Mezi lidmi bylo hodně konfliktů.}
\dicEntry[viðsjáll] \dicTerm{við··sjáll} \dicIPA{{v}{\textsci}{ð}{s}{j}{au}{\textsubring{d}}{\textsubring{l}}} \dicPos{adj}[8]\dicFlx{}[-1] \dicSynonym{varasamur} \dicDirectTranslationCS{záludný, ošemetný, ošidný}
\dicEntry[viðsjárverður] \dicTerm{við·sjár··verður} \dicIPA{{v}{\textsci}{ð}{s}{j}{au}{r}{v}{\textepsilon}{r}{ð}{\textscy}{\textsubring{r}}} \dicPos{adj}[2]\dicFlx{}[-1] \dicSynonym{varhugaverður} \dicDirectTranslationCS{záludný, zrádný, ošemetný}
\dicEntry[viðsk.] \dicTerm{viðsk.} \dicPos{zkr} \dicPhraseIS{viðskeyti} \dicFieldCat{jaz.} \dicDirectTranslationCS{přípona}
\dicEntry[viðskeytamál] \dicTerm{við·skeyta··mál} \dicIPA{{v}{\textsci}{ð}{s}{\r{\textObardotlessj}}{ei}{\textsubring{d}}{a}{m}{au}{\textsubring{l}}} \dicPos{n}[2] \dicFlx{(‑s, ‑)}[5] \dicFieldCat{jaz.} \dicDirectTranslationCS{aglutinační jazyk}
\dicEntry[viðskeyti] \dicTerm{við··skeyti} \dicIPA{{v}{\textsci}{ð}{s}{\r{\textObardotlessj}}{ei}{\textsubring{d}}{\textsci}} \dicPos{n}[2] \dicFlx{(‑s, ‑)}[14] \dicFieldCat{jaz.} \dicDirectTranslationCS{přípona, sufix}
\dicEntry[viðskeyttur] \dicTerm{við··skeyttur} \dicIPA{{v}{\textsci}{ð}{s}{\r{\textObardotlessj}}{ei}{h}{\textsubring{d}}{\textscy}{\textsubring{r}}} \dicPos{adj}[1]\dicFlx{}[-13] \dicFieldCat{jaz.} \dicDirectTranslationCS{připojený};  \dicPhraseIS{viðskeyttur greinir} \dicFieldCat{jaz.} \dicDirectTranslationCS{připojený člen}
\dicEntry[viðskila] \dicTerm{við··skila} \dicIPA{{v}{\textsci}{ð}{s}{\r{\textObardotlessj}}{\textsci}{l}{a}} \dicPos{adj}[13] \dicFlx{indecl}[1] \dicDirectTranslationCS{odloučený, separovaný};  \dicPhraseIS{verða viðskila við e‑n} \dicDirectTranslationCS{být odloučený od (koho)}
\dicEntry[viðskilnaður] \dicTerm{við··skil·nað|ur} \dicIPA{{v}{\textsci}{ð}{s}{\r{\textObardotlessj}}{\textsci}{l}{n}{a}{ð}{\textscy}{\textsubring{r}}} \dicPos{m}[10] \dicFlx{(‑ar)}[7] \textbf{1.} \dicSynonym{brottför} \dicDirectTranslationCS{opuštění, odchod}  \textbf{2.} \dicSynonym{aðskilnaður} \dicDirectTranslationCS{odloučení, separace, rozluka}  \textbf{3.} \dicSynonym{fráfall} \dicDirectTranslationCS{odchod, skon, úmrtí}
\begin{xtolerant}{}{1pt}
\dicEntry[viðskiptabanki] \dicTerm{við·skipta··bank|i}\addthinS\dicIPA{{v}{\textsci}{ð}{s}{\r{\textObardotlessj}}{\textsci}{f}{\textsubring{d}}{a}{\textsubring{b}}{au}{\r{\textltailn}}{\r{\textObardotlessj}}{\textsci}}\addthinS\dicPos{m}[1]\addthinS\dicFlx{(‑a, ‑ar)}[8] \dicDirectTranslationCS{obchodní\,/\addthin komerční banka}
\end{xtolerant}
\dicEntry[viðskiptabann] \dicTerm{við·skipta··bann} \dicIPA{{v}{\textsci}{ð}{s}{\r{\textObardotlessj}}{\textsci}{f}{\textsubring{d}}{a}{\textsubring{b}}{a}{\textsubring{n}}} \dicPos{n}[2] \dicFlx{(‑s)}[2] \dicDirectTranslationCS{obchodní embargo}
\dicEntry[viðskiptabréf] \dicTerm{við·skipta··bréf} \dicIPA{{v}{\textsci}{ð}{s}{\r{\textObardotlessj}}{\textsci}{f}{\textsubring{d}}{a}{\textsubring{b}}{r}{j}{\textepsilon}{f}} \dicPos{n}[2] \dicFlx{(‑s, ‑)}[5] \textbf{1.} \dicSynonym{verðbréf} \dicDirectTranslationCS{cenný papír}  \textbf{2.} \dicSynonym{sendibréf} \dicDirectTranslationCS{obchodní dopis}
\dicEntry[viðskiptadeild] \dicTerm{við·skipta··deild} \dicIPA{{v}{\textsci}{ð}{s}{\r{\textObardotlessj}}{\textsci}{f}{\textsubring{d}}{a}{\textsubring{d}}{ei}{l}{\textsubring{d}}} \dicPos{f}[7] \dicFlx{(‑ar, ‑ir)}[1] \dicFieldCat{škol.} \dicDirectTranslationCS{obchodní fakulta}
\dicEntry[viðskiptafrelsi] \dicTerm{við·skipta··frelsi} \dicIPA{{v}{\textsci}{ð}{s}{\r{\textObardotlessj}}{\textsci}{f}{\textsubring{d}}{a}{f}{r}{\textepsilon}{l}{s}{\textsci}} \dicPos{n}[2] \dicFlx{(‑s)}[20] \dicDirectTranslationCS{svoboda obchodování}
\dicEntry[viðskiptafræði] \dicTerm{við·skipta··fræð|i} \dicIPA{{v}{\textsci}{ð}{s}{\r{\textObardotlessj}}{\textsci}{f}{\textsubring{d}}{a}{f}{r}{a}{i}{ð}{\textsci}} \dicPos{f}[3] \dicFlx{(‑i)}[3] \dicSynonym{hagfræði} \dicDirectTranslationCS{obchodní věda, ekonomie}
\dicEntry[viðskiptafræðingur] \dicTerm{við·skipta·fræð··ing|ur} \dicIPA{{v}\-{\textsci}\-{ð}\-{s}\-{\r{\textObardotlessj}}\-{\textsci}\-{f}\-{\textsubring{d}}\-{a}\-{f}\-{r}\-{a}\-{i}\-{ð}\-{i}\-{\ng}\-{\r{g}}\-{\textscy}\-{\textsubring{r}}\-} \dicPos{m}[6] \dicFlx{(‑s, ‑ar)}[8] \dicDirectTranslationCS{ekonom(ka)}
\dicEntry[viðskiptahagsmunir] \dicTerm{við·skipta··hags·munir} \dicIPA{{v}\-{\textsci}\-{ð}\-{s}\-{\r{\textObardotlessj}}\-{\textsci}\-{f}\-{\textsubring{d}}\-{a}\-{h}\-{a}\-{x}\-{s}\-{m}\-{\textscy}\-{n}\-{\textsci}\-{\textsubring{r}}\-} \dicPos{m}[10] \dicFlx{pl}[1] \dicDirectTranslationCS{obchodní zájmy}
\dicEntry[viðskiptahalli] \dicTerm{við·skipta··hall|i} \dicIPA{{v}{\textsci}{ð}{s}{\r{\textObardotlessj}}{\textsci}{f}{\textsubring{d}}{a}{h}{a}{\textsubring{d}}{l}{\textsci}} \dicPos{m}[1] \dicFlx{(‑a)}[3] \dicFieldCat{ekon.} \dicDirectTranslationCS{deficit obchodní bilance}
\dicEntry[viðskiptaháskóli] \dicTerm{við·skipta··há·skól|i} \dicIPA{{v}\-{\textsci}\-{ð}\-{s}\-{\r{\textObardotlessj}}\-{\textsci}\-{f}\-{\textsubring{d}}\-{a}\-{h}\-{au}\-{s}\-{\r{g}}\-{ou}\-{l}\-{\textsci}\-} \dicPos{m}[1] \dicFlx{(‑a, ‑ar)}[1] \dicFieldCat{škol.} \dicDirectTranslationCS{vysoká škola ekonomická}
\dicEntry[viðskiptahættir] \dicTerm{við·skipta··hættir} \dicIPA{{v}{\textsci}{ð}{s}{\r{\textObardotlessj}}{\textsci}{f}{\textsubring{d}}{a}{h}{a}{i}{h}{\textsubring{d}}{\textsci}{\textsubring{r}}} \dicPos{m}[12] \dicFlx{pl}[6] \dicDirectTranslationCS{obchodní praktiky}
\dicEntry[viðskiptajöfnuður] \dicTerm{við·skipta··|jöfn·uður} \dicIPA{{v}\-{\textsci}\-{ð}\-{s}\-{\r{\textObardotlessj}}\-{\textsci}\-{f}\-{\textsubring{d}}\-{a}\-{j}\-{\oe}\-{\textsubring{b}}\-{n}\-{\textscy}\-{ð}\-{\textscy}\-{\textsubring{r}}\-} \dicPos{m}[10] \dicFlx{(‑jafnaðar\,/\addthin ‑jöfnuðar)}[43] \dicFieldCat{ekon.} \dicDirectTranslationCS{obchodní bilance}
\dicEntry[viðskiptakjör] \dicTerm{við·skipta··kjör} \dicIPA{{v}{\textsci}{ð}{s}{\r{\textObardotlessj}}{\textsci}{f}{\textsubring{d}}{a}{c\smash{\textsuperscript{h}}}{\oe}{\textsubring{r}}} \dicPos{n}[2] \dicFlx{pl}[9] \dicFieldCat{ekon.} \dicDirectTranslationCS{obchodní podmínky}
\dicEntry[viðskiptamaður] \dicTerm{við·skipta··|maður} \dicIPA{{v}{\textsci}{ð}{s}{\r{\textObardotlessj}}{\textsci}{f}{\textsubring{d}}{a}{m}{a}{ð}{\textscy}{\textsubring{r}}} \dicPos{m}[13] \dicFlx{(‑manns, ‑menn)}[2] \textbf{1.} \dicSynonym{viðskiptavinur} \dicDirectTranslationCS{zákazník, zákaznice}  \textbf{2.} \dicSynonym*{sem stundar viðskipti} \dicDirectTranslationCS{obchodník, obchodnice}
\dicEntry[viðskiptaráðherra] \dicTerm{við·skipta··ráð·herr|a} \dicIPA{{v}\-{\textsci}\-{ð}\-{s}\-{\r{\textObardotlessj}}\-{\textsci}\-{f}\-{\textsubring{d}}\-{a}\-{r}\-{au}\-{\texttheta}\-{h}\-{\textepsilon}\-{r}\-{a}\-} \dicPos{m}[1] \dicFlx{(‑a, ‑ar)}[17] \dicDirectTranslationCS{ministr(yně) průmyslu a~obchodu}
\dicEntry[viðskiptaráðuneyti] \dicTerm{við·skipta··ráðu·neyti} \dicIPA{{v}\-{\textsci}\-{ð}\-{s}\-{\r{\textObardotlessj}}\-{\textsci}\-{f}\-{\textsubring{d}}\-{a}\-{r}\-{au}\-{ð}\-{\textscy}\-{n}\-{ei}\-{\textsubring{d}}\-{\textsci}\-} \dicPos{n}[2] \dicFlx{(‑s, ‑)}[14] \dicDirectTranslationCS{ministerstvo průmyslu a~obchodu}
\dicEntry[viðskiptareikningur] \textls[15]{\dicTerm{við·skipta··reikn·ing|ur} \dicIPA{{v}\-{\textsci}\-{ð}\-{s}\-{\r{\textObardotlessj}}\-{\textsci}\-{f}\-{\textsubring{d}}\-{a}\-{r}\-{ei}\-{h}\-{\r{g}}\-{n}\-{i}\-{\ng}\-{\r{g}}\-{\textscy}\-{\textsubring{r}}}} \dicPos{m}[6] \dicFlx{(‑s, ‑ar)}[8] \dicFieldCat{ekon.} \dicDirectTranslationCS{kreditní\,/\addthin úvěrový účet}
\dicEntry[viðskiptasamband] \textls[15]{\dicTerm{við·skipta··sam·|band} \dicIPA{{v}\-{\textsci}\-{ð}\-{s}\-{\r{\textObardotlessj}}\-{\textsci}\-{f}\-{\textsubring{d}}\-{a}\-{s}\-{a}\-{m}\-{\textsubring{b}}\-{a}\-{n}\-{\textsubring{d}}\-} \dicPos{n}[2] \dicFlx{(‑bands, ‑bönd)}[8]} \dicDirectTranslationCS{obchodní styk\,/\addthin kontakt}
\dicEntry[viðskiptasamningur] \dicTerm{við·skipta··samn·ing|ur} \dicIPA{{v}\-{\textsci}\-{ð}\-{s}\-{\r{\textObardotlessj}}\-{\textsci}\-{f}\-{\textsubring{d}}\-{a}\-{s}\-{a}\-{m}\-{n}\-{i}\-{\ng}\-{\r{g}}\-{\textscy}\-{\textsubring{r}}\-} \dicPos{m}[6] \dicFlx{(‑s, ‑ar)}[8] \dicFieldCat{ekon.} \dicDirectTranslationCS{obchodní smlouva}
\dicEntry[viðskiptaskrá] \dicTerm{við·skipta··skrá} \dicIPA{{v}{\textsci}{ð}{s}{\r{\textObardotlessj}}{\textsci}{f}{\textsubring{d}}{a}{s}{\r{g}}{r}{au}} \dicPos{f}[4] \dicFlx{(‑r\,/\addthin ‑ar, ‑r)}[21] \dicFieldCat{ekon.} \dicDirectTranslationCS{obchodní rejstřík}
\dicEntry[viðskiptavild] \dicTerm{við·skipta··vild} \dicIPA{{v}{\textsci}{ð}{s}{\r{\textObardotlessj}}{\textsci}{f}{\textsubring{d}}{a}{v}{\textsci}{l}{\textsubring{d}}} \dicPos{f}[4] \dicFlx{(‑ar)}[3] \dicFieldCat{ekon.} \dicDirectTranslationCS{renomé, reputace, dobré jméno}
\dicEntry[viðskiptavinur] \textls[10]{\dicTerm{við·skipta··vin|ur} \dicsymFrequent\  \dicIPA{{v}{\textsci}{ð}{s}{\r{\textObardotlessj}}{\textsci}{f}{\textsubring{d}}{a}{v}{\textsci}{n}{\textscy}{\textsubring{r}}} \dicPos{m}[10] \dicFlx{(‑ar, ‑ir)}[17] \textbf{1.} \dicDirectTranslationCS{klient(ka), zákazník, zákaznice}  \textbf{2.} \dicPhraseIS{viðskiptavinir} \dicFlx{pl} \dicDirectTranslationCS{klientela, zákazníci} \dicExampleIS{eiga samskipti við viðskiptavini} \dicExampleCS{spolupracovat se zákazníky}}
\dicEntry[viðskipti] \dicTerm{við··skipti} \dicsymFrequent\  \dicIPA{{v}{\textsci}{ð}{s}{\r{\textObardotlessj}}{\textsci}{f}{\textsubring{d}}{\textsci}} \dicPos{n}[2] \dicFlx{(‑s, ‑)}[14] \dicPhraseIS{viðskipti} \dicFlx{pl} \dicDirectTranslationCS{obchod, byznys, komerce, obchodování} \dicExampleIS{viðskipti með fisk} \dicExampleCS{obchodování s~rybami};  \dicPhraseIS{eiga viðskipti við e‑n} \dicDirectTranslationCS{obchodovat s~(kým)} \dicExampleIS{eiga viðskipti við Íslendinga} \dicExampleCS{obchodovat s~Islanďany};  \dicIdiom{viðskipti}{ \dicPhraseIS{vera (erfiður) viðskiptis}} \dicDirectTranslationCS{být (těžký) k~vypořádání (problém ap.)}
\dicEntry[viðskotaillur] \dicTerm{við·skota··|illur} \dicIPA{{v}{\textsci}{ð}{s}{\r{g}}{\textopeno}{\textsubring{d}}{a}{\textsci}{\textsubring{d}}{l}{\textscy}{\textsubring{r}}} \dicPos{adj}[11] \dicFlx{(comp ‑verri, sup ‑verst\-ur)}[7] \dicSynonym{önugur} \dicDirectTranslationCS{popudlivý, prchlivý, cholerický}
\dicEntry[viðspyrna] \dicTerm{við··spyrn|a} \dicIPA{{v}{\textsci}{ð}{s}{\textsubring{b}}{\textsci}{r}{\textsubring{d}}{n}{a}} \dicPos{f}[1] \dicFlx{(‑u, ‑ur)}[7] \textbf{1.} \dicDirectTranslationCS{opora, pevné místo} \dicExampleIS{góð viðspyrna} \dicExampleCS{dobrá opora}  \textbf{2.} \dicSynonym{mótstaða} \dicDirectTranslationCS{odpor};  \dicPhraseIS{veita e‑u viðspyrnu} \dicDirectTranslationCS{klást (čemu) odpor, odporovat (čemu)}
\dicEntry[viðstaddur] \dicTerm{við··|staddur} \dicsymFrequent\  \dicIPA{{v}{\textsci}{ð}{s}{\textsubring{d}}{a}{\textsubring{d}}{\textscy}{\textsubring{r}}} \dicPos{adj}[2] \dicFlx{(f ‑stödd)}[20] \dicDirectTranslationCS{přítomný, zúčastněný} \dicExampleIS{kenna öllum viðstöddum að dansa} \dicExampleCS{učit všechny přítomné tancovat}
\dicEntry[viðstaða] \dicTerm{við··|staða} \dicIPA{{v}{\textsci}{ð}{s}{\textsubring{d}}{a}{ð}{a}} \dicPos{f}[1] \dicFlx{(‑stöðu)}[2] \textbf{1.} \dicSynonym{viðnám} \dicDirectTranslationCS{odpor, rezistence} \dicExampleIS{Herinn veitti viðstöðu í virkjum við landamærin.} \dicExampleCS{Armáda kladla odpor v~pevnostech při hranici.}  \textbf{2.} \dicSynonym{dvöl} \dicDirectTranslationCS{zastávka, zastavení} \dicExampleIS{hafa viðstöðu þar} \dicExampleCS{udělat si tam zastávku}
\dicEntry[viðstöðulaus] \dicTerm{við·stöðu··laus} \dicIPA{{v}{\textsci}{ð}{s}{\textsubring{d}}{\oe}{ð}{\textscy}{l}{\oe i}{s}} \dicPos{adj}[5]\dicFlx{}[-1] \dicDirectTranslationCS{ustavičný, nepřetržitý}
\dicEntry[viðstöðulaust] \dicTerm{við·stöðu··laust} \dicIPA{{v}{\textsci}{ð}{s}{\textsubring{d}}{\oe}{ð}{\textscy}{l}{\oe i}{s}{\textsubring{d}}} \dicPos{adv} \dicDirectTranslationCS{ustavičně, nepřetržitě}
\dicEntry[viðtaka] \dicTerm{við··|taka} \dicIPA{{v}{\textsci}{ð}{t\smash{\textsuperscript{h}}}{a}{\r{g}}{a}} \dicPos{f}[1] \dicFlx{(‑töku, ‑tökur)}[20] \textbf{1.} \dicSynonym{móttaka} \dicDirectTranslationCS{přijetí, převzetí} \dicExampleIS{veita e‑u viðtöku} \dicExampleCS{převzít (co)}  \textbf{2.} \dicPhraseIS{viðtökur} \dicFlx{pl} \dicDirectTranslationCS{přijetí, uznání, akceptace} \dicExampleIS{fá góðar viðtökur} \dicExampleCS{být dobře přijat}
\dicEntry[viðtakandi] \dicTerm{við··tak·|andi} \dicIPA{{v}{\textsci}{ð}{t\smash{\textsuperscript{h}}}{a}{\r{g}}{a}{n}{\textsubring{d}}{\textsci}} \dicPos{m}[2] \dicFlx{(‑anda, ‑endur)}[1] \textbf{1.} \dicSynonym{móttakandi} \dicDirectTranslationCS{příjemce, příjemkyně}  \textbf{2.} \dicDirectTranslationCS{adresát(ka)} \dicExampleIS{sendandi og viðtakandi} \dicExampleCS{odesílatel a~adresát}
\dicEntry[viðtal] \dicTerm{við··|tal} \dicsymFrequent\  \dicIPA{{v}{\textsci}{ð}{t\smash{\textsuperscript{h}}}{a}{\textsubring{l}}} \dicPos{n}[2] \dicFlx{(‑tals, ‑töl)}[8] \textbf{1.} \dicSynonym{blaðaviðtal} \dicDirectTranslationCS{(formální) rozhovor, pohovor, interview} \dicExampleIS{taka viðtal við e‑n} \dicExampleCS{udělat s~(kým) interview}  \textbf{2.} \dicSynonym*{viðtal við sérfræðing} \dicDirectTranslationCS{(odborná) porada, konzultace}  \textbf{3.} \dicSynonym{áheyrn} \dicDirectTranslationCS{slyšení, audience}  \textbf{4.} \dicSynonym{samtal} \dicDirectTranslationCS{rozhovor, rozmluva}
\dicEntry[viðtalstími] \dicTerm{við·tals··tím|i} \dicIPA{{v}{\textsci}{ð}{t\smash{\textsuperscript{h}}}{a}{l}{s}{t\smash{\textsuperscript{h}}}{i}{m}{\textsci}} \dicPos{m}[1] \dicFlx{(‑a, ‑ar)}[1] \textbf{1.} \dicSynonym*{viðtalstími forstjóra} \dicDirectTranslationCS{konzultační hodiny}  \textbf{2.} \dicSynonym*{viðtalstími læknis} \dicDirectTranslationCS{ordinační hodiny}
\dicEntry[viðtekinn] \dicTerm{við··tekinn} \dicIPA{{v}{\textsci}{ð}{t\smash{\textsuperscript{h}}}{\textepsilon}{\r{\textObardotlessj}}{\textsci}{\textsubring{n}}} \dicPos{adj}[6]\dicFlx{}[-2] \dicSynonym{hefðbundinn} \dicDirectTranslationCS{přijímaný, (obecně) přijatý} \dicExampleIS{viðtekin venja} \dicExampleCS{zvyklost}
\begin{xtolerant}{}{1pt}
\dicEntry[viðtengingarháttur] \dicTerm{við·tengingar··|háttur} \dicIPA{{v}\-{\textsci}\-{ð}\-{t\smash{\textsuperscript{h}}}\-{ei}\-{\textltailn}\-{\r{\textObardotlessj}}\-{i}\-{\ng}\-{\r{g}}\-{a}\-{\textsubring{r}}\-{h}\-{au}\-{h}\-{\textsubring{d}}\-{\textscy}\-{\textsubring{r}}}\addthinS\dicPos{m}[12]\addthinS\dicFlx{(‑hátt\-ar, ‑hættir)}[7] \dicFieldCat{jaz.} \dicDirectTranslationCS{spojovací způsob, konjunktiv}
\end{xtolerant}
\dicEntry[viðtæki] \dicTerm{við··tæki} \dicIPA{{v}{\textsci}{ð}{t\smash{\textsuperscript{h}}}{a}{i}{\r{\textObardotlessj}}{\textsci}} \dicPos{n}[2] \dicFlx{(‑s, ‑)}[16] \dicDirectTranslationCS{(rozhlasový) přijímač}
\dicEntry[viðtökur] \dicTerm{við··tökur} \dicIPA{{v}{\textsci}{ð}{t\smash{\textsuperscript{h}}}{\oe}{\r{g}}{\textscy}{\textsubring{r}}} \dicPos{f} \dicFlx{pl} \dicLink{viðtaka}
\dicEntry[viðunandi] \textls[15]{\dicTerm{við··un·andi} \dicIPA{{v}{\textsci}{\textlengthmark}{ð}{\textscy}{n}{a}{n}{\textsubring{d}}{\textsci}} \dicPos{adj}[13] \dicFlx{indecl}[1] \dicSynonym{þolanlegur}} \dicDirectTranslationCS{uspokojivý, přijatelný, akceptovatelný} \dicExampleIS{viðunandi verð} \dicExampleCS{přijatelná cena}
\dicEntry[viðundur] \dicTerm{við··undur} \dicIPA{{v}{\textsci}{\textlengthmark}{ð}{\textscy}{n}{\textsubring{d}}{\textscy}{\textsubring{r}}} \dicPos{n}[2] \dicFlx{(‑s, ‑)}[25] \dicSynonym{undur\smash{\textsuperscript{1}}} \dicDirectTranslationCS{(po)divnost};  \dicPhraseIS{gera sig að viðundri} \dicLangCat{přen.} \dicDirectTranslationCS{dělat ze sebe šaška}
\dicEntry[viður] \dicTerm{við|ur} \dicsymFrequent\  \dicIPA{{v}{\textsci}{\textlengthmark}{ð}{\textscy}{\textsubring{r}}} \dicPos{m}[10] \dicFlx{(‑ar, ‑ir)}[19] \textbf{1.} \dicSynonym{tré} \dicDirectTranslationCS{strom, stromek} \dicIndirectTranslationCS{(ve složeninách)} \dicExampleIS{ýviður} \dicExampleCS{tis}  \textbf{2.} \dicSynonym{timbur} \dicDirectTranslationCS{dřevo, dříví} \dicIndirectTranslationCS{(stavební materiál)};  \dicIdiom{viður}{ \dicPhraseIS{sólin gengur\,/\addthin hnígur til viðar}} \dicDirectTranslationCS{slunce klesá k~obzoru}
\dicEntry[viðureign] \dicTerm{viður··eign} \dicIPA{{v}{\textsci}{\textlengthmark}{ð}{\textscy}{r}{ei}{\r{g}}{\textsubring{n}}} \dicPos{f}[7] \dicFlx{(‑ar, ‑ir)}[1] \dicDirectTranslationCS{střetnutí, utkání} \dicExampleIS{viðureign liðanna} \dicExampleCS{utkání týmů};  \dicPhraseIS{e‑að er (erfitt) viðureignar} \dicDirectTranslationCS{s~(čím) se dá (obtížně) vypořádat}
\dicEntry[viðurkenna] \dicTerm{viður··kenn|a} \dicsymFrequent\  \dicIPA{{v}{\textsci}{\textlengthmark}{ð}{\textscy}{\textsubring{r}}{c\smash{\textsuperscript{h}}}{\textepsilon}{n}{a}} \dicPos{v}[2] \dicFlx{(‑di, ‑t)}[141] \dicFlx{acc} \textbf{1.} \dicSynonym{játa} \dicDirectTranslationCS{přiznat, přiznávat, doznat, doznávat, uznat, uznávat, připustit, připouštět} \dicExampleIS{viðurkenna þjófnaðinn} \dicExampleCS{přiznat krádež}  \textbf{2.} \dicSynonym*{taka gilt} \dicDirectTranslationCS{uznat, uznávat, akceptovat, schválit, schvalovat (skutečnost ap.)} \dicExampleIS{viðurkenna sjálfstæði} \dicExampleCS{uznat nezávislost}
\dicEntry[viðurkenndur] \dicTerm{viður··kenndur} \dicIPA{{v}{\textsci}{\textlengthmark}{ð}{\textscy}{\textsubring{r}}{c\smash{\textsuperscript{h}}}{\textepsilon}{n}{\textsubring{d}}{\textscy}{\textsubring{r}}} \dicPos{adj}[2]\dicFlx{}[-14] \dicDirectTranslationCS{uznaný, akceptovaný} \dicExampleIS{viðurkenndar kennisetningar} \dicExampleCS{uznané teorémy}
\dicEntry[viðurkenning] \dicTerm{viður··kenn·ing} \dicsymFrequent\  \dicIPA{{v}{\textsci}{\textlengthmark}{ð}{\textscy}{\textsubring{r}}{c\smash{\textsuperscript{h}}}{\textepsilon}{n}{i}{\ng}{\r{g}}} \dicPos{f}[4] \dicFlx{(‑ar, ‑ar)}[5] \textbf{1.} \dicSynonym{játning} \dicDirectTranslationCS{přiznání, doznání, připuštění}  \textbf{2.} \dicSynonym*{almennt samþykki} \dicDirectTranslationCS{uznání, přijetí, akceptace}  \textbf{3.} \dicSynonym{verðlaun} \dicDirectTranslationCS{uznání, ocenění} \dicExampleIS{viðurkenning á verkum hans} \dicExampleCS{uznání jeho práce}
\dicEntry[viðurkenningartenging] \dicTerm{viður·kenningar··teng·ing} \dicIPA{{v}\-{\textsci}\-{\textlengthmark}\-{ð}\-{\textscy}\-{\textsubring{r}}\-{c\smash{\textsuperscript{h}}}\-{\textepsilon}\-{n}\-{i}\-{\ng}\-{\r{g}}\-{a}\-{\textsubring{r}}\-{t\smash{\textsuperscript{h}}}\-{ei}\-{\textltailn}\-{\r{\textObardotlessj}}\-{i}\-{\ng}\-{\r{g}}\-} \dicPos{f}[4] \dicFlx{(‑ar, ‑ar)}[5] \dicFieldCat{jaz.} \dicDirectTranslationCS{spojka přípustková}
\dicEntry[viðurkvæmilegur] \dicTerm{viður··kvæmi·legur} \dicIPA{{v}{\textsci}{\textlengthmark}{ð}{\textscy}{\textsubring{r}}{k\smash{\textsuperscript{h}}}{v}{a}{i}{m}{\textsci}{l}{\textepsilon}{\textbabygamma}{\textscy}{\textsubring{r}}} \dicPos{adj}[1]\dicFlx{}[-8] \dicSynonym{sæmilegur} \dicDirectTranslationCS{vhodný, odpovídající, adekvátní}
\dicEntry[viðurlag] \dicTerm{viður··|lag} \dicIPA{{v}{\textsci}{\textlengthmark}{ð}{\textscy}{r}{l}{a}{x}} \dicPos{n}[2] \dicFlx{(‑lags, ‑lög)}[8] \textbf{1.} \dicFieldCat{jaz.} \dicDirectTranslationCS{přístavek, apozice}  \textbf{2.} \dicPhraseIS{viðurlög} \dicFlx{pl} \dicFieldCat{práv.} \dicDirectTranslationCS{sankce, postih}
\dicEntry[viðurnafn] \dicTerm{viður··|nafn} \dicIPA{{v}{\textsci}{\textlengthmark}{ð}{\textscy}{r}{n}{a}{\textsubring{b}}{\textsubring{n}}} \dicPos{n}[2] \dicFlx{(‑nafns, ‑nöfn)}[8] \dicLink{viðurnefni}
\dicEntry[viðurnefni] \dicTerm{viður··nefni}\dicTerm{, viðurnafn} \dicIPA{{v}\-{\textsci}\-{\textlengthmark}\-{ð}\-{\textscy}\-{r}\-{n}\-{\textepsilon}\-{\textsubring{b}}\-{n}\-{\textsci}\-} \dicPos{n}[2] \dicFlx{(‑s, ‑)}[14] \dicDirectTranslationCS{přezdívka, přízvisko}
\dicEntry[viðurstyggð] \dicTerm{viður··styggð} \dicIPA{{v}{\textsci}{\textlengthmark}{ð}{\textscy}{\textsubring{r}}{s}{\textsubring{d}}{\textsci}{\textbabygamma}{\texttheta}} \dicPos{f}[4] \dicFlx{(‑ar)}[3] \dicSynonym{andstyggð} \dicDirectTranslationCS{ohyzdnost, ohavnost, odpornost}
\dicEntry[viðurstyggilegur] \dicTerm{viður·styggi··legur} \dicIPA{{v}{\textsci}{\textlengthmark}{ð}{\textscy}{\textsubring{r}}{s}{\textsubring{d}}{\textsci}{\r{\textObardotlessj}}{\textsci}{l}{\textepsilon}{\textbabygamma}{\textscy}{\textsubring{r}}} \dicPos{adj}[1]\dicFlx{}[-8] \dicSynonym{andstyggilegur} \dicDirectTranslationCS{ohyzdný, ohavný, odporný} \dicExampleIS{Lyktin af hákarlinum er viðurstyggileg.} \dicExampleCS{Zápach žraloka je odporný.}
\dicEntry[viðurvist] \dicTerm{viður··vist} \dicIPA{{v}{\textsci}{\textlengthmark}{ð}{\textscy}{r}{v}{\textsci}{s}{\textsubring{d}}} \dicPos{f}[7] \dicFlx{(‑ar)}[3] \dicSynonym{nærvera} \dicDirectTranslationCS{přítomnost, prezence} \dicExampleIS{haga sér eins og fífl í viðurvist foreldra sinna} \dicExampleCS{chovat se jako hlupák v~přítomnosti vlastních rodičů}
\dicEntry[viðurværi] \dicTerm{viður··væri} \dicIPA{{v}{\textsci}{\textlengthmark}{ð}{\textscy}{r}{v}{a}{i}{r}{\textsci}} \dicPos{n}[2] \dicFlx{(‑s)}[20] \dicDirectTranslationCS{obživa, živobytí};  \dicPhraseIS{hafa viðurværi sitt af e‑u} \dicDirectTranslationCS{živit se (čím)} \dicExampleIS{hafa viðurværi sitt af lausavinnu} \dicExampleCS{živit se příležitostnou prací}
\dicEntry[viðutan] \dicTerm{við··utan} \dicIPA{{v}{\textsci}{\textlengthmark}{ð}{\textscy}{\textsubring{d}}{a}{\textsubring{n}}} \dicPos{adj}[13] \dicFlx{indecl}[2] \dicSynonym*{annars hugar} \dicDirectTranslationCS{(duchem) nepřítomný, zamyšlený}
\dicEntry[viðvaningslega] \dicTerm{við·vanings··lega} \dicIPA{{v}{\textsci}{ð}{v}{a}{n}{i}{\ng}{s}{l}{\textepsilon}{\textbabygamma}{a}} \dicPos{adv} \dicDirectTranslationCS{amatérsky, diletantsky, neprofesionálně}
\dicEntry[viðvaningslegur] \dicTerm{við·vanings··legur} \dicIPA{{v}{\textsci}{ð}{v}{a}{n}{i}{\ng}{s}{l}{\textepsilon}{\textbabygamma}{\textscy}{\textsubring{r}}} \dicPos{adj}[1]\dicFlx{}[-8] \dicDirectTranslationCS{amatérský, diletantský, neprofesionální}
\dicEntry[viðvaningur] \dicTerm{við··van·ing|ur} \dicIPA{{v}{\textsci}{ð}{v}{a}{n}{i}{\ng}{\r{g}}{\textscy}{\textsubring{r}}} \dicPos{m}[6] \dicFlx{(‑s, ‑ar)}[8] \dicSynonym{byrjandi} \dicDirectTranslationCS{začátečník, začátečnice, nováček}
\dicEntry[viðvarandi] \dicTerm{við··var·andi} \dicIPA{{v}{\textsci}{ð}{v}{a}{r}{a}{n}{\textsubring{d}}{\textsci}} \dicPos{adj}[13] \dicFlx{indecl}[1] \dicSynonym{varanlegur} \dicDirectTranslationCS{trvající, přetrvávající} \dicExampleIS{viðvarandi atvinnuleysi} \dicExampleCS{trvající nezaměstnanost}
\dicEntry[viðvart] \dicTerm{við··vart} \dicIPA{{v}{\textsci}{ð}{v}{a}{\textsubring{r}}{\textsubring{d}}} \dicPos{adj}[13]\dicFlx{}[-4] \dicPhraseIS{gera e‑m viðvart um e‑ð} \dicDirectTranslationCS{upozornit (koho) na (co), varovat (koho) před (čím)}
\dicEntry[viðveik] \dicTerm{við··veik} \dicIPA{{v}{\textsci}{ð}{v}{ei}{\r{g}}} \dicPos{v} \dicFlx{ind pf sg 1 pers} \dicLink{viðvíkja}
\dicEntry[viðvék] \dicTerm{við··vék} \dicIPA{{v}{\textsci}{ð}{v}{j}{\textepsilon}{\r{g}}} \dicPos{v} \dicFlx{ind pf sg 1 pers} \dicLink{viðvíkja}
\dicEntry[viðvik] \dicTerm{við··vik} \dicIPA{{v}{\textsci}{ð}{v}{\textsci}{\r{g}}} \dicPos{n}[2] \dicFlx{(‑s, ‑)}[5] \dicSynonym*{smávinna} \dicDirectTranslationCS{drobná práce\,/\addthin služba}
\dicEntry[viðviki] \dicTerm{við··viki} \dicIPA{{v}{\textsci}{ð}{v}{\textsci}{\r{\textObardotlessj}}{\textsci}} \dicPos{v} \dicFlx{con pf sg 1 pers} \dicLink{viðvíkja}
\dicEntry[viðvikið] \dicTerm{við··vikið} \dicIPA{{v}{\textsci}{ð}{v}{\textsci}{\r{\textObardotlessj}}{\textsci}{\texttheta}} \dicPos{v} \dicFlx{supin} \dicLink{viðvíkja}
\dicEntry[viðvíkja] \dicTerm{við··|víkja} \dicIPA{{v}{\textsci}{ð}{v}{i}{\r{\textObardotlessj}}{a}} \dicPos{v}[6] \dicFlx{(‑víkur, ‑vék\,/\addthin ‑veik, ‑viki, ‑vikið)}[62] \dicFlx{impers} \dicPhraseIS{hvað e‑u viðvíkur} \dicDirectTranslationCS{co se (čeho) týká, pokud jde o~(co)}
\dicEntry[viðvíkjandi] \dicTerm{við·víkj··andi} \dicIPA{{v}{\textsci}{ð}{v}{i}{\r{\textObardotlessj}}{a}{n}{\textsubring{d}}{\textsci}} \dicPos{adj}[13] \dicFlx{indecl}[1] \dicPhraseIS{viðvíkjandi e‑u} \dicDirectTranslationCS{týkající se (čeho)}
\dicEntry[viðvíkur] \dicTerm{við··víkur} \dicIPA{{v}{\textsci}{ð}{v}{i}{\r{g}}{\textscy}{\textsubring{r}}} \dicPos{v} \dicFlx{ind praes sg 1 pers} \dicLink{viðvíkja}
\dicEntry[viðvörun] \dicTerm{við··|vörun} \dicIPA{{v}{\textsci}{ð}{v}{\oe}{r}{\textscy}{\textsubring{n}}} \dicPos{f}[7] \dicFlx{(‑vörunar, ‑varanir)}[11] \dicSynonym{aðvörun} \dicDirectTranslationCS{varování, výstraha, upozornění} \dicExampleIS{viðvörun um flóð} \dicExampleCS{výstraha o~nebezpečí povodní}
\dicEntry[viðvörunarljós] \dicTerm{við·vörunar··ljós} \dicIPA{{v}{\textsci}{ð}{v}{\oe}{r}{\textscy}{n}{a}{r}{l}{j}{ou}{s}} \dicPos{n}[2] \dicFlx{(‑s, ‑)}[5] \dicDirectTranslationCS{výstražné světlo}
\dicEntry[viðvörunarmerki] \dicTerm{við·vörunar··merki} \dicIPA{{v}{\textsci}{ð}{v}{\oe}{r}{\textscy}{n}{a}{r}{m}{\textepsilon}{\textsubring{r}}{\r{\textObardotlessj}}{\textsci}} \dicPos{n}[2] \dicFlx{(‑s, ‑)}[16] \dicDirectTranslationCS{výstraž\-né znamení, výstražná značka}
\dicEntry[viðþolslaus] \dicTerm{við·þols··laus} \dicIPA{{v}{\textsci}{\texttheta}{\textlengthmark}{\textopeno}{l}{s}{l}{\oe i}{s}} \dicPos{adj}[5]\dicFlx{}[-1] \dicPhraseIS{vera viðþolslaus af e‑u} \dicDirectTranslationCS{být nepříčetný z~(čeho), nebýt už schopný (co) vydržet (bolest ap.)}
\dicEntry[vifta] \dicTerm{vift|a} \dicIPA{{v}{\textsci}{f}{\textsubring{d}}{a}} \dicPos{f}[1] \dicFlx{(‑u, ‑ur)}[7] \textbf{1.} \dicSynonym*{vifta til loftkælingar} \dicDirectTranslationCS{větrák, ventilátor}  \textbf{2.} \dicSynonym*{eldhúsvifta} \dicDirectTranslationCS{odsavač (par), digestoř}
\dicEntry[vigt] \dicTerm{vigt} \dicIPA{{v}{\textsci}{x}{\textsubring{d}}} \dicPos{f}[7] \dicFlx{(‑ar, ‑ir)}[1] \textbf{1.} \dicSynonym{þyngd} \dicDirectTranslationCS{váha} \dicIndirectTranslationCS{(hmotnost)};  \dicPhraseIS{selja e‑ð eftir vigt} \dicDirectTranslationCS{prodávat (co) na váhu}  \textbf{2.} \dicSynonym{vog} \dicDirectTranslationCS{váha} \dicIndirectTranslationCS{(přístroj na vážení)};  \dicIdiom{vigt}{ \dicPhraseIS{e‑að hefur mikla vigt}} \dicLangCat{přen.} \dicDirectTranslationCS{(co) má velkou váhu}
\dicEntry[vigta] \dicTerm{vigt|a} \dicIPA{{v}{\textsci}{x}{\textsubring{d}}{a}} \dicPos{v}[1] \dicFlx{(‑aði)}[1] \dicFlx{acc} \dicSynonym{vega\smash{\textsuperscript{1}}} \dicDirectTranslationCS{(z)vážit} \dicIndirectTranslationCS{(určit váhu něčeho)} \dicExampleIS{vigta vöruna} \dicExampleCS{zvážit výrobek}
\dicEntry[vik] \dicTerm{vik} \dicsymFrequent\  \dicIPA{{v}{\textsci}{\textlengthmark}{\r{g}}} \dicPos{n}[2] \dicFlx{(‑s, ‑)}[5] \textbf{1.} \dicLangCat{zast.} \dicSynonym{viðvik} \dicDirectTranslationCS{výpomoc, drobná práce}  \textbf{2.} \dicSynonym{skarð} \dicDirectTranslationCS{zářez, vrub};  \dicIdiom{vik}{ \dicPhraseIS{eiga ekki hægt um vik}} \dicLangCat{přen.} \dicDirectTranslationCS{nemít to lehké, mít obtíže};  \dicPhraseIS{fyrir vikið} \dicFlx{adv} \dicDirectTranslationCS{kvůli tomu}
\dicEntry[vika] \dicTerm{vik|a} \dicsymFrequent\  \dicIPA{{v}{\textsci}{\textlengthmark}{\r{g}}{a}} \dicPos{f}[1] \dicFlx{(‑u, ‑ur)}[13] \dicDirectTranslationCS{týden} \dicExampleIS{fara í sund þrisvar í viku} \dicExampleCS{chodit plavat třikrát týdně}
\dicEntry[vikadrengur] \dicTerm{vika··dreng|ur} \dicIPA{{v}{\textsci}{\textlengthmark}{\r{g}}{a}{\textsubring{d}}{r}{ei}{\ng}{\r{g}}{\textscy}{\textsubring{r}}} \dicPos{m}[9] \dicFlx{(‑s, ‑ir)}[15] \dicDirectTranslationCS{poslíček, pikolo}
\dicEntry[viki] \dicTerm{viki} \dicIPA{{v}{\textsci}{\textlengthmark}{\r{\textObardotlessj}}{\textsci}} \dicPos{v} \dicFlx{con pf sg 1 pers} \dicLink{víkja}
\dicEntry[vikið] \dicTerm{vikið} \dicIPA{{v}{\textsci}{\textlengthmark}{\r{\textObardotlessj}}{\textsci}{\texttheta}} \dicPos{v} \dicFlx{supin} \dicLink{víkja}
\dicEntry[vikivaki] \dicTerm{viki··vak|i} \dicIPA{{v}{\textsci}{\textlengthmark}{\r{\textObardotlessj}}{\textsci}{v}{a}{\r{\textObardotlessj}}{\textsci}} \dicPos{m}[1] \dicFlx{(‑a, ‑ar)}[8] \dicIndirectTranslationCS{starý islandský národní tanec}
\dicEntry[vikna] \dicTerm{vikn|a} \dicIPA{{v}{\textsci}{h}{\r{g}}{n}{a}} \dicPos{v}[1] \dicFlx{(‑aði)}[1] \dicDirectTranslationCS{být dojat, (z)jihnout, dojmout se} \dicExampleIS{vikna af geðshræringu} \dicExampleCS{zjihnout dojetím}
\dicEntry[Viktoríuvatn] \dicTerm{Viktoríu··vatn} \dicIPA{{v}{\textsci}{x}{\textsubring{d}}{\textopeno}{r}{i}{j}{\textscy}{v}{a}{h}{\textsubring{d}}{\textsubring{n}}} \dicPos{n}[2] \dicFlx{(‑s)}[4] \dicFieldCat{geog.} \dicDirectTranslationCS{Viktoriino jezero} \dicIndirectTranslationCS{(jezero ve východní Africe)}
\dicEntry[vikublað] \dicTerm{viku··|blað} \dicIPA{{v}{\textsci}{\textlengthmark}{\r{g}}{\textscy}{\textsubring{b}}{l}{a}{\texttheta}} \dicPos{n}[2] \dicFlx{(‑blaðs, ‑blöð)}[8] \dicDirectTranslationCS{týdeník}
\dicEntry[vikudagur] \dicTerm{viku··dag|ur} \dicIPA{{v}{\textsci}{\textlengthmark}{\r{g}}{\textscy}{\textsubring{d}}{a}{\textbabygamma}{\textscy}{\textsubring{r}}} \dicPos{m}[6] \dicFlx{(‑s, ‑ar)}[62] \dicDirectTranslationCS{den} \dicIndirectTranslationCS{(v~týdnu)}
\dicEntry[vikukaup] \dicTerm{viku··kaup} \dicIPA{{v}{\textsci}{\textlengthmark}{\r{g}}{\textscy}{k\smash{\textsuperscript{h}}}{\oe i}{\textsubring{b}}} \dicPos{n}[2] \dicFlx{(‑s)}[2] \dicDirectTranslationCS{týdenní plat\,/\addthin mzda}
\dicEntry[vikulaun] \dicTerm{viku··laun} \dicIPA{{v}{\textsci}{\textlengthmark}{\r{g}}{\textscy}{l}{\oe i}{\textsubring{n}}} \dicPos{n}[2] \dicFlx{pl}[1] \dicDirectTranslationCS{týdenní plat\,/\addthin mzda}
\dicEntry[vikulega] \dicTerm{viku··lega} \dicIPA{{v}{\textsci}{\textlengthmark}{\r{g}}{\textscy}{l}{\textepsilon}{\textbabygamma}{a}} \dicPos{adv} \dicDirectTranslationCS{týdně, každý týden, jednou za týden}
\dicEntry[vikulegur] \dicTerm{viku··legur} \dicIPA{{v}{\textsci}{\textlengthmark}{\r{g}}{\textscy}{l}{\textepsilon}{\textbabygamma}{\textscy}{\textsubring{r}}} \dicPos{adj}[1]\dicFlx{}[-6] \dicDirectTranslationCS{týdenní}
\dicEntry[vikum] \dicTerm{vikum} \dicIPA{{v}{\textsci}{\textlengthmark}{\r{g}}{\textscy}{\textsubring{m}}} \dicPos{v} \dicFlx{ind pf pl 1 pers} \dicLink{víkja}
\dicEntry[vikur] \dicTerm{vik|ur} \dicIPA{{v}{\textsci}{\textlengthmark}{\r{g}}{\textscy}{\textsubring{r}}} \dicPos{m}[5] \dicFlx{(‑urs)}[2] \dicFieldCat{geol.} \dicDirectTranslationCS{pemza}
\dicEntry[vil] \dicTerm{vil\smash{\textsuperscript{1}}} \dicIPA{{v}{\textsci}{\textlengthmark}{\textsubring{l}}} \dicPos{n}[2] \dicFlx{(‑s)}[2] \dicDirectTranslationCS{prospěch, dobro};  \dicPhraseIS{í vil} \dicFlx{adv} \dicDirectTranslationCS{k~dobru, ku prospěchu} \dicExampleIS{Hvernig er staðan í leiknum? -- 2:0 Íslendingum í vil.} \dicExampleCS{Jaký je stav hry? -- 2:0 pro Islanďany.};  \dicPhraseIS{e‑m gengur allt í vil} \dicFlx{impers} \dicDirectTranslationCS{(kdo) má úspěch, (komu) se vše daří}
\dicEntry[vil] \dicTerm{vil\smash{\textsuperscript{2}}} \dicIPA{{v}{\textsci}{\textlengthmark}{\textsubring{l}}} \dicPos{v} \dicFlx{ind praes sg 1 pers} \dicLink{vilja}
\dicEntry[vild] \dicTerm{vild} \dicIPA{{v}{\textsci}{l}{\textsubring{d}}} \dicPos{f}[4] \dicFlx{(‑ar)}[3] \dicPhraseIS{að vild} \dicFlx{adv} \dicDirectTranslationCS{podle libosti, jak je libo}
\dicEntry[vildarkjör] \dicTerm{vildar··kjör} \dicIPA{{v}{\textsci}{l}{\textsubring{d}}{a}{\textsubring{r}}{c\smash{\textsuperscript{h}}}{\oe}{\textsubring{r}}} \dicPos{n}[2] \dicFlx{pl}[9] \dicDirectTranslationCS{příznivé\,/\addthin výhodné podmínky}
\dicEntry[vildarvinur] \dicTerm{vildar··vin|ur} \dicIPA{{v}{\textsci}{l}{\textsubring{d}}{a}{r}{v}{\textsci}{n}{\textscy}{\textsubring{r}}} \dicPos{m}[10] \dicFlx{(‑ar, ‑ir)}[17] \dicDirectTranslationCS{důvěrný kamarád, důvěrná kamarádka, blízký přítel, blízká přítelkyně}
\dicEntry[vildi] \dicTerm{vildi} \dicIPA{{v}{\textsci}{l}{\textsubring{d}}{\textsci}} \dicPos{v} \dicFlx{ind\,/\addthin con pf sg 1 pers} \dicLink{vilja}
\dicEntry[vildum] \dicTerm{vildum} \dicIPA{{v}{\textsci}{l}{\textsubring{d}}{\textscy}{\textsubring{m}}} \dicPos{v} \dicFlx{ind pf pl 1 pers} \dicLink{vilja}
\begin{xtolerant}{}{1pt}
\dicEntry[vilhallur] \dicTerm{vil··|hallur} \dicIPA{{v}{\textsci}{\textlengthmark}{\textsubring{l}}{h}{a}{\textsubring{d}}{l}{\textscy}{\textsubring{r}}} \dicPos{adj}[1] \dicFlx{(f ‑höll)}[2] \dicSynonym{hlutdrægur} \dicDirectTranslationCS{zaujatý, předpojatý, neobjektivní} \dicExampleIS{vilhallur e‑m} \dicExampleCS{předpojatý ke (komu)}
\end{xtolerant}
\dicEntry[vilja] \dicTerm{vilja} \dicsymFrequent\  \dicIPA{{v}{\textsci}{l}{j}{a}} \dicPos{v}[4] \dicFlx{(vil, vildi, vildum, vildi, viljað)}[39] \dicFlx{acc\,/\addthin dat} \textbf{1.} \dicFlx{acc} \dicSynonym*{langa í} \dicDirectTranslationCS{chtít, přát si, prahnout};  \dicPhraseIS{vilja e‑ð} \dicDirectTranslationCS{chtít (co)} \dicExampleIS{vilja ekki gjöfina} \dicExampleCS{nechtít dárek};  \dicPhraseIS{vilja e‑ð frekar en e‑ð annað} \dicDirectTranslationCS{dávat (čemu) přednost před (čím), upřednostňovat (co) před (čím)}  \textbf{2.} \dicSynonym{óska} \dicDirectTranslationCS{přát si, mít přání} \dicExampleIS{vilja þetta endilega} \dicExampleCS{moc si to přát}  \textbf{3.} \dicDirectTranslationCS{moct, moci} \dicIndirectTranslationCS{(vyjadřuje požadavek nebo příkaz)} \dicExampleIS{Viltu koma hingað?} \dicExampleCS{Můžeš sem přijít?}  \textbf{4.} \dicFlx{dat} \dicSynonym*{eiga erindi} \dicDirectTranslationCS{chtít} \dicIndirectTranslationCS{(mít něco na vyřízení)} \dicExampleIS{Hvað viltu mér?} \dicExampleCS{Co mi chceš?}  \textbf{5.} \dicDirectTranslationCS{mít tendenci, chtít} \dicIndirectTranslationCS{(užívané jako pomocné sloveso)} \dicExampleIS{það vill fara svona} \dicExampleCS{má to tendenci takhle se chovat};  \dicIdiom{vilja}[til]{ \dicPhraseIS{það vill þannig til að (e‑að gerist)}} \dicFlx{impers} \dicDirectTranslationCS{stává se, že se ((co) stane); přihodí se, že se ((co) stane)};  \dicPhraseIS{það vildi e‑m til happs að (e‑að gerðist)} \dicFlx{impers} \dicDirectTranslationCS{(kdo) měl štěstí, že se ((co) stalo)};  \dicPhraseIS{svo vel vildi til} \dicFlx{adv} \dicDirectTranslationCS{naštěstí, náhoda tomu chtěla};  \dicIdiom{vilja}{ \dicPhraseIS{ef til vill}} \dicFlx{adv} \dicSynonym{kannski} \dicDirectTranslationCS{snad, možná, asi}; { \dicPhraseIS{vilja e‑m vel}} \dicDirectTranslationCS{přát (komu) jen to dobré}; { \dicPhraseIS{vilja halda e‑u fram}} \dicDirectTranslationCS{prohlásit (co), tvrdit (co)}
\dicEntry[viljað] \dicTerm{viljað} \dicIPA{{v}{\textsci}{l}{j}{a}{\texttheta}} \dicPos{v} \dicFlx{supin} \dicLink{vilja}
\dicEntry[viljafastur] \dicTerm{vilja··|fastur} \dicIPA{{v}{\textsci}{l}{j}{a}{f}{a}{s}{\textsubring{d}}{\textscy}{\textsubring{r}}} \dicPos{adj}[1] \dicFlx{(f ‑föst)}[11] \dicDirectTranslationCS{mající pevnou vůli, rozhodný}
\dicEntry[viljafesta] \dicTerm{vilja··fest|a} \dicIPA{{v}{\textsci}{l}{j}{a}{f}{\textepsilon}{s}{\textsubring{d}}{a}} \dicPos{f}[1] \dicFlx{(‑u)}[5] \dicDirectTranslationCS{pevná vůle, rozhodnost}
\dicEntry[viljalaus] \dicTerm{vilja··laus} \dicIPA{{v}{\textsci}{l}{j}{a}{l}{\oe i}{s}} \dicPos{adj}[5]\dicFlx{}[-1] \dicDirectTranslationCS{mající slabou vůli}
\dicEntry[viljaleysi] \dicTerm{vilja··leysi} \dicIPA{{v}{\textsci}{l}{j}{a}{l}{ei}{s}{\textsci}} \dicPos{n}[2] \dicFlx{(‑s)}[20] \dicDirectTranslationCS{nedostatek vůle}
\dicEntry[viljandi] \dicTerm{vilj··andi\smash{\textsuperscript{1}}} \dicIPA{{v}{\textsci}{l}{j}{a}{n}{\textsubring{d}}{\textsci}} \dicPos{adj}[13] \dicFlx{indecl}[1] \dicDirectTranslationCS{úmyslný, záměrný}
\dicEntry[viljandi] \dicTerm{vilj··andi\smash{\textsuperscript{2}}} \dicIPA{{v}{\textsci}{l}{j}{a}{n}{\textsubring{d}}{\textsci}} \dicPos{adv} \dicDirectTranslationCS{úmyslně, záměrně, schválně} \dicExampleIS{Hann braut glasið viljandi.} \dicExampleCS{Rozbil tu skleničku úmyslně.}
\dicEntry[viljasterkur] \dicTerm{vilja··sterkur} \dicIPA{{v}{\textsci}{l}{j}{a}{s}{\textsubring{d}}{\textepsilon}{\textsubring{r}}{\r{g}}{\textscy}{\textsubring{r}}} \dicPos{adj}[1]\dicFlx{}[-1] \dicSynonym{viljafastur} \dicDirectTranslationCS{mající silnou vůli, odhodlaný}
\dicEntry[viljastyrkur] \dicTerm{vilja··styrk|ur} \dicIPA{{v}{\textsci}{l}{j}{a}{s}{\textsubring{d}}{\textsci}{\textsubring{r}}{\r{g}}{\textscy}{\textsubring{r}}} \dicPos{m}[9] \dicFlx{(‑s, ‑ir)}[15] \dicSynonym{viljaþrek} \dicDirectTranslationCS{síla vůle}
\dicEntry[viljaþrek] \dicTerm{vilja··þrek} \dicIPA{{v}{\textsci}{l}{j}{a}{\texttheta}{r}{\textepsilon}{\r{g}}} \dicPos{n}[2] \dicFlx{(‑s)}[2] \dicDirectTranslationCS{síla vůle}
\dicEntry[vilji] \dicTerm{vilj|i} \dicsymFrequent\  \dicIPA{{v}{\textsci}{l}{j}{\textsci}} \dicPos{m}[1] \dicFlx{(‑a)}[3] \textbf{1.} \dicDirectTranslationCS{vůle} \dicIndirectTranslationCS{(vytrvalá snaha o~dosažení cíle)} \dicExampleIS{hafa veikan vilja} \dicExampleCS{mít slabou vůli}  \textbf{2.} \dicDirectTranslationCS{vůle, přání};  \dicPhraseIS{af fúsum vilja} \dicFlx{adv} \dicDirectTranslationCS{ochotně, z~vlastní vůle};  \dicPhraseIS{fara að vilja e‑rs} \dicDirectTranslationCS{jednat podle (čí) vůle};  \dicPhraseIS{fá vilja sínum framgengt} \dicDirectTranslationCS{prosadit svou vůli};  \dicPhraseIS{(gera e‑ð) gegn vilja e‑rs} \dicDirectTranslationCS{(udělat (co)) proti (čí) vůli};  \dicPhraseIS{koma fram vilja sínum} \dicDirectTranslationCS{prosadit svou vůli};  \dicPhraseIS{sýna vilja til e‑s} \dicDirectTranslationCS{ukázat ochotu k~(čemu)}
\dicEntry[viljugur] \dicTerm{vilj··ugur} \dicIPA{{v}{\textsci}{l}{j}{\textscy}{\textbabygamma}{\textscy}{\textsubring{r}}} \dicPos{adj}[1]\dicFlx{}[-8] \textbf{1.} \dicSynonym{fús} \dicDirectTranslationCS{ochotný, ochoten} \dicExampleIS{Hún er viljug að hjálpa mömmu sinni.} \dicExampleCS{Je ochotná pomáhat své mamince.}  \textbf{2.} \dicSynonym{fjörugur} \dicDirectTranslationCS{čilý, živý, hbitý, agilní} \dicExampleIS{viljugur hestur} \dicExampleCS{živý kůň};  \dicPhraseIS{(gera e‑ð) nauðugur viljugur} \dicDirectTranslationCS{(udělat (co)) chtě nechtě}
\dicEntry[villa] \dicTerm{vill|a\smash{\textsuperscript{1}}} \dicIPA{{v}{\textsci}{\textsubring{d}}{l}{a}} \dicPos{f}[1] \dicFlx{(‑u, ‑ur)}[13] \textbf{1.} \dicSynonym{mistök} \dicDirectTranslationCS{chyba} \dicExampleIS{villa í reikningi} \dicExampleCS{chyba ve výpočtu}  \textbf{2.} \dicSynonym*{röng hugmynd} \dicDirectTranslationCS{omyl, mýlka}  \textbf{3.} \dicSynonym*{það að villast} \dicDirectTranslationCS{(za)bloudění}  \textbf{4.} \dicSynonym{villutrú} \dicDirectTranslationCS{zbloudění, hereze};  \dicIdiom{villa}{ \dicPhraseIS{snúa frá villu síns vegar}} \dicLangCat{přen.} \dicDirectTranslationCS{obrátit se na správnou cestu}; { \dicPhraseIS{vaða í villu og svíma}} \dicLangCat{přen.} \dicDirectTranslationCS{být na omylu}
\dicEntry[villa] \dicTerm{vill|a\smash{\textsuperscript{2}}} \dicIPA{{v}{\textsci}{l}{\textlengthmark}{a}} \dicPos{f}[1] \dicFlx{(‑u, ‑ur)}[13] \dicSynonym{einbýlishús} \dicDirectTranslationCS{vila}
\dicEntry[villa] \dicTerm{vill|a\smash{\textsuperscript{3}}} \dicsymFrequent\  \dicIPA{{v}{\textsci}{\textsubring{d}}{l}{a}} \dicPos{v}[2] \dicFlx{(‑ti, ‑t)}[63] \dicPhraseIS{villa e‑m sýn} \dicDirectTranslationCS{(z)mýlit (koho), oklamat (koho)};  \dicIdiom{villa}[á]{ \dicPhraseIS{villa á sér heimildir}} \dicDirectTranslationCS{vydávat se za někoho jiného, vystupovat pod falešnou identitou};  \dicIdiom{villa}[um]{ \dicPhraseIS{villa um fyrir e‑m}} \dicDirectTranslationCS{oklamat (koho), uvést (koho) v~omyl};  \dicIdiom{villast}{ \dicPhraseIS{villast}} \dicFlx{refl} \dicDirectTranslationCS{ztratit se, ztrácet se, (za)bloudit, pomýlit si cestu} \dicExampleIS{villast í þokunni} \dicExampleCS{ztratit se v~mlze}; { \dicPhraseIS{villast af leið}} \dicFlx{refl} \dicDirectTranslationCS{zmýlit si cestu, ztratit orientaci}; { \dicPhraseIS{það er ekki um að villast}} \dicFlx{refl} \dicDirectTranslationCS{není o~tom pochyb};  \dicIdiom{villast}[á]{ \dicPhraseIS{villast á e‑u\,/\addthin e‑m}} \dicFlx{refl} \dicDirectTranslationCS{splést si (co\,/\addthin koho), poplést (co\,/\addthin koho)} \dicExampleIS{villast á tvíburunum} \dicExampleCS{splést si dvojčata}
\dicEntry[villi] \dicTerm{villi-} \dicIPA{{v}{\textsci}{\textsubring{d}}{l}{\textsci}} \dicPos{predp} \dicDirectTranslationCS{divoký}
\dicEntry[villiblóm] \dicTerm{villi··blóm} \dicIPA{{v}{\textsci}{\textsubring{d}}{l}{\textsci}{\textsubring{b}}{l}{ou}{\textsubring{m}}} \dicPos{n}[2] \dicFlx{(‑s, ‑)}[5] \dicDirectTranslationCS{divoká\,/\addthin polní květina}
\dicEntry[villidýr] \dicTerm{villi··dýr} \dicIPA{{v}{\textsci}{\textsubring{d}}{l}{\textsci}{\textsubring{d}}{i}{\textsubring{r}}} \dicPos{n}[2] \dicFlx{(‑s, ‑)}[5] \textbf{1.} \dicSynonym*{villt dýr} \dicDirectTranslationCS{divoké zvíře, divoká zvěř}  \textbf{2.} \dicSynonym{óhemja} \dicDirectTranslationCS{divoch, divoška, divočák} \dicIndirectTranslationCS{(o~člověku)}
\dicEntry[villigöltur] \dicTerm{villi··|göltur} \dicIPA{{v}{\textsci}{\textsubring{d}}{l}{\textsci}{\r{g}}{\oe}{\textsubring{l}}{\textsubring{d}}{\textscy}{\textsubring{r}}} \dicPos{m}[11] \dicFlx{(‑galtar, ‑geltir)}[5] \dicDirectTranslationCS{divočák, kanec}
\dicEntry[villigötur] \dicTerm{villi··götur} \dicIPA{{v}{\textsci}{\textsubring{d}}{l}{\textsci}{\r{g}}{\oe}{\textsubring{d}}{\textscy}{\textsubring{r}}} \dicPos{f}[1] \dicFlx{pl}[15] \dicDirectTranslationCS{scestí};  \dicPhraseIS{leiða e‑n á villigötur} \dicDirectTranslationCS{vést (koho) na scestí};  \dicPhraseIS{vera á villigötum} \dicDirectTranslationCS{být na scestí, octnout se na scestí}
\dicEntry[villiköttur] \dicTerm{villi··|köttur} \dicIPA{{v}{\textsci}{\textsubring{d}}{l}{\textsci}{k\smash{\textsuperscript{h}}}{\oe}{h}{\textsubring{d}}{\textscy}{\textsubring{r}}} \dicPos{m}[11] \dicFlx{(‑kattar, ‑kettir)}[5] \dicDirectTranslationCS{divoká kočka}
\dicEntry[villimaður] \dicTerm{villi··|maður} \dicIPA{{v}{\textsci}{\textsubring{d}}{l}{\textsci}{m}{a}{ð}{\textscy}{\textsubring{r}}} \dicPos{m}[13] \dicFlx{(‑manns, ‑menn)}[2] \textbf{1.} \dicSynonym*{villtur maður} \dicDirectTranslationCS{divoch, divoška, barbar(ka)} \dicIndirectTranslationCS{(necivilizovaný člověk)}  \textbf{2.} \dicSynonym*{illa siðaður maður} \dicDirectTranslationCS{divoch, divoška} \dicIndirectTranslationCS{(neukázněný člověk)}
\dicEntry[villimannlegur] \dicTerm{villi·mann··legur}\dicTerm{, villimannslegur} \dicIPA{{v}\-{\textsci}\-{\textsubring{d}}\-{l}\-{\textsci}\-{m}\-{a}\-{n}\-{l}\-{\textepsilon}\-{\textbabygamma}\-{\textscy}\-{\textsubring{r}}\-} \dicPos{adj}[1]\dicFlx{}[-8] \dicDirectTranslationCS{brutální, bestiální}
\dicEntry[villimannslegur] \dicTerm{villi·manns··legur} \dicIPA{{v}{\textsci}{\textsubring{d}}{l}{\textsci}{m}{a}{n}{s}{l}{\textepsilon}{\textbabygamma}{\textscy}{\textsubring{r}}} \dicPos{adj}[1]\dicFlx{}[-8] \dicLink{villimannlegur}
\dicEntry[villimennska] \dicTerm{villi··mennsk|a} \dicIPA{{v}{\textsci}{\textsubring{d}}{l}{\textsci}{m}{\textepsilon}{n}{s}{\r{g}}{a}} \dicPos{f}[1] \dicFlx{(‑u)}[5] \textbf{1.} \dicDirectTranslationCS{divošství, barbarství} \dicIndirectTranslationCS{(kulturní zaostalost)}  \textbf{2.} \dicSynonym*{ósæmileg hegðun} \dicDirectTranslationCS{divošství, barbarství} \dicIndirectTranslationCS{(hrubost)}
\dicEntry[villingur] \dicTerm{vill··ing|ur} \dicIPA{{v}{\textsci}{\textsubring{d}}{l}{i}{\ng}{\r{g}}{\textscy}{\textsubring{r}}} \dicPos{m}[6] \dicFlx{(‑s, ‑ar)}[8] \dicDirectTranslationCS{divoch, divoška, výtržník, výtržnice}
\dicEntry[villisvín] \dicTerm{villi··svín} \dicIPA{{v}{\textsci}{\textsubring{d}}{l}{\textsci}{s}{v}{i}{\textsubring{n}}} \dicPos{n}[2] \dicFlx{(‑s, ‑)}[5] \dicFieldCat{zool.} \dicDirectTranslationCS{prase divoké} \textit{(l.~{\textLA{Sus scrofa}})}  \dicsymPhoto\ 
\dicFigure{109669.jpg}{Villisvín}{Villisvín - Velechovský Marek, Biolib, Copyright/CC-BY-NC}
\dicEntry[villtur] \dicTerm{villtur} \dicsymFrequent\  \dicIPA{{v}{\textsci}{\textsubring{l}}{\textsubring{d}}{\textscy}{\textsubring{r}}} \dicPos{adj}[1]\dicFlx{}[-10] \textbf{1.} \dicSynonym{ótaminn} \dicDirectTranslationCS{divoký, neochočený} \dicExampleIS{villt spendýr á Íslandi} \dicExampleCS{divocí savci na Islandu}  \textbf{2.} \dicSynonym{áttavilltur} \dicDirectTranslationCS{ztracený, zbloudilý}  \textbf{3.} \dicSynonym{hömlulaus} \dicDirectTranslationCS{divoký, bujný, neukázněný}
\dicEntry[villugjarn] \dicTerm{villu··|gjarn} \dicIPA{{v}{\textsci}{\textsubring{d}}{l}{\textscy}{\r{\textObardotlessj}}{a}{r}{\textsubring{d}}{\textsubring{n}}} \dicPos{adj}[5] \dicFlx{(f ‑gjörn)}[6] \dicDirectTranslationCS{zrádný na orientaci, snadný k~zabloudění} \dicExampleIS{villugjarn vegur} \dicExampleCS{cesta zrádná na orientaci}
\dicEntry[villukenning] \dicTerm{villu··kenn·ing} \dicIPA{{v}{\textsci}{\textsubring{d}}{l}{\textscy}{c\smash{\textsuperscript{h}}}{\textepsilon}{n}{i}{\ng}{\r{g}}} \dicPos{f}[4] \dicFlx{(‑ar, ‑ar)}[5] \dicFieldCat{náb.} \dicDirectTranslationCS{kacířské učení}
\dicEntry[villulaus] \dicTerm{villu··laus} \dicIPA{{v}{\textsci}{\textsubring{d}}{l}{\textscy}{l}{\oe i}{s}} \dicPos{adj}[5]\dicFlx{}[-1] \dicDirectTranslationCS{bezchybný, (jsoucí) bez chyby} \dicExampleIS{villulaus stíll} \dicExampleCS{bezchybný styl}
\dicEntry[villulaust] \dicTerm{villu··laust} \dicIPA{{v}{\textsci}{\textsubring{d}}{l}{\textscy}{l}{\oe i}{s}{\textsubring{d}}} \dicPos{adv} \dicDirectTranslationCS{bezchybně}
\dicEntry[villuleit] \dicTerm{villu··leit} \dicIPA{{v}{\textsci}{\textsubring{d}}{l}{\textscy}{l}{ei}{\textsubring{d}}} \dicPos{f}[7] \dicFlx{(‑ar)}[3] \dicDirectTranslationCS{vyhledání chyb, kontrola pravopisu}
\dicEntry[villuleitarforrit] \dicTerm{villu·leitar··for·rit} \dicIPA{{v}\-{\textsci}\-{\textsubring{d}}\-{l}\-{\textscy}\-{l}\-{ei}\-{\textsubring{d}}\-{a}\-{\textsubring{r}}\-{f}\-{\textopeno}\-{r}\-{\textlengthmark}\-{\textsci}\-{\textsubring{d}}\-} \dicPos{n}[2] \dicFlx{(‑s, ‑)}[5] \dicFieldCat{poč.} \dicDirectTranslationCS{kontrola pravopisu, (automatický) korektor pravopisu} \dicIndirectTranslationCS{(počítačový program)}
\dicEntry[villuljós] \dicTerm{villu··ljós} \dicIPA{{v}{\textsci}{\textsubring{d}}{l}{\textscy}{l}{j}{ou}{s}} \dicPos{n}[2] \dicFlx{(‑s, ‑)}[5] \dicDirectTranslationCS{bludička}
\dicEntry[villur] \dicTerm{villur} \dicIPA{{v}{\textsci}{\textsubring{d}}{l}{\textscy}{\textsubring{r}}} \dicPos{adj}[1]\dicFlx{}[-1] \dicDirectTranslationCS{ztracený, zbloudilý};  \dicPhraseIS{fara villur vegar} \dicDirectTranslationCS{zabloudit, ztratit se}
\dicEntry[villuráfandi] \dicTerm{villu··ráf·andi} \dicIPA{{v}{\textsci}{\textsubring{d}}{l}{\textscy}{r}{au}{v}{a}{n}{\textsubring{d}}{\textsci}} \dicPos{adj}[13] \dicFlx{indecl}[1] \dicDirectTranslationCS{zbloudilý, ztracený} \dicExampleIS{villuráfandi sálir} \dicExampleCS{zbloudilé duše}
\dicEntry[villutrú] \dicTerm{villu··trú} \dicIPA{{v}{\textsci}{\textsubring{d}}{l}{\textscy}{t\smash{\textsuperscript{h}}}{r}{u}} \dicPos{f}[4] \dicFlx{(‑ar)}[27] \dicDirectTranslationCS{bludařství, kacířství, hereze}
\begin{xtolerant}{}{1pt}
\dicEntry[villutrúarmaður] \dicTerm{villu·trúar··|maður}\addthinSS\dicIPA{{v}{\textsci}{\textsubring{d}}{l}{\textscy}{t\smash{\textsuperscript{h}}}{r}{u}{a}{r}{m}{a}{ð}{\textscy}{\textsubring{r}}}\addthinS\dicPos{m}[13]\addthinS\dicFlx{(‑manns,\addthinSS‑menn)}[2] \dicDirectTranslationCS{kacíř(ka), heretik, heretička}
\end{xtolerant}
\dicEntry[Vilníus] \dicTerm{Vilníus} \dicIPA{{v}{\textsci}{l}{n}{i}{j}{\textscy}{s}} \dicPos{subs} \dicFlx{indecl} \dicFieldCat{geog.} \dicDirectTranslationCS{Vilnius} \dicIndirectTranslationCS{(hlavní město Litvy)}
\dicEntry[vilyrði] \dicTerm{vil··yrði} \dicIPA{{v}{\textsci}{\textlengthmark}{l}{\textsci}{r}{ð}{\textsci}} \dicPos{n}[2] \dicFlx{(‑s, ‑)}[14] \dicDirectTranslationCS{příslib, přislíbení};  \dicPhraseIS{gefa vilyrði fyrir e‑u} \dicDirectTranslationCS{přislíbit (co)}
\dicEntry[vin] \dicTerm{vin\smash{\textsuperscript{1}}} \dicIPA{{v}{\textsci}{\textlengthmark}{\textsubring{n}}} \dicPos{m} \dicFlx{sg acc} \dicLink{vinur}
\dicEntry[vin] \dicTerm{vin\smash{\textsuperscript{2}}} \dicIPA{{v}{\textsci}{\textlengthmark}{\textsubring{n}}} \dicPos{f}[4] \dicFlx{(‑jar, ‑jar)}[9] \dicDirectTranslationCS{oáza}
\dicEntry[vina] \dicTerm{vin|a} \dicIPA{{v}{\textsci}{\textlengthmark}{n}{a}} \dicPos{f}[1] \dicFlx{(‑u, ‑ur)}[7] \dicSynonym{vinkona} \dicDirectTranslationCS{přítelkyně, kamarádka}
\dicEntry[vinafólk] \dicTerm{vina··fólk} \dicIPA{{v}{\textsci}{\textlengthmark}{n}{a}{f}{ou}{\textsubring{l}}{\r{g}}} \dicPos{n}[2] \dicFlx{(‑s)}[2] \dicDirectTranslationCS{přátelé, kamarádi}
\dicEntry[vinalaus] \dicTerm{vina··laus} \dicIPA{{v}{\textsci}{\textlengthmark}{n}{a}{l}{\oe i}{s}} \dicPos{adj}[5]\dicFlx{}[-4] \dicDirectTranslationCS{(jsoucí) bez přátel\,/\addthin kamarádů, nemající přátele\,/\addthin kamarády}
\dicEntry[vinalega] \dicTerm{vina··lega} \dicIPA{{v}{\textsci}{\textlengthmark}{n}{a}{l}{\textepsilon}{\textbabygamma}{a}} \dicPos{adv} \dicDirectTranslationCS{přátelsky, kamarádsky, vlídně}
\dicEntry[vinalegur] \dicTerm{vina··legur} \dicIPA{{v}{\textsci}{\textlengthmark}{n}{a}{l}{\textepsilon}{\textbabygamma}{\textscy}{\textsubring{r}}} \dicPos{adj}[1]\dicFlx{}[-8] \dicSynonym{vinsamlegur} \dicDirectTranslationCS{přátelský, kamarádský, vlídný};  \dicPhraseIS{vera vinalegur við e‑n} \dicDirectTranslationCS{chovat se ke (komu) přátelsky}
\dicEntry[vinarbragð] \dicTerm{vinar··|bragð} \dicIPA{{v}{\textsci}{n}{a}{r}{\textsubring{b}}{r}{a}{\textbabygamma}{\texttheta}} \dicPos{n}[2] \dicFlx{(‑bragðs, ‑brögð)}[8] \dicSynonym{vinátta} \dicDirectTranslationCS{přátelskost, přátelství}
\dicEntry[vinarhugur] \dicTerm{vinar··hug|ur} \dicIPA{{v}{\textsci}{\textlengthmark}{n}{a}{\textsubring{r}}{h}{\textscy}{\textbabygamma}{\textscy}{\textsubring{r}}} \dicPos{m}[10] \dicFlx{(‑ar)}[12] \dicDirectTranslationCS{přátelskost, kamarádskost, vlídnost}
\dicEntry[vinarþel] \dicTerm{vinar··þel} \dicIPA{{v}{\textsci}{\textlengthmark}{n}{a}{\textsubring{r}}{\texttheta}{\textepsilon}{\textsubring{l}}} \dicPos{n}[2] \dicFlx{(‑s)}[2] \dicSynonym{vinarhugur} \dicDirectTranslationCS{přátelskost, kamarádskost, vlídnost}
\dicEntry[vinátta] \dicTerm{vin··átt|a} \dicsymFrequent\  \dicIPA{{v}{\textsci}{\textlengthmark}{n}{au}{h}{\textsubring{d}}{a}} \dicPos{f}[1] \dicFlx{(‑u)}[5] \dicDirectTranslationCS{přátelství, kamarádství} \dicExampleIS{Með þeim tókst vinátta sem stóð meðan báðir lifðu.} \dicExampleCS{Vzniklo mezi nimi přátelství, které trvalo, dokud oba žili.}
\dicEntry[vináttuheimsókn] \dicTerm{vin·áttu··heim·sókn} \dicIPA{{v}{\textsci}{\textlengthmark}{n}{au}{h}{\textsubring{d}}{\textscy}{h}{ei}{m}{s}{ou}{h}{\r{g}}{\textsubring{n}}} \dicPos{f}[7] \dicFlx{(‑ar, ‑ir)}[1] \dicDirectTranslationCS{přátelská\,/\addthin oficiální návštěva}
\dicEntry[vináttuvottur] \dicTerm{vin·áttu··vott|ur} \dicIPA{{v}{\textsci}{\textlengthmark}{n}{au}{h}{\textsubring{d}}{\textscy}{v}{\textopeno}{h}{\textsubring{d}}{\textscy}{\textsubring{r}}} \dicPos{m}[6] \dicFlx{(‑s)}[7] \dicDirectTranslationCS{důkaz přátelství}
\dicEntry[vind] \dicTerm{vind} \dicIPA{{v}{\textsci}{n}{\textsubring{d}}} \dicPos{v} \dicFlx{ind praes sg 1 pers} \dicLink{vinda\smash{\textsuperscript{2}}}
\dicEntry[vinda] \dicTerm{vind|a\smash{\textsuperscript{1}}} \dicIPA{{v}{\textsci}{n}{\textsubring{d}}{a}} \dicPos{f}[1] \dicFlx{(‑u, ‑ur)}[19] \textbf{1.} \dicFieldCat{techn.} \dicSynonym{spil} \dicDirectTranslationCS{naviják, rumpál}  \textbf{2.} \dicDirectTranslationCS{ždímačka}
\dicEntry[vinda] \dicTerm{vinda\smash{\textsuperscript{2}}} \dicIPA{{v}{\textsci}{n}{\textsubring{d}}{a}} \dicPos{v}[6] \dicFlx{(vind, vatt, undum, yndi, undið)}[38] \dicFlx{acc\,/\addthin dat} \textbf{1.} \dicFlx{acc} \dicSynonym*{snúa upp á} \dicDirectTranslationCS{(vy)ždímat (prádlo ap.)} \dicExampleIS{vinda þvott} \dicExampleCS{vyždímat prádlo}  \textbf{2.} \dicFlx{acc} \dicSynonym{flétta\smash{\textsuperscript{2}}} \dicDirectTranslationCS{navinout, navíjet, svinout, svíjet, smotat, smotávat} \dicExampleIS{vinda band upp á hnykilinn} \dicExampleCS{svinout provázek do klubíčka};  \dicIdiom{vinda}[að]{ \dicPhraseIS{vinda sér að e‑m}} \dicDirectTranslationCS{otočit se ke (komu)};  \dicIdiom{vinda}[fram]{ \dicPhraseIS{e‑u vindur fram}} \dicFlx{impers} \dicDirectTranslationCS{(co) se odvíjí, (co) pokračuje (kurz ap.)};  \dicIdiom{vinda}[í]{ \dicPhraseIS{vinda sér í að (gera e‑ð)}} \dicDirectTranslationCS{dát\,/\addthin pustit se do (dělání (čeho))};  \dicIdiom{vinda}[upp]{ \dicPhraseIS{vinda upp segl}} \dicDirectTranslationCS{svinout plachty}
\dicEntry[vindasamur] \dicTerm{vinda··|samur} \dicIPA{{v}{\textsci}{n}{\textsubring{d}}{a}{s}{a}{m}{\textscy}{\textsubring{r}}} \dicPos{adj}[1] \dicFlx{(f ‑söm)}[2] \dicDirectTranslationCS{větrný} \dicExampleIS{vindasamur staður} \dicExampleCS{větrné místo};  \dicPhraseIS{það er vindasamt e‑s staðar} \dicFlx{impers} \dicDirectTranslationCS{(kde) je větrno}
\dicEntry[vindátt] \dicTerm{vind··átt} \dicIPA{{v}{\textsci}{n}{\textsubring{d}}{au}{h}{\textsubring{d}}} \dicPos{f}[7] \dicFlx{(‑ar, ‑ir)}[1] \dicDirectTranslationCS{směr větru}
\dicEntry[vindgangur] \dicTerm{vind··gang|ur} \dicIPA{{v}{\textsci}{n}{\textsubring{d}}{\r{g}}{au}{\ng}{\r{g}}{\textscy}{\textsubring{r}}} \dicPos{m}[6] \dicFlx{(‑s)}[7] \dicSynonym*{ókyrrð í maganum} \dicDirectTranslationCS{větry, nadýmání, plynatost}
\dicEntry[vindhani] \dicTerm{vind··han|i} \dicIPA{{v}{\textsci}{n}{\textsubring{d}}{h}{a}{n}{\textsci}} \dicPos{m}[1] \dicFlx{(‑a, ‑ar)}[8] \dicDirectTranslationCS{větrná korouhev\,/\addthin korouhvička}
\dicEntry[vindheldur] \dicTerm{vind··heldur} \dicIPA{{v}{\textsci}{n}{\textsubring{d}}{h}{\textepsilon}{l}{\textsubring{d}}{\textscy}{\textsubring{r}}} \dicPos{adj}[2]\dicFlx{}[-14] \dicDirectTranslationCS{odolný proti větru}
\dicEntry[vindhraði] \dicTerm{vind··hrað|i} \dicIPA{{v}{\textsci}{n}{\textsubring{d}}{\textsubring{r}}{a}{ð}{\textsci}} \dicPos{m}[1] \dicFlx{(‑a)}[3] \dicFieldCat{meteo.} \dicDirectTranslationCS{rychlost větru}
\dicEntry[vindhviða] \dicTerm{vind··hvið|a} \dicIPA{{v}{\textsci}{n}{\textsubring{d}}{k\smash{\textsuperscript{h}}}{v}{\textsci}{ð}{a}} \dicPos{f}[1] \dicFlx{(‑u, ‑ur)}[19] \dicDirectTranslationCS{poryv\,/\addthin náraz větru}
\dicEntry[vindhögg] \dicTerm{vind··högg} \dicIPA{{v}{\textsci}{n}{\textsubring{d}}{h}{\oe}{\r{g}}} \dicPos{n}[2] \dicFlx{(‑s, ‑)}[5] \dicDirectTranslationCS{rána vedle, rána naprázdno};  \dicPhraseIS{slá vindhögg} \dicLangCat{přen.} \dicDirectTranslationCS{šlápnout vedle} \dicIndirectTranslationCS{(zmýlit se)}
\dicEntry[vindill] \dicTerm{vind|ill} \dicIPA{{v}{\textsci}{n}{\textsubring{d}}{\textsci}{\textsubring{d}}{\textsubring{l}}} \dicPos{m}[6] \dicFlx{(‑ils, ‑lar)}[35] \dicDirectTranslationCS{doutník}
\dicEntry[vindjakki] \dicTerm{vind··jakk|i} \dicIPA{{v}{\textsci}{n}{\textsubring{d}}{j}{a}{h}{\r{\textObardotlessj}}{\textsci}} \dicPos{m}[1] \dicFlx{(‑a, ‑ar)}[8] \dicDirectTranslationCS{větrovka}
\dicEntry[vindknúinn] \dicTerm{vind··knúinn} \dicIPA{{v}{\textsci}{n}{\textsubring{d}}{k\smash{\textsuperscript{h}}}{n}{u}{\textsci}{\textsubring{n}}} \dicPos{adj}[6]\dicFlx{}[-6] \dicDirectTranslationCS{poháněný větrem}
\dicEntry[vindlingur] \dicTerm{vind··ling|ur} \dicIPA{{v}{\textsci}{n}{\textsubring{d}}{l}{i}{\ng}{\r{g}}{\textscy}{\textsubring{r}}} \dicPos{m}[6] \dicFlx{(‑s, ‑ar)}[8] \dicSynonym{sígaretta} \dicDirectTranslationCS{cigareta}
\dicEntry[vindmylla] \dicTerm{vind··myll|a} \dicIPA{{v}{\textsci}{n}{\textsubring{d}}{m}{\textsci}{l}{\textlengthmark}{a}} \dicPos{f}[1] \dicFlx{(‑u, ‑ur)}[19] \dicDirectTranslationCS{větrný mlýn};  \dicPhraseIS{berjast við vindmyllur} \dicFlx{refl} \dicLangCat{přen.} \dicDirectTranslationCS{bojovat s~větrnými mlýny}
\dicEntry[vindmælir] \dicTerm{vind··mæl|ir} \dicIPA{{v}{\textsci}{n}{\textsubring{d}}{m}{a}{i}{l}{\textsci}{\textsubring{r}}} \dicPos{m}[7] \dicFlx{(‑is, ‑ar)}[1] \dicFieldCat{meteo.} \dicDirectTranslationCS{anemometr, větroměr}
\dicEntry[vindorka] \dicTerm{vind··ork|a} \dicIPA{{v}{\textsci}{n}{\textsubring{d}}{\textopeno}{\textsubring{r}}{\r{g}}{a}} \dicPos{f}[1] \dicFlx{(‑u)}[5] \dicDirectTranslationCS{větrná síla\,/\addthin energie}
\dicEntry[vindstig] \dicTerm{vind··stig} \dicIPA{{v}{\textsci}{n}{\textsubring{d}}{s}{\textsubring{d}}{\textsci}{x}} \dicPos{n}[2] \dicFlx{(‑s, ‑)}[5] \dicFieldCat{meteo.} \dicDirectTranslationCS{síla větru} \dicIndirectTranslationCS{(udávaná v~metrech za sekundu)}
\dicEntry[vindsæng] \dicTerm{vind··sæng} \dicIPA{{v}{\textsci}{n}{\textsubring{d}}{s}{a}{i}{\ng}{\r{g}}} \dicPos{f}[10] \dicFlx{(‑ur, ‑ur)}[1] \dicDirectTranslationCS{nafukovací matrace}
\dicEntry[vindubrú] \dicTerm{vindu··|brú} \dicIPA{{v}{\textsci}{n}{\textsubring{d}}{\textscy}{\textsubring{b}}{r}{u}} \dicPos{f}[9] \dicFlx{(‑brúar, ‑brýr)}[4] \dicDirectTranslationCS{padací most}
\dicEntry[vindur] \dicTerm{vind|ur} \dicsymFrequent\  \dicIPA{{v}{\textsci}{n}{\textsubring{d}}{\textscy}{\textsubring{r}}} \dicPos{m}[6] \dicFlx{(‑s, ‑ar)}[4] \textbf{1.} \dicDirectTranslationCS{vítr} \dicExampleIS{Hann hreyfir vind.} \dicExampleCS{Fouká.}  \textbf{2.} \dicSynonym*{loft í innyflum} \dicDirectTranslationCS{nadýmání, větry};  \dicPhraseIS{leysa vind} \dicLangCat{přen.} \dicSynonym{freta} \dicDirectTranslationCS{pustit větry, uprdnout si}  \textbf{3.} \dicSynonym{mont} \dicDirectTranslationCS{nadřazenost, arogance};  \dicIdiom{vindur}{ \dicPhraseIS{e‑að fer\,/\addthin hverfur\,/\addthin fýkur út í veður og vind}} \dicLangCat{přen.} \dicDirectTranslationCS{(co) mizí do ztracena}; { \dicPhraseIS{fá vind í seglin}} \dicLangCat{přen.} \dicDirectTranslationCS{dostat vítr do plachet};  \dicPhraseIS{hafa vindinn í fangið} \dicLangCat{přen.} \dicDirectTranslationCS{narazit na odpor, setkat se s~odporem};  \dicPhraseIS{vita\,/\addthin skilja af hvaða átt\,/\addthin hvaðan vindurinn blæs} \dicLangCat{přen.} \dicDirectTranslationCS{vědět, odkud vítr fouká\,/\addthin vane};  \dicPhraseIS{það er allur vindur úr e‑m} \dicFlx{impers} \dicLangCat{přen.} \dicDirectTranslationCS{(kdo) ztrácí poslední zbytek sil}
\dicEntry[vindustigi] \dicTerm{vindu··stig|i} \dicIPA{{v}{\textsci}{n}{\textsubring{d}}{\textscy}{s}{\textsubring{d}}{i}{j}{\textsci}} \dicPos{m}[1] \dicFlx{(‑a, ‑ar)}[1] \dicDirectTranslationCS{točité schody}
\dicEntry[vindþéttur] \dicTerm{vind··þéttur} \dicIPA{{v}{\textsci}{n}{\textsubring{d}}{\texttheta}{j}{\textepsilon}{h}{\textsubring{d}}{\textscy}{\textsubring{r}}} \dicPos{adj}[1]\dicFlx{}[-10] \dicDirectTranslationCS{odolný proti větru}
\dicEntry[vinfastur] \dicTerm{vin··|fastur} \dicIPA{{v}{\textsci}{n}{f}{a}{s}{\textsubring{d}}{\textscy}{\textsubring{r}}} \dicPos{adj}[1] \dicFlx{(f ‑föst)}[11] \dicDirectTranslationCS{věrný (přítel ap.)}
\dicEntry[vingjarnlega] \dicTerm{vin·gjarn··lega} \dicIPA{{v}{\textsci}{n}{\r{\textObardotlessj}}{a}{r}{\textsubring{d}}{n}{l}{\textepsilon}{\textbabygamma}{a}} \dicPos{adv} \dicDirectTranslationCS{přátelsky, přívětivě, vlídně}
\dicEntry[vingjarnlegur] \dicTerm{vin·gjarn··legur} \dicsymFrequent\  \dicIPA{{v}{\textsci}{n}{\r{\textObardotlessj}}{a}{r}{\textsubring{d}}{n}{l}{\textepsilon}{\textbabygamma}{\textscy}{\textsubring{r}}} \dicPos{adj}[1]\dicFlx{}[-8] \dicDirectTranslationCS{přátelský, přívětivý, vlídný} \dicExampleIS{Hann var vingjarnlegur og glaðlegur náungi.} \dicExampleCS{Byl to přátelský a~veselý chlapík.}
\dicEntry[vingjarnleiki] \dicTerm{vin·gjarn··leik|i} \dicIPA{{v}{\textsci}{n}{\r{\textObardotlessj}}{a}{r}{\textsubring{d}}{n}{l}{ei}{\r{\textObardotlessj}}{\textsci}} \dicPos{m}[1] \dicFlx{(‑a)}[3] \dicDirectTranslationCS{přátelskost, přívětivost, vlídnost}
\dicEntry[vingóður] \dicTerm{vin··|góður} \dicIPA{{v}{\textsci}{n}{\r{g}}{ou}{ð}{\textscy}{\textsubring{r}}} \dicPos{adj}[11] \dicFlx{(comp ‑betri, sup ‑bestur)}[1] \dicPhraseIS{eiga vingott við e‑n} {\textbf{a.}} \dicSynonym*{vera vinur e‑rs} \dicDirectTranslationCS{být s~(kým) zadobře, přátelit se s~(kým)};  {\textbf{b.}} \dicSynonym*{hafa kynmök við e‑n} \dicDirectTranslationCS{mít aférku s~(kým)}
\dicEntry[vingsa] \dicTerm{vings|a} \dicIPA{{v}{i}{\ng}{s}{a}} \dicPos{v}[1] \dicFlx{(‑aði)}[1] \dicFlx{dat} \dicSynonym{sveifla\smash{\textsuperscript{2}}} \dicDirectTranslationCS{kývat, klátit, mávat} \dicExampleIS{vingsa handleggnum} \dicExampleCS{mávat rukama}
\dicEntry[vingull] \dicTerm{ving|ull} \dicIPA{{v}{i}{\ng}{\r{g}}{\textscy}{\textsubring{d}}{\textsubring{l}}} \dicPos{m}[6] \dicFlx{(‑uls, ‑lar)}[31] \dicFieldCat{bot.} \dicDirectTranslationCS{kostřava} \textit{(l.~{\textLA{Festuca}})}  \dicsymPhoto\ 
\dicFigure{82936.jpg}{Vingull}{Vingull - Antušek Ivo, Biolib, PD/CC-BY-NC}
\dicEntry[vingulslegur] \dicTerm{vinguls··legur} \dicIPA{{v}{i}{\ng}{\r{g}}{\textscy}{l}{s}{l}{\textepsilon}{\textbabygamma}{\textscy}{\textsubring{r}}} \dicPos{adj}[1]\dicFlx{}[-8] \dicDirectTranslationCS{nestálý, přelétavý, vrtkavý}
\dicEntry[vink] \dicTerm{vink} \dicIPA{{v}{i}{\r{\ng}}{\r{g}}} \dicPos{n}[2] \dicFlx{(‑s, ‑)}[5] \dicDirectTranslationCS{mávnutí, zamávání}
\dicEntry[vinka] \dicTerm{vink|a} \dicIPA{{v}{i}{\r{\ng}}{\r{g}}{a}} \dicPos{v}[1] \dicFlx{(‑aði)}[1] \dicSynonym{veifa\smash{\textsuperscript{2}}} \dicDirectTranslationCS{(za)mávat} \dicExampleIS{vinka e‑m bless} \dicExampleCS{mávat (komu) na rozloučenou}
\dicEntry[vinkill] \dicTerm{vink|ill} \dicIPA{{v}{i}{\r{\textltailn}}{\r{\textObardotlessj}}{\textsci}{\textsubring{d}}{\textsubring{l}}} \dicPos{m}[6] \dicFlx{(‑ils, ‑lar)}[35] \textbf{1.} \dicSynonym{horn} \dicDirectTranslationCS{pravý úhel, vinkl, vingl}  \textbf{2.} \dicFieldCat{techn.} \dicDirectTranslationCS{úhelník} \dicIndirectTranslationCS{(nástroj)}  \textbf{3.} \dicSynonym{sjónarhorn} \dicDirectTranslationCS{úhel pohledu}
\dicEntry[vinkona] \dicTerm{vin··kon|a} \dicsymFrequent\  \dicIPA{{v}{\textsci}{n}{k\smash{\textsuperscript{h}}}{\textopeno}{n}{a}} \dicPos{f}[1] \dicFlx{(‑u, ‑ur)}[27] \dicDirectTranslationCS{kamarádka, přítelkyně} \dicExampleIS{Ertu búin að tala við vinkonur hennar?} \dicExampleCS{Mluvila jsi s~jejími kamarádkami?}
\dicEntry[vinmargur] \dicTerm{vin··|margur} \dicIPA{{v}{\textsci}{n}{m}{a}{r}{\r{g}}{\textscy}{\textsubring{r}}} \dicPos{adj}[11] \dicFlx{(f ‑mörg, comp ‑fleiri, sup ‑flestur)}[8] \dicDirectTranslationCS{mající mnoho kamarádů\,/\addthin přátel}
\dicEntry[vinn] \dicTerm{vinn} \dicIPA{{v}{\textsci}{n}{\textlengthmark}} \dicPos{v} \dicFlx{ind praes sg 1 pers} \dicLink{vinna\smash{\textsuperscript{2}}}
\dicEntry[vinna] \dicTerm{vinn|a\smash{\textsuperscript{1}}} \dicsymFrequent\  \dicIPA{{v}{\textsci}{n}{\textlengthmark}{a}} \dicPos{f}[1] \dicFlx{(‑u)}[5] \textbf{1.} \dicSynonym{verk} \dicDirectTranslationCS{práce, úkol, úsilí, námaha} \dicExampleIS{Þetta verður mikil vinna.} \dicExampleCS{To bude velká práce.}  \textbf{2.} \dicSynonym{atvinna} \dicDirectTranslationCS{práce, zaměstnání} \dicExampleIS{hafa vinnu} \dicExampleCS{mít práci};  \dicPhraseIS{eftir vinnu} \dicFlx{adv} \dicDirectTranslationCS{po práci\,/\addthin zaměstnání}  \textbf{3.} \dicSynonym{handbragð} \dicDirectTranslationCS{práce, vykonání, provedení}  \textbf{4.} \dicFieldCat{fyz.} \dicDirectTranslationCS{práce}
\dicEntry[vinna] \dicTerm{vinna\smash{\textsuperscript{2}}} \dicsymFrequent\  \dicIPA{{v}{\textsci}{n}{\textlengthmark}{a}} \dicPos{v}[8] \dicFlx{(vinn, vann, unnum, ynni, unnið)}[6] \dicFlx{acc} \textbf{1.} \dicDirectTranslationCS{pracovat, dělat, být zaměstnán} \dicExampleIS{vinna á skrifstofu} \dicExampleCS{pracovat v~kanceláři};  \dicPhraseIS{vinna eins og skepna} \dicLangCat{přen.} \dicDirectTranslationCS{dřít jako kůň};  \dicPhraseIS{vinna í fiski} \dicDirectTranslationCS{zpracovávat ryby}  \textbf{2.} \dicFlx{acc} \dicDirectTranslationCS{vykonat, vykonávat} \dicExampleIS{vinna frábært starf} \dicExampleCS{vykonat výbornou práci}  \textbf{3.} \dicFlx{dat + acc} \dicDirectTranslationCS{přinést, způsobit (užitek ap.)}  \textbf{4.} \dicFlx{acc} \dicDirectTranslationCS{zpracovat, zpracovávat (surovinu ap.)}  \textbf{5.} \dicFlx{acc} \dicDirectTranslationCS{vyhrát, vyhrávat, (z)vítězit} \dicExampleIS{vinna keppni} \dicExampleCS{vyhrát závody};  \dicPhraseIS{vinna sigur} \dicDirectTranslationCS{zvítězit};  \dicIdiom{vinna}[að]{ \dicPhraseIS{vinna að e‑u}} \dicDirectTranslationCS{pracovat na (čem), dělat na (čem)} \dicExampleIS{vinna að orðabókinni} \dicExampleCS{pracovat na slovníku};  \dicIdiom{vinna}[af]{ \dicPhraseIS{vinna e‑ð af sér}} \dicDirectTranslationCS{odpracovat si (co)} \dicExampleIS{vinna tvo daga af sér} \dicExampleCS{odpracovat si dva dny};  \dicIdiom{vinna}[á]{ \dicPhraseIS{e‑að vinnur ekki á e‑u}} \dicDirectTranslationCS{(co) nefunguje na (co)}; { \dicPhraseIS{vinna á}} \dicDirectTranslationCS{zlepšit se, stát se lepším (kniha ke konci ap.)}; { \dicPhraseIS{vinna bug á e‑u}} \dicDirectTranslationCS{vyhrát nad (čím), zvítězit nad (čím), pokořit (co)} \dicExampleIS{vinna bug á atvinnuleysi} \dicExampleCS{zvítězit nad nezaměstnaností};  \dicIdiom{vinna}[fyrir]{ \dicPhraseIS{vinna fyrir e‑u}} \dicDirectTranslationCS{pracovat za (co) (plat ap.)}; { \dicPhraseIS{vinna fyrir sér}} \dicDirectTranslationCS{vydělat si, vydělávat si} \dicIndirectTranslationCS{ (na živobytí)};  \dicIdiom{vinna}[gegn]{ \dicPhraseIS{vinna gegn e‑u}} \dicDirectTranslationCS{pracovat proti (čemu), mařit (co)};  \dicIdiom{vinna}[inn]{ \dicPhraseIS{vinna sér inn (peninga)}} \dicDirectTranslationCS{vydělávat (peníze)}; { \dicPhraseIS{vinna sér inn fyrir e‑u}} \dicDirectTranslationCS{vydělávat na (co), šetřit si na (co)};  \dicIdiom{vinna}[í]{ \dicPhraseIS{vinna í e‑u}} \dicDirectTranslationCS{pracovat na (čem), pracovat s~(čím)};  \dicIdiom{vinna}[til]{ \dicPhraseIS{vilja allt til vinna}} \dicLangCat{přen.} \dicDirectTranslationCS{být připraven udělat vše}; { \dicPhraseIS{vinna til e‑s}} \dicDirectTranslationCS{získat (co), vyhrát (co)} \dicExampleIS{Hann hefur unnið til verðlauna.} \dicExampleCS{Vyhrál medaili.};  \dicIdiom{vinna}[upp]{ \dicPhraseIS{vinna e‑ð upp}} \dicDirectTranslationCS{dohnat (co) (zpoždění ap.)}; { \dicPhraseIS{vinna sig upp}} \dicDirectTranslationCS{vypracovat se};  \dicIdiom{vinna}[úr]{ \dicPhraseIS{vinna e‑ð úr e‑u}} \dicDirectTranslationCS{získat (co) z~(čeho), extrahovat (co) z~(čeho)}; { \dicPhraseIS{vinna e‑ð úr jörðu}} \dicDirectTranslationCS{vydobýt (co) ze země}; { \dicPhraseIS{vinna úr e‑u}} {\textbf{a.}} \dicDirectTranslationCS{zpracovat (co), zpracovávat (co), analyzovat (co)} \dicExampleIS{vinna úr öllum þessum gögnum} \dicExampleCS{zpracovat všechna ta data};  {\textbf{b.}} \dicDirectTranslationCS{překonat (co) (trauma ap.)};  \dicIdiom{vinnast}{ \dicPhraseIS{vinnast}} \dicFlx{refl} {\textbf{a.}} \dicDirectTranslationCS{postupovat, jít kupředu} \dicExampleIS{Verkið vannst fljótt.} \dicExampleCS{Práce šla rychle kupředu.};  {\textbf{b.}} \dicDirectTranslationCS{být získán\,/\addthin vyhrán (vítězství ap.)}; { \dicPhraseIS{e‑m vinnst ekki tími til að (gera e‑ð)}} \dicFlx{refl impers} \dicDirectTranslationCS{(kdo) se nedostane k~(udělání (čeho)), (kdo) nemá dostatek času k~(udělání (čeho))}; { \dicPhraseIS{það vinnst ekki tími til e‑s}} \dicFlx{refl impers} \dicDirectTranslationCS{není dostatek času k~(čemu)}; { \dicPhraseIS{það vinnst (ekkert) með e‑u}} \dicFlx{refl impers} \dicDirectTranslationCS{(nic) se (čím) nezíská}; { \dicPhraseIS{ef tími vinnst til}} \dicFlx{refl} \dicDirectTranslationCS{pokud čas dovolí}
\dicEntry[vinnandi] \dicTerm{vinn··andi} \dicIPA{{v}{\textsci}{n}{\textlengthmark}{a}{n}{\textsubring{d}}{\textsci}} \dicPos{adj}[13] \dicFlx{indecl}[1] \dicDirectTranslationCS{pracující} \dicExampleIS{vinnandi fólk} \dicExampleCS{pracující lidé}
\dicEntry[vinnanlegur] \dicTerm{vinnan··legur} \dicIPA{{v}{\textsci}{n}{\textlengthmark}{a}{n}{l}{\textepsilon}{\textbabygamma}{\textscy}{\textsubring{r}}} \dicPos{adj}[1]\dicFlx{}[-8] \dicDirectTranslationCS{zpracovatelný, dobyvatelný}
\dicEntry[vinningur] \dicTerm{vinn··ing|ur} \dicIPA{{v}{\textsci}{n}{\textlengthmark}{i}{\ng}{\r{g}}{\textscy}{\textsubring{r}}} \dicPos{m}[6] \dicFlx{(‑s, ‑ar)}[8] \textbf{1.} \dicSynonym{gróði} \dicDirectTranslationCS{výhra} \dicIndirectTranslationCS{(peníze nebo předmět získaný ve hře)}  \textbf{2.} \dicSynonym{sigur} \dicDirectTranslationCS{výhra, vítězství} \dicExampleIS{hlaupa til vinnings} \dicExampleCS{běžet si pro vítězství};  \dicPhraseIS{hafa vinninginn} \dicDirectTranslationCS{zvítězit, vyhrát, odnést si vítězství}
\dicEntry[vinnsla] \dicTerm{vinnsl|a} \dicIPA{{v}{\textsci}{n}{s}{\textsubring{d}}{l}{a}} \dicPos{f}[1] \dicFlx{(‑u, ‑ur)}[19] \dicDirectTranslationCS{zpracování, opracování, obrobení} \dicExampleIS{vinnsla á hráefninu} \dicExampleCS{zpracování surovin}
\dicEntry[vinnsluaðferð] \dicTerm{vinnslu··að·ferð} \dicIPA{{v}{\textsci}{n}{s}{\textsubring{d}}{l}{\textscy}{a}{ð}{f}{\textepsilon}{r}{\texttheta}} \dicPos{f}[7] \dicFlx{(‑ar, ‑ir)}[1] \dicDirectTranslationCS{metoda zpracování}
\dicEntry[vinnslukostnaður] \dicTerm{vinnslu··kost·nað|ur} \dicIPA{{v}{\textsci}{n}{s}{\textsubring{d}}{l}{\textscy}{k\smash{\textsuperscript{h}}}{\textopeno}{s}{\textsubring{d}}{n}{a}{ð}{\textscy}{\textsubring{r}}} \dicPos{m}[10] \dicFlx{(‑ar)}[9] \dicDirectTranslationCS{náklady na zpracování}
\dicEntry[vinnsluminni] \dicTerm{vinnslu··minni} \dicIPA{{v}{\textsci}{n}{s}{\textsubring{d}}{l}{\textscy}{m}{\textsci}{n}{\textsci}} \dicPos{n}[2] \dicFlx{(‑s, ‑)}[14] \dicFieldCat{poč.} \dicDirectTranslationCS{operační paměť, RAM} \dicAntonym{lesminni}
\dicEntry[vinnslustöð] \dicTerm{vinnslu··stöð} \dicIPA{{v}{\textsci}{n}{s}{\textsubring{d}}{l}{\textscy}{s}{\textsubring{d}}{\oe}{\texttheta}} \dicPos{f}[6] \dicFlx{(‑var, ‑var)}[1] \dicDirectTranslationCS{výrobna, zpracovatelský závod}
\dicEntry[vinnuaðferð] \dicTerm{vinnu··að·ferð} \dicIPA{{v}{\textsci}{n}{\textlengthmark}{\textscy}{a}{ð}{f}{\textepsilon}{r}{\texttheta}} \dicPos{f}[7] \dicFlx{(‑ar, ‑ir)}[1] \dicDirectTranslationCS{pracovní postup\,/\addthin metoda}
\dicEntry[vinnuafl] \dicTerm{vinnu··afl} \dicIPA{{v}{\textsci}{n}{\textlengthmark}{\textscy}{a}{\textsubring{b}}{\textsubring{l}}} \dicPos{n}[2] \dicFlx{(‑s)}[2] \dicDirectTranslationCS{pracovní síla} \dicIndirectTranslationCS{(pracovníci)}
\dicEntry[vinnuaflsfrekur] \dicTerm{vinnu·afls··frekur} \dicIPA{{v}{\textsci}{n}{\textlengthmark}{\textscy}{a}{\textsubring{b}}{\textsubring{l}}{s}{f}{r}{\textepsilon}{\r{g}}{\textscy}{\textsubring{r}}} \dicPos{adj}[1]\dicFlx{}[-1] \dicDirectTranslationCS{náročný na pracovní sílu}
\dicEntry[vinnuálag] \dicTerm{vinnu··á·lag} \dicIPA{{v}{\textsci}{n}{\textlengthmark}{\textscy}{au}{l}{a}{x}} \dicPos{n}[2] \dicFlx{(‑s)}[2] \dicDirectTranslationCS{pracovní zátěž}
\dicEntry[vinnuborð] \dicTerm{vinnu··borð} \dicIPA{{v}{\textsci}{n}{\textlengthmark}{\textscy}{\textsubring{b}}{\textopeno}{r}{\texttheta}} \dicPos{n}[2] \dicFlx{(‑s, ‑)}[5] \dicDirectTranslationCS{pracovní stůl}
\dicEntry[vinnubók] \dicTerm{vinnu··|bók} \dicIPA{{v}{\textsci}{n}{\textlengthmark}{\textscy}{\textsubring{b}}{ou}{\r{g}}} \dicPos{f}[8] \dicFlx{(‑bókar, ‑bækur)}[5] \dicDirectTranslationCS{pracovní sešit, cvičebnice}
\dicEntry[vinnubrögð] \dicTerm{vinnu··brögð} \dicIPA{{v}{\textsci}{n}{\textlengthmark}{\textscy}{\textsubring{b}}{r}{\oe}{\textbabygamma}{\texttheta}} \dicPos{n}[2] \dicFlx{pl}[9] \dicDirectTranslationCS{pracovní přístup, styl práce} \dicExampleIS{góð vinnubrögð} \dicExampleCS{dobrý pracovní přístup}
\dicEntry[vinnudagur] \dicTerm{vinnu··dag|ur} \dicIPA{{v}{\textsci}{n}{\textlengthmark}{\textscy}{\textsubring{d}}{a}{\textbabygamma}{\textscy}{\textsubring{r}}} \dicPos{m}[6] \dicFlx{(‑s, ‑ar)}[62] \dicDirectTranslationCS{pracovní den} \dicExampleIS{tíu tíma vinnudagur} \dicExampleCS{desetihodinový pracovní den}
\dicEntry[vinnudeila] \dicTerm{vinnu··deil|a} \dicIPA{{v}{\textsci}{n}{\textlengthmark}{\textscy}{\textsubring{d}}{ei}{l}{a}} \dicPos{f}[1] \dicFlx{(‑u, ‑ur)}[13] \dicDirectTranslationCS{pracovní spor}
\dicEntry[vinnudýr] \dicTerm{vinnu··dýr} \dicIPA{{v}{\textsci}{n}{\textlengthmark}{\textscy}{\textsubring{d}}{i}{\textsubring{r}}} \dicPos{n}[2] \dicFlx{(‑s, ‑)}[5] \dicDirectTranslationCS{pracovní zvíře, zvíře využívané na práci}
\dicEntry[vinnuflokkur] \dicTerm{vinnu··flokk|ur} \dicIPA{{v}{\textsci}{n}{\textlengthmark}{\textscy}{f}{l}{\textopeno}{h}{\r{g}}{\textscy}{\textsubring{r}}} \dicPos{m}[6] \dicFlx{(‑s, ‑ar)}[8] \dicDirectTranslationCS{pracovní kolektiv\,/\addthin skupina\,/\addthin četa}
\dicEntry[vinnufólk] \dicTerm{vinnu··fólk} \dicIPA{{v}{\textsci}{n}{\textlengthmark}{\textscy}{f}{ou}{\textsubring{l}}{\r{g}}} \dicPos{n}[2] \dicFlx{(‑s)}[2] \dicDirectTranslationCS{pracovníci} \dicIndirectTranslationCS{(zejména na farmě)}
\dicEntry[vinnufrekur] \dicTerm{vinnu··frekur} \dicIPA{{v}{\textsci}{n}{\textlengthmark}{\textscy}{f}{r}{\textepsilon}{\r{g}}{\textscy}{\textsubring{r}}} \dicPos{adj}[1]\dicFlx{}[-1] \dicDirectTranslationCS{pracný, náročný (na práci)}
\dicEntry[vinnufriður] \dicTerm{vinnu··frið|ur} \dicIPA{{v}{\textsci}{n}{\textlengthmark}{\textscy}{f}{r}{\textsci}{ð}{\textscy}{\textsubring{r}}} \dicPos{m}[10] \dicFlx{(‑ar)}[7] \dicDirectTranslationCS{klid na práci} \dicExampleIS{fá vinnufrið} \dicExampleCS{mít klid na práci}
\dicEntry[vinnufús] \dicTerm{vinnu··fús} \dicIPA{{v}{\textsci}{n}{\textlengthmark}{\textscy}{f}{u}{s}} \dicPos{adj}[5]\dicFlx{}[-1] \dicDirectTranslationCS{ochotný pracovat}
\dicEntry[vinnufær] \dicTerm{vinnu··fær} \dicIPA{{v}{\textsci}{n}{\textlengthmark}{\textscy}{f}{a}{i}{\textsubring{r}}} \dicPos{adj}[5] \dicFlx{(f ‑)}[8] \dicDirectTranslationCS{práceschopný, schopný práce}
\dicEntry[vinnuföt] \dicTerm{vinnu··föt} \dicIPA{{v}{\textsci}{n}{\textlengthmark}{\textscy}{f}{\oe}{\textsubring{d}}} \dicPos{n}[2] \dicFlx{pl}[9] \dicDirectTranslationCS{pracovní oblečení}
\dicEntry[vinnuhanski] \dicTerm{vinnu··hansk|i} \dicIPA{{v}{\textsci}{n}{\textlengthmark}{\textscy}{h}{a}{n}{s}{\r{\textObardotlessj}}{\textsci}} \dicPos{m}[1] \dicFlx{(‑a, ‑ar)}[8] \dicDirectTranslationCS{pracovní rukavice}
\dicEntry[vinnuherbergi] \dicTerm{vinnu··her·bergi} \dicIPA{{v}{\textsci}{n}{\textlengthmark}{\textscy}{h}{\textepsilon}{r}{\textsubring{b}}{\textepsilon}{r}{\r{\textObardotlessj}}{\textsci}} \dicPos{n}[2] \dicFlx{(‑s, ‑)}[16] \dicDirectTranslationCS{pracovna, pracovní místnost}
\dicEntry[vinnuhópur] \dicTerm{vinnu··hóp|ur} \dicIPA{{v}{\textsci}{n}{\textlengthmark}{\textscy}{h}{ou}{\textsubring{b}}{\textscy}{\textsubring{r}}} \dicPos{m}[6] \dicFlx{(‑s, ‑ar)}[22] \dicDirectTranslationCS{pracovní skupina\,/\addthin tým}
\dicEntry[vinnukona] \dicTerm{vinnu··kon|a} \dicIPA{{v}{\textsci}{n}{\textlengthmark}{\textscy}{k\smash{\textsuperscript{h}}}{\textopeno}{n}{a}} \dicPos{f}[1] \dicFlx{(‑u, ‑ur)}[27] \dicDirectTranslationCS{služebná, služka, komorná}
\dicEntry[vinnulaun] \dicTerm{vinnu··laun} \dicIPA{{v}{\textsci}{n}{\textlengthmark}{\textscy}{l}{\oe i}{\textsubring{n}}} \dicPos{n}[2] \dicFlx{pl}[1] \dicSynonym*{laun} \dicDirectTranslationCS{plat, mzda}
\dicEntry[vinnulaus] \dicTerm{vinnu··laus} \dicIPA{{v}{\textsci}{n}{\textlengthmark}{\textscy}{l}{\oe i}{s}} \dicPos{adj}[5]\dicFlx{}[-1] \textbf{1.} \dicSynonym{atvinnulaus} \dicDirectTranslationCS{nezaměstnaný, (jsoucí) bez práce\,/\addthin zaměstnání}  \textbf{2.} \dicSynonym{iðjulaus} \dicDirectTranslationCS{nečinný, nezaměstnaný}
\dicEntry[vinnumaður] \dicTerm{vinnu··|maður} \dicIPA{{v}{\textsci}{n}{\textlengthmark}{\textscy}{m}{a}{ð}{\textscy}{\textsubring{r}}} \dicPos{m}[13] \dicFlx{(‑manns, ‑menn)}[2] \dicDirectTranslationCS{čeledín, služebník}
\dicEntry[vinnumarkaður] \textls[15]{\dicTerm{vinnu··mark·að|ur}\dicTerm{, atvinnumarkaður} \dicIPA{{v}\-{\textsci}\-{n}\-{\textlengthmark}\-{\textscy}\-{m}\-{a}\-{\textsubring{r}}\-{\r{g}}\-{a}\-{ð}\-{\textscy}\-{\textsubring{r}}}} \dicPos{m}[10] \dicFlx{(‑ar\,/\addthin ‑s, ‑ir)}[31] \dicFieldCat{ekon.} \dicDirectTranslationCS{trh práce}
\dicEntry[vinnumiðlun] \dicTerm{vinnu··miðl|un} \dicIPA{{v}{\textsci}{n}{\textlengthmark}{\textscy}{m}{\textsci}{ð}{l}{\textscy}{\textsubring{n}}} \dicPos{f}[7] \dicFlx{(‑unar, ‑anir)}[8] \textbf{1.} \dicDirectTranslationCS{zprostředkování práce}  \textbf{2.} \dicDirectTranslationCS{úřad práce, zprostředkovatelna práce}
\dicEntry[vinnupallur] \dicTerm{vinnu··pall|ur} \dicIPA{{v}{\textsci}{n}{\textlengthmark}{\textscy}{p\smash{\textsuperscript{h}}}{a}{\textsubring{d}}{l}{\textscy}{\textsubring{r}}} \dicPos{m}[6] \dicFlx{(‑s, ‑ar)}[5] \dicDirectTranslationCS{lešení}
\dicEntry[vinnuréttur] \dicTerm{vinnu··rétt|ur} \dicIPA{{v}{\textsci}{n}{\textlengthmark}{\textscy}{r}{j}{\textepsilon}{h}{\textsubring{d}}{\textscy}{\textsubring{r}}} \dicPos{m}[10] \dicFlx{(‑ar)}[7] \dicFieldCat{práv.} \dicDirectTranslationCS{pracovní právo}
\dicEntry[vinnusamningur] \dicTerm{vinnu··samn·ing|ur} \dicIPA{{v}{\textsci}{n}{\textlengthmark}{\textscy}{s}{a}{m}{n}{i}{\ng}{\r{g}}{\textscy}{\textsubring{r}}} \dicPos{m}[6] \dicFlx{(‑s, ‑ar)}[8] \dicFieldCat{práv.} \dicDirectTranslationCS{pracovní smlouva}
\dicEntry[vinnusamur] \dicTerm{vinnu··|samur} \dicIPA{{v}{\textsci}{n}{\textlengthmark}{\textscy}{s}{a}{m}{\textscy}{\textsubring{r}}} \dicPos{adj}[1] \dicFlx{(f ‑söm)}[2] \dicDirectTranslationCS{pracovitý, pilný}
\dicEntry[vinnuseðill] \dicTerm{vinnu··seð|ill} \dicIPA{{v}{\textsci}{n}{\textlengthmark}{\textscy}{s}{\textepsilon}{ð}{\textsci}{\textsubring{d}}{\textsubring{l}}} \dicPos{m}[6] \dicFlx{(‑ils, ‑lar)}[35] \dicDirectTranslationCS{pracovní výkaz}
\dicEntry[vinnusemi] \dicTerm{vinnu··sem|i} \dicIPA{{v}{\textsci}{n}{\textlengthmark}{\textscy}{s}{\textepsilon}{m}{\textsci}} \dicPos{f}[2] \dicFlx{(‑i)}[2] \dicDirectTranslationCS{pracovitost, píle}
\dicEntry[vinnusjúkdómur] \dicTerm{vinnu··sjúk·dóm|ur} \dicIPA{{v}{\textsci}{n}{\textlengthmark}{\textscy}{s}{j}{u}{\r{g}}{\textsubring{d}}{ou}{m}{\textscy}{\textsubring{r}}} \dicPos{m}[6] \dicFlx{(‑s, ‑ar)}[10] \dicDirectTranslationCS{nemoc z~povolání}
\dicEntry[vinnuskilyrði] \dicTerm{vinnu··skil·yrði} \dicIPA{{v}{\textsci}{n}{\textlengthmark}{\textscy}{s}{\r{\textObardotlessj}}{\textsci}{l}{\textsci}{r}{ð}{\textsci}} \dicPos{n}[2] \dicFlx{(‑s, ‑)}[14] \dicDirectTranslationCS{pracovní podmínky}
\dicEntry[vinnuskylda] \dicTerm{vinnu··skyld|a} \dicIPA{{v}{\textsci}{n}{\textlengthmark}{\textscy}{s}{\r{\textObardotlessj}}{\textsci}{l}{\textsubring{d}}{a}} \dicPos{f}[1] \dicFlx{(‑u)}[5] \dicDirectTranslationCS{pracovní závazek}
\dicEntry[vinnuslys] \dicTerm{vinnu··slys} \dicIPA{{v}{\textsci}{n}{\textlengthmark}{\textscy}{s}{\textsubring{d}}{l}{\textsci}{s}} \dicPos{n}[2] \dicFlx{(‑s, ‑)}[5] \dicDirectTranslationCS{pracovní nehoda\,/\addthin úraz}
\dicEntry[vinnustaður] \dicTerm{vinnu··stað|ur} \dicsymFrequent\  \dicIPA{{v}{\textsci}{n}{\textlengthmark}{\textscy}{s}{\textsubring{d}}{a}{ð}{\textscy}{\textsubring{r}}} \dicPos{m}[10] \dicFlx{(‑ar, ‑ir)}[14] \dicDirectTranslationCS{pracoviště, místo práce} \dicExampleIS{spenna og álag á vinnustað} \dicExampleCS{napětí a~zátěž na pracovišti}
\dicEntry[vinnustofa] \dicTerm{vinnu··stof|a} \dicIPA{{v}{\textsci}{n}{\textlengthmark}{\textscy}{s}{\textsubring{d}}{\textopeno}{v}{a}} \dicPos{f}[1] \dicFlx{(‑u, ‑ur)}[7] \textbf{1.} \dicSynonym*{vinnuhúsnæði} \dicDirectTranslationCS{pracovna, ateliér, dílna}  \textbf{2.} \dicSynonym{fræðsla} \dicDirectTranslationCS{pracovní seminář, dílna, workshop}
\dicEntry[vinnustund] \dicTerm{vinnu··stund} \dicIPA{{v}{\textsci}{n}{\textlengthmark}{\textscy}{s}{\textsubring{d}}{\textscy}{n}{\textsubring{d}}} \dicPos{f}[7] \dicFlx{(‑ar, ‑ir)}[1] \dicDirectTranslationCS{pracovní hodina}
\dicEntry[vinnustúlka] \dicTerm{vinnu··stúlk|a} \dicIPA{{v}{\textsci}{n}{\textlengthmark}{\textscy}{s}{\textsubring{d}}{u}{\textsubring{l}}{\r{g}}{a}} \dicPos{f}[1] \dicFlx{(‑u, ‑ur)}[13] \dicDirectTranslationCS{služebná, služka}
\dicEntry[vinnustöð] \dicTerm{vinnu··stöð} \dicIPA{{v}{\textsci}{n}{\textlengthmark}{\textscy}{s}{\textsubring{d}}{\oe}{\texttheta}} \dicPos{f}[6] \dicFlx{(‑var, ‑var)}[1] \dicFieldCat{poč.} \dicDirectTranslationCS{pracovní stanice}
\dicEntry[vinnustöðvun] \dicTerm{vinnu··stöðv|un} \dicIPA{{v}{\textsci}{n}{\textlengthmark}{\textscy}{s}{\textsubring{d}}{\oe}{ð}{v}{\textscy}{\textsubring{n}}} \dicPos{f}[7] \dicFlx{(‑unar, ‑anir)}[8] \dicSynonym{verkfall} \dicDirectTranslationCS{přerušení práce, stávka}
\dicEntry[vinnutilgáta] \dicTerm{vinnu··til·gát|a} \dicIPA{{v}{\textsci}{n}{\textlengthmark}{\textscy}{t\smash{\textsuperscript{h}}}{\textsci}{l}{\r{g}}{au}{\textsubring{d}}{a}} \dicPos{f}[1] \dicFlx{(‑u, ‑ur)}[19] \dicDirectTranslationCS{pracovní hypotéza}
\dicEntry[vinnutími] \dicTerm{vinnu··tím|i} \dicIPA{{v}{\textsci}{n}{\textlengthmark}{\textscy}{t\smash{\textsuperscript{h}}}{i}{m}{\textsci}} \dicPos{m}[1] \dicFlx{(‑a, ‑ar)}[1] \dicDirectTranslationCS{pracovní doba} \dicExampleIS{í vinnutímanum} \dicExampleCS{v~pracovní době}
\begin{xtolerant}{}{1pt}
\dicEntry[vinnuveitandi] \dicTerm{vinnu··veit·\hspace{-0.06667em}|andi}\addthinSS\dicIPA{{v}{\textsci}{n}{\textlengthmark}{\textscy}{v}{ei}{\textsubring{d}}{a}{n}{\textsubring{d}}{\textsci}}\addthinS\dicPos{m}[2]\addthinS\dicFlx{(‑anda,\addthinSS‑endur)}[1]\addthinSS\dicSynonym{atvin\-nurekandi} \dicDirectTranslationCS{zaměstnavatel(ka), poskytovatel(ka) práce}
\end{xtolerant}
\dicEntry[vinnuvika] \dicTerm{vinnu··vik|a} \dicIPA{{v}{\textsci}{n}{\textlengthmark}{\textscy}{v}{\textsci}{\r{g}}{a}} \dicPos{f}[1] \dicFlx{(‑u, ‑ur)}[13] \dicDirectTranslationCS{pracovní týden}
\dicEntry[vinnuþegi] \dicTerm{vinnu··þeg|i} \dicIPA{{v}{\textsci}{n}{\textlengthmark}{\textscy}{\texttheta}{ei}{\textsci}} \dicPos{m}[1] \dicFlx{(‑a, ‑ar)}[1] \dicSynonym{launþegi} \dicDirectTranslationCS{zaměstnanec, zaměstnankyně}
\dicEntry[vinnuþjarkur] \dicTerm{vinnu··þjark|ur} \dicIPA{{v}{\textsci}{n}{\textlengthmark}{\textscy}{\texttheta}{j}{a}{\textsubring{r}}{\r{g}}{\textscy}{\textsubring{r}}} \dicPos{m}[6] \dicFlx{(‑s, ‑ar)}[5] \dicDirectTranslationCS{pracant(ka), dříč(ka)}
\dicEntry[vinnuþol] \dicTerm{vinnu··þol} \dicIPA{{v}{\textsci}{n}{\textlengthmark}{\textscy}{\texttheta}{\textopeno}{\textsubring{l}}} \dicPos{n}[2] \dicFlx{(‑s)}[2] \dicDirectTranslationCS{pracovní výdrž\,/\addthin vytrvalost}
\dicEntry[vinnuþrek] \dicTerm{vinnu··þrek} \dicIPA{{v}{\textsci}{n}{\textlengthmark}{\textscy}{\texttheta}{r}{\textepsilon}{\r{g}}} \dicPos{n}[2] \dicFlx{(‑s)}[2] \dicSynonym{starfsorka} \dicDirectTranslationCS{pracovní schopnost\,/\addthin kapacita}
\dicEntry[vinsa] \dicTerm{vins|a} \dicIPA{{v}{\textsci}{n}{s}{a}} \dicPos{v}[1] \dicFlx{(‑aði)}[1] \dicFlx{acc} \dicPhraseIS{vinsa e‑ð úr e‑u} \dicDirectTranslationCS{vybrat\,/\addthin vylovit (co) z~(čeho)} \dicExampleIS{vinsa það besta úr dótinu} \dicExampleCS{vybrat to nejlepší z~věcí}
\dicEntry[vinsamlega] \dicTerm{vin·sam··lega} \dicIPA{{v}{\textsci}{n}{s}{a}{m}{l}{\textepsilon}{\textbabygamma}{a}} \dicPos{adv} \textbf{1.} \dicDirectTranslationCS{laskavě, přátelsky}  \textbf{2.} \dicLangCat{form.} \dicDirectTranslationCS{prosím}
\dicEntry[vinsamlegast] \dicTerm{vin·sam··legast} \dicIPA{{v}{\textsci}{n}{s}{a}{m}{l}{\textepsilon}{\textbabygamma}{a}{s}{\textsubring{d}}} \dicPos{adv} \dicLangCat{form.} \dicDirectTranslationCS{prosím} \dicExampleIS{Vinsamlegast ekki snerta sýningargripina!} \dicExampleCS{Nedotýkat se prosím exponátů!}
\dicEntry[vinsamlegur] \dicTerm{vin·sam··legur} \dicIPA{{v}{\textsci}{n}{s}{a}{m}{l}{\textepsilon}{\textbabygamma}{\textscy}{\textsubring{r}}} \dicPos{adj}[1]\dicFlx{}[-8] \dicSynonym{vingjarnlegur} \dicDirectTranslationCS{přátelský, přívětivý, vlídný}
\dicEntry[vinsemd] \dicTerm{vin··semd} \dicIPA{{v}{\textsci}{n}{s}{\textepsilon}{m}{\textsubring{d}}} \dicPos{f}[7] \dicFlx{(‑ar)}[3] \dicDirectTranslationCS{přátelskost, přívětivost} \dicExampleIS{sýna e‑m vinsemd} \dicExampleCS{být ke (komu) přátelský}
\dicEntry[vinskapur] \dicTerm{vin··skap|ur} \dicIPA{{v}{\textsci}{n}{s}{\r{g}}{a}{\textsubring{b}}{\textscy}{\textsubring{r}}} \dicPos{m}[10] \dicFlx{(‑ar)}[15] \dicDirectTranslationCS{přátelství, kamarádství} \dicExampleIS{Það er mikill vinskapur á milli þeirra.} \dicExampleCS{Jsou to velcí kamarádi.}
\dicEntry[vinslit] \dicTerm{vin··slit} \dicIPA{{v}{\textsci}{n}{s}{\textsubring{d}}{l}{\textsci}{\textsubring{d}}} \dicPos{n}[2] \dicFlx{pl}[1] \dicDirectTranslationCS{konec\,/\addthin rozpad přátelství\,/\addthin kamarádství}
\dicEntry[vinstri] \dicTerm{vinstr|i} \dicsymFrequent\  \dicIPA{{v}{\textsci}{n}{s}{\textsubring{d}}{r}{\textsci}} \dicPos{adj}[12] \dicFlx{(n ‑a)}[3] \textbf{1.} \dicDirectTranslationCS{levý} \dicAntonym{hægri};  \dicPhraseIS{á vinstri hönd} \dicFlx{adv} \dicDirectTranslationCS{nalevo, po levé ruce};  \dicPhraseIS{til vinstri} \dicFlx{adv} \dicDirectTranslationCS{doleva, vlevo}  \textbf{2.} \dicDirectTranslationCS{levicový};  \dicPhraseIS{vinstri maður} \dicDirectTranslationCS{levicově orientovaný člověk}
\dicEntry[vinstriflokkur] \dicTerm{vinstri··flokk|ur} \dicIPA{{v}{\textsci}{n}{s}{\textsubring{d}}{r}{\textsci}{f}{l}{\textopeno}{h}{\r{g}}{\textscy}{\textsubring{r}}} \dicPos{m}[6] \dicFlx{(‑s, ‑ar)}[8] \dicDirectTranslationCS{levicová strana} \dicAntonym{hægriflokkur}
\dicEntry[vinstrimaður] \dicTerm{vinstri··|maður} \dicIPA{{v}{\textsci}{n}{s}{\textsubring{d}}{r}{\textsci}{m}{a}{ð}{\textscy}{\textsubring{r}}} \dicPos{m}[13] \dicFlx{(‑manns, ‑menn)}[2] \dicDirectTranslationCS{levicově orientovaný člověk, stoupenec\,/\addthin stoupenkyně levice, levičák, levičačka} \dicAntonym{hægrimaður}
\dicEntry[vinstrisinnaður] \dicTerm{vinstri··sinn·|aður} \dicIPA{{v}{\textsci}{n}{s}{\textsubring{d}}{r}{\textsci}{s}{\textsci}{n}{a}{ð}{\textscy}{\textsubring{r}}} \dicPos{adj}[3] \dicFlx{(f ‑uð)}[3] \dicDirectTranslationCS{levicový, levicově smýšlející\,/\addthin orientovaný} \dicAntonym{hægrisinnaður}
\dicEntry[vinstrisinni] \dicTerm{vinstri··sinn|i} \dicIPA{{v}{\textsci}{n}{s}{\textsubring{d}}{r}{\textsci}{s}{\textsci}{n}{\textsci}} \dicPos{m}[1] \dicFlx{(‑a, ‑ar)}[1] \dicDirectTranslationCS{levicově orientovaný člověk, stoupenec\,/\addthin stoupenkyně levice} \dicAntonym{hægrisinni}
\dicEntry[vinstristjórn] \dicTerm{vinstri··stjórn}\dicTerm{, vinstri stjórn} \dicIPA{{v}\-{\textsci}\-{n}\-{s}\-{\textsubring{d}}\-{r}\-{\textsci}\-{s}\-{\textsubring{d}}\-{j}\-{ou}\-{r}\-{\textsubring{d}}\-{\textsubring{n}}\-} \dicPos{f}[7] \dicFlx{(‑ar, ‑ir)}[1] \dicDirectTranslationCS{levicová vláda} \dicAntonym{hægristjórn}
\dicEntry[vinstúlka] \dicTerm{vin··stúlk|a} \dicIPA{{v}{\textsci}{n}{s}{\textsubring{d}}{u}{\textsubring{l}}{\r{g}}{a}} \dicPos{f}[1] \dicFlx{(‑u, ‑ur)}[13] \dicSynonym{vinkona} \dicDirectTranslationCS{kamarádka, přítelkyně}
\dicEntry[vinsæld] \dicTerm{vin··sæld} \dicIPA{{v}{\textsci}{n}{s}{a}{i}{l}{\textsubring{d}}} \dicPos{f}[7] \dicFlx{(‑ar, ‑ir)}[1] \dicPhraseIS{vinsældir} \dicFlx{pl} \dicDirectTranslationCS{popularita, oblíbenost}
\dicEntry[vinsæll] \dicTerm{vin··sæll} \dicsymFrequent\  \dicIPA{{v}{\textsci}{n}{s}{a}{i}{\textsubring{d}}{\textsubring{l}}} \dicPos{adj}[8]\dicFlx{}[-1] \dicDirectTranslationCS{populární, oblíbený} \dicExampleIS{vinsæll leikari} \dicExampleCS{oblíbený herec}
\dicEntry[vinur] \dicTerm{vin|ur} \dicsymFrequent\  \dicIPA{{v}{\textsci}{\textlengthmark}{n}{\textscy}{\textsubring{r}}} \dicPos{m}[10] \dicFlx{(‑ar, ‑ir)}[17] \dicDirectTranslationCS{přítel, kamarád} \dicExampleIS{Við erum góðir vinir.} \dicExampleCS{Jsme dobří kamarádi.};  \dicPhraseIS{e‑jum er vel til vina} \dicFlx{impers} \dicDirectTranslationCS{(kdo) jsou dobří kamarádi}
\dicEntry[vinveittur] \dicTerm{vin··veittur} \dicIPA{{v}{\textsci}{n}{v}{ei}{h}{\textsubring{d}}{\textscy}{\textsubring{r}}} \dicPos{adj}[1]\dicFlx{}[-10] \dicDirectTranslationCS{přátelský, přívětivý, přátelsky naladěný};  \dicPhraseIS{vera e‑m vinveittur} \dicDirectTranslationCS{být (komu) přátelsky nakloněn}
\dicEntry[vippa] \dicTerm{vipp|a} \dicIPA{{v}{\textsci}{h}{\textsubring{b}}{a}} \dicPos{v}[1] \dicFlx{(‑aði)}[5] \dicFlx{dat} \dicDirectTranslationCS{přehoupnout, přeskočit};  \dicPhraseIS{vippa sér yfir e‑ð} \dicDirectTranslationCS{přehoupnout se přes (co)}
\dicEntry[virða] \dicTerm{vir|ða} \dicsymFrequent\  \dicIPA{{v}{\textsci}{r}{ð}{a}} \dicPos{v}[2] \dicFlx{(‑ti, ‑t)}[51] \dicFlx{acc (+ gen)} \textbf{1.} \dicFlx{acc} \dicSynonym*{meta til verðs} \dicDirectTranslationCS{ocenit, ohodnotit, určit cenu (domu ap.)} \dicExampleIS{virða húsið á tvær milljónir} \dicExampleCS{ocenit dům na dva milióny}  \textbf{2.} \dicFlx{acc (+ gen)} \dicSynonym*{sýna virðingu} \dicDirectTranslationCS{cenit si, vážit si, respektovat, uznávat} \dicExampleIS{virða hann mikils} \dicExampleCS{moc si ho vážit};  \dicPhraseIS{virða e‑n ekki viðlits} \dicDirectTranslationCS{nevěnovat (komu) pozornost, ignorovat (koho)};  \dicPhraseIS{virða e‑n ekki svars} \dicDirectTranslationCS{neobtěžovat se (komu) odpovědět};  \dicIdiom{virða}[að]{ \dicPhraseIS{virða e‑ð að vettugi}} \dicDirectTranslationCS{ignorovat (co), nevšímat si (čeho), nebrat na (co) ohled};  \dicIdiom{virða}[fyrir]{ \dicPhraseIS{virða e‑ð fyrir sér}} \dicDirectTranslationCS{prohlédnout\,/\addthin prohlížet si (co), (pro)zkoumat (co)} \dicExampleIS{virða myndina vandlega fyrir sér} \dicExampleCS{důkladně si prohlédnout obrázek};  \dicIdiom{virða}[við]{ \dicPhraseIS{virða e‑ð við e‑n}} \dicDirectTranslationCS{cenit si (co\,/\addthin čeho) na (kom), respektovat (koho) pro (co)};  \dicIdiom{virðast}{ \dicPhraseIS{e‑að\,/\addthin e‑r virðist}} \dicFlx{refl} \dicDirectTranslationCS{(co\,/\addthin kdo) se zdá, (co\,/\addthin kdo) se jeví, (co\,/\addthin kdo) vypadá} \dicExampleIS{Hann virðist vera hræddur.} \dicExampleCS{Vypadá to, že se bojí.}; { \dicPhraseIS{e‑m virðist}} \dicFlx{refl impers} \dicDirectTranslationCS{(komu) se zdá, (komu) se jeví}
\dicEntry[virði] \dicTerm{virði} \dicsymFrequent\  \dicIPA{{v}{\textsci}{r}{ð}{\textsci}} \dicPos{n}[2] \dicFlx{(‑s, ‑)}[14] \dicSynonym{gildi} \dicDirectTranslationCS{hodnota, cena};  \dicPhraseIS{einskis virði} \dicFlx{adj} \dicDirectTranslationCS{bezcenný, nemající žádnou hodnotu}
\dicEntry[virðing] \dicTerm{virð··ing} \dicsymFrequent\  \dicIPA{{v}{\textsci}{r}{ð}{i}{\ng}{\r{g}}} \dicPos{f}[4] \dicFlx{(‑ar, ‑ar)}[5] \textbf{1.} \dicSynonym{mat\smash{\textsuperscript{1}}} \dicDirectTranslationCS{ocenění, odhad ceny}  \textbf{2.} \dicDirectTranslationCS{úcta, respekt , uznání} \dicExampleIS{virðing fyrir lífinu} \dicExampleCS{úcta k~životu};  \dicPhraseIS{e‑að er fyrir neðan virðingu e‑rs} \dicDirectTranslationCS{(co) je pod (čí) úroveň};  \dicPhraseIS{bera virðingu fyrir e‑m} \dicDirectTranslationCS{vážit si (koho), respektovat (koho), mít (koho) v~úctě};  \dicPhraseIS{með fullri virðingu fyrir e‑m} \dicDirectTranslationCS{při vší úctě ke (komu)}
\dicEntry[virðingarleysi] \dicTerm{virðingar··leysi} \dicIPA{{v}{\textsci}{r}{ð}{i}{\ng}{\r{g}}{a}{r}{l}{ei}{s}{\textsci}} \dicPos{n}[2] \dicFlx{(‑s)}[20] \dicDirectTranslationCS{neúcta, despekt}
\dicEntry[virðingarsess] \dicTerm{virðingar··sess} \dicIPA{{v}{\textsci}{r}{ð}{i}{\ng}{\r{g}}{a}{\textsubring{r}}{s}{\textepsilon}{s}} \dicPos{m}[4] \dicFlx{(‑, ‑ar)}[35] \dicDirectTranslationCS{čestné místo, čestná pozice (v~úřadě ap.)}
\dicEntry[virðingarverður] \dicTerm{virðingar··verður} \dicIPA{{v}{\textsci}{r}{ð}{i}{\ng}{\r{g}}{a}{r}{v}{\textepsilon}{r}{ð}{\textscy}{\textsubring{r}}} \dicPos{adj}[2]\dicFlx{}[-1] \textbf{1.} \dicDirectTranslationCS{úctyhodný, vážený, respektovaný}  \textbf{2.} \dicDirectTranslationCS{úctyhodný, chvályhodný} \dicExampleIS{virðingarverð viðleitni} \dicExampleCS{úctyhodná snaha}
\dicEntry[virðingarvottur] \dicTerm{virðingar··vott|ur} \dicIPA{{v}{\textsci}{r}{ð}{i}{\ng}{\r{g}}{a}{r}{v}{\textopeno}{h}{\textsubring{d}}{\textscy}{\textsubring{r}}} \dicPos{m}[6] \dicFlx{(‑s)}[7] \dicDirectTranslationCS{uznání, pocta, počest}
\dicEntry[virðisaukaskattur] \dicTerm{virðis·auka··skatt|ur} \dicIPA{{v}\-{\textsci}\-{r}\-{ð}\-{\textsci}\-{s}\-{\oe i}\-{\r{g}}\-{a}\-{s}\-{\r{g}}\-{a}\-{h}\-{\textsubring{d}}\-{\textscy}\-{\textsubring{r}}\-} \dicPos{m}[6] \dicFlx{(‑s)}[7] \dicFieldCat{ekon.} \dicDirectTranslationCS{daň z~přidané hodnoty}
\dicEntry[virðisauki] \dicTerm{virðis··auk|i} \dicIPA{{v}{\textsci}{r}{ð}{\textsci}{s}{\oe i}{\r{\textObardotlessj}}{\textsci}} \dicPos{m}[1] \dicFlx{(‑a)}[3] \dicFieldCat{ekon.} \dicDirectTranslationCS{přidaná hodnota}
\dicEntry[virðulega] \dicTerm{virðu··lega} \dicIPA{{v}{\textsci}{r}{ð}{\textscy}{l}{\textepsilon}{\textbabygamma}{a}} \dicPos{adv} \dicDirectTranslationCS{s~úctou\,/\addthin respektem}
\dicEntry[virðulegur] \dicTerm{virðu··legur} \dicIPA{{v}{\textsci}{r}{ð}{\textscy}{l}{\textepsilon}{\textbabygamma}{\textscy}{\textsubring{r}}} \dicPos{adj}[1]\dicFlx{}[-8] \dicDirectTranslationCS{vážený, respektovaný, ctěný, ctihodný} \dicExampleIS{virðulegu gestir} \dicExampleCS{vážení hosté};  \dicPhraseIS{virðulegi forseti} \dicDirectTranslationCS{vážený pane prezidente}
\dicEntry[virðuleiki] \dicTerm{virðu··leik|i} \dicIPA{{v}{\textsci}{r}{ð}{\textscy}{l}{ei}{\r{\textObardotlessj}}{\textsci}} \dicPos{m}[1] \dicFlx{(‑a)}[3] \dicDirectTranslationCS{váženost, respekt, důstojnost, ctihodnost}
\dicEntry[virka] \dicTerm{virk|a} \dicIPA{{v}{\textsci}{\textsubring{r}}{\r{g}}{a}} \dicPos{v}[1] \dicFlx{(‑aði)}[1] \dicDirectTranslationCS{fungovat, pracovat, jít} \dicExampleIS{Síminn virkar ekki.} \dicExampleCS{Telefon nefunguje.};  \dicPhraseIS{virka á e‑n} \dicDirectTranslationCS{fungovat na (koho), mít vliv na (koho)}
\dicEntry[virki] \dicTerm{virk|i\smash{\textsuperscript{1}}} \dicIPA{{v}{\textsci}{\textsubring{r}}{\r{\textObardotlessj}}{\textsci}} \dicPos{m}[1] \dicFlx{(‑ja, ‑jar)}[14] \textbf{1.} \dicFieldCat{mat.} \dicDirectTranslationCS{operátor}  \textbf{2.} \dicIndirectTranslationCS{přípona označující profesi} \dicExampleIS{rafvirki} \dicExampleCS{elektrikář}
\dicEntry[virki] \dicTerm{virki\smash{\textsuperscript{2}}} \dicIPA{{v}{\textsci}{\textsubring{r}}{\r{\textObardotlessj}}{\textsci}} \dicPos{n}[2] \dicFlx{(‑s, ‑)}[16] \dicSynonym{vígi} \dicDirectTranslationCS{pevnost, citadela, tvrz}
\dicEntry[virkilega] \dicTerm{virki··lega} \dicsymFrequent\  \dicIPA{{v}{\textsci}{\textsubring{r}}{\r{\textObardotlessj}}{\textsci}{l}{\textepsilon}{\textbabygamma}{a}} \dicPos{adv} \textbf{1.} \dicSynonym{raunverulega} \dicDirectTranslationCS{skutečně, opravdu, vskutku} \dicExampleIS{Ætlarðu virkilega að flækjast um Afríku?} \dicExampleCS{Opravdu se chceš toulat po Africe?}  \textbf{2.} \dicSynonym*{til áherslu} \dicDirectTranslationCS{skutečně, opravdu} \dicIndirectTranslationCS{(k~zdůraznění)}
\dicEntry[virkilegur] \dicTerm{virki··legur} \dicIPA{{v}{\textsci}{\textsubring{r}}{\r{\textObardotlessj}}{\textsci}{l}{\textepsilon}{\textbabygamma}{\textscy}{\textsubring{r}}} \dicPos{adj}[1]\dicFlx{}[-8] \dicDirectTranslationCS{skutečný, opravdový, faktický}
\dicEntry[virkja] \dicTerm{virkj|a} \dicIPA{{v}{\textsci}{\textsubring{r}}{\r{\textObardotlessj}}{a}} \dicPos{v}[1] \dicFlx{(‑aði)}[1] \dicFlx{acc} \textbf{1.} \dicDirectTranslationCS{spoutat, využít, využívat (sílu přírody na výrobu elektřiny ap.)} \dicExampleIS{virkja foss} \dicExampleCS{využívat vodopád na výrobu energie}  \textbf{2.} \dicDirectTranslationCS{zaměstnat, zaneprázdnit} \dicIndirectTranslationCS{(přidělením úkolu)}
\dicEntry[virkjun] \dicTerm{virkj|un} \dicIPA{{v}{\textsci}{\textsubring{r}}{\r{\textObardotlessj}}{\textscy}{\textsubring{n}}} \dicPos{f}[7] \dicFlx{(‑unar, ‑anir)}[8] \dicSynonym{orkuver} \dicDirectTranslationCS{elektrárna}
\dicEntry[virkni] \dicTerm{virkn|i} \dicIPA{{v}{\textsci}{\textsubring{r}}{\r{g}}{n}{\textsci}} \dicPos{f}[3] \dicFlx{(‑i)}[3] \textbf{1.} \dicSynonym*{það hvernig e‑að virkar} \dicDirectTranslationCS{účinek, efekt}  \textbf{2.} \dicSynonym*{það að vera virkur} \dicDirectTranslationCS{aktivita, činnost, činorodost}
\dicEntry[virkt] \dicTerm{virkt} \dicIPA{{v}{\textsci}{\textsubring{r}}{\textsubring{d}}} \dicPos{f}[7] \dicFlx{(‑ar, ‑ir)}[1] \dicPhraseIS{með virktum} \dicFlx{adv} \dicDirectTranslationCS{srdečně, vřele}
\dicEntry[virkur] \dicTerm{virkur} \dicsymFrequent\  \dicIPA{{v}{\textsci}{\textsubring{r}}{\r{g}}{\textscy}{\textsubring{r}}} \dicPos{adj}[1]\dicFlx{}[-1] \textbf{1.} \dicSynonym*{með virkni} \dicDirectTranslationCS{účinný, efektivní}  \textbf{2.} \dicSynonym{starfsamur} \dicDirectTranslationCS{aktivní, činný, činorodý} \dicExampleIS{virkur þátttakandi} \dicExampleCS{aktivní účastník};  \textls[15]{\dicPhraseIS{taka virkan þátt í e‑u} \dicDirectTranslationCS{aktivně se zúčastnit (čeho)}  \textbf{3.} \dicDirectTranslationCS{pracovní} \dicIndirectTranslationCS{(o~dni)} \dicExampleIS{virkur dagur} \dicExampleCS{\mbox{pracovní den}}}
\dicEntry[virtur] \dicTerm{virtur} \dicIPA{{v}{\textsci}{\textsubring{r}}{\textsubring{d}}{\textscy}{\textsubring{r}}} \dicPos{adj}[1]\dicFlx{}[-10] \dicDirectTranslationCS{uznávaný, ctěný, respektovaný}
\dicEntry[visinn] \dicTerm{visinn} \dicIPA{{v}{\textsci}{\textlengthmark}{s}{\textsci}{\textsubring{n}}} \dicPos{adj}[6]\dicFlx{}[-2] \textbf{1.} \dicDirectTranslationCS{povadlý, zvadlý, uschlý (list ap.)} \dicExampleIS{visið laufblað} \dicExampleCS{zvadlý list}  \textbf{2.} \dicDirectTranslationCS{vetchý, sešlý, zesláblý (stařec ap.)}
\dicEntry[viska] \dicTerm{visk|a\smash{\textsuperscript{1}}} \dicIPA{{v}{\textsci}{s}{\r{g}}{a}} \dicPos{f}[1] \dicFlx{(‑u)}[5] \dicSynonym{vísdómur} \dicDirectTranslationCS{moudrost}
\dicEntry[viska] \dicTerm{visk|a\smash{\textsuperscript{2}}} \dicIPA{{v}{\textsci}{s}{\r{g}}{a}} \dicPos{v}[1] \dicFlx{(‑aði)}[1] \dicFlx{acc} \dicPhraseIS{viska út} \dicLangCat{zast.} \dicSynonym*{þurrka út} \dicDirectTranslationCS{vymazat, smazat}
\dicEntry[viskí] \dicTerm{viskí} \dicIPA{{v}{\textsci}{s}{\r{\textObardotlessj}}{i}} \dicPos{n}[2] \dicFlx{(‑s)}[2] \dicDirectTranslationCS{whisky, viska}
\dicEntry[viskusteinn] \dicTerm{visku··stein|n} \dicIPA{{v}{\textsci}{s}{\r{g}}{\textscy}{s}{\textsubring{d}}{ei}{\textsubring{d}}{\textsubring{n}}} \dicPos{m}[6] \dicFlx{(‑s, ‑ar)}[42] \dicDirectTranslationCS{kámen mudrců} \textit{(l.~{\textLA{lapis philosophorum}})}
\dicEntry[viskustykki] \dicTerm{visku··stykki} \dicIPA{{v}{\textsci}{s}{\r{g}}{\textscy}{s}{\textsubring{d}}{\textsci}{h}{\r{\textObardotlessj}}{\textsci}} \dicPos{n}[2] \dicFlx{(‑s, ‑)}[16] \dicDirectTranslationCS{utěrka}
\dicEntry[viskutré] \dicTerm{visku··tré} \dicIPA{{v}{\textsci}{s}{\r{g}}{\textscy}{t\smash{\textsuperscript{h}}}{r}{j}{\textepsilon}} \dicPos{n}[2] \dicFlx{(‑s, ‑)}[36] \dicFieldCat{náb.} \dicDirectTranslationCS{strom poznání} \textit{(l.~{\textLA{lignum scientiae}})}
\dicEntry[Visla] \dicTerm{Visl|a} \dicIPA{{v}{\textsci}{s}{\textsubring{d}}{l}{a}} \dicPos{f}[1] \dicFlx{(‑u)}[6] \dicFieldCat{geog.} \dicDirectTranslationCS{Visla} \dicIndirectTranslationCS{(řeka v~Polsku)}
\dicEntry[visna] \dicTerm{visn|a} \dicIPA{{v}{\textsci}{s}{\textsubring{d}}{n}{a}} \dicPos{v}[1] \dicFlx{(‑aði)}[1] \dicSynonym*{verða visinn} \dicDirectTranslationCS{(z)vadnout, chřadnout, uschnout, usychat} \dicExampleIS{Grasið visnar fljótt.} \dicExampleCS{Tráva rychle vadne.}
\dicEntry[visnaður] \textls[15]{\dicTerm{visn··|aður} \dicIPA{{v}{\textsci}{s}{\textsubring{d}}{\textsubring{n}}{a}{ð}{\textscy}{\textsubring{r}}} \dicPos{adj}[3] \dicFlx{(f ‑uð)}[4] \dicDirectTranslationCS{zvadlý, uvadlý, uschlý (list ap.)}}
\dicEntry[visnun] \dicTerm{visn|un} \dicIPA{{v}{\textsci}{s}{\textsubring{d}}{n}{\textscy}{\textsubring{n}}} \dicPos{f}[7] \dicFlx{(‑unar)}[9] \dicDirectTranslationCS{(z)vadnutí, chřadnutí, usychání}
\dicEntry[viss] \dicTerm{viss} \dicsymFrequent\  \dicIPA{{v}{\textsci}{s}{\textlengthmark}} \dicPos{adj}[5]\dicFlx{}[-1] \textbf{1.} \dicSynonym{öruggur} \dicDirectTranslationCS{jistý, zaručený, nepochybný};  \dicPhraseIS{það er vissara að (gera e‑ð)} \dicDirectTranslationCS{je jistější (udělat (co))}  \textbf{2.} \dicSynonym{ákveðinn} \dicDirectTranslationCS{jistý, určitý, stanovený, domluvený};  \dicPhraseIS{á vissan hátt} \dicFlx{adv} \dicDirectTranslationCS{určitým způsobem};  \dicPhraseIS{við vissar aðstæður} \dicFlx{adv} \dicDirectTranslationCS{za jistých podmínek}
\dicEntry[vissa] \dicTerm{viss|a} \dicsymFrequent\  \dicIPA{{v}{\textsci}{s}{\textlengthmark}{a}} \dicPos{f}[1] \dicFlx{(‑u)}[5] \dicDirectTranslationCS{jistota, určitost} \dicExampleIS{vita e‑ð með vissu} \dicExampleCS{vědět (co) s~jistotou};  \dicPhraseIS{hafa vissu fyrir e‑u} \dicDirectTranslationCS{mít jistotu o~(čem)}
\dicEntry[vissi] \dicTerm{vissi} \dicIPA{{v}{\textsci}{s}{\textlengthmark}{\textsci}} \dicPos{v} \dicFlx{ind\,/\addthin con pf sg 1 pers} \dicLink{vita\smash{\textsuperscript{1}}}
\dicEntry[vissulega] \dicTerm{vissu··lega} \dicsymFrequent\  \dicIPA{{v}{\textsci}{s}{\textlengthmark}{\textscy}{l}{\textepsilon}{\textbabygamma}{a}} \dicPos{adv} \dicDirectTranslationCS{jistě, určitě, rozhodně, nepochybně} \dicIndirectTranslationCS{(k~zdůraznění)} \dicExampleIS{Það er vissulega rétt.} \dicExampleCS{Určitě je to správně.}
\dicEntry[vissum] \dicTerm{vissum} \dicIPA{{v}{\textsci}{s}{\textlengthmark}{\textscy}{\textsubring{m}}} \dicPos{v} \dicFlx{ind pf pl 1 pers} \dicLink{vita\smash{\textsuperscript{1}}}
\dicEntry[vist] \dicTerm{vist\smash{\textsuperscript{1}}} \dicIPA{{v}{\textsci}{s}{\textsubring{d}}} \dicPos{f}[7] \dicFlx{(‑ar, ‑ir)}[1] \textbf{1.} \dicSynonym{dvöl} \dicDirectTranslationCS{pobyt, přebývání}  \textbf{2.} \dicSynonym*{vinnukonustarf} \dicDirectTranslationCS{výpomoc, pomocná práce, drobná práce v~domácnosti} \dicIndirectTranslationCS{(práce na dobu určitou v~soukromém domě s~úkoly jako je uklízení, hlídání dětí ap.)};  \dicPhraseIS{vera í vist hjá e‑m} \dicDirectTranslationCS{vypomáhat (komu), být na výpomoc u~(koho)}  \textbf{3.} \dicPhraseIS{vistir} \dicFlx{pl} \dicDirectTranslationCS{zásoby (potravin), proviant}
\dicEntry[vist] \dicTerm{vist\smash{\textsuperscript{2}}} \dicIPA{{v}{\textsci}{s}{\textsubring{d}}} \dicPos{f}[7] \dicFlx{(‑ar)}[3] \dicDirectTranslationCS{whist}
\dicEntry[vista] \dicTerm{vist|a} \dicIPA{{v}{\textsci}{s}{\textsubring{d}}{a}} \dicPos{v}[1] \dicFlx{(‑aði)}[1] \dicFlx{acc} \textbf{1.} \dicDirectTranslationCS{umístit, umísťovat} \dicIndirectTranslationCS{(zejména o~úřadech)} \dicExampleIS{vista e‑n e‑s staðar} \dicExampleCS{umístit (koho kde)}  \textbf{2.} \dicFieldCat{poč.} \dicDirectTranslationCS{uložit, ukládat};  \dicIdiom{vistast}{ \dicPhraseIS{vistast}} \dicFlx{refl} {\textbf{a.}} \dicDirectTranslationCS{žít, bydlet, být ubytován};  {\textbf{b.}} \dicFieldCat{poč.} \dicDirectTranslationCS{uložit se, ukládat se}
\dicEntry[vistfang] \dicTerm{vist··|fang} \dicIPA{{v}{\textsci}{s}{\textsubring{d}}{f}{au}{\ng}{\r{g}}} \dicPos{n}[2] \dicFlx{(‑fangs, ‑föng)}[8] \dicFieldCat{poč.} \dicDirectTranslationCS{adresa}
\dicEntry[vistfræði] \dicTerm{vist··fræð|i} \dicIPA{{v}{\textsci}{s}{\textsubring{d}}{f}{r}{a}{i}{ð}{\textsci}} \dicPos{f}[3] \dicFlx{(‑i)}[3] \dicDirectTranslationCS{ekologie}
\dicEntry[vistfræðingur] \dicTerm{vist·fræð··ing|ur} \dicIPA{{v}{\textsci}{s}{\textsubring{d}}{f}{r}{a}{i}{ð}{i}{\ng}{\r{g}}{\textscy}{\textsubring{r}}} \dicPos{m}[6] \dicFlx{(‑s, ‑ar)}[8] \dicDirectTranslationCS{ekolog, ekoložka}
\dicEntry[vistheimili] \dicTerm{vist··heimili} \dicIPA{{v}{\textsci}{s}{\textsubring{d}}{h}{ei}{m}{\textsci}{l}{\textsci}} \dicPos{n}[2] \dicFlx{(‑s, ‑)}[14] \dicDirectTranslationCS{ústav, (pečovatelský) domov, léčebna} \dicExampleIS{vistheimili drykkjusjúklinga} \dicExampleCS{protialkoholní léčebna}
\dicEntry[vistkerfi] \dicTerm{vist··kerfi} \dicIPA{{v}{\textsci}{s}{\textsubring{d}}{c\smash{\textsuperscript{h}}}{\textepsilon}{r}{v}{\textsci}} \dicPos{n}[2] \dicFlx{(‑s, ‑)}[14] \dicDirectTranslationCS{ekosystém} \dicExampleIS{vistkerfi sjávar} \dicExampleCS{ekosystém moře}
\dicEntry[vistlegur] \dicTerm{vist··legur} \dicIPA{{v}{\textsci}{s}{\textsubring{d}}{l}{\textepsilon}{\textbabygamma}{\textscy}{\textsubring{r}}} \dicPos{adj}[1]\dicFlx{}[-8] \dicDirectTranslationCS{pohodlný, útulný, příjemný}
\dicEntry[vistmaður] \dicTerm{vist··|maður} \dicIPA{{v}{\textsci}{s}{\textsubring{d}}{m}{a}{ð}{\textscy}{\textsubring{r}}} \dicPos{m}[13] \dicFlx{(‑manns, ‑menn)}[2] \dicDirectTranslationCS{chovanec, chovankyně (ústavu ap.)}
\dicEntry[vistun] \dicTerm{vist|un} \dicIPA{{v}{\textsci}{s}{\textsubring{d}}{\textscy}{\textsubring{n}}} \dicPos{f}[7] \dicFlx{(‑unar, ‑anir)}[8] \textbf{1.} \dicDirectTranslationCS{umístění, ubytování (pacienta ap.)}  \textbf{2.} \dicDirectTranslationCS{uložení, archivace (informací ap.)}
\dicEntry[vistvænn] \dicTerm{vist··vænn} \dicIPA{{v}{\textsci}{s}{\textsubring{d}}{v}{a}{i}{\textsubring{d}}{\textsubring{n}}} \dicPos{adj}[7]\dicFlx{}[-1] \dicDirectTranslationCS{ekologický, šetrný k~životnímu prostředí} \dicExampleIS{vistvænt farartæki} \dicExampleCS{ekologický dopravní prostředek}
\dicEntry[vit] \dicTerm{vit} \dicIPA{{v}{\textsci}{\textlengthmark}{\textsubring{d}}} \dicPos{n}[2] \dicFlx{(‑s, ‑)}[5] \textbf{1.} \dicDirectTranslationCS{rozum, inteligence, znalost, vědomost}  \textbf{2.} \dicPhraseIS{vit} \dicFlx{pl} \dicDirectTranslationCS{nos a~pusa};  \dicPhraseIS{taka fyrir vitin} \dicDirectTranslationCS{zacpat si nos a~pusu};  \dicIdiom{vit}{ \dicPhraseIS{að (mínu) viti}} \dicFlx{adv} \dicDirectTranslationCS{podle (mého) názoru}; { \dicPhraseIS{ganga á vit e‑s}} \dicDirectTranslationCS{vydat se na (co) (dobrodružství ap.)};  \dicPhraseIS{hafa allt sitt vit úr e‑u} \dicDirectTranslationCS{mít všechny znalosti z~(čeho), mít vše z~(čeho) (knih ap.)};  \dicPhraseIS{hafa vit á e‑u} \dicDirectTranslationCS{mít smysl pro (co), mít oko na (co)};  \dicPhraseIS{hafa vit fyrir e‑m} \dicDirectTranslationCS{dohlížet na (koho), mít (koho) na očích};  \dicPhraseIS{koma vitinu fyrir e‑n} \dicDirectTranslationCS{přivést (koho) k~rozumu};  \dicPhraseIS{missa vitið} \dicDirectTranslationCS{zbláznit se, přijít o~rozum};  \dicPhraseIS{stíga ekki í vitið} \dicDirectTranslationCS{být mdlého rozumu};  \dicPhraseIS{vera genginn af vitinu} \dicDirectTranslationCS{přijít o~rozum, zbláznit se};  \dicPhraseIS{vera kominn til vits og ára} \dicDirectTranslationCS{být starším a~moudřejším};  \dicPhraseIS{vera vel viti borinn} \dicDirectTranslationCS{být chytrý\,/\addthin inteligentní}; { \dicPhraseIS{vera viti sínu fjær}} \dicLangCat{přen.} \dicDirectTranslationCS{být jak smyslů zbavený};  \dicPhraseIS{vita sínu viti} \dicDirectTranslationCS{být dostatečně moudrý};  \dicPhraseIS{það er ekki\,/\addthin ekkert vit í e‑u} \dicFlx{impers} \dicDirectTranslationCS{(co) postrádá smysl}
\dicEntry[vita] \dicTerm{vita\smash{\textsuperscript{1}}} \dicsymFrequent\  \dicIPA{{v}{\textsci}{\textlengthmark}{\textsubring{d}}{a}} \dicPos{v}[7] \dicFlx{(veit, vissi, vissum, vissi, vitað)}[10] \dicFlx{acc} \textbf{1.} \dicSynonym{þekkja} \dicDirectTranslationCS{vědět, znát} \dicExampleIS{Fólkið vissi ekkert.} \dicExampleCS{Lidé nic nevěděli.}  \textbf{2.} \dicDirectTranslationCS{dozvědět se, podívat se, zkontrolovat};  \dicPhraseIS{vita e‑ð fyrir víst} \dicDirectTranslationCS{vědět (co) jistě\,/\addthin najisto}  \textbf{3.} \dicSynonym{stefna\smash{\textsuperscript{2}}} \dicDirectTranslationCS{ležet, mířit, být obrácen\,/\addthin orientován (směrem k)} \dicExampleIS{Glugginn veit að garðinum.} \dicExampleCS{Okno je obrácené k~zahradě.};  \dicIdiom{vita}[af]{ \dicPhraseIS{vita ekki af e‑u}} \dicDirectTranslationCS{nevědět o~(čem), nebýt si vědom (čeho)}; { \dicPhraseIS{vita ekki af sér}} \dicDirectTranslationCS{nic si nepamatovat, nevědět o~sobě};  \dicIdiom{vita}[á]{ \dicPhraseIS{e‑að veit á e‑ð}} \dicDirectTranslationCS{(co) věští (co), (co) je znamením (čeho)} \dicExampleIS{Draumurinn veit á gott.} \dicExampleCS{Sen věští dobré věci.};  \dicIdiom{vita}[fyrir]{ \dicPhraseIS{vita e‑ð fyrir}} \dicDirectTranslationCS{vědět o~(čem) dopředu};  \dicIdiom{vita}[til]{ \dicPhraseIS{vita ekki til þess}} \dicDirectTranslationCS{nevědět o~tom};  \dicIdiom{vita}[um]{ \dicPhraseIS{vita um e‑n}} \dicDirectTranslationCS{vědět o~(kom)} \dicIndirectTranslationCS{(co druhý člověk dělá ap.)};  \dicIdiom{vita}{ \dicPhraseIS{að því er ég best veit}} \dicDirectTranslationCS{podle toho, co vím; podle mého nejlepšího vědomí}; { \dicPhraseIS{hver veit?}} \dicLangCat{přen.} \dicDirectTranslationCS{kdo ví?};  \dicPhraseIS{láta e‑n vita um e‑ð} \dicDirectTranslationCS{uvědomit (koho) o~(čem), informovat (koho) o~(čem)};  \dicPhraseIS{veistu hvað?} \dicDirectTranslationCS{víš co?, chceš něco říct?};  \dicPhraseIS{vita ekkert í sinn haus} \dicLangCat{přen.} \dicDirectTranslationCS{být praštěný}; { \dicPhraseIS{viti menn!}} \dicFlx{con} \dicDirectTranslationCS{ale!, jaké překvapení!};  \dicPhraseIS{það er aldrei að vita} \dicLangCat{přen.} \dicDirectTranslationCS{člověk nikdy neví};  \dicPhraseIS{það er vitað að} \dicDirectTranslationCS{je známo, že}; { \dicPhraseIS{það má guð\,/\addthin hamingjan vita}} \dicLangCat{přen.} \dicDirectTranslationCS{to ví jen Bůh}
\dicEntry[vita] \dicTerm{vita\smash{\textsuperscript{2}}} \dicIPA{{v}{\textsci}{\textlengthmark}{\textsubring{d}}{a}} \dicPos{adv} \dicDirectTranslationCS{naprosto, zcela, úplně}
\dicEntry[vita] \dicTerm{vita-\smash{\textsuperscript{3}}} \dicIPA{{v}{\textsci}{\textlengthmark}{\textsubring{d}}{a}} \dicPos{predp} \dicSynonym{alveg} \dicDirectTranslationCS{naprosto, zcela, úplně}
\dicEntry[vitað] \dicTerm{vitað} \dicIPA{{v}{\textsci}{\textlengthmark}{\textsubring{d}}{a}{\texttheta}} \dicPos{v} \dicFlx{supin} \dicLink{vita\smash{\textsuperscript{1}}}
\begin{xtolerant}{}{1pt}
\dicEntry[vitagagnslaus] \dicTerm{vita··gagns·laus} \dicIPA{{v}{\textsci}{\textlengthmark}{\textsubring{d}}{a}{\r{g}}{a}{x}{s}{l}{\oe i}{s}} \dicPos{adj}[5]\dicFlx{}[-4] \dicDirectTranslationCS{zcela\,/\addthin komplet\-ně neúčinný\,/\addthin neefektivní}
\end{xtolerant}
\dicEntry[vitandi] \dicTerm{vit··andi} \dicIPA{{v}{\textsci}{\textlengthmark}{\textsubring{d}}{a}{n}{\textsubring{d}}{\textsci}} \dicPos{adj}[13] \dicFlx{indecl}[1] \dicPhraseIS{vitandi vits} \dicFlx{adv} \dicDirectTranslationCS{vědomě, záměrně, úmyslně}
\dicEntry[vitanlega] \dicTerm{vitan··lega} \dicsymFrequent\  \dicIPA{{v}{\textsci}{\textlengthmark}{\textsubring{d}}{a}{n}{l}{\textepsilon}{\textbabygamma}{a}} \dicPos{adv} \dicSynonym{auðvitað} \dicDirectTranslationCS{samozřejmě, přirozeně, pochopitelně} \dicExampleIS{Það gerðist vitanlega ekki.} \dicExampleCS{To se pochopitelně nestalo.};  \dicPhraseIS{mér vitanlega} \dicDirectTranslationCS{co (já) vím, co je mi známo}
\dicEntry[vitanlegur] \dicTerm{vitan··legur} \dicIPA{{v}{\textsci}{\textlengthmark}{\textsubring{d}}{a}{n}{l}{\textepsilon}{\textbabygamma}{\textscy}{\textsubring{r}}} \dicPos{adj}[1]\dicFlx{}[-8] \dicSynonym{kunnur} \dicDirectTranslationCS{zřejmý, patrný, evidentní}
\dicEntry[vitaskuld] \dicTerm{vita··skuld} \dicsymFrequent\  \dicIPA{{v}{\textsci}{\textlengthmark}{\textsubring{d}}{a}{s}{\r{g}}{\textscy}{l}{\textsubring{d}}} \dicPos{adv} \dicSynonym*{að sjálfsögðu} \dicDirectTranslationCS{samo sebou, samozřejmě, ovšem, jistě, pochopitelně} \dicExampleIS{Hann var vitaskuld ánægður í leikslok.} \dicExampleCS{Byl samo sebou na konci hry spokojený.}
\dicEntry[vitavarðar] \dicTerm{vita··varðar} \dicIPA{{v}{\textsci}{\textlengthmark}{\textsubring{d}}{a}{v}{a}{r}{ð}{a}{\textsubring{r}}} \dicPos{m} \dicFlx{sg gen} \dicLink{vitavörður}
\dicEntry[vitaverðir] \dicTerm{vita··verðir} \dicIPA{{v}{\textsci}{\textlengthmark}{\textsubring{d}}{a}{v}{\textepsilon}{r}{ð}{\textsci}{\textsubring{r}}} \dicPos{m} \dicFlx{pl nom} \dicLink{vitavörður}
\dicEntry[vitavörður] \dicTerm{vita··|vörður} \dicIPA{{v}{\textsci}{\textlengthmark}{\textsubring{d}}{a}{v}{\oe}{r}{ð}{\textscy}{\textsubring{r}}} \dicPos{m}[11] \dicFlx{(‑varðar, ‑verðir)}[5] \dicDirectTranslationCS{strážce\,/\addthin strážkyně majáku}
\dicEntry[vitfirring] \dicTerm{vit··firr·ing} \dicIPA{{v}{\textsci}{\textlengthmark}{\textsubring{d}}{f}{\textsci}{r}{i}{\ng}{\r{g}}} \dicPos{f}[4] \dicFlx{(‑ar)}[7] \dicDirectTranslationCS{šílenství, šílenost, bláznovství, nepříčetnost}
\dicEntry[vitfirringur] \dicTerm{vit··firr·ing|ur} \dicIPA{{v}{\textsci}{\textlengthmark}{\textsubring{d}}{f}{\textsci}{r}{i}{\ng}{\r{g}}{\textscy}{\textsubring{r}}} \dicPos{m}[6] \dicFlx{(‑s, ‑ar)}[8] \dicDirectTranslationCS{šílenec, šílenkyně, pošetilec, pošetilka, ztřeštěnec, ztřeštěnkyně, ztřeštěnka}
\dicEntry[vitfirrtur] \dicTerm{vit··firrtur} \dicIPA{{v}{\textsci}{\textlengthmark}{\textsubring{d}}{f}{\textsci}{\textsubring{r}}{\textsubring{d}}{\textscy}{\textsubring{r}}} \dicPos{adj}[1]\dicFlx{}[-10] \dicSynonym{vitskertur} \dicDirectTranslationCS{šílený, bláznivý, nepříčetný}
\dicEntry[vitgrannur] \dicTerm{vit··|grannur} \dicIPA{{v}{\textsci}{\textlengthmark}{\textsubring{d}}{\r{g}}{r}{a}{n}{\textscy}{\textsubring{r}}} \dicPos{adj}[10] \dicFlx{(f ‑grönn, comp ‑grennri, sup ‑grennstur)}[2] \dicSynonym{heimskur} \dicDirectTranslationCS{(jsoucí) mdlého rozumu, zabedněný, natvrdlý}
\dicEntry[viti] \dicTerm{vit|i} \dicIPA{{v}{\textsci}{\textlengthmark}{\textsubring{d}}{\textsci}} \dicPos{m}[1] \dicFlx{(‑a, ‑ar)}[1] \dicDirectTranslationCS{maják}
\dicEntry[vitiborinn] \dicTerm{viti··borinn} \dicIPA{{v}{\textsci}{\textlengthmark}{\textsubring{d}}{\textsci}{\textsubring{b}}{\textopeno}{r}{\textsci}{\textsubring{n}}} \dicPos{adj}[6]\dicFlx{}[-6] \dicSynonym{hygginn} \dicDirectTranslationCS{inteligentní, moudrý}
\dicEntry[vitja] \dicTerm{vitj|a} \dicIPA{{v}{\textsci}{\textlengthmark}{\textsubring{d}}{j}{a}} \dicPos{v}[1] \dicFlx{(‑aði)}[1] \dicFlx{gen} \dicDirectTranslationCS{navštívit, stavit se, zkontrolovat} \dicExampleIS{vitja sjúklings} \dicExampleCS{navštívit nemocného};  \dicPhraseIS{vitja um e‑ð} \dicDirectTranslationCS{(z)kontrolovat (co), dohlídnout\,/\addthin dohlížet na (co)} \dicExampleIS{vitja um hestana} \dicExampleCS{zkontrolovat koně}
\dicEntry[vitjun] \dicTerm{vitj|un} \dicIPA{{v}{\textsci}{\textlengthmark}{\textsubring{d}}{j}{\textscy}{\textsubring{n}}} \dicPos{f}[7] \dicFlx{(‑unar, ‑anir)}[8] \dicDirectTranslationCS{návštěva, navštívení, zkontrolování} \dicIndirectTranslationCS{(zejména o~domácí návštěvě doktora u~pacienta)} \dicExampleIS{sjúkravitjun} \dicExampleCS{návštěva nemocného}
\dicEntry[vitlaus] \dicTerm{vit··laus} \dicsymFrequent\  \dicIPA{{v}{\textsci}{h}{\textsubring{d}}{l}{\oe i}{s}} \dicPos{adj}[5]\dicFlx{}[-1] \textbf{1.} \dicSynonym{æstur} \dicDirectTranslationCS{zuřivý, vzteklý};  \dicPhraseIS{gera allt vitlaust} \dicLangCat{přen.} \dicDirectTranslationCS{vyvolat rozruch, dostat vše do varu};  \dicPhraseIS{það verður allt vitlaust} \dicLangCat{přen.} \dicDirectTranslationCS{vypukne (hotové) peklo}  \textbf{2.} \dicSynonym{heimskur} \dicDirectTranslationCS{hloupý, blbý, stupidní}  \textbf{3.} \dicSynonym{skakkur} \dicDirectTranslationCS{chybný, nesprávný} \dicExampleIS{vitlaust númer} \dicExampleCS{omyl (chybné telefonní číslo)}  \textbf{4.} \dicSynonym{fíkinn} \dicDirectTranslationCS{zblázněný, šílený};  \dicPhraseIS{vera vitlaus í e‑ð} \dicDirectTranslationCS{být do (čeho) blázen}  \textbf{5.} \dicDirectTranslationCS{zuřivý, prudký (bouře ap.)};  \dicIdiom{vitlaus}{ \dicPhraseIS{vitlausa beinið}} \dicDirectTranslationCS{brňavka}
\dicEntry[vitlausrahæli] \dicTerm{vit·lausra··hæli} \dicIPA{{v}{\textsci}{h}{\textsubring{d}}{l}{\oe i}{s}{r}{a}{h}{a}{i}{l}{\textsci}} \dicPos{n}[2] \dicFlx{(‑s, ‑)}[14] \dicLangCat{han.} \dicDirectTranslationCS{cvokárna, cvokhaus, blázinec}
\dicEntry[vitlegur] \dicTerm{vit··legur} \dicIPA{{v}{\textsci}{\textlengthmark}{\textsubring{d}}{l}{\textepsilon}{\textbabygamma}{\textscy}{\textsubring{r}}} \dicPos{adj}[1]\dicFlx{}[-8] \dicSynonym{skynsamlegur} \dicDirectTranslationCS{moudrý, soudný}
\dicEntry[vitleysa] \dicTerm{vit··leys|a} \dicsymFrequent\  \dicIPA{{v}{\textsci}{h}{\textsubring{d}}{l}{ei}{s}{a}} \dicPos{f}[1] \dicFlx{(‑u, ‑ur)}[7] \textbf{1.} \dicSynonym{villa\smash{\textsuperscript{1}}} \dicDirectTranslationCS{chyba, mýlka}  \textbf{2.} \dicSynonym{rugl} \dicDirectTranslationCS{nesmysl, hloupost};  \dicPhraseIS{hvaða vitleysa!} \dicDirectTranslationCS{nesmysl!}
\dicEntry[vitleysingur] \dicTerm{vit··leys·ing|ur} \dicIPA{{v}{\textsci}{h}{\textsubring{d}}{l}{ei}{s}{i}{\ng}{\r{g}}{\textscy}{\textsubring{r}}} \dicPos{m}[6] \dicFlx{(‑s, ‑ar)}[8] \textbf{1.} \dicSynonym{heimskingi} \dicDirectTranslationCS{hlupák, hlupačka, blbec, idiot(ka)}  \textbf{2.} \dicSynonym{vitfirringur} \dicDirectTranslationCS{šílenec, šílenkyně, ztřeštěnec, ztřeštěnkyně, ztřeštěnka}
\dicEntry[vitmaður] \dicTerm{vit··|maður} \dicIPA{{v}{\textsci}{\textlengthmark}{\textsubring{d}}{m}{a}{ð}{\textscy}{\textsubring{r}}} \dicPos{m}[13] \dicFlx{(‑manns, ‑menn)}[2] \textbf{1.} \dicSynonym*{vitur maður} \dicDirectTranslationCS{moudrý člověk}  \textbf{2.} \dicFieldCat{hist.\,/\addthin biol.} \dicDirectTranslationCS{člověk moudrý} \textit{(l.~{\textLA{Homo sapiens}})}
\dicEntry[vitna] \dicTerm{vitn|a} \dicsymFrequent\  \dicIPA{{v}{\textsci}{h}{\textsubring{d}}{n}{a}} \dicPos{v}[1] \dicFlx{(‑aði)}[22] \dicSynonym*{bera vitni} \dicDirectTranslationCS{(do)svědčit, vypovídat (u~soudu ap.)};  \dicPhraseIS{vitna fyrir rétti} \dicDirectTranslationCS{svědčit před soudem};  \dicIdiom{vitna}[gegn]{ \dicPhraseIS{vitna gegn e‑m}} \dicDirectTranslationCS{svědčit proti (komu)};  \dicIdiom{vitna}[í]{ \dicPhraseIS{vitna í e‑n\,/\addthin e‑ð}} \dicDirectTranslationCS{citovat (koho\,/\addthin co)};  \dicIdiom{vitna}[til]{ \dicPhraseIS{vitna til e‑s}} \dicDirectTranslationCS{odvolat\,/\addthin odvolávat se na (co)};  \dicIdiom{vitna}[um]{ \dicPhraseIS{vitna um e‑ð}} {\textbf{a.}} \dicDirectTranslationCS{dosvědčit (co), vypovídat o~(čem) (u~soudu ap.)};  {\textbf{b.}} \dicDirectTranslationCS{svědčit o~(čem), ukazovat na (co)};  \dicIdiom{vitnast}{ \dicPhraseIS{vitnast}} \dicFlx{refl} \dicDirectTranslationCS{rozšířit se, dostat se na veřejnost (zpráva ap.)}
\dicEntry[vitnastúka] \dicTerm{vitna··stúk|a} \dicIPA{{v}{\textsci}{h}{\textsubring{d}}{n}{a}{s}{\textsubring{d}}{u}{\r{g}}{a}} \dicPos{f}[1] \dicFlx{(‑u, ‑ur)}[19] \dicDirectTranslationCS{lavice svědků}
\dicEntry[vitnavernd] \dicTerm{vitna··vernd} \dicIPA{{v}{\textsci}{h}{\textsubring{d}}{n}{a}{v}{\textepsilon}{r}{n}{\textsubring{d}}} \dicPos{f}[4] \dicFlx{(‑ar)}[3] \dicFieldCat{práv.} \dicDirectTranslationCS{ochrana svědka}
\dicEntry[vitneskja] \dicTerm{vit··neskj|a} \dicsymFrequent\  \dicIPA{{v}{\textsci}{h}{\textsubring{d}}{n}{\textepsilon}{s}{\r{\textObardotlessj}}{a}} \dicPos{f}[1] \dicFlx{(‑u)}[5] \dicSynonym{upplýsing} \dicDirectTranslationCS{informace, údaje, sdělení} \dicExampleIS{fá vitneskju} \dicExampleCS{dostat informace}
\dicEntry[vitni] \dicTerm{vitni} \dicsymFrequent\  \dicIPA{{v}{\textsci}{h}{\textsubring{d}}{n}{\textsci}} \dicPos{n}[2] \dicFlx{(‑s, ‑)}[14] \textbf{1.} \dicSynonym{vitnisburður} \dicDirectTranslationCS{svědectví, výpověď};  \dicPhraseIS{bera vitni} \dicDirectTranslationCS{svědčit, vypovídat (u~soudu ap.)}  \textbf{2.} \dicSynonym*{maður sem vitnar} \dicDirectTranslationCS{svědek, svědkyně};  \dicPhraseIS{leiða fram vitni} \dicDirectTranslationCS{předvolat svědka};  \dicPhraseIS{vera vitni að e‑u} \dicDirectTranslationCS{být svědkem (čeho)};  \dicIdiom{vitni}{ \dicPhraseIS{eins og raun ber vitni}} \dicFlx{adv} \dicDirectTranslationCS{jak zkušenost ukazuje};  \dicPhraseIS{það þarf ekki frekar vitnanna við} \dicLangCat{přen.} \dicDirectTranslationCS{další důkazy již nejsou potřeba}
\dicEntry[vitnisburður] \dicTerm{vitnis··burð|ur} \dicIPA{{v}{\textsci}{h}{\textsubring{d}}{n}{\textsci}{s}{\textsubring{b}}{\textscy}{r}{ð}{\textscy}{\textsubring{r}}} \dicPos{m}[10] \dicFlx{(‑ar, ‑ir)}[4] \textbf{1.} \dicSynonym*{framburður vitnis} \dicDirectTranslationCS{(svědecká) výpověď, svědectví}  \textbf{2.} \dicSynonym{umsögn} \dicDirectTranslationCS{svědectví, prohlášení, (závažná) zpráva}  \textbf{3.} \dicFieldCat{škol.} \dicSynonym*{námseinkunn} \dicDirectTranslationCS{osvědčení, reference}
\dicEntry[vitrast] \dicTerm{vitr|ast} \dicIPA{{v}{\textsci}{\textlengthmark}{\textsubring{d}}{r}{a}{s}{\textsubring{d}}} \dicPos{v}[1] \dicFlx{(‑aðist)}[99] \dicFlx{refl impers} \dicPhraseIS{e‑m vitrast e‑að} \dicDirectTranslationCS{(komu) se zjeví (co), (komu) se zjevuje (co)}
\dicEntry[vitringur] \textls[10]{\dicTerm{vitr··ing|ur} \dicIPA{{v}{\textsci}{\textlengthmark}{\textsubring{d}}{r}{i}{\ng}{\r{g}}{\textscy}{\textsubring{r}}} \dicPos{m}[6] \dicFlx{(‑s, ‑ar)}[8]} \textbf{1.} \dicSynonym{spekingur} \dicDirectTranslationCS{mudrc};  \dicPhraseIS{vitringarnir þrír} \dicFieldCat{náb.} \dicDirectTranslationCS{tři králové}  \textbf{2.} \dicSynonym*{gáfaður maður} \dicDirectTranslationCS{génius} \dicIndirectTranslationCS{(ironicky)}
\dicEntry[vitrun] \dicTerm{vitr|un} \dicIPA{{v}{\textsci}{\textlengthmark}{\textsubring{d}}{r}{\textscy}{\textsubring{n}}} \dicPos{f}[7] \dicFlx{(‑unar, ‑anir)}[8] \dicSynonym{opinberun} \dicDirectTranslationCS{zjevení, vize} \dicExampleIS{fá vitrun} \dicExampleCS{dostat zjevení, mít vizi}
\dicEntry[vitrænn] \dicTerm{vit··rænn} \dicIPA{{v}{\textsci}{\textlengthmark}{\textsubring{d}}{r}{a}{i}{\textsubring{d}}{\textsubring{n}}} \dicPos{adj}[7]\dicFlx{}[-1] \dicDirectTranslationCS{rozumový, racionální, myšlenkový}
\dicEntry[vitskertur] \dicTerm{vit··skertur} \dicIPA{{v}{\textsci}{\textlengthmark}{\textsubring{d}}{s}{\r{\textObardotlessj}}{\textepsilon}{\textsubring{r}}{\textsubring{d}}{\textscy}{\textsubring{r}}} \dicPos{adj}[1]\dicFlx{}[-10] \dicSynonym{brjálaður} \dicDirectTranslationCS{ztřeštěný, šílený, nepříčetný}
\dicEntry[vitsmunalíf] \dicTerm{vits·muna··líf} \dicIPA{{v}{\textsci}{\textlengthmark}{\textsubring{d}}{s}{m}{\textscy}{n}{a}{l}{i}{f}} \dicPos{n}[2] \dicFlx{(‑s)}[2] \dicDirectTranslationCS{rozumný život}
\dicEntry[vitsmunir] \dicTerm{vits··munir} \dicIPA{{v}{\textsci}{\textlengthmark}{\textsubring{d}}{s}{m}{\textscy}{n}{\textsci}{\textsubring{r}}} \dicPos{m}[10] \dicFlx{pl}[1] \dicSynonym{vit} \dicDirectTranslationCS{rozum, intelekt, důvtip}
\dicEntry[vitstola] \dicTerm{vit··stola} \dicIPA{{v}{\textsci}{\textlengthmark}{\textsubring{d}}{s}{\textsubring{d}}{\textopeno}{l}{a}} \dicPos{adj}[13] \dicFlx{indecl}[1] \dicSynonym{vitfirrtur} \dicDirectTranslationCS{šílený, ztřeštěný} \dicExampleIS{verða vitstola af hræðslu} \dicExampleCS{šílet strachy}
\dicEntry[vitund] \dicTerm{vit··und} \dicsymFrequent\  \dicIPA{{v}{\textsci}{\textlengthmark}{\textsubring{d}}{\textscy}{n}{\textsubring{d}}} \dicPos{f}[7] \dicFlx{(‑ar)}[3] \textbf{1.} \dicSynonym{meðvitund} \dicDirectTranslationCS{(po)vědomí, uvědomování si};  \dicPhraseIS{vakna til vitundar um e‑ð} \dicDirectTranslationCS{uvědomit si (co)}  \textbf{2.} \dicSynonym{vitneskja} \dicDirectTranslationCS{(po)vědomí, vědění};  \dicPhraseIS{án e‑rs vitundar} \dicDirectTranslationCS{bez vědomí (koho)} \dicExampleIS{Þetta gerðist án hennar vitundar.} \dicExampleCS{Stalo se to bez jejího vědomí.};  \dicPhraseIS{með vitund og vilja e‑rs} \dicDirectTranslationCS{se souhlasem (koho)};  \dicIdiom{vitund}{ \dicPhraseIS{ekki vitund}} \dicFlx{adv} \dicDirectTranslationCS{ani v~nejmenším, ani trochu}
\dicEntry[vitundarlaus] \dicTerm{vitundar··laus} \dicIPA{{v}{\textsci}{\textlengthmark}{\textsubring{d}}{\textscy}{n}{\textsubring{d}}{a}{r}{l}{\oe i}{s}} \dicPos{adj}[5]\dicFlx{}[-4] \dicSynonym{meðvitundarlaus} \dicDirectTranslationCS{bezvědomý}
\dicEntry[vitundarstreymi] \dicTerm{vitundar··streymi} \dicIPA{{v}{\textsci}{\textlengthmark}{\textsubring{d}}{\textscy}{n}{\textsubring{d}}{a}{\textsubring{r}}{s}{\textsubring{d}}{r}{ei}{m}{\textsci}} \dicPos{n}[2] \dicFlx{(‑s)}[20] \dicFieldCat{lit.} \dicDirectTranslationCS{proud vědomí}
\dicEntry[vitundarvottur] \dicTerm{vitundar··vott|ur} \dicIPA{{v}{\textsci}{\textlengthmark}{\textsubring{d}}{\textscy}{n}{\textsubring{d}}{a}{r}{v}{\textopeno}{h}{\textsubring{d}}{\textscy}{\textsubring{r}}} \dicPos{m}[6] \dicFlx{(‑s, ‑ar)}[4] \dicFieldCat{práv.} \dicDirectTranslationCS{očitý svědek, očitá svědkyně}
\dicEntry[vitur] \dicTerm{vitur} \dicIPA{{v}{\textsci}{\textlengthmark}{\textsubring{d}}{\textscy}{\textsubring{r}}} \dicPos{adj}[9] \dicFlx{(f ‑)}[1] \dicSynonym{vís} \dicDirectTranslationCS{moudrý, rozumný}
\dicEntry[viturlega] \dicTerm{vitur··lega} \dicIPA{{v}{\textsci}{\textlengthmark}{\textsubring{d}}{\textscy}{r}{l}{\textepsilon}{\textbabygamma}{a}} \dicPos{adv} \dicDirectTranslationCS{moudře, inteligentně}
\dicEntry[viturlegur] \dicTerm{vitur··legur} \dicIPA{{v}{\textsci}{\textlengthmark}{\textsubring{d}}{\textscy}{r}{l}{\textepsilon}{\textbabygamma}{\textscy}{\textsubring{r}}} \dicPos{adj}[1]\dicFlx{}[-8] \dicDirectTranslationCS{moudrý, inteligentní} \dicExampleIS{viturleg tillaga} \dicExampleCS{moudrý návrh}
\dicEntry[viturleiki] \dicTerm{vitur··leik|i}\dicTerm{, viturleikur} \dicIPA{{v}\-{\textsci}\-{\textlengthmark}\-{\textsubring{d}}\-{\textscy}\-{r}\-{l}\-{ei}\-{\r{\textObardotlessj}}\-{\textsci}\-} \dicPos{m}[1] \dicFlx{(‑a)}[3] \dicDirectTranslationCS{moudrost, rozumnost, rozvážnost}
\dicEntry[viturleikur] \dicTerm{vitur··leik|ur} \dicIPA{{v}{\textsci}{\textlengthmark}{\textsubring{d}}{\textscy}{r}{l}{ei}{\r{g}}{\textscy}{\textsubring{r}}} \dicPos{m}[6] \dicFlx{(‑s)}[17] \dicLink{viturleiki}
\dicEntry[víbrafónn] \dicTerm{víbra··fón|n} \dicIPA{{v}{i}{\textlengthmark}{\textsubring{b}}{r}{a}{f}{ou}{\textsubring{d}}{\textsubring{n}}} \dicPos{m}[6] \dicFlx{(‑s, ‑ar)}[42] \dicFieldCat{hud.} \dicDirectTranslationCS{vibrafon}
\dicEntry[vídd] \dicTerm{vídd} \dicIPA{{v}{i}{\textsubring{d}}{\textlengthmark}} \dicPos{f}[7] \dicFlx{(‑ar, ‑ir)}[1] \textbf{1.} \dicDirectTranslationCS{šířka} \dicExampleIS{vídd kjóls} \dicExampleCS{šířka šatů}  \textbf{2.} \dicSynonym{víðátta} \dicDirectTranslationCS{šíře, rozloha, rozlehlost} \dicExampleIS{vídd öræfanna} \dicExampleCS{rozlehlost pouště}  \textbf{3.} \dicFieldCat{mat.} \dicDirectTranslationCS{rozměr, dimenze}  \textbf{4.} \dicLangCat{přen.} \dicDirectTranslationCS{rozměr, dimenze (života ap.)} \dicExampleIS{nýjar víddir} \dicExampleCS{nové rozměry}
\dicEntry[vídeó] \dicTerm{vídeó} \dicIPA{{v}{i}{\textlengthmark}{\textsubring{d}}{\textepsilon}{ou}} \dicPos{n}[2] \dicFlx{(‑s, ‑)}[30] \dicSynonym{myndband} \dicDirectTranslationCS{video} \dicIndirectTranslationCS{(snímek)}
\dicEntry[vídeótæki] \dicTerm{vídeó··tæki} \dicIPA{{v}{i}{\textlengthmark}{\textsubring{d}}{\textepsilon}{ou}{t\smash{\textsuperscript{h}}}{a}{i}{\r{\textObardotlessj}}{\textsci}} \dicPos{n}[2] \dicFlx{(‑s, ‑)}[16] \dicDirectTranslationCS{video, videorekordér}
\dicEntry[víð] \dicTerm{víð} \dicIPA{{v}{i}{\textlengthmark}{\texttheta}} \dicPos{f}[4] \dicFlx{(‑ar)}[3] \dicPhraseIS{á víð og dreif} \dicFlx{adv} \dicDirectTranslationCS{všude (možně)}
\dicEntry[víða] \dicTerm{víð|a} \dicsymFrequent\  \dicIPA{{v}{i}{\textlengthmark}{ð}{a}} \dicPos{adv} \dicFlx{(comp ‑ar, sup ‑ast)} \dicDirectTranslationCS{široko daleko, všude možně} \dicExampleIS{Ég hef víða verið.} \dicExampleCS{Byl jsem všude možně.};  \dicPhraseIS{víðar} \dicFlx{comp} \dicDirectTranslationCS{jinde, na jiných místech}
\dicEntry[víðar] \dicTerm{víðar} \dicIPA{{v}{i}{\textlengthmark}{ð}{a}{\textsubring{r}}} \dicPos{adv} \dicFlx{comp} \textbf{1.} \dicLink{vítt}  \textbf{2.} \dicLink{víða}
\dicEntry[víðara] \dicTerm{víðara} \dicIPA{{v}{i}{\textlengthmark}{ð}{a}{r}{a}} \dicPos{adv} \dicFlx{comp} \dicLink{vítt}
\dicEntry[víðast] \dicTerm{víðast} \dicIPA{{v}{i}{\textlengthmark}{ð}{a}{s}{\textsubring{d}}} \dicPos{adv} \dicFlx{sup} \textbf{1.} \dicLink{vítt}  \textbf{2.} \dicLink{víða}
\begin{xtolerant}{}{1pt}
\dicEntry[víðavangshlaup] \dicTerm{víða·vangs··hlaup} \dicIPA{{v}{i}{\textlengthmark}{ð}{a}{v}{au}{\ng}{s}{\textsubring{l}}{\oe i}{\textsubring{b}}} \dicPos{n}[2] \dicFlx{(‑s, ‑)}[5] \dicFieldCat{sport.} \dicDirectTranslationCS{přes\-polní\,/\addthin terénní běh}
\end{xtolerant}
\dicEntry[víðavangur] \dicTerm{víða··vang|ur} \dicIPA{{v}{i}{\textlengthmark}{ð}{a}{v}{au}{\ng}{\r{g}}{\textscy}{\textsubring{r}}} \dicPos{m}[6] \dicFlx{(‑s, ‑ar)}[25] \dicPhraseIS{(úti) á víðavangi} \dicFlx{adv} \dicDirectTranslationCS{(venku) v~přírodě\,/\addthin divočině (kempovat ap.)}
\dicEntry[víðátta] \dicTerm{víð··átt|a} \dicIPA{{v}{i}{\textlengthmark}{ð}{au}{h}{\textsubring{d}}{a}} \dicPos{f}[1] \dicFlx{(‑u, ‑ur)}[7] \dicSynonym{vídd} \dicDirectTranslationCS{šíře, rozlehlost, prostor} \dicExampleIS{endalaus víðátta} \dicExampleCS{nekonečný prostor}
\dicEntry[víðáttumikill] \dicTerm{víð·áttu··|mikill} \dicIPA{{v}{i}{\textlengthmark}{ð}{au}{h}{\textsubring{d}}{\textscy}{m}{\textsci}{\r{\textObardotlessj}}{\textsci}{\textsubring{d}}{\textsubring{l}}} \dicPos{adj}[11] \dicFlx{(comp ‑meiri, sup ‑mestur)}[12] \dicDirectTranslationCS{rozlehlý, prostorný, širý} \dicExampleIS{víðáttumikil slétta} \dicExampleCS{rozlehlá planina}
\dicEntry[víðfeðmur] \dicTerm{víð··feðmur} \dicIPA{{v}{i}{ð}{f}{\textepsilon}{ð}{m}{\textscy}{\textsubring{r}}} \dicPos{adj}[1]\dicFlx{}[-1] \dicSynonym{víðtækur} \dicDirectTranslationCS{rozsáhlý, rozlehlý, prostorný}
\dicEntry[víðfrægur] \dicTerm{víð··frægur} \dicIPA{{v}{i}{ð}{f}{r}{a}{i}{\textbabygamma}{\textscy}{\textsubring{r}}} \dicPos{adj}[1]\dicFlx{}[-1] \dicDirectTranslationCS{proslulý, proslavený, vě\-hlasný}
\dicEntry[víðförull] \dicTerm{víð··förull} \dicIPA{{v}{i}{ð}{f}{\oe}{r}{\textscy}{\textsubring{d}}{\textsubring{l}}} \dicPos{adj}[8]\dicFlx{}[-5] \dicDirectTranslationCS{scestovalý}
\dicEntry[víðir] \dicTerm{víð|ir} \dicIPA{{v}{i}{\textlengthmark}{ð}{\textsci}{\textsubring{r}}} \dicPos{m}[7] \dicFlx{(‑is)}[2] \dicFieldCat{bot.} \dicDirectTranslationCS{vrba} \textit{(l.~{\textLA{Salix}})}  \dicsymPhoto\ 
\dicFigure{ds_image_vidir_0_1.jpg}{Víðir}{Víðir - Demeester, CC BY-SA 3.0}
\dicEntry[víðkunnur] \dicTerm{víð··kunnur} \dicIPA{{v}{i}{ð}{k\smash{\textsuperscript{h}}}{\textscy}{n}{\textscy}{\textsubring{r}}} \dicPos{adj}[1]\dicFlx{}[-1] \dicSynonym{frægur} \dicDirectTranslationCS{známý, proslulý}
\dicEntry[víðlendur] \dicTerm{víð··lendur} \dicIPA{{v}{i}{ð}{l}{\textepsilon}{n}{\textsubring{d}}{\textscy}{\textsubring{r}}} \dicPos{adj}[2]\dicFlx{}[-14] \dicSynonym{víðáttumikill} \dicDirectTranslationCS{rozsáhlý, rozlehlý, širý} \dicExampleIS{víðlent ríki} \dicExampleCS{rozlehlá země}
\dicEntry[víðlesinn] \dicTerm{víð··lesinn} \dicIPA{{v}{i}{ð}{l}{\textepsilon}{s}{\textsci}{\textsubring{n}}} \dicPos{adj}[6]\dicFlx{}[-2] \textbf{1.} \dicDirectTranslationCS{čtený (širokou veřejností ap.)} \dicExampleIS{víðlesnasta vikublaðið} \dicExampleCS{nejčtenější týdeník}  \textbf{2.} \dicDirectTranslationCS{sečtělý}
\dicEntry[víðnet] \dicTerm{víð··net} \dicIPA{{v}{i}{ð}{n}{\textepsilon}{\textsubring{d}}} \dicPos{n}[2] \dicFlx{(‑s, ‑)}[5] \dicFieldCat{poč.} \dicDirectTranslationCS{WAN, rozsáhlá síť} \dicIndirectTranslationCS{(Wide Area Network)}
\dicEntry[víðsýni] \dicTerm{víð··sýni} \dicIPA{{v}{i}{ð}{s}{i}{n}{\textsci}} \dicPos{n}[2] \dicFlx{(‑s)}[20] \dicDirectTranslationCS{tolerance, (myšlenková) otevřenost, nepředpojatost}
\dicEntry[víðsýnn] \dicTerm{víð··sýnn} \dicIPA{{v}{i}{ð}{s}{i}{\textsubring{d}}{\textsubring{n}}} \dicPos{adj}[7]\dicFlx{}[-1] \textbf{1.} \dicDirectTranslationCS{tolerantní, (myšlenkově) otevřený, nepředpojatý, (jsoucí) bez předsudků} \dicExampleIS{vera víðsýnn í trúarefnum} \dicExampleCS{být otevřený v~otázkách víry}  \textbf{2.} \dicDirectTranslationCS{(jsoucí) s~dobrým výhledem}
\dicEntry[víðtækur] \dicTerm{víð··tækur} \dicIPA{{v}{i}{ð}{t\smash{\textsuperscript{h}}}{a}{i}{\r{g}}{\textscy}{\textsubring{r}}} \dicPos{adj}[1]\dicFlx{}[-1] \dicSynonym{yfirgripsmikill} \dicDirectTranslationCS{rozsáhlý, obsáhlý} \dicExampleIS{víðtæk þekking} \dicExampleCS{rozsáhlá znalost}
\dicEntry[víður] \dicTerm{víður} \dicsymFrequent\  \dicIPA{{v}{i}{\textlengthmark}{ð}{\textscy}{\textsubring{r}}} \dicPos{adj}[2]\dicFlx{}[-6] \textbf{1.} \dicSynonym{breiður} \dicDirectTranslationCS{široký} \dicExampleIS{vítt op} \dicExampleCS{široká mezera}  \textbf{2.} \dicSynonym{rúmur} \dicDirectTranslationCS{široký, volný (oděv ap.)} \dicExampleIS{víðir skór} \dicExampleCS{volné boty}
\dicEntry[Víetnam] \dicTerm{Víetnam} \dicIPA{{v}{i}{j}{\textepsilon}{\textsubring{d}}{n}{a}{\textsubring{m}}} \dicPos{n}[2] \dicFlx{(‑s)}[32] \dicFieldCat{geog.} \dicDirectTranslationCS{Vietnam}
\dicEntry[Víetnami] \dicTerm{Víetnam|i} \dicIPA{{v}{i}{j}{\textepsilon}{\textsubring{d}}{n}{a}{m}{\textsci}} \dicPos{m}[1] \dicFlx{(‑a, ‑ar)}[8] \dicDirectTranslationCS{Vietnamec, Vietnamka}
\dicEntry[víetnamska] \dicTerm{víetnamska} \dicIPA{{v}{i}{j}{\textepsilon}{\textsubring{d}}{n}{a}{m}{s}{\r{g}}{a}} \dicPos{f}[1] \dicFlx{(víetnömsku)}[2] \dicDirectTranslationCS{vietnamština}
\dicEntry[víetnamskur] \dicTerm{víetnamskur} \dicIPA{{v}{i}{j}{\textepsilon}{\textsubring{d}}{n}{a}{m}{s}{\r{g}}{\textscy}{\textsubring{r}}} \dicPos{adj}[1] \dicFlx{(f víetnömsk)}[3] \dicDirectTranslationCS{vietnamský}
\dicEntry[víf] \dicTerm{víf} \dicIPA{{v}{i}{\textlengthmark}{f}} \dicPos{n}[2] \dicFlx{(‑s, ‑)}[5] \dicLangCat{básn.} \dicSynonym{kona} \dicDirectTranslationCS{žena, manželka}
\dicEntry[víg] \dicTerm{víg} \dicIPA{{v}{i}{\textlengthmark}{x}} \dicPos{n}[2] \dicFlx{(‑s, ‑)}[5] \textbf{1.} \dicSynonym{dráp} \dicDirectTranslationCS{zabití}  \textbf{2.} \dicLangCat{zast.} \dicSynonym{bardagi} \dicDirectTranslationCS{boj, bitva};  \dicPhraseIS{standa vel\,/\addthin illa að vígi} \dicDirectTranslationCS{být dobře\,/\addthin špatně připraven}
\dicEntry[vígahnöttur] \dicTerm{víga··|hnöttur} \dicIPA{{v}{i}{\textlengthmark}{\textbabygamma}{a}{\textsubring{n}}{\oe}{h}{\textsubring{d}}{\textscy}{\textsubring{r}}} \dicPos{m}[11] \dicFlx{(‑hnattar, ‑hnettir)}[5] \dicDirectTranslationCS{meteor, povětroň, létavice}
\dicEntry[vígahugur] \dicTerm{víga··hug|ur} \dicIPA{{v}{i}{\textlengthmark}{\textbabygamma}{a}{h}{\textscy}{\textbabygamma}{\textscy}{\textsubring{r}}} \dicPos{m}[10] \dicFlx{(‑ar)}[12] \dicDirectTranslationCS{bojovná nálada, bojovnost, výbojnost};  \dicPhraseIS{vera í vígahug} \dicDirectTranslationCS{být v~bojovné náladě}
\dicEntry[vígalegur] \dicTerm{víga··legur} \dicIPA{{v}{i}{\textlengthmark}{\textbabygamma}{a}{l}{\textepsilon}{\textbabygamma}{\textscy}{\textsubring{r}}} \dicPos{adj}[1]\dicFlx{}[-8] \dicDirectTranslationCS{bojovný, bojechtivý}
\dicEntry[vígamaður] \dicTerm{víga··|maður} \dicIPA{{v}{i}{\textlengthmark}{\textbabygamma}{a}{m}{a}{ð}{\textscy}{\textsubring{r}}} \dicPos{m}[13] \dicFlx{(‑manns, ‑menn)}[2] \dicDirectTranslationCS{zabiják, vražedný útočník\,/\addthin útočnice}
\dicEntry[vígasveit] \dicTerm{víga··sveit} \dicIPA{{v}{i}{\textlengthmark}{\textbabygamma}{a}{s}{v}{ei}{\textsubring{d}}} \dicPos{f}[7] \dicFlx{(‑ar, ‑ir)}[1] \dicDirectTranslationCS{guerilla, gerila, guerillový oddíl}
\dicEntry[vígbúnaður] \dicTerm{víg··bú·nað|ur} \dicIPA{{v}{i}{\textbabygamma}{\textsubring{b}}{u}{n}{a}{ð}{\textscy}{\textsubring{r}}} \dicPos{m}[10] \dicFlx{(‑ar)}[9] \dicDirectTranslationCS{zbrojení, ozbrojování} \dicExampleIS{vígbúnaðarkapphlaup stórveldanna} \dicExampleCS{závod světových velmocí ve zbrojení}
\dicEntry[víggirða] \dicTerm{víg··gir|ða} \dicIPA{{v}{i}{\textbabygamma}{\r{\textObardotlessj}}{\textsci}{r}{ð}{a}} \dicPos{v}[2] \dicFlx{(‑ti, ‑t)}[50] \dicFlx{acc} \dicDirectTranslationCS{opevnit, opevňovat, fortifikovat} \dicExampleIS{Borg var vel víggirt.} \dicExampleCS{Město bylo dobře opevněno.}
\dicEntry[víggirðing] \dicTerm{víg··girð·ing} \dicIPA{{v}{i}{\textbabygamma}{\r{\textObardotlessj}}{\textsci}{r}{ð}{i}{\ng}{\r{g}}} \dicPos{f}[4] \dicFlx{(‑ar, ‑ar)}[5] \dicDirectTranslationCS{opevnění, hradba, fortifikace}
\dicEntry[víggirtur] \dicTerm{víg··girtur} \dicIPA{{v}{i}{\textbabygamma}{\r{\textObardotlessj}}{\textsci}{\textsubring{r}}{\textsubring{d}}{\textscy}{\textsubring{r}}} \dicPos{adj}[1]\dicFlx{}[-13] \dicDirectTranslationCS{opevněný, fortifikovaný}
\dicEntry[vígi] \dicTerm{vígi} \dicIPA{{v}{i}{j}{\textlengthmark}{\textsci}} \dicPos{n}[2] \dicFlx{(‑s, ‑)}[16] \textbf{1.} \dicSynonym{virki\smash{\textsuperscript{2}}} \dicDirectTranslationCS{pevnost, opevnění} \dicExampleIS{óvinnandi vígi} \dicExampleCS{nedobytná pevnost}  \textbf{2.} \dicSynonym*{miðstöð hugmynda} \dicDirectTranslationCS{centrum, bašta (názorového směru ap.)}
\dicEntry[vígja] \dicTerm{víg|ja} \dicIPA{{v}{i}{j}{\textlengthmark}{a}} \dicPos{v}[2] \dicFlx{(‑ði, ‑t)}[89] \dicFlx{acc} \textbf{1.} \dicFieldCat{náb.} \dicDirectTranslationCS{vysvětit, konsekrovat} \dicIndirectTranslationCS{(církevním obřadem učinit oprávněným ke kněžskému úřadu)} \dicExampleIS{vígja prest} \dicExampleCS{vysvětit kněze}  \textbf{2.} \dicDirectTranslationCS{slavnostně otevřít, inaugurovat} \dicExampleIS{vígja brúna} \dicExampleCS{slavnostně otevřít most}
\dicEntry[víglína] \dicTerm{víg··lín|a} \dicIPA{{v}{i}{\textbabygamma}{l}{i}{n}{a}} \dicPos{f}[1] \dicFlx{(‑u, ‑ur)}[7] \dicFieldCat{voj.} \dicDirectTranslationCS{fronta, (bojová) linie}
\dicEntry[vígorð] \dicTerm{víg··orð} \dicIPA{{v}{i}{\textlengthmark}{\textbabygamma}{\textopeno}{r}{\texttheta}} \dicPos{n}[2] \dicFlx{(‑s, ‑)}[5] \dicSynonym{slagorð} \dicDirectTranslationCS{heslo, slogan}
\dicEntry[vígsla] \dicTerm{vígsl|a} \dicIPA{{v}{i}{x}{s}{\textsubring{d}}{l}{a}} \dicPos{f}[1] \dicFlx{(‑u, ‑ur)}[13] \textbf{1.} \dicFieldCat{náb.} \dicDirectTranslationCS{vysvěcení (kněze ap.)}  \textbf{2.} \dicSynonym{opnun} \dicDirectTranslationCS{slavnostní otevření, inaugurace}
\dicEntry[vígtönn] \dicTerm{víg··|tönn} \dicIPA{{v}{i}{x}{t\smash{\textsuperscript{h}}}{\oe}{\textsubring{n}}} \dicPos{f}[8] \dicFlx{(‑tannar, ‑tennur)}[2] \dicDirectTranslationCS{tesák, špičák} \dicIndirectTranslationCS{(zub)}
\dicEntry[vígvöllur] \dicTerm{víg··|völlur} \dicIPA{{v}{i}{\textbabygamma}{v}{\oe}{\textsubring{d}}{l}{\textscy}{\textsubring{r}}} \dicPos{m}[11] \dicFlx{(‑vallar, ‑vellir)}[5] \dicDirectTranslationCS{bojiště, bitevní pole}
\dicEntry[vík] \dicTerm{vík\smash{\textsuperscript{1}}} \dicsymFrequent\  \dicIPA{{v}{i}{\textlengthmark}{\r{g}}} \dicPos{f}[10] \dicFlx{(‑ur, ‑ur)}[1] \dicDirectTranslationCS{záliv, zátoka} \dicExampleIS{í litlum víkum} \dicExampleCS{v~malých zálivech}
\dicEntry[vík] \dicTerm{vík\smash{\textsuperscript{2}}} \dicIPA{{v}{i}{\textlengthmark}{\r{g}}} \dicPos{v} \dicFlx{ind praes sg 1 pers} \dicLink{víkja}
\dicEntry[víking] \dicTerm{vík··ing} \dicIPA{{v}{i}{\textlengthmark}{\r{\textObardotlessj}}{i}{\ng}{\r{g}}} \dicPos{f}[4] \dicFlx{(‑ar)}[7] \dicFieldCat{hist.} \dicDirectTranslationCS{vikinská výprava};  \dicPhraseIS{leggjast í víking} \dicFlx{refl} \dicDirectTranslationCS{vydat se na vikinskou výpravu}
\dicEntry[víkingaskip] \dicTerm{víkinga··skip} \dicIPA{{v}{i}{\textlengthmark}{\r{\textObardotlessj}}{i}{\ng}{\r{g}}{a}{s}{\r{\textObardotlessj}}{\textsci}{\textsubring{b}}} \dicPos{n}[2] \dicFlx{(‑s, ‑)}[5] \dicDirectTranslationCS{vikinská loď}
\dicEntry[víkingur] \dicTerm{vík··ing|ur} \dicIPA{{v}{i}{\textlengthmark}{\r{\textObardotlessj}}{i}{\ng}{\r{g}}{\textscy}{\textsubring{r}}} \dicPos{m}[6] \dicFlx{(‑s, ‑ar)}[8] \textbf{1.} \dicFieldCat{hist.} \dicDirectTranslationCS{viking}  \textbf{2.} \dicSynonym*{dugnaðarforkur} \dicDirectTranslationCS{dříč(ka), tahoun(ka)}
\dicEntry[víkja] \dicTerm{víkja} \dicsymFrequent\  \dicIPA{{v}{i}{\textlengthmark}{\r{\textObardotlessj}}{a}} \dicPos{v}[6] \dicFlx{(vík, vék, vikum, viki, vikið)}[61] \dicFlx{dat} \textbf{1.} \dicSynonym{snúa} \dicDirectTranslationCS{uhnout, uhýbat, ustoupit, ustupovat} \dicExampleIS{víkja til hliðar} \dicExampleCS{uhnout stranou};  \dicPhraseIS{víkja fyrir e‑m} \dicDirectTranslationCS{uhnout (komu), ustoupit (komu)} \dicIndirectTranslationCS{(nepřekážet v~pohybu)} \dicExampleIS{víkja úr vegi fyrir e‑m} \dicExampleCS{uhnout (komu) z~cesty}  \textbf{2.} \dicFlx{dat} \dicDirectTranslationCS{propustit, propouštět, vyhodit, vyhazovat (ze zaměstnání ap.)} \dicExampleIS{víkja e‑m úr starfi} \dicExampleCS{propustit (koho) z~práce}  \textbf{3.} \dicPhraseIS{nú víkur sögunni til e‑rs\,/\addthin e‑s} \dicFlx{impers} \dicDirectTranslationCS{a~nyní se příběh obrací ke (komu\,/\addthin čemu)};  \dicIdiom{víkja}[að]{ \dicPhraseIS{víkja e‑u að e‑m}} \dicDirectTranslationCS{dát (co komu), věnovat (co komu)} \dicExampleIS{víkja nokkrum krónum börnunum} \dicExampleCS{dát dětem nějaké peníze}; { \dicPhraseIS{víkja (orðum) að e‑u}} \dicDirectTranslationCS{zmínit se o~(čem), podotknout (co)}; { \dicPhraseIS{víkja sér að e‑m}} \dicDirectTranslationCS{obrátit\,/\addthin otočit se ke (komu)};  \dicIdiom{víkja}[frá]{ \dicPhraseIS{víkja frá e‑m}} \dicDirectTranslationCS{odklonit se od (koho), odchýlit se od (koho)}; { \dicPhraseIS{víkja frá e‑u}} \dicDirectTranslationCS{odklonit se od (čeho), upustit od (čeho)};  \dicIdiom{víkja}[undan]{ \dicPhraseIS{víkja sér undan e‑u}} \dicDirectTranslationCS{ustoupit od (čeho), couvnout od (čeho)};  \dicIdiom{víkjast}[undan]{ \dicPhraseIS{víkjast undan e‑u}} \dicFlx{refl} \dicDirectTranslationCS{vyhnout se (čemu), couvnout od (čeho)}
\dicEntry[víkjandi] \dicTerm{víkj··andi} \dicIPA{{v}{i}{\textlengthmark}{\r{\textObardotlessj}}{a}{n}{\textsubring{d}}{\textsci}} \dicPos{adj}[13] \dicFlx{indecl}[1] \dicDirectTranslationCS{recesivní, ustupující};  \dicPhraseIS{víkjandi gen} \dicFieldCat{biol.} \dicDirectTranslationCS{recesivní gen}
\dicEntry[víkka] \dicTerm{víkk|a} \dicIPA{{v}{i}{h}{\r{g}}{a}} \dicPos{v}[1] \dicFlx{(‑aði)}[1] \dicFlx{acc} \textbf{1.} \dicFlx{acc} \dicSynonym*{gera víðari} \dicDirectTranslationCS{rozšířit, rozšiřovat} \dicExampleIS{víkka buxurnar} \dicExampleCS{rozšířit kalhoty};  \dicPhraseIS{víkka sjóndeildarhring e‑rs} \dicDirectTranslationCS{rozšířit (čí) obzory}  \textbf{2.} \dicSynonym*{verða víðari} \dicDirectTranslationCS{rozšířit se, rozšiřovat se, zvětšit se, zvětšovat se} \dicExampleIS{Gatið víkkaði.} \dicExampleCS{Díra se zvětšila.}
\dicEntry[víkkun] \dicTerm{víkk|un} \dicIPA{{v}{i}{h}{\r{g}}{\textscy}{\textsubring{n}}} \dicPos{f}[7] \dicFlx{(‑unar)}[9] \dicDirectTranslationCS{rozšíření}
\dicEntry[víla] \dicTerm{víl|a} \dicIPA{{v}{i}{\textlengthmark}{l}{a}} \dicPos{v}[1] \dicFlx{(‑aði)}[1] \dicFlx{acc} \dicPhraseIS{víla e‑ð ekki fyrir sér} \dicDirectTranslationCS{nezaleknout se (čeho)}
\dicEntry[víma] \dicTerm{vím|a} \dicIPA{{v}{i}{\textlengthmark}{m}{a}} \dicPos{f}[1] \dicFlx{(‑u, ‑ur)}[19] \dicDirectTranslationCS{opilost, omámenost, opojení, intoxikace} \dicIndirectTranslationCS{(alkoholem, drogami)}
\dicEntry[vímuefni] \dicTerm{vímu··efni} \dicIPA{{v}{i}{\textlengthmark}{m}{\textscy}{\textepsilon}{\textsubring{b}}{n}{\textsci}} \dicPos{n}[2] \dicFlx{(‑s, ‑)}[14] \dicDirectTranslationCS{omamná látka, droga} \dicExampleIS{ólögleg vímuefni} \dicExampleCS{ilegální omamné látky}
\dicEntry[vímuvarnir] \dicTerm{vímu··varnir} \dicIPA{{v}{i}{\textlengthmark}{m}{\textscy}{v}{a}{r}{\textsubring{d}}{n}{\textsci}{\textsubring{r}}} \dicPos{f}[7] \dicFlx{pl}[18] \dicDirectTranslationCS{protialkoholní\,/\addthin protidrogová opatření}
\dicEntry[Vín] \dicTerm{Vín\smash{\textsuperscript{1}}} \dicIPA{{v}{i}{\textlengthmark}{\textsubring{n}}} \dicPos{f} \dicFlx{(‑ar)} \dicFieldCat{geog.} \dicSynonym{Vínarborg} \dicDirectTranslationCS{Vídeň} \dicIndirectTranslationCS{(hlavní město Rakouska)}
\dicEntry[vín] \dicTerm{vín\smash{\textsuperscript{2}}} \dicsymFrequent\  \dicIPA{{v}{i}{\textlengthmark}{\textsubring{n}}} \dicPos{n} \dicFlx{(‑s, ‑)} \textbf{1.} \dicSynonym{léttvín} \dicDirectTranslationCS{víno} \dicExampleIS{rautt og hvítt vín} \dicExampleCS{červené a~bílé víno}  \textbf{2.} \dicSynonym{áfengi} \dicDirectTranslationCS{alkohol}
\dicEntry[vínandi] \dicTerm{vín··|andi} \dicIPA{{v}{i}{\textlengthmark}{n}{a}{n}{\textsubring{d}}{\textsci}} \dicPos{m}[2] \dicFlx{(‑anda)}[2] \dicSynonym{spíritus} \dicDirectTranslationCS{líh, etylalkohol, etanol}
\dicEntry[Vínarborg] \dicTerm{Vínar··borg} \dicIPA{{v}{i}{\textlengthmark}{n}{a}{r}{\textsubring{b}}{\textopeno}{r}{\r{g}}} \dicPos{f}[7] \dicFlx{(‑ar)}[4] \dicFieldCat{geog.} \dicSynonym{Vín\smash{\textsuperscript{1}}} \dicDirectTranslationCS{Vídeň} \dicIndirectTranslationCS{(hlavní město Rakouska)}
\dicEntry[vínarbrauð] \dicTerm{vínar··brauð} \dicIPA{{v}{i}{\textlengthmark}{n}{a}{r}{\textsubring{b}}{r}{\oe i}{\texttheta}} \dicPos{n}[2] \dicFlx{(‑s, ‑)}[5] \dicFieldCat{kulin.} \dicIndirectTranslationCS{koláč z~listového těsta zpravidla polévaný polevou o~různé chuti} \dicsymPhoto\ 
\dicFigure{ds_image_vinarbraud_0_3.jpg}{Vínarbrauð}{Vínarbrauð - hvalur.org, CC BY 3.0}
\dicEntry[vínarpylsa] \dicTerm{vínar··pyls|a} \dicIPA{{v}{i}{\textlengthmark}{n}{a}{\textsubring{r}}{p\smash{\textsuperscript{h}}}{\textsci}{l}{s}{a}} \dicPos{f}[1] \dicFlx{(‑u, ‑ur)}[13] \dicFieldCat{kulin.} \dicDirectTranslationCS{vídeňský párek}
\dicEntry[vínbann] \dicTerm{vín··|bann} \dicIPA{{v}{i}{n}{\textsubring{b}}{a}{\textsubring{n}}} \dicPos{n}[2] \dicFlx{(‑banns, ‑bönn)}[8] \dicSynonym{bann} \dicDirectTranslationCS{prohibice}
\dicEntry[vínber] \dicTerm{vín··ber} \dicIPA{{v}{i}{n}{\textsubring{b}}{\textepsilon}{\textsubring{r}}} \dicPos{n}[2] \dicFlx{(‑s, ‑)}[13] \dicDirectTranslationCS{hroznové víno} \dicIndirectTranslationCS{(plod)}
\dicEntry[vínberjaklasi] \dicTerm{vín·berja··klas|i} \dicIPA{{v}{i}{n}{\textsubring{b}}{\textepsilon}{r}{j}{a}{k\smash{\textsuperscript{h}}}{l}{a}{s}{\textsci}} \dicPos{m}[1] \dicFlx{(‑a, ‑ar)}[8] \dicDirectTranslationCS{hrozen} \dicIndirectTranslationCS{(soubor plodů vinné révy)}
\dicEntry[vínbúð] \dicTerm{vín··búð} \dicIPA{{v}{i}{n}{\textsubring{b}}{u}{\texttheta}} \dicPos{f}[7] \dicFlx{(‑ar, ‑ir)}[1] \dicDirectTranslationCS{obchod s~alkoholem};  \dicPhraseIS{Vínbúð} \dicFlx{prop} \dicIndirectTranslationCS{státní obchod s~alkoholem mající monopol na Islandu}
\dicEntry[vínedik] \dicTerm{vín··edik} \dicIPA{{v}{i}{\textlengthmark}{n}{\textepsilon}{\textsubring{d}}{\textsci}{\r{g}}} \dicPos{n}[2] \dicFlx{(‑s)}[2] \dicDirectTranslationCS{vinný ocet}
\dicEntry[vínflaska] \dicTerm{vín··|flaska} \dicIPA{{v}{i}{n}{f}{l}{a}{s}{\r{g}}{a}} \dicPos{f}[1] \dicFlx{(‑flösku, ‑flöskur)}[14] \textbf{1.} \dicSynonym*{vín í flösku} \dicDirectTranslationCS{láhev vína}  \textbf{2.} \dicSynonym*{flaska undir vín} \dicDirectTranslationCS{láhev na víno}
\dicEntry[vínföng] \dicTerm{vín··föng} \dicIPA{{v}{i}{n}{f}{\oe i}{\ng}{\r{g}}} \dicPos{n}[2] \dicFlx{pl}[9] \dicDirectTranslationCS{alkoholické nápoje}
\dicEntry[vínglas] \dicTerm{vín··|glas} \dicIPA{{v}{i}{n}{\r{g}}{l}{a}{s}} \dicPos{n}[2] \dicFlx{(‑glass, ‑glös)}[8] \textbf{1.} \dicSynonym*{glas af víni} \dicDirectTranslationCS{sklenice vína}  \textbf{2.} \dicSynonym*{glas fyrir vín} \dicDirectTranslationCS{sklenice na víno}
\dicEntry[vínkjallari] \dicTerm{vín··kjall·ar|i} \dicIPA{{v}{i}{n}{c\smash{\textsuperscript{h}}}{a}{\textsubring{d}}{\textsubring{l}}{a}{r}{\textsci}} \dicPos{m}[1] \dicFlx{(‑a, ‑ar)}[10] \dicDirectTranslationCS{vinný sklep}
\dicEntry[vínrækt] \dicTerm{vín··rækt} \dicIPA{{v}{i}{n}{r}{a}{i}{x}{\textsubring{d}}} \dicPos{f}[4] \dicFlx{(‑ar)}[3] \dicSynonym{vínyrkja} \dicDirectTranslationCS{vinařství, vinohradnictví}
\dicEntry[víntegund] \dicTerm{vín··tegund} \dicIPA{{v}{i}{n}{t\smash{\textsuperscript{h}}}{\textepsilon}{\textbabygamma}{\textscy}{n}{\textsubring{d}}} \dicPos{f}[7] \dicFlx{(‑ar, ‑ir)}[1] \dicDirectTranslationCS{druh vína}
\dicEntry[vínveitingar] \dicTerm{vín··veit·ingar} \dicIPA{{v}{i}{n}{v}{ei}{\textsubring{d}}{i}{\ng}{\r{g}}{a}{\textsubring{r}}} \dicPos{f}[4] \dicFlx{pl}[6] \dicDirectTranslationCS{prodej alkoholu} \dicExampleIS{Það eru engar vínveitingar á þessum stað.} \dicExampleCS{Zde se alkohol neprodává.}
\dicEntry[vínviður] \dicTerm{vín··við|ur} \dicIPA{{v}{i}{n}{v}{\textsci}{ð}{\textscy}{\textsubring{r}}} \dicPos{m}[10] \dicFlx{(‑ar, ‑ir)}[19] \dicFieldCat{bot.} \dicDirectTranslationCS{vinná réva} \textit{(l.~{\textLA{Vitis vinifera}})}  \dicsymPhoto\ 
\dicFigure{ds_image_vinvidur_0_1.jpg}{Vínviður}{Vínviður - Fir0002, GFDL}
\dicEntry[vínyrkja] \dicTerm{vín··yrkj|a} \dicIPA{{v}{i}{\textlengthmark}{n}{\textsci}{\textsubring{r}}{\r{\textObardotlessj}}{a}} \dicPos{f}[1] \dicFlx{(‑u)}[5] \dicDirectTranslationCS{vinařství, vinohradnictví}
\dicEntry[víóla] \dicTerm{víól|a} \dicIPA{{v}{i}{j}{ou}{l}{a}} \dicPos{f}[1] \dicFlx{(‑u, ‑ur)}[19] \dicFieldCat{hud.} \dicDirectTranslationCS{viola}
\dicEntry[víóluleikari] \dicTerm{víólu··leik·ar|i} \dicIPA{{v}{i}{j}{ou}{l}{\textscy}{l}{ei}{\r{g}}{a}{r}{\textsci}} \dicPos{m}[1] \dicFlx{(‑a, ‑ar)}[13] \dicDirectTranslationCS{violista, violistka}
\dicEntry[vír] \dicTerm{vír} \dicsymFrequent\  \dicIPA{{v}{i}{\textlengthmark}{\textsubring{r}}} \dicPos{m}[4] \dicFlx{(‑s, ‑ar)}[14] \dicDirectTranslationCS{drát} \dicExampleIS{mjór vír} \dicExampleCS{tenký drát}
\dicEntry[víravirki] \dicTerm{víra··virki} \dicIPA{{v}{i}{\textlengthmark}{r}{a}{v}{\textsci}{\textsubring{r}}{\r{\textObardotlessj}}{\textsci}} \dicPos{n}[2] \dicFlx{(‑s, ‑)}[16] \dicDirectTranslationCS{filigrán} \dicIndirectTranslationCS{(drobný umělecký předmět)}
\dicEntry[vírnet] \dicTerm{vír··net} \dicIPA{{v}{i}{r}{n}{\textepsilon}{\textsubring{d}}} \dicPos{n}[2] \dicFlx{(‑s, ‑)}[5] \dicDirectTranslationCS{drátěné pletivo}
\dicEntry[vírus] \dicTerm{vírus} \dicIPA{{v}{i}{\textlengthmark}{r}{\textscy}{s}} \dicPos{m}[4] \dicFlx{(‑s, ‑ar)}[25] \textbf{1.} \dicFieldCat{med.} \dicSynonym{veira} \dicDirectTranslationCS{virus, vir}  \textbf{2.} \dicFieldCat{poč.} \dicSynonym*{tölvuvírus} \dicDirectTranslationCS{(počítačový) virus\,/\addthin vir}
\dicEntry[vírussjúkdómur] \dicTerm{vírus··sjúk·dóm|ur} \dicIPA{{v}{i}{\textlengthmark}{r}{\textscy}{s}{j}{u}{\r{g}}{\textsubring{d}}{ou}{m}{\textscy}{\textsubring{r}}} \dicPos{m}[6] \dicFlx{(‑s, ‑ar)}[10] \dicFieldCat{med. } \dicDirectTranslationCS{virové onemocnění}
\dicEntry[vís] \dicTerm{vís} \dicsymFrequent\  \dicIPA{{v}{i}{\textlengthmark}{s}} \dicPos{adj}[5]\dicFlx{}[-1] \textbf{1.} \dicSynonym{viss} \dicDirectTranslationCS{jistý, zaručený};  \dicPhraseIS{á vísum stað} \dicFlx{adv} \dicDirectTranslationCS{na bezpečném místě};  \dicPhraseIS{ganga að því vísu að (gera e‑ð)} \dicDirectTranslationCS{pokládat za jisté (udělání (čeho))};  \dicPhraseIS{vera vís til þess að (gera e‑ð)} \dicDirectTranslationCS{zaručeně\,/\addthin jistě (udělat (co))};  \dicPhraseIS{það er nokkuð víst að (e‑að gerist)} \dicDirectTranslationCS{je celkem jisté, že (se (co) stane)}  \textbf{2.} \dicSynonym{vitur} \dicDirectTranslationCS{moudrý, rozumný, znalý};  \dicPhraseIS{verða e‑s vísari} \dicDirectTranslationCS{dozvědět se (co), zjistit (co)};  \dicIdiom{vís}{ \dicPhraseIS{að vísu}} \dicFlx{adv} \dicSynonym{vissulega} \dicDirectTranslationCS{jistě, určitě, zajisté};  \dicPhraseIS{það er víst} \dicFlx{adv} \dicDirectTranslationCS{jistě, určitě}
\dicEntry[vísa] \dicTerm{vís|a\smash{\textsuperscript{1}}} \dicsymFrequent\  \dicIPA{{v}{i}{\textlengthmark}{s}{a}} \dicPos{f}[1] \dicFlx{(‑u, ‑ur)}[13] \textbf{1.} \dicFieldCat{lit.} \dicSynonym{staka} \dicDirectTranslationCS{verš}  \textbf{2.} \dicSynonym{aðferð} \dicDirectTranslationCS{způsob, metoda} \dicExampleIS{á mína vísu} \dicExampleCS{mým způsobem}  \textbf{3.} \dicSynonym{vegabréfsáritun} \dicDirectTranslationCS{vízum};  \dicProverb\  \dicPhraseIS{Aldrei er góð vísa of oft kveðin.} \dicLangCat{přís.} \dicDirectTranslationCS{Nikdy není dobrý verš recitovaný příliš často.} \dicIndirectTranslationCS{(používané, když si lidé stěžují na opakované výtky)}
\dicEntry[vísa] \dicTerm{vís|a\smash{\textsuperscript{2}}} \dicsymFrequent\  \dicIPA{{v}{i}{\textlengthmark}{s}{a}} \dicPos{v}[1] \dicFlx{(‑aði)}[1] \dicFlx{dat (+ acc)} \textbf{1.} \dicDirectTranslationCS{ukázat, ukazovat}  \textbf{2.} \dicDirectTranslationCS{odvolávat se, odkazovat se (na zdroj ap.)};  \dicPhraseIS{vísa e‑m leið\,/\addthin veginn} \dicDirectTranslationCS{ukázat (komu) cestu};  \dicIdiom{vísa}[á]{ \dicPhraseIS{vísa e‑m á dyr}} \dicDirectTranslationCS{vykázat (koho) ze dveří}; { \dicPhraseIS{vísa e‑m á e‑ð}} \dicDirectTranslationCS{ukázat (komu) na (co)}; { \dicPhraseIS{vísa e‑m á e‑n}} \dicDirectTranslationCS{odkázat (koho) na (koho)}; { \dicPhraseIS{vísa e‑u á bug}} \dicDirectTranslationCS{odmítnout (co), popřít (co)};  \dicIdiom{vísa}[burt]{ \dicPhraseIS{vísa e‑m burt}} \dicDirectTranslationCS{vykázat (koho)};  \dicIdiom{vísa}[frá]{ \dicPhraseIS{vísa e‑u\,/\addthin e‑m frá}} \dicDirectTranslationCS{odmítnout (co\,/\addthin koho)};  \dicIdiom{vísa}[í]{ \dicPhraseIS{vísa í e‑ð}} \dicDirectTranslationCS{odvolávat se na (co), odkazovat se na (co), poukazovat na (co)} \dicExampleIS{vísa í fræðirit} \dicExampleCS{odvolávat se na vědeckou publikaci}; { \dicPhraseIS{vísa í heimildir}} \dicDirectTranslationCS{odvolávat se na prameny};  \dicIdiom{vísa}[til]{ \dicPhraseIS{vísa e‑m til sætis}} \dicDirectTranslationCS{ukázat (komu) jeho místo}; { \dicPhraseIS{vísa e‑m til vegar}} \dicDirectTranslationCS{ukázat (komu) cestu}; { \dicPhraseIS{vísa til e‑s}} \dicDirectTranslationCS{odvolat\,/\addthin odvolávat se na (co)};  \dicIdiom{vísa}[úr]{ \dicPhraseIS{vísa e‑m úr skóla}} \dicDirectTranslationCS{vyloučit (koho) ze školy};  \dicIdiom{vísast}{ \dicPhraseIS{vísast}} \dicFlx{refl} \dicDirectTranslationCS{odkazovat se, být odkazován}
\dicEntry[vísbending] \dicTerm{vís··bend·ing} \dicsymFrequent\  \dicIPA{{v}{i}{s}{\textsubring{b}}{\textepsilon}{n}{\textsubring{d}}{i}{\ng}{\r{g}}} \dicPos{f}[4] \dicFlx{(‑ar, ‑ar)}[5] \dicDirectTranslationCS{vodítko, stopa, nápověda, indicie} \dicExampleIS{gefa e‑m vísbendingu um e‑ð} \dicExampleCS{poskytnout (komu) vodítko k~(čemu)}
\dicEntry[vísdómstönn] \dicTerm{vís·dóms··|tönn} \dicIPA{{v}{i}{s}{\textsubring{d}}{ou}{m}{s}{t\smash{\textsuperscript{h}}}{\oe}{\textsubring{n}}} \dicPos{f}[8] \dicFlx{(‑tannar, ‑tennur)}[2] \dicDirectTranslationCS{zub moudrosti}
\dicEntry[vísdómur] \dicTerm{vís··dóm|ur} \dicIPA{{v}{i}{s}{\textsubring{d}}{ou}{m}{\textscy}{\textsubring{r}}} \dicPos{m}[6] \dicFlx{(‑s)}[12] \dicDirectTranslationCS{moudrost}
\dicEntry[vísifingur] \dicTerm{vísi··fing|ur} \dicsymFrequent\  \dicIPA{{v}{i}{\textlengthmark}{s}{\textsci}{f}{i}{\ng}{\r{g}}{\textscy}{\textsubring{r}}} \dicPos{m}[5] \dicFlx{(‑urs, ‑ur)}[13] \dicDirectTranslationCS{ukazovák, ukazováček} \dicExampleIS{banka á það með vísifingri} \dicExampleCS{poklepat na to ukazováčkem}
\dicEntry[vísindagrein] \dicTerm{vísinda··grein} \dicIPA{{v}{i}{\textlengthmark}{s}{\textsci}{n}{\textsubring{d}}{a}{\r{g}}{r}{ei}{\textsubring{n}}} \dicPos{f}[4] \dicFlx{(‑ar, ‑ar)}[1] \dicDirectTranslationCS{vědní obor\,/\addthin odvětví}
\dicEntry[vísindakenning] \dicTerm{vísinda··kenn·ing} \dicIPA{{v}{i}{\textlengthmark}{s}{\textsci}{n}{\textsubring{d}}{a}{c\smash{\textsuperscript{h}}}{\textepsilon}{n}{i}{\ng}{\r{g}}} \dicPos{f}[4] \dicFlx{(‑ar, ‑ar)}[5] \dicDirectTranslationCS{vědecká teorie}
\dicEntry[vísindalegur] \dicTerm{vísinda··legur} \dicsymFrequent\  \dicIPA{{v}{i}{\textlengthmark}{s}{\textsci}{n}{\textsubring{d}}{a}{l}{\textepsilon}{\textbabygamma}{\textscy}{\textsubring{r}}} \dicPos{adj}[1]\dicFlx{}[-8] \dicDirectTranslationCS{vědecký} \dicExampleIS{breiðara vísindalegt sjónarmið} \dicExampleCS{širší vědecké hledisko}
\dicEntry[vísindamaður] \dicTerm{vísinda··|maður} \dicsymFrequent\  \dicIPA{{v}{i}{\textlengthmark}{s}{\textsci}{n}{\textsubring{d}}{a}{m}{a}{ð}{\textscy}{\textsubring{r}}} \dicPos{m}[13] \dicFlx{(‑manns, ‑menn)}[2] \dicDirectTranslationCS{vědec, vědkyně} \dicExampleIS{hópur af erlendum vísindamönnum} \dicExampleCS{skupina zahraničních vědkyň}
\dicEntry[vísindarit] \dicTerm{vísinda··rit} \dicIPA{{v}{i}{\textlengthmark}{s}{\textsci}{n}{\textsubring{d}}{a}{r}{\textsci}{\textsubring{d}}} \dicPos{n}[2] \dicFlx{(‑s, ‑)}[5] \textbf{1.} \dicDirectTranslationCS{vědecký časopis}  \textbf{2.} \dicDirectTranslationCS{vědecká publikace}

\dicEntry[vísindaskáldsaga] \dicTerm{vísinda··skáld·|saga} \dicIPA{{v}{i}{\textlengthmark}{s}{\textsci}{n}{\textsubring{d}}{a}{s}{\r{g}}{au}{l}{\textsubring{d}}{s}{a}{\textbabygamma}{a}} \dicPos{f}[1] \dicFlx{(‑sögu, ‑sög\-ur)}[14] \dicFieldCat{lit.} \dicDirectTranslationCS{sci-fi, science fiction, vědecko-fantastický příběh}

\begin{xtolerant}{}{1pt}
\dicEntry[vísindaskáldskapur] \dicTerm{vísinda··skáld·skap|ur}\addthinS\dicIPA{{v}\-{i}\-{\textlengthmark}\-{s}\-{\textsci}\-{n}\-{\textsubring{d}}\-{a}\-{s}\-{\r{g}}\-{au}\-{l}\-{\textsubring{d}}\-{s}\-{\r{g}}\-{a}\-{\textsubring{b}}\-{\textscy}\-{\textsubring{r}}} \dicPos{m}[10]\addthinS\dicFlx{(‑ar)}[15] \dicFieldCat{lit.} \dicDirectTranslationCS{sci-fi, science fiction, vědecko-fantastická literatura}
\end{xtolerant}
\dicEntry[vísindaskyn] \dicTerm{vísinda··skyn} \dicIPA{{v}{i}{\textlengthmark}{s}{\textsci}{n}{\textsubring{d}}{a}{s}{\r{\textObardotlessj}}{\textsci}{\textsubring{n}}} \dicPos{n}[2] \dicFlx{(‑s)}[2] \dicPhraseIS{í vísindaskyni} \dicFlx{adv} \dicDirectTranslationCS{pro výzkumné\,/\addthin vědecké účely}
\dicEntry[vísindi] \dicTerm{vís··indi} \dicIPA{{v}{i}{\textlengthmark}{s}{\textsci}{n}{\textsubring{d}}{\textsci}} \dicPos{n}[2] \dicFlx{pl}[19] \dicDirectTranslationCS{věda} \dicExampleIS{stunda vísindi} \dicExampleCS{zabývat se vědou}
\dicEntry[vísir] \dicTerm{vís|ir} \dicIPA{{v}{i}{\textlengthmark}{s}{\textsci}{\textsubring{r}}} \dicPos{m}[7] \dicFlx{(‑is, ‑ar)}[1] \textbf{1.} \dicSynonym{spíra\smash{\textsuperscript{1}}} \dicDirectTranslationCS{výhonek, šlahoun};  \dicPhraseIS{e‑að er vísir að e‑u} \dicDirectTranslationCS{(co) je zárodkem\,/\addthin počátkem (čeho)}  \textbf{2.} \dicSynonym{armur\smash{\textsuperscript{1}}} \dicDirectTranslationCS{ručička, ukazatel} \dicIndirectTranslationCS{(u~hodin)}
\dicEntry[vísitala] \dicTerm{vísi··|tala} \dicIPA{{v}{i}{\textlengthmark}{s}{\textsci}{t\smash{\textsuperscript{h}}}{a}{l}{a}} \dicPos{f}[1] \dicFlx{(‑tölu, ‑tölur)}[14] \dicDirectTranslationCS{index, koeficient} \dicExampleIS{vísitala gengis} \dicExampleCS{index směnného kurzu}
\dicEntry[vísitölubatt] \dicTerm{vísi·tölu··batt} \dicIPA{{v}{i}{\textlengthmark}{s}{\textsci}{t\smash{\textsuperscript{h}}}{\oe}{l}{\textscy}{\textsubring{b}}{a}{h}{\textsubring{d}}} \dicPos{v} \dicFlx{ind pf sg 1 pers} \dicLink{vísitölubinda}
\dicEntry[vísitölubinda] \dicTerm{vísi·tölu··|binda} \dicIPA{{v}{i}{\textlengthmark}{s}{\textsci}{t\smash{\textsuperscript{h}}}{\oe}{l}{\textscy}{\textsubring{b}}{\textsci}{n}{\textsubring{d}}{a}} \dicPos{v}[6] \dicFlx{(‑bind, ‑batt, ‑bundum, ‑byndi, ‑bundið)}[132] \dicFlx{acc} \dicFieldCat{ekon.} \dicDirectTranslationCS{indexovat, indexově vázat}
\dicEntry[vísitölubinding] \dicTerm{vísi·tölu··bind·ing} \dicIPA{{v}{i}{\textlengthmark}{s}{\textsci}{t\smash{\textsuperscript{h}}}{\oe}{l}{\textscy}{\textsubring{b}}{\textsci}{n}{\textsubring{d}}{i}{\ng}{\r{g}}} \dicPos{f}[4] \dicFlx{(‑ar)}[7] \dicFieldCat{ekon.} \dicDirectTranslationCS{indexace}
\dicEntry[vísitölubundið] \dicTerm{vísi·tölu··bundið} \dicIPA{{v}{i}{\textlengthmark}{s}{\textsci}{t\smash{\textsuperscript{h}}}{\oe}{l}{\textscy}{\textsubring{b}}{\textscy}{n}{\textsubring{d}}{\textsci}{\texttheta}} \dicPos{v} \dicFlx{supin} \dicLink{vísitölubinda}
\dicEntry[vísitölubyndi] \dicTerm{vísi·tölu··byndi} \dicIPA{{v}{i}{\textlengthmark}{s}{\textsci}{}{t\smash{\textsuperscript{h}}}{\oe}{l}{\textscy}{\textsubring{b}}{\textsci}{n}{\textsubring{d}}{\textsci}} \dicPos{v} \dicFlx{con pf sg 1 pers} \dicLink{vísitölubinda}
\dicEntry[vísla] \dicTerm{vísl|a} \dicIPA{{v}{i}{s}{\textsubring{d}}{l}{a}} \dicPos{f}[1] \dicFlx{(‑u, ‑ur)}[13] \dicFieldCat{zool.} \dicSynonym{hreysiköttur} \dicDirectTranslationCS{hranostaj, lasice hranostaj} \textit{(l.~{\textLA{Mustela erminea}})} 
\dicEntry[víst] \dicTerm{vís|t} \dicsymFrequent\  \dicIPA{{v}{i}{s}{\textsubring{d}}} \dicPos{adv} \dicFlx{(sup ‑ast)} \textbf{1.} \dicSynonym{örugglega} \dicDirectTranslationCS{jistě, najisto, určitě, zaručeně} \dicExampleIS{Það er víst rétt.} \dicExampleCS{Je to určitě pravda.};  \dicPhraseIS{vita e‑ð fyrir víst} \dicDirectTranslationCS{vědět (co) najisto}  \textbf{2.} \dicSynonym*{sem staðfesting} \dicDirectTranslationCS{jistě, určitě} \dicIndirectTranslationCS{(jako potvrzení)}  \textbf{3.} \dicSynonym{einmitt} \dicDirectTranslationCS{ale ano, určitě jo} \dicIndirectTranslationCS{(jako protimluva)}
\dicEntry[vísun] \dicTerm{vís|un} \dicIPA{{v}{i}{\textlengthmark}{s}{\textscy}{\textsubring{n}}} \dicPos{f}[7] \dicFlx{(‑unar, ‑anir)}[8] \textbf{1.} \dicDirectTranslationCS{odkaz, odvolávání se} \dicExampleIS{vísun til heimilda} \dicExampleCS{odvolávání se na prameny}  \textbf{2.} \dicFieldCat{lit.} \dicDirectTranslationCS{aluze}
\dicEntry[vísundur] \dicTerm{vísund|ur} \dicIPA{{v}{i}{\textlengthmark}{s}{\textscy}{n}{\textsubring{d}}{\textscy}{\textsubring{r}}} \dicPos{m}[6] \dicFlx{(‑s\,/\addthin ‑ar, ‑ar)}[63] \dicFieldCat{zool.} \dicDirectTranslationCS{bizon} \textit{(l.~{\textLA{Bison}})}  \dicsymPhoto\ 
\dicFigure{1333.jpg}{Vísundur}{Vísundur - Zicha Ondřej, Biolib, Copyright/CC-BY-NC}
\dicEntry[vísuorð] \dicTerm{vísu··orð} \dicIPA{{v}{i}{\textlengthmark}{s}{\textscy}{\textopeno}{r}{\texttheta}} \dicPos{n}[2] \dicFlx{(‑s, ‑)}[5] \dicFieldCat{lit.} \dicSynonym{braglína} \dicDirectTranslationCS{verš}
\dicEntry[vísvitandi] \dicTerm{vís··vit·andi\smash{\textsuperscript{1}}} \dicIPA{{v}{i}{s}{v}{\textsci}{\textsubring{d}}{a}{n}{\textsubring{d}}{\textsci}} \dicPos{adj}[13] \dicFlx{indecl}[1] \dicDirectTranslationCS{vědomý, úmyslný, záměrný} \dicExampleIS{Þetta eru vísvitandi lygar.} \dicExampleCS{To jsou vědomé lži.}
\dicEntry[vísvitandi] \dicTerm{vís··vit·andi\smash{\textsuperscript{2}}} \dicIPA{{v}{i}{s}{v}{\textsci}{\textsubring{d}}{a}{n}{\textsubring{d}}{\textsci}} \dicPos{adv} \dicDirectTranslationCS{vědomě, úmyslně, záměrně, schválně} \dicExampleIS{brjóta samninginn vísvitandi} \dicExampleCS{porušit úmyslně smlouvu}
\dicEntry[víta] \dicTerm{vít|a} \dicIPA{{v}{i}{\textlengthmark}{\textsubring{d}}{a}} \dicPos{v}[2] \dicFlx{(‑ti, ‑t)}[57] \dicFlx{acc} \dicSynonym{ávíta} \dicDirectTranslationCS{(po)kárat, napomenout, napomínat};  \dicPhraseIS{víta e‑n fyrir e‑ð} \dicDirectTranslationCS{pokárat (koho) za (co)}
\dicEntry[vítahringur] \dicTerm{víta··hring|ur} \dicIPA{{v}{i}{\textlengthmark}{\textsubring{d}}{a}{\textsubring{r}}{i}{\ng}{\r{g}}{\textscy}{\textsubring{r}}} \dicPos{m}[6] \dicFlx{(‑s)}[17] \textbf{1.} \dicLangCat{přen.} \dicDirectTranslationCS{začarovaný kruh}  \textbf{2.} \dicFieldCat{filos.} \dicDirectTranslationCS{bludný kruh} \textit{(l.~{\textLA{circulus vitiosus}})}
\dicEntry[vítakast] \dicTerm{víta··|kast} \dicIPA{{v}{i}{\textlengthmark}{\textsubring{d}}{a}{k\smash{\textsuperscript{h}}}{a}{s}{\textsubring{d}}} \dicPos{n}[2] \dicFlx{(‑kasts, ‑köst)}[8] \dicFieldCat{sport.} \dicDirectTranslationCS{trestný hod}
\dicEntry[vítamín] \dicTerm{vítamín} \dicIPA{{v}{i}{\textlengthmark}{\textsubring{d}}{a}{m}{i}{\textsubring{n}}} \dicPos{n}[2] \dicFlx{(‑s, ‑)}[5] \dicDirectTranslationCS{vitamín}
\dicEntry[vítaspyrna] \dicTerm{víta··spyrn|a} \dicIPA{{v}{i}{\textlengthmark}{\textsubring{d}}{a}{s}{\textsubring{b}}{\textsci}{r}{\textsubring{d}}{n}{a}} \dicPos{f}[1] \dicFlx{(‑u, ‑ur)}[7] \dicFieldCat{sport.} \dicDirectTranslationCS{pokutový kop, penalta}
\dicEntry[vítaverður] \dicTerm{víta··verður} \dicIPA{{v}{i}{\textlengthmark}{\textsubring{d}}{a}{v}{\textepsilon}{r}{ð}{\textscy}{\textsubring{r}}} \dicPos{adj}[2]\dicFlx{}[-1] \dicSynonym{ámælisverður} \dicDirectTranslationCS{trestuhodný, zavrženíhodný, odsouzeníhodný} \dicExampleIS{vítavert athæfi} \dicExampleCS{trestuhodné jednání}
\dicEntry[víti] \dicTerm{víti} \dicIPA{{v}{i}{\textlengthmark}{\textsubring{d}}{\textsci}} \dicPos{n}[2] \dicFlx{(‑s, ‑)}[14] \textbf{1.} \dicFieldCat{sport.} \dicSynonym*{vítaspyrna\,/\addthin vítakast} \dicDirectTranslationCS{pokutový kop\,/\addthin hod} \dicExampleIS{skora úr vítinu} \dicExampleCS{skórovat z~pokutového kopu}  \textbf{2.} \dicSynonym{helvíti\smash{\textsuperscript{1}}} \dicDirectTranslationCS{peklo}
\dicEntry[vítt] \dicTerm{vítt} \dicIPA{{v}{i}{h}{\textsubring{d}}} \dicPos{adv} \dicFlx{(comp víðar\,/\addthin víðara, sup víðast)} \dicDirectTranslationCS{široko};  \dicPhraseIS{vítt og breitt} \dicFlx{adv} \dicDirectTranslationCS{široko daleko}
\dicEntry[vítur] \dicTerm{vítur} \dicIPA{{v}{i}{\textlengthmark}{\textsubring{d}}{\textscy}{\textsubring{r}}} \dicPos{f}[12] \dicFlx{pl}[7] \dicDirectTranslationCS{důtka, napomenutí}
\dicEntry[víxill] \dicTerm{víx|ill} \dicIPA{{v}{i}{x}{s}{\textsci}{\textsubring{d}}{\textsubring{l}}} \dicPos{m}[6] \dicFlx{(‑ils, ‑lar)}[35] \dicFieldCat{ekon.} \dicDirectTranslationCS{směnka} \dicExampleIS{afsegja víxil} \dicExampleCS{odmítnout směnku}
\dicEntry[víxillán] \dicTerm{víxil··lán} \dicIPA{{v}{i}{x}{s}{\textsci}{l}{au}{\textsubring{n}}} \dicPos{n}[2] \dicFlx{(‑s, ‑)}[5] \dicFieldCat{ekon.} \dicDirectTranslationCS{půjčka na směnku}
\dicEntry[víxl] \dicTerm{víxl} \dicIPA{{v}{i}{x}{s}{\textsubring{d}}{\textsubring{l}}} \dicPos{n}[2] \dicFlx{(‑s)}[2] \dicDirectTranslationCS{záměna, výměna, prohození};  \dicPhraseIS{á víxl} \dicFlx{adv} \dicDirectTranslationCS{střídavě, proměnlivě}
\dicEntry[víxla] \dicTerm{víxl|a} \dicIPA{{v}{i}{x}{s}{\textsubring{d}}{l}{a}} \dicPos{v}[1] \dicFlx{(‑aði)}[1] \dicFlx{dat} \textbf{1.} \dicSynonym*{skipta á} \dicDirectTranslationCS{zaměnit, zaměňovat, prohodit} \dicExampleIS{víxla stöfum í e‑u} \dicExampleCS{zaměnit písmena v~(čem)}  \textbf{2.} \dicSynonym*{skipta um} \dicDirectTranslationCS{směnit, směňovat, vyměnit, vyměňovat (cizí měnu ap.)} \dicExampleIS{víxla peningum} \dicExampleCS{směnit peníze};  \dicIdiom{víxlast}{ \dicPhraseIS{víxlast}} \dicFlx{refl} \dicDirectTranslationCS{zaměnit se, prohodit se}
\dicEntry[víxlrím] \dicTerm{víxl··rím} \dicIPA{{v}{i}{x}{s}{\textsubring{d}}{l}{r}{i}{\textsubring{m}}} \dicPos{n}[2] \dicFlx{(‑s, ‑)}[5] \dicFieldCat{lit.} \dicDirectTranslationCS{střídavý rým}
\dicEntry[víxlverkun] \dicTerm{víxl··verk|un} \dicIPA{{v}{i}{x}{s}{\textsubring{d}}{l}{v}{\textepsilon}{\textsubring{r}}{\r{g}}{\textscy}{\textsubring{n}}} \dicPos{f}[7] \dicFlx{(‑unar, ‑anir)}[8] \dicDirectTranslationCS{interakce, vzájemné působení}
\dicEntry[vodka] \dicTerm{vodka}\dicTerm{, vodki} \dicIPA{{v}\-{\textopeno}\-{\textsubring{d}}\-{\r{g}}\-{a}\-} \dicPos{n}[4] \dicFlx{indecl}[1] \dicDirectTranslationCS{vodka}
\dicEntry[vodki] \dicTerm{vodk|i} \dicIPA{{v}{\textopeno}{\textsubring{d}}{\r{\textObardotlessj}}{\textsci}} \dicPos{m}[1] \dicFlx{(‑a)}[3] \dicLink{vodka}
\dicEntry[voð] \dicTerm{voð} \dicIPA{{v}{\textopeno}{\textlengthmark}{\texttheta}} \dicPos{f}[7] \dicFlx{(‑ar, ‑ir)}[1] \dicDirectTranslationCS{sukno, tkanina} \dicExampleIS{vefa voðina} \dicExampleCS{tkát sukno}
\dicEntry[voða] \dicTerm{voða} \dicIPA{{v}{\textopeno}{\textlengthmark}{ð}{a}} \dicPos{adv} \dicSynonym{ægilega} \dicDirectTranslationCS{strašně, příšerně, fantasticky} \dicExampleIS{voða sætur} \dicExampleCS{strašně roztomilý}
\dicEntry[voðalega] \dicTerm{voða··lega} \dicIPA{{v}{\textopeno}{\textlengthmark}{ð}{a}{l}{\textepsilon}{\textbabygamma}{a}} \dicPos{adv} \dicDirectTranslationCS{hrozně, příšerně, strašně} \dicExampleIS{voðalega fallegur} \dicExampleCS{hrozně krásný}
\dicEntry[voðalegur] \dicTerm{voða··legur} \dicIPA{{v}{\textopeno}{\textlengthmark}{ð}{a}{l}{\textepsilon}{\textbabygamma}{\textscy}{\textsubring{r}}} \dicPos{adj}[1]\dicFlx{}[-8] \textbf{1.} \dicSynonym{ægilegur} \dicDirectTranslationCS{příšerný, hrozný, strašný} \dicIndirectTranslationCS{(ke zdůraznění)} \dicExampleIS{voðalegur hávaði} \dicExampleCS{příšerný hluk}  \textbf{2.} \dicSynonym{hræðilegur} \dicDirectTranslationCS{strašlivý, děsivý}
\dicEntry[voði] \dicTerm{voð|i} \dicsymFrequent\  \dicIPA{{v}{\textopeno}{\textlengthmark}{ð}{\textsci}} \dicPos{m}[1] \dicFlx{(‑a, ‑ar)}[1] \dicSynonym{hætta\smash{\textsuperscript{1}}} \dicDirectTranslationCS{nebezpečí, ohrožení, hrozba};  \dicPhraseIS{fara sér að voða} \dicDirectTranslationCS{vystavit se nebezpečí, riskovat};  \dicPhraseIS{stofna e‑u í voða} \dicDirectTranslationCS{ohrožovat (co), být hrozbou pro (co)}
\dicEntry[vofa] \dicTerm{vof|a\smash{\textsuperscript{1}}} \dicIPA{{v}{\textopeno}{\textlengthmark}{v}{a}} \dicPos{f}[1] \dicFlx{(‑u, ‑ur)}[19] \dicDirectTranslationCS{duch, strašidlo, přízrak, zjevení}
\dicEntry[vofa] \dicTerm{vof|a\smash{\textsuperscript{2}}} \dicIPA{{v}{\textopeno}{\textlengthmark}{v}{a}} \dicPos{v}[2] \dicFlx{(‑ði, ‑að)}[128] \dicSynonym{svífa} \dicDirectTranslationCS{vznášet se, viset};  \dicPhraseIS{vofa yfir e‑u} \dicDirectTranslationCS{vznášet se nad (čím)} \dicExampleIS{Ógn vofir yfir þjóðinni.} \dicExampleCS{Strach se vznáší nad národem.}
\dicEntry[voff] \dicTerm{voff\smash{\textsuperscript{1}}}\dicTerm{, vóff} \dicIPA{{v}\-{\textopeno}\-{f}\-{\textlengthmark}\-} \dicPos{n}[2] \dicFlx{(‑s)}[2] \dicDirectTranslationCS{hafání, bafání, štěkání}
\dicEntry[voff] \dicTerm{voff\smash{\textsuperscript{2}}} \dicIPA{{v}{\textopeno}{f}{\textlengthmark}} \dicPos{inter} \dicDirectTranslationCS{haf, baf, ňaf}
\dicEntry[voffi] \dicTerm{voff|i} \dicIPA{{v}{\textopeno}{f}{\textlengthmark}{\textsci}} \dicPos{m}[1] \dicFlx{(‑a, ‑ar)}[1] \dicLangCat{dět.} \dicSynonym{hundur} \dicDirectTranslationCS{hafánek, hafík, pejsek, psíček}
\dicEntry[vog] \dicTerm{vog} \dicIPA{{v}{\textopeno}{\textlengthmark}{x}} \dicPos{f}[7] \dicFlx{(‑ar, ‑ir)}[1] \textbf{1.} \dicSynonym{vigt} \dicDirectTranslationCS{váha} \dicIndirectTranslationCS{(přístroj na vážení)}  \textbf{2.} \dicDirectTranslationCS{Váhy} \dicIndirectTranslationCS{(znamení zvěrokruhu)} \textit{(l.~{\textLA{Libra}})}
\dicEntry[voga] \dicTerm{vog|a} \dicIPA{{v}{\textopeno}{\textlengthmark}{\textbabygamma}{a}} \dicPos{v}[1] \dicFlx{(‑aði)}[1] \dicFlx{dat} \dicDirectTranslationCS{odvážit se, troufnout si} \dicIndirectTranslationCS{(nejčastěji se záporem)};  \dicPhraseIS{voga sér ekki að (gera e‑ð)} \dicDirectTranslationCS{neodvážit se ((co) udělat)}
\dicEntry[vogarmerki] \dicTerm{vogar··merki} \dicIPA{{v}{\textopeno}{\textlengthmark}{\textbabygamma}{a}{r}{m}{\textepsilon}{\textsubring{r}}{\r{\textObardotlessj}}{\textsci}} \dicPos{n}[2] \dicFlx{(‑s)}[20] \dicSynonym{vog} \dicDirectTranslationCS{znamení Vah}
\dicEntry[vogarskál] \dicTerm{vogar··skál} \dicIPA{{v}{\textopeno}{\textlengthmark}{\textbabygamma}{a}{\textsubring{r}}{s}{\r{g}}{au}{\textsubring{l}}} \dicPos{f}[4] \dicFlx{(‑ar, ‑ar\,/\addthin ‑ir)}[30] \dicDirectTranslationCS{miska vah};  \dicPhraseIS{leggja (sitt) lóð á vogarskál} \dicLangCat{přen.} \dicDirectTranslationCS{přispět svým dílem}
\dicEntry[vogarstöng] \dicTerm{vogar··|stöng} \dicIPA{{v}{\textopeno}{\textlengthmark}{\textbabygamma}{a}{\textsubring{r}}{s}{\textsubring{d}}{\oe i}{\ng}{\r{g}}} \dicPos{f}[8] \dicFlx{(‑stangar, ‑stengur\,/\addthin ‑stangir)}[9] \dicDirectTranslationCS{páka} \dicIndirectTranslationCS{(jednoduchý stroj)}
\dicEntry[vogrís] \dicTerm{vog··rís} \dicIPA{{v}{\textopeno}{\textlengthmark}{\r{g}}{r}{i}{s}} \dicPos{m}[4] \dicFlx{(‑s, ‑ar)}[14] \dicFieldCat{med. } \dicDirectTranslationCS{ječné zrno} \textit{(l.~{\textLA{hordeolum}})}
\dicEntry[vogskorinn] \dicTerm{vog··skorinn} \dicIPA{{v}{\textopeno}{x}{s}{\r{g}}{\textopeno}{r}{\textsci}{\textsubring{n}}} \dicPos{adj}[6]\dicFlx{}[-2] \dicDirectTranslationCS{členitý (pobřeží ap.)} \dicExampleIS{vogskorin strandlengja} \dicExampleCS{členité pobřeží}
\dicEntry[vogun] \dicTerm{vog|un} \dicIPA{{v}{\textopeno}{\textlengthmark}{\textbabygamma}{\textscy}{\textsubring{n}}} \dicPos{f}[7] \dicFlx{(‑unar)}[9] \dicDirectTranslationCS{risk, riskování, riziko};  \dicPhraseIS{vogun vinnur, vogun tapar} \dicLangCat{přen.} \dicDirectTranslationCS{risk je zisk}
\dicEntry[vogur] \dicTerm{vog|ur} \dicIPA{{v}{\textopeno}{\textlengthmark}{\textbabygamma}{\textscy}{\textsubring{r}}} \dicPos{m}[6] \dicFlx{(‑s, ‑ar)}[19] \dicSynonym*{vík} \dicDirectTranslationCS{zátoka, zátočina}
\dicEntry[vol] \dicTerm{vol} \dicIPA{{v}{\textopeno}{\textlengthmark}{\textsubring{l}}} \dicPos{n}[2] \dicFlx{(‑s)}[2] \dicDirectTranslationCS{fňukání, naříkání, skuhrání}
\dicEntry[vola] \dicTerm{vol|a} \dicIPA{{v}{\textopeno}{\textlengthmark}{l}{a}} \dicPos{v}[1] \dicFlx{(‑aði)}[1] \dicSynonym{kjökra} \dicDirectTranslationCS{fňukat, naříkat, skuhrat} \dicExampleIS{vola út af e‑u} \dicExampleCS{fňukat kvůli (čemu)}
\dicEntry[voldugur] \dicTerm{vold··ugur} \dicsymFrequent\  \dicIPA{{v}{\textopeno}{l}{\textsubring{d}}{\textscy}{\textbabygamma}{\textscy}{\textsubring{r}}} \dicPos{adj}[1]\dicFlx{}[-8] \textbf{1.} \dicSynonym{máttugur} \dicDirectTranslationCS{mocný, vlivný, silný} \dicExampleIS{voldugt fyrirtæki} \dicExampleCS{vlivná společnost}  \textbf{2.} \dicSynonym{stór} \dicDirectTranslationCS{mohutný, masivní, robustní}
\dicEntry[volfram] \dicTerm{volfram} \dicIPA{{v}{\textopeno}{l}{f}{r}{a}{\textsubring{m}}} \dicPos{n}[2] \dicFlx{(‑s)}[2] \dicFieldCat{chem.} \dicDirectTranslationCS{wolfram} \textit{(l.~{\textLA{W, Wolframium}})}
\dicEntry[Volga] \dicTerm{Volg|a} \dicIPA{{v}{\textopeno}{l}{\r{g}}{a}} \dicPos{f}[1] \dicFlx{(‑u)}[6] \dicFieldCat{geog.} \dicDirectTranslationCS{Volha} \dicIndirectTranslationCS{(řeka v~evropské části Ruska)}
\dicEntry[volgna] \dicTerm{volgn|a} \dicIPA{{v}{\textopeno}{l}{n}{a}} \dicPos{v}[1] \dicFlx{(‑aði)}[45] \dicSynonym{hlýna} \dicDirectTranslationCS{ohřát se, ohřívat se, nahřát se, nahřívat se, oteplit se, oteplovat se} \dicExampleIS{Ofninn volgnar.} \dicExampleCS{Trouba se ohřívá.}
\dicEntry[volgur] \dicTerm{volgur} \dicsymFrequent\  \dicIPA{{v}{\textopeno}{l}{\r{g}}{\textscy}{\textsubring{r}}} \dicPos{adj}[1]\dicFlx{}[-1] \textbf{1.} \dicSynonym{hlýr} \dicDirectTranslationCS{vlažný, vlahý, (mírně) teplý} \dicExampleIS{Brauðið er ennþá volgt.} \dicExampleCS{Chléb je ještě teplý.}  \textbf{2.} \dicDirectTranslationCS{vlažný, nedostatečně horlivý, nepříliš nadšený (návrhem ap.)}
\dicEntry[volk] \dicTerm{volk} \dicIPA{{v}{\textopeno}{\textsubring{l}}{\r{g}}} \dicPos{n}[2] \dicFlx{(‑s)}[2] \dicSynonym{slark} \dicDirectTranslationCS{útrapy, strádání, trampoty (na cestách ap.)}
\dicEntry[voll] \dicTerm{voll} \dicIPA{{v}{\textopeno}{l}{\textlengthmark}} \dicPos{n}[2] \dicFlx{(‑s)}[2] \dicPhraseIS{það er allt í volli} \dicLangCat{přen.} \dicDirectTranslationCS{všechno je vzhůru nohama}
\dicEntry[volt] \dicTerm{volt} \dicIPA{{v}{\textopeno}{\textsubring{l}}{\textsubring{d}}} \dicPos{n}[2] \dicFlx{(‑s, ‑)}[5] \dicFieldCat{fyz.} \dicDirectTranslationCS{volt}
\dicEntry[volæði] \dicTerm{vol··æði} \dicIPA{{v}{\textopeno}{\textlengthmark}{l}{a}{i}{ð}{\textsci}} \dicPos{n}[2] \dicFlx{(‑s)}[20] \dicSynonym{eymd} \dicDirectTranslationCS{mizérie, bída}
\dicEntry[von] \dicTerm{von} \dicsymFrequent\  \dicIPA{{v}{\textopeno}{\textlengthmark}{\textsubring{n}}} \dicPos{f}[7] \dicFlx{(‑ar, ‑ir)}[1] \dicDirectTranslationCS{naděje, očekávání};  \dicIdiom{von}{ \dicPhraseIS{að vonum}} \dicFlx{adv} \dicDirectTranslationCS{podle očekávání};  \dicPhraseIS{binda vonir við e‑ð} \dicDirectTranslationCS{vkládat\,/\addthin skládat naděje v~(co)};  \dicPhraseIS{eiga von á barni} \dicDirectTranslationCS{očekávat dítě}; { \dicPhraseIS{eiga von á e‑u\,/\addthin e‑m}} \dicDirectTranslationCS{očekávat (co\,/\addthin koho)};  \dicPhraseIS{framar vonum} \dicFlx{adv} \dicDirectTranslationCS{nad očekávání};  \dicPhraseIS{gefa upp alla von} \dicDirectTranslationCS{vzdát se vší naděje};  \dicPhraseIS{gera sér vonir um e‑ð} \dicDirectTranslationCS{dělat si naděje na (co)};  \dicPhraseIS{halda í vonina} \dicDirectTranslationCS{udržovat si naději};  \dicPhraseIS{í von um að (e‑að gerist)} \dicDirectTranslationCS{v~očekávání, že (se (co) stane)};  \dicPhraseIS{lifa í voninni} \dicDirectTranslationCS{žít v~naději};  \dicPhraseIS{milli vonar og ótta} \dicFlx{adv} \dicDirectTranslationCS{v~nejistotě};  \dicPhraseIS{til vonar og vara} \dicFlx{adv} \dicDirectTranslationCS{pro jistotu};  \dicPhraseIS{upp á von og óvon} \dicFlx{adv} \dicDirectTranslationCS{kdyby náhodou};  \dicPhraseIS{vera úrkula vonar} \dicDirectTranslationCS{ztratit naději, vzdát se naděje};  \dicPhraseIS{það er borin von að (e‑að gerist)} \dicFlx{impers} \dicDirectTranslationCS{je naděje, že (se (co) stane)};  \dicPhraseIS{það er ekki von á góðu} \dicFlx{impers} \dicDirectTranslationCS{nevypadá to nadějně};  \dicPhraseIS{það er ekki öll von úti} \dicFlx{impers} \dicDirectTranslationCS{ještě zbývá naděje, ještě není vše ztraceno};  \dicPhraseIS{það standa vonir til þess að (e‑að gerist)} \dicFlx{impers} \dicDirectTranslationCS{je naděje, že (se (co) stane)}
\dicEntry[vona] \dicTerm{von|a} \dicsymFrequent\  \dicIPA{{v}{\textopeno}{\textlengthmark}{n}{a}} \dicPos{v}[1] \dicFlx{(‑aði)}[1] \dicDirectTranslationCS{doufat, mít naději} \dicExampleIS{Ég vona að þið hafið notið kvöldsins.} \dicExampleCS{Doufám, že jste si užili večer.};  \dicPhraseIS{vona hið besta} \dicDirectTranslationCS{doufat v~to nejlepší};  \dicIdiom{vonast}[eftir]{ \dicPhraseIS{vonast eftir e‑u}} \dicFlx{refl} \dicDirectTranslationCS{doufat v~(co)};  \dicIdiom{vonast}[til]{ \dicPhraseIS{vonast til e‑s}} \dicFlx{refl} \dicDirectTranslationCS{doufat v~(co)}
\dicEntry[vonandi] \dicTerm{von··andi} \dicsymFrequent\  \dicIPA{{v}{\textopeno}{\textlengthmark}{n}{a}{n}{\textsubring{d}}{\textsci}} \dicPos{adv} \dicSynonym{sennilega} \dicDirectTranslationCS{snad, doufejme} \dicExampleIS{Vonandi verður það til einhverrar hjálpar.} \dicExampleCS{Snad to nějak pomůže.}
\dicEntry[vonaraugu] \dicTerm{vonar··augu} \dicIPA{{v}{\textopeno}{\textlengthmark}{n}{a}{r}{\oe i}{\textbabygamma}{\textscy}} \dicPos{n}[1] \dicFlx{pl}[3] \dicPhraseIS{horfa vonaraugum til e‑s} \dicFlx{adv} \dicDirectTranslationCS{vyhlížet (co) s~nadějí v~očích}
\dicEntry[vonarglæta] \dicTerm{vonar··glæt|a} \dicIPA{{v}{\textopeno}{\textlengthmark}{n}{a}{r}{\r{g}}{l}{a}{i}{\textsubring{d}}{a}} \dicPos{f}[1] \dicFlx{(‑u)}[5] \dicDirectTranslationCS{jiskra\,/\addthin jiskřička naděje, záblesk naděje}
\dicEntry[vonarneisti] \dicTerm{vonar··neist|i} \dicIPA{{v}{\textopeno}{\textlengthmark}{n}{a}{r}{n}{ei}{s}{\textsubring{d}}{\textsci}} \dicPos{m}[1] \dicFlx{(‑a)}[3] \dicDirectTranslationCS{jiskra\,/\addthin jiskřička naděje, záblesk naděje}
\dicEntry[vonarvölur] \dicTerm{vonar··|völur} \dicIPA{{v}{\textopeno}{\textlengthmark}{n}{a}{r}{v}{\oe}{l}{\textscy}{\textsubring{r}}} \dicPos{m}[11] \dicFlx{(‑valar)}[8] \dicDirectTranslationCS{žebrácká hůl};  \dicPhraseIS{vera kominn á vonarvöl} \dicDirectTranslationCS{dostat se na mizinu}
\dicEntry[vonbrigði] \dicTerm{von··brigði} \dicsymFrequent\  \dicIPA{{v}{\textopeno}{n}{\textsubring{b}}{r}{\textsci}{\textbabygamma}{ð}{\textsci}} \dicPos{n}[2] \dicFlx{pl}[19] \dicDirectTranslationCS{zklamání, rozčarování, deziluze};  \dicPhraseIS{valda e‑m vonbrigðum} \dicDirectTranslationCS{způsobit (komu) zklamání, zklamat (koho)};  \dicPhraseIS{verða fyrir vonbrigðum} \dicDirectTranslationCS{být zklamaný, zklamat se, zažít zklamání}
\dicEntry[vondaufur] \dicTerm{von··daufur} \dicIPA{{v}{\textopeno}{n}{\textsubring{d}}{\oe i}{v}{\textscy}{\textsubring{r}}} \dicPos{adj}[1]\dicFlx{}[-1] \dicDirectTranslationCS{(jsoucí) s~malou nadějí, nemající velká očekávání, nedělající si velké naděje}
\dicEntry[vondur] \dicTerm{vondur} \dicsymFrequent\  \dicIPA{{v}{\textopeno}{n}{\textsubring{d}}{\textscy}{\textsubring{r}}} \dicPos{adj}[11] \dicFlx{(comp verri, sup verstur)}[5] \textbf{1.} \dicSynonym{illur} \dicDirectTranslationCS{špatný, zlý};  \dicPhraseIS{vera vondur við e‑n} \dicDirectTranslationCS{být na (koho) zlý, chovat se ke (komu) špatně}  \textbf{2.} \dicSynonym{reiður} \dicDirectTranslationCS{rozzlobený, rozezlený};  \dicPhraseIS{vondur út í e‑n} \dicDirectTranslationCS{rozzlobený na (koho)}  \textbf{3.} \dicSynonym{slæmur} \dicDirectTranslationCS{špatný, mizerný (počasí ap.)} \dicExampleIS{vondur vegur} \dicExampleCS{špatná cesta}  \textbf{4.} \dicSynonym{lélegur} \dicDirectTranslationCS{špatný (řidič ap.)};  \dicIdiom{vondur}{ \dicPhraseIS{gera illt verra}} \dicDirectTranslationCS{zhoršit, udělat to horší};  \dicPhraseIS{hafa verra af} \dicLangCat{přen.} \dicDirectTranslationCS{mít po legraci};  \dicPhraseIS{það er verri sagan að (e‑að gerðist)} \dicFlx{impers} \dicDirectTranslationCS{horší je, že (se (co) stalo)}
\dicEntry[vonglaður] \dicTerm{von··|glaður} \dicIPA{{v}{\textopeno}{n}{\r{g}}{l}{a}{ð}{\textscy}{\textsubring{r}}} \dicPos{adj}[2] \dicFlx{(f ‑glöð)}[7] \dicDirectTranslationCS{(jsoucí) plný naděje, šťastný}
\dicEntry[vongóður] \dicTerm{von··|góður} \dicIPA{{v}{\textopeno}{n}{\r{g}}{ou}{ð}{\textscy}{\textsubring{r}}} \dicPos{adj}[11] \dicFlx{(comp ‑betri, sup ‑bestur)}[1] \dicDirectTranslationCS{nadějný, slibný, perspektivní}
\dicEntry[vonlaus] \textls[10]{\dicTerm{von··laus} \dicsymFrequent\  \dicIPA{{v}{\textopeno}{n}{l}{\oe i}{s}} \dicPos{adj}[5]\dicFlx{}[-1] \textbf{1.} \dicSynonym*{án vonar} \dicDirectTranslationCS{beznadějný, zoufalý}  \textbf{2.} \dicSynonym{lélegur} \dicDirectTranslationCS{beznadějný (kuchař ap.)}  \textbf{3.} \dicSynonym{ómögulegur} \dicDirectTranslationCS{beznadějný, marný, nemožný} \dicExampleIS{vonlaust verk} \dicExampleCS{marná práce};  \dicPhraseIS{það er vonlaust að (gera e‑ð)} \dicDirectTranslationCS{je beznadějné (dělat (co))}}
\dicEntry[vonleysi] \dicTerm{von··leysi} \dicIPA{{v}{\textopeno}{n}{l}{ei}{s}{\textsci}} \dicPos{n}[2] \dicFlx{(‑s)}[20] \dicDirectTranslationCS{beznaděj, zoufalství, marnost}
\dicEntry[vonska] \dicTerm{vonsk|a} \dicIPA{{v}{\textopeno}{n}{s}{\r{g}}{a}} \dicPos{f}[1] \dicFlx{(‑u)}[5] \textbf{1.} \dicSynonym{illska} \dicDirectTranslationCS{špatnost, zlo}  \textbf{2.} \dicSynonym{reiði\smash{\textsuperscript{2}}} \dicDirectTranslationCS{zlost, zloba, hněv} \dicExampleIS{Það hljóp í hann vonska.} \dicExampleCS{Rozhněval se.}
\dicEntry[vonsvikinn] \dicTerm{von··svikinn} \dicIPA{{v}{\textopeno}{n}{s}{v}{\textsci}{\r{\textObardotlessj}}{\textsci}{\textsubring{n}}} \dicPos{adj}[6]\dicFlx{}[-2] \dicDirectTranslationCS{zklamaný, rozčarovaný} \dicExampleIS{vonsvikinn með samstarfið} \dicExampleCS{zklamaný ze spolupráce}
\dicEntry[vopn] \dicTerm{vopn} \dicsymFrequent\  \dicIPA{{v}{\textopeno}{h}{\textsubring{b}}{\textsubring{n}}} \dicPos{n}[2] \dicFlx{(‑s, ‑)}[5] \dicDirectTranslationCS{zbraň};  \dicPhraseIS{grípa til vopna} \dicDirectTranslationCS{chopit se zbraně};  \dicPhraseIS{leggja niður vopn(in)} \dicDirectTranslationCS{složit zbraně};  \dicPhraseIS{slá vopnin úr höndum e‑rs} \dicDirectTranslationCS{odzbrojit (koho)}
\dicEntry[vopna] \dicTerm{vopn|a} \dicIPA{{v}{\textopeno}{h}{\textsubring{b}}{n}{a}} \dicPos{v}[1] \dicFlx{(‑aði)}[1] \dicFlx{acc} \dicDirectTranslationCS{(o)zbrojit, vyzbrojit} \dicExampleIS{vopna e‑n} \dicExampleCS{vyzbrojit (koho)};  \dicIdiom{vopnast}{ \dicPhraseIS{vopnast}} \dicFlx{refl} \dicDirectTranslationCS{ozbrojit se, vyzbrojit se};  \dicPhraseIS{vopnast e‑u} \dicFlx{refl} \dicDirectTranslationCS{vyzbrojit se (čím)}
\dicEntry[vopnabúnaður] \dicTerm{vopna··bú·nað|ur} \dicIPA{{v}{\textopeno}{h}{\textsubring{b}}{n}{a}{\textsubring{b}}{u}{n}{a}{ð}{\textscy}{\textsubring{r}}} \dicPos{m}[10] \dicFlx{(‑ar)}[9] \dicDirectTranslationCS{výzbroj}
\dicEntry[vopnabúr] \dicTerm{vopna··búr} \dicIPA{{v}{\textopeno}{h}{\textsubring{b}}{n}{a}{\textsubring{b}}{u}{\textsubring{r}}} \dicPos{n}[2] \dicFlx{(‑s, ‑)}[5] \dicDirectTranslationCS{skladiště zbraní, zbrojnice}
\dicEntry[vopnaður] \dicTerm{vopn··|aður} \dicIPA{{v}{\textopeno}{h}{\textsubring{b}}{n}{a}{ð}{\textscy}{\textsubring{r}}} \dicPos{adj}[3] \dicFlx{(f ‑uð)}[3] \dicDirectTranslationCS{ozbrojený, vyzbrojený};  \dicPhraseIS{vera vopnaður e‑u} \dicDirectTranslationCS{být ozbrojený (čím)};  \dicPhraseIS{vopnað rán} \dicDirectTranslationCS{ozbrojená loupež}
\dicEntry[vopnahlé] \dicTerm{vopna··hlé} \dicIPA{{v}{\textopeno}{h}{\textsubring{b}}{n}{a}{\textsubring{l}}{j}{\textepsilon}} \dicPos{n}[2] \dicFlx{(‑s, ‑)}[5] \dicDirectTranslationCS{příměří, klid zbraní} \dicExampleIS{semja um vopnahlé} \dicExampleCS{vyjednat příměří}
\dicEntry[vopnlaus] \dicTerm{vopn··laus} \dicIPA{{v}{\textopeno}{h}{\textsubring{b}}{n}{l}{\oe i}{s}} \dicPos{adj}[5]\dicFlx{}[-1] \dicDirectTranslationCS{neozbrojený, nevyzbrojený}
\dicEntry[vor] \dicTerm{vor\smash{\textsuperscript{1}}} \dicsymFrequent\  \dicIPA{{v}{\textopeno}{\textlengthmark}{\textsubring{r}}} \dicPos{n}[2] \dicFlx{(‑s, ‑)}[5] \dicDirectTranslationCS{jaro};  \dicPhraseIS{á vorin} \dicFlx{adv} \dicDirectTranslationCS{na jaře};  \dicPhraseIS{í vor} \dicFlx{adv} \dicDirectTranslationCS{na jaře};  \dicPhraseIS{með vorinu} \dicFlx{adv} \dicDirectTranslationCS{na jaře, během jara};  \dicPhraseIS{um vorið} \dicFlx{adv} \dicDirectTranslationCS{na jaře, během jara}
\dicEntry[vor] \dicTerm{vor\smash{\textsuperscript{2}}} \dicIPA{{v}{\textopeno}{\textlengthmark}{\textsubring{r}}} \dicPos{pron}[4] \dicFlx{pers gen pl} \dicLink{vér}
\dicEntry[vor] \dicTerm{vor\smash{\textsuperscript{3}}} \dicIPA{{v}{\textopeno}{\textlengthmark}{\textsubring{r}}} \dicPos{pron}[4] \dicFlx{poss} \dicLangCat{kniž.} \dicDirectTranslationCS{náš, naše}
\dicEntry[vora] \dicTerm{vor|a} \dicIPA{{v}{\textopeno}{\textlengthmark}{r}{a}} \dicPos{v}[1] \dicFlx{(‑aði)}[83] \dicFlx{impers} \dicPhraseIS{það vorar} \dicDirectTranslationCS{přichází jaro, blíží se jaro}
\dicEntry[vorboði] \dicTerm{vor··boð|i} \dicIPA{{v}{\textopeno}{r}{\textsubring{b}}{\textopeno}{ð}{\textsci}} \dicPos{m}[1] \dicFlx{(‑a, ‑ar)}[1] \dicDirectTranslationCS{předzvěst jara, první příznaky jara};  \dicPhraseIS{vorboðinn ljúfi} \dicSynonym{þröstur} \dicDirectTranslationCS{drozd}
\dicEntry[vordagur] \dicTerm{vor··dag|ur} \dicIPA{{v}{\textopeno}{r}{\textsubring{d}}{a}{\textbabygamma}{\textscy}{\textsubring{r}}} \dicPos{m}[6] \dicFlx{(‑s, ‑ar)}[62] \dicDirectTranslationCS{jarní den};  \dicPhraseIS{á vordögum} \dicFlx{adv} \dicDirectTranslationCS{na jaře}
\dicEntry[vorjafndægur] \dicTerm{vor··jafn·dægur} \dicIPA{{v}{\textopeno}{r}{j}{a}{\textsubring{b}}{\textsubring{n}}{\textsubring{d}}{a}{i}{\textbabygamma}{\textscy}{\textsubring{r}}} \dicPos{n}[2] \dicFlx{(‑s, ‑)}[25] \dicDirectTranslationCS{jarní rovnodennost} \dicAntonym{haustjafndægur}
\dicEntry[vorkenna] \dicTerm{vor··kenn|a} \dicIPA{{v}{\textopeno}{\textsubring{r}}{c\smash{\textsuperscript{h}}}{\textepsilon}{n}{a}} \dicPos{v}[2] \dicFlx{(‑di, ‑t)}[160] \dicFlx{dat} \dicSynonym*{kenna í brjósti um} \dicDirectTranslationCS{(po)litovat, soucítit} \dicExampleIS{vorkenna honum} \dicExampleCS{litovat ho}
\dicEntry[vorkunn] \dicTerm{vor··kunn} \dicIPA{{v}{\textopeno}{\textsubring{r}}{k\smash{\textsuperscript{h}}}{\textscy}{\textsubring{n}}} \dicPos{f}[7] \dicFlx{(‑ar)}[3] \dicSynonym{samúð} \dicDirectTranslationCS{lítost, (po)litování, soucit} \dicExampleIS{sýna e‑m vorkunn} \dicExampleCS{litovat (koho)};  \dicPhraseIS{e‑m er vorkunn} \dicFlx{impers} \dicDirectTranslationCS{člověk musí soucítit s~(kým)};  \dicPhraseIS{virða e‑m e‑ð til vorkunnar} \dicDirectTranslationCS{přičíst (co komu) k~dobrému}
\dicEntry[vorkunnsamur] \dicTerm{vor·kunn··|samur} \dicIPA{{v}{\textopeno}{\textsubring{r}}{k\smash{\textsuperscript{h}}}{\textscy}{n}{s}{a}{m}{\textscy}{\textsubring{r}}} \dicPos{adj}[1] \dicFlx{(f ‑söm)}[2] \dicDirectTranslationCS{soucitný, soucítící}
\dicEntry[vorkunnsemi] \dicTerm{vor·kunn··sem|i} \dicIPA{{v}{\textopeno}{\textsubring{r}}{k\smash{\textsuperscript{h}}}{\textscy}{n}{s}{\textepsilon}{m}{\textsci}} \dicPos{f}[2] \dicFlx{(‑i)}[2] \dicDirectTranslationCS{soucit, soucítění}
\dicEntry[vorlegur] \dicTerm{vor··legur} \dicIPA{{v}{\textopeno}{r}{l}{\textepsilon}{\textbabygamma}{\textscy}{\textsubring{r}}} \dicPos{adj}[1]\dicFlx{}[-8] \dicDirectTranslationCS{jarní, připomínající jaro}
\dicEntry[vormisseri] \dicTerm{vor··misseri} \dicIPA{{v}{\textopeno}{r}{m}{\textsci}{s}{\textepsilon}{r}{\textsci}} \dicPos{n}[2] \dicFlx{(‑s, ‑)}[14] \dicFieldCat{škol.} \dicDirectTranslationCS{jarní\,/\addthin letní semestr} \dicAntonym{haustmisseri}
\dicEntry[vorum] \dicTerm{vorum} \dicIPA{{v}{\textopeno}{\textlengthmark}{r}{\textscy}{\textsubring{m}}} \dicPos{v} \dicFlx{ind pf pl 1 pers} \dicLink{vera\smash{\textsuperscript{2}}}
\dicEntry[vorönn] \dicTerm{vor··|önn} \dicIPA{{v}{\textopeno}{\textlengthmark}{r}{\oe}{\textsubring{n}}} \dicPos{f}[7] \dicFlx{(‑annar, ‑annir)}[16] \dicFieldCat{škol.} \dicDirectTranslationCS{jarní\,/\addthin letní semestr} \dicAntonym{haustönn}
\dicEntry[vos] \dicTerm{vos} \dicIPA{{v}{\textopeno}{\textlengthmark}{s}} \dicPos{n}[2] \dicFlx{(‑s, ‑)}[5] \dicDirectTranslationCS{mokro a~zima (při pobytu v~přírodě ap.)}
\dicEntry[vosbúð] \dicTerm{vos··búð} \dicIPA{{v}{\textopeno}{\textlengthmark}{s}{\textsubring{b}}{u}{\texttheta}} \dicPos{f}[7] \dicFlx{(‑ar)}[3] \dicDirectTranslationCS{útrapy\,/\addthin strádání (zima a~vlhko) na cestách díky špatnému počasí }
\dicEntry[vothey] \dicTerm{vot··hey} \dicIPA{{v}{\textopeno}{\textlengthmark}{\textsubring{d}}{h}{ei}} \dicPos{n}[2] \dicFlx{(‑s, ‑)}[13] \dicDirectTranslationCS{siláž} \dicIndirectTranslationCS{(zkvašené krmivo)}
\dicEntry[votheysturn] \dicTerm{vot·heys··turn} \dicIPA{{v}{\textopeno}{\textlengthmark}{\textsubring{d}}{h}{ei}{s}{t\smash{\textsuperscript{h}}}{\textscy}{r}{\textsubring{d}}{\textsubring{n}}} \dicPos{m}[4] \dicFlx{(‑s, ‑ar)}[3] \dicDirectTranslationCS{silo}
\dicEntry[votheysveiki] \dicTerm{vot·heys··veik|i} \dicIPA{{v}{\textopeno}{\textlengthmark}{\textsubring{d}}{h}{ei}{s}{v}{ei}{\r{\textObardotlessj}}{\textsci}} \dicPos{f}[3] \dicFlx{(‑i)}[3] \dicLink{Hvanneyrarveiki}
\dicEntry[votlendi] \dicTerm{vot··lendi} \dicIPA{{v}{\textopeno}{\textlengthmark}{\textsubring{d}}{l}{\textepsilon}{n}{\textsubring{d}}{\textsci}} \dicPos{n}[2] \dicFlx{(‑s)}[20] \dicDirectTranslationCS{mokřina, mokřisko, mokřad}
\dicEntry[votlendur] \dicTerm{vot··lendur} \dicIPA{{v}{\textopeno}{\textlengthmark}{\textsubring{d}}{l}{\textepsilon}{n}{\textsubring{d}}{\textscy}{\textsubring{r}}} \dicPos{adj}[2]\dicFlx{}[-14] \dicDirectTranslationCS{mokřinatý}
\dicEntry[votta] \dicTerm{vott|a} \dicsymFrequent\  \dicIPA{{v}{\textopeno}{h}{\textsubring{d}}{a}} \dicPos{v}[1] \dicFlx{(‑aði)}[7] \dicFlx{(dat +) acc} \textbf{1.} \dicFlx{acc} \dicDirectTranslationCS{dosvědčit, dosvědčovat, stvrdit, stvrzovat}  \textbf{2.} \dicFlx{dat + acc} \dicDirectTranslationCS{vyjádřit, vyjadřovat, projevit, projevovat} \dicExampleIS{votta honum virðingu sína} \dicExampleCS{vyjádřit mu svoji úctu};  \dicIdiom{votta}[fyrir]{ \dicPhraseIS{það vottar fyrir e‑u}} \dicFlx{impers} \dicDirectTranslationCS{svědčí to o~(čem), ukazuje to na (co)} \dicExampleIS{Það vottar aðeins fyrir bólgu í fætinum.} \dicExampleCS{Svědčí to pouze o~zánětu v~noze.}
\dicEntry[vottfastur] \dicTerm{vott··|fastur} \dicIPA{{v}{\textopeno}{h}{\textsubring{d}}{f}{a}{s}{\textsubring{d}}{\textscy}{\textsubring{r}}} \dicPos{adj}[1] \dicFlx{(f ‑föst)}[12] \dicDirectTranslationCS{stvrzený, potvrzený, osvědčený}
\dicEntry[vottfesta] \dicTerm{vott··fest|a} \dicIPA{{v}{\textopeno}{h}{\textsubring{d}}{f}{\textepsilon}{s}{\textsubring{d}}{a}} \dicPos{v}[2] \dicFlx{(‑i, ‑)}[12] \dicFlx{acc} \dicDirectTranslationCS{stvrdit, stvrzovat, osvědčit, osvědčovat}
\dicEntry[vottorð] \dicTerm{vott··orð} \dicIPA{{v}{\textopeno}{h}{\textsubring{d}}{\textopeno}{r}{\texttheta}} \dicPos{n}[2] \dicFlx{(‑s, ‑)}[5] \textbf{1.} \dicSynonym{læknisvottorð} \dicDirectTranslationCS{lékařské osvědčení\,/\addthin potvrzení}  \textbf{2.} \dicSynonym{skjal} \dicDirectTranslationCS{osvědčení, certifikát, diplom} \dicExampleIS{vottorð um námsárangur} \dicExampleCS{osvědčení o~studijních výsledcích}
\dicEntry[vottun] \dicTerm{vott|un} \dicIPA{{v}{\textopeno}{h}{\textsubring{d}}{\textscy}{\textsubring{n}}} \dicPos{f}[7] \dicFlx{(‑unar)}[9] \dicDirectTranslationCS{osvědčení, certifikát}
\dicEntry[vottur] \dicTerm{vott|ur} \dicIPA{{v}{\textopeno}{h}{\textsubring{d}}{\textscy}{\textsubring{r}}} \dicPos{m}[6] \dicFlx{(‑s, ‑ar)}[4] \textbf{1.} \dicSynonym{vitni} \dicDirectTranslationCS{svědek, svědkyně} \dicExampleIS{Vottar Jehóva} \dicExampleCS{svědkové Jehovovi};  \dicPhraseIS{í votta viðurvist} \dicFlx{adv} \dicDirectTranslationCS{v~přítomnosti svědků};  \dicPhraseIS{vera vottur að e‑u} \dicDirectTranslationCS{být svědkem (čeho)}  \textbf{2.} \dicSynonym{ögn\smash{\textsuperscript{1}}} \dicDirectTranslationCS{troška, špetka} \dicIndirectTranslationCS{(nejčastěji se záporem)} \dicExampleIS{Ég á ekki vott af kaffi.} \dicExampleCS{Nemám ani trochu kafe.};  \dicIdiom{vottur}{ \dicPhraseIS{e‑að ber vott um e‑ð}} \dicDirectTranslationCS{(co) nese svědectví o~(čem)}
\dicEntry[votur] \dicTerm{votur} \dicIPA{{v}{\textopeno}{\textlengthmark}{\textsubring{d}}{\textscy}{\textsubring{r}}} \dicPos{adj}[1]\dicFlx{}[-1] \dicSynonym{blautur} \dicDirectTranslationCS{vlhký, mokrý};  \dicPhraseIS{vera votur í fæturna} \dicDirectTranslationCS{mít mokré nohy}
\dicEntry[voveiflega] \dicTerm{vo·veif··lega} \dicIPA{{v}{\textopeno}{\textlengthmark}{v}{ei}{v}{l}{\textepsilon}{\textbabygamma}{a}} \dicPos{adv} \dicDirectTranslationCS{tragicky, neblaze}
\dicEntry[voveiflegur] \dicTerm{vo·veif··legur} \dicIPA{{v}{\textopeno}{\textlengthmark}{v}{ei}{v}{l}{\textepsilon}{\textbabygamma}{\textscy}{\textsubring{r}}} \dicPos{adj}[1]\dicFlx{}[-8] \dicSynonym{hræðilegur} \dicDirectTranslationCS{tragic\-ký, neblahý} \dicExampleIS{deyja á voveiflegan hátt} \dicExampleCS{zemřít tragickým způsobem}
\dicEntry[vó] \dicTerm{vó} \dicIPA{{v}{ou}{\textlengthmark}} \dicPos{v} \dicFlx{ind pf sg 1 pers} \textbf{1.} \dicLink{vega\smash{\textsuperscript{1}}}  \textbf{2.} \dicLink{vega\smash{\textsuperscript{2}}}
\dicEntry[vóff] \dicTerm{vóff} \dicIPA{{v}{ou}{f}{\textlengthmark}} \dicPos{inter} \dicLink{voff}
\dicEntry[vógum] \dicTerm{vógum} \dicIPA{{v}{ou}{\textlengthmark}{\textbabygamma}{\textscy}{\textsubring{m}}} \dicPos{v} \dicFlx{ind pf pl 1 pers} \textbf{1.} \dicLink{vega\smash{\textsuperscript{1}}}  \textbf{2.} \dicLink{vega\smash{\textsuperscript{2}}}
\dicEntry[VS] \dicTerm{VS} \dicPos{zkr} \dicPhraseIS{vatnssalerni} \dicDirectTranslationCS{WC, záchodky}
\dicEntry[VSK] \dicTerm{VSK} \dicPos{zkr} \dicPhraseIS{virðisaukaskattur} \dicFieldCat{ekon.} \dicDirectTranslationCS{DPH}
\dicEntry[væði] \dicTerm{væði} \dicIPA{{v}{a}{i}{\textlengthmark}{ð}{\textsci}} \dicPos{v} \dicFlx{con pf sg 1 pers} \dicLink{vaða\smash{\textsuperscript{2}}}
\dicEntry[væður] \dicTerm{væður} \dicIPA{{v}{a}{i}{\textlengthmark}{ð}{\textscy}{\textsubring{r}}} \dicPos{adj}[2]\dicFlx{}[-6] \dicDirectTranslationCS{přebroditelný, (jsoucí) k~přebrodění} \dicExampleIS{Áin er ekki væð.} \dicExampleCS{Řeka se nedá přebrodit.}
\dicEntry[væfi] \dicTerm{væfi} \dicIPA{{v}{a}{i}{\textlengthmark}{v}{\textsci}} \dicPos{v} \dicFlx{con pf sg 1 pers} \dicLink{vefa}
\dicEntry[vægð] \dicTerm{vægð} \dicIPA{{v}{a}{i}{\textbabygamma}{\texttheta}} \dicPos{f}[7] \dicFlx{(‑ar, ‑ir)}[1] \dicSynonym{náð\smash{\textsuperscript{1}}} \dicDirectTranslationCS{shovívavost, slitování, milost};  \dicPhraseIS{biðjast vægðar} \dicFlx{refl} \dicDirectTranslationCS{prosit o~shovívavost}
\dicEntry[vægðarlaus] \dicTerm{vægðar··laus} \dicIPA{{v}{a}{i}{\textbabygamma}{ð}{a}{r}{l}{\oe i}{s}} \dicPos{adj}[5]\dicFlx{}[-1] \dicSynonym{miskunnarlaus} \dicDirectTranslationCS{nemilosrdný, nelítostný, neúprosný} \dicExampleIS{vægðarlaus gagnrýni} \dicExampleCS{nelítostná kritika}
\dicEntry[vægðarlaust] \dicTerm{vægðar··laust} \dicIPA{{v}{a}{i}{\textbabygamma}{ð}{a}{r}{l}{\oe i}{s}{\textsubring{d}}} \dicPos{adv} \dicDirectTranslationCS{nemilosrdně, nelítostně, neúprosně}
\dicEntry[vægðarleysi] \dicTerm{vægðar··leysi} \dicIPA{{v}{a}{i}{\textbabygamma}{ð}{a}{r}{l}{ei}{s}{\textsci}} \dicPos{n}[2] \dicFlx{(‑s)}[20] \dicDirectTranslationCS{nemilosrdnost, nelítostnost, neúprosnost}
\dicEntry[vægi] \dicTerm{vægi\smash{\textsuperscript{1}}} \dicIPA{{v}{a}{i}{j}{\textlengthmark}{\textsci}} \dicPos{n}[2] \dicFlx{(‑s)}[20] \dicDirectTranslationCS{váha, význam, podíl} \dicExampleIS{vægi ritgerðarinnar á prófi} \dicExampleCS{podíl písemné práce na zkoušce}
\dicEntry[vægi] \dicTerm{vægi\smash{\textsuperscript{2}}} \dicIPA{{v}{a}{i}{j}{\textlengthmark}{\textsci}} \dicPos{v} \dicFlx{con pf sg 1 pers} \textbf{1.} \dicLink{vega\smash{\textsuperscript{1}}}  \textbf{2.} \dicLink{vega\smash{\textsuperscript{2}}}
\dicEntry[vægilegur] \dicTerm{vægi··legur} \dicIPA{{v}{a}{i}{j}{\textlengthmark}{\textsci}{l}{\textepsilon}{\textbabygamma}{\textscy}{\textsubring{r}}} \dicPos{adj}[1]\dicFlx{}[-8] \dicSynonym{vægur} \dicDirectTranslationCS{mírný, shovívavý}
\dicEntry[væginn] \dicTerm{væginn} \dicIPA{{v}{a}{i}{j}{\textlengthmark}{\textsci}{\textsubring{n}}} \dicPos{adj}[6]\dicFlx{}[-2] \dicSynonym{vægilegur} \dicDirectTranslationCS{mírný, shovívavý}
\dicEntry[vægja] \dicTerm{væg|ja} \dicIPA{{v}{a}{i}{j}{\textlengthmark}{a}} \dicPos{v}[2] \dicFlx{(‑ði, ‑t)}[97] \dicFlx{dat} \dicSynonym{hlífa} \dicDirectTranslationCS{slitovat se, smilovat se, ušetřit} \dicExampleIS{vægja e‑m} \dicExampleCS{slitovat se nad (kým)};  \dicProverb\  \dicPhraseIS{Sá vægir\,/\addthin skal vægja sem vitið hefur meira.} \dicLangCat{přís.} \dicDirectTranslationCS{Moudřejší ustoupí.}
\dicEntry[vægur] \dicTerm{vægur} \dicIPA{{v}{a}{i}{\textlengthmark}{\textbabygamma}{\textscy}{\textsubring{r}}} \dicPos{adj}[1]\dicFlx{}[-1] \textbf{1.} \dicSynonym*{hófstilltur} \dicDirectTranslationCS{mírný, lehký} \dicExampleIS{vægt blótsyrði} \dicExampleCS{mírná nadávka};  \dicPhraseIS{vægast sagt} \dicFlx{adv} \dicDirectTranslationCS{mírně řečeno}  \textbf{2.} \dicSynonym{miskunnsamur} \dicDirectTranslationCS{shovívavý, mírný} \dicExampleIS{vera of vægur við hann} \dicExampleCS{být k~němu příliš shovívavý}
\dicEntry[væl] \dicTerm{væl} \dicIPA{{v}{a}{i}{\textlengthmark}{\textsubring{l}}} \dicPos{n}[2] \dicFlx{(‑s, ‑)}[5] \textbf{1.} \dicSynonym*{grátur} \dicDirectTranslationCS{fňukání, kňourání}  \textbf{2.} \dicSynonym*{væl í sírenu} \dicDirectTranslationCS{houkání, kvílení, vytí (sirény ap.)}  \textbf{3.} \dicSynonym*{ugluvæl} \dicDirectTranslationCS{houkání} \dicIndirectTranslationCS{(o~zvuku, který vydává sova)}
\dicEntry[væla] \dicTerm{væl|a} \dicIPA{{v}{a}{i}{\textlengthmark}{l}{a}} \dicPos{v}[2] \dicFlx{(‑di, ‑t)}[160] \textbf{1.} \dicSynonym*{gráta lágt} \dicDirectTranslationCS{fňukat, kňourat} \dicExampleIS{Hættu nú að væla.} \dicExampleCS{Přestaň fňukat.}  \textbf{2.} \dicDirectTranslationCS{houkat, kvílet, výt (siréna ap.)}  \textbf{3.} \dicDirectTranslationCS{houkat} \dicIndirectTranslationCS{(o~sovách)}
\dicEntry[vælulegur] \dicTerm{vælu··legur} \dicIPA{{v}{a}{i}{\textlengthmark}{l}{\textscy}{l}{\textepsilon}{\textbabygamma}{\textscy}{\textsubring{r}}} \dicPos{adj}[1]\dicFlx{}[-8] \dicDirectTranslationCS{ufňukaný, ukňouraný}
\dicEntry[væluskjóða] \textls[15]{\dicTerm{vælu··skjóð|a} \dicIPA{{v}{a}{i}{\textlengthmark}{l}{\textscy}{s}{\r{\textObardotlessj}}{ou}{ð}{a}}} \dicPos{f}[1] \dicFlx{(‑u, ‑ur)}[19] \dicDirectTranslationCS{fňukal(ka), kňoural(ka), uplakánek}
\dicEntry[væminn] \dicTerm{væminn} \dicIPA{{v}{a}{i}{\textlengthmark}{m}{\textsci}{\textsubring{n}}} \dicPos{adj}[6]\dicFlx{}[-2] \textbf{1.} \dicSynonym{tilfinningasamur} \dicDirectTranslationCS{sentimentální, dojemný}  \textbf{2.} \dicSynonym*{sem veldur klígju} \dicDirectTranslationCS{způsobující nevolnost, zvedající žaludek (jídlo ap.)}
\dicEntry[væmni] \dicTerm{væmn|i} \dicIPA{{v}{a}{i}{m}{n}{\textsci}} \dicPos{f}[3] \dicFlx{(‑i)}[3] \dicDirectTranslationCS{sentimentalita, dojemnost}
\dicEntry[væna] \dicTerm{væn|a} \dicIPA{{v}{a}{i}{\textlengthmark}{n}{a}} \dicPos{v}[2] \dicFlx{(‑di, ‑t)}[156] \dicFlx{acc} \dicSynonym{ásaka} \dicDirectTranslationCS{nařknout, obvinit, obviňovat, osočit, osočovat};  \dicPhraseIS{væna e‑n um e‑ð} \dicDirectTranslationCS{nařknout (koho) z~(čeho)} \dicExampleIS{væna hann um lygi} \dicExampleCS{nařknout ho ze lži}
\dicEntry[vænd] \dicTerm{vænd} \dicIPA{{v}{a}{i}{n}{\textsubring{d}}} \dicPos{f}[7] \dicFlx{(‑ar, ‑ir)}[1] \dicSynonym{von} \dicDirectTranslationCS{očekávání, naděje};  \dicPhraseIS{e‑að er í vændum} \dicDirectTranslationCS{(co) se očekává};  \dicPhraseIS{eiga e‑ð í vændum} \dicDirectTranslationCS{očekávat (co), vyhlížet (co)}
\dicEntry[vændi] \dicTerm{vændi} \dicIPA{{v}{a}{i}{n}{\textsubring{d}}{\textsci}} \dicPos{n}[2] \dicFlx{(‑s)}[20] \dicDirectTranslationCS{prostituce}
\dicEntry[vændishús] \dicTerm{vændis··hús} \dicIPA{{v}{a}{i}{n}{\textsubring{d}}{\textsci}{s}{h}{u}{s}} \dicPos{n}[2] \dicFlx{(‑s, ‑)}[5] \dicDirectTranslationCS{veřejný dům, nevěstinec, bordel}
\dicEntry[vændiskona] \dicTerm{vændis··kon|a} \dicIPA{{v}{a}{i}{n}{\textsubring{d}}{\textsci}{s}{k\smash{\textsuperscript{h}}}{\textopeno}{n}{a}} \dicPos{f}[1] \dicFlx{(‑u, ‑ur)}[27] \dicSynonym{gleðikona} \dicDirectTranslationCS{prostitutka}
\dicEntry[vængbrotinn] \dicTerm{væng··brotinn} \dicIPA{{v}{a}{i}{\ng}{\r{g}}{\textsubring{b}}{r}{\textopeno}{\textsubring{d}}{\textsci}{\textsubring{n}}} \dicPos{adj}[6]\dicFlx{}[-6] \textbf{1.} \dicSynonym*{með brotinn væng} \dicDirectTranslationCS{(jsoucí) s~polámanými křídly (pták ap.)}  \textbf{2.} \dicSynonym*{sem á erfitt} \dicDirectTranslationCS{bezradný, rozpačitý}
\dicEntry[vænghaf] \dicTerm{væng··|haf} \dicIPA{{v}{a}{i}{\ng}{\r{g}}{h}{a}{f}} \dicPos{n}[2] \dicFlx{(‑hafs, ‑höf)}[8] \dicDirectTranslationCS{rozpětí křídel}
\dicEntry[vængjaður] \dicTerm{vængj··|aður} \dicIPA{{v}{a}{i}{\textltailn}{\r{\textObardotlessj}}{a}{ð}{\textscy}{\textsubring{r}}} \dicPos{adj}[3] \dicFlx{(f ‑uð)}[3] \dicDirectTranslationCS{okřídlený};  \dicPhraseIS{vængjuð orð} \dicLangCat{přen.} \dicDirectTranslationCS{okřídlená slova}
\dicEntry[vængjahurð] \textls[15]{\dicTerm{vængja··hurð} \dicIPA{{v}{a}{i}{\textltailn}{\r{\textObardotlessj}}{a}{h}{\textscy}{r}{\texttheta}} \dicPos{f}[7] \dicFlx{(‑ar, ‑ir)}[1] \dicDirectTranslationCS{kyvné\,/\addthin létací dveře, lítačky}}
\dicEntry[vængur] \dicTerm{væng|ur} \dicsymFrequent\  \dicIPA{{v}{a}{i}{\ng}{\r{g}}{\textscy}{\textsubring{r}}} \dicPos{m}[9] \dicFlx{(‑s\,/\addthin ‑jar, ‑ir)}[26] \textbf{1.} \dicSynonym*{fluglimur} \dicDirectTranslationCS{křídlo, křidélko} \dicIndirectTranslationCS{(pohyblivý orgán ptáků)} \dicExampleIS{með útbreidda vængi} \dicExampleCS{s~roztaženými křídly}  \textbf{2.} \dicSynonym{kantur} \dicDirectTranslationCS{křídlo} \dicIndirectTranslationCS{(skupina lidí lišící se názorově od většiny)};  \dicIdiom{vængur}{ \dicPhraseIS{stíga í vænginn við e‑n}} \dicDirectTranslationCS{dvořit se (komu), nadbíhat (komu)}
\dicEntry[vænisýki] \dicTerm{væni··sýk|i} \dicIPA{{v}{a}{i}{\textlengthmark}{n}{\textsci}{s}{i}{\r{\textObardotlessj}}{\textsci}} \dicPos{f}[3] \dicFlx{(‑i)}[3] \dicFieldCat{med.\,/\addthin psych.} \dicSynonym{ofsóknaræði} \dicDirectTranslationCS{paranoia}
\dicEntry[vænlega] \dicTerm{væn··lega} \dicIPA{{v}{a}{i}{n}{l}{\textepsilon}{\textbabygamma}{a}} \dicPos{adv} \dicDirectTranslationCS{příznivě, slibně}
\dicEntry[vænlegur] \dicTerm{væn··legur} \dicIPA{{v}{a}{i}{n}{l}{\textepsilon}{\textbabygamma}{\textscy}{\textsubring{r}}} \dicPos{adj}[1]\dicFlx{}[-8] \dicDirectTranslationCS{příznivý, slibný} \dicExampleIS{vænleg staða} \dicExampleCS{příznivá pozice}
\dicEntry[vænn] \dicTerm{vænn} \dicsymFrequent\  \dicIPA{{v}{a}{i}{\textsubring{d}}{\textsubring{n}}} \dicPos{adj}[7]\dicFlx{}[-1] \textbf{1.} \dicSynonym{góður} \dicDirectTranslationCS{hodný, dobrý};  \dicPhraseIS{vera (hinn) vænsti maður} \dicDirectTranslationCS{být ten nejhodnější člověk};  \dicPhraseIS{viltu vera svo vænn að (gera e‑ð)} \dicDirectTranslationCS{byl bys tak hodný a~(udělal (co))}  \textbf{2.} \dicSynonym*{um stóran hluta} \dicDirectTranslationCS{pořádný, pěkný (suma ap.)}  \textbf{3.} \dicSynonym{allstór} \dicDirectTranslationCS{pořádný, náramný} \dicExampleIS{vænn kjötbiti} \dicExampleCS{pořádný kus masa};  \dicIdiom{vænn}{ \dicPhraseIS{það er ekki seinna vænna að (gera e‑ð)}} \dicDirectTranslationCS{je nejvyšší čas (udělat (co))}; { \dicPhraseIS{þykja vænt um e‑n}} \dicDirectTranslationCS{milovat (koho), mít (koho) rád}
\dicEntry[vænta] \dicTerm{vænt|a} \dicsymFrequent\  \dicIPA{{v}{a}{i}{\textsubring{n}}{\textsubring{d}}{a}} \dicPos{v}[2] \dicFlx{(‑i, ‑)}[20] \dicFlx{gen} \dicDirectTranslationCS{očekávat, předpokládat} \dicExampleIS{vænta hans á morgun} \dicExampleCS{očekávat ho zítra};  \dicIdiom{vænta}{ \dicPhraseIS{vænta sín}} \dicDirectTranslationCS{očekávat dítě} \dicExampleIS{vænta sín í næsta mánuði} \dicExampleCS{očekávat dítě příští měsíc}; { \dicPhraseIS{ekki vænti ég að (þú hafir e‑ð)}} \dicDirectTranslationCS{(nemáte) náhodou (co)} \dicIndirectTranslationCS{(zdvořilá otázka)}
\dicEntry[væntanlega] \dicTerm{væntan··lega} \dicsymFrequent\  \dicIPA{{v}{a}{i}{\textsubring{n}}{\textsubring{d}}{a}{n}{l}{\textepsilon}{\textbabygamma}{a}} \dicPos{adv} \dicSynonym{líklega} \dicDirectTranslationCS{pravděpodobně, patrně, nejspíš} \dicExampleIS{Það verður væntanlega skemmtilegt.} \dicExampleCS{Bude to pravděpodobně zajímavé.}
\dicEntry[væntanlegur] \dicTerm{væntan··legur} \dicsymFrequent\  \dicIPA{{v}{a}{i}{\textsubring{n}}{\textsubring{d}}{a}{n}{l}{\textepsilon}{\textbabygamma}{\textscy}{\textsubring{r}}} \dicPos{adj}[1]\dicFlx{}[-8] \dicDirectTranslationCS{očekávaný, předpokládaný} \dicExampleIS{Hann er væntanlegur á morgun.} \dicExampleCS{Očekává se zítra.}
\dicEntry[vænting] \dicTerm{vænt··ing} \dicIPA{{v}{a}{i}{\textsubring{n}}{\textsubring{d}}{i}{\ng}{\r{g}}} \dicPos{f}[4] \dicFlx{(‑ar, ‑ar)}[5] \dicPhraseIS{væntingar} \dicFlx{pl} \dicDirectTranslationCS{očekávání, naděje};  \dicPhraseIS{gera sér væntingar} \dicDirectTranslationCS{mít očekávání};  \dicPhraseIS{standa undir væntingum} \dicDirectTranslationCS{splnit očekávání}
\dicEntry[væntumþykja] \dicTerm{vænt·um··þykj|a} \dicIPA{{v}{a}{i}{\textsubring{n}}{\textsubring{d}}{\textscy}{m}{\texttheta}{\textsci}{\r{\textObardotlessj}}{a}} \dicPos{f}[1] \dicFlx{(‑u)}[5] \dicSynonym{elska\smash{\textsuperscript{1}}} \dicDirectTranslationCS{náklonnost, zalíbení, láska}
\dicEntry[vær] \dicTerm{vær} \dicIPA{{v}{a}{i}{\textlengthmark}{\textsubring{r}}} \dicPos{adj}[5] \dicFlx{(f ‑)}[8] \dicSynonym{kyrr} \dicDirectTranslationCS{(po)klidný, pokojný} \dicExampleIS{vær svefn} \dicExampleCS{poklidný spánek};  \dicPhraseIS{vakna af værum blundi} \dicDirectTranslationCS{probudit se ze sladkého spánku};  \dicPhraseIS{e‑m er ekki vært e‑s staðar} \dicFlx{impers} \dicDirectTranslationCS{(kdo) nemůže vydržet (kde)}
\dicEntry[værð] \dicTerm{værð} \dicIPA{{v}{a}{i}{r}{\texttheta}} \dicPos{f}[7] \dicFlx{(‑ar, ‑ir)}[1] \dicSynonym{kyrrð} \dicDirectTranslationCS{klid, pokoj} \dicExampleIS{Það sótti á hann værð.} \dicExampleCS{Zmocnil se ho klid.}
\dicEntry[værðarvoð] \dicTerm{værðar··voð} \dicIPA{{v}{a}{i}{r}{ð}{a}{r}{v}{\textopeno}{\texttheta}} \dicPos{f}[7] \dicFlx{(‑ar, ‑ir)}[1] \dicDirectTranslationCS{deka, přikrývka}
\dicEntry[væri] \dicTerm{væri} \dicIPA{{v}{a}{i}{\textlengthmark}{r}{\textsci}} \dicPos{v} \dicFlx{con pf sg 1 pers} \dicLink{vera\smash{\textsuperscript{2}}}
\dicEntry[værukær] \dicTerm{væru··kær} \dicIPA{{v}{a}{i}{\textlengthmark}{r}{\textscy}{c\smash{\textsuperscript{h}}}{a}{i}{\textsubring{r}}} \dicPos{adj}[5] \dicFlx{(f ‑)}[8] \dicDirectTranslationCS{pohodlný, lenivý}
\dicEntry[væsa] \dicTerm{væs|a} \dicIPA{{v}{a}{i}{\textlengthmark}{s}{a}} \dicPos{v}[2] \dicFlx{(‑ti, ‑t)}[84] \dicPhraseIS{það væsir ekki um e‑n} \dicFlx{impers} \dicLangCat{přen.} \dicDirectTranslationCS{(kdo) se má dobře, (kdo) si nemá na co stěžovat}
\dicEntry[væskill] \dicTerm{væsk|ill} \dicIPA{{v}{a}{i}{s}{\r{\textObardotlessj}}{\textsci}{\textsubring{d}}{\textsubring{l}}} \dicPos{m}[6] \dicFlx{(‑ils, ‑lar)}[35] \dicDirectTranslationCS{slaboch, slaboška, padavka}
\dicEntry[væta] \dicTerm{væt|a\smash{\textsuperscript{1}}} \dicIPA{{v}{a}{i}{\textlengthmark}{\textsubring{d}}{a}} \dicPos{f}[1] \dicFlx{(‑u, ‑ur)}[7] \textbf{1.} \dicSynonym{bleyta\smash{\textsuperscript{1}}} \dicDirectTranslationCS{vlhko, mokro}  \textbf{2.} \dicSynonym{rigning} \dicDirectTranslationCS{mrholení, mžení}
\dicEntry[væta] \dicTerm{væt|a\smash{\textsuperscript{2}}} \dicIPA{{v}{a}{i}{\textlengthmark}{\textsubring{d}}{a}} \dicPos{v}[2] \dicFlx{(‑ti, ‑t)}[56] \dicFlx{acc} \dicDirectTranslationCS{namočit, namáčet, navlhčit, navlhčovat, zvlhčit, zvlhčovat} \dicExampleIS{væta klútinn í volgu vatni} \dicExampleCS{namočit hadřík ve vlažné vodě};  \dicIdiom{væta}{ \dicPhraseIS{væta rúmið}} \dicDirectTranslationCS{počůrat se v~posteli};  \dicPhraseIS{það vætir} \dicFlx{impers} \dicDirectTranslationCS{mrholí, mží}
\dicEntry[vætla] \dicTerm{vætl|a} \dicIPA{{v}{a}{i}{h}{\textsubring{d}}{l}{a}} \dicPos{v}[1] \dicFlx{(‑aði)}[75] \dicSynonym{drjúpa} \dicDirectTranslationCS{kapat, odkapávat, (pomalu) stékat} \dicExampleIS{Það vætlar úr krananum.} \dicExampleCS{Z~kohoutku kape.}
\dicEntry[vættur] \dicTerm{vætt|ur\smash{\textsuperscript{1}}} \dicIPA{{v}{a}{i}{h}{\textsubring{d}}{\textscy}{\textsubring{r}}} \dicPos{m}[10] \dicFlx{(‑ar, ‑ir)}[4] \dicLink{vættur\smash{\textsuperscript{2}}}
\dicEntry[vættur] \dicTerm{vætt|ur\smash{\textsuperscript{2}}}\dicTerm{, vættur\smash{\textsuperscript{1}}} \dicIPA{{v}\-{a}\-{i}\-{h}\-{\textsubring{d}}\-{\textscy}\-{\textsubring{r}}\-} \dicPos{f}[7] \dicFlx{(‑ar, ‑ir)}[1] \dicDirectTranslationCS{duch, nadpřirozená bytost} \dicExampleIS{vættir í fjöllum} \dicExampleCS{nadpřirozené bytosti z~hor}
\dicEntry[vöð] \dicTerm{vöð} \dicIPA{{v}{\oe}{\textlengthmark}{\texttheta}} \dicPos{n} \dicFlx{pl nom} \dicLink{vað}
\dicEntry[vöðla] \dicTerm{vöðl|a} \dicIPA{{v}{\oe}{ð}{l}{a}} \dicPos{v}[1] \dicFlx{(‑aði)}[1] \dicFlx{dat} \dicSynonym*{hnoða saman} \dicDirectTranslationCS{(z)muchlat, (z)mačkat} \dicExampleIS{vöðla saman fötum} \dicExampleCS{zmuchlat šaty dohromady}
\dicEntry[vöðu] \dicTerm{vöðu} \dicIPA{{v}{\oe}{\textlengthmark}{ð}{\textscy}} \dicPos{f} \dicFlx{sg gen} \dicLink{vaða\smash{\textsuperscript{1}}}
\dicEntry[vöðuselur] \dicTerm{vöðu··sel|ur} \dicIPA{{v}{\oe}{\textlengthmark}{ð}{\textscy}{s}{\textepsilon}{l}{\textscy}{\textsubring{r}}} \dicPos{m}[9] \dicFlx{(‑s, ‑ir)}[8] \dicFieldCat{zool.} \dicDirectTranslationCS{tuleň grónský} \textit{(l.~{\textLA{Pagophilus groenlandicus}})}  \dicsymPhoto\ 
\dicFigure{ds_image_voduselur_0_2.jpg}{Vöðuselur}{Vöðuselur - Matthieu Godbout, GFDL}
\dicEntry[vöðvakrampi] \dicTerm{vöðva··kramp|i} \dicIPA{{v}{\oe}{ð}{v}{a}{k\smash{\textsuperscript{h}}}{r}{a}{\textsubring{m}}{\textsubring{b}}{\textsci}} \dicPos{m}[1] \dicFlx{(‑a, ‑ar)}[8] \dicDirectTranslationCS{svalová křeč}
\dicEntry[vöðvamikill] \dicTerm{vöðva··|mikill} \dicIPA{{v}{\oe}{ð}{v}{a}{m}{\textsci}{\r{\textObardotlessj}}{\textsci}{\textsubring{d}}{\textsubring{l}}} \dicPos{adj}[11] \dicFlx{(comp ‑meiri, sup ‑mestur)}[12] \dicDirectTranslationCS{svalnatý}
\dicEntry[vöðvarýrnun] \dicTerm{vöðva··rýrn|un} \dicIPA{{v}{\oe}{ð}{v}{a}{r}{i}{r}{\textsubring{d}}{n}{\textscy}{\textsubring{n}}} \dicPos{f}[7] \dicFlx{(‑unar)}[9] \dicFieldCat{med. } \dicDirectTranslationCS{svalová dystrofie}
\dicEntry[vöðvi] \dicTerm{vöðv|i} \dicsymFrequent\  \dicIPA{{v}{\oe}{ð}{v}{\textsci}} \dicPos{m}[1] \dicFlx{(‑a, ‑ar)}[1] \dicFieldCat{anat.} \dicDirectTranslationCS{sval} \dicExampleIS{Þessi æfing styrkir vöðvana.} \dicExampleCS{Toto cvičení posiluje svaly.}
\dicEntry[vöf] \dicTerm{vöf} \dicIPA{{v}{\oe}{\textlengthmark}{f}} \dicPos{n} \dicFlx{pl nom} \dicLink{vaf}
\dicEntry[vöfðum] \dicTerm{vöfðum} \dicIPA{{v}{\oe}{v}{ð}{\textscy}{\textsubring{m}}} \dicPos{v} \dicFlx{ind pf pl 1 pers} \dicLink{vefja}
\dicEntry[vöff] \dicTerm{vöff} \dicIPA{{v}{\oe}{f}{\textlengthmark}} \dicPos{n} \dicFlx{pl nom} \dicLink{vaff}
\dicEntry[vöfflu] \dicTerm{vöfflu} \dicIPA{{v}{\oe}{f}{l}{\textscy}} \dicPos{f} \dicFlx{sg gen} \dicLink{vaffla}
\dicEntry[vöfflujárn] \dicTerm{vöfflu··járn} \dicIPA{{v}{\oe}{f}{l}{\textscy}{j}{au}{r}{\textsubring{d}}{\textsubring{n}}} \dicPos{n}[2] \dicFlx{(‑s, ‑)}[5] \dicDirectTranslationCS{vaflovač}
\dicEntry[vöflur] \dicTerm{vöflur} \dicIPA{{v}{\oe}{\textsubring{b}}{l}{\textscy}{\textsubring{r}}} \dicPos{f}[12] \dicFlx{pl}[4] \dicSynonym{vandræði} \dicDirectTranslationCS{rozpaky, otálení};  \dicPhraseIS{hafa engar vöflur á e‑u} \dicDirectTranslationCS{nerozpakovat se nad (čím)}
\dicEntry[vöggu] \dicTerm{vöggu} \dicIPA{{v}{\oe}{\r{g}}{\textlengthmark}{\textscy}} \dicPos{f} \dicFlx{sg gen} \dicLink{vagga\smash{\textsuperscript{1}}}
\dicEntry[vöggugjöf] \dicTerm{vöggu··|gjöf} \dicIPA{{v}{\oe}{\r{g}}{\textlengthmark}{\textscy}{\r{\textObardotlessj}}{\oe}{f}} \dicPos{f}[7] \dicFlx{(‑gjafar, ‑gjafir)}[16] \dicPhraseIS{fá e‑ð í vöggugjöf} \dicLangCat{přen.} \dicDirectTranslationCS{dostat (co) do vínku}
\dicEntry[vögguprent] \dicTerm{vöggu··prent} \dicIPA{{v}{\oe}{\r{g}}{\textlengthmark}{\textscy}{p\smash{\textsuperscript{h}}}{r}{\textepsilon}{\textsubring{n}}{\textsubring{d}}} \dicPos{n}[2] \dicFlx{(‑s)}[2] \dicDirectTranslationCS{inkunábule, prvotisk}
\dicEntry[vöggustofa] \dicTerm{vöggu··stof|a} \dicIPA{{v}{\oe}{\r{g}}{\textlengthmark}{\textscy}{s}{\textsubring{d}}{\textopeno}{v}{a}} \dicPos{f}[1] \dicFlx{(‑u, ‑ur)}[7] \dicFieldCat{škol.} \dicDirectTranslationCS{jesle}
\dicEntry[vögguvísa] \dicTerm{vöggu··vís|a} \dicIPA{{v}{\oe}{\r{g}}{\textlengthmark}{\textscy}{v}{i}{s}{a}} \dicPos{f}[1] \dicFlx{(‑u, ‑ur)}[13] \dicDirectTranslationCS{ukolébavka, uspávanka}
\dicEntry[vök] \dicTerm{vök} \dicIPA{{v}{\oe}{\textlengthmark}{\r{g}}} \dicPos{f}[7] \dicFlx{(vakar, vakir)}[16] \dicDirectTranslationCS{díra v~ledu};  \dicPhraseIS{eiga í vök að verjast} \dicLangCat{přen.} \dicDirectTranslationCS{být v~prekérní situaci}
\dicEntry[vökna] \dicTerm{vökn|a} \dicIPA{{v}{\oe}{h}{\r{g}}{n}{a}} \dicPos{v}[1] \dicFlx{(‑aði)}[37] \dicSynonym{blotna} \dicDirectTranslationCS{(z)vlhnout, namoknout} \dicExampleIS{Hann vöknaði í fæturna.} \dicExampleCS{Promokly mu nohy.};  \dicPhraseIS{e‑m vöknar um augu} \dicFlx{impers} \dicDirectTranslationCS{(komu) vlhnou oči}
\dicEntry[vöknuð] \dicTerm{vöknuð} \dicIPA{{v}{\oe}{h}{\r{g}}{n}{\textscy}{\texttheta}} \dicPos{adj} \dicFlx{f nom\,/\addthin acc} \dicLink{vaknaður}
\dicEntry[vöktum] \dicTerm{vöktum} \dicIPA{{v}{\oe}{x}{\textsubring{d}}{\textscy}{\textsubring{m}}} \dicPos{v} \dicFlx{ind pf pl 1 pers} \dicLink{vekja}
\dicEntry[vöktun] \dicTerm{vökt|un} \dicIPA{{v}{\oe}{x}{\textsubring{d}}{\textscy}{\textsubring{n}}} \dicPos{f}[7] \dicFlx{(‑unar)}[12] \dicDirectTranslationCS{monitoring, monitorování, sledování}
\dicEntry[vökudeild] \dicTerm{vöku··deild} \dicIPA{{v}{\oe}{\textlengthmark}{\r{g}}{\textscy}{\textsubring{d}}{ei}{l}{\textsubring{d}}} \dicPos{f}[7] \dicFlx{(‑ar, ‑ir)}[1] \dicDirectTranslationCS{novorozenecké oddělení}
\dicEntry[vökull] \dicTerm{vökull} \dicIPA{{v}{\oe}{\textlengthmark}{\r{g}}{\textscy}{\textsubring{d}}{\textsubring{l}}} \dicPos{adj}[8]\dicFlx{}[-4] \dicDirectTranslationCS{bdělý, ostražitý}
\dicEntry[vökva] \dicTerm{vökv|a} \dicIPA{{v}{\oe}{\textlengthmark}{\r{g}}{v}{a}} \dicPos{v}[1] \dicFlx{(‑aði)}[1] \dicFlx{acc} \dicDirectTranslationCS{zalít, zalévat, zavlažit, zavlažovat} \dicExampleIS{vökva garðinn} \dicExampleCS{zalévat zahradu}
\dicEntry[vökvakerfi] \dicTerm{vökva··kerfi} \dicIPA{{v}{\oe}{\textlengthmark}{\r{g}}{v}{a}{c\smash{\textsuperscript{h}}}{\textepsilon}{r}{v}{\textsci}} \dicPos{n}[2] \dicFlx{(‑s, ‑)}[14] \dicDirectTranslationCS{hydraulický systém}
\dicEntry[vökvalyfta] \dicTerm{vökva··lyft|a} \dicIPA{{v}{\oe}{\textlengthmark}{\r{g}}{v}{a}{l}{\textsci}{f}{\textsubring{d}}{a}} \dicPos{f}[1] \dicFlx{(‑u, ‑ur)}[7] \dicDirectTranslationCS{hydraulický zvedák, hever}
\dicEntry[vökvastýri] \dicTerm{vökva··stýri} \dicIPA{{v}{\oe}{\textlengthmark}{\r{g}}{v}{a}{s}{\textsubring{d}}{i}{r}{\textsci}} \dicPos{n}[2] \dicFlx{(‑s, ‑)}[14] \dicSynonym{aflstýri} \dicDirectTranslationCS{posilovač řízení}

\dicEntry[vökvaþrýstingur] \dicTerm{vökva··þrýst·ing|ur} \dicIPA{{v}{\oe}{\textlengthmark}{\r{g}}{v}{a}{\texttheta}{r}{i}{s}{\textsubring{d}}{i}{\ng}{\r{g}}{\textscy}{\textsubring{r}}} \dicPos{m}[6] \dicFlx{(‑s)}[9] \dicFieldCat{fyz.} \dicDirectTranslationCS{vod\-ní\,/\addthin hydraulický tlak}

\dicEntry[vökvi] \dicTerm{vökv|i} \dicsymFrequent\  \dicIPA{{v}{\oe}{\textlengthmark}{\r{g}}{v}{\textsci}} \dicPos{m}[1] \dicFlx{(‑a, ‑ar)}[1] \dicDirectTranslationCS{tekutina, kapalina} \dicExampleIS{glær vökvi} \dicExampleCS{průzračná tekutina}
\dicEntry[vökvun] \dicTerm{vökv|un} \dicIPA{{v}{\oe}{\textlengthmark}{\r{g}}{v}{\textscy}{\textsubring{n}}} \dicPos{f}[7] \dicFlx{(‑unar)}[12] \dicDirectTranslationCS{zalití, zalévání, zavlažování}
\dicEntry[völ] \dicTerm{völ\smash{\textsuperscript{1}}} \dicIPA{{v}{\oe}{\textlengthmark}{\textsubring{l}}} \dicPos{f}[7] \dicFlx{(valar)}[19] \dicSynonym{val} \dicDirectTranslationCS{výběr, volba} \dicExampleIS{eiga völ á e‑u} \dicExampleCS{mít na výběr (co)};  \dicPhraseIS{eiga engra kosta völ} \dicLangCat{přen.} \dicDirectTranslationCS{nemít nic jiného na výběr};  \dicProverb\  \dicPhraseIS{Sá á kvölina sem á völina.} \dicLangCat{přís.} \dicDirectTranslationCS{Bývá obtížné si vybrat.}
\dicEntry[völ] \dicTerm{völ\smash{\textsuperscript{2}}} \dicIPA{{v}{\oe}{\textlengthmark}{\textsubring{l}}} \dicPos{n} \dicFlx{pl nom} \dicLink{val}
\dicEntry[völd] \dicTerm{völd\smash{\textsuperscript{1}}} \dicIPA{{v}{\oe}{l}{\textsubring{d}}} \dicPos{n} \dicFlx{pl nom} \dicLink{vald}
\dicEntry[völd] \dicTerm{völd\smash{\textsuperscript{2}}} \dicIPA{{v}{\oe}{l}{\textsubring{d}}} \dicPos{adj} \dicFlx{f sg nom pos} \dicLink{valdur}
\dicEntry[völdum] \dicTerm{völdum} \dicIPA{{v}{\oe}{l}{\textsubring{d}}{\textscy}{\textsubring{m}}} \dicPos{v} \dicFlx{ind pf pl 1 pers} \dicLink{velja}
\dicEntry[völlur] \dicTerm{völlur} \dicsymFrequent\  \dicIPA{{v}{\oe}{\textsubring{d}}{l}{\textscy}{\textsubring{r}}} \dicPos{m}[11] \dicFlx{(vallar, vellir)}[5] \textbf{1.} \dicSynonym{flöt} \dicDirectTranslationCS{pláň, planina}  \textbf{2.} \dicSynonym{íþróttavöllur} \dicDirectTranslationCS{stadion, hřiště}  \textbf{3.} \dicSynonym{flugvöllur} \dicDirectTranslationCS{letiště};  \dicPhraseIS{fara út á völl} \dicDirectTranslationCS{jet na letiště};  \dicIdiom{völlur}{ \dicPhraseIS{e‑að heldur velli}} \dicDirectTranslationCS{(co) vytrvá, (co) vydrží (vláda ap.)}; { \dicPhraseIS{leggja e‑n að velli}} \dicDirectTranslationCS{porazit (koho)};  \dicPhraseIS{þéttur að velli} \dicFlx{adj} \dicDirectTranslationCS{statný, silný}
\dicEntry[völt] \dicTerm{völt} \dicIPA{{v}{\oe}{\textsubring{l}}{\textsubring{d}}} \dicPos{adj} \dicFlx{f sg nom pos} \dicLink{valtur}
\dicEntry[völundarhús] \dicTerm{völundar··hús} \dicIPA{{v}{\oe}{\textlengthmark}{l}{\textscy}{n}{\textsubring{d}}{a}{\textsubring{r}}{h}{u}{s}} \dicPos{n}[2] \dicFlx{(‑s, ‑)}[5] \textbf{1.} \dicDirectTranslationCS{bludiště, labyrint}  \textbf{2.} \dicFieldCat{anat.} \dicDirectTranslationCS{labyrint} \dicIndirectTranslationCS{(část vnitřního ucha)}
\dicEntry[völundur] \dicTerm{völund|ur} \dicIPA{{v}{\oe}{\textlengthmark}{l}{\textscy}{n}{\textsubring{d}}{\textscy}{\textsubring{r}}} \dicPos{m}[6] \dicFlx{(‑s\,/\addthin ‑ar, ‑ar)}[63] \dicDirectTranslationCS{geniální řemeslník\,/\addthin řemeslnice}
\dicEntry[Völuspá] \dicTerm{Völu··spá} \dicIPA{{v}{\oe}{\textlengthmark}{l}{\textscy}{s}{\textsubring{b}}{au}} \dicPos{f}[4] \dicFlx{(‑r)}[20] \dicFieldCat{lit.} \dicDirectTranslationCS{Vědmina píseň} \dicIndirectTranslationCS{(část Poetické Eddy)}
\dicEntry[völva] \dicTerm{völv|a} \dicIPA{{v}{\oe}{l}{v}{a}} \dicPos{f}[1] \dicFlx{(‑u, ‑ur)}[7] \dicSynonym*{forspá kona} \dicDirectTranslationCS{věštkyně, prorokyně}
\dicEntry[vömb] \dicTerm{vömb} \dicIPA{{v}{\oe}{m}{\textsubring{b}}} \dicPos{f}[7] \dicFlx{(vambar, vambir)}[16] \textbf{1.} \dicSynonym{magi} \dicDirectTranslationCS{bachor, panděro} \dicIndirectTranslationCS{(břicho)};  \dicPhraseIS{kýla vömbina} \dicLangCat{přen.} \dicDirectTranslationCS{nacpat si panděro};  \dicPhraseIS{dæma eins og vömbin veltist} \dicLangCat{přen.} \dicDirectTranslationCS{střílet od oka} \dicIndirectTranslationCS{(hádat)}  \textbf{2.} \dicFieldCat{zool.} \dicDirectTranslationCS{bachor} \dicIndirectTranslationCS{(první část složeného žaludku přežvýkavců)}
\dicEntry[vön] \dicTerm{vön} \dicIPA{{v}{\oe}{\textlengthmark}{\textsubring{n}}} \dicPos{adj} \dicFlx{f sg nom pos} \dicLink{vanur\smash{\textsuperscript{2}}}
\dicEntry[vönd] \dicTerm{vönd} \dicIPA{{v}{\oe}{n}{\textsubring{d}}} \dicPos{adj} \dicFlx{f sg nom pos} \dicLink{vandur}
\dicEntry[vönduð] \dicTerm{vönduð} \dicIPA{{v}{\oe}{n}{\textsubring{d}}{\textscy}{\texttheta}} \dicPos{adj} \dicFlx{f sg nom pos} \dicLink{vandaður}
\dicEntry[vöndull] \dicTerm{vönd|ull} \dicIPA{{v}{\oe}{n}{\textsubring{d}}{\textscy}{\textsubring{d}}{\textsubring{l}}} \dicPos{m}[6] \dicFlx{(‑uls, ‑lar)}[31] \dicDirectTranslationCS{balík (knih ap.)}
\dicEntry[vöndum] \dicTerm{vöndum} \dicIPA{{v}{\oe}{n}{\textsubring{d}}{\textscy}{\textsubring{m}}} \dicPos{v} \dicFlx{ind pf pl 1 pers} \dicLink{venja}
\dicEntry[vöndun] \dicTerm{vönd|un} \dicIPA{{v}{\oe}{n}{\textsubring{d}}{\textscy}{\textsubring{n}}} \dicPos{f}[7] \dicFlx{(‑unar)}[12] \dicDirectTranslationCS{pečlivost, preciznost}
\dicEntry[vöndur] \dicTerm{vöndur} \dicIPA{{v}{\oe}{n}{\textsubring{d}}{\textscy}{\textsubring{r}}} \dicPos{m}[11] \dicFlx{(vandar, vendir)}[5] \textbf{1.} \dicSynonym{sópur} \dicDirectTranslationCS{metla, koště}  \textbf{2.} \dicSynonym*{greinaknippi} \dicDirectTranslationCS{metla, rákoska} \dicExampleIS{hýða börnin með vendi} \dicExampleCS{dát dětem na zadek rákoskou}  \textbf{3.} \dicSynonym{blómvöndur} \dicDirectTranslationCS{kytice, pugét}
\dicEntry[vönkuð] \dicTerm{vönkuð} \dicIPA{{v}{\oe}{\r{\ng}}{\r{g}}{\textscy}{\texttheta}} \dicPos{adj} \dicFlx{f sg nom pos} \dicLink{vankaður}
\dicEntry[vöntun] \dicTerm{vönt|un} \dicIPA{{v}{\oe}{\textsubring{n}}{\textsubring{d}}{\textscy}{\textsubring{n}}} \dicPos{f}[7] \dicFlx{(‑unar)}[12] \dicSynonym{skortur} \dicDirectTranslationCS{nedostatek, deficit} \dicExampleIS{Það er vöntun á hráefni.} \dicExampleCS{Je nedostatek surovin.}
\dicEntry[vönuð] \dicTerm{vönuð} \dicIPA{{v}{\oe}{\textlengthmark}{n}{\textscy}{\texttheta}} \dicPos{adj} \dicFlx{f nom\,/\addthin acc} \dicLink{vanaður}
\dicEntry[vönun] \dicTerm{vönun} \dicIPA{{v}{\oe}{\textlengthmark}{n}{\textscy}{\textsubring{n}}} \dicPos{f}[7] \dicFlx{(‑ar, vananir)}[16] \dicSynonym{gelding} \dicDirectTranslationCS{kastrace}
\dicEntry[vör] \dicTerm{vör\smash{\textsuperscript{1}}} \dicsymFrequent\  \dicIPA{{v}{\oe}{\textlengthmark}{\textsubring{r}}} \dicPos{f}[7] \dicFlx{(varar, varir)}[16] \textbf{1.} \dicDirectTranslationCS{ret};  \dicPhraseIS{efri vör} \dicDirectTranslationCS{horní ret};  \dicPhraseIS{neðri vör} \dicDirectTranslationCS{dolní ret}  \textbf{2.} \dicDirectTranslationCS{místo přistání (lodi ap.)};  \dicIdiom{vör}{ \dicPhraseIS{e‑að er á allra vörum}} \dicDirectTranslationCS{o~(čem) všichni mluví}; { \dicPhraseIS{e‑að er komið fram á varirnar á e‑m}} \dicLangCat{přen.} \dicDirectTranslationCS{(co komu) přichází na jazyk}
\dicEntry[vör] \dicTerm{vör\smash{\textsuperscript{2}}} \dicIPA{{v}{\oe}{\textlengthmark}{\textsubring{r}}} \dicPos{adj} \dicFlx{f sg nom pos} \dicLink{var\smash{\textsuperscript{1}}}
\dicEntry[vörðum] \dicTerm{vörðum} \dicIPA{{v}{\oe}{r}{ð}{\textscy}{\textsubring{m}}} \dicPos{v} \dicFlx{ind pf pl 1 pers} \textbf{1.} \dicLink{verja\smash{\textsuperscript{2}}}  \textbf{2.} \dicLink{verja\smash{\textsuperscript{3}}}
\dicEntry[vörður] \dicTerm{vörður} \dicsymFrequent\  \dicIPA{{v}{\oe}{r}{ð}{\textscy}{\textsubring{r}}} \dicPos{m}[11] \dicFlx{(varðar, verðir)}[5] \dicSynonym{gæslumaður} \dicDirectTranslationCS{stráž, strážce, strážkyně, strážný, strážná, hlídač(ka)};  \dicPhraseIS{halda vörð um e‑ð} \dicDirectTranslationCS{hlídat (co), strážit (co)};  \dicPhraseIS{standa vörð um e‑ð} \dicLangCat{přen.} \dicDirectTranslationCS{chránit (co), střežit (co)};  \dicPhraseIS{standa vörð} \dicDirectTranslationCS{stát na stráži};  \dicPhraseIS{vera á  verði} \dicDirectTranslationCS{hlídat, držet hlídku\,/\addthin stráž};  \dicPhraseIS{verðir laganna} \dicDirectTranslationCS{strážci práva}
\dicEntry[vörm] \dicTerm{vörm} \dicIPA{{v}{\oe}{r}{\textsubring{m}}} \dicPos{adj} \dicFlx{f sg nom pos} \dicLink{varmur}
\dicEntry[vörn] \dicTerm{vörn} \dicsymFrequent\  \dicIPA{{v}{\oe}{r}{\textsubring{d}}{\textsubring{n}}} \dicPos{f}[7] \dicFlx{(varnar, varnir)}[16] \textbf{1.} \dicSynonym*{það að verjast} \dicDirectTranslationCS{obrana, bránění, defenziva};  \dicPhraseIS{snúa vörn í sókn} \dicLangCat{přen.} \dicDirectTranslationCS{přejít z~obrany do útoku};  \dicPhraseIS{snúast til varnar fyrir e‑n} \dicFlx{refl} \dicDirectTranslationCS{bránit (koho), zastávat se (koho)}  \textbf{2.} \dicFieldCat{sport.} \dicDirectTranslationCS{obrana} \dicAntonym{sókn}  \textbf{3.} \dicDirectTranslationCS{ochrana (před slunečními paprsky ap.)}  \textbf{4.} \dicFieldCat{práv.} \dicDirectTranslationCS{obhajoba}
\dicEntry[vörpu] \dicTerm{vörpu} \dicIPA{{v}{\oe}{\textsubring{r}}{\textsubring{b}}{\textscy}} \dicPos{f} \dicFlx{sg gen} \dicLink{varpa\smash{\textsuperscript{1}}}
\dicEntry[vörpun] \dicTerm{vörpun} \dicIPA{{v}{\oe}{\textsubring{r}}{\textsubring{b}}{\textscy}{\textsubring{n}}} \dicPos{f}[7] \dicFlx{(‑ar, varpanir)}[16] \textbf{1.} \dicDirectTranslationCS{vrh, vrhnutí}  \textbf{2.} \dicFieldCat{mat.} \dicDirectTranslationCS{zobrazení}
\dicEntry[vörslu] \dicTerm{vörslu} \dicIPA{{v}{\oe}{\textsubring{r}}{s}{\textsubring{d}}{l}{\textscy}} \dicPos{f} \dicFlx{sg gen} \dicLink{varsla}
\dicEntry[vörtu] \dicTerm{vörtu} \dicIPA{{v}{\oe}{\textsubring{r}}{\textsubring{d}}{\textscy}} \dicPos{f} \dicFlx{sg gen} \dicLink{varta}
\dicEntry[vörtusvín] \dicTerm{vörtu··svín} \dicIPA{{v}{\oe}{\textsubring{r}}{\textsubring{d}}{\textscy}{s}{v}{i}{\textsubring{n}}} \dicPos{n}[2] \dicFlx{(‑s, ‑)}[5] \dicFieldCat{zool.} \dicDirectTranslationCS{prase savanové} \textit{(l.~{\textLA{Phacochoerus africanus}})}  \dicsymPhoto\ 
\dicFigure{ds_image_vortusvin_0_1.jpg}{Vörtusvín}{Vörtusvín - Ikiwaner, GFDL 1.2}
\dicEntry[vöru] \dicTerm{vöru} \dicIPA{{v}{\oe}{\textlengthmark}{r}{\textscy}} \dicPos{f} \dicFlx{sg gen} \dicLink{vara\smash{\textsuperscript{1}}}
\dicEntry[vörubifreið] \dicTerm{vöru··bif·reið} \dicIPA{{v}{\oe}{\textlengthmark}{r}{\textscy}{\textsubring{b}}{\textsci}{v}{r}{ei}{\texttheta}} \dicPos{f}[4] \dicFlx{(‑ar, ‑ar\,/\addthin ‑ir)}[30] \dicLangCat{form.} \dicSynonym{vörubíll} \dicDirectTranslationCS{nákladní auto}
\dicEntry[vörubirgðir] \dicTerm{vöru··birgðir} \dicIPA{{v}{\oe}{\textlengthmark}{r}{\textscy}{\textsubring{b}}{\textsci}{r}{ð}{\textsci}{\textsubring{r}}} \dicPos{f}[12] \dicFlx{pl}[10] \dicDirectTranslationCS{sklad, zásoby (zboží)}
\dicEntry[vörubíll] \dicTerm{vöru··bíl|l}\dicTerm{, vörubifreið} \dicIPA{{v}\-{\oe}\-{\textlengthmark}\-{r}\-{\textscy}\-{\textsubring{b}}\-{i}\-{\textsubring{d}}\-{\textsubring{l}}\-} \dicPos{m}[6] \dicFlx{(‑s, ‑ar)}[48] \dicDirectTranslationCS{nákladní auto}
\dicEntry[vörubílstjóri] \dicTerm{vöru·bíl··stjór|i} \dicIPA{{v}{\oe}{\textlengthmark}{r}{\textscy}{\textsubring{b}}{i}{l}{s}{\textsubring{d}}{j}{ou}{r}{\textsci}} \dicPos{m}[1] \dicFlx{(‑a, ‑ar)}[1] \dicDirectTranslationCS{kamionista, kamionistka, řidič(ka) kamionu\,/\addthin nákladního auta}
\dicEntry[vörubretti] \dicTerm{vöru··bretti} \dicIPA{{v}{\oe}{\textlengthmark}{r}{\textscy}{\textsubring{b}}{r}{\textepsilon}{h}{\textsubring{d}}{\textsci}} \dicPos{n}[2] \dicFlx{(‑s, ‑)}[14] \dicSynonym{bretti} \dicDirectTranslationCS{(transportní) paleta}
\dicEntry[vörudreifing] \dicTerm{vöru··dreif·ing} \dicIPA{{v}{\oe}{\textlengthmark}{r}{\textscy}{\textsubring{d}}{r}{ei}{v}{i}{\ng}{\r{g}}} \dicPos{f}[4] \dicFlx{(‑ar)}[7] \dicDirectTranslationCS{distribuce}
\dicEntry[vöruflokkur] \dicTerm{vöru··flokk|ur} \dicIPA{{v}{\oe}{\textlengthmark}{r}{\textscy}{f}{l}{\textopeno}{h}{\r{g}}{\textscy}{\textsubring{r}}} \dicPos{m}[6] \dicFlx{(‑s, ‑ar)}[8] \dicDirectTranslationCS{druh zboží}
\dicEntry[vöruflutningabíll] \dicTerm{vöru·flutninga··bíl|l} \dicIPA{{v}\-{\oe}\-{\textlengthmark}\-{r}\-{\textscy}\-{f}\-{l}\-{\textscy}\-{h}\-{\textsubring{d}}\-{n}\-{i}\-{\ng}\-{\r{g}}\-{a}\-{\textsubring{b}}\-{i}\-{\textsubring{d}}\-{\textsubring{l}}\-} \dicPos{m}[6] \dicFlx{(‑s, ‑ar)}[48] \dicDirectTranslationCS{kamion, nákladní auto}
\dicEntry[vöruflutningar] \dicTerm{vöru··flut·ningar} \dicIPA{{v}{\oe}{\textlengthmark}{r}{\textscy}{f}{l}{\textscy}{h}{\textsubring{d}}{n}{i}{\ng}{\r{g}}{a}{\textsubring{r}}} \dicPos{m}[6] \dicFlx{pl}[2] \dicDirectTranslationCS{přeprava\,/\addthin transport zboží}
\dicEntry[vöruflutningavél] \dicTerm{vöru·flutninga··vél} \dicIPA{{v}{\oe}{\textlengthmark}{r}{\textscy}{f}{l}{\textscy}{h}{\textsubring{d}}{n}{i}{\ng}{\r{g}}{a}{v}{j}{\textepsilon}{\textsubring{l}}} \dicPos{f}[4] \dicFlx{(‑ar, ‑ar)}[1] \dicDirectTranslationCS{nákladní letadlo, letadlo přepravující zboží}
\dicEntry[vörugeymsla] \dicTerm{vöru··geymsl|a} \dicIPA{{v}{\oe}{\textlengthmark}{r}{\textscy}{\r{\textObardotlessj}}{ei}{m}{s}{\textsubring{d}}{l}{a}} \dicPos{f}[1] \dicFlx{(‑u, ‑ur)}[13] \dicDirectTranslationCS{skladiště, sklad zboží}
\dicEntry[vörugjald] \dicTerm{vöru··|gjald} \dicIPA{{v}{\oe}{\textlengthmark}{r}{\textscy}{\r{\textObardotlessj}}{a}{l}{\textsubring{d}}} \dicPos{n}[2] \dicFlx{(‑gjalds, ‑gjöld)}[8] \dicFieldCat{ekon.} \dicDirectTranslationCS{spotřební daň}
\dicEntry[vörugæði] \dicTerm{vöru··gæði} \dicIPA{{v}{\oe}{\textlengthmark}{r}{\textscy}{\r{\textObardotlessj}}{a}{i}{ð}{\textsci}} \dicPos{n}[2] \dicFlx{pl}[19] \dicDirectTranslationCS{kvalita\,/\addthin jakost zboží}
\dicEntry[vöruheiti] \dicTerm{vöru··heiti} \dicIPA{{v}{\oe}{\textlengthmark}{r}{\textscy}{h}{ei}{\textsubring{d}}{\textsci}} \dicPos{n}[2] \dicFlx{(‑s, ‑)}[14] \dicDirectTranslationCS{název zboží, jméno výrobku}
\dicEntry[vöruhús] \dicTerm{vöru··hús} \dicIPA{{v}{\oe}{\textlengthmark}{r}{\textscy}{h}{u}{s}} \dicPos{n}[2] \dicFlx{(‑s, ‑)}[5] \textbf{1.} \dicSynonym{lager} \dicDirectTranslationCS{sklad, skladiště}  \textbf{2.} \dicSynonym*{stór verslun} \dicDirectTranslationCS{velkoobchod, hypermarket}
\dicEntry[vörulán] \dicTerm{vöru··lán} \dicIPA{{v}{\oe}{\textlengthmark}{r}{\textscy}{l}{au}{\textsubring{n}}} \dicPos{n}[2] \dicFlx{(‑s, ‑)}[5] \dicFieldCat{ekon.} \dicDirectTranslationCS{dodavatelský úvěr}
\dicEntry[vörulisti] \dicTerm{vöru··list|i} \dicIPA{{v}{\oe}{\textlengthmark}{r}{\textscy}{l}{\textsci}{s}{\textsubring{d}}{\textsci}} \dicPos{m}[1] \dicFlx{(‑a, ‑ar)}[1] \dicDirectTranslationCS{katalog zboží}
\dicEntry[vörulyfta] \dicTerm{vöru··lyft|a} \dicIPA{{v}{\oe}{\textlengthmark}{r}{\textscy}{l}{\textsci}{f}{\textsubring{d}}{a}} \dicPos{f}[1] \dicFlx{(‑u, ‑ur)}[7] \dicDirectTranslationCS{nákladní výtah}
\dicEntry[vörumerki] \dicTerm{vöru··merki} \dicIPA{{v}{\oe}{\textlengthmark}{r}{\textscy}{m}{\textepsilon}{\textsubring{r}}{\r{\textObardotlessj}}{\textsci}} \dicPos{n}[2] \dicFlx{(‑s, ‑)}[16] \dicDirectTranslationCS{obchodní značka, ochranná známka} \dicExampleIS{skrásett vörumerki} \dicExampleCS{registrovaná obchodní značka}
\dicEntry[vörumiði] \dicTerm{vöru··mið|i} \dicIPA{{v}{\oe}{\textlengthmark}{r}{\textscy}{m}{\textsci}{ð}{\textsci}} \dicPos{m}[1] \dicFlx{(‑a, ‑ar)}[1] \dicDirectTranslationCS{cenovka}
\dicEntry[vörunúmer] \dicTerm{vöru··númer} \dicIPA{{v}{\oe}{\textlengthmark}{r}{\textscy}{n}{u}{m}{\textepsilon}{\textsubring{r}}} \dicPos{n}[2] \dicFlx{(‑s, ‑)}[5] \dicDirectTranslationCS{číslo výrobku\,/\addthin zboží}
\dicEntry[vörureikningur] \dicTerm{vöru··reikn·ing|ur} \dicIPA{{v}{\oe}{\textlengthmark}{r}{\textscy}{r}{ei}{h}{\r{g}}{n}{i}{\ng}{\r{g}}{\textscy}{\textsubring{r}}} \dicPos{m}[6] \dicFlx{(‑s, ‑ar)}[8] \dicFieldCat{ekon.} \dicDirectTranslationCS{faktura, účet za zboží}
\dicEntry[vörurýrnun] \dicTerm{vöru··rýrn|un} \dicIPA{{v}{\oe}{\textlengthmark}{r}{\textscy}{r}{i}{r}{\textsubring{d}}{n}{\textscy}{\textsubring{n}}} \dicPos{f}[7] \dicFlx{(‑unar)}[9] \dicDirectTranslationCS{úbytek zboží}
\dicEntry[vörusending] \dicTerm{vöru··send·ing} \dicIPA{{v}{\oe}{\textlengthmark}{r}{\textscy}{s}{\textepsilon}{n}{\textsubring{d}}{i}{\ng}{\r{g}}} \dicPos{f}[4] \dicFlx{(‑ar, ‑ar)}[5] \dicDirectTranslationCS{zásilka zboží}
\dicEntry[vöruskemma] \textls[15]{\dicTerm{vöru··skemm|a} \dicIPA{{v}{\oe}{\textlengthmark}{r}{\textscy}{s}{\r{\textObardotlessj}}{\textepsilon}{m}{a}} \dicPos{f}[1] \dicFlx{(‑u, ‑ur)}[7] \dicSynonym{birgðageymsla} \dicDirectTranslationCS{sklad, skladiště}}
\dicEntry[vöruskemmdir] \dicTerm{vöru··skemmdir} \dicIPA{{v}{\oe}{\textlengthmark}{r}{\textscy}{s}{\r{\textObardotlessj}}{\textepsilon}{m}{\textsubring{d}}{\textsci}{\textsubring{r}}} \dicPos{f}[7] \dicFlx{pl}[2] \dicDirectTranslationCS{škody na zboží}
\dicEntry[vöruskiptahalli] \dicTerm{vöru·skipta··hall|i} \dicIPA{{v}{\oe}{\textlengthmark}{r}{\textscy}{s}{\r{\textObardotlessj}}{\textsci}{f}{\textsubring{d}}{a}{h}{a}{\textsubring{d}}{l}{\textsci}} \dicPos{m}[1] \dicFlx{(‑a)}[3] \dicFieldCat{ekon.} \dicDirectTranslationCS{deficit obchodní bilance}
\begin{xtolerant}{}{1pt}
\dicEntry[vöruskiptajöfnuður] \dicTerm{vöru·skipta··|jöfn·uður} \dicIPA{{v}\-{\oe}\-{\textlengthmark}\-{r}\-{\textscy}\-{s}\-{\r{\textObardotlessj}}\-{\textsci}\-{f}\-{\textsubring{d}}\-{a}\-{j}\-{\oe}\-{\textsubring{b}}\-{n}\-{\textscy}\-{ð}\-{\textscy}\-{\textsubring{r}}}\addthinS\dicPos{m}[10]\addthinS\dicFlx{(‑jafn\-aðar\,/\addthin ‑jöfnuðar)}[43] \dicFieldCat{ekon.} \dicDirectTranslationCS{obchodní bilance}
\end{xtolerant}
\dicEntry[vöruskipti] \dicTerm{vöru··skipti} \dicIPA{{v}{\oe}{\textlengthmark}{r}{\textscy}{s}{\r{\textObardotlessj}}{\textsci}{f}{\textsubring{d}}{\textsci}} \dicPos{n}[2] \dicFlx{pl}[19] \dicFieldCat{ekon.} \dicDirectTranslationCS{výměna zboží, barter}
\dicEntry[vöruskrá] \dicTerm{vöru··skrá} \dicIPA{{v}{\oe}{\textlengthmark}{r}{\textscy}{s}{\r{g}}{r}{au}} \dicPos{f}[4] \dicFlx{(‑r\,/\addthin ‑ar, ‑r)}[21] \dicDirectTranslationCS{seznam\,/\addthin inventář zboží}
\dicEntry[vörusýning] \dicTerm{vöru··sýn·ing} \dicIPA{{v}{\oe}{\textlengthmark}{r}{\textscy}{s}{i}{n}{i}{\ng}{\r{g}}} \dicPos{f}[4] \dicFlx{(‑ar, ‑ar)}[5] \dicDirectTranslationCS{veletrh}
\dicEntry[vörutalning] \dicTerm{vöru··tal·ning} \dicIPA{{v}{\oe}{\textlengthmark}{r}{\textscy}{t\smash{\textsuperscript{h}}}{a}{l}{n}{i}{\ng}{\r{g}}} \dicPos{f}[4] \dicFlx{(‑ar, ‑ar)}[5] \dicDirectTranslationCS{inventář zboží} \dicExampleIS{framkvæma vörutalning} \dicExampleCS{inventarizovat zboží}
\dicEntry[vörutegund] \dicTerm{vöru··tegund} \dicIPA{{v}{\oe}{\textlengthmark}{r}{\textscy}{t\smash{\textsuperscript{h}}}{\textepsilon}{\textbabygamma}{\textscy}{n}{\textsubring{d}}} \dicPos{f}[7] \dicFlx{(‑ar, ‑ir)}[1] \dicDirectTranslationCS{druh zboží, artikl, komodita}
\dicEntry[vörutollur] \dicTerm{vöru··toll|ur} \dicIPA{{v}{\oe}{\textlengthmark}{r}{\textscy}{t\smash{\textsuperscript{h}}}{\textopeno}{\textsubring{d}}{l}{\textscy}{\textsubring{r}}} \dicPos{m}[6] \dicFlx{(‑s, ‑ar)}[4] \dicFieldCat{ekon.} \dicDirectTranslationCS{dovozní clo}
\dicEntry[vöruúrval] \dicTerm{vöru··úr·val} \dicIPA{{v}{\oe}{\textlengthmark}{r}{\textscy}{u}{r}{v}{a}{\textsubring{l}}} \dicPos{n}[2] \dicFlx{(‑s)}[2] \dicDirectTranslationCS{výběr zboží, sortiment}
\dicEntry[vöruverð] \dicTerm{vöru··verð} \dicIPA{{v}{\oe}{\textlengthmark}{r}{\textscy}{v}{\textepsilon}{r}{\texttheta}} \dicPos{n}[2] \dicFlx{(‑s)}[2] \dicDirectTranslationCS{cena zboží}
\dicEntry[vösk] \dicTerm{vösk} \dicIPA{{v}{\oe}{s}{\r{g}}} \dicPos{adj} \dicFlx{f sg nom pos} \dicLink{vaskur}
\dicEntry[vötn] \dicTerm{vötn} \dicIPA{{v}{\oe}{h}{\textsubring{d}}{\textsubring{n}}} \dicPos{n} \dicFlx{pl nom} \dicLink{vatn}
\dicEntry[Vötnin miklu] \dicTerm{Vötnin miklu} \dicIPA{{v}{\oe}{h}{\textsubring{d}}{n}{\textsci}{n}{m}{\textsci}{h}{\r{g}}{l}{\textscy}} \dicPos{n}[2] \dicFlx{pl}[11] \dicFieldCat{geog.} \dicDirectTranslationCS{Velká jezera} \dicIndirectTranslationCS{(skupina jezer v~Severní Americe)}
\dicEntry[vött] \dicTerm{vött} \dicIPA{{v}{\oe}{h}{\textsubring{d}}} \dicPos{n} \dicFlx{pl nom} \dicLink{vatt\smash{\textsuperscript{1}}}
\dicEntry[vöxtur] \dicTerm{vöxtur} \dicsymFrequent\  \dicIPA{{v}{\oe}{x}{s}{\textsubring{d}}{\textscy}{\textsubring{r}}} \dicPos{m}[11] \dicFlx{(vaxtar, vextir)}[5] \textbf{1.} \dicSynonym*{það að vaxa} \dicDirectTranslationCS{(ná)růst}  \textbf{2.} \dicSynonym{stærð} \dicDirectTranslationCS{vzrůst, výška};  \dicIdiom{vöxtur}{ \dicPhraseIS{e‑að er í vexti}} \dicDirectTranslationCS{(co) stoupá (hladina řeky ap.)};  \dicPhraseIS{e‑að er mikið að vöxtum} \dicDirectTranslationCS{(co) je velké\,/\addthin rozsáhlé (sbírka ap.)}; { \dicPhraseIS{e‑að færist í vöxt}} \dicFlx{refl} \dicDirectTranslationCS{(co) vzrůstá, (co) je stále běžnější (zájem o~sport ap.)};  \dicPhraseIS{vera hár\,/\addthin lágur vexti} \dicDirectTranslationCS{být vysoký\,/\addthin nízký vzrůstem}

\else
\makeatletter
\@for\next:=\alphabetF\do{\input{letters/\next.tex}}
\makeatother 
\fi
\restoregeometry
\pagestyle{plain}

\fi

% ================================================================ PHONETICS =

\ifinputphon

\cleardoublepage

\chapter{Seznam islandských fonémů}              \label{sec:phon_phonems}

\begin{table}[h]
\begin{tabular}{lllll} \toprule

\textbf{foném} & \textbf{vysvětlivka} & \textbf{příklad} & \textbf{fonetický přepis} & \textbf{překlad} \\
\midrule
{\textipa{[{a}]}} & otevřená přední nezaokrouhlená samohl. & rata & {\textipa{[{r}{a}{\textlengthmark}{\textsubring{d}}{a}]}} & trefit \\ 
{\textipa{[{\textepsilon}]}} & polootevřená přední nezaokrouhlená samohl. & meta & {\textipa{[{m}{\textepsilon}{\textlengthmark}{\textsubring{d}}{a}]}} & ocenit \\ 
{\textipa{[{e}]}} & polozavřená přední nezaokrouhlená samohl. & heiður & {\textipa{[{h}{ei}{\textlengthmark}{ð}{\textscy}{\textsubring{r}}]}} & čest \\
{\textipa{[{\textsci}]}} & téměř zavřená téměř přední nezaokrouhlená samohl. & vinna & {\textipa{[{v}{\textsci}{n}{\textlengthmark}{a}]}} & pracovat \\ 
{\textipa{[{i}]}} & zavřená přední nezaokrouhlená samohl. & fínn & {\textipa{[{f}{i}{\textsubring{d}}{\textsubring{n}}]}} & jemný \\ 
{\textipa{[{\textopeno}]}} & polootevřená zadní zaokrouhlená samohl. & lofa & {\textipa{[{l}{\textopeno}{\textlengthmark}{v}{a}]}} & slibovat \\ 
{\textipa{[{o}]}} & polozavřená zadní zaokrouhlená samohl. & bóndi & {\textipa{[{\textsubring{b}}{ou}{n}{\textsubring{d}}{\textsci}]}} & farmář \\ 
{\textipa{[{\textscy}]}} & téměř zavřená téměř přední zaokrouhlená samohl. & hundur & {\textipa{[{h}{\textscy}{n}{\textsubring{d}}{\textscy}{\textsubring{r}}]}} & pes \\ 
{\textipa{[{u}]}} & zavřená zadní zaokrouhlená samohl. & súkkulaði & {\textipa{[{s}{u}{h}{\r{g}}{\textscy}{l}{a}{ð}{\textsci}]}} & čokoláda \\ 
{\textipa{[{\oe}]}} & téměř otevřená přední nezaokrouhlená samohl. & lönd & {\textipa{[{l}{\oe}{n}{\textsubring{d}}]}} & země \\ 
{\textipa{[{m}]}} & bilabiální nazála & mýri & {\textipa{[{m}{i}{\textlengthmark}{r}{\textsci}]}} & mokřina \\ 
{\textipa{[{\textsubring{m}}]}} & neznělá bilabiální nazála & þyrmt & {\textipa{[{\texttheta}{\textsci}{\textsubring{r}}{\textsubring{m}}{\textsubring{d}}]}} & {\textit{$\shortrightarrow$ þyrma}} \\  
{\textipa{[{n}]}} & alveolární nazála & nemandi & {\textipa{[{n}{\textepsilon}{\textlengthmark}{m}{a}{n}{\textsubring{d}}{\textsci}]}} & žák \\ 
{\textipa{[{\textsubring{n}}]}} & neznělá alveolární nazála & mennt & {\textipa{[{m}{\textepsilon}{\textsubring{n}}{\textsubring{d}}]}} & vzdělání \\ 
{\textipa{[{\textltailn}]}} & palatální nazála & angi & {\textipa{[{au}{\textltailn}{\r{\textObardotlessj}}{\textsci}]}} & drobek \\ 
{\textipa{[{\r{\textltailn}}]}} & neznělá palatální nazála & banki & {\textipa{[{\textsubring{b}}{au}{\r{\textltailn}}{\r{\textObardotlessj}}{\textsci}]}} & banka \\
{\textipa{[{\ng}]}} & velární nazála & skyggndur & {\textipa{[{s}{\r{\textObardotlessj}}{\textsci}{\ng}{\textsubring{d}}{\textscy}{\textsubring{r}}]}} & {\textit{$\shortrightarrow$ skyggna}} \\ 
{\textipa{[{\r{\ng}}]}} & neznělá velární nazála & skyggnt & {\textipa{[{s}{\r{\textObardotlessj}}{\textsci}{\r{\ng}}{\textsubring{d}}]}} & {\textit{$\shortrightarrow$ skyggna}} \\ 
{\textipa{[{\textsubring{b}}]}} & neznělá bilabiální ploziva & bara & {\textipa{[{\textsubring{b}}{a}{\textlengthmark}{r}{a}]}} & jenom \\ 
{\textipa{[{p\textsuperscript{h}}]}} & neznělá bilabiální ploziva s přídechem & prestur & {\textipa{[{p\textsuperscript{h}}{r}{\textepsilon}{s}{\textsubring{d}}{\textscy}{\textsubring{r}}]}} & kněz \\ 
{\textipa{[{\textsubring{d}}]}} & neznělá alveolární ploziva & dæmi & {\textipa{[{\textsubring{d}}{a}{i}{\textlengthmark}{m}{\textsci}]}} & příklad \\ 
{\textipa{[{t\textsuperscript{h}}]}} & neznělá alveolární ploziva s přídechem & tómur & {\textipa{[{t\textsuperscript{h}}{ou}{\textlengthmark}{m}{\textscy}{\textsubring{r}}]}} & prázdný \\ 
{\textipa{[{\r{\textObardotlessj}}]}} & neznělá palatální ploziva & gefa & {\textipa{[{\r{\textObardotlessj}}{\textepsilon}{\textlengthmark}{v}{a}]}} & dát \\ 
{\textipa{[{c\textsuperscript{h}}]}} & neznělá palatální ploziva s přídechem & kenna & {\textipa{[{c\textsuperscript{h}}{\textepsilon}{n}{\textlengthmark}{a}]}} & učit \\
{\textipa{[{\r{g}}]}} & neznělá velární ploziva & gata & {\textipa{[{\r{g}}{a}{\textlengthmark}{\textsubring{d}}{a}]}} & ulice \\ 
{\textipa{[{k\textsuperscript{h}}]}} & neznělá velární ploziva s přídechem & koma & {\textipa{[{k\textsuperscript{h}}{\textopeno}{\textlengthmark}{m}{a}]}} & přijít \\
{\textipa{[{v}]}} & znělá labiodentální frikativa & hafa & {\textipa{[{h}{a}{\textlengthmark}{v}{a}]}} & mít \\ 
{\textipa{[{f}]}} & neznělá labiodentální frikativa & finna & {\textipa{[{f}{\textsci}{n}{\textlengthmark}{a}]}} & najít \\ 
{\textipa{[{ð}]}} & znělá dentální frikativa & borða & {\textipa{[{\textsubring{b}}{\textopeno}{r}{ð}{a}]}} & jíst \\ 
{\textipa{[{\texttheta}]}} & neznělá dentální frikativa & þola & {\textipa{[{\texttheta}{\textopeno}{\textlengthmark}{l}{a}]}} & snést \\
{\textipa{[{s}]}} & neznělá alveolární frikativa & systir & {\textipa{[{s}{\textsci}{s}{\textsubring{d}}{\textsci}{\textsubring{r}}]}} & sestra \\ 
{\textipa{[{\c{c}}]}} & neznělá palatální frikativa & hjóla & {\textipa{[{\c{c}}{ou}{\textlengthmark}{l}{a}]}} & jet na kole \\ 
{\textipa{[{\textbabygamma}]}} & znělá velární frikativa & sagði & {\textipa{[{s}{a}{\textbabygamma}{ð}{\textsci}]}} & řekl \\ 
{\textipa{[{x}]}} & neznělá velární frikativa & lag & {\textipa{[{l}{a}{\textlengthmark}{x}]}} & píseň \\
{\textipa{[{h}]}} & neznělá glotální frikativa & heima & {\textipa{[{h}{ei}{\textlengthmark}{m}{a}]}} & doma \\ 
{\textipa{[{l}]}} & alveolární laterální aproximanta & loforð & {\textipa{[{l}{\textopeno}{\textlengthmark}{v}{\textopeno}{r}{\texttheta}]}} & slib \\ 
{\textipa{[{\textsubring{l}}]}} & neznělá alveolární laterální aproximanta & hlutur & {\textipa{[{\textsubring{l}}{\textscy}{\textlengthmark}{\textsubring{d}}{\textscy}{\textsubring{r}}]}} & věc \\  
{\textipa{[{j}]}} & palatální aproximanta & jól & {\textipa{[{j}{ou}{\textlengthmark}{\textsubring{l}}]}} & Vánoce \\ 
{\textipa{[{r}]}} & alveolární vibranta & reynsla & {\textipa{[{r}{ei}{n}{s}{\textsubring{d}}{l}{a}]}} & zkušenost \\ 
{\textipa{[{\textsubring{r}}]}} & neznělá alveolární vibranta & hraun & {\textipa{[{\textsubring{r}}{\oe i}{\textlengthmark}{\textsubring{n}}]}} & láva \\ 
\bottomrule
\end{tabular}
\end{table}

\clearpage

%\section{Výslovnost islandských samohlásek}      \label{sec:phon_vowels}
%\LTXtable{\columnwidth}{phon/vowels}
%\clearpage

%\section{Výslovnost islandských souhlásek}       \label{sec:phon_consonants}
%\LTXtable{\columnwidth}{phon/consonants}
%\clearpage
\fi

% =============================================================== MORPHOLOGY =

\ifinputmorpho

\cleardoublepage

\chapter{Morfologický klíč}                      \label{sec:morpho}

\section{Podstatné jméno, rod mužský}            \label{sec:morpho_m}
{\small\LTXfw{morpho/m}}
\clearpage

\section{Podstatné jméno, rod ženský}            \label{sec:morpho_f}
{\small\LTXfw{morpho/f}}
\clearpage

\section{Podstatné jméno, rod střední}           \label{sec:morpho_n}
{\small\LTXfw{morpho/n}}
\clearpage

\section{Přídavné jméno}                         \label{sec:morpho_adj}
{\small\LTXfw{morpho/adj}}
\clearpage

\section{Zájmeno}                                \label{sec:morpho_pron}
{\small\LTXfw{morpho/pron}}
 \begin{center}
\begin{minipage}[t]{.45\textwidth}
\begin{tabular}{l>{\footnotesize\itshape}lll}
 {\textbf{\textit{třída}}} & {\textit{pád}}   & \textit{sg} & \textit{pl}  \\ 
 \hline
\multirow{3}{*}{{{\textbf{pron} \Large{\textbf{17}}}}}  &  nom & ég & við   \\*
 & acc &  mig  & okkur  \\*
 & dat & mér & okkur   \\*
 & gen & mín  & okkar  \\
\hline

\multirow{3}{*}{{{\textbf{pron} \Large{\textbf{17}}}}}  &  nom & hann & þeir   \\*
 & acc &  hann  & þá  \\*
 & dat & honum & þeim   \\*
 & gen & hans  & þeirra  \\
\hline

\multirow{3}{*}{{{\textbf{pron} \Large{\textbf{17}}}}}  &  nom & hún & þær   \\*
 & acc &  hana  & þær  \\*
 & dat & henni & þeim   \\*
 & gen & hennar  & þeirra  \\
\hline

\multirow{3}{*}{{{\textbf{pron} \Large{\textbf{17}}}}}  &  nom &  & vér   \\*
 & acc &    & oss  \\*
 & dat &  & oss   \\*
 & gen &   & vor  \\
\hline

\multirow{3}{*}{{{\textbf{pron} \Large{\textbf{17}}}}}  &  nom &  & þér   \\*
 & acc &    & yður  \\*
 & dat &  & yður   \\*
 & gen &   & yðar  \\
\hline

\multirow{3}{*}{{{\textbf{pron} \Large{\textbf{17}}}}}  &  nom & þú & þið   \\*
 & acc &  þig  & ykkur  \\*
 & dat & þér & ykkur   \\*
 & gen & þín  & ykkar  \\
\hline

\multirow{3}{*}{{{\textbf{pron} \Large{\textbf{17}}}}}  &  nom & það & þau   \\*
 & acc &  það  & þau  \\*
 & dat & því & þeim   \\*
 & gen & þess  & þeirra  \\
\hline
\end{tabular}
\end{minipage}\hfil
\begin{minipage}[t]{.45\textwidth}
\begin{tabular}{l>{\footnotesize\itshape}lll}
\hline
\multirow{3}{*}{{{\textbf{pron} \Large{\textbf{18}}}}}  &  nom &  &    \\*
 & acc &  sig  &   \\*
 & dat & sér &    \\*
 & gen & sín  &   \\
 \hline
\end{tabular}
\end{minipage}
\end{center}

\clearpage

\section{Číslovka}                               \label{sec:morpho_num}
 {\small\LTXfw{morpho/num}}
\clearpage

\begin{landscape}
\section{Sloveso}                                \label{sec:morpho_v}
{\small\LTXfw{morpho/v}}
\end{landscape}

\clearpage

\fi

% ============================================================== BACK MATTER =

\ifmakebackmatter

% ------------------------------------------------ Photographs authors index -

\cleardoublepage
%\pagestyle{empty}
%\pagestyle{indexstyle}
%\phantomsection
%\addcontentsline{toc}{chapter}
%  {Seznam autorů fotografií a ilustrací}        
\label{sec:photo}
\printindex[figures]

% ------------------------------------------------------------- Bibliography -
\cleardoublepage
\chapter{Bibliografie}
%\twocolumn
%\phantomsection
%\addcontentsline{toc}{chapter}{Bibliografie}
\begin{multicols}{2}
\nocite{*}\printbibliography[heading=none]
\end{multicols}

% ------------------------------------------------------------------ Apendix -
\cleardoublepage
\onecolumn
\ifPDF
\else
\chapter{Apendix: Islandská gramatika}
\clearpage
\covergeometry  
\makeatletter
\includepdf[
  pages={1,2,3,4,5,6,7,8,9,10,11,12,13,14,15,16,17,18,19,20,21,22,23,24,25,26,27,28,29,30,31,32,33,34,35,36,37,38,39,40,41,42,43,44,45,46,47,48,49,50,51,52,53,54,55,56,57,58,59,60,61,62,63,64,65,66,67,68,69,70,71,72,73},
  turn=false,
  scale=0.79,
  keepaspectratio,
  offset={\dimexpr(\paperwidth-\csname Gm@layoutwidth\endcsname)/2-\csname Gm@layouthoffset\endcsname\relax}
  {\dimexpr(\paperheight-\csname  Gm@layoutheight\endcsname)/2-\csname Gm@layoutvoffset\endcsname\relax}
   ]{gram.pdf}
\makeatother
\restoregeometry  
\fi

% ------------------------------------------------------------------ Barcode -
\pagenumbering{gobble}
\ifPDF
\else
\cleardoublepage\null\clearpage % clear to verso
\vspace*\fill
\EANisbn[SC3]
\fi
\fi

% --------------------------------------------------------------- Back cover -

\ifmakecovers
  \cleardoublepage\null\clearpage % inner side of the cover
\covergeometry  
  \makecoverwith\backcoverimages
\restoregeometry  
\fi

\end{document}

